% Generated mechanically from the frozen etale-cohomology.tex target.
% Do not edit this generated file; edit the bound locale source instead.

% OK, start here.
%

\setcounter{chapter}{58}
\chapter{\'Etale 上同调}



\phantomsection
\label{etale-cohomology-section-phantom}



\section{引言}
\label{etale-cohomology-section-introduction}

\noindent
本章是关于概形 \'etale 上同调系列章节的第一章。本章讨论 \'etale
拓扑以及该拓扑中阿贝尔层上同调的最基本内容。所讨论的许多主题在首次阅读时
可以放心略过；关于如何决定略过哪些内容，请参阅下一节的建议。

\medskip\noindent
本章的初版源自 Johan de Jong 于 2009 年秋季在哥伦比亚大学讲授的
一门 \'etale 上同调课程第一部分的讲义。最初记录讲义的是
Thibaut Pugin、Zachary Maddock 和 Min Lee。课程第二部分见迹公式一章，参见
The Trace Formula，第 \ref{trace-section-introduction} 节。




\section{首次阅读时可略过哪些小节？}
\label{etale-cohomology-section-skip}

\noindent
我们希望将本章材料用于发展与代数空间、Deligne--Mumford 叠、代数叠等
有关的理论。因此，我们在原本关于概形 \'etale 上同调的阐述中加入了一些
相当技术性的材料。读者可以通过 ``topos'' 一词出现的频繁程度、与集合论
有关的讨论，或处理概形态射非常一般性质的证明来辨认这些材料。其中一些
讨论在首次阅读时可以略过。

\medskip\noindent
特别地，我们建议读者略过下列小节：
\begin{enumerate}
\item 比较大拓扑斯与小拓扑斯，
第 \ref{etale-cohomology-section-compare} 节。
\item 恢复态射，
第 \ref{etale-cohomology-section-morphisms} 节。
\item 推与拉，
第 \ref{etale-cohomology-section-monomorphisms} 节。
\item 性质 (A)，
第 \ref{etale-cohomology-section-A} 节。
\item 性质 (B)，
第 \ref{etale-cohomology-section-B} 节。
\item 性质 (C)，
第 \ref{etale-cohomology-section-C} 节。
\item 小 \'etale 网站的拓扑不变性，
第 \ref{etale-cohomology-section-topological-invariance} 节。
\item 整的处处单射态射，
第 \ref{etale-cohomology-section-integral-universally-injective} 节。
\item 大网站与前推，
第 \ref{etale-cohomology-section-big} 节。
\item 大下撇函子的正合性，
第 \ref{etale-cohomology-section-exactness-lower-shriek} 节。
\end{enumerate}
除这些小节外，还有一些零散结果也可以略过；读者可根据上述关键词辨认它们。



%9.08.09
\section{序言}
\label{etale-cohomology-section-prologue}

\noindent
本讲义讨论另一种上同调理论。首先要指出，Zariski 拓扑并不完全令人满意。
它无法给出我们所希望结果的主要原因之一是：若 $X$ 是复簇且
$\mathcal{F}$ 是常值层，则
$$
H^i(X, \mathcal{F}) = 0, \quad \text{ 对所有 } i > 0.
$$
原因如下。在不可约概形（特别是簇）中，任意两个非空开子集都有交，因此
常值层的限制映射是满射的。我们称这样的层为 {\it flasque}（柔软层）。
此时所有高阶 {\v C}ech 上同调群都消失，所有高阶 Zariski 上同调群也都消失。
换言之，Zariski 拓扑中的开集 ``不够多''，无法检测这种高阶上同调。

\medskip\noindent
另一方面，若 $X$ 是光滑射影复簇，则
$$
H_{Betti}^{2 \dim X}(X (\mathbf{C}), \Lambda) = \Lambda \quad \text{ 对于 }
\Lambda = \mathbf{Z}, \ \mathbf{Z}/n\mathbf{Z},
$$
其中 $X(\mathbf{C})$ 表示 $X$ 的复点集。这是代数几何中值得复现的性质，
尤其是在正特征情形下。




\section{\'Etale 拓扑}
\label{etale-cohomology-section-etale-topology}

\noindent
要仅仅通过 ``增加'' 额外开集来细化 Zariski 拓扑是很困难的。定义拓扑的一种
有效方法是不只考虑开集，还考虑位于开集之上的某些概形。为了定义 \'etale
拓扑，我们考虑所有 \'etale 的态射 $\varphi : U \to X$。若 $X$ 是
$\mathbf{C}$ 上的光滑射影簇，这意味着
\begin{enumerate}
\item $U$ 是光滑簇的不交并，且
\item $\varphi$（从解析意义上说）局部为同构。
\end{enumerate}
这里 ``解析地'' 指的是 $\mathbf{C}$ 上通常的（超越的）拓扑。因此第二个
条件意味着 $\varphi$ 的导数处处满秩（特别地，$U$ 的每个连通分支都与
$X$ 具有相同维数）。

\medskip\noindent
例如，一个二重覆盖——粗略地说，即簇之间次数为 $2$ 的有限映射——
$$
\Spec(\mathbf{C}[t])
\longrightarrow
\Spec(\mathbf{C}[t]),
\quad t \longmapsto t^2
$$
若其纤维含有单个点，就不是 \'etale 态射。在该例中，这发生在 $t = 0$ 时。
要使 $\mathbf{C}$ 上的簇之间有限映射为 \'etale，所有纤维都应具有相同的点数。
从上述映射的源中去掉点 $t = 0$ 会使态射变为 \'etale。但也可以从源中去掉
其他点，所得态射仍然是 \'etale。为了考虑 \'etale 拓扑，我们必须考察所有
这类态射。与 Zariski 拓扑不同，它们不必只是 $X$ 的开子集，尽管它们的像
总是开子集。

\begin{definition}
\label{etale-cohomology-definition-etale-covering-initial}
若态射族 $\{ \varphi_i : U_i \to X\}_{i \in I}$ 中每个 $\varphi_i$ 都是
\'etale 态射，且它们的像覆盖 $X$，即
$X = \bigcup_{i \in I} \varphi_i(U_i)$.
\end{definition}

\noindent
这就 ``定义'' 了 \'etale 拓扑。换言之，现在可以说明层是什么。一个
{\it \'etale 层} $\mathcal{F}$（集合层，分别地，阿贝尔群层、向量空间层等）
在 $X$ 上是如下数据：
\begin{enumerate}
\item 对每个 \'etale 态射 $\varphi : U \to X$，给出一个集合
（分别地，阿贝尔群、向量空间等）$\mathcal{F}(U)$；
\item 对每一对 \'etale 概形 $U, \ U'$（它们位于 $X$ 上），以及每个态射
$U \to U'$（它位于 $X$ 上且自动是 \'etale 的），给出限制映射
$\rho^{U'}_U : \mathcal{F}(U') \to \mathcal{F}(U)$
\end{enumerate}
这些数据必须满足：对于恒等态射，有 $\rho^U_U = \text{id}$，其中恒等态射为
$U \to U$；有 $\rho^{U'}_U \circ \rho^{U''}_{U'} = \rho^{U''}_U$，当有态射
$U \to U' \to U''$（它们是位于 $X$ 上的 \'etale 概形之间的态射）时；此外还满足如下
{\it 层公理}：
\begin{itemize}
\item[(*)] 对每个 \'etale 覆盖 $\{ \varphi_i : U_i \to U\}_{i \in
I}$，图
$$
\xymatrix{
\mathcal{F} (U) \ar[r] &
\Pi_{i \in I} \mathcal{F} (U_i) \ar@<1ex>[r] \ar@<-1ex>[r] &
\Pi_{i, j \in I} \mathcal{F} (U_i \times_U U_j)
}
$$
在集合（分别地，阿贝尔群、向量空间等）范畴中是等化子图。
\end{itemize}

\begin{remark}
\label{etale-cohomology-remark-i-is-j}
在最后一个陈述中，不能遗漏 $i = j$ 的情形；这通常是一个高度非平凡的
条件（不同于 Zariski 拓扑）。事实上，许多重要覆盖只有一个元素。
\end{remark}

\noindent
由于恒等态射是 \'etale 态射，我们可以计算 \'etale 层的整体截面，而上同调
就只是相应的右导出函子。换言之，待进一步发展理论并使陈述精确化之后，
定义上同调将不再有障碍。




\section{\'Etale 拓扑的性质}
\label{etale-cohomology-section-feats}

\noindent
对自然数 $n \in \mathbf{N} = \{1, 2, 3, 4, \ldots\}$，有
$$
H_\etale^2 (\mathbf{P}^1_\mathbf{C}, \mathbf{Z}/n\mathbf{Z}) =
\mathbf{Z}/n\mathbf{Z}.
$$
更一般地，若 $X$ 是复簇，则其取有限域系数的 \'etale Betti 数与
$X(\mathbf{C})$ 的通常 Betti 数相同，即
$$
\dim_{\mathbf{F}_q} H_\etale^{2i} (X, \mathbf{F}_q) =
\dim_{\mathbf{F}_q} H_{Betti}^{2i} (X(\mathbf{C}), \mathbf{F}_q).
$$
这令人非常满意。然而，这些等式只对挠系数成立，并非一般成立。对整数系数，有
$$
H_\etale^2 (\mathbf{P}^1_\mathbf{C}, \mathbf{Z}) = 0.
$$
相反，$H_{Betti}^2(\mathbf{P}^1(\mathbf{C}), \mathbf{Z}) = \mathbf{Z}$，
因为拓扑空间 $\mathbf{P}^1(\mathbf{C})$ 同胚于 $2$-球面。
通过极限过程，可以从挠系数回到非挠系数；我们很快就会讨论这种方法。




\section{一个计算}
\label{etale-cohomology-section-computation}

\noindent
如何计算 $\mathbf{P}^1_\mathbf{C}$ 取系数
$\Lambda = \mathbf{Z}/n\mathbf{Z}$ 时的上同调？
我们使用 {\v C}ech 上同调。$\mathbf{P}^1_\mathbf{C}$ 的一个覆盖由两个标准开集
$U_0, U_1$ 给出；二者都同构于 $\mathbf{A}^1_\mathbf{C}$，而它们的交同构于
$\mathbf{A}^1_\mathbf{C} \setminus \{0\} = \mathbf{G}_{m, \mathbf{C}}$。
事实证明，Mayer--Vietoris 序列在 \'etale 上同调中成立，从而得到正合序列
$$
H_\etale^{i-1}(U_0\cap U_1, \Lambda) \to
H_\etale^i(\mathbf{P}^1_\mathbf{C}, \Lambda) \to
H_\etale^i(U_0, \Lambda) \oplus
H_\etale^i(U_1, \Lambda) \to H_\etale^i(U_0\cap U_1,
\Lambda).
$$
为了得到预期答案，我们需要证明第三项中的直和消失。事实上，与通常拓扑一样，有
$$
H_\etale^q (\mathbf{A}^1_\mathbf{C}, \Lambda) = 0
\quad \text{ 当 } q \geq 1,
$$
并且
$$
H_\etale^q (\mathbf{A}^1_\mathbf{C} \setminus \{0\}, \Lambda) = \left\{
\begin{matrix}
\Lambda & \text{ 当 }q = 1\text{ 时；} \\
0 & \text{ 当 }q \geq 2\text{ 时。}
\end{matrix}
\right.
$$
这些结果已经相当困难（是否存在初等证明？）。下面说明一旦掌握 \'etale
上同调的工具，我们将如何进行这个计算。

\medskip\noindent
{\bf 高阶上同调。}下面的一般事实处理了这一点：若 $X$ 是 $\mathbf{C}$ 上的
仿射曲线，则
$$
H_\etale^q (X, \mathbf{Z}/n\mathbf{Z}) = 0 \quad \text{ 当 } q \geq 2.
$$
这可通过考察曲线的泛点并进行一些 Galois 上同调来证明。因此只需考虑一阶上同调。

\medskip\noindent
{\bf 一阶上同调。}我们使用如下识别：
\begin{eqnarray*}
H_\etale^1 (X, \mathbf{Z}/n\mathbf{Z}) = \left\{
\begin{matrix}
\text{集合层 }\mathcal{F}\text{，它位于 \'etale 网站 }X_\etale
\text{ 上并带有} \\
\text{作用 }\mathbf{Z}/n\mathbf{Z} \times \mathcal{F} \to \mathcal{F}
\text{，使得 }\mathcal{F}\text{ 是一个 }\mathbf{Z}/n\mathbf{Z}\text{-挠子。}
\end{matrix}
\right\}
\Big/ \cong
\\
 = \left\{
\begin{matrix}
\text{态射 }Y \to X\text{，它们是有限 \'etale 的，并带有} \\
\text{自由的 }\mathbf{Z}/n\mathbf{Z}\text{ 作用，使得 }
X = Y/(\mathbf{Z}/n\mathbf{Z}).
\end{matrix}
\right\}
\Big/ \cong.
\end{eqnarray*}
第一个识别非常一般（对网站上的任意上同调理论都成立），与 \'etale 拓扑无关。
第二个识别是下降理论的结果。最后一个集合描述了一族我们可以实际操作的几何对象。

\medskip\noindent
曲线 $\mathbf{A}^1_\mathbf{C}$ 没有非平凡有限 \'etale 覆盖，因此
$H_\etale^1 (\mathbf{A}^1_\mathbf{C}, \mathbf{Z}/n\mathbf{Z}) = 0$.
这既可从拓扑上看出，也可使用下一段中的论证。

\medskip\noindent
下面描述有限 \'etale 覆盖
$\varphi : Y \to \mathbf{A}^1_\mathbf{C} \setminus \{0\}$。
只需考虑 $Y$ 连通的情形，以下即作此假设。我们将应用 Riemann--Hurwitz
公式来确定 $Y$ 的可能形式（当然这有些迂回；如果愿意，也可以直接略过下一节）。
设该态射是 $n$ 对 1，并考虑射影紧化
$$
\xymatrix{
{Y\ } \ar@{^{(}->}[r] \ar[d]^\varphi &
{\bar Y} \ar[d]^{\bar\varphi} \\
{\mathbf{A}^1_\mathbf{C} \setminus \{0\}} \ar@{^{(}->}[r] &
{\mathbf{P}^1_\mathbf{C}}
}
$$
尽管 $\varphi$ 是 \'etale 且不分歧，$\bar{\varphi}$ 可能在 0 和 $\infty$ 处分歧。
设 0 的原像是点 $y_1,
\ldots, y_r$，分歧指数为 $e_1, \ldots e_r$；
$\infty$ 的原像是点 $y_1', \ldots, y_s'$，分歧指数为 $d_1, \ldots d_s$。
特别地，$\sum e_i = n = \sum d_j$。应用 Riemann--Hurwitz 公式得到
$$
2 g_Y - 2 = -2n + \sum (e_i - 1) + \sum (d_j - 1)
$$
因此 $g_Y = 0$、$r = s = 1$ 且 $e_1 = d_1 = n$。
于是 $Y \cong {\mathbf{A}^1_\mathbf{C} \setminus \{0\}}$，容易看出
$\varphi(z) = \lambda z^n$，其中 $\lambda \in \mathbf{C}^*$。
重新参数化 $Y$ 后可假设 $\lambda = 1$。因此该覆盖就是取 $n$ 次根，
作用在 $\mathbf{A}^1_{\mathbf{C}} \setminus \{0\}$ 上的坐标。

\medskip\noindent
请记住，我们需要分类
${\mathbf{A}^1_\mathbf{C} \setminus \{0\}}$ 的覆盖，并同时考虑其上的自由
$\mathbf{Z}/n\mathbf{Z}$-作用。在本例中，任意这样的作用都对应于 $Y$ 的一个
自同构，将 $z$ 送到 $\zeta_n z$，其中 $\zeta_n$ 是本原 $n$ 次单位根。
这样的作用有 $\phi(n)$ 个（这里 $\phi(n)$ 表示 Euler 函数）。因此，连通有限 \'etale
覆盖恰有 $\phi(n)$ 个；给定自由 $\mathbf{Z}/n\mathbf{Z}$-作用时，每个对应一个本原
$n$ 次单位根。留给读者验证：非连通有限 \'etale 次数为 $n$ 的覆盖
$\mathbf{A}^1_{\mathbf{C}} \setminus \{0\}$，带有给定自由
$\mathbf{Z}/n\mathbf{Z}$-作用，与非本原的 $n$ 次单位根（它们是 $1$ 的根）一一对应。
换言之，这个计算表明
$$
H_\etale^1 (\mathbf{A}^1_\mathbf{C} \setminus \{0\},
\mathbf{Z}/n\mathbf{Z}) =
\Hom(\mu_n(\mathbf{C}), \mathbf{Z}/n\mathbf{Z}) \cong \mathbf{Z}/n\mathbf{Z}.
$$
第一个识别是典范的，第二个不是；见备注 \ref{etale-cohomology-remark-normalize-H1-Gm}。
由于 Riemann--Hurwitz 公式的证明不使用上同调计算，上述论证实际上构成一个证明
（只要补充消失等事项的细节）。




\section{非挠系数}
\label{etale-cohomology-section-nontorsion}

\noindent
为研究非挠系数，作如下定义：
$$
H_\etale^i (X, \mathbf{Q}_\ell) :=
\left( \lim_n H_\etale^i(X, \mathbf{Z}/\ell^n\mathbf{Z}) \right)
\otimes_{\mathbf{Z}_\ell} \mathbf{Q}_\ell.
$$
符号 $\lim_n$ 表示上同调群系统
$H_\etale^i(X, \mathbf{Z}/\ell^n\mathbf{Z})$ 按 $n$ 编号的 {\it 极限}，见
Categories, Section \ref{categories-section-posets-limits}。
因此，我们需要研究满足某些相容条件的层系统。




\section{层理论}
\label{etale-cohomology-section-sheaf-theory}
%9.10.09

\noindent
从这里开始，我们正式讨论网站与层。接下来的几节可以容纳数量惊人的有用
抽象材料。其中一部分已经在前面的章节中展开，例如网站、网站上的模以及
网站上的上同调等章节。这里尽量不加入过多材料，只保留足以使本章后续内容
能够理解的部分。




\section{预层}
\label{etale-cohomology-section-presheaves}

\noindent
本节的参考资料是
Sites，第 \ref{sites-section-presheaves} 节。

\begin{definition}
\label{etale-cohomology-definition-presheaf}
设 $\mathcal{C}$ 是一个范畴。$\mathcal{C}$ 上的{\it 集合预层}（分别地，
{\it 阿贝尔预层}）是函子 $\mathcal{C}^{opp} \to
\textit{Sets}$（分别地，$\textit{Ab}$）。
\end{definition}

\noindent
{\bf 术语。}若 $U \in \Ob(\mathcal{C})$，则 $\mathcal{F}(U)$ 的元素称为
$\mathcal{F}$ 在 $U$ 上的{\it 截面}。对 $\varphi : V \to U$（其中它是
$\mathcal{C}$ 中的态射），映射 $\mathcal{F}(\varphi) : \mathcal{F}(U) \to \mathcal{F}(V)$
称为{\it 限制映射}，通常记作 $s \mapsto s|_V$，有时也记作
$s \mapsto \varphi^*s$。记号 $s|_V$ 有歧义，因为限制映射取决于 $\varphi$，
但这是约定俗成的记号滥用。我们还使用记号
$\Gamma(U, \mathcal{F}) = \mathcal{F}(U)$。

\medskip\noindent
说 $\mathcal{F}$ 是函子，意味着如果 $W \to V \to U$ 是
$\mathcal{C}$ 中的态射，且 $s \in \Gamma(U, \mathcal{F})$，那么
$(s|_V)|_W = s |_W$，这里仍使用上述记号滥用。此外，与恒等态射
$\text{id}_U : U \to U$ 对应的限制映射就是恒等映射。

\medskip\noindent
在 $\mathcal{C}$ 上的集合预层（分别地，阿贝尔预层）的范畴记作
$\textit{PSh} (\mathcal{C})$（分别地，$\textit{PAb}
(\mathcal{C})$）。它是从 $\mathcal{C}^{opp}$ 到
$\textit{Sets}$（分别地，$\textit{Ab}$）的函子范畴，也就是说，预层之间的
态射是函子的自然变换。我们只考虑
$\textit{PSh}(\mathcal{C})$ 和 $\textit{PAb}(\mathcal{C})$，并且此处范畴
$\mathcal{C}$ 为小范畴。
（我们的约定是：除非另有说明，范畴都是小的；若范畴不是小的，应列在
Categories，备注 \ref{categories-remark-big-categories} 中。）

\begin{example}
\label{etale-cohomology-example-representable-presheaf}
给定对象 $X \in \Ob(\mathcal{C})$，考虑如下规则：
$$
\begin{matrix}
U & \longmapsto & h_X(U) = \Mor_\mathcal{C}(U, X) \\
V \xrightarrow{\varphi} U & \longmapsto &
(\psi \mapsto \psi \circ \varphi) : h_X(U) \to h_X(V).
\end{matrix}
$$
这定义了函子 $h_X : \mathcal{C}^{opp} \to \textit{Sets}$，因而定义了一个预层。
它称为与 $X$ 关联的{\it 可表预层}。并非在每个网站的每个拓扑中，
可表预层都是层。
\end{example}

\begin{lemma}[Yoneda]
\label{etale-cohomology-lemma-yoneda}
\begin{slogan}
对象之间的态射与它们所表示的函子之间的自然变换构成一一对应。
\end{slogan}
设 $\mathcal{C}$ 是一个范畴，且 $X, Y \in
\Ob(\mathcal{C})$。存在自然双射
$$
\begin{matrix}
\Mor_\mathcal{C}(X, Y) &
\longrightarrow &
\Mor_{\textit{PSh}(\mathcal{C})} (h_X, h_Y) \\
\psi &
\longmapsto &
h_\psi = \psi \circ - : h_X \to h_Y.
\end{matrix}
$$
\end{lemma}

\begin{proof}
见
Categories，引理 \ref{categories-lemma-yoneda}。
\end{proof}




\section{网站}
\label{etale-cohomology-section-sites}


\begin{definition}
\label{etale-cohomology-definition-family-morphisms-fixed-target}
设 $\mathcal{C}$ 是一个范畴。一个以 U 为固定目标的{\it 态射族}
$\mathcal{U} = \{\varphi_i : U_i \to U\}_{i\in I}$ 由以下数据组成：
\begin{enumerate}
\item 对象 $U \in \mathcal{C}$；
\item 集合 $I$（可以为空）；以及
\item 对每个 $i\in I$，给出态射 $\varphi_i : U_i \to U$，它属于
$\mathcal{C}$ 且以 $U$ 为目标。
\end{enumerate}
\end{definition}

\noindent
我们有以固定目标的态射族之间的{\it 态射}这一概念；其特殊情形就是
{\it 加细}。这部分材料的参考资料是
Sites，第 \ref{sites-section-refinements} 节。

\begin{definition}
\label{etale-cohomology-definition-site}
一个{\it 网站}\footnote{我们称为 site 的对象，在
\cite[Expos\'e II, D\'efinition 1.3]{SGA4} 中称为带有预拓扑的范畴；
在 \cite{ArtinTopologies} 中称为带有 Grothendieck 拓扑的范畴。}
由范畴 $\mathcal{C}$ 及集合 $\text{Cov}(\mathcal{C})$ 组成；后者由以固定目标的
态射族构成，称为{\it 覆盖}，并满足：
\begin{enumerate}
\item （同构）若 $\varphi : V \to U$ 是 $\mathcal{C}$ 中的同构，
则 $\{\varphi : V \to U\}$ 是覆盖；
\item （局部性）若 $\{\varphi_i : U_i \to U\}_{i\in I}$ 是覆盖，并且
对每个 $i \in I$ 给定覆盖
$\{\psi_{ij} : U_{ij} \to U_i \}_{j\in I_i}$，则
$$
\{
\varphi_i \circ \psi_{ij} : U_{ij} \to U
\}_{(i, j)\in \prod_{i\in I} \{i\} \times I_i}
$$
也是覆盖；以及
\item （基变换）若 $\{U_i \to U\}_{i\in I}$
是覆盖，且 $V \to U$ 是 $\mathcal{C}$ 中的态射，则
\begin{enumerate}
\item 对每个 $i \in I$，纤维积 $U_i \times_U V$ 在 $\mathcal{C}$ 中存在；并且
\item $\{U_i \times_U V \to V\}_{i\in I}$ 是覆盖。
\end{enumerate}
\end{enumerate}
\end{definition}

\noindent
对我们而言，网站的底层范畴总是“``小''”的，即其对象的集合构成一个集合，
网站的覆盖集合也同样是一个集合（如上面的定义所述）。在本章中，我们通常省略
那些把对象族和覆盖族限制为集合的论证。进一步讨论见
Sites，备注 \ref{sites-remark-no-big-sites}。

\begin{example}
\label{etale-cohomology-example-site-topological-space}
若 $X$ 是拓扑空间，则有如下定义的关联网站 $X_{Zar}$：$X_{Zar}$ 的对象是
$X$ 的开子集，它们之间的态射是包含映射，覆盖是通常的拓扑（满射）覆盖。
注意，若 $U, V \subset W \subset X$ 是开子集，则 $U \times_W V = U \cap V$
存在；也就是说，这个范畴有纤维积。所有验证都是平凡的，结果正如预期。
\end{example}




\section{层}
\label{etale-cohomology-section-sheaves}

\begin{definition}
\label{etale-cohomology-definition-sheaf}
集合预层 $\mathcal{F}$（分别地，阿贝尔预层）在网站 $\mathcal{C}$ 上称为
{\it 分离预层}，如果对每个覆盖
$\{\varphi_i : U_i \to U\}_{i\in I} \in \text{Cov} (\mathcal{C})$，映射
$$
\mathcal{F}(U) \longrightarrow \prod\nolimits_{i\in I} \mathcal{F}(U_i)
$$
是单射。这里该映射为 $s \mapsto (s|_{U_i})_{i\in I}$。
预层 $\mathcal{F}$ 若对每个覆盖 $\{\varphi_i : U_i \to U\}_{i\in I} \in \text{Cov} (\mathcal{C})$，图
\begin{equation}
\label{etale-cohomology-equation-sheaf-axiom}
\xymatrix{
\mathcal{F}(U) \ar[r] &
\prod_{i\in I} \mathcal{F}(U_i) \ar@<1ex>[r] \ar@<-1ex>[r] &
\prod_{i, j \in I} \mathcal{F}(U_i \times_U U_j),
}
\end{equation}
其中第一个映射是 $s \mapsto (s|_{U_i})_{i\in I}$，右侧两个映射分别为
$(s_i)_{i\in I} \mapsto (s_i |_{U_i \times_U U_j})$ 和
$(s_i)_{i\in I} \mapsto (s_j |_{U_i \times_U U_j})$；该图在集合范畴
（分别地，阿贝尔群范畴）中是等化子图。
\end{definition}

\begin{remark}
\label{etale-cohomology-remark-empty-covering}
对于空覆盖（即 $I = \emptyset$），这意味着 $\mathcal{F}(\emptyset)$ 是空积，
它是相应范畴中的终对象（对 $\textit{Sets}$ 和 $\textit{Ab}$ 都是单元素集）。
\end{remark}

\begin{example}
\label{etale-cohomology-example-sheaf-site-space}
对拓扑空间关联的网站 $X_{Zar}$ 作上述展开（见例 \ref{etale-cohomology-example-site-topological-space}），
得到通常的层概念。
\end{example}

\begin{definition}
\label{etale-cohomology-definition-category-sheaves}
我们用 $\Sh(\mathcal{C})$（分别地，$\textit{Ab}(\mathcal{C})$）表示
$\textit{PSh}(\mathcal{C})$（分别地，$\textit{PAb}(\mathcal{C})$）的满子范畴，
其对象是层。这就是 $\mathcal{C}$ 上的{\it 集合层范畴}
（分别地，{\it 阿贝尔层范畴}）。
\end{definition}




\section{G-集的例子}
\label{etale-cohomology-section-G-sets}

\noindent
设 $G$ 是群，并如下定义网站 $\mathcal{T}_G$：其底层范畴是 $G$-集范畴，
即对象是带有左 $G$-作用的集合，态射是等变映射；$\mathcal{T}_G$ 的覆盖是态射族
$\{\varphi_i : U_i \to U\}_{i\in I}$，满足
$U = \bigcup_{i\in I} \varphi_i(U_i)$。

\medskip\noindent
网站 $\mathcal{T}_G$ 中有一个特殊对象，即带有左平移自然作用的 $G$-集 $G$。
记作 ${}_G G$。注意存在自然群同构
$$
\begin{matrix}
\rho : & G^{opp} & \longrightarrow & \text{Aut}_{G\textit{-Sets}}({}_G G) \\
& g & \longmapsto & (h \mapsto hg).
\end{matrix}
$$
特别地，对任意预层 $\mathcal{F}$，集合 $\mathcal{F}({}_G G)$ 继承一个
$G$-作用（通过 $\rho$）。（注意，由于 $\mathcal{F}$ 是反变的，集合 $\mathcal{F}({}_G G)$
仍是左 $G$-集。）事实上，函子
$$
\begin{matrix}
\Sh(\mathcal{T}_G) & \longrightarrow & G\textit{-Sets} \\
\mathcal{F} & \longmapsto & \mathcal{F}({}_G G)
\end{matrix}
$$
是范畴等价。其拟逆函子是 $X \mapsto
h_X$。这里不提供完整证明（可见
Sites，第 \ref{sites-section-example-sheaf-G-sets} 节），而尝试解释其原因。
\begin{enumerate}
\item
若 $S$ 是 $G$-集，可以将其分解为轨道 $S = \coprod_{i\in I}
O_i$。
覆盖 $\{O_i \to S\}_{i\in I}$ 的层公理说明
$$
\xymatrix{
\mathcal{F}(S) \ar[r] &
\prod_{i\in I} \mathcal{F}(O_i) \ar@<1ex>[r] \ar@<-1ex>[r] &
\prod_{i, j \in I} \mathcal{F}(O_i \times_S O_j)
}
$$
是等化子。注意 $G\textit{-Sets}$ 中的纤维积由 $\textit{Sets}$ 中的纤维积诱导，
并利用 $\mathcal{F}(\emptyset)$ 是单元素 $G$-集这一事实，可得
$$
\prod_{i, j \in I} \mathcal{F}(O_i \times_S O_j) = \prod_{i \in I}
\mathcal{F}(O_i)
$$
且上面的两个映射实际上相同。因此层公理仅仅说明
$\mathcal{F}(S) = \prod_{i\in I} \mathcal{F}(O_i)$。
\item
若 $S$ 是 $G$-集 $S= G/H$，且 $\mathcal{F}$ 是 $\mathcal{T}_G$ 上的层，
则有
$$
\mathcal{F}(G/H) = \mathcal{F}({}_G G)^H
$$
特别地，$\mathcal{F}(\{*\}) = \mathcal{F}({}_G G)^G$。为此，使用覆盖
$\{ {}_G G \to G/H \}$ 的层公理（该覆盖覆盖 $S$）。我们有
\begin{eqnarray*}
{}_G G \times_{G/H} {}_G G & \cong & G \times H \\
(g_1, g_2) & \longmapsto & (g_1, g_1^{-1} g_2)
\end{eqnarray*}
是 ${}_G G$ 的不交并（作为 $G$-集）。因此层公理变为
$$
\xymatrix{
\mathcal{F} (G/H) \ar[r] &
\mathcal{F}({}_G G) \ar@<1ex>[r] \ar@<-1ex>[r] &
\prod_{h\in H} \mathcal{F}({}_G G)
}
$$
其中右侧两个映射为 $s \mapsto (s)_{h \in H}$ 和 $s \mapsto
(hs)_{h \in H}$。因此 $\mathcal{F}(G/H) = \mathcal{F}({}_G G)^H$，如所声称。
\end{enumerate}
这还不能完全证明所声称的范畴等价，但至少说明层 $\mathcal{F}$ 完全由其在
${}_G G$ 上的截面决定。细节（以及集合论方面的说明）见
Sites，第 \ref{sites-section-example-sheaf-G-sets} 节。




\section{层化}
\label{etale-cohomology-section-sheafification}

\begin{definition}
\label{etale-cohomology-definition-0-cech}
设 $\mathcal{F}$ 是网站 $\mathcal{C}$ 上的预层，且
$\mathcal{U} = \{U_i \to U\} \in \text{Cov} (\mathcal{C})$。
定义 $\mathcal{F}$ 关于 $\mathcal{U}$ 的{\it 零阶 {\v C}ech 上同调群}为
$$
\check H^0 (\mathcal{U}, \mathcal{F}) =
\left\{
(s_i)_{i\in I} \in \prod\nolimits_{i\in I }\mathcal{F}(U_i)
\text{ 使得 }
s_i|_{U_i \times_U U_j} = s_j |_{U_i \times_U U_j}
\right\}.
$$
\end{definition}

\noindent
存在典范映射
$\mathcal{F}(U) \to \check H^0 (\mathcal{U}, \mathcal{F})$，
$s \mapsto (s |_{U_i})_{i\in I}$。
从覆盖 $\mathcal{V} = \{V_j \to V\}_{j \in J}$ 到 $\mathcal{U}$ 的
{\it 覆盖态射}定义为三元组 $(\chi, \alpha, \chi_j)$，其中
$\chi : V \to U$ 是态射，$\alpha : J \to I$ 是集合映射，并且对每个
$j \in J$，态射 $\chi_j$ 使下图交换：
$$
\xymatrix{
V_j \ar[rr]_{\chi_j} \ar[d] & & U_{\alpha(j)} \ar[d] \\
V \ar[rr]^\chi & & U.
}
$$
给定数据 $\chi, \alpha, \{\chi_j\}_{j \in J}$，定义
\begin{eqnarray*}
\check H^0(\mathcal{U}, \mathcal{F}) & \longrightarrow &
\check H^0(\mathcal{V}, \mathcal{F}) \\
(s_i)_{i\in I} & \longmapsto &
\left(\chi_j^*\left(s_{\alpha(j)}\right)\right)_{j\in J}.
\end{eqnarray*}
我们断言：
\begin{enumerate}
\item 该映射定义良好；并且
\item 它只依赖于 $\chi$，与 $\alpha, \{\chi_j\}_{j \in J}$ 的选择无关。
\end{enumerate}
第一点的证明从略。为说明第 (2) 点，考虑另一个三元组
$(\psi, \beta, \psi_j)$，满足 $\chi = \psi$。于是有交换图
$$
\xymatrix{
V_j \ar[rrr]_{(\chi_j, \psi_j)} \ar[dd] & & &
U_{\alpha(j)} \times_U U_{\beta(j)} \ar[dl] \ar[dr] \\
& & U_{\alpha(j)} \ar[dr] & &
U_{\beta(j)} \ar[dl] \\
V \ar[rrr]^{\chi = \psi} & & & U.
}
$$
给定截面 $s \in \mathcal{F}(\mathcal{U})$，其在 $\mathcal{F}(V_j)$ 中、由
$(\chi, \alpha, \{\chi_j\}_{j \in J})$ 给出的映射下的像为
$\chi_j^*s_{\alpha(j)}$；在由
$(\psi, \beta, \{\psi_j\}_{j \in J})$ 给出的映射下，其像为
$\psi_j^*s_{\beta(j)}$。二者相等，因为按假设
$s \in \check H^0(\mathcal{U}, \mathcal{F})$，所以二者都等于共同值
$$
s_{\alpha(j)}|_{U_{\alpha(j)} \times_U U_{\beta(j)}} =
s_{\beta(j)}|_{U_{\alpha(j)} \times_U U_{\beta(j)}}
$$
沿图中的映射 $(\chi_j, \psi_j)$ 的拉回。

\begin{theorem}
\label{etale-cohomology-theorem-sheafification}
设 $\mathcal{C}$ 是网站，$\mathcal{F}$ 是 $\mathcal{C}$ 上的预层。
\begin{enumerate}
\item 规则
$$
U \mapsto \mathcal{F}^+(U) :=
\colim_{\mathcal{U} \text{ 为 }U\text{ 的覆盖}}
\check H^0(\mathcal{U}, \mathcal{F})
$$
是预层，并且该余极限是有向余极限。
\item 存在预层的典范映射 $\mathcal{F} \to \mathcal{F}^+$。
\item 若 $\mathcal{F}$ 是分离预层，则 $\mathcal{F}^+$ 是层，且 (2) 中的映射是单射。
\item $\mathcal{F}^+$ 是分离预层。
\item $\mathcal{F}^\# = (\mathcal{F}^+)^+$ 是层，并且典范映射诱导函子同构
$$
\Hom_{\textit{PSh}(\mathcal{C})}(\mathcal{F}, \mathcal{G}) =
\Hom_{\Sh(\mathcal{C})}(\mathcal{F}^\#, \mathcal{G})
$$
对任意 $\mathcal{G} \in \Sh(\mathcal{C})$ 都成立。
\end{enumerate}
\end{theorem}

\begin{proof}
见 Sites，定理 \ref{sites-theorem-plus}。
\end{proof}

\noindent
换言之，这意味着自然映射 $\mathcal{F} \to \mathcal{F}^\#$ 是遗忘函子
$\Sh(\mathcal{C}) \to \textit{PSh}(\mathcal{C})$ 的左伴随。




\section{上同调}
\label{etale-cohomology-section-cohomology}

\noindent
下面的基本结果使我们能够定义网站上阿贝尔层的上同调。

\begin{theorem}
\label{etale-cohomology-theorem-enough-injectives}
网站上的阿贝尔层范畴是一个具有足够内射对象的阿贝尔范畴。
\end{theorem}

\begin{proof}
见
Modules on Sites，引理 \ref{sites-modules-lemma-abelian-abelian}，以及
Injectives，定理 \ref{injectives-theorem-sheaves-injectives}。
\end{proof}

\noindent
因此，可以将上同调定义为截面函子的右导出函子：若 $U \in \Ob(\mathcal{C})$ 且
$\mathcal{F} \in \textit{Ab}(\mathcal{C})$，则
$$
H^p(U, \mathcal{F}) :=
R^p\Gamma(U, \mathcal{F}) =
H^p(\Gamma(U, \mathcal{I}^\bullet))
$$
其中 $\mathcal{F} \to \mathcal{I}^\bullet$ 是内射分解。为此需要检查函子
$\Gamma(U, -)$ 是左正合的。这一点成立，也是范畴 $\textit{Ab}(\mathcal{C})$
成为阿贝尔范畴的原因之一，见
Modules on Sites，引理 \ref{sites-modules-lemma-abelian-abelian}。
关于网站上上同调的更一般讨论（包括整体截面函子及其右导出函子），见
Cohomology on Sites，第 \ref{sites-cohomology-section-cohomology-sheaves} 节。



\section{fpqc 拓扑}
\label{etale-cohomology-section-fpqc}
%9.15.09

\noindent
在讨论 \'etale 上同调之前，我们先研究一些 fpqc 拓扑，因为它对拟凝聚层很有效。

\begin{definition}
\label{etale-cohomology-definition-fpqc-covering}
设 $T$ 是概形。$T$ 的{\it fpqc 覆盖}是满足以下条件的族
$\{ \varphi_i : T_i \to T\}_{i \in I}$：
\begin{enumerate}
\item 每个 $\varphi_i$ 都是平坦态射，且
$\bigcup_{i\in I} \varphi_i(T_i) = T$，并且
\item 对每个仿射开集 $U \subset T$，存在有限集合 $K$、映射
$\mathbf{i} : K \to I$ 以及仿射开集
$U_{\mathbf{i}(k)} \subset T_{\mathbf{i}(k)}$，使得
$U = \bigcup_{k \in K} \varphi_{\mathbf{i}(k)}(U_{\mathbf{i}(k)})$.
\end{enumerate}
\end{definition}

\begin{remark}
\label{etale-cohomology-remark-fpqc}
第一个条件对应 fp，它代表法语 {\it fid\`element plat}（忠实平坦）；
第二个条件对应 qc，即{\it quasi-compact}（拟紧）。当第二个条件成立时，
第一个条件的后半部分是不必要的。
\end{remark}

\begin{example}
\label{etale-cohomology-example-fpqc-coverings}
fpqc 覆盖的例子。
\begin{enumerate}
\item $T$ 的任意 Zariski 开覆盖都是 fpqc 覆盖。
\item 族 $\{\Spec(B) \to \Spec(A)\}$ 是 fpqc 覆盖，当且仅当 $A \to B$
是忠实平坦的环映射。
\item 若 $f: X \to Y$ 平坦、满射且拟紧，则 $\{ f: X\to
Y\}$ 是 fpqc 覆盖。
\item 态射
$\varphi :
\coprod_{x \in \mathbf{A}^1_k} \Spec(\mathcal{O}_{\mathbf{A}^1_k, x})
\to \mathbf{A}^1_k$,
其中 $k$ 是域；该态射平坦且满射，但不是拟紧的，事实上族 $\{\varphi\}$
不是 fpqc 覆盖。
\item 记
$\mathbf{A}^2_k = \Spec(k[x, y])$。记 $i_x : D(x) \to \mathbf{A}^2_k$
和 $i_y : D(y) \to \mathbf{A}^2_k$ 为标准开子集。
那么族
$\{i_x, i_y, \Spec(k[[x, y]]) \to \mathbf{A}^2_k\}$
以及
$\{i_x, i_y, \Spec(\mathcal{O}_{\mathbf{A}^2_k, 0}) \to \mathbf{A}^2_k\}$
都是 fpqc 覆盖。
\end{enumerate}
\end{example}

\begin{lemma}
\label{etale-cohomology-lemma-site-fpqc}
概形范畴上的 fpqc 覆盖族满足网站公理。
\end{lemma}

\begin{proof}
见 Topologies，引理 \ref{topologies-lemma-fpqc}。
\end{proof}

\noindent
看起来，这个引理允许我们定义概形范畴的 fpqc 网站。然而，考虑 fpqc 拓扑时会
出现集合论问题，见 Topologies，第 \ref{topologies-section-fpqc} 节。
问题源于我们要求网站是“``小''”的，但没有任何小的概形范畴能够包含给定非空概形的
fpqc 覆盖的余最终系统。严格说来，这并不妨碍定义“部分”的 fpqc 网站，
但这样做似乎并不明智。变通办法是允许使用 fpqc 拓扑的层概念（见下文），
但禁止考虑所有 fpqc 层所成的范畴。

\begin{definition}
\label{etale-cohomology-definition-sheaf-property-fpqc}
设 $S$ 是概形。$S$ 上概形的范畴记作
$\Sch/S$。考虑函子
$\mathcal{F} : (\Sch/S)^{opp} \to \textit{Sets}$，即集合预层。称 $\mathcal{F}$
{\it 满足 fpqc 拓扑的层性质}，如果对每个 fpqc 覆盖
$\{U_i \to U\}_{i \in I}$（其覆盖 $S$ 上的概形），图 (\ref{etale-cohomology-equation-sheaf-axiom}) 是等化子图。
\end{definition}

\noindent
类似地，称 $\mathcal{F}$ {\it 满足 Zariski 拓扑的层性质}，如果对每个开覆盖
$U = \bigcup_{i \in I} U_i$，图 (\ref{etale-cohomology-equation-sheaf-axiom}) 是等化子图。见
Schemes，定义 \ref{schemes-definition-representable-by-open-immersions}。
显然，这等价于：对每个概形 $T$（位于 $S$ 上），将 $\mathcal{F}$ 限制到 $T$ 的开集上
所得的是一个（通常意义下的）层。

\begin{lemma}
\label{etale-cohomology-lemma-fpqc-sheaves}
设 $\mathcal{F}$ 是 $\Sch/S$ 上的预层。则 $\mathcal{F}$ 满足 fpqc 拓扑的层性质
当且仅当
\begin{enumerate}
\item $\mathcal{F}$ 满足 Zariski 拓扑的层性质；并且
\item 对每个忠实平坦态射 $\Spec(B) \to \Spec(A)$（它是 $S$ 上仿射概形之间的态射），
覆盖 $\{\Spec(B) \to \Spec(A)\}$ 满足层公理。也就是说，
$$
\xymatrix{
\mathcal{F}(\Spec(A)) \ar[r] &
\mathcal{F}(\Spec(B)) \ar@<1ex>[r] \ar@<-1ex>[r] &
\mathcal{F}(\Spec(B \otimes_A B))
}
$$
是等化子图。
\end{enumerate}
\end{lemma}

\begin{proof}
见 Topologies，引理 \ref{topologies-lemma-sheaf-property-fpqc}。
\end{proof}

\noindent
对于满足 fpqc 拓扑层条件的预层 $\mathcal{F}$（它定义在 $\Sch/S$ 上），
另一种理解方式是以下数据：
\begin{enumerate}
\item 对每个 $T/S$，给出通常意义（即 Zariski）的层 $\mathcal{F}_T$（位于 $T_{Zar}$ 上）；
\item 对每个映射 $f : T' \to T$（位于 $S$ 上），给出限制映射
$f^{-1}\mathcal{F}_T \to \mathcal{F}_{T'}$；
\end{enumerate}
满足
\begin{enumerate}
\item[(a)] 限制映射满足函子性；
\item[(b)] 若 $f : T' \to T$ 是开浸入，则限制映射
$f^{-1}\mathcal{F}_T \to \mathcal{F}_{T'}$ 是同构；并且
\item[(c)] 对每个忠实平坦态射
$\Spec(B) \to \Spec(A)$（位于 $S$ 上），图
$$
\xymatrix{
\mathcal{F}_{\Spec(A)}(\Spec(A)) \ar[r] &
\mathcal{F}_{\Spec(B)}(\Spec(B)) \ar@<1ex>[r] \ar@<-1ex>[r] &
\mathcal{F}_{\Spec(B \otimes_A B)}(\Spec(B \otimes_A B))
}
$$
是等化子。
\end{enumerate}
数据 (1)、(2) 以及条件 (a)、(b) 给出 $\Sch/S$ 上满足 Zariski 拓扑层条件的预层。
由引理 \ref{etale-cohomology-lemma-fpqc-sheaves}，条件 (c) 即足以得到 fpqc 拓扑的层条件。

\begin{example}
\label{etale-cohomology-example-quasi-coherent}
考虑预层
$$
\begin{matrix}
\mathcal{F} : & (\Sch/S)^{opp} & \longrightarrow & \textit{Ab} \\
& T/S & \longmapsto & \Gamma(T, \Omega_{T/S}).
\end{matrix}
$$
微分与局部化的相容性说明，$\mathcal{F}$ 是 Zariski 网站上的层。
然而，它不满足 fpqc 拓扑的层条件。具体地，考虑
$S = \Spec(\mathbf{F}_p)$ 以及态射
$$
\varphi :
V = \Spec(\mathbf{F}_p[v])
\to
U = \Spec(\mathbf{F}_p[u])
$$
它把 $u$ 映到 $v^p$。族 $\{\varphi\}$ 是 fpqc 覆盖，
但限制映射
$\mathcal{F}(U) \to \mathcal{F}(V)$
把生成元 $\text{d}u$ 映到 $\text{d}(v^p) = 0$，所以它是零映射，而图
$$
\xymatrix{
\mathcal{F}(U) \ar[r]^{0} &
\mathcal{F}(V) \ar@<1ex>[r] \ar@<-1ex>[r] &
\mathcal{F}(V \times_U V)
}
$$
不是等化子。稍后将看到，$\mathcal{F}$ 的确在 \'etale 网站与光滑网站上
给出一个层。
\end{example}

\begin{lemma}
\label{etale-cohomology-lemma-representable-sheaf-fpqc}
$\Sch/S$ 上任意可表预层都满足 fpqc 拓扑的层条件。
\end{lemma}

\begin{proof}
见 Descent，引理 \ref{descent-lemma-fpqc-universal-effective-epimorphisms}。
\end{proof}

\noindent
稍后还会回到这一点，因为该事实的证明要用到拟凝聚层的下降，而这将在下一节
讨论。用一种较为精致的说法，该引理断言
{\it fpqc 拓扑弱于典范拓扑}，或者说 fpqc 拓扑是{\it 次典范的}。
在网站的框架下，这一点见 Sites，第 \ref{sites-section-representable-sheaves} 节。

\begin{remark}
\label{etale-cohomology-remark-fpqc-finest}
fpqc 拓扑比 Zariski、\'etale、光滑、syntomic 与 fppf 拓扑都更细。因此，
任意满足 fpqc 拓扑层条件的预层，在 Zariski、\'etale、光滑、syntomic 与
fppf 网站上都是层。特别地，例如可表预层在概形的 \'etale 网站上是层。
\end{remark}

\begin{example}
\label{etale-cohomology-example-additive-group-sheaf}
设 $S$ 是概形。考虑加法群概形
$\mathbf{G}_{a, S} = \mathbf{A}^1_S$（在 $S$ 上），见
Groupoids，例 \ref{groupoids-example-additive-group}。
关联的可表预层由下式给出：
$$
h_{\mathbf{G}_{a, S}}(T) =
\Mor_S(T, \mathbf{G}_{a, S}) =
\Gamma(T, \mathcal{O}_T).
$$
由上文可知，这是满足 fpqc 拓扑层条件的集合预层。另一方面，它显然也是环预层。
因此可以将其看作函子
$$
\begin{matrix}
\mathcal{O} : &
(\Sch/S)^{opp} &
\longrightarrow &
\textit{Rings} \\
&
T/S &
\longmapsto &
\Gamma(T, \mathcal{O}_T)
\end{matrix}
$$
且满足 fpqc 拓扑的层条件。相应地，还可以定义 $\mathcal{O}$-模，等等。
\end{example}




\section{忠实平坦下降}
\label{etale-cohomology-section-fpqc-descent}

\noindent
本节讨论拟凝聚模的忠实平坦下降。更准确地说，我们将证明：对于 fpqc 覆盖，
拟凝聚模满足有效下降。

\begin{definition}
\label{etale-cohomology-definition-descent-datum}
设 $\mathcal{U} = \{ t_i : T_i \to T\}_{i \in I}$ 是以固定目标为目标的概形态射族。
关于 $\mathcal{U}$ 的拟凝聚层{\it 下降数据}是集合
$((\mathcal{F}_i)_{i \in I}, (\varphi_{ij})_{i, j \in I})$，其中
\begin{enumerate}
\item $\mathcal{F}_i$ 是 $T_i$ 上的拟凝聚层；
\item $\varphi_{ij} : \text{pr}_0^* \mathcal{F}_i \to
\text{pr}_1^* \mathcal{F}_j$ 是 $T_i \times_T T_j$ 上的模同构；
\end{enumerate}
并满足{\it 循环条件}：下图
$$
\xymatrix{
\text{pr}_0^*\mathcal{F}_i \ar[dr]_{\text{pr}_{02}^*\varphi_{ik}}
\ar[rr]^{\text{pr}_{01}^*\varphi_{ij}} & &
\text{pr}_1^*\mathcal{F}_j \ar[dl]^{\text{pr}_{12}^*\varphi_{jk}} \\
& \text{pr}_2^*\mathcal{F}_k
}
$$
在 $T_i \times_T T_j \times_T T_k$ 上交换。
若存在拟凝聚层 $\mathcal{F}$ 位于 $T$ 上，以及
$\mathcal{O}_{T_i}$-模同构 $\varphi_i : t_i^* \mathcal{F} \cong \mathcal{F}_i$，
且它们与映射 $\varphi_{ij}$ 相容，即
$$
\varphi_{ij} = \text{pr}_1^* (\varphi_j) \circ \text{pr}_0^* (\varphi_i)^{-1}.
$$
\end{definition}

\noindent
本节及下一节讨论下述定理证明中的一些成分，以及相关材料。

\begin{theorem}
\label{etale-cohomology-theorem-descent-quasi-coherent}
若 $\mathcal{V} = \{T_i \to T\}_{i\in I}$ 是 fpqc 覆盖，则关于 $\mathcal{V}$ 的
所有拟凝聚层下降数据都是有效的。
\end{theorem}

\begin{proof}
见
Descent，命题 \ref{descent-proposition-fpqc-descent-quasi-coherent}。
\end{proof}

\noindent
换言之，拟凝聚层的纤维范畴是 fpqc 网站上的一个叠。
定理的证明分两步。第一步是注意到对 Zariski 覆盖，这一点利用层的标准粘合（见
Sheaves，第 \ref{sheaves-section-glueing-sheaves} 节）以及拟凝聚性的局部性，
很容易（或众所周知）。第二步是处理形如 $\{\Spec(B) \to \Spec(A)\}$ 的 fpqc 覆盖，
其中 $A \to B$ 是忠实平坦的环映射。这是一个代数引理，现予以陈述。

\medskip\noindent
{\bf 模的下降。}若 $A \to B$ 是环映射，考虑复形
$$
(B/A)_\bullet : B \to B \otimes_A B \to B \otimes_A B \otimes_A B \to \ldots
$$
其中 $B$ 位于次数 0，$B \otimes_A B$ 位于次数 1，依此类推；映射由下式给出：
\begin{eqnarray*}
b & \mapsto & 1 \otimes b - b \otimes 1, \\
b_0 \otimes b_1 & \mapsto & 1 \otimes b_0 \otimes b_1 - b_0 \otimes 1 \otimes
b_1 + b_0 \otimes b_1 \otimes 1, \\
& \text{etc.}
\end{eqnarray*}

\begin{lemma}
\label{etale-cohomology-lemma-algebra-descent}
若 $A \to B$ 忠实平坦，则复形 $(B/A)_\bullet$ 在正次数上正合，且
$H^0((B/A)_\bullet) = A$。
\end{lemma}

\begin{proof}
见 Descent，引理 \ref{descent-lemma-ff-exact}。
\end{proof}

\noindent
Grothendieck 分三步证明这一点。首先，他假设映射 $A
\to B$ 有截面，并构造到如下复形的显式同伦：只有次数 0 的 $A$ 项非零。其次，
他指出，只需在忠实平坦基变换 $A \to
A'$ 后证明结果，此时将 $B$ 替换为 $B' = B \otimes_A A'$。第三，他取忠实平坦基变换
$A \to A' = B$，并注意映射
$A' = B \to B' = B \otimes_A B$ 有自然截面。

\medskip\noindent
同样的策略证明以下引理。

\begin{lemma}
\label{etale-cohomology-lemma-descent-modules}
若 $A \to B$ 忠实平坦且 $M$ 是 $A$-模，则复形
$(B/A)_\bullet \otimes_A M$ 在正次数上正合，且
$H^0((B/A)_\bullet \otimes_A M) = M$。
\end{lemma}

\begin{proof}
见 Descent，引理 \ref{descent-lemma-ff-exact}。
\end{proof}

\begin{definition}
\label{etale-cohomology-definition-descent-datum-modules}
设 $A \to B$ 是环映射，$N$ 是 $B$-模。$N$ 关于 $A \to B$ 的
一个{\it 下降数据}，是 $B \otimes_A B$-模同构
$\varphi : N \otimes_A B \cong B \otimes_A N$，使得下列
$B \otimes_A B \otimes_A B$-模图
$$
\xymatrix{
{N \otimes_A B \otimes_A B} \ar[dr]_{\varphi_{02}} \ar[rr]^{\varphi_{01}} & &
{B \otimes_A N \otimes_A B} \ar[dl]^{\varphi_{12}} \\
& {B \otimes_A B \otimes_A N}
}
$$
交换，其中 $\varphi_{01} = \varphi \otimes \text{id}_B$，而
$\varphi_{12}$ 与 $\varphi_{02}$ 以类似方式定义。
\end{definition}

\noindent
若 $N' = B \otimes_A M$，其中 M 是某个 $A$-模，则该模具有由下列映射
给出的典范下降数据：
$$
\begin{matrix}
\varphi_\text{can}: & N' \otimes_A B & \to & B \otimes_A N' \\
& b_0 \otimes m \otimes b_1 & \mapsto & b_0 \otimes b_1 \otimes m.
\end{matrix}
$$

\begin{definition}
\label{etale-cohomology-definition-effective-modules}
若存在 $A$-模 $M$，使得
$(N, \varphi) \cong (B \otimes_A M, \varphi_\text{can})$，则称下降数据
$(N, \varphi)$ 是{\it 有效的}；这里采用显然的下降数据同构概念。
\end{definition}

\noindent
定理 \ref{etale-cohomology-theorem-descent-quasi-coherent} 是下述结果的推论。

\begin{theorem}
\label{etale-cohomology-theorem-descent-modules}
若 $A \to B$ 忠实平坦，则关于 $A\to B$ 的下降数据都是有效的。
\end{theorem}

\begin{proof}
见 Descent，命题 \ref{descent-proposition-descent-module}。
另见 Descent，评注
\ref{descent-remark-homotopy-equivalent-cosimplicial-algebras}，其中给出证明的另一种看法。
\end{proof}

\begin{remarks}
\label{etale-cohomology-remarks-theorem-modules-exactness}
模的下降结果有若干应用：
\begin{enumerate}
\item 对覆盖 $\{\Spec(B) \to \Spec(A)\}$，其中 $A \to B$ 忠实平坦，
{\v C}ech 复形在正次数上正合。这将给出若干上同调消失结果。
\item 若 $(N, \varphi)$ 是关于忠实平坦映射 $A \to B$ 的下降数据，则相应的
$A$-模由下式给出：
$$
M = \Ker \left(
\begin{matrix}
N & \longrightarrow & B \otimes_A N \\
n & \longmapsto & 1 \otimes n - \varphi(n \otimes 1)
\end{matrix}
\right).
$$
见 Descent，命题 \ref{descent-proposition-descent-module}。
\end{enumerate}
\end{remarks}




%9.17.09
\section{拟凝聚层}
\label{etale-cohomology-section-quasi-coherent}

\noindent
我们可以利用模的下降来研究拟凝聚层。

\begin{proposition}
\label{etale-cohomology-proposition-quasi-coherent-sheaf-fpqc}
对任意拟凝聚层 $\mathcal{F}$（定义在 $S$ 上），预层
$$
\begin{matrix}
\mathcal{F}^a : & \Sch/S & \to & \textit{Ab}\\
& (f: T \to S) & \mapsto & \Gamma(T, f^*\mathcal{F})
\end{matrix}
$$
是满足 fpqc 拓扑层条件的 $\mathcal{O}$-模。
\end{proposition}

\begin{proof}
这在 Descent，引理 \ref{descent-lemma-sheaf-condition-holds} 中证明。
这里给出证明梗概。由引理 \ref{etale-cohomology-lemma-fpqc-sheaves} 可知，
只需检查仿射概形的 Zariski 覆盖和忠实平坦态射上的层性质。
Zariski 覆盖的层性质是标准概形理论，因为
$\Gamma(U, i^\ast \mathcal{F}) = \mathcal{F}(U)$，其中
$i : U \hookrightarrow S$ 是开浸入。

\medskip\noindent
对于 $\left\{\Spec(B)\to \Spec(A)\right\}$，若 $A\to B$ 忠实
平坦且 $\mathcal{F}|_{\Spec(A)} = \widetilde{M}$，这对应于
$M = H^0\left((B/A)_\bullet \otimes_A M \right)$，也就是说
\begin{align*}
0 \to M \to B \otimes_A M \to B \otimes_A B \otimes_A M
\end{align*}
由引理 \ref{etale-cohomology-lemma-descent-modules} 可知是正合的。
\end{proof}

\noindent
在环化网站上还有拟凝聚层的抽象概念。这里简要介绍它；
更多信息见 Modules on Sites，第 \ref{sites-modules-section-local} 节。
设 $\mathcal{C}$ 是一个范畴，$U$ 是 $\mathcal{C}$ 的对象。
记 $\mathcal{C}/U$ 为位于 $U$ 之上的对象所成的范畴，见
Categories，例 \ref{categories-example-category-over-X}。
若 $\mathcal{C}$ 是网站，则 $\mathcal{C}/U$ 也是网站；具体地，
$V/U$ 的覆盖是形如 $\{V_i/U \to V/U\}$ 的
$\mathcal{C}/U$ 态射族，并满足 $\{V_i \to V\}$ 是
$\mathcal{C}$ 的覆盖。此外，给定任意层 $\mathcal{F}$（定义在
$\mathcal{C}$ 上），其{\it 限制}
$\mathcal{F}|_{\mathcal{C}/U}$（按显然方式定义）也是层。详见
Sites，第 \ref{sites-section-localize} 节。

\begin{definition}
\label{etale-cohomology-definition-ringed-site}
设 $\mathcal{C}$ 是一个{\it 环化网站}，即带有环层
$\mathcal{O}$ 的网站。称 $\mathcal{O}$-模层 $\mathcal{F}$
在 $\mathcal{C}$ 上是{\it 拟凝聚的}，如果对所有
$U \in \Ob(\mathcal{C})$，都存在由
$\{U_i \to U\}_{i\in I}$ 给出的 $\mathcal{C}$ 覆盖，使得限制
$\mathcal{F}|_{\mathcal{C}/U_i}$ 同构于自由 $\mathcal{O}$-模之间某个
$\mathcal{O}$-线性映射的余核
$$
\bigoplus\nolimits_{k \in K} \mathcal{O}|_{\mathcal{C}/U_i}
\longrightarrow
\bigoplus\nolimits_{l \in L} \mathcal{O}|_{\mathcal{C}/U_i}.
$$
对 $K$ 的直和是与预层
$V \mapsto \bigoplus_{k \in K} \mathcal{O}(V)$ 关联的层，另一侧同理。
\end{definition}

\noindent
虽然能够给出上述一般定义很有用，但这个概念一般而言性质并不好。

\begin{remark}
\label{etale-cohomology-remark-final-object}
若 $\mathcal{C}$ 有终对象（例如 $S$），则上述定义中的条件只需
对 $U = S$ 检查即可。见 Modules on Sites，引理
\ref{sites-modules-lemma-local-final-object}。
\end{remark}

\begin{theorem}[拟凝聚层元定理]
\label{etale-cohomology-theorem-quasi-coherent}
设 $S$ 是概形。设 $\mathcal{C}$ 是一个网站，并假设
\begin{enumerate}
\item 底层范畴 $\mathcal{C}$ 是 $\Sch/S$ 的全子范畴；
\item $T \in \Ob(\mathcal{C})$ 的任意 Zariski 覆盖都可由
$\mathcal{C}$ 的某个覆盖细化；
\item $S/S$ 是 $\mathcal{C}$ 的对象；
\item $\mathcal{C}$ 的每个覆盖都是概形的 fpqc 覆盖。
\end{enumerate}
则预层 $\mathcal{O}$ 是 $\mathcal{C}$ 上的层，并且
任意拟凝聚 $\mathcal{O}$-模在 $(\mathcal{C}, \mathcal{O})$ 上都具有
$\mathcal{F}^a$ 的形式，其中 $\mathcal{F}$ 是定义在 $S$ 上的某个
拟凝聚层。
\end{theorem}

\begin{proof}
经过一些形式论证，这正是定理 \ref{etale-cohomology-theorem-descent-quasi-coherent}。
细节略去。在 Descent，命题
\ref{descent-proposition-equivalence-quasi-coherent} 中，我们对 $S$ 的
大 Zariski、fppf、\'etale、光滑和 syntomic 网站，以及 $S$ 的小 Zariski
和 \'etale 网站证明了该定理的更精确版本。
\end{proof}

\noindent
换言之，$S$ 上的拟凝聚模与定理中这类网站上的拟凝聚
$\mathcal{O}$-模（网站记为 $\mathcal{C}$）之间没有区别。关于大、小网站
$(\Sch/S)_{fppf}$、$S_\etale$ 等的更精确陈述，可见 Descent 的
以下各节：
\ref{descent-section-quasi-coherent-sheaves},
\ref{descent-section-quasi-coherent-cohomology}，和
\ref{descent-section-quasi-coherent-sheaves-bis}。
在本章中，为了方便陈述适用于上述任一情形的结果，
我们有时将其称为“定理 \ref{etale-cohomology-theorem-quasi-coherent} 中的这类网站”。






\section{{\v C}ech 上同调}
\label{etale-cohomology-section-cech-cohomology}

\noindent
下一目标是用下降理论证明等式
$H^i(\mathcal{C}, \mathcal{F}^a) = H_{Zar}^i(S, \mathcal{F})$
对所有拟凝聚层 $\mathcal{F}$（定义在 $S$ 上）以及定理
\ref{etale-cohomology-theorem-quasi-coherent} 中的任意网站 $\mathcal{C}$ 都成立。
为此，我们引入网站上的 {\v C}ech 上同调。更多细节见
\cite{ArtinTopologies}，以及 Cohomology on Sites 的各节
\ref{sites-cohomology-section-cech}、
\ref{sites-cohomology-section-cech-functor}
和 \ref{sites-cohomology-section-cech-cohomology-cohomology}。

\begin{definition}
\label{etale-cohomology-definition-cech-complex}
设 $\mathcal{C}$ 是一个范畴。设
$\mathcal{U} = \{U_i \to U\}_{i \in I}$ 是 $\mathcal{C}$ 中具有共同靶
的态射族。假设所有纤维积
$U_{i_0} \times_U \ldots \times_U U_{i_p}$ 都存在于 $\mathcal{C}$ 中。
设 $\mathcal{F} \in \textit{PAb}(\mathcal{C})$ 是阿贝尔预层。定义
{\it {\v C}ech 复形} $\check{\mathcal{C}}^\bullet(\mathcal{U}, \mathcal{F})$ 为
$$
\prod_{i_0 \in I} \mathcal{F}(U_{i_0}) \to
\prod_{i_0, i_1 \in I} \mathcal{F}(U_{i_0} \times_U U_{i_1}) \to
\prod_{i_0, i_1, i_2 \in I}
\mathcal{F}(U_{i_0} \times_U U_{i_1} \times_U U_{i_2}) \to \ldots
$$
其中第一项位于次数 0，各映射取通常的映射。定义
{\it {\v C}ech 上同调群}为
$$
\check{H}^p(\mathcal{U}, \mathcal{F}) =
H^p(\check{\mathcal{C}}^\bullet(\mathcal{U}, \mathcal{F})).
$$
\end{definition}

\noindent
在上述定义中，允许指标 $(i_0, \ldots, i_p)$ 重复至关重要。

\begin{lemma}
\label{etale-cohomology-lemma-cech-presheaves}
记号和假设同定义 \ref{etale-cohomology-definition-cech-complex}。函子
$\check{\mathcal{C}}^\bullet(\mathcal{U}, -)$ 在范畴
$\textit{PAb}(\mathcal{C})$ 上正合。
\end{lemma}

\noindent
换言之，若 $0\to \mathcal{F}_1\to \mathcal{F}_2\to \mathcal{F}_3\to 0$
是阿贝尔群预层的短正合列，则
$$
0 \to \check{\mathcal{C}}^\bullet\left(\mathcal{U}, \mathcal{F}_1\right)
\to\check{\mathcal{C}}^\bullet(\mathcal{U}, \mathcal{F}_2) \to
\check{\mathcal{C}}^\bullet(\mathcal{U}, \mathcal{F}_3)\to 0
$$
是复形的短正合列。

\begin{proof}
这由预层短正合列的定义立即得到。事实上，阿贝尔预层范畴是某个范畴到
$\textit{Ab}$ 的函子范畴，因而自动成为阿贝尔范畴：列
$\mathcal{F}_1\to \mathcal{F}_2\to \mathcal{F}_3$ 在
$\textit{PAb}$ 中正合，当且仅当对所有 $U \in \Ob(\mathcal{C})$，列
$\mathcal{F}_1(U) \to \mathcal{F}_2(U) \to \mathcal{F}_3(U)$ 在
$\textit{Ab}$ 中正合。因此，上述复形在每个次数上不过是短正合列的积。
另见 Cohomology on Sites，引理
\ref{sites-cohomology-lemma-cech-exact-presheaves}。
\end{proof}

\noindent
这说明 $\check{H}^\bullet(\mathcal{U}, -)$ 是 $\delta$-函子。
下面证明它是泛 $\delta$-函子。为此，需要证明它是{\it 可抹的}函子。
我们先回顾 Yoneda 引理。

\begin{lemma}[Yoneda 引理]
\label{etale-cohomology-lemma-yoneda-presheaf}
对任意预层 $\mathcal{F}$（定义在范畴 $\mathcal{C}$ 上）以及
$U \in \Ob(\mathcal{C})$，有函子性同构
$$
\Hom_{\textit{PSh}(\mathcal{C})}(h_U, \mathcal{F}) =
\mathcal{F}(U).
$$
\end{lemma}

\begin{proof}
见 Categories，引理 \ref{categories-lemma-yoneda}。
\end{proof}

\noindent
给定集合 $E$，本节记 $\mathbf{Z}[E]$ 为由 $E$ 生成的自由阿贝尔群。
用公式表示即为 $\mathbf{Z}[E] = \bigoplus_{e \in E} \mathbf{Z}$；也就是说，
$\mathbf{Z}[E]$ 是一个自由 $\mathbf{Z}$-模，其基由 $E$ 的元素组成。
利用这一记号，我们引入集合预层上的自由阿贝尔预层。

\begin{definition}
\label{etale-cohomology-definition-free-abelian-presheaf}
设 $\mathcal{C}$ 是一个范畴。给定集合预层 $\mathcal{G}$，定义
{\it $\mathcal{G}$ 上的自由阿贝尔预层}，记作
$\mathbf{Z}_\mathcal{G}$，规则为
$$
\mathbf{Z}_\mathcal{G}(U)
=
\mathbf{Z}[\mathcal{G}(U)]
$$
其中 $U \in \Ob(\mathcal{C})$，限制映射由 $\mathcal{G}$ 的限制映射诱导。
在特殊情形 $\mathcal{G} = h_U$ 下，简记为
$\mathbf{Z}_U = \mathbf{Z}_{h_U}$。
\end{definition}

\noindent
函子 $\mathcal{G} \mapsto \mathbf{Z}_\mathcal{G}$ 是遗忘函子
$\textit{PAb}(\mathcal{C}) \to \textit{PSh}(\mathcal{C})$ 的左伴随。
因此，对任意预层 $\mathcal{F}$，有典范同构
$$
\Hom_{\textit{PAb}(\mathcal{C})}(\mathbf{Z}_U, \mathcal{F})
=
\Hom_{\textit{PSh}(\mathcal{C})}(h_U, \mathcal{F})
=
\mathcal{F}(U)
$$
最后一个等式由 Yoneda 引理得到。特别地，有以下结果。

\begin{lemma}
\label{etale-cohomology-lemma-cech-complex-describe}
记号和假设同定义 \ref{etale-cohomology-definition-cech-complex}。{\v C}ech 复形
$\check{\mathcal{C}}^\bullet(\mathcal{U}, \mathcal{F})$ 可显式描述如下：
\begin{eqnarray*}
\check{\mathcal{C}}^\bullet(\mathcal{U}, \mathcal{F})
& = &
\left(
\prod_{i_0 \in I}
\Hom_{\textit{PAb}(\mathcal{C})}(\mathbf{Z}_{U_{i_0}}, \mathcal{F}) \to
\prod_{i_0, i_1 \in I}
\Hom_{\textit{PAb}(\mathcal{C})}(
\mathbf{Z}_{U_{i_0} \times_U U_{i_1}}, \mathcal{F}) \to \ldots
\right) \\
& = &
\Hom_{\textit{PAb}(\mathcal{C})}\left(
\left(
\bigoplus_{i_0 \in I} \mathbf{Z}_{U_{i_0}} \leftarrow
\bigoplus_{i_0, i_1 \in I} \mathbf{Z}_{U_{i_0} \times_U U_{i_1}} \leftarrow
\ldots
\right), \mathcal{F}\right)
\end{eqnarray*}
\end{lemma}

\begin{proof}
这由上式得到。见 Cohomology on Sites，引理
\ref{sites-cohomology-lemma-cech-map-into}。
\end{proof}

\noindent
这使我们只需研究最后一个 $\Hom$ 的第一变元中的复形。

\begin{lemma}
\label{etale-cohomology-lemma-exact}
记号和假设同定义 \ref{etale-cohomology-definition-cech-complex}。阿贝尔预层复形
\begin{align*}
\mathbf{Z}_\mathcal{U}^\bullet \quad : \quad
\bigoplus_{i_0 \in I} \mathbf{Z}_{U_{i_0}} \leftarrow
\bigoplus_{i_0, i_1 \in I} \mathbf{Z}_{U_{i_0} \times_U U_{i_1}} \leftarrow
\bigoplus_{i_0, i_1, i_2 \in I}
\mathbf{Z}_{U_{i_0} \times_U U_{i_1} \times_U U_{i_2}} \leftarrow
\ldots
\end{align*}
除次数 $0$ 外，该复形在 $\textit{PAb}(\mathcal{C})$ 中的所有次数上正合。
\end{lemma}

\begin{proof}
对任意 $V\in \Ob(\mathcal{C})$，阿贝尔群复形
$\mathbf{Z}_\mathcal{U}^\bullet(V)$ 为
$$
\begin{matrix}
\mathbf{Z}\left[
\coprod_{i_0\in I} \Mor_\mathcal{C}(V, U_{i_0})\right]
\leftarrow
\mathbf{Z}\left[
\coprod_{i_0, i_1 \in I}
\Mor_\mathcal{C}(V, U_{i_0} \times_U U_{i_1})\right]
\leftarrow \ldots = \\
\bigoplus_{\varphi : V \to U}
\left(
\mathbf{Z}\left[\coprod_{i_0 \in I} \Mor_\varphi(V, U_{i_0})\right]
\leftarrow
\mathbf{Z}\left[\coprod_{i_0, i_1\in I} \Mor_\varphi(V, U_{i_0}) \times
\Mor_\varphi(V, U_{i_1})\right]
\leftarrow
\ldots
\right)
\end{matrix}
$$
其中
$$
\Mor_{\varphi}(V, U_i)
=
\{ V \to U_i \text{，使得 } V \to U_i \to U \text{ 等于 } \varphi \}.
$$
令 $S_\varphi = \coprod_{i\in I} \Mor_\varphi(V, U_i)$，于是
$$
\mathbf{Z}_\mathcal{U}^\bullet(V)
=
\bigoplus_{\varphi : V \to U}
\left(
\mathbf{Z}[S_\varphi] \leftarrow
\mathbf{Z}[S_\varphi \times S_\varphi] \leftarrow
\mathbf{Z}[S_\varphi \times S_\varphi \times S_\varphi] \leftarrow
\ldots \right).
$$
因此，只需证明对每个 $S = S_\varphi$，复形
\begin{align*}
\mathbf{Z}[S] \leftarrow
\mathbf{Z}[S \times S] \leftarrow
\mathbf{Z}[S \times S \times S] \leftarrow \ldots
\end{align*}
在负次数上正合。为此，可以给出一个显式同伦。固定 $s\in S$，并定义
$K: n_{(s_0, \ldots, s_p)} \mapsto n_{(s, s_0,
\ldots, s_p)}.$ 容易验证，$K$ 是下述算子的零同伦：
$$
\delta :
\eta_{(s_0, \ldots, s_p)}
\mapsto
\sum\nolimits_{i = 0}^p (-1)^p \eta_{(s_0, \ldots, \hat s_i, \ldots, s_p)}.
$$
更多细节见 Cohomology on Sites，引理
\ref{sites-cohomology-lemma-homology-complex}。
\end{proof}

\begin{lemma}
\label{etale-cohomology-lemma-hom-injective}
记号和假设同定义 \ref{etale-cohomology-definition-cech-complex}。若 $\mathcal{I}$ 是
$\textit{PAb}(\mathcal{C})$ 的内射对象，则
$\check H^p(\mathcal{U}, \mathcal{I}) = 0$ 对所有 $p > 0$ 成立。
\end{lemma}

\begin{proof}
{\v C}ech 复形是将函子
$\Hom_{\textit{PAb}(\mathcal{C})}(-, \mathcal{I}) $ 作用于复形 $
\mathbf{Z}^\bullet_\mathcal{U} $ 所得，也就是说
$$
\check H^p(\mathcal{U}, \mathcal{I}) = H^p
(\Hom_{\textit{PAb}(\mathcal{C})} (\mathbf{Z}^\bullet_\mathcal{U},
\mathcal{I})).
$$
但我们刚刚看到 $\mathbf{Z}^\bullet_\mathcal{U}$ 在负次数上正合，而函子
$\Hom_{\textit{PAb}(\mathcal{C})}(-,
\mathcal{I})$ 正合，因此 $\Hom_{\textit{PAb}(\mathcal{C})}
(\mathbf{Z}^\bullet_\mathcal{U}, \mathcal{I})$ 在正次数上正合。
\end{proof}

\begin{theorem}
\label{etale-cohomology-theorem-cech-derived}
记号和假设同定义 \ref{etale-cohomology-definition-cech-complex}。在
$\textit{PAb}(\mathcal{C})$ 上，函子 $\check{H}^p(\mathcal{U}, -)$ 是
$\check{H}^0(\mathcal{U}, -)$ 的右导出函子。
\end{theorem}

\begin{proof}
由引理 \ref{etale-cohomology-lemma-hom-injective}，函子
$\check H^p(\mathcal{U}, -)$ 因为可抹而成为泛 $\delta$-函子。
$\check H^0(\mathcal{U}, -)$ 的右导出函子也是如此。由于二者在次数
$0$ 上一致，泛 $\delta$-函子的泛性质说明它们一致。更多细节见
Cohomology on Sites，引理
\ref{sites-cohomology-lemma-cech-cohomology-derived-presheaves}。
\end{proof}

\begin{remark}
\label{etale-cohomology-remark-presheaves-no-topology}
注意，前面所有陈述都只涉及预层，因此尚未使用拓扑。
\end{remark}




\section{{\v C}ech 到上同调谱序列}
\label{etale-cohomology-section-cech-ss}

\noindent
这个谱序列是证明层上同调基本结果的基础工具。

\begin{lemma}
\label{etale-cohomology-lemma-forget-injectives}
遗忘函子 $\textit{Ab}(\mathcal{C})\to \textit{PAb}(\mathcal{C})$
将内射对象映为内射对象。
\end{lemma}

\begin{proof}
这是形式上的：遗忘函子有左伴随，即层化，而层化是正合函子。
更多细节见 Cohomology on Sites，引理
\ref{sites-cohomology-lemma-injective-abelian-sheaf-injective-presheaf}。
\end{proof}

\begin{theorem}
\label{etale-cohomology-theorem-cech-ss}
设 $\mathcal{C}$ 是网站。对任意覆盖
$\mathcal{U} = \{U_i \to U\}_{i \in I}$（它覆盖
$U \in \Ob(\mathcal{C})$）以及任意阿贝尔层 $\mathcal{F}$（在
$\mathcal{C}$ 上），存在谱序列
$$
E_2^{p, q}
=
\check H^p(\mathcal{U}, \underline{H}^q(\mathcal{F}))
\Rightarrow
H^{p+q}(U, \mathcal{F}),
$$
其中 $\underline{H}^q(\mathcal{F})$ 是阿贝尔预层
$V \mapsto H^q(V, \mathcal{F})$。
\end{theorem}

\begin{proof}
选取内射消解 $\mathcal{F}\to \mathcal{I}^\bullet$（位于
$\textit{Ab}(\mathcal{C})$ 中），并考虑双复形
$\check{\mathcal{C}}^\bullet(\mathcal{U}, \mathcal{I}^\bullet)$ 及映射
$$
\xymatrix{
\Gamma(U, I^\bullet) \ar[r] &
\check{\mathcal{C}}^\bullet(\mathcal{U}, \mathcal{I}^\bullet) \\
& \check{\mathcal{C}}^\bullet(\mathcal{U}, \mathcal{F}) \ar[u]
}
$$
这里，水平映射是到最左列的自然映射
$\Gamma(U, I^\bullet) \to
\check{\mathcal{C}}^0(\mathcal{U}, \mathcal{I}^\bullet)$，垂直映射由
$\mathcal{F}\to \mathcal{I}^0$ 诱导，并落在最下行。根据假设，
$\mathcal{I}^\bullet$ 是 $\textit{Ab}(\mathcal{C})$ 中的内射对象复形，
故由引理 \ref{etale-cohomology-lemma-forget-injectives}，它也是
$\textit{PAb}(\mathcal{C})$ 中的内射对象复形。因此，双复形的各行在
正次数上正合（引理 \ref{etale-cohomology-lemma-hom-injective}），并且
$\check{\mathcal{C}}^0(\mathcal{U}, \mathcal{I}^\bullet)
\to \check{\mathcal{C}}^1(\mathcal{U}, \mathcal{I}^\bullet)$ 的核等于
$\Gamma(U, \mathcal{I}^\bullet)$，因为 $\mathcal{I}^\bullet$
是层的复形。特别地，总复形的上同调就是全局截面函子
$H^0(U, \mathcal{F})$ 的标准上同调。

\medskip\noindent
在垂直方向上，第 $q$ 个上同调群（取自第 $p$ 列）为
$$
\prod_{i_0, \ldots, i_p}
H^q(U_{i_0} \times_U \ldots \times_U U_{i_p}, \mathcal{F})
=
\prod_{i_0, \ldots, i_p}
\underline{H}^q(\mathcal{F})(U_{i_0} \times_U \ldots \times_U U_{i_p})
$$
它位于 $E_1^{p, q}$ 项。因此，这是标准的双复形谱序列，且
$E_2$-页如所述。更多细节见 Cohomology on Sites，引理
\ref{sites-cohomology-lemma-cech-spectral-sequence}。
\end{proof}

\begin{remark}
\label{etale-cohomology-remark-grothendieck-ss}
这是下述函子复合的 Grothendieck 谱序列：
$$
\textit{Ab}(\mathcal{C}) \longrightarrow
\textit{PAb}(\mathcal{C}) \xrightarrow{\check H^0} \textit{Ab}.
$$
\end{remark}








\section{概形的大网站与小网站}
\label{etale-cohomology-section-big-small}

\noindent
设 $S$ 是概形。设 $\tau$ 是我们将讨论的拓扑之一。因此
$\tau \in \{fppf, syntomic, smooth, \etale, Zariski\}$。
当然，如果只关心 \'etale 拓扑，则可始终假设 $\tau = \etale$。
此外，从第 \ref{etale-cohomology-section-etale-site} 节起，我们将更详细地讨论
\'etale 态射、\'etale 覆盖和 \'etale 网站。为了继续讨论拟凝聚层的
上同调，引入大 $\tau$-网站很方便；当
$\tau \in \{\etale, Zariski\}$ 时，还要引入小 $\tau$-网站，即 $S$ 的相应小网站。
为此，先引入 $\tau$-覆盖的概念。

\begin{definition}
\label{etale-cohomology-definition-tau-covering}
（见 Topologies，定义
\ref{topologies-definition-fppf-covering},
\ref{topologies-definition-syntomic-covering},
\ref{topologies-definition-smooth-covering},
\ref{topologies-definition-etale-covering}，和
\ref{topologies-definition-zariski-covering}。）
设 $\tau \in \{fppf, syntomic, smooth, \etale, Zariski\}$。
称具有共同靶的概形态射族 $\{f_i : T_i \to T\}_{i \in I}$ 是
{\it $\tau$-覆盖}，当且仅当每个 $f_i$ 分别是有限表示平坦态射、
syntomic 态射、光滑态射、\'etale 态射或开浸入，并且
$\bigcup f_i(T_i) = T$。
\end{definition}

\noindent
所有 $\tau$-覆盖所成的类满足定义 \ref{etale-cohomology-definition-site}（我们对网站的定义）
中的公理 (1)、(2) 和 (3)，见 Topologies，引理
\ref{topologies-lemma-fppf},
\ref{topologies-lemma-syntomic},
\ref{topologies-lemma-smooth},
\ref{topologies-lemma-etale}，和
\ref{topologies-lemma-zariski}。

\medskip\noindent
下面引入我们将使用的网站。与 \cite{SGA4} 中的做法不同，我们不想选取一个
宇宙，而是选取一个“部分宇宙”（即 Sets，第
\ref{sets-section-sets-hierarchy} 节所述的一个适当大的集合），并考虑包含在
这个集合中的所有概形。当然，我们确保所选基概形 $S$ 包含在部分宇宙中。
选定底层范畴后，再选取一个适当大的 $\tau$-覆盖集合，使之成为网站。
细节见概形拓扑一章；这些选择有很大自由度，但最终的具体选择不会影响
$S$ 的 \'etale（或其他）上同调（正如在 \cite{SGA4} 中，宇宙的具体选择
最终并不重要）。此外，本材料的写法也允许愿意使用强不可达基数（即宇宙）
的读者以之替代部分宇宙。

\begin{definition}
\label{etale-cohomology-definition-tau-site}
设 $S$ 是概形。
设 $\tau \in \{fppf, syntomic, smooth, \etale, \linebreak[0] Zariski\}$。
\begin{enumerate}
\item {\it 大 $\tau$-网站（基为 $S$）}是按上述方式构造的任一网站
$(\Sch/S)_\tau$；更详细的构造见 Topologies，定义
\ref{topologies-definition-big-small-fppf},
\ref{topologies-definition-big-small-syntomic},
\ref{topologies-definition-big-small-smooth},
\ref{topologies-definition-big-small-etale}，和
\ref{topologies-definition-big-small-Zariski}.
\item 若 $\tau \in \{\etale, Zariski\}$，则{\it 小 $\tau$-网站（基为 $S$）}
是全子范畴 $S_\tau$，位于 $(\Sch/S)_\tau$ 中；其对象是概形 $T$（位于 $S$ 上），
其结构态射 $T \to S$ 分别为 \'etale 态射或开浸入。$S_\tau$ 中的覆盖是族
$\{U_i \to U\}$，它在 $(\Sch/S)_\tau$ 中构成覆盖，其中 $U$ 是 $S_\tau$ 的对象。
\end{enumerate}
\end{definition}

\noindent
网站 $(\Sch/S)_\tau$ 的底层范畴具有合理的“封闭性”。例如，给定其中的概形
$T$，$T$ 的任意局部闭子概形都同构于 $(\Sch/S)_\tau$ 的某个对象。其他这类
封闭性包括：对概形的纤维积封闭；可取可数不交并；可取给定概形上的有限型
概形；给定仿射概形 $\Spec(R)$，可对 $R$ 作完备化、局部化或关于某个理想取商，
而仍留在该范畴中；等等。另一方面，例如对 $(\Sch/S)_\tau$ 中的概形取任意
不交并，则可能离开这个范畴。还要注意，给定对象 $T$（位于 $(\Sch/S)_\tau$ 中），可能
存在 $\tau$-覆盖 $\{T_i \to T\}_{i \in I}$（如定义
\ref{etale-cohomology-definition-tau-covering} 所述），但它不是 $(\Sch/S)_\tau$ 中的覆盖；例如，
概形 $T_i$ 可能不是范畴 $(\Sch/S)_\tau$ 的对象。不过，我们对网站
$(\Sch/S)_\tau$ 的选取保证，总存在覆盖 $\{U_j \to T\}_{j \in J}$
（属于 $(\Sch/S)_\tau$），它加细覆盖 $\{T_i \to T\}_{i \in I}$，见
Topologies，引理
\ref{topologies-lemma-fppf-induced},
\ref{topologies-lemma-syntomic-induced},
\ref{topologies-lemma-smooth-induced},
\ref{topologies-lemma-etale-induced}，和
\ref{topologies-lemma-zariski-induced}.
本章大多不再理会这些问题。

\medskip\noindent
若 $\mathcal{F}$ 是 $(\Sch/S)_\tau$ 或 $S_\tau$ 上的层，则以
$$
H^p_\tau(U, \mathcal{F}), \text{ 特别地 }
H^p_\tau(S, \mathcal{F})
$$
表示 $\mathcal{F}$ 在网站对象 $U$ 上的上同调群，见第
\ref{etale-cohomology-section-cohomology} 节。因此有
$H^p_{fppf}(S, \mathcal{F})$,
$H^p_{syntomic}(S, \mathcal{F})$,
$H^p_{smooth}(S, \mathcal{F})$,
$H^p_\etale(S, \mathcal{F})$，和
$H^p_{Zar}(S, \mathcal{F})$。后两种记号可能有歧义，因为它们既可能指大
\'etale 或 Zariski 网站，也可能指相应的小网站。不过，由下述引理可知，
这种歧义无害。

\begin{lemma}
\label{etale-cohomology-lemma-compare-cohomology-big-small}
设 $\tau \in \{\etale, Zariski\}$。若 $\mathcal{F}$ 是定义在
$(\Sch/S)_\tau$ 上的阿贝尔层，则 $\mathcal{F}$ 在 $S$ 上的上同调群与
$\mathcal{F}|_{S_\tau}$ 在 $S$ 上的上同调群相同。
\end{lemma}

\begin{proof}
由 Topologies，引理 \ref{topologies-lemma-at-the-bottom} 和
\ref{topologies-lemma-at-the-bottom-etale}
可知，函子 $S_\tau \to (\Sch/S)_\tau$ 满足 Sites，引理
\ref{sites-lemma-bigger-site} 的假设。因此，本引理由 Cohomology on Sites，
引理 \ref{sites-cohomology-lemma-cohomology-bigger-site} 得出。
\end{proof}

\noindent
$S$ 的大或小 \'etale 网站上的层范畴，只依赖于 $(\Sch/S)_\etale$ 或
$S_\etale$ 中由仿射对象组成的全子范畴；并且只需考虑这些对象之间的标准
\'etale 覆盖（定义见下文）。由此得到网站 $(\textit{Aff}/S)_\etale$ 和
$S_{affine, \etale}$，见 Topologies，定义
\ref{topologies-definition-big-small-etale}。比较结果见 Topologies，引理
\ref{topologies-lemma-affine-big-site-etale} 和
\ref{topologies-lemma-alternative}。下面定义本章将考虑的若干拓扑中的标准覆盖。

\begin{definition}
\label{etale-cohomology-definition-standard-tau}
（见 Topologies，定义
\ref{topologies-definition-standard-fppf},
\ref{topologies-definition-standard-syntomic},
\ref{topologies-definition-standard-smooth},
\ref{topologies-definition-standard-etale}，和
\ref{topologies-definition-standard-Zariski}.)
设 $\tau \in \{fppf, syntomic, smooth, \etale, Zariski\}$，并设 $T$ 是仿射概形。
一个{\it 标准 $\tau$-覆盖（靶为 $T$）}是一个族
$\{f_j : U_j \to T\}_{j = 1, \ldots, m}$，其中每个 $U_j$ 都是仿射的，
每个 $f_j$ 分别是有限表示平坦态射、标准 syntomic 态射、标准光滑态射、
\'etale 态射或 $T$ 中标准主开集的浸入，并且 $T = \bigcup f_j(U_j)$。
\end{definition}

\begin{lemma}
\label{etale-cohomology-lemma-tau-affine}
设 $\tau \in \{fppf, syntomic, smooth, \etale, Zariski\}$。
仿射概形的任意 $\tau$-覆盖都可由某个标准 $\tau$-覆盖加细。
\end{lemma}

\begin{proof}
见 Topologies，引理
\ref{topologies-lemma-fppf-affine},
\ref{topologies-lemma-syntomic-affine},
\ref{topologies-lemma-smooth-affine},
\ref{topologies-lemma-etale-affine}，和
\ref{topologies-lemma-zariski-affine}.
\end{proof}

\noindent
为完整起见，我们陈述并证明这里所考虑的上同调群不依赖于部分宇宙的选取。
下文将在引理 \ref{etale-cohomology-lemma-etale-topos-independent-partial-universe} 中证明小
\'etale 拓扑斯的不变性。本引理使用的记号和术语见 Topologies，第
\ref{topologies-section-change-alpha} 节。

\begin{lemma}
\label{etale-cohomology-lemma-cohomology-enlarge-partial-universe}
设 $\tau \in \{fppf, syntomic, smooth, \etale, Zariski\}$。设 $S$ 是概形。
设 $(\Sch/S)_\tau$ 和 $(\Sch'/S)_\tau$ 是两个大 $\tau$-网站（基为 $S$），
并假设前者包含于后者。此时：
\begin{enumerate}
\item 对任意阿贝尔层 $\mathcal{F}'$（定义在 $(\Sch'/S)_\tau$ 上），以及
任意对象 $U$（属于 $(\Sch/S)_\tau$），有
$$
H^p_\tau(U, \mathcal{F}'|_{(\Sch/S)_\tau}) =
H^p_\tau(U, \mathcal{F}')
$$
换言之，在较大网站中计算的 $\mathcal{F}'$ 在 $U$ 上的上同调，与把
$\mathcal{F}'$ 限制到较小网站后在 $U$ 上的上同调相同。
\item 对任意阿贝尔层 $\mathcal{F}$（在 $(\Sch/S)_\tau$ 上），存在阿贝尔层
$\mathcal{F}'$（在 $(\Sch/S)_\tau'$ 上），其到 $(\Sch/S)_\tau$ 的限制
同构于 $\mathcal{F}$。
\end{enumerate}
\end{lemma}

\begin{proof}
由 Topologies，引理 \ref{topologies-lemma-change-alpha} 可知，包含函子
$(\Sch/S)_\tau \to (\Sch'/S)_\tau$ 满足 Sites，引理
\ref{sites-lemma-bigger-site} 的假设。这推出 (2)，而 (1) 由
Cohomology on Sites，引理 \ref{sites-cohomology-lemma-cohomology-bigger-site}
得出。
\end{proof}




\section{\'Etale 拓扑斯}
\label{etale-cohomology-section-etale-topos}

\noindent
{\it 拓扑斯}是某个网站上的集合层所成的范畴，见 Sites，定义
\ref{sites-definition-topos}。因此，通常用“概形的 \'etale 拓扑斯”这一说法
指概形的小 \'etale 网站上的层范畴。正式定义如下。

\begin{definition}
\label{etale-cohomology-definition-etale-topos}
设 $S$ 是概形。
\begin{enumerate}
\item $S$ 的{\it \'etale 拓扑斯}，或称{\it 小 \'etale 拓扑斯}，是集合层范畴
$\Sh(S_\etale)$，这些层位于 $S$ 的小 \'etale 网站上。
\item $S$ 的{\it Zariski 拓扑斯}，或称{\it 小 Zariski 拓扑斯}，是集合层范畴
$\Sh(S_{Zar})$，这些层位于 $S$ 的小 Zariski 网站上。
\item 对 $\tau \in \{fppf, syntomic, smooth, \etale, Zariski\}$，
{\it 大 $\tau$-拓扑斯}是某个大 $\tau$-拓扑斯（基为 $S$）上的集合层范畴。
\end{enumerate}
\end{definition}

\noindent
注意，$S$ 的小 Zariski 拓扑斯就是 $S$ 的底层拓扑空间上的集合层范畴，见
Topologies，引理 \ref{topologies-lemma-Zariski-usual}。小 \'etale 拓扑斯
不依赖于构造小 \'etale 网站时所作的选择，但一般而言，大拓扑斯确实依赖于
这些选择。

\medskip\noindent
事实上，大或小 \'etale 拓扑斯只依赖于 $(\Sch/S)_\etale$ 或 $S_\etale$
中由仿射对象组成的全子范畴，见 Topologies，引理
\ref{topologies-lemma-affine-big-site-etale} 和
\ref{topologies-lemma-alternative}。例如，在下述引理的证明中将用到这一点。

\begin{lemma}
\label{etale-cohomology-lemma-etale-topos-independent-partial-universe}
设 $S$ 是概形。$S$ 的 \'etale 拓扑斯不依赖于定义
\ref{etale-cohomology-definition-tau-site} 中小 \'etale 网站的构造（在典范等价意义下）。
\end{lemma}

\begin{proof}
需要证明：给定两个大 \'etale 网站 $\Sch_\etale$ 和 $\Sch_\etale'$，
它们都包含 $S$；采用显然的记号，有
$\Sh(S_\etale) \cong \Sh(S_\etale')$。由 Topologies，引理
\ref{topologies-lemma-contained-in}，可假设
$\Sch_\etale \subset \Sch_\etale'$。由 Sets，引理
\ref{sets-lemma-what-is-in-it}，$S$ 上任意 \'etale 的仿射概形都同构于
$\Sch_\etale$ 和 $\Sch_\etale'$ 二者中的对象。因此，诱导函子
$S_{affine, \etale} \to S_{affine, \etale}'$
是等价。此外，显然该函子及其某个拟逆都把标准 \'etale 覆盖变为标准
\'etale 覆盖。因此，结论由 Topologies，引理
\ref{topologies-lemma-alternative} 得出。
\end{proof}





\section{拟凝聚层的上同调}
\label{etale-cohomology-section-cohomology-quasi-coherent}
%9.22.09

\noindent
先从一个简单引理开始（它在比所述情形更一般的范围内成立）。该引理说明，
标准覆盖的 {\v C}ech 复形等于形如 $\{\Spec(B) \to \Spec(A)\}$ 的 fpqc
覆盖的 {\v C}ech 复形，其中 $A \to B$ 忠实平坦。

\begin{lemma}
\label{etale-cohomology-lemma-cech-complex}
设 $\tau \in \{fppf, syntomic, smooth, \etale, Zariski\}$。设 $S$ 是概形。
设 $\mathcal{F}$ 是 $(\Sch/S)_\tau$ 上的阿贝尔层；也可设它是
$S_\tau$ 上的阿贝尔层（当 $\tau = \etale$ 时）。再设
$\mathcal{U} = \{U_i \to U\}_{i \in I}$
是该网站的标准 $\tau$-覆盖。令 $V = \coprod_{i \in I} U_i$。则：
\begin{enumerate}
\item $V$ 是仿射概形；
\item $\mathcal{V} = \{V \to U\}$ 是 fpqc 覆盖；它也是 $\tau$-覆盖，
除非 $\tau = Zariski$；
\item {\v C}ech 复形 $\check{\mathcal{C}}^\bullet (\mathcal{U}, \mathcal{F})$
与 $\check{\mathcal{C}}^\bullet (\mathcal{V}, \mathcal{F})$ 相同。
\end{enumerate}
\end{lemma}

\begin{proof}
标准 $\tau$-覆盖的定义见 Topologies，定义
\ref{topologies-definition-standard-Zariski},
\ref{topologies-definition-standard-etale},
\ref{topologies-definition-standard-smooth},
\ref{topologies-definition-standard-syntomic}，和
\ref{topologies-definition-standard-fppf}.
由定义，每个概形 $U_i$ 都是仿射的，且 $I$ 是有限集。因此 $V$ 是仿射概形。
显然 $V \to U$ 平坦且满射，故 $\mathcal{V}$ 是 fpqc 覆盖，见例
\ref{etale-cohomology-example-fpqc-coverings}。除 Zariski 情形外，覆盖 $\mathcal{V}$ 也是
$\tau$-覆盖，见 Topologies，定义
\ref{topologies-definition-etale-covering},
\ref{topologies-definition-smooth-covering},
\ref{topologies-definition-syntomic-covering}，和
\ref{topologies-definition-fppf-covering}.

\medskip\noindent
注意，$\mathcal{U}$ 是 $\mathcal{V}$ 的加细，因而有 {\v C}ech 复形的映射
$\check{\mathcal{C}}^\bullet (\mathcal{V}, \mathcal{F}) \to
\check{\mathcal{C}}^\bullet (\mathcal{U}, \mathcal{F})$，见
Cohomology on Sites，等式
(\ref{sites-cohomology-equation-map-cech-complexes})。接着注意，若
$T = \coprod_{j \in J} T_j$ 是 $\mathcal{F}$ 所定义的网站中若干概形的不交并，
则具有共同靶的态射族 $\{T_j \to T\}_{j \in J}$ 是 Zariski 覆盖，故
\begin{equation}
\label{etale-cohomology-equation-sheaf-coprod}
\mathcal{F}(T) =
\mathcal{F}(\coprod\nolimits_{j \in J} T_j) =
\prod\nolimits_{j \in J} \mathcal{F}(T_j)
\end{equation}
由 $\mathcal{F}$ 的层条件成立。这说明上面的 {\v C}ech 复形映射在每一次数上
都是同构，因为作为概形有
$$
V \times_U \ldots \times_U V
=
\coprod\nolimits_{i_0, \ldots i_p} U_{i_0} \times_U \ldots \times_U U_{i_p}
$$
\end{proof}

\noindent
注意，等式 (\ref{etale-cohomology-equation-sheaf-coprod}) 对一般预层不成立。即使对层，它也不在
任意网站上都成立，因为余积未必产生覆盖，也未必是不交的。不过，对所有常用
网站（至少对我们将研究的全部网站），该等式成立。

\begin{remark}
\label{etale-cohomology-remark-refinement}
在引理 \ref{etale-cohomology-lemma-cech-complex} 的陈述中，覆盖 $\mathcal{U}$ 是
$\mathcal{V}$ 的加细，反向却不成立。形如 $\{V \to U\}$ 的覆盖并不构成
$U$ 的所有覆盖所成范畴的始子范畴。不过，仍可借助这些覆盖计算 {\v C}ech
上同调 $\check H^n(U, \mathcal{F})$；这里它定义为：在覆盖 $\mathcal{U}$
（其靶为 $U$）所成范畴的反范畴上，对 $\mathcal{F}$ 关于 $\mathcal{U}$ 的
{\v C}ech 上同调群取余极限，而计算可借助覆盖 $\{V \to U\}$。若以后需要，
我们将陈述一个精确引理（它只对层成立）并补在此处。
\end{remark}

\begin{lemma}[上同调的局部性]
\label{etale-cohomology-lemma-locality-cohomology}
设 $\mathcal{C}$ 是网站，$\mathcal{F}$ 是 $\mathcal{C}$ 上的阿贝尔层，
$U$ 是 $\mathcal{C}$ 的对象，$p > 0$ 是整数，并且 $\xi \in
H^p(U, \mathcal{F})$。则存在覆盖
$\mathcal{U} = \{U_i \to U\}_{i \in I}$；它是 $U$ 在 $\mathcal{C}$ 中的
覆盖，并且 $\xi |_{U_i} = 0$ 对所有 $i \in I$ 都成立。
\end{lemma}

\begin{proof}
选取内射消解 $\mathcal{F} \to \mathcal{I}^\bullet$。于是，$\xi$ 由上闭链
$\tilde{\xi} \in \mathcal{I}^p(U)$ 表示，并且
$d^p(\tilde{\xi}) = 0$。根据假设，序列
$\mathcal{I}^{p - 1} \to \mathcal{I}^p \to \mathcal{I}^{p + 1}$ 在
$\textit{Ab}(\mathcal{C})$ 中正合。这意味着存在覆盖
$\mathcal{U} = \{U_i \to U\}_{i \in I}$，使得
$\tilde{\xi}|_{U_i} = d^{p - 1}(\xi_i)$，其中
$\xi_i \in \mathcal{I}^{p-1}(U_i)$。上同调类 $\xi|_{U_i}$ 由上闭链
$\tilde{\xi}|_{U_i}$ 表示，而后者是上边界，故该上同调类为零。更多细节见
Cohomology on Sites，引理
\ref{sites-cohomology-lemma-kill-cohomology-class-on-covering}。
\end{proof}

\begin{theorem}
\label{etale-cohomology-theorem-zariski-fpqc-quasi-coherent}
设 $S$ 是概形，$\mathcal{F}$ 是拟凝聚 $\mathcal{O}_S$-模。设
$\mathcal{C}$ 或者是 $(\Sch/S)_\tau$，其中
$\tau \in \{fppf, syntomic, smooth, \etale, Zariski\}$；或者是
$S_\etale$。则
$$
H^p(S, \mathcal{F}) = H^p_\tau(S, \mathcal{F}^a)
$$
对所有 $p \geq 0$ 都成立，其中：
\begin{enumerate}
\item 左端表示层 $\mathcal{F}$ 在概形 $S$ 的底层拓扑空间上的通常上同调；
\item 右端表示阿贝尔层 $\mathcal{F}^a$（见命题
\ref{etale-cohomology-proposition-quasi-coherent-sheaf-fpqc}）在网站 $\mathcal{C}$ 上的上同调。
\end{enumerate}
\end{theorem}

\begin{proof}
我们将证明
$H^p(U, f^*\mathcal{F}) = H^p_\tau(U, \mathcal{F}^a)$；这里
$f : U \to S$ 是网站 $\mathcal{C}$ 的任意对象。由层性质，结论在
$p = 0$ 时成立。

\medskip\noindent
假设 $U$ 是仿射的。我们要证明
$H^p_\tau(U, \mathcal{F}^a) = 0$ 对所有 $p > 0$ 都成立。对 $p$ 作归纳。
\begin{enumerate}
\item[p = 1]
取 $\xi \in H^1_\tau(U, \mathcal{F}^a)$。由引理
\ref{etale-cohomology-lemma-locality-cohomology}，存在 fpqc 覆盖
$\mathcal{U} = \{U_i \to U\}_{i \in I}$，使得 $\xi|_{U_i} = 0$ 对所有
$i \in I$ 都成立。加细 $\mathcal{U}$ 后，可假设 $\mathcal{U}$ 是标准
$\tau$-覆盖。应用定理 \ref{etale-cohomology-theorem-cech-ss} 的谱序列可知，$\xi$ 来自上同调类
$\check \xi \in \check H^1(\mathcal{U}, \mathcal{F}^a)$。考虑覆盖
$\mathcal{V} = \{\coprod_{i\in I} U_i \to U\}$。由引理
\ref{etale-cohomology-lemma-cech-complex}，
$\check H^\bullet(\mathcal{U}, \mathcal{F}^a) =
\check H^\bullet(\mathcal{V}, \mathcal{F}^a)$.
另一方面，因为 $\mathcal{V}$ 是形如 $\{\Spec(B) \to \Spec(A)\}$ 的覆盖，
且 $f^*\mathcal{F} = \widetilde{M}$；这里存在 $A$-模 $M$。所以 {\v C}ech 复形
$\check{\mathcal{C}}^\bullet(\mathcal{V}, \mathcal{F})$
就是复形 $(B/A)_\bullet \otimes_A M$。由引理
\ref{etale-cohomology-lemma-descent-modules}，有
$H^p((B/A)_\bullet \otimes_A M) = 0$（当 $p > 0$ 时），故
$\check \xi = 0$，进而
$\xi = 0$。
\item[p > 1]
取 $\xi \in H^p_\tau(U, \mathcal{F}^a)$。由引理
\ref{etale-cohomology-lemma-locality-cohomology}，存在 fpqc 覆盖
$\mathcal{U} = \{U_i \to U\}_{i \in I}$，使得 $\xi|_{U_i} = 0$ 对所有
$i \in I$ 都成立。加细 $\mathcal{U}$ 后，可假设 $\mathcal{U}$ 是标准
$\tau$-覆盖。应用定理 \ref{etale-cohomology-theorem-cech-ss} 的谱序列。注意，各个交
$U_{i_0} \times_U \ldots \times_U U_{i_p}$ 都是仿射的，故由归纳假设，上同调群
$$
E_2^{p, q} = \check H^p(\mathcal{U}, \underline{H}^q(\mathcal{F}^a))
$$
对所有 $0 < q < p$ 都消失。由此可见，$\xi$ 必来自某个
$\check \xi \in \check H^p(\mathcal{U}, \mathcal{F}^a)$。把
$\mathcal{U}$ 换成只含一个态射的覆盖 $\mathcal{V}$，并再次使用引理
\ref{etale-cohomology-lemma-descent-modules}，便知 {\v C}ech 上同调类 $\check \xi$ 必为零，
故 $\xi = 0$。
\end{enumerate}
接着假设 $U$ 是分离的。选取仿射开覆盖
$U = \bigcup_{i \in I} U_i$（它覆盖 $U$）。于是族
$\mathcal{U} = \{U_i \to U\}_{i \in I}$ 是 fpqc 覆盖；并且所有交
$U_{i_0} \times_U \ldots \times_U U_{i_p}$ 都是仿射的，因为 $U$ 分离。因此，定理
\ref{etale-cohomology-theorem-cech-ss} 的谱序列除第零行外各行均为零。所以
$$
H^p_\tau(U, \mathcal{F}^a) =
\check H^p(\mathcal{U}, \mathcal{F}^a) =
\check H^p(\mathcal{U}, \mathcal{F}) = H^p(U, \mathcal{F})
$$
；最后一个等式来自标准概形理论，见 Cohomology of Schemes，引理
\ref{coherent-lemma-cech-cohomology-quasi-coherent}。

\medskip\noindent
一般情形较为技术性；若要延伸这里给出的证明，还需讨论谱序列之间的映射，
故我们不作处理。它由 Descent，命题
\ref{descent-proposition-same-cohomology-quasi-coherent}（其证明采用稍有不同的
方法）与 Cohomology on Sites，引理
\ref{sites-cohomology-lemma-cohomology-of-open} 结合得出。
\end{proof}

\begin{remark}
\label{etale-cohomology-remark-right-derived-global-sections}
对定理 \ref{etale-cohomology-theorem-zariski-fpqc-quasi-coherent} 作一点说明。由于 $S$ 是范畴
$\mathcal{C}$ 的终对象，右端的上同调群就是整体截面函子的右导出函子。事实上，
证明表明
$H^p(U, f^*\mathcal{F}) = H^p_\tau(U, \mathcal{F}^a)$
对任意对象 $f : U \to S$（属于网站 $\mathcal{C}$）都成立。
\end{remark}





\section{层的例子}
\label{etale-cohomology-section-examples-sheaves}

\noindent
设 $S$ 和 $\tau$ 如第 \ref{etale-cohomology-section-big-small} 节所述。前面已经看到，任意可表
预层都是 $(\Sch/S)_\tau$ 或 $S_\tau$ 上的层，见引理
\ref{etale-cohomology-lemma-representable-sheaf-fpqc} 和注
\ref{etale-cohomology-remark-fpqc-finest}。下面给出一些特殊情形。

\begin{definition}
\label{etale-cohomology-definition-additive-sheaf}
在第 \ref{etale-cohomology-section-big-small} 节的任一网站 $(\Sch/S)_\tau$ 或 $S_\tau$ 上：
\begin{enumerate}
\item 层 $T \mapsto \Gamma(T, \mathcal{O}_T)$ 记作 $\mathcal{O}_S$ 或
$\mathbf{G}_a$；若要标明基概形，则记作 $\mathbf{G}_{a, S}$。
\item 类似地，层 $T \mapsto \Gamma(T, \mathcal{O}^*_T)$ 记作
$\mathcal{O}_S^*$ 或 $\mathbf{G}_m$；若要标明基概形，则记作
$\mathbf{G}_{m, S}$。
\item 任意网站上的{\it 常值层} $\underline{\mathbf{Z}/n\mathbf{Z}}$，是常值预层
$U \mapsto \mathbf{Z}/n\mathbf{Z}$ 的层化。
\end{enumerate}
\end{definition}

\noindent
例如，由定理 \ref{etale-cohomology-theorem-quasi-coherent} 可知第一个是层。第二个是第一个的
子预层，而且容易看出它本身也是层。第三个按定义即为层。注意，这些层都可表。
前两个分别由概形 $\mathbf{G}_{a, S}$ 和 $\mathbf{G}_{m, S}$ 表示，见
Groupoids，第 \ref{groupoids-section-group-schemes} 节。第三个由有限 \'etale
群概形 $\mathbf{Z}/n\mathbf{Z}_S$ 表示，它有时记作
$(\mathbf{Z}/n\mathbf{Z})_S$；这就是 $n$ 个 $S$ 的副本，并赋予 $S$ 上显然的
群概形结构，见 Groupoids，例 \ref{groupoids-example-constant-group} 及下述注。

\begin{remark}
\label{etale-cohomology-remark-constant-locally-constant-maps}
设 $G$ 是抽象群。在第 \ref{etale-cohomology-section-big-small} 节的任一网站
$(\Sch/S)_\tau$ 或 $S_\tau$ 上，层化 $\underline{G}$（即与 $G$ 相伴的
常值预层在该网站的{\it Zariski 拓扑}中的层化）已经给出
$$
\Gamma(U, \underline{G}) =
\{\text{Zariski 局部常值映射 }U \to G\}
$$
。由 Groupoids，例 \ref{groupoids-example-constant-group}，这个 Zariski 层
由群概形 $G_S$ 表示。由引理 \ref{etale-cohomology-lemma-representable-sheaf-fpqc}，任意可表预层
也满足 $\tau$-拓扑的层条件，故上述 Zariski 层化 $\underline{G}$ 同时也是
$\tau$-层化。
\end{remark}

\begin{definition}
\label{etale-cohomology-definition-structure-sheaf}
设 $S$ 是概形。$S$ 的{\it 结构层}是环层 $\mathcal{O}_S$，它位于上述任一网站
$S_{Zar}$、$S_\etale$ 或 $(\Sch/S)_\tau$ 上。
\end{definition}

\noindent
若所使用的网站可能产生混淆，我们就用下标加以标明。例如，可用
$\mathcal{O}_{S_\etale}$ 强调我们是在 $S$ 的小 \'etale 网站上工作。

\begin{remark}
\label{etale-cohomology-remark-special-case-fpqc-cohomology-quasi-coherent}
用上面引入的术语，定理 \ref{etale-cohomology-theorem-zariski-fpqc-quasi-coherent} 的一个
特殊情形是
$$
H_{fppf}^p(X, \mathbf{G}_a) =
H_\etale^p(X, \mathbf{G}_a) =
H_{Zar}^p(X, \mathbf{G}_a) =
H^p(X, \mathcal{O}_X)
$$
，对所有 $p \geq 0$ 成立。此外，可以用记号
$H^p_{fppf}(X, \mathcal{O}_X)$ 表示结构层在 $X$ 的大 fppf 网站上的上同调。
\end{remark}




\section{Picard 群}
\label{etale-cohomology-section-picard-groups}

\noindent
下述定理有时称为“Hilbert 定理 90”。

\begin{theorem}
\label{etale-cohomology-theorem-picard-group}
对任意概形 $X$，有典范等同
\begin{align*}
H_{fppf}^1(X, \mathbf{G}_m) & = H^1_{syntomic}(X, \mathbf{G}_m) \\
& = H^1_{smooth}(X, \mathbf{G}_m) \\
& = H_\etale^1(X, \mathbf{G}_m) \\
& = H^1_{Zar}(X, \mathbf{G}_m) \\
& = \Pic(X) \\
& = H^1(X, \mathcal{O}_X^*)
\end{align*}
\end{theorem}

\begin{proof}
设 $\tau$ 是第 \ref{etale-cohomology-section-big-small} 节所考虑的拓扑之一。由
Cohomology on Sites，引理
\ref{sites-cohomology-lemma-h1-invertible} 可知
$H^1_\tau(X, \mathbf{G}_m) =
H^1_\tau(X, \mathcal{O}_\tau^*) =
\Pic(\mathcal{O}_\tau)$
，其中 $\mathcal{O}_\tau$ 是网站 $(\Sch/X)_\tau$ 的结构层。于是，可逆
$\mathcal{O}_\tau$-模是拟凝聚 $\mathcal{O}_\tau$-模。由定理
\ref{etale-cohomology-theorem-quasi-coherent}，或更精确的 Descent，命题
\ref{descent-proposition-equivalence-quasi-coherent}，可知
$\Pic(\mathcal{O}_\tau) = \Pic(X)$。最后一个等式以同样方式证明。
\end{proof}






\section{\'Etale 网站}
\label{etale-cohomology-section-etale-site}

\noindent
从这里开始，我们更详细地研究概形的 \'etale 网站。第一步是简要讨论
\'etale 态射的概念。





\section{\'Etale 态射}
\label{etale-cohomology-section-etale-morphism}

\noindent
正式定义见 Morphisms，第 \ref{morphisms-section-etale} 节；\'etale 态射的
若干有趣性质综述见 \'Etale Morphisms，诸节
\ref{etale-section-etale-morphisms},
\ref{etale-section-structure-etale-map},
\ref{etale-section-etale-smooth},
\ref{etale-section-topological-etale},
\ref{etale-section-functorial-etale}，和
\ref{etale-section-properties-permanence}。

\medskip\noindent
回忆：代数 $A$（在代数闭域 $k$ 上）称为{\it 光滑的}，如果它是有限型的，且微分模
$\Omega_{A/k}$ 是秩等于维数的有限局部自由模。概形 $X$（在 $k$ 上）称为在 $k$
上{\it 光滑}，如果它局部为有限型，并且每个仿射开集都是某个光滑 $k$-代数的
谱。若 $k$ 不是代数闭的，则称 $k$-代数 $A$ 是光滑 $k$-代数，如果
$A \otimes_k \overline{k}$ 是光滑 $\overline{k}$-代数。环映射 $A \to B$
称为光滑的，如果它平坦、有限表示，并且对所有素理想
$\mathfrak p \subset A$，纤维环 $\kappa(\mathfrak p) \otimes_A B$ 在剩余域
$\kappa(\mathfrak p)$ 上光滑。更一般地，概形态射称为{\it 光滑的}，如果它
平坦、局部有限表示，并且各几何纤维光滑。

\medskip\noindent
关于这些事实，见 Morphisms，第 \ref{morphisms-section-smooth} 节。由此可如下
定义 \'etale 态射。

\begin{definition}
\label{etale-cohomology-definition-etale-morphism}
若概形态射是相对维数为 0 的光滑态射，则称它是{\it \'etale 的}。
\end{definition}

\noindent
特别地，若概形态射 $X \to S$ 光滑且 $\Omega_{X/S} = 0$，则它是 \'etale 的。

\begin{proposition}
\label{etale-cohomology-proposition-etale-morphisms}
关于 \'etale 态射的若干事实。
\begin{enumerate}
\item 设 $k$ 是域。概形态射 $U \to \Spec(k)$ 是 \'etale 的，当且仅当
$U \cong \coprod_{i \in I} \Spec(k_i)$，其中对每个 $i \in I$，环 $k_i$ 是域，
并且是 $k$ 的有限可分扩张。
\item 设 $\varphi : U \to S$ 是概形态射。下列条件等价：
\begin{enumerate}
\item $\varphi$ 是 \'etale 的；
\item $\varphi$ 局部有限表示且平坦，并且它的所有纤维都是 \'etale 的；
\item $\varphi$ 平坦、非分歧且局部有限表示。
\end{enumerate}
\item 环映射 $A \to B$ 是 \'etale 的，当且仅当
$B \cong A[x_1, \ldots, x_n]/(f_1, \ldots, f_n)$
，并且 $\Delta = \det \left( \frac{\partial f_i}{\partial x_j} \right)$
在 $B$ 中可逆。
\item \'Etale 态射的基变换仍是 \'etale 态射。
\item \'Etale 态射的复合仍是 \'etale 态射。
\item \'Etale 态射的纤维积和乘积仍是 \'etale 态射。
\item \'Etale 态射的相对维数为 0。
\item 设 $Y \to X$ 是 \'etale 态射。若 $X$ 是既约的（相应地，正则的），
则 $Y$ 也是如此。
\item \'Etale 态射是开的。
\item 若 $X \to S$ 和 $Y \to S$ 都是 \'etale 的，则任意 $S$-态射
$X \to Y$ 也是 \'etale 的。
\end{enumerate}
\end{proposition}

\begin{proof}
我们已在前面各章证明了这些事实以及更多结果。参考文献如下：
(1) Morphisms，引理 \ref{morphisms-lemma-etale-over-field}。
(2) Morphisms，引理 \ref{morphisms-lemma-etale-flat-etale-fibres} 和
\ref{morphisms-lemma-flat-unramified-etale}。
(3) Algebra，引理 \ref{algebra-lemma-etale-standard-smooth}。
(4) Morphisms，引理 \ref{morphisms-lemma-base-change-etale}。
(5) Morphisms，引理 \ref{morphisms-lemma-composition-etale}。
(6) 由 (4) 和 (5) 形式地得出。
(7) Morphisms，引理 \ref{morphisms-lemma-etale-locally-quasi-finite} 和
\ref{morphisms-lemma-locally-quasi-finite-rel-dimension-0}。
(8) 见 Algebra，引理 \ref{algebra-lemma-reduced-goes-up} 和
\ref{algebra-lemma-Rk-goes-up}；此类更多结果还见 \'Etale Morphisms，第
\ref{etale-section-properties-permanence} 节。
(9) 见 Morphisms，引理 \ref{morphisms-lemma-fppf-open} 和
\ref{morphisms-lemma-etale-flat}。
(10) 见 Morphisms，引理 \ref{morphisms-lemma-etale-permanence}。
\end{proof}

\begin{definition}
\label{etale-cohomology-definition-standard-etale}
若环映射 $A \to B$ 满足
$B \cong \left(A[t]/(f)\right)_g$，其中 $f, g \in A[t]$、$f$ 首一，且
$\text{d}f/\text{d}t$ 在 $B$ 中可逆，则称它是{\it 标准 \'etale 的}。
\end{definition}

\noindent
标准 \'etale 环映射确实是 \'etale 的。事实上，假设
$B = \left(A[t]/(f)\right)_g$，其中 $f, g \in A[t]$、$f$ 首一，且
$\text{d}f/\text{d}t$ 在 $B$ 中可逆。则 $A[t]/(f)$ 是有限自由 $A$-模，
其秩等于首一多项式 $f$ 的次数。因此，$B$ 作为这个自由代数的局部化，
在 $A$ 上有限表示且平坦。要完成 $B$ 为 \'etale 的证明，只需说明纤维环
$$
\kappa(\mathfrak p) \otimes_A B
\cong
\kappa(\mathfrak p) \otimes_A (A[t]/(f))_g
\cong
\kappa(\mathfrak p)[t, 1/\overline{g}]/(\overline{f})
$$
是有限可分域扩张的有限乘积。这里
$\overline{f}, \overline{g} \in \kappa(\mathfrak p)[t]$ 分别是 $f$ 和 $g$ 的像。
令
$$
\overline{f} = \overline{f}_1 \ldots \overline{f}_a
\overline{f}_{a + 1}^{e_1} \ldots \overline{f}_{a + b}^{e_b}
$$
是 $\overline{f}$ 分解成两两不同的不可约首一因子 $\overline{f}_i$ 的幂，
其中 $e_1, \ldots, e_b > 1$。根据假设，
$\text{d}\overline{f}/\text{d}t$ 在
$\kappa(\mathfrak p)[t, 1/\overline{g}]$ 中可逆。因此至少所有
$\overline{f}_i$（其中 $i > a$）都可逆。由此得出
$$
\kappa(\mathfrak p)[t, 1/\overline{g}]/(\overline{f})
\cong
\prod\nolimits_{i \in I} \kappa(\mathfrak p)[t]/(\overline{f}_i)
$$
，其中 $I \subset \{1, \ldots, a\}$ 是指标 $i$ 的子集，所取指标满足
$\overline{f}_i$ 不整除 $\overline{g}$。此外，
$\text{d}\overline{f}/\text{d}t$ 在因子
$\kappa(\mathfrak p)[t]/(\overline{f}_i)$ 中的像，显然等于某个单位乘以
$\text{d}\overline{f}_i/\text{d}t$。因此，
$\kappa_i = \kappa(\mathfrak p)[t]/(\overline{f}_i)$ 是
$\kappa(\mathfrak p)$ 的有限域扩张，由一个极小多项式可分的元素生成；也就是说，
域扩张 $\kappa_i/\kappa(\mathfrak p)$ 如所需是有限可分的。

\medskip\noindent
事实上，任意 \'etale 环映射都局部标准 \'etale。为陈述这一点，引入下述记号。
称环映射 $A \to B${\it 在素理想 $\mathfrak q$（它是 $B$ 的素理想）处
\'etale}，如果存在 $h \in B$、$h \not \in \mathfrak q$，使得
$A \to B_h$ 是 \'etale 的。结果如下。

\begin{theorem}
\label{etale-cohomology-theorem-standard-etale}
环映射 $A \to B$ 在素理想 $\mathfrak q$ 处 \'etale，当且仅当存在
$g \in B$、$g \not \in \mathfrak q$，使得 $B_g$ 在 $A$ 上标准 \'etale。
\end{theorem}

\begin{proof}
见 Algebra，命题 \ref{algebra-proposition-etale-locally-standard}。
\end{proof}





\section{\'Etale 覆盖}
\label{etale-cohomology-section-etale-covering}

\noindent
回忆其定义。

\begin{definition}
\label{etale-cohomology-definition-etale-covering}
概形 $U$ 的一个{\it \'etale 覆盖}是概形态射族
$\{\varphi_i : U_i \to U\}_{i \in I}$，满足：
\begin{enumerate}
\item 每个 $\varphi_i$ 都是 \'etale 态射；
\item 诸 $U_i$ 覆盖 $U$，即 $U = \bigcup_{i\in I}\varphi_i(U_i)$。
\end{enumerate}
\end{definition}

\begin{lemma}
\label{etale-cohomology-lemma-etale-fpqc}
任意 \'etale 覆盖都是 fpqc 覆盖。
\end{lemma}

\begin{proof}
（另见 Topologies，引理
\ref{topologies-lemma-zariski-etale-smooth-syntomic-fppf-fpqc}。）设
$\{\varphi_i : U_i \to U\}_{i \in I}$ 是 \'etale 覆盖。由于 \'etale 态射平坦，
且覆盖的成员覆盖其靶，性质 fp（忠实平坦）成立。为验证性质 qc（拟紧），设
$V \subset U$ 是仿射开集，并写成
$\varphi_i^{-1}(V) = \bigcup_{j \in J_i} V_{ij}$，其中
$V_{ij} \subset U_i$ 是某些仿射开集。因为 $\varphi_i$ 是开态射
（\'etale 态射是开的），所以
$V = \bigcup_{i\in I} \bigcup_{j \in J_i} \varphi_i(V_{ij})$
是 $V$ 的开覆盖。又因 $V$ 拟紧，该覆盖有有限加细。
\end{proof}

\noindent
因此，对 fpqc 覆盖成立的任意陈述，{\it 更加}对 \'etale 覆盖成立。例如，
\'etale 网站是次典范的。

\begin{definition}
\label{etale-cohomology-definition-big-etale-site}
（更多细节见第 \ref{etale-cohomology-section-big-small} 节，或 Topologies，第
\ref{topologies-section-etale} 节。）设 $S$ 是概形。{\it $S$ 上的大 \'etale
网站}是网站 $(\Sch/S)_\etale$，见定义 \ref{etale-cohomology-definition-tau-site}。
{\it $S$ 上的小 \'etale 网站}是网站 $S_\etale$，见定义
\ref{etale-cohomology-definition-tau-site}。类似地，定义 $S$ 上的{\it 大}和{\it 小 Zariski
网站}，分别记作 $(\Sch/S)_{Zar}$ 和 $S_{Zar}$。
\end{definition}

\noindent
粗略地说，$S$ 的大 \'etale 网站由 $S$ 上的概形组成，其覆盖为 \'etale 覆盖；
$S$ 的小 \'etale 网站由在 $S$ 上 \'etale 的概形组成，其覆盖仍为 \'etale 覆盖。
事实上，由命题 \ref{etale-cohomology-proposition-etale-morphisms}，$S_\etale$ 的对象之间的任意态射
都是 \'etale 的。因此，要验证 $\{U_i \to U\}_{i \in I}$（位于 $S_\etale$ 中）
是覆盖，只需验证 $\coprod U_i \to U$ 满射。

\medskip\noindent
小 \'etale 网站的对象少于大 \'etale 网站；它只包含 $S$ 的 \'etale 拓扑中的
“开集”。它是大 \'etale 网站的全子范畴，其拓扑由大网站上的拓扑诱导。因此，
从大 \'etale 网站到小网站的限制函子正合，并把内射对象映为内射对象。由此得到
下述结论。

\begin{proposition}
\label{etale-cohomology-proposition-cohomology-restrict-small-site}
设 $S$ 是概形，$\mathcal{F}$ 是 $(\Sch/S)_\etale$ 上的阿贝尔层。则
$\mathcal{F}|_{S_\etale}$ 是 $S_\etale$ 上的层，并且
$$
H^p_\etale(S, \mathcal{F}|_{S_\etale}) =
H^p_\etale(S, \mathcal{F})
$$
对所有 $p \geq 0$ 都成立。
\end{proposition}

\begin{proof}
这是引理 \ref{etale-cohomology-lemma-compare-cohomology-big-small} 的特殊情形。
\end{proof}

\noindent
按照第 \ref{etale-cohomology-section-big-small} 节引入的一般记号，把上述上同调群记作
$H_\etale^p(S, \mathcal{F})$。





%9.24.09
\section{Kummer 理论}
\label{etale-cohomology-section-kummer}

\noindent
设 $n \in \mathbf{N}$，并考虑如下定义的函子 $\mu_n$：
$$
\begin{matrix}
\Sch^{opp} & \longrightarrow & \textit{Ab} \\
S & \longmapsto &
\mu_n(S)
=
\{t \in \Gamma(S, \mathcal{O}_S^*) \mid t^n = 1 \}.
\end{matrix}
$$
由 Groupoids，例 \ref{groupoids-example-roots-of-unity}，这是可表函子，表示它的
概形也记作 $\mu_n$。由引理 \ref{etale-cohomology-lemma-representable-sheaf-fpqc}，该函子满足
fpqc 拓扑的层条件（特别地，它也满足 \'etale、Zariski 等拓扑的层条件）。

\begin{lemma}
\label{etale-cohomology-lemma-kummer-sequence}
若 $n\in \mathcal{O}_S^*$，则
$$
0 \to
\mu_{n, S} \to
\mathbf{G}_{m, S} \xrightarrow{(\cdot)^n}
\mathbf{G}_{m, S} \to 0
$$
在 $S$ 的小 \'etale 网站和大 \'etale 网站上都是层的短正合列。
\end{lemma}

\begin{proof}
按定义，层 $\mu_{n, S}$ 是映射 $(\cdot)^n$ 的核。因此，只需证明最后一个映射
满射。设 $U$ 是 $S$ 上的概形。设
$f \in \mathbf{G}_m(U) = \Gamma(U, \mathcal{O}_U^*)$。需要证明：可找到 $U$ 的
一个 \'etale 覆盖，使得 $f$ 限制到每个覆盖成员后都是 $n$ 次幂。令
$$
U' =
\underline{\Spec}_U(\mathcal{O}_U[T]/(T^n-f))
\xrightarrow{\pi}
U.
$$
（关于相对谱，见 Constructions，第
\ref{constructions-section-spec-via-glueing} 或
\ref{constructions-section-spec} 节。）设 $\Spec(A) \subset U$ 是仿射开集，
并设 $f|_{\Spec(A)}$ 对应于单位 $a \in A^*$。则
$\pi^{-1}(\Spec(A)) = \Spec(B)$，其中 $B = A[T]/(T^n - a)$。环映射
$A \to B$ 是秩为 $n$ 的有限自由映射，因而忠实平坦，所以
$\Spec(B) \to \Spec(A)$ 满射。该结论对 $U$ 中每个仿射开集都成立，故 $\pi$
满射。此外，$n$ 和 $T^{n - 1}$ 在 $B$ 中可逆，所以
$nT^{n-1} \in B^*$，并且环映射 $A \to B$ 标准 \'etale，特别地是 \'etale 的。
该结论对 $U$ 中每个仿射开集都成立，故 $\pi$ 是 \'etale 的。因此
$\mathcal{U} = \{\pi : U' \to U\}$ 是 \'etale 覆盖。此外，
$f|_{U'} = (f')^n$，其中 $f'$ 是 $T$ 在
$\Gamma(U', \mathcal{O}_{U'}^*)$ 中的类，故 $\mathcal{U}$ 具有所需性质。
\end{proof}

\begin{remark}
\label{etale-cohomology-remark-no-kummer-sequence-zariski}
若把引理 \ref{etale-cohomology-lemma-kummer-sequence} 中的“\'etale”换成“Zariski”，则引理不成立。
由于 \'etale 拓扑比光滑拓扑粗，见 Topologies，引理
\ref{topologies-lemma-zariski-etale-smooth}，可知该序列在光滑拓扑中也正合。
\end{remark}

\noindent
由定理 \ref{etale-cohomology-theorem-picard-group}、引理 \ref{etale-cohomology-lemma-kummer-sequence} 以及上同调的
一般性质，得到上同调长正合列
$$
\xymatrix{
0 \ar[r] &
H_\etale^0(S, \mu_{n, S}) \ar[r] &
\Gamma(S, \mathcal{O}_S^*) \ar[r]^{(\cdot)^n} &
\Gamma(S, \mathcal{O}_S^*) \ar@(rd, ul)[rdllllr]
\\
& H_\etale^1(S, \mu_{n, S}) \ar[r] &
\Pic(S) \ar[r]^{(\cdot)^n} &
\Pic(S) \ar@(rd, ul)[rdllllr] \\
& H_\etale^2(S, \mu_{n, S}) \ar[r] &
\ldots
}
$$
；至少当 $n$ 在 $S$ 上可逆时如此。当 $n$ 在 $S$ 上不可逆时，可应用下述引理。

\begin{lemma}
\label{etale-cohomology-lemma-kummer-sequence-syntomic}
对任意 $n \in \mathbf{N}$，序列
$$
0 \to
\mu_{n, S} \to
\mathbf{G}_{m, S} \xrightarrow{(\cdot)^n}
\mathbf{G}_{m, S} \to 0
$$
在网站 $(\Sch/S)_{fppf}$ 和 $(\Sch/S)_{syntomic}$ 上都是层的短正合列。
\end{lemma}

\begin{proof}
按定义，层 $\mu_{n, S}$ 是映射 $(\cdot)^n$ 的核。因此，只需证明最后一个映射
满射。由于 syntomic 拓扑弱于 fppf 拓扑，见 Topologies，引理
\ref{topologies-lemma-zariski-etale-smooth-syntomic-fppf}，只需对 syntomic 拓扑
证明此事。设 $U$ 是 $S$ 上的概形。设
$f \in \mathbf{G}_m(U) = \Gamma(U, \mathcal{O}_U^*)$。需要证明：可找到 $U$ 的
一个 syntomic 覆盖，使得 $f$ 限制到每个覆盖成员后都是 $n$ 次幂。令
$$
U' =
\underline{\Spec}_U(\mathcal{O}_U[T]/(T^n-f))
\xrightarrow{\pi}
U.
$$
（关于相对谱，见 Constructions，第
\ref{constructions-section-spec-via-glueing} 或
\ref{constructions-section-spec} 节。）设 $\Spec(A) \subset U$ 是仿射开集，
并设 $f|_{\Spec(A)}$ 对应于单位 $a \in A^*$。则
$\pi^{-1}(\Spec(A)) = \Spec(B)$，其中 $B = A[T]/(T^n - a)$。环映射
$A \to B$ 是秩为 $n$ 的有限自由映射，因而忠实平坦，所以
$\Spec(B) \to \Spec(A)$ 满射。该结论对 $U$ 中每个仿射开集都成立，故 $\pi$
满射。此外，$B$ 在 $A$ 上是整体相对完全交，所以环映射 $A \to B$ 标准
syntomic，特别地是 syntomic 的。该结论对 $U$ 中每个仿射开集都成立，故 $\pi$
是 syntomic 的。因此 $\mathcal{U} = \{\pi : U' \to U\}$ 是 syntomic 覆盖。
此外，$f|_{U'} = (f')^n$，其中 $f'$ 是 $T$ 在
$\Gamma(U', \mathcal{O}_{U'}^*)$ 中的类，故 $\mathcal{U}$ 具有所需性质。
\end{proof}

\begin{remark}
\label{etale-cohomology-remark-no-kummer-sequence-smooth-etale-zariski}
引理 \ref{etale-cohomology-lemma-kummer-sequence-syntomic} 对光滑、\'etale 或 Zariski 拓扑不成立。
\end{remark}

\noindent
由定理 \ref{etale-cohomology-theorem-picard-group}、引理
\ref{etale-cohomology-lemma-kummer-sequence-syntomic} 以及上同调的一般性质，得到上同调长正合列
$$
\xymatrix{
0 \ar[r] &
H_{fppf}^0(S, \mu_{n, S}) \ar[r] &
\Gamma(S, \mathcal{O}_S^*) \ar[r]^{(\cdot)^n} &
\Gamma(S, \mathcal{O}_S^*) \ar@(rd, ul)[rdllllr]
\\
& H_{fppf}^1(S, \mu_{n, S}) \ar[r] &
\Pic(S) \ar[r]^{(\cdot)^n} &
\Pic(S) \ar@(rd, ul)[rdllllr] \\
& H_{fppf}^2(S, \mu_{n, S}) \ar[r] &
\ldots
}
$$
，它对任意概形 $S$ 和任意整数 $n$ 成立。当然，对 syntomic 上同调也有类似序列。

\medskip\noindent
设 $n \in \mathbf{N}$，并设 $S$ 是任意概形。还有一种更直接的方法描述取值于
$\mu_n$ 的第一上同调群。考虑序对 $(\mathcal{L}, \alpha)$，其中
$\mathcal{L}$ 是 $S$ 上的可逆层，而
$\alpha : \mathcal{L}^{\otimes n} \to \mathcal{O}_S$ 是第 $n$ 次张量幂的平凡化，
这一张量幂来自 $\mathcal{L}$。设 $(\mathcal{L}', \alpha')$ 是另一这样的序对。同构
$\varphi : (\mathcal{L}, \alpha) \to (\mathcal{L}', \alpha')$ 是可逆层的同构
$\varphi : \mathcal{L} \to \mathcal{L}'$，并使下图
$$
\xymatrix{
\mathcal{L}^{\otimes n} \ar[d]_{\varphi^{\otimes n}} \ar[r]_\alpha &
\mathcal{O}_S \ar[d]^1 \\
(\mathcal{L}')^{\otimes n} \ar[r]^{\alpha'} &
\mathcal{O}_S \\
}
$$
交换。因此有
\begin{equation}
\label{etale-cohomology-equation-isomorphisms-pairs}
\mathit{Isom}_S((\mathcal{L}, \alpha), (\mathcal{L}', \alpha'))
=
\left\{
\begin{matrix}
\emptyset & \text{若} & \text{二者不同构} \\
H^0(S, \mu_{n, S})\cdot \varphi & \text{若} &
\varphi \text{ 是序对同构}
\end{matrix}
\right.
\end{equation}
此外，给定两个序对 $(\mathcal{L}, \alpha)$、$(\mathcal{L}', \alpha')$，其张量积
$$
(\mathcal{L}, \alpha) \otimes (\mathcal{L}', \alpha')
=
(\mathcal{L} \otimes \mathcal{L}', \alpha \otimes \alpha')
$$
仍是一个序对。序对 $(\mathcal{O}_S, 1)$ 是该张量积运算的单位元，而逆元由
$$
(\mathcal{L}, \alpha)^{-1} = (\mathcal{L}^{\otimes -1}, \alpha^{\otimes -1}).
$$
给出。因此，序对的同构类集合构成阿贝尔群。注意，
$$
(\mathcal{L}, \alpha)^{\otimes n}
=
(\mathcal{L}^{\otimes n}, \alpha^{\otimes n})
\xrightarrow{\alpha}
(\mathcal{O}_S, 1)
$$
是同构，故该群中每个元素的阶都整除 $n$。提醒读者：一般而言，这个群{\bf 并非}
$n$-挠子群（位于 $\Pic(S)$ 中）。

\begin{lemma}
\label{etale-cohomology-lemma-describe-h1-mun}
设 $S$ 是概形。有典范等同
$$
H_\etale^1(S, \mu_n) =
\text{上述序对 }(\mathcal{L}, \alpha)\text{ 的同构类群}
$$
，条件是 $n$ 在 $S$ 上可逆。一般而言，有
$$
H_{fppf}^1(S, \mu_n) =
\text{上述序对 }(\mathcal{L}, \alpha)\text{ 的同构类群}.
$$
把 fppf 换成 syntomic 后，同一结果仍成立。
\end{lemma}

\begin{proof}
先证明第二个同构。设 $(\mathcal{L}, \alpha)$ 是上述序对。选取仿射开覆盖
$S = \bigcup U_i$，使得 $\mathcal{L}|_{U_i} \cong \mathcal{O}_{U_i}$。设
$s_i \in \mathcal{L}(U_i)$ 是生成元。于是
$\alpha(s_i^{\otimes n}) = f_i \in \mathcal{O}_S^*(U_i)$。写
$U_i = \Spec(A_i)$，则存在整体相对完全交
$A_i \to B_i = A_i[T]/(T^n - f_i)$，使 $f_i$ 的像是 $n$ 次幂（在 $B_i$ 中）。
换言之，令 $V_i = \Spec(B_i)$，便得到 syntomic 覆盖
$\mathcal{V} = \{V_i \to S\}_{i \in I}$ 和平凡化
$\varphi_i : (\mathcal{L}, \alpha)|_{V_i} \to (\mathcal{O}_{V_i}, 1)$。

\medskip\noindent
我们用这个结果（覆盖 $\mathcal{V}$ 的存在性）把
$H^1_{syntomic}(S, \mu_{n, S})$ 中的一个上同调类与该序对相联系。给出两种
（等价的）构造。

\medskip\noindent
第一种构造：使用 {\v C}ech 上同调。在二重交 $V_i \times_S V_j$ 上，有序对的同构
$$
(\mathcal{O}_{V_i \times_S V_j}, 1)
\xrightarrow{\text{pr}_0^*\varphi_i^{-1}}
(\mathcal{L}|_{V_i \times_S V_j}, \alpha|_{V_i \times_S V_j})
\xrightarrow{\text{pr}_1^*\varphi_j}
(\mathcal{O}_{V_i \times_S V_j}, 1)
$$
。由 (\ref{etale-cohomology-equation-isomorphisms-pairs})，它由元素
$\zeta_{ij} \in \mu_n(V_i \times_S V_j)$ 给出。略去如下验证：这些
$\zeta_{ij}$ 给出一个 $1$-上闭链，也就是说，给出元素
$(\zeta_{i_0i_1}) \in \check C(\mathcal{V}, \mu_n)$，并且
$d(\zeta_{i_0i_1}) = 0$。因此，它的类是
$\check H^1(\mathcal{V}, \mu_n)$ 中的元素，并且由定理
\ref{etale-cohomology-theorem-cech-ss}，它映到 $H^1_{syntomic}(S, \mu_{n, S})$ 中的上同调类。

\medskip\noindent
第二种构造：使用挠子。考虑预层
$$
\mu_n(\mathcal{L}, \alpha) :
U
\longmapsto
\mathit{Isom}_U((\mathcal{O}_U, 1), (\mathcal{L}, \alpha)|_U)
$$
，它位于 $(\Sch/S)_{syntomic}$ 上。可把它看成
$\SheafHom_\mathcal{O}(\mathcal{O}, \mathcal{L})$ 的子预层（这是内 Hom 层，
见 Modules on Sites，第 \ref{sites-modules-section-internal-hom} 节）。由于定义
该子预层的条件是局部的，它是层。由
(\ref{etale-cohomology-equation-isomorphisms-pairs})，这个层带有层 $\mu_{n, S}$ 的自由作用。
因此，只需验证它局部有截面。平凡化覆盖 $\mathcal{V}$ 的存在性说明了这一点。
所以 $\mu_n(\mathcal{L}, \alpha)$ 是 $\mu_{n, S}$-挠子；由
Cohomology on Sites，引理 \ref{sites-cohomology-lemma-torsors-h1}，得到
$H^1_{syntomic}(S, \mu_{n, S})$ 中相应的元素。

\medskip\noindent
现在还需证明下列各点：
\begin{enumerate}
\item 两种构造给出相同的上同调类。
\item 同构的序对给出相同的上同调类。
\item
$(\mathcal{L}, \alpha) \otimes (\mathcal{L}', \alpha')$
的上同调类，是 $(\mathcal{L}, \alpha)$ 与 $(\mathcal{L}', \alpha')$ 的上同调类之和。
\item 若上同调类平凡，则序对平凡。
\item $H^1_{syntomic}(S, \mu_{n, S})$ 的任意元素都是某个序对的上同调类。
\end{enumerate}
略去 (1) 的证明。(2) 由第二种构造显然成立，因为同构挠子给出相同的上同调类。
(3) 由第一种构造显然成立，因为所得 {\v C}ech 类相加。(4) 由第二种构造显然
成立，因为挠子平凡当且仅当它有整体截面，见 Cohomology on Sites，引理
\ref{sites-cohomology-lemma-trivial-torsor}。

\medskip\noindent
(5) 可如下看出（虽然直接证明会更好）。假设
$\xi \in H^1_{syntomic}(S, \mu_{n, S})$。则 $\xi$ 映到元素
$\overline{\xi} \in H^1_{syntomic}(S, \mathbf{G}_{m, S})$，且
$n \overline{\xi} = 0$。由定理 \ref{etale-cohomology-theorem-picard-group} 可知，
$\overline{\xi}$ 对应于可逆层 $\mathcal{L}$，其第 $n$ 次张量幂同构于
$\mathcal{O}_S$。因此，存在序对 $(\mathcal{L}, \alpha')$，其上同调类 $\xi'$
有同样的像 $\overline{\xi'}$，该像位于
$H^1_{syntomic}(S, \mathbf{G}_{m, S})$ 中。所以只需证明 $\xi - \xi'$ 是某个
序对的类。由构造及上面的上同调长正合列可知，
$\xi - \xi' = \partial(f)$，其中某个 $f \in H^0(S, \mathcal{O}_S^*)$。
考虑序对 $(\mathcal{O}_S, f)$。略去如下验证：
该序对的上同调类是 $\partial(f)$。这就完成第一个等同（把 fppf 换成 syntomic）
的证明。

\medskip\noindent
为看出第一个等同，注意：若 $n$ 在 $S$ 上可逆，则证明第一部分构造的覆盖
$\mathcal{V}$ 实际上是 \'etale 覆盖（比较引理
\ref{etale-cohomology-lemma-kummer-sequence} 的证明）。证明的其余部分与拓扑无关，只有最后一个
论证使用了 Kummer 序列的正合性，即使用了引理
\ref{etale-cohomology-lemma-kummer-sequence}。
\end{proof}






\section{邻域、茎与点}
\label{etale-cohomology-section-stalks}

\noindent
可为 $S$ 的每个几何点配上一个正合的茎函子。$S_\etale$ 上的层态射是同构，
当且仅当它在所有这些茎上都是同构。阿贝尔层复形正合，当且仅当其茎复形在
每个几何点处都正合。合起来说，这意味着概形 $S$ 的小 \'etale 网站有足够多的点。
而且，$S$ 的小 \'etale 拓扑斯的每个点（这是一个抽象概念）都由一个几何点给出。
因此，在某种意义上，可用几何点及其邻域来理解 $S$ 的小 \'etale 拓扑斯。

\begin{definition}
\label{etale-cohomology-definition-geometric-point}
设 $S$ 是概形。
\begin{enumerate}
\item $S$ 的一个{\it 几何点}是态射
$\Spec(k) \to S$，其中 $k$ 是代数闭域。这种点通常记作 $\overline{s}$，
即用带上划线的小写字母表示。我们常用 $\overline{s}$ 同时表示概形
$\Spec(k)$ 和该态射，并用 $\kappa(\overline{s})$ 表示 $k$。
\item 说 $\overline{s}${\it 位于} $s${\it 之上}，意指
$s \in S$ 是 $\overline{s}$ 的像。
\item 几何点 $\overline{s}$ 在 $S$ 中的一个{\it \'etale 邻域}是交换图
$$
\xymatrix{
& U \ar[d]^\varphi \\
{\overline{s}} \ar[r]^{\overline{s}} \ar[ur]^{\bar u} & S
}
$$
，其中 $\varphi$ 是概形的 \'etale 态射。把它记作
$(U, \overline{u}) \to (S, \overline{s})$。
\item {\it \'Etale 邻域的态射}
$(U, \overline{u}) \to (U', \overline{u}')$
是一个 $S$-态射 $h: U \to U'$，满足
$\overline{u}' = h \circ \overline{u}$。
\end{enumerate}
\end{definition}

\begin{remark}
\label{etale-cohomology-remark-etale-between-etale}
由于 $U$ 和 $U'$ 在 $S$ 上 \'etale，它们之间的任意 $S$-态射也都是
\'etale 的，见命题 \ref{etale-cohomology-proposition-etale-morphisms}。特别地，
\'etale 邻域的所有态射都是 \'etale 的。
\end{remark}

\begin{remark}
\label{etale-cohomology-remark-etale-neighbourhoods}
设 $S$ 是概形，$s \in S$ 是一点。在 More on Morphisms，定义
\ref{more-morphisms-definition-etale-neighbourhood} 中，我们定义了
\'etale 邻域 $(U, u) \to (S, s)$ 作为 $(S, s)$ 的邻域这一概念。若
$\overline{s}$ 是 $S$ 中位于 $s$ 之上的几何点，则任意 \'etale 邻域
$(U, \overline{u}) \to (S, \overline{s})$ 都给出 $(U, u)$，即
$(S, s)$ 的一个 \'etale 邻域：取 $u \in U$ 为 $U$ 中使
$\overline{u}$ 位于 $u$ 之上的唯一一点。反过来，给定
$(U, u)$，即 $(S, s)$ 的一个 \'etale 邻域，剩余域扩张
$\kappa(u)/\kappa(s)$ 有限可分（见命题
\ref{etale-cohomology-proposition-etale-morphisms}），因而可找到嵌入
$\kappa(u) \subset \kappa(\overline{s})$，它位于 $\kappa(s)$ 上。
换言之，可找到几何点 $\overline{u}$，它是 $U$ 中位于 $u$ 之上的点，
使得 $(U, \overline{u})$ 是 $(S, \overline{s})$ 的 \'etale 邻域。
我们将用这些观察在两类
\'etale 邻域之间转换。
\end{remark}

\begin{lemma}
\label{etale-cohomology-lemma-cofinal-etale}
设 $S$ 是概形，$\overline{s}$ 是 $S$ 的几何点。\'Etale 邻域的范畴是余滤过的。
更确切地说：
\begin{enumerate}
\item 设 $(U_i, \overline{u}_i)_{i = 1, 2}$ 是 $\overline{s}$ 在 $S$
中的两个 \'etale 邻域。则存在第三个 \'etale 邻域 $(U, \overline{u})$
以及态射
$(U, \overline{u}) \to (U_i, \overline{u}_i)$，$i = 1, 2$。
\item 设
$h_1, h_2: (U, \overline{u}) \to (U', \overline{u}')$
是 $\overline{s}$ 的 \'etale 邻域之间的两个态射。则存在 \'etale 邻域
$(U'', \overline{u}'')$ 和态射
$h : (U'', \overline{u}'') \to (U, \overline{u})$，使 $h_1$ 与 $h_2$
相等，即 $h_1 \circ h = h_2 \circ h$。
\end{enumerate}
\end{lemma}

\begin{proof}
对于 (1)，考虑纤维积 $U = U_1 \times_S U_2$。由于 \'etale 态射在基变换下
保持，见命题 \ref{etale-cohomology-proposition-etale-morphisms}，它在 $U_1$ 和 $U_2$ 上都
\'etale。映射 $\overline{s} \to U$（由
$(\overline{u}_1, \overline{u}_2)$ 定义）赋予它一个映到 $U_1$ 和
$U_2$ 的 \'etale 邻域结构。
对于 (2)，把 $U''$ 定义为纤维积
$$
\xymatrix{
U'' \ar[r] \ar[d] & U \ar[d]^{(h_1, h_2)} \\
U' \ar[r]^-\Delta & U' \times_S U'.
}
$$
由于 $\overline{u}$ 和 $\overline{u}'$ 在 $S$ 上都与 $\overline{s}$
相容，可知 $\overline{u}'' = (\overline{u}, \overline{u}')$ 是 $U''$
的几何点。特别地，$U'' \not = \emptyset$。此外，由于 $U'$ 在 $S$ 上
\'etale，纤维积 $U'\times_S U'$ 也如此（见命题
\ref{etale-cohomology-proposition-etale-morphisms}）。故由上面的注
\ref{etale-cohomology-remark-etale-between-etale}，竖直箭头 $(h_1, h_2)$ 是 \'etale 的。
于是经基变换，$U''$ 在 $U'$ 上 \'etale，因而也在 $S$ 上 \'etale
（因为 \'etale 态射的复合仍是 \'etale 的）。所以
$(U'', \overline{u}'')$ 给出了所求解答。
\end{proof}

\begin{lemma}
\label{etale-cohomology-lemma-geometric-lift-to-cover}
设 $S$ 是概形，$\overline{s}$ 是 $S$ 的几何点，
$(U, \overline{u})$ 是 $\overline{s}$ 的 \'etale 邻域，并设
$\mathcal{U} = \{\varphi_i : U_i \to U \}_{i\in I}$ 是 \'etale 覆盖。
则存在 $i \in I$ 和 $\overline{u}_i : \overline{s} \to U_i$，使得
$\varphi_i : (U_i, \overline{u}_i) \to (U, \overline{u})$
是 \'etale 邻域的态射。
\end{lemma}

\begin{proof}
因为 $U = \bigcup_{i\in I} \varphi_i(U_i)$，纤维积
$\overline{s} \times_{\overline{u}, U, \varphi_i} U_i$ 对某个 $i$ 非空。
考察笛卡尔图
$$
\xymatrix{
\overline{s} \times_{\overline{u}, U, \varphi_i} U_i
\ar[d]^{\text{pr}_1} \ar[r]_-{\text{pr}_2} & U_i
\ar[d]^{\varphi_i} \\
\Spec(k) = \overline{s} \ar@/^1pc/[u]^\sigma
\ar[r]^-{\overline{u}} & U
}
$$
。投影 $\text{pr}_1$ 是 \'etale 态射的基变换，故为 \'etale 的，见命题
\ref{etale-cohomology-proposition-etale-morphisms}。因此，由命题
\ref{etale-cohomology-proposition-etale-morphisms}，
$\overline{s} \times_{\overline{u}, U, \varphi_i} U_i$
是 $k$ 的有限可分扩张的谱的不交并；这里
$\overline{s} = \Spec(k)$。但 $k$ 代数闭，所以这些扩张全都平凡，
从而存在截面 $\sigma$，它是 $\text{pr}_1$ 的截面。复合
$\text{pr}_2 \circ \sigma$ 给出与 $\overline{u}$ 相容的映射。
\end{proof}

\begin{definition}
\label{etale-cohomology-definition-stalk}
设 $S$ 是概形，$\mathcal{F}$ 是 $S_\etale$ 上的预层，而
$\overline{s}$ 是 $S$ 的几何点。$\mathcal{F}$ 在 $\overline{s}$ 处的
{\it 茎}是
$$
\mathcal{F}_{\overline{s}}
=
\colim_{(U, \overline{u})} \mathcal{F}(U)
$$
，其中 $(U, \overline{u})$ 遍历 $\overline{s}$ 在 $S$ 中的所有
\'etale 邻域。
\end{definition}

\noindent
由引理 \ref{etale-cohomology-lemma-cofinal-etale}，该余极限取遍一个滤过指标范畴，即
\'etale 邻域范畴的反范畴。换言之，可把
$\mathcal{F}_{\overline{s}}$ 的元素看作三元组
$(U, \overline{u}, \sigma)$，其中 $\sigma \in \mathcal{F}(U)$。
两个三元组 $(U, \overline{u}, \sigma)$、
$(U', \overline{u}', \sigma')$ 定义茎中的同一元素，当且仅当存在第三个
\'etale 邻域 $(U'', \overline{u}'')$ 以及 \'etale 邻域态射
$h : (U'', \overline{u}'') \to (U, \overline{u})$、
$h' : (U'', \overline{u}'') \to (U', \overline{u}')$，使得
$h^*\sigma = (h')^*\sigma'$ 在 $\mathcal{F}(U'')$ 中成立。见
Categories，第 \ref{categories-section-directed-colimits} 节。

\begin{lemma}
\label{etale-cohomology-lemma-stalk-gives-point}
设 $S$ 是概形，$\overline{s}$ 是 $S$ 的几何点。考虑函子
\begin{align*}
u : S_\etale & \longrightarrow \textit{Sets}, \\
U & \longmapsto
|U_{\overline{s}}|
=
\{\overline{u} \text{，使得 }(U, \overline{u})
\text{ 是 }\overline{s}\text{ 的一个 \'etale 邻域}\}.
\end{align*}
这里 $|U_{\overline{s}}|$ 表示几何纤维的底集。则 $u$ 定义一个点
$p$，它是网站 $S_\etale$ 的点（Sites，定义
\ref{sites-definition-point}），
而与之相伴的茎函子 $\mathcal{F} \mapsto \mathcal{F}_p$
（Sites，等式 \ref{sites-equation-stalk}）就是上面定义的函子
$\mathcal{F} \mapsto \mathcal{F}_{\overline{s}}$。
\end{lemma}

\begin{proof}
在引理 \ref{etale-cohomology-lemma-geometric-lift-to-cover} 的证明中已经看到，概形
$U_{\overline{s}}$ 是与 $\overline{s}$ 同构的概形的不交并。因此，也可把
$|U_{\overline{s}}|$ 看作 $U$ 中位于 $\overline{s}$ 之上的几何点集合，
即所有能嵌入定义 \ref{etale-cohomology-definition-geometric-point} 中图式的态射
$\overline{u} : \overline{s} \to U$ 所成的集合。由此可知 $u(S)$ 是单点集，并且
$u(U \times_V W) = u(U) \times_{u(V)} u(W)$
，只要 $U \to V$ 和 $W \to V$ 是 $S_\etale$ 中的态射。又给定覆盖
$\{U_i \to U\}_{i \in I}$，它位于 $S_\etale$ 中，由引理
\ref{etale-cohomology-lemma-geometric-lift-to-cover} 可知
$\coprod u(U_i) \to u(U)$ 满射。因此可应用 Sites，命题
\ref{sites-proposition-point-limits}，故 $p$ 是网站 $S_\etale$ 的一个点。
最后，我们的函子 $\mathcal{F} \mapsto \mathcal{F}_{\overline{s}}$
与 Sites，等式 \ref{sites-equation-stalk} 中的函子
$\mathcal{F} \mapsto \mathcal{F}_p$（它和 $p$ 相伴）由完全相同的
余极限给出，这就证明了
最后一个断言。
\end{proof}

\begin{remark}
\label{etale-cohomology-remark-map-stalks}
设 $S$ 是概形，并设 $\overline{s} : \Spec(k) \to S$ 和
$\overline{s}' : \Spec(k') \to S$ 是 $S$ 的两个几何点。几何点的一个
{\it 态射 $a : \overline{s} \to \overline{s}'$} 就是一个态射
$a : \Spec(k) \to \Spec(k')$，它满足
$\overline{s}' \circ a = \overline{s}$。给定这样的态射，得到从
$\overline{s}'$ 的 \'etale 邻域范畴到 $\overline{s}$ 的 \'etale 邻域范畴的
函子，其规则为
$(U, \overline{u}') \mapsto (U, \overline{u}' \circ a)$。因此得到典范映射
$$
\mathcal{F}_{\overline{s}'}
=
\colim_{(U, \overline{u}')} \mathcal{F}(U)
\longrightarrow
\colim_{(U, \overline{u})} \mathcal{F}(U)
=
\mathcal{F}_{\overline{s}}
$$
，见 Categories，引理 \ref{categories-lemma-functorial-colimit}。利用把茎中元素
描述为三元组的方式，该映射把 $\mathcal{F}_{\overline{s}'}$ 中由三元组
$(U, \overline{u}', \sigma)$ 表示的元素，映到
$\mathcal{F}_{\overline{s}}$ 中由三元组
$(U, \overline{u}' \circ a, \sigma)$ 表示的元素。由于上述函子显然是等价，
可知这个典范映射是茎函子的同构。

\medskip\noindent
我们来确认，与 $a$ 对应的茎映射方向无误。注意，按 Sites，定义
\ref{sites-definition-morphism-points}，上述内容意味着 $a$ 定义态射
$a : p \to p'$；这是点 $p, p'$ 之间的态射，它们是网站
$S_\etale$ 的点，而这些点由
引理 \ref{etale-cohomology-lemma-stalk-gives-point} 分别与
$\overline{s}, \overline{s}'$ 相伴。还有更一般的点态射（对应于 $S$ 中点的
特殊化），我们稍后会描述；这些态射将不是同构，见第
\ref{etale-cohomology-section-specialization} 节。
\end{remark}

\begin{lemma}
\label{etale-cohomology-lemma-stalk-exact}
设 $S$ 是概形，$\overline{s}$ 是 $S$ 的几何点。
\begin{enumerate}
\item 茎函子
$\textit{PAb}(S_\etale) \to \textit{Ab}$,
$\mathcal{F}  \mapsto  \mathcal{F}_{\overline{s}}$
是正合的。
\item 有 $(\mathcal{F}^\#)_{\overline{s}} = \mathcal{F}_{\overline{s}}$，
其中 $\mathcal{F}$ 是 $S_\etale$ 上任意集合预层。
\item 函子
$\textit{Ab}(S_\etale) \to \textit{Ab}$,
$\mathcal{F} \mapsto \mathcal{F}_{\overline{s}}$ 是正合的。
\item 类似地，函子
$\textit{PSh}(S_\etale) \to \textit{Sets}$ 以及
$\Sh(S_\etale) \to \textit{Sets}$（由茎函子
$\mathcal{F} \mapsto \mathcal{F}_{\overline{s}}$ 给出）也是正合的（见
Categories，定义
\ref{categories-definition-exact}），并与任意余极限交换。
\end{enumerate}
\end{lemma}

\begin{proof}
在说明如何直接证明之前，先注意该结果可由 Modules on Sites，第
\ref{sites-modules-section-stalks} 节的一般材料得出。之所以如此，是因为
$\mathcal{F} \mapsto \mathcal{F}_{\overline{s}}$ 来自 $S$ 的小 \'etale
网站的一个点，见引理 \ref{etale-cohomology-lemma-stalk-gives-point}。我们只直接证明
(1)、(2) 和 (3)，略去 (4) 的直接证明。

\medskip\noindent
它作为 $\textit{PAb}(S_\etale)$ 上函子的正合性，形式地来自如下事实：
有向余极限与所有余极限及有限极限交换。(2) 中茎的等同由映射
$$
\kappa :
\mathcal{F}_{\overline{s}}
\longrightarrow
(\mathcal{F}^\#)_{\overline{s}}
$$
给出；该映射由自然态射 $\mathcal{F}\to \mathcal{F}^\#$ 诱导，见定理
\ref{etale-cohomology-theorem-sheafification}。我们断言这个映射是阿贝尔群的同构。
下面证明单射性，并略去满射性的证明。

\medskip\noindent
设 $\sigma\in \mathcal{F}_{\overline{s}}$。存在 \'etale 邻域
$(U, \overline{u})\to (S, \overline{s})$，使 $\sigma$ 是某个截面
$s \in \mathcal{F}(U)$ 的像。若 $\kappa(\sigma) = 0$ 在
$(\mathcal{F}^\#)_{\overline{s}}$ 中成立，则存在 \'etale 邻域态射
$(U', \overline{u}')\to (U, \overline{u})$，使 $s|_{U'}$ 在
$\mathcal{F}^\#(U')$ 中为零。由此存在 \'etale 覆盖
$\{U_i'\to U'\}_{i\in I}$，使得 $s|_{U_i'}=0$ 在
$\mathcal{F}(U_i')$ 中成立，对所有 $i$ 都如此。由引理
\ref{etale-cohomology-lemma-geometric-lift-to-cover}，
存在 $i \in I$ 和态射 $\overline{u}_i': \overline{s} \to U_i'$，使
$(U_i', \overline{u}_i') \to (U', \overline{u}')\to (U, \overline{u})$
是 \'etale 邻域的态射。于是 $\sigma = 0$，因为
$(U_i', \overline{u}_i') \to (U, \overline{u})$ 是 \'etale 邻域态射，
并且 $s|_{U'_i}=0$。这证明了 $\kappa$ 是单射。

\medskip\noindent
为证明函子 $\textit{Ab}(S_\etale) \to \textit{Ab}$ 正合，考虑
$\textit{Ab}(S_\etale)$ 中任意短正合列：
$
0\to \mathcal{F}\to \mathcal{G}\to \mathcal H \to 0.
$
它给出预层的正合列
$$
0 \to \mathcal{F} \to \mathcal{G} \to \mathcal H \to
\mathcal H/^p\mathcal{G} \to 0,
$$
，其中 $/^p$ 表示 $\textit{PAb}(S_\etale)$ 中的商。在
$\overline{s}$ 处取茎，可知
$(\mathcal H /^p\mathcal{G})_{\bar{s}} =
(\mathcal H /\mathcal{G})_{\bar{s}} = 0$，因为
$\mathcal H/^p\mathcal{G}$ 的层化是 $0$。因此，
$$
0\to \mathcal{F}_{\overline{s}} \to \mathcal{G}_{\overline{s}} \to
\mathcal{H}_{\overline{s}} \to 0 = (\mathcal H/^p\mathcal{G})_{\overline{s}}
$$
正合，因为取茎作为预层上的函子是正合的。
\end{proof}

\begin{theorem}
\label{etale-cohomology-theorem-exactness-stalks}
设 $S$ 是概形。集合层的态射 $a : \mathcal{F} \to \mathcal{G}$ 是单射
（相应地，满射），当且仅当茎上的映射
$a_{\overline{s}} : \mathcal{F}_{\overline{s}} \to \mathcal{G}_{\overline{s}}$
对 $S$ 的每个几何点都是单射（相应地，满射）。$S_\etale$ 上的阿贝尔层序列正合，当且仅当它在
$S$ 的几何点处的所有茎上都正合。
\end{theorem}

\begin{proof}
茎上正合的必要性由引理 \ref{etale-cohomology-lemma-stalk-exact} 得出。反过来，只需证明：
层态射满射（相应地，单射），当且仅当它在所有茎上满射（相应地，单射）。
我们证明满射情形，略去单射情形的证明。

\medskip\noindent
设 $\alpha : \mathcal{F} \to \mathcal{G}$ 是层态射，并且对所有几何点，
$\mathcal{F}_{\overline{s}} \to \mathcal{G}_{\overline{s}}$ 都满射。固定
$U \in \Ob(S_\etale)$ 和 $s \in \mathcal{G}(U)$。对每个 $u \in U$，
选取某个 $\overline{u} \to U$（它位于 $u$ 之上）以及 \'etale 邻域
$(V_u , \overline{v}_u) \to (U, \overline{u})$，使得
$s|_{V_u} = \alpha(s_{V_u})$ 对某个
$s_{V_u} \in \mathcal{F}(V_u)$ 成立。这是可能的，因为 $\alpha$ 在茎上满射。
于是 $\{V_u \to U\}_{u \in U}$ 是一个 \'etale 覆盖，在其上 $s$ 的限制属于
映射 $\alpha$ 的像。因此 $\alpha$ 满射，见 Sites，第
\ref{sites-section-sheaves-injective} 节。
\end{proof}

\begin{remarks}
\label{etale-cohomology-remarks-enough-points}
关于几何网站的点。
\begin{enumerate}
\item 定理 \ref{etale-cohomology-theorem-exactness-stalks} 说明：$S_\etale$ 的点族
（它由 $S$ 的几何点给出，见引理 \ref{etale-cohomology-lemma-stalk-gives-point}）是保守的，见
Sites，定义 \ref{sites-definition-enough-points}。特别地，$S_\etale$
有足够多的点。
\item 假设 $\mathcal{F}$ 是 $S$ 的大 \'etale 网站上的层
\label{etale-cohomology-item-stalks-big}
。设 $T \to S$ 是 $S$ 的大 \'etale 网站的对象，并设
$\overline{t}$ 是 $T$ 的几何点。把 $\mathcal{F}_{\overline{t}}$ 定义为
限制 $\mathcal{F}|_{T_\etale}$ 的茎；这是把 $\mathcal{F}$ 限制到
$T$ 的小 \'etale 网站所得的限制。换言之，可定义 $\mathcal{F}$ 在任意概形
$T/S \in \Ob((\Sch/S)_\etale)$ 的任意几何点处的茎。
\item 考虑该网站所有对象的所有几何点，可知 $S$ 的大 \'etale 网站也有足够多的点，
见 (\ref{etale-cohomology-item-stalks-big})。
\end{enumerate}
\end{remarks}

\noindent
初次阅读时应跳过下述引理。

\begin{lemma}
\label{etale-cohomology-lemma-points-small-etale-site}
设 $S$ 是概形。
\begin{enumerate}
\item 设 $p$ 是小 \'etale 网站
$S_\etale$ 的一个点；该网站属于概形 $S$，且该点由函子
$u : S_\etale \to \textit{Sets}$ 给出。则存在一个几何点 $\overline{s}$，
它是 $S$ 的几何点，使得 $p$ 与 $S_\etale$ 的如下点同构：引理
\ref{etale-cohomology-lemma-stalk-gives-point} 中与 $\overline{s}$ 相伴的点。
\item 设 $p : \Sh(pt) \to \Sh(S_\etale)$ 是 $S$ 的小 \'etale 拓扑斯的一个点。
则 $p$ 来自 $S$ 的某个几何点；也就是说，茎函子
$\mathcal{F} \mapsto \mathcal{F}_p$ 与定义
\ref{etale-cohomology-definition-stalk} 所定义的某个茎函子同构。
\end{enumerate}
\end{lemma}

\begin{proof}
由 Sites，引理 \ref{sites-lemma-point-site-topos}，网站的点与相伴拓扑斯的点
一一对应，故只需证明 (1)。由 Sites，命题
\ref{sites-proposition-point-limits}，函子 $u$ 具有下列性质：
(a) $u(S) = \{*\}$；(b) $u(U \times_V W) = u(U) \times_{u(V)} u(W)$；
并且 (c) 若 $\{U_i \to U\}$ 是一个 \'etale 覆盖，则
$\coprod u(U_i) \to u(U)$ 满射。
特别地，若 $U' \subset U$ 是开子概形，则 $u(U') \subset u(U)$。此外，由
Sites，引理 \ref{sites-lemma-point-site-topos}，可写成
$u(U) = p^{-1}(h_U^\#)$；换言之，$u(U)$ 是可表层 $h_U$ 的茎。若
$U = V \amalg W$，则 $h_U = (h_V \amalg h_W)^\#$，并且得到
$u(U) = u(V) \amalg u(W)$，因为 $p^{-1}$ 正合。

\medskip\noindent
考虑 $u$ 到 $S_{Zar}$ 的限制。由 Sites，例
\ref{sites-example-point-topological} 和
\ref{sites-example-point-topology}，存在唯一一点 $s \in S$，使得对开集
$S' \subset S$，$u(S') = \{*\}$ 当且仅当 $s \in S'$，而
$u(S') = \emptyset$ 当且仅当 $s \not \in S'$。注意，若
$\varphi : U \to S$ 是 $S_\etale$ 的对象，则
$\varphi(U) \subset S$ 是开集（见命题
\ref{etale-cohomology-proposition-etale-morphisms}），而
$\{U \to \varphi(U)\}$ 是 \'etale 覆盖。因此可知
$u(U) = \emptyset \Leftrightarrow s \not \in \varphi(U)$。

\medskip\noindent
取一个几何点 $\overline{s} : \overline{s} \to S$，它位于 $s$ 上方；关于记号的惯常滥用，
见定义 \ref{etale-cohomology-definition-geometric-point}。设 $\varphi : U \to S$ 是
$S_\etale$ 的对象，且 $U$ 仿射。注意，$\varphi$ 是分离的，而纤维
$U_s$（即 $\varphi$ 在 $s$ 上的纤维）是 $\Spec(\kappa(s))$ 上的仿射概形；
它是一个有限乘积之谱，其中各因子 $k_i$ 是 $\kappa(s)$ 的有限可分扩张。因此可以应用
\'Etale Morphisms，引理 \ref{etale-lemma-etale-etale-local-technical}，
得到一个 \'etale 邻域 $(V, \overline{v})$，它属于 $(S, \overline{s})$，并使得
$$
U \times_S V = U_1 \amalg \ldots \amalg U_n \amalg W
$$
其中 $U_i \to V$ 是同构，且 $W$ 没有位于 $\overline{v}$ 上方的点。因此可得
$$
u(U) \times u(V) =
u(U \times_S V) =
u(U_1) \amalg \ldots \amalg u(U_n) \amalg u(W)
$$
当然还有 $u(U_i) = u(V)$。将 $V$ 略微缩小后，可假定 $V$ 恰有一个位于
$s$ 上方的点，因而 $W$ 没有位于 $s$ 上方的点。由上述结论，这给出
$u(W) = \emptyset$。于是得到
$$
u(U) \times u(V) =
u(U_1) \amalg \ldots \amalg u(U_n) =
\coprod\nolimits_{i = 1, \ldots, n} u(V)
$$
其中各等号都与到集合 $u(V)$ 的映射相容。注意 $u(V) \not = \emptyset$，
因为 $s$ 属于 $V \to S$ 的像。特别地，在这种情形下 $u(U)$ 是一个含
$n$ 个元素的有限集：左边的纤维为 $u(U)$，它位于 $u(V)$ 的任一元素上方，
而右边的纤维基数为 $n$。

\medskip\noindent
考虑极限
$$
\lim_{(V, \overline{v})} u(V)
$$
这里极限取遍 \'etale 邻域 $(V, \overline{v})$ 所成的范畴；这些邻域属于
$\overline{s}$。
显然，只在 $(V, \overline{v})$ 所成、其中 $V$ 仿射的子范畴上取极限，所得值相同。
由上一段（交换 $V$ 与 $U$ 的角色）可知，在这种情形下 $u(V)$ 总是有限非空集。
此外，该极限是余滤过的，见引理 \ref{etale-cohomology-lemma-cofinal-etale}。因此由 Categories，第
\ref{categories-section-codirected-limits} 节，该极限非空。从中取一个元素 $x$。
这意味着我们得到一个 $x_{V, \overline{v}} \in u(V)$；这里
$(V, \overline{v})$ 可取为 $(S, \overline{s})$ 的任一 \'etale 邻域，并且对
\'etale 邻域的每个态射
$\varphi : (V', \overline{v}') \to (V, \overline{v})$，都有
$u(\varphi)(x_{V', \overline{v}'}) = x_{V, \overline{v}}$。

\medskip\noindent
我们将利用 $x$ 的选取来构造函子性双射
$$
c : |U_{\overline{s}}| \longrightarrow u(U)
$$
（其中 $U \in \Ob(S_\etale)$），这将完成证明。关于 $|U_{\overline{s}}|$ 的描述，
见引理 \ref{etale-cohomology-lemma-stalk-gives-point} 及其证明。首先声称，只需在 $U$ 仿射时构造该映射；
我们略去此断言的证明。设 $U \to S$ 是 $S_\etale$ 中的对象，且 $U$ 仿射；
并设 $\overline{u} : \overline{s} \to U$ 是 $|U_{\overline{s}}|$ 的一个元素。
选取 $(V, \overline{v})$，使 $U \times_S V$ 按证明第三段所述分解。
于是序对 $(\overline{u}, \overline{v})$ 给出 $U \times_S V$ 的一个几何点，
该点位于 $\overline{v}$ 上方，并确定一个分支 $U_i$，它是 $U \times_S V$ 的分支。
更确切地，存在一个截面 $\sigma : V \to U \times_S V$，它是投影
$\text{pr}_U$ 的截面，并使得
$(\overline{u}, \overline{v}) = \sigma \circ \overline{v}$。令
$c(\overline{u}) = u(\text{pr}_U)(u(\sigma)(x_{V, \overline{v}})) \in u(U)$。
需要验证它与 $(V, \overline{v})$ 的选取无关。由引理
\ref{etale-cohomology-lemma-cofinal-etale}，\'etale 邻域所成的范畴是余滤过的。因此只需证明：
给定 \'etale 邻域的态射
$\varphi : (V', \overline{v}') \to (V, \overline{v})$，并选取投影的截面
$\sigma' : V' \to U \times_S V'$，使得
$(\overline{u}, \overline{v'}) = \sigma' \circ \overline{v}'$，则有
$u(\sigma')(x_{V', \overline{v}'}) = u(\sigma)(x_{V, \overline{v}})$。
考虑图表
$$
\xymatrix{
V' \ar[d]^{\sigma'} \ar[r]_\varphi & V \ar[d]^\sigma \\
U \times_S V' \ar[r]^{1 \times \varphi} &
U \times_S V
}
$$
该图表未必交换。这是因为概形 $V'$ 和 $V$ 可能不连通，因而上面用来构造
$\sigma'$ 与 $\sigma$ 的分解可能不唯一。不过由构造确有
$\sigma \circ \varphi \circ \overline{v}' =
(1 \times \varphi) \circ \sigma' \circ \overline{v}'$
。因此，由于 $U \times_S V$ 在 $S$ 上是 \'etale 的，存在
$V'' \subset V'$；这是 $\overline{v'}$ 的一个开邻域，并使得把该图表限制到 $V''$
后它交换，见 Morphisms，引理 \ref{morphisms-lemma-value-at-one-point}。
这意味着可将上图扩充为
$$
\xymatrix{
V'' \ar[r] \ar[d]^{\sigma'|_{V''}} &
V' \ar[d]^{\sigma'} \ar[r]_\varphi &
V \ar[d]^\sigma \\
U \times_S V'' \ar[r] &
U \times_S V' \ar[r]^{1 \times \varphi} &
U \times_S V
}
$$
使得左方块与外围矩形都交换。由于 $u$ 是函子，这蕴含
$x_{V'', \overline{v}'}$ 沿图表中的任一条路径，都映到
$u(U \times_S V)$ 中的同一个元素。另一方面，它分别映到元素
$x_{V', \overline{v}'}$ 与 $x_{V, \overline{v}}$，它们分别属于
$u(V')$ 与 $u(V)$。这就蕴含所需等式
$u(\sigma')(x_{V', \overline{v}'}) = u(\sigma)(x_{V, \overline{v}})$。

\medskip\noindent
类似地可证明，构造 $c : |U_{\overline{s}}| \to u(U)$ 关于 $U$ 是函子性的；
细节略去。最后，由第三段的结果，映射 $c$ 显然为双射。这就完成引理的证明。
\end{proof}







\section{其他拓扑中的点}
\label{etale-cohomology-section-points-topologies}

\noindent
本节简要讨论概形的 \'etale 网站以外若干网站的点是否存在。关于本节所用术语，参见
Sites，第 \ref{sites-section-sites-enough-points} 节，以及
Topologies，第 \ref{topologies-section-procedure} 节及后续各节。
所有几何网站都有足够多的点。

\begin{lemma}
\label{etale-cohomology-lemma-points-fppf}
设 $S$ 是概形。下列所有网站都有足够多的点：
$S_{affine, Zar}$, $S_{Zar}$, $S_{affine, \etale}$, $S_\etale$,
$(\Sch/S)_{Zar}$, $(\textit{Aff}/S)_{Zar}$,
$(\Sch/S)_\etale$, $(\textit{Aff}/S)_\etale$,
$(\Sch/S)_{smooth}$, $(\textit{Aff}/S)_{smooth}$,
$(\Sch/S)_{syntomic}$, $(\textit{Aff}/S)_{syntomic}$,
$(\Sch/S)_{fppf}$，以及 $(\textit{Aff}/S)_{fppf}$。
\end{lemma}

\begin{proof}
对每个大网站，其相伴拓扑斯等价于网站 $(\textit{Aff}/S)_\tau$ 所定义的拓扑斯，见
Topologies，引理 \ref{topologies-lemma-affine-big-site-Zariski}、
\ref{topologies-lemma-affine-big-site-etale},
\ref{topologies-lemma-affine-big-site-smooth},
\ref{topologies-lemma-affine-big-site-syntomic} 和
\ref{topologies-lemma-affine-big-site-fppf}。网站 $(\textit{Aff}/S)_\tau$
的结论立即由 Deligne 的结果 Sites，引理 \ref{sites-lemma-criterion-points} 得出。

\medskip\noindent
关于 $S_{Zar}$ 的结论是显然的。关于 $S_{affine, Zar}$ 的结论由 Deligne 的结果得出。
关于 $S_\etale$ 的结论，或者由定理 \ref{etale-cohomology-theorem-exactness-stalks}（的证明）得出，
或者由 Topologies，引理 \ref{topologies-lemma-alternative} 以及将 Deligne 的结果应用于
$S_{affine, \etale}$ 得出。
\end{proof}

\noindent
上述引理保证点的存在性，却没有说明这些点的形貌。可以如下显式构造{\it 一些}点。
设 $\overline{s} : \Spec(k) \to S$ 是几何点，且 $k$ 代数闭。考虑函子
$$
u : (\Sch/S)_{fppf} \longrightarrow \textit{Sets},
\quad
u(U) = U(k) = \Mor_S(\Spec(k), U).
$$
注意 $U \mapsto U(k)$ 与有限极限交换，因为 $S(k) = \{\overline{s}\}$，且
$(U_1 \times_U U_2)(k) = U_1(k) \times_{U(k)} U_2(k)$.
此外，若 $\{U_i \to U\}$ 是 fppf 覆盖，则 $\coprod U_i(k) \to U(k)$ 满射。
由 Sites，命题 \ref{sites-proposition-point-limits}，可知 $u$ 定义一个点 $p$；
这是 $(\Sch/S)_{fppf}$ 的点，其茎为
$$
\mathcal{F}_p = \colim_{(U, x)} \mathcal{F}(U)
$$
其中照例对序对 $U \to S$、$x \in U(k)$ 取余极限。然而，这个范畴有一个始对象，
即 $(\Spec(k), \text{id})$，因此可知
$$
\mathcal{F}_p = \mathcal{F}(\Spec(k))
$$
，这并没有多大意思！事实上，一般而言，这些点不会组成保守点族。下述评注描述一类更有意思的点。

\begin{remark}
\label{etale-cohomology-remark-points-fppf-site}
\begin{reference}
这一点在 \cite{Schroeer} 中讨论。
\end{reference}
设 $S = \Spec(A)$ 是仿射概形。设 $(p, u)$ 是网站
$(\textit{Aff}/S)_{fppf}$ 的一个点，见 Sites，第 \ref{sites-section-points} 节和
第 \ref{sites-section-construct-points} 节。令 $B = \mathcal{O}_p$ 为结构层在点
$p$ 处的茎。回顾
$$
B = \colim_{(U, x)} \mathcal{O}(U) =
\colim_{(\Spec(C), x_C)} C
$$
，其中 $x_C \in u(\Spec(C))$。可能发生如下情形：$\Spec(B)$ 是
$(\textit{Aff}/S)_{fppf}$ 的对象，并且存在元素 $x_B \in u(\Spec(B))$，
它映到相容族 $x_C$。此时邻域系有始对象，因而有
$\mathcal{F}_p = \mathcal{F}(\Spec(B))$；其中 $\mathcal{F}$ 是
$(\textit{Aff}/S)_{fppf}$ 上的任意层。不难看出，若
$\mathcal{F} \mapsto \mathcal{F}(\Spec(B))$ 定义
$\Sh((\textit{Aff}/S)_{fppf})$ 的一个点，则 $B$ 必须是局部 $A$-代数，
使得每个忠实平坦、有限表示的环映射 $B \to B'$ 都有一个截面 $B' \to B$。
反之，对任意这样的 $A$-代数 $B$，函子
$\mathcal{F} \mapsto \mathcal{F}(\Spec(B))$ 是某个点的茎函子。细节略去。
网站 $(\textit{Aff}/S)_{fppf}$ 的一般点是什么样的，尚不清楚。
\end{remark}











\section{阿贝尔层的支集}
\label{etale-cohomology-section-support}

\noindent
先讨论局部截面的支集。

\begin{lemma}
\label{etale-cohomology-lemma-support-subsheaf-final}
设 $S$ 是概形。设 $\mathcal{F}$ 是 $S$ 的 \'etale 拓扑斯之终对象的子层（见
Sites，例 \ref{sites-example-singleton-sheaf}）。则存在唯一开集
$W \subset S$，使得 $\mathcal{F} = h_W$。
\end{lemma}

\begin{proof}
该条件意味着 $\mathcal{F}(U)$ 是单点集或空集；这里
$\varphi : U \to S$ 遍历 $\Ob(S_\etale)$。特别地，局部截面总能黏合。若
$\mathcal{F}(U) \not = \emptyset$，则 $\mathcal{F}(\varphi(U)) \not = \emptyset$，
因为 $\{\varphi : U \to \varphi(U)\}$ 是覆盖。因此可取
$W = \bigcup_{\varphi : U \to S, \mathcal{F}(U) \not = \emptyset} \varphi(U)$.
\end{proof}

\begin{lemma}
\label{etale-cohomology-lemma-zero-over-image}
设 $S$ 是概形。设 $\mathcal{F}$ 是 $S_\etale$ 上的阿贝尔层。
设 $\sigma \in \mathcal{F}(U)$ 是局部截面。存在开子集 $W \subset U$，使得
\begin{enumerate}
\item $W \subset U$ 是 $U$ 中使 $\sigma|_W = 0$ 的最大 Zariski 开子集；
\item 对每个 $\varphi : V \to U$（它属于 $S_\etale$），有
$$
\sigma|_V = 0 \Leftrightarrow \varphi(V) \subset W,
$$
\item 对每个几何点 $\overline{u}$（它是 $U$ 的几何点），有
$$
(U, \overline{u}, \sigma) = 0\text{ 于 }\mathcal{F}_{\overline{s}}
\Leftrightarrow
\overline{u} \in W
$$
，其中 $\overline{s} = (U \to S) \circ \overline{u}$。
\end{enumerate}
\end{lemma}

\begin{proof}
由于 $\mathcal{F}$ 是 \'etale 拓扑中的层，$\mathcal{F}$ 在 $U_{Zar}$ 上的限制是 $U$ 的
Zariski 拓扑中的层。因此存在满足 (1) 的 Zariski 开集 $W$，见 Modules，引理
\ref{modules-lemma-support-section-closed}。设 $\varphi : V \to U$ 是
$S_\etale$ 中的箭头。注意，$\varphi(V) \subset U$ 是开子集，并且
$\{V \to \varphi(V)\}$ 是 \'etale 覆盖。因此，若 $\sigma|_V = 0$，则由
$\mathcal{F}$ 的层条件可知 $\sigma|_{\varphi(V)} = 0$。这证明了 (2)。
为证明 (3)，需要证明：若 $(U, \overline{u}, \sigma)$ 定义
$\mathcal{F}_{\overline{s}}$ 的零元素，则 $\overline{u} \in W$。这是因为该假设意味着
存在 \'etale 邻域的一个态射 $(V, \overline{v}) \to (U, \overline{u})$，使得
$\sigma|_V = 0$。故由 (2) 可知 $V \to U$ 映入 $W$，从而
$\overline{u} \in W$。
\end{proof}

\noindent
设 $S$ 是概形，并设 $s \in S$。设 $\mathcal{F}$ 是 $S_\etale$ 上的层。
由评注 \ref{etale-cohomology-remark-map-stalks}，层 $\mathcal{F}$ 在位于 $s$ 上方的几何点处之茎的
同构类是良定义的。

\begin{definition}
\label{etale-cohomology-definition-support}
设 $S$ 是概形。设 $\mathcal{F}$ 是 $S_\etale$ 上的阿贝尔层。
\begin{enumerate}
\item {\it $\mathcal{F}$ 的支集}是所有点 $s \in S$ 组成的集合；它们满足
$\mathcal{F}_{\overline{s}} \not = 0$，其中 $\overline{s}$ 是位于 $s$ 上方的
任一（等价地，某个）几何点。
\item 设 $\sigma \in \mathcal{F}(U)$ 是截面。{\it $\sigma$ 的支集}是闭子集
$U \setminus W$，其中 $W \subset U$ 是 $U$ 中使 $\sigma$ 限制为零的最大开子集
（见引理 \ref{etale-cohomology-lemma-zero-over-image}）。
\end{enumerate}
\end{definition}

\noindent
一般而言，阿贝尔层的支集不是闭集。例如，设 $S = \mathbf{A}^1_{\mathbf{C}}$。
令 $i_t : \Spec(\mathbf{C}) \to S$ 是点 $t \in \mathbf{C}$ 的嵌入。
稍后将看到，$\mathcal{F}_t = i_{t, *}(\mathbf{Z}/2\mathbf{Z})$ 是支集恰为
$\{t\}$ 的阿贝尔层，见第 \ref{etale-cohomology-section-closed-immersions} 节。于是
$$
\bigoplus\nolimits_{n \in \mathbf{N}} \mathcal{F}_n
$$
是支集为 $\{1, 2, 3, \ldots\} \subset S$ 的阿贝尔层。这是因为取茎与余极限交换，
见引理 \ref{etale-cohomology-lemma-stalk-exact}。这就给出了支集不闭的阿贝尔层之例。
下面是关于层与截面的支集的一些基本事实。

\begin{lemma}
\label{etale-cohomology-lemma-support-section-closed}
设 $S$ 是概形。设 $\mathcal{F}$ 是 $S_\etale$ 上的阿贝尔层。
设 $U \in \Ob(S_\etale)$ 且 $\sigma \in \mathcal{F}(U)$。
\begin{enumerate}
\item $\sigma$ 的支集在 $U$ 中是闭的。
\item $\sigma + \sigma'$ 的支集包含于 $\sigma, \sigma' \in \mathcal{F}(U)$
二者支集的并。
\item 若 $\varphi : \mathcal{F} \to \mathcal{G}$ 是 $S_\etale$ 上阿贝尔层的映射，
则 $\varphi(\sigma)$ 的支集包含于 $\sigma \in \mathcal{F}(U)$ 的支集。
\item $\mathcal{F}$ 的支集是 $\mathcal{F}$ 的所有局部截面之支集的像的并。
\item 若 $\mathcal{F} \to \mathcal{G}$ 满射，则 $\mathcal{G}$ 的支集是
$\mathcal{F}$ 的支集的子集。
\item 若 $\mathcal{F} \to \mathcal{G}$ 单射，则 $\mathcal{F}$ 的支集是
$\mathcal{G}$ 的支集的子集。
\end{enumerate}
\end{lemma}

\begin{proof}
(1) 由定义成立。(2) 与 (3) 成立，因为对 $\mathcal{F}$ 与 $\mathcal{G}$ 在
$U_{Zar}$ 上的限制，相应结论成立，见 Modules，引理
\ref{modules-lemma-support-section-closed}。(4) 是引理
\ref{etale-cohomology-lemma-zero-over-image} 的 (3) 的直接推论。(5) 与 (6) 由其他各项得出。
\end{proof}

\begin{lemma}
\label{etale-cohomology-lemma-support-sheaf-rings-closed}
$S_\etale$ 上环层的支集是闭集。
\end{lemma}

\begin{proof}
这是因为（按照我们的约定）环是 $0$ 当且仅当 $1 = 0$，所以环层的支集就是单位截面的支集。
\end{proof}




\section{Hensel 环}
\label{etale-cohomology-section-henselian-ring}

\noindent
先陈述一个在 Stacks project 中已经多次使用的定理。该结果有许多版本；这里仅陈述其代数版本。

\begin{theorem}
\label{etale-cohomology-theorem-quasi-finite-etale-locally}
设 $A\to B$ 是有限型环映射，且 $\mathfrak p \subset A$ 是素理想。则存在
\'etale 环映射 $A \to A'$ 和素理想 $\mathfrak p' \subset A'$；后者位于
$\mathfrak p$ 上方，并且
\begin{enumerate}
\item
$\kappa(\mathfrak p) = \kappa(\mathfrak p')$,
\item
$ B \otimes_A A' = B_1\times \ldots \times B_r \times C$,
\item
$ A'\to B_i$ 是有限的，并且存在唯一素理想 $q_i\subset B_i$ 位于
$\mathfrak p'$ 上方；以及
\item 纤维 $\Spec(C \otimes_{A'} \kappa(\mathfrak p'))$ 的所有不可约分支的维数都至少为 1；
这是 $C$ 在 $\mathfrak p'$ 上的纤维。
\end{enumerate}
\end{theorem}

\begin{proof}
见 Algebra，引理 \ref{algebra-lemma-etale-makes-quasi-finite-finite}，或见
\cite[Th\'eor\`eme 18.12.1]{EGA4}。关于用概形态射表述的一系列版本，见
More on Morphisms，第 \ref{more-morphisms-section-etale-localization} 节。
\end{proof}

\noindent
回顾 Hensel 引理。该引理有许多版本。下面列出两个：
\begin{enumerate}
\item[(f)] 若 $f\in \mathbf{Z}_p[T]$ 首一，且
$f \bmod p = g_0 h_0$，其中 $gcd(g_0, h_0) = 1$，则 $f$ 可分解为
$f = gh$，其中 $\bar g = g_0$ 且 $\bar h = h_0$；
\item[(r)] 若 $f \in \mathbf{Z}_p[T]$ 首一，且 $a_0 \in \mathbf{F}_p$ 满足
$\bar f(a_0) =0$ 而 $\bar f'(a_0) \neq 0$，则存在 $a \in \mathbf{Z}_p$，使得
$f(a) = 0$ 且 $\bar a = a_0$。
\end{enumerate}
两个版本都成立（稍后会看到这一点）。第一个版本要求把互素因子分解提升，
第二个版本要求把模极大理想的单根提升。事实证明，对一般局部环要求这两个条件是等价的，
并且还等价于许多其他条件。我们用根提升性质来定义 Hensel 局部环，因为它通常最容易检验。

%10.01.09
\begin{definition}
\label{etale-cohomology-definition-henselian}
（见 Algebra，定义 \ref{algebra-definition-henselian}。）称局部环
$(R, \mathfrak m, \kappa)$ 是{\it Hensel 的}，如果对所有首一
$f \in R[T]$ 以及所有 $a_0 \in \kappa$，只要
$\bar f(a_0) = 0$ 且 $\bar f'(a_0) \neq 0$，就存在 $a \in R$，使得
$f(a) = 0$ 且 $a \bmod \mathfrak m = a_0$。
\end{definition}

\noindent
应当记住的 Hensel 局部环的一个好例子是完备局部环。回顾（Algebra，定义
\ref{algebra-definition-complete-local-ring}），完备局部环是局部环
$(R, \mathfrak m)$，它满足 $R \cong \lim_n R/\mathfrak m^n$；换言之，
它关于 $\mathfrak m$-进拓扑是完备且分离的。

\begin{theorem}
\label{etale-cohomology-theorem-hensel}
完备局部环是 Hensel 的。
\end{theorem}

\begin{proof}
使用 Newton 方法。见 Algebra，引理 \ref{algebra-lemma-complete-henselian}。
\end{proof}

\begin{theorem}
\label{etale-cohomology-theorem-henselian}
设 $(R, \mathfrak m, \kappa)$ 是局部环。下列条件等价：
\begin{enumerate}
\item $R$ 是 Hensel 的；
\item 对任意 $f\in R[T]$ 以及任意分解 $\bar f = g_0 h_0$（它在
$\kappa[T]$ 中且满足 $\gcd(g_0, h_0)=1$），存在分解 $f = gh$（它在
$R[T]$ 中），并满足 $\bar g = g_0$ 且 $\bar h = h_0$；
\item 任意有限 $R$-代数 $S$ 都同构于若干在 $R$ 上有限的局部环的有限乘积；
\item 任意有限型 $R$-代数 $A$ 都同构于一个乘积
$A \cong A' \times C$，其中 $A' \cong A_1 \times \ldots \times A_r$
是有限局部 $R$-代数的乘积，并且 $C \otimes_R \kappa$ 的所有不可约分支的维数至少为 1；
\item 若 $A$ 是 \'etale $R$-代数，且 $\mathfrak n$ 是 $A$ 的一个位于
$\mathfrak m$ 上方的极大理想，满足 $\kappa \cong A/\mathfrak n$，则存在同构
$\varphi : A \cong R \times A'$，使得
$\varphi(\mathfrak n) = \mathfrak m \times A' \subset R \times A'$。
\end{enumerate}
\end{theorem}

\begin{proof}
这只是 Algebra，引理 \ref{algebra-lemma-characterize-henselian} 中若干结果的子集。
注意，上面的 (5) 对应于 Algebra，引理
\ref{algebra-lemma-characterize-henselian} 的 (8)，但表述略有不同。
\end{proof}

\begin{lemma}
\label{etale-cohomology-lemma-finite-over-henselian}
若 $R$ 是 Hensel 的，且 $A$ 是有限 $R$-代数，则 $A$ 是 Hensel 局部环的有限乘积。
\end{lemma}

\begin{proof}
见 Algebra，引理 \ref{algebra-lemma-finite-over-henselian}。
\end{proof}

\begin{definition}
\label{etale-cohomology-definition-strictly-henselian}
若局部环 $R$ 是 Hensel 的，且其剩余域可分闭，则称其为{\it 严格 Hensel 的}。
\end{definition}

\begin{example}
\label{etale-cohomology-example-powerseries}
在 $R = \mathbf{C}[[t]]$ 的情形，\'etale $R$-代数是平凡扩张 $R \to R$ 与扩张
$R \to R[X, X^{-1}]/(X^n-t)$ 的有限乘积。后一类扩张经开集
$D(t) \subset \Spec(R)$ 分解，故任意 \'etale 覆盖都可由覆盖
$\{\text{id} : \Spec(R) \to \Spec(R)\}$ 加细。下面将看到，对 Hensel 环之谱的
\'etale 覆盖而言，这是一个颇为一般的事实。由此将说明，严格 Hensel 环之谱的高阶
\'etale 上同调为零。
\end{example}

\begin{theorem}
\label{etale-cohomology-theorem-henselization}
设 $(R, \mathfrak m, \kappa)$ 是局部环，且
$\kappa\subset\kappa^{sep}$ 是可分代数闭包。存在典范平坦局部环映射
$R \to R^h \to R^{sh}$，其中
\begin{enumerate}
\item $R^h$、$R^{sh}$ 是 \'etale $R$-代数的滤过余极限；
\item $R^h$ 是 Hensel 的，$R^{sh}$ 是严格 Hensel 的；
\item $\mathfrak m R^h$（相应地 $\mathfrak m R^{sh}$）是 $R^h$（相应地
$R^{sh}$）的极大理想；以及
\item $\kappa = R^h/\mathfrak m R^h$，且
$\kappa^{sep} = R^{sh}/\mathfrak m R^{sh}$；这些等式按 $\kappa$ 的扩张来理解。
\end{enumerate}
\end{theorem}

\begin{proof}
$R^h$ 与 $R^{sh}$ 的结构见 Algebra，引理 \ref{algebra-lemma-henselization} 和
\ref{algebra-lemma-strict-henselization}。
\end{proof}

\noindent
定理 \ref{etale-cohomology-theorem-henselization} 中构造的环分别称为局部环 $R$ 的
{\it Hensel 化}和{\it 严格 Hensel 化}，见 Algebra，定义
\ref{algebra-definition-henselization}。$R$ 的许多性质会反映在其（严格）Hensel 化中，
见 More on Algebra，第 \ref{more-algebra-section-permanence-henselization} 节。




\section{结构层的茎}
\label{etale-cohomology-section-stalks-structure-sheaf}

\noindent
本节把结构层在几何点处的茎与相应“通常”点处局部环的严格 Hensel 化等同起来。

\begin{lemma}
\label{etale-cohomology-lemma-describe-etale-local-ring}
\begin{slogan}
概形结构层在 etale 拓扑中的茎是严格 Hensel 化。
\end{slogan}
设 $S$ 是概形。设 $\overline{s}$ 是 $S$ 的几何点，位于 $s \in S$ 上方。
令 $\kappa = \kappa(s)$，并令
$\kappa \subset \kappa^{sep} \subset \kappa(\overline{s})$ 表示
$\kappa$ 在 $\kappa(\overline{s})$ 中的可分代数闭包。则有典范等同
$$
(\mathcal{O}_{S, s})^{sh}
\cong
(\mathcal{O}_S)_{\overline{s}}
$$
，其中左边是定理 \ref{etale-cohomology-theorem-henselization} 所述局部环
$\mathcal{O}_{S, s}$ 的严格 Hensel 化，右边是结构层 $\mathcal{O}_S$
（定义在 $S_\etale$ 上）在几何点 $\overline{s}$ 处的茎。
\end{lemma}

\begin{proof}
设 $\Spec(A) \subset S$ 是 $s$ 的仿射邻域。设 $\mathfrak p \subset A$
是对应于 $s$ 的素理想。作出这些选择后，有典范同构
$\mathcal{O}_{S, s} = A_{\mathfrak p}$ 和 $\kappa(s) = \kappa(\mathfrak p)$。
因而有
$\kappa(\mathfrak p) \subset \kappa^{sep} \subset \kappa(\overline{s})$.
回顾
$$
(\mathcal{O}_S)_{\overline{s}} =
\colim_{(U, \overline{u})} \mathcal{O}(U)
$$
，其中极限取遍 $(S, \overline{s})$ 的 \'etale 邻域。一个余共尾系由如下
\'etale 邻域 $(U, \overline{u})$ 给出：$U$ 仿射，且 $U \to S$ 经
$\Spec(A)$ 分解。换言之，可知
$$
(\mathcal{O}_S)_{\overline{s}} = \colim_{(B, \mathfrak q, \phi)} B
$$
，其中余极限取遍 \'etale $A$-代数 $B$；它们配备素理想 $\mathfrak q$，
该理想位于 $\mathfrak p$ 上方，并配备一个 $\kappa(\mathfrak p)$-代数映射
$\phi : \kappa(\mathfrak q) \to \kappa(\overline{s})$。注意，由于
$\kappa(\mathfrak q)$ 在 $\kappa(\mathfrak p)$ 上有限可分，$\phi$ 的像包含于
$\kappa^{sep}$。经这些转译，本引理的结论等价于 Algebra，引理
\ref{algebra-lemma-strict-henselization-different} 的结论。
\end{proof}

\begin{definition}
\label{etale-cohomology-definition-etale-local-rings}
设 $S$ 是概形。设 $\overline{s}$ 是 $S$ 的几何点，位于点 $s \in S$ 上方。
\begin{enumerate}
\item {\it $S$ 在 $\overline{s}$ 处的 \'etale 局部环}是结构层 $\mathcal{O}_S$
（定义在 $S_\etale$ 上）在 $\overline{s}$ 处的茎。有时称它为
{\it $\mathcal{O}_{S, s}$ 的严格 Hensel 化}，这里是相对于几何点
$\overline{s}$ 而言。
所用记号为 $\mathcal{O}_{S, \overline{s}}^{sh}$。
\item {\it $\mathcal{O}_{S, s}$ 的 Hensel 化}是 $S$ 在 $s$ 处的局部环的 Hensel 化。
见 Algebra，定义 \ref{algebra-definition-henselization}，以及定理
\ref{etale-cohomology-theorem-henselization}。记为 $\mathcal{O}_{S, s}^h$。
\item {\it $S$ 在 $\overline{s}$ 处的严格 Hensel 化}是概形
$\Spec(\mathcal{O}_{S, \overline{s}}^{sh})$。
\item {\it $S$ 在 $s$ 处的 Hensel 化}是概形 $\Spec(\mathcal{O}_{S, s}^h)$。
\end{enumerate}
\end{definition}

\noindent
设 $f : T \to S$ 是概形态射。设 $\overline{t}$ 是 $T$ 的几何点，其像为
$\overline{s}$，后者位于 $S$ 中。设 $t \in T$ 和 $s \in S$ 为它们的像。则得到典范交换图
$$
\xymatrix{
\Spec(\mathcal{O}^{sh}_{T, \overline{t}}) \ar[r] \ar[d] &
\Spec(\mathcal{O}^h_{T, t}) \ar[r] \ar[d] &
T \ar[d]^f \\
\Spec(\mathcal{O}^{sh}_{S, \overline{s}}) \ar[r] &
\Spec(\mathcal{O}^h_{S, s}) \ar[r] &
S
}
$$
；这是 $T$ 与 $S$ 的 Hensel 化和严格 Hensel 化所成的图。可通过选取 $t$ 和 $s$
的仿射邻域，并使用 Algebra，引理
\ref{algebra-lemma-henselian-functorial-improve} 和
\ref{algebra-lemma-strictly-henselian-functorial-improve} 给出的（严格）Hensel 化的函子性来证明。

\begin{lemma}
\label{etale-cohomology-lemma-describe-henselization}
设 $S$ 是概形。设 $s \in S$。则有
$$
\mathcal{O}_{S, s}^h =
\colim_{(U, u)} \mathcal{O}(U)
$$
，其中余极限取遍滤过范畴中的 \'etale 邻域 $(U, u)$；它们是
$(S, s)$ 的邻域，并满足 $\kappa(s) = \kappa(u)$。
\end{lemma}

\begin{proof}
本引理是 More on Morphisms，引理
\ref{more-morphisms-lemma-describe-henselization} 的一个副本。
\end{proof}

\begin{remark}
\label{etale-cohomology-remark-henselization-Noetherian}
设 $S$ 是概形。设 $s \in S$。若 $S$ 局部 Noether，则
$\mathcal{O}_{S, s}^h$ 也 Noether，并且二者有相同的完备化：
$$
\widehat{\mathcal{O}_{S, s}} \cong \widehat{\mathcal{O}_{S, s}^h}.
$$
特别地，
$\mathcal{O}_{S, s} \subset
\mathcal{O}_{S, s}^h \subset
\widehat{\mathcal{O}_{S, s}}$.
$\mathcal{O}_{S, s}$ 的 Hensel 化一般远小于它的完备化，并继承它的许多性质。
例如，若 $\mathcal{O}_{S, s}$ 既约，则 $\mathcal{O}_{S, s}^h$ 也既约；
但对完备化而言，一般并非如此。此处待补充未来的参考文献。
\end{remark}

\begin{lemma}
\label{etale-cohomology-lemma-etale-site-locally-ringed}
设 $S$ 是概形。小 \'etale 网站 $S_\etale$ 配备其结构层 $\mathcal{O}_S$ 后，
是局部环化网站，见 Modules on Sites，定义
\ref{sites-modules-definition-locally-ringed}。
\end{lemma}

\begin{proof}
这是因为各茎 $(\mathcal{O}_S)_{\overline{s}} = \mathcal{O}^{sh}_{S, \overline{s}}$
都是局部环，且 $S_\etale$ 有足够多的点；见引理
\ref{etale-cohomology-lemma-describe-etale-local-ring}、定理 \ref{etale-cohomology-theorem-exactness-stalks} 和评注
\ref{etale-cohomology-remarks-enough-points}。关于这些事实蕴含小 \'etale 网站是局部环化网站，见
Modules on Sites，引理 \ref{sites-modules-lemma-locally-ringed-stalk} 和
\ref{sites-modules-lemma-ringed-stalk-not-zero}。
\end{proof}






\section{小 \'etale 拓扑斯的函子性}
\label{etale-cohomology-section-functoriality}

\noindent
到目前为止，我们尚未讨论 \'etale 网站的函子性，换言之，给定概形态射时会发生什么。
精确的形式化讨论见 Topologies，第 \ref{topologies-section-etale} 节。
本节和随后几节将在小 \'etale 网站的框架下简要讨论这些内容。

\medskip\noindent
设 $f : X \to Y$ 是概形态射。得到函子
\begin{equation}
\label{etale-cohomology-equation-functorial}
u : Y_\etale \longrightarrow X_\etale, \quad
V/Y \longmapsto X \times_Y V/X.
\end{equation}
该函子具有下列重要性质：
\begin{enumerate}
\item $u(\text{终对象}) = \text{终对象}$；
\item $u$ 保持纤维积；
\item 若 $\{V_j \to V\}$ 是 $Y_\etale$ 中的覆盖，则
$\{u(V_j) \to u(V)\}$ 是 $X_\etale$ 中的覆盖。
\end{enumerate}
每一项都容易检验（略去）。因此得到所谓的{\it 网站态射}
$$
f_{small} : X_\etale \longrightarrow Y_\etale,
$$
，见 Sites，定义 \ref{sites-definition-morphism-sites} 和 Sites，命题
\ref{sites-proposition-get-morphism}。处理 \'etale 层与 \'etale 上同调时，
不必详细掌握这个抽象概念。通常只需知道 \'etale 层上存在函子
$f_{small, *}$（推前）和 $f_{small}^{-1}$（拉回），并了解它们的一些简单性质。
随后几节将讨论这些性质；不过有时我们会引用较抽象的材料来作证明，因为那里往往是证明它们的自然框架。


\section{正像}
\label{etale-cohomology-section-direct-image}

\noindent
定义预层的推前。

\begin{definition}
\label{etale-cohomology-definition-direct-image-presheaf}
设 $f: X\to Y$ 是概形态射。设 $\mathcal{F} $ 是 $X_\etale$ 上的集合预层。
$\mathcal{F}$（在 $f$ 下）的{\it 正像}，或称{\it 推前}，是
$$
f_*\mathcal{F} : Y_\etale^{opp} \longrightarrow \textit{Sets}, \quad
(V/Y) \longmapsto \mathcal{F}(X \times_Y V/X).
$$
有时写作 $f_* = f_{small, *}$，以区别于其他正像函子（例如通常的 Zariski 推前或
$f_{big, *}$）。
\end{definition}

\noindent
这是良定义的 \'etale 预层，因为 \'etale 态射的基变换仍是 \'etale 的。
更范畴化地说，$f_*\mathcal{F}$ 是函子的复合 $\mathcal{F} \circ u$，其中
$u$ 如方程 (\ref{etale-cohomology-equation-functorial}) 所示。这清楚地说明，该构造关于预层
$\mathcal{F}$ 是函子性的，故得到函子
$$
f_* = f_{small, *} :
\textit{PSh}(X_\etale)
\longrightarrow
\textit{PSh}(Y_\etale)
$$
。注意，若 $\mathcal{F}$ 是阿贝尔群预层，则 $f_*\mathcal{F}$ 也是阿贝尔群预层，
并得到
$$
f_* = f_{small, *} :
\textit{PAb}(X_\etale)
\longrightarrow
\textit{PAb}(Y_\etale)
$$
，与前面相同（即由完全相同的规则定义）。

\begin{remark}
\label{etale-cohomology-remark-direct-image-sheaf}
我们断言，层的正像仍是层。事实上，若 $\{V_j \to V\}$ 是 $Y_\etale$ 中的
\'etale 覆盖，则 $\{X \times_Y V_j \to X \times_Y V\}$ 是
$X_\etale$ 中的 \'etale 覆盖。因此，$\mathcal{F}$ 关于
$\{X \times_Y V_i \to X \times_Y V\}$ 的层条件，等价于
$f_*\mathcal{F}$ 关于 $\{V_i \to V\}$ 的层条件。故若 $\mathcal{F}$ 是层，
则 $f_*\mathcal{F}$ 也是层。
\end{remark}

\begin{definition}
\label{etale-cohomology-definition-direct-image-sheaf}
设 $f: X\to Y$ 是概形态射。设 $\mathcal{F} $ 是 $X_\etale$ 上的集合层。
$\mathcal{F}$（在 $f$ 下）的{\it 正像}，或称{\it 推前}，是
$$
f_*\mathcal{F} : Y_\etale^{opp} \longrightarrow \textit{Sets}, \quad
(V/Y) \longmapsto \mathcal{F}(X \times_Y V/X)
$$
；由评注 \ref{etale-cohomology-remark-direct-image-sheaf}，这是层。有时写作
$f_* = f_{small, *}$，以区别于其他正像函子（例如通常的 Zariski 推前或
$f_{big, *}$）。
\end{definition}

\noindent
完全相同的讨论仍然适用，并得到函子
$$
f_* = f_{small, *} :
\Sh(X_\etale)
\longrightarrow
\Sh(Y_\etale)
$$
以及
$$
f_* = f_{small, *} :
\textit{Ab}(X_\etale)
\longrightarrow
\textit{Ab}(Y_\etale)
$$
，仍称为{\it 正像}。

\medskip\noindent
阿贝尔层上的函子 $f_*$ 是左正合的。（关于阿贝尔范畴之间的函子为左正合是什么意思，
见 Homology，第 \ref{homology-section-functors} 节。）事实上，若
$0 \to \mathcal{F}_1 \to \mathcal{F}_2 \to \mathcal{F}_3$
在 $X_\etale$ 上正合，则对每个
$U/X \in \Ob(X_\etale)$
，阿贝尔群序列
$0 \to \mathcal{F}_1(U) \to \mathcal{F}_2(U) \to \mathcal{F}_3(U)$
正合。因此，对每个 $V/Y \in \Ob(Y_\etale)$，阿贝尔群序列
$0 \to f_*\mathcal{F}_1(V) \to f_*\mathcal{F}_2(V) \to f_*\mathcal{F}_3(V)$
正合，因为这就是前一个序列取 $U = X \times_Y V$ 所得的序列。

\begin{definition}
\label{etale-cohomology-definition-higher-direct-images}
设 $f: X \to Y$ 是概形态射。右导出函子 $\{R^pf_*\}_{p \geq 1}$（它们来自函子
$f_* : \textit{Ab}(X_\etale) \to \textit{Ab}(Y_\etale)$）称为{\it 高阶正像}。
\end{definition}

\noindent
关于高阶正像及其导出范畴版本，将在（此处待补充未来的参考文献）中作更详细的讨论。



\section{逆像}
\label{etale-cohomology-section-inverse-image}

\noindent
本节简要讨论小 \'etale 网站上层的拉回。其精确构造见 Topologies，第
\ref{topologies-section-etale} 节。

\begin{definition}
\label{etale-cohomology-definition-inverse-image}
设 $f: X\to Y$ 是概形态射。{\it 逆像}，或称{\it 拉回}\footnote{对于集合层或
阿贝尔群层的拉回，我们使用记号 $f^{-1}$；而把 $f^*$ 留给通过环化网站/拓扑斯态射所得的
模层拉回。}函子是
$$
f^{-1} = f_{small}^{-1} :
\Sh(Y_\etale)
\longrightarrow
\Sh(X_\etale)
$$
以及
$$
f^{-1} = f_{small}^{-1} :
\textit{Ab}(Y_\etale)
\longrightarrow
\textit{Ab}(X_\etale)
$$
；它们左伴随于 $f_* = f_{small, *}$。因此，$f^{-1}$ 由下列事实刻画：
$$
\Hom_{{\Sh(X_\etale)}} (f^{-1}\mathcal{G}, \mathcal{F})
=
\Hom_{\Sh(Y_\etale)} (\mathcal{G}, f_*\mathcal{F})
$$
对任意 $\mathcal{F} \in \Sh(X_\etale)$ 和
$\mathcal{G} \in \Sh(Y_\etale)$ 函子性地成立。类似地，有
$$
\Hom_{{\textit{Ab}(X_\etale)}} (f^{-1}\mathcal{G}, \mathcal{F})
=
\Hom_{\textit{Ab}(Y_\etale)} (\mathcal{G}, f_*\mathcal{F})
$$
，其中 $\mathcal{F} \in \textit{Ab}(X_\etale)$ 且
$\mathcal{G} \in \textit{Ab}(Y_\etale)$。
\end{definition}

\noindent
这种伴随是否存在并不显然。另一方面，它在相当一般的框架中存在，见下面的评注
\ref{etale-cohomology-remark-functoriality-general}。一般机制说明，$f^{-1}\mathcal{G}$ 是下述预层的相伴层：
\begin{equation}
\label{etale-cohomology-equation-pullback}
U/X
\longmapsto
\colim_{U \to X \times_Y V} \mathcal{G}(V/Y)
\end{equation}
，其中余极限取遍序对
$(V/Y, \varphi : U/X \to X \times_Y V/X)$ 所成的范畴。为看出这一点，
将 Sites，命题 \ref{sites-proposition-get-morphism} 应用于方程
(\ref{etale-cohomology-equation-functorial}) 的函子 $u$，并使用 Sites，第
\ref{sites-section-continuous-functors} 节和第
\ref{sites-section-functoriality-PSh} 节对 $u_s = (u_p\ )^\#$ 的描述。
我们偶尔会使用这个拉回公式来证明它的一些基本性质。

\begin{lemma}
\label{etale-cohomology-lemma-stalk-pullback}
设 $f : X \to Y$ 是概形态射。
\begin{enumerate}
\item 函子
$f^{-1} : \textit{Ab}(Y_\etale) \to \textit{Ab}(X_\etale)$
是正合的。
\item 函子
$f^{-1} : \Sh(Y_\etale) \to \Sh(X_\etale)$
是正合的；换言之，它与有限极限和余极限交换，见 Categories，定义
\ref{categories-definition-exact}。
\item 设 $\overline{x} \to X$ 是几何点。设 $\mathcal{G}$ 是 $Y_\etale$ 上的层。
则有典范等同
$$
(f^{-1}\mathcal{G})_{\overline{x}} = \mathcal{G}_{\overline{y}}.
$$
，其中 $\overline{y} = f \circ \overline{x}$。
\item 对任意 \'etale 态射 $V \to Y$，有 $f^{-1}h_V = h_{X \times_Y V}$。
\end{enumerate}
\end{lemma}

\begin{proof}
$f^{-1}$ 在集合层上的正合性，是把 Sites，命题
\ref{sites-proposition-get-morphism} 应用于方程
(\ref{etale-cohomology-equation-functorial}) 中的函子 $u$ 所得的推论。
事实上，拉回的正合性是拓扑斯态射（也可说网站态射）定义的一部分。因此 (2) 成立。
它蕴含 (1)，因为给定阿贝尔层 $\mathcal{G}$（在 $Y_\etale$ 上），
$f^{-1}\mathcal{F}$ 的底层集合层等于 $f^{-1}$ 应用于 $\mathcal{F}$ 的底层集合层，
见 Sites，第 \ref{sites-section-sheaves-algebraic-structures} 节。另见 Modules on Sites，
引理 \ref{sites-modules-lemma-flat-pullback-exact}。文献中有时通过定理
\ref{etale-cohomology-theorem-exactness-stalks} 从 (3) 推出 (1) 与 (2)。

\medskip\noindent
(3) 是关于拉回之茎的一般事实，见 Sites，引理
\ref{sites-lemma-point-morphism-sites}。下面也直接证明 (3)。注意，由引理
\ref{etale-cohomology-lemma-stalk-exact}，取茎与层化交换。回顾 $f^{-1}\mathcal{G}$ 是下述预层的相伴层：
$$
U \longrightarrow \colim_{U \to X \times_Y V} \mathcal{G}(V),
$$
见方程 (\ref{etale-cohomology-equation-pullback})。因此有
\begin{align*}
(f^{-1}\mathcal{G})_{\overline{x}}
& = \colim_{(U, \overline{u})} f^{-1}\mathcal{G}(U) \\
& = \colim_{(U, \overline{u})}
\colim_{a : U \to X \times_Y V} \mathcal{G}(V) \\
& = \colim_{(V, \overline{v})} \mathcal{G}(V) \\
& = \mathcal{G}_{\overline{y}}
\end{align*}
；在第三个等号中，序对 $(U, \overline{u})$ 和映射
$a : U \to X \times_Y V$ 对应于序对 $(V, a \circ \overline{u})$。

\medskip\noindent
(4) 可用类似方法证明，只需等同定义 $f^{-1}h_V$ 的余极限。也可以使用 Yoneda 引理
（Categories，引理 \ref{categories-lemma-yoneda}）和函子性等式
$$
\Mor_{\Sh(X_\etale)}(f^{-1}h_V, \mathcal{F}) =
\Mor_{\Sh(Y_\etale)}(h_V, f_*\mathcal{F}) =
f_*\mathcal{F}(V) = \mathcal{F}(X \times_Y V)
$$
，再结合可表预层是层这一事实。完全一般的结果另见 Sites，引理
\ref{sites-lemma-pullback-representable-sheaf}。
\end{proof}

\noindent
函子对 $(f_*, f^{-1})$ 定义一个小 \'etale 拓扑斯态射
$$
f_{small} :
\Sh(X_\etale)
\longrightarrow
\Sh(Y_\etale)
$$
。关于层上同调的许多一般结论对拓扑斯与拓扑斯态射都成立。我们将尽量指出哪些结果是一般的，
哪些是 \'etale 拓扑斯所特有的。

\begin{remark}
\label{etale-cohomology-remark-functoriality-general}
更一般地，设 $\mathcal{C}_1, \mathcal{C}_2$ 是网站，并假定它们有终对象和纤维积。
设 $u: \mathcal{C}_2 \to \mathcal{C}_1$ 是满足下列条件的函子：
\begin{enumerate}
\item 若 $\{V_i \to V\}$ 是 $\mathcal{C}_2$ 的覆盖，则
$\{u(V_i) \to u(V)\}$ 是 $\mathcal{C}_1$ 的覆盖（此时称 $u$ 是{\it 连续的}）；以及
\item $u$ 与有限极限交换（即 $u$ 左正合；也即 $u$ 保持纤维积与终对象）。
\end{enumerate}
于是可以定义
$f_*: \Sh(\mathcal{C}_1) \to \Sh(\mathcal{C}_2)$
，令 $ f_* \mathcal{F}(V) = \mathcal{F}(u(V))$。此外，存在正合函子
$f^{-1}$，它左伴随于 $f_*$，见 Sites，定义
\ref{sites-definition-morphism-sites} 和命题
\ref{sites-proposition-get-morphism}。警告：仅要求 $u$ 连续且与纤维积交换，
不足以得到拓扑斯态射。
\end{remark}




\section{大拓扑斯的函子性}
\label{etale-cohomology-section-functoriality-big-topoi}

\noindent
给定概形态射 $f : X \to Y$，有大量与 $f$ 相伴的拓扑斯态射；其列表见
Topologies，第 \ref{topologies-section-change-topologies} 节。也许最常用的是拓扑斯态射
$$
f_{big} = f_{big, \tau} :
\Sh((\Sch/X)_\tau)
\longrightarrow
\Sh((\Sch/Y)_\tau)
$$
，其中 $\tau \in \{Zariski, \etale, smooth, syntomic, fppf\}$。
它们分别对应于连续函子
$$
(\Sch/Y)_\tau \longrightarrow (\Sch/X)_\tau, \quad
V/Y \longmapsto X \times_Y V/X
$$
；该函子保持终对象、纤维积与覆盖，因而定义网站态射
$$
f_{big} : (\Sch/X)_\tau \longrightarrow (\Sch/Y)_\tau.
$$
见 Topologies，第 \ref{topologies-section-zariski} 节、
\ref{topologies-section-etale},
\ref{topologies-section-smooth},
\ref{topologies-section-syntomic} 和
\ref{topologies-section-fppf}。特别地，沿 $f_{big}$ 的推前由规则
$$
(f_{big, *}\mathcal{F})(V/Y) = \mathcal{F}(X \times_Y V/X)
$$
给出。事实证明，这些拓扑斯态射有一个容易描述的逆像函子 $f_{big}^{-1}$。即有
$$
(f_{big}^{-1}\mathcal{G})(U/X) = \mathcal{G}(U/Y)
$$
，其中 $U/Y$ 的结构态射是结构态射 $U \to X$ 与 $f$ 的复合，见 Topologies，引理
\ref{topologies-lemma-morphism-big}、
\ref{topologies-lemma-morphism-big-etale},
\ref{topologies-lemma-morphism-big-smooth},
\ref{topologies-lemma-morphism-big-syntomic} 和
\ref{topologies-lemma-morphism-big-fppf}.







\section{函子性与模层}
\label{etale-cohomology-section-morphisms-modules}

\noindent
本节将在 \'etale 拓扑的框架下，重新表述 Descent，第
\ref{descent-section-quasi-coherent-sheaves} 节、
\ref{descent-section-quasi-coherent-cohomology} 节和
\ref{descent-section-quasi-coherent-sheaves-bis}
所说明的一些材料。设 $f : X \to Y$ 是概形态射。前面已经看到（见第
\ref{etale-cohomology-section-functoriality}、\ref{etale-cohomology-section-direct-image} 和
\ref{etale-cohomology-section-inverse-image} 节），它诱导小 \'etale 网站的态射 $f_{small}$。
在 Descent，评注 \ref{descent-remark-change-topologies-ringed} 中已经看到，
$f$ 还诱导自然映射
$$
f_{small}^\sharp :
\mathcal{O}_{Y_\etale}
\longrightarrow
f_{small, *}\mathcal{O}_{X_\etale}
$$
；这是 $Y_\etale$ 上环层的映射，并且 $(f_{small}, f_{small}^\sharp)$
是环化网站的态射。环化网站态射的定义见 Modules on Sites，定义
\ref{sites-modules-definition-ringed-site}。这里只回顾，$f_{small}^\sharp$ 由相容映射系
$$
\text{pr}_V^\sharp : \mathcal{O}(V) \longrightarrow \mathcal{O}(X \times_Y V)
$$
定义，其中 $V$ 遍历 $Y_\etale$ 的对象。

\medskip\noindent
显然，该构造与概形态射的复合相容。更确切地，若 $f : X \to Y$ 和
$g : Y \to Z$ 是概形态射，则有
$$
(g_{small}, g_{small}^\sharp)
\circ
(f_{small}, f_{small}^\sharp)
=
((g \circ f)_{small}, (g \circ f)_{small}^\sharp)
$$
，这是环化拓扑斯态射的等式。此外，由 Modules on Sites，定义
\ref{sites-modules-definition-pushforward} 可知，给定概形态射 $f : X \to Y$，
可得到良定义的拉回与正像函子
\begin{align*}
f_{small}^* :
\textit{Mod}(\mathcal{O}_{Y_\etale})
\longrightarrow
\textit{Mod}(\mathcal{O}_{X_\etale}), \\
f_{small, *} :
\textit{Mod}(\mathcal{O}_{X_\etale})
\longrightarrow
\textit{Mod}(\mathcal{O}_{Y_\etale})
\end{align*}
，它们按通常方式互为伴随。若 $g : Y \to Z$ 是另一个概形态射，则有
$(g \circ f)_{small}^* = f_{small}^* \circ g_{small}^*$
以及 $(g \circ f)_{small, *} = g_{small, *} \circ f_{small, *}$，
这是由上面关于复合的讨论得出的。

\medskip\noindent
由所有 $\mathcal{O}_X$-模（定义在 $X$ 上）组成的范畴，与由所有
$\mathcal{O}_{X_\etale}$-模（定义在 $X_\etale$ 上）组成的范畴之间，有相当大的差别。
但是 Descent，第
\ref{descent-section-quasi-coherent-sheaves} 节、
\ref{descent-section-quasi-coherent-cohomology} 节和
\ref{descent-section-quasi-coherent-sheaves-bis}
的结果告诉我们，考虑 $S$ 上的拟凝聚模与考虑 $S_\etale$ 上的拟凝聚模之间，
并没有太大差别。（例如，已经在定理 \ref{etale-cohomology-theorem-quasi-coherent} 中看到这一点。）
特别地，若 $f : X \to Y$ 是任意概形态射，则对拟凝聚层而言，拉回函子
$f_{small}^*$ 与 $f^*$ 相符，见 Descent，命题
\ref{descent-proposition-equivalence-quasi-coherent-functorial}。
此外，只要 $f$ 拟紧且拟分离，对推前也有相同结论，见 Descent，引理
\ref{descent-lemma-higher-direct-images-small-etale}。

\medskip\noindent
最后简述大网站上结构层的函子性。设 $f : X \to Y$ 是概形态射。任取拓扑
$\tau \in \{Zariski,\linebreak[0]
\etale,\linebreak[0] smooth,\linebreak[0] syntomic,\linebreak[0]
fppf\}$。于是态射
$f_{big} : (\Sch/X)_\tau \to (\Sch/Y)_\tau$
通过映射
$$
f_{big}^\sharp : \mathcal{O}_Y \longrightarrow f_{big, *}\mathcal{O}_X
$$
成为环化网站态射，见 Descent，评注
\ref{descent-remark-change-topologies-ringed}。事实上，它由上面小网站情形中的相同构造给出。












\section{比较拓扑}
\label{etale-cohomology-section-compare-topologies}

\noindent
本节开始研究比较不同拓扑下的层时会发生什么。

\begin{lemma}
\label{etale-cohomology-lemma-where-sections-are-equal}
设 $S$ 是概形。设 $\mathcal{F}$ 是 $S_\etale$ 上的集合层。设
$s, t \in \mathcal{F}(S)$。则存在开集 $W \subset S$，它由下列性质刻画：
态射 $f : T \to S$ 经 $W$ 分解，当且仅当 $s|_T = t|_T$（限制是通过
$f_{small}$ 的拉回）。
\end{lemma}

\begin{proof}
考虑如下预层：对 $U \in \Ob(S_\etale)$，若 $s|_U \not = t|_U$，则赋予空集；
否则赋予单点集。显然，这是 $\Sh(S_\etale)$ 的终对象的子层。由引理
\ref{etale-cohomology-lemma-support-subsheaf-final}，找到表示该预层的开集 $W \subset S$。
对几何点 $\overline{x}$（它位于 $S$ 上），可知 $\overline{x} \in W$ 当且仅当
$s$ 与 $t$ 在 $\overline{x}$ 处的茎相同。由引理
\ref{etale-cohomology-lemma-stalk-pullback} 对拉回之茎的描述，可知 $W$ 具有所需性质。
\end{proof}

\begin{lemma}
\label{etale-cohomology-lemma-describe-pullback}
设 $S$ 是概形。设 $\tau \in \{Zariski, \etale\}$。考虑态射
$$
\pi_S : (\Sch/S)_\tau \longrightarrow S_\tau
$$
，它见于 Topologies，引理 \ref{topologies-lemma-at-the-bottom} 或
\ref{topologies-lemma-at-the-bottom-etale}。设 $\mathcal{F}$ 是 $S_\tau$ 上的层。
则 $\pi_S^{-1}\mathcal{F}$ 由规则
$$
(\pi_S^{-1}\mathcal{F})(T) = \Gamma(T_\tau, f_{small}^{-1}\mathcal{F})
$$
给出，其中 $f : T \to S$。此外，$\pi_S^{-1}\mathcal{F}$ 关于 fpqc 覆盖满足层条件。
\end{lemma}

\begin{proof}
注意，有态射 $i_f : \Sh(T_\tau) \to \Sh(\Sch/S)_\tau)$，使得
$\pi_S \circ i_f = f_{small}$；这里把它们视为态射 $T_\tau \to S_\tau$。见 Topologies，引理
\ref{topologies-lemma-put-in-T}、
\ref{topologies-lemma-morphism-big-small},
\ref{topologies-lemma-put-in-T-etale} 和
\ref{topologies-lemma-morphism-big-small-etale}.
由于拉回具有传递性，可知
$i_f^{-1} \pi_S^{-1}\mathcal{F} = f_{small}^{-1}\mathcal{F}$，正如所需。

\medskip\noindent
设 $\{g_i : T_i \to T\}_{i \in I}$ 是 fpqc 覆盖。最后一个断言的含义如下：
给定层 $\mathcal{G}$（定义在 $T_\tau$ 上），并给定截面
$s_i \in \Gamma(T_i, g_{i, small}^{-1}\mathcal{G})$，它们到
$T_i \times_T T_j$ 的拉回相同，则存在唯一截面 $s$；它是 $\mathcal{G}$ 在
$T$ 上的截面，且其到 $T_i$ 的拉回与 $s_i$ 相同。

\medskip\noindent
设 $V \to T$ 是 $T_\tau$ 的对象，并设 $t \in \mathcal{G}(V)$。
对每个 $i$，存在最大的开集 $W_i \subset T_i \times_T V$，使得
$s_i$ 与 $t$ 的拉回作为 $\mathcal{G}$ 到 $W_i \subset T_i \times_T V$
之拉回的截面相同；见引理 \ref{etale-cohomology-lemma-where-sections-are-equal}。
由于 $s_i$ 与 $s_j$ 在 $T_i \times_T T_j$ 上相同，可知
$W_i$ 与 $W_j$ 拉回到 $T_i \times_T T_j \times_T V$ 后给出同一个开集。
由 Descent，引理 \ref{descent-lemma-open-fpqc-covering}，得到开集
$W \subset V$，它在 $T_i \times_T V$ 中的逆像恢复 $W_i$。

\medskip\noindent
根据 $g_{i, small}^{-1}\mathcal{G}$ 的构造，存在一个
$\tau$-覆盖 $\{T_{ij} \to T_i\}_{j \in J_i}$；对每个 $j$，存在开浸入或
\'etale 态射 $V_{ij} \to T$、截面 $t_{ij} \in \mathcal{G}(V_{ij})$，以及交换图
$$
\xymatrix{
T_{ij} \ar[r] \ar[d] & V_{ij} \ar[d] \\
T_i \ar[r] & T
}
$$
使得 $s_i|_{T_{ij}}$ 是 $t_{ij}$ 的拉回。换言之，在将覆盖
$\{T_i \to T\}$ 替换为 $\{T_{ij} \to T\}$ 后，可以假设存在分解
$T_i \to V_i \to T$，其中 $V_i \in \Ob(T_\tau)$，并且截面
$t_i \in \mathcal{G}(V_i)$ 拉回后等于 $s_i$（在 $T_i$ 上）。
由上一段的结果，得到开集 $W_i \subset V_i$，使得 $t_i|_{W_i}$ 与每个
$s_j$ 在 $T_j \times_T W_i$ 上``相同''。注意 $T_i \to V_i$ 经
$W_i$ 分解。因此 $\{W_i \to T\}$ 是一个 $\tau$-覆盖，引理得证。
\end{proof}

\begin{lemma}
\label{etale-cohomology-lemma-sections-upstairs}
设 $S$ 是概形。设 $f : T \to S$ 是满足下列条件的态射：
\begin{enumerate}
\item $f$ 平坦且拟紧；
\item $f$ 的几何纤维连通。
\end{enumerate}
设 $\mathcal{F}$ 是 $S_\etale$ 上的层。则
$\Gamma(S, \mathcal{F}) = \Gamma(T, f^{-1}_{small}\mathcal{F})$。
\end{lemma}

\begin{proof}
存在典范映射
$\Gamma(S, \mathcal{F}) \to \Gamma(T, f_{small}^{-1}\mathcal{F})$.
由于 $f$ 是满射（因为它的纤维连通），可知此映射是单射。

\medskip\noindent
为证明该映射是满射，设
$\alpha \in \Gamma(T, f_{small}^{-1}\mathcal{F})$。
由于 $\{T \to S\}$ 是 fpqc 覆盖，可以用引理 \ref{etale-cohomology-lemma-describe-pullback}
看出，只需证明 $\alpha$ 通过两个投影拉回到 $T \times_S T$ 后给出同一个截面。
设 $\overline{s} \to S$ 是几何点。只需证明这种相同性在
$(T \times_S T)_{\overline{s}}$ 上成立，因为 $T \times_S T$ 的每个几何点
都包含在这些几何纤维之一中。换言之，我们要证明
$\alpha|_{T_{\overline{s}}}$ 拉回到
$$
(T \times_S T)_{\overline{s}} =
T_{\overline{s}} \times_{\overline{s}} T_{\overline{s}}
$$
后，经到 $T_{\overline{s}}$ 的两个投影所得的截面相同。
然而，$\mathcal{F}|_{T_{\overline{s}}}$ 是 $\mathcal{F}|_{\overline{s}}$ 的拉回，
故它是取值于 $\mathcal{F}_{\overline{s}}$ 的常值层。根据假设，
$T_{\overline{s}}$ 连通，所以常值层的任一截面都是常值的。因此
$\alpha|_{T_{\overline{s}}}$ 对应于 $\mathcal{F}_{\overline{s}}$ 的一个元素。
于是，到 $(T \times_S T)_{\overline{s}}$ 的两个拉回均对应于这个相同元素，结论得证。
\end{proof}

\noindent
下面给出引理 \ref{etale-cohomology-lemma-sections-upstairs} 的一个版本，其中不假设该态射平坦。

\begin{lemma}
\label{etale-cohomology-lemma-sections-upstairs-submersive}
设 $S$ 是概形。设 $f : X \to S$ 是满足下列条件的态射：
\begin{enumerate}
\item $f$ 是商映射；
\item $f$ 的几何纤维连通。
\end{enumerate}
设 $\mathcal{F}$ 是 $S_\etale$ 上的层。则
$\Gamma(S, \mathcal{F}) = \Gamma(X, f^{-1}_{small}\mathcal{F})$。
\end{lemma}

\begin{proof}
存在典范映射
$\Gamma(S, \mathcal{F}) \to \Gamma(X, f_{small}^{-1}\mathcal{F})$.
由于 $f$ 是满射（因为它的纤维连通），可知此映射是单射。

\medskip\noindent
为证明该映射是满射，设
$\tau \in \Gamma(X, f_{small}^{-1}\mathcal{F})$。
只需找到一个 \'etale 覆盖 $\{U_i \to S\}$ 和截面
$\sigma_i \in \mathcal{F}(U_i)$，使得 $\sigma_i$ 拉回后等于
$\tau|_{X \times_S U_i}$。事实上，上面证明的单射性保证 $\sigma_i$ 与
$\sigma_j$ 在 $\mathcal{F}$ 到 $U_i \times_S U_j$ 上的限制相同。
因此得到唯一截面 $\sigma \in \mathcal{F}(S)$，其限制为 $\sigma_i$
（在 $U_i$ 上）。于是 $\sigma$ 到 $X$ 的拉回就是 $\tau$，因为这在局部成立。

\medskip\noindent
设 $\overline{x}$ 是 $X$ 的几何点，像为 $\overline{s}$（在 $S$ 中）。
考虑 $\tau$ 在下列茎中的像：
$$
(f_{small}^{-1}\mathcal{F})_{\overline{x}} = \mathcal{F}_{\overline{s}}
$$
见引理 \ref{etale-cohomology-lemma-stalk-pullback}。可以找到一个 \'etale 邻域
$U \to S$（它是 $\overline{s}$ 的邻域）和截面 $\sigma \in \mathcal{F}(U)$，
它映到茎中的这个像。因此，在对 $S$ 以 $U$ 作替换，并对 $X$ 以
$X \times_S U$ 作替换后，可以假设存在截面 $\sigma$，它是 $\mathcal{F}$
在 $S$ 上的截面，并且它在
$(f_{small}^{-1}\mathcal{F})_{\overline{x}}$ 中的像与 $\tau$ 相同。

\medskip\noindent
由引理 \ref{etale-cohomology-lemma-where-sections-are-equal}，存在最大的开集 $W \subset X$，
使得 $f_{small}^{-1}\sigma$ 与 $\tau$ 在 $W$ 上相同，且 $W$ 的构造与进一步的拉回可交换。
注意，$\mathcal{F}$ 到几何纤维 $X_{\overline{s}}$ 的拉回，是把
$\mathcal{F}_{\overline{s}}$ 看作 $\overline{s}$ 上的层后经
$X_{\overline{s}} \to \overline{s}$ 所作的拉回。因此，$\tau$ 与 $\sigma$
给出取值于 $\mathcal{F}_{\overline{s}}$ 的常值层的截面（该层在 $X_{\overline{s}}$ 上），
并且它们在一点处相同。根据假设 $X_{\overline{s}}$ 连通，所以 $W$ 包含
$X_s$。对不同几何纤维作同样的论证可知，$W$ 包含它所遇到的每条纤维。
由于 $f$ 是商映射，可知 $W$ 是 $s$ 的某个开邻域（在 $S$ 中）的逆像。
证明完成。
\end{proof}

\begin{lemma}
\label{etale-cohomology-lemma-sections-base-field-extension}
设 $K/k$ 是域扩张，且 $k$ 是可分代数闭域。设 $S$ 是 $k$ 上的概形。记
$p : S_K = S \times_{\Spec(k)} \Spec(K) \to S$ 为投影。设 $\mathcal{F}$ 是
$S_\etale$ 上的层。则
$\Gamma(S, \mathcal{F}) = \Gamma(S_K, p^{-1}_{small}\mathcal{F})$。
\end{lemma}

\begin{proof}
这由引理 \ref{etale-cohomology-lemma-sections-upstairs} 得出。事实上，$p$ 作为
$\Spec(K) \to \Spec(k)$ 的基变换，显然平坦且拟紧。另一方面，若
$\overline{s} : \Spec(L) \to S$ 是几何点，则 $p$ 在 $\overline{s}$ 上的纤维
是 $K \otimes_k L$ 的谱；由 Algebra，引理
\ref{algebra-lemma-separably-closed-irreducible}，它不可约，因而连通。
\end{proof}









\section{恢复态射}
\label{etale-cohomology-section-morphisms}

\noindent
本节证明：把概形对应到其局部环化小 \'etale 拓扑斯的规则，在适当意义下是全忠实的；
见定理 \ref{etale-cohomology-theorem-fully-faithful}。

\begin{lemma}
\label{etale-cohomology-lemma-morphism-locally-ringed}
设 $f : X \to Y$ 是概形态射。环化网站态射
$(f_{small}, f_{small}^\sharp)$（它与 $f$ 对应）是局部环化网站的态射；见 Modules on Sites，
定义 \ref{sites-modules-definition-morphism-locally-ringed-topoi}。
\end{lemma}

\begin{proof}
注意，该断言有意义，因为我们已经看到
$(X_\etale, \mathcal{O}_{X_\etale})$ 和
$(Y_\etale, \mathcal{O}_{Y_\etale})$ 是局部环化网站；见引理
\ref{etale-cohomology-lemma-etale-site-locally-ringed}。此外，我们知道 $X_\etale$ 有足够多的点；
见定理 \ref{etale-cohomology-theorem-exactness-stalks} 和注记 \ref{etale-cohomology-remarks-enough-points}。
因此，只需证明 $(f_{small}, f_{small}^\sharp)$ 满足 Modules on Sites，
引理 \ref{sites-modules-lemma-locally-ringed-morphism} 的条件 (3)。
为此取一点 $p$（它属于 $X_\etale$）。由引理
\ref{etale-cohomology-lemma-points-small-etale-site}，$p$ 对应于几何点 $\overline{x}$（它位于 $X$ 上）。
由引理 \ref{etale-cohomology-lemma-stalk-pullback}，点 $q = f_{small} \circ p$ 对应于几何点
$\overline{y} = f \circ \overline{x}$（它位于 $Y$ 上）。因此，我们要证明的断言是：
所诱导的茎映射
$$
(\mathcal{O}_Y)_{\overline{y}} \longrightarrow (\mathcal{O}_X)_{\overline{x}}
$$
是局部环同态。设 $a \in (\mathcal{O}_Y)_{\overline{y}}$ 是左侧的一个元素，
它映到右侧的极大理想中的元素。设 $a$ 是三元组 $(V, \overline{v}, a)$ 的等价类，
其中 $V \to Y$ 是 \'etale 态射，$\overline{v} : \overline{x} \to V$ 是 $Y$ 上的态射，
且 $a \in \mathcal{O}(V)$。它映到下列三元组的等价类：
$(X \times_Y V, \overline{x} \times \overline{v}, \text{pr}_V^\sharp(a))$
（它位于局部环 $(\mathcal{O}_X)_{\overline{x}}$ 中）。显然，属于极大理想意味着：
把 $\text{pr}_V^\sharp(a)$ 拉回为 $\kappa(\overline{x})$ 的一个元素后得到零。
因而，把 $a$ 拉回到 $\kappa(\overline{x})$ 也得到零。这意味着 $a$ 属于
$(\mathcal{O}_Y)_{\overline{y}}$ 的极大理想。
\end{proof}

\begin{lemma}
\label{etale-cohomology-lemma-2-morphism}
设 $X$、$Y$ 是概形。设 $f : X \to Y$ 是概形态射。设 $t$ 是一个
$2$-态射，其源和目标均为 $(f_{small}, f_{small}^\sharp)$；见 Modules on Sites，
定义 \ref{sites-modules-definition-2-morphism-ringed-topoi}。则 $t = \text{id}$。
\end{lemma}

\begin{proof}
这意味着 $t : f^{-1}_{small} \to f^{-1}_{small}$ 是一个函子变换，并且下图交换：
$$
\xymatrix{
f_{small}^{-1}\mathcal{O}_Y
\ar[rd]_{f_{small}^\sharp}  & &
f_{small}^{-1}\mathcal{O}_Y \ar[ll]^t \ar[ld]^{f_{small}^\sharp} \\
& \mathcal{O}_X
}
$$
设 $V \to Y$ 是 \'etale 态射，且 $V$ 仿射。由 Morphisms，引理
\ref{morphisms-lemma-quasi-affine-finite-type-over-S}，可以选择浸入
$i : V \to \mathbf{A}^n_Y$（在 $Y$ 上）。用层来表述，这意味着 $i$ 诱导层的单射
$h_i : h_V \to \prod_{j = 1, \ldots, n} \mathcal{O}_Y$。
基变换 $i'$（它把 $i$ 基变换到 $X$ 得到）是浸入
（Schemes，引理 \ref{schemes-lemma-base-change-immersion}）。因此
$i' : X \times_Y V \to \mathbf{A}^n_X$ 是浸入，这又意味着
$h_{i'} : h_{X \times_Y V} \to \prod_{j = 1, \ldots, n} \mathcal{O}_X$
是层的单射。借助引理 \ref{etale-cohomology-lemma-stalk-pullback} 的等同
$f_{small}^{-1}h_V = h_{X \times_Y V}$，映射 $h_{i'}$ 等于
$$
\xymatrix{
f_{small}^{-1}h_V \ar[r]^-{f^{-1}h_i} &
\prod_{j = 1, \ldots, n} f_{small}^{-1}\mathcal{O}_Y
\ar[r]^{\prod f^\sharp} &
\prod_{j = 1, \ldots, n} \mathcal{O}_X
}
$$
（验证略去）。这意味着映射 $t : f_{small}^{-1}h_V \to f_{small}^{-1}h_V$
可置于下列交换图中：
$$
\xymatrix{
f_{small}^{-1}h_V \ar[r]^-{f^{-1}h_i} \ar[d]^t &
\prod_{j = 1, \ldots, n} f_{small}^{-1}\mathcal{O}_Y
\ar[r]^-{\prod f^\sharp} \ar[d]^{\prod t} &
\prod_{j = 1, \ldots, n} \mathcal{O}_X \ar[d]^{\text{id}}\\
f_{small}^{-1}h_V \ar[r]^-{f^{-1}h_i} &
\prod_{j = 1, \ldots, n} f_{small}^{-1}\mathcal{O}_Y
\ar[r]^-{\prod f^\sharp} &
\prod_{j = 1, \ldots, n} \mathcal{O}_X
}
$$
右方块的交换性由上面说明的关于 $t$ 的假设得出。
根据前面的讨论，横箭头的复合是单射，故左边的竖箭头也是恒等映射。
$Y_\etale$ 上的任何集合层都接受一个满射，其源是形如 $h_V$ 的层的（巨大）余积，
其中 $V$ 仿射（合用 Topologies，引理 \ref{topologies-lemma-alternative} 与
Sites，引理 \ref{sites-lemma-sheaf-coequalizer-representable}）。因此
$t : f_{small}^{-1} \to f_{small}^{-1}$ 是所需的恒等变换。
\end{proof}

\begin{lemma}
\label{etale-cohomology-lemma-faithful}
设 $X$、$Y$ 是概形。任何两个概形态射 $a, b : X \to Y$，如果存在一个 $2$-同构
$(a_{small}, a_{small}^\sharp) \cong (b_{small}, b_{small}^\sharp)$
（它位于环化拓扑斯的 $2$-范畴中），则二者相等。
\end{lemma}

\begin{proof}
让我们仔细论证，因为这里有些容易混淆。设
$t : a_{small}^{-1} \to b_{small}^{-1}$ 为该 $2$-同构。考虑任意开集
$V \subset Y$。注意 $h_V$ 是终层 $*$ 的子层。因此
$a_{small}^{-1}h_V = h_{a^{-1}(V)}$ 和
$b_{small}^{-1}h_V = h_{b^{-1}(V)}$ 都是终层的子层。故同构
$$
t : a_{small}^{-1}h_V = h_{a^{-1}(V)} \to b_{small}^{-1}h_V = h_{b^{-1}(V)}
$$
必为恒等，并且 $a^{-1}(V) = b^{-1}(V)$。因此 $a$ 与 $b$ 在底拓扑空间上相等。
接着取截面 $f \in \mathcal{O}_Y(V)$。它决定且被下列集合层映射决定：
$f : h_V \to \mathcal{O}_Y$。将其拉回并作用 $t$，得到交换图
$$
\xymatrix{
h_{b^{-1}(V)} \ar@{=}[r] &
b_{small}^{-1}h_V \ar[d]^{b_{small}^{-1}(f)} & &
a_{small}^{-1}h_V \ar[d]^{a_{small}^{-1}(f)} \ar[ll]^t &
h_{a^{-1}(V)} \ar@{=}[l]
\\
& b_{small}^{-1}\mathcal{O}_Y
\ar[rd]_{b^\sharp}  & &
a_{small}^{-1}\mathcal{O}_Y \ar[ll]^t \ar[ld]^{a^\sharp} \\
& & \mathcal{O}_X
}
$$
其中三角形的交换性来自 Modules on Sites，第
\ref{sites-modules-section-2-category} 节中 $2$-同构的定义。
上面已看到，顶部水平箭头的复合来自恒等式 $a^{-1}(V) = b^{-1}(V)$。
因此图的交换性给出
$a_{small}^\sharp(f) = b_{small}^\sharp(f)$ 在
$\mathcal{O}_X(a^{-1}(V)) = \mathcal{O}_X(b^{-1}(V))$ 中成立。
由于这对每个开集 $V$ 和每个 $f \in \mathcal{O}_Y(V)$ 都成立，可知概形态射
$a = b$。
\end{proof}

\begin{lemma}
\label{etale-cohomology-lemma-morphism-ringed-etale-topoi-affines}
设 $X$、$Y$ 是仿射概形。设
$$
(g, g^\#) :
(\Sh(X_\etale), \mathcal{O}_X)
\longrightarrow
(\Sh(Y_\etale), \mathcal{O}_Y)
$$
是局部环化拓扑斯的态射。则存在唯一的概形态射 $f : X \to Y$，使得
$(g, g^\#)$ $2$-同构于 $(f_{small}, f_{small}^\sharp)$；见 Modules on Sites，
定义 \ref{sites-modules-definition-2-morphism-ringed-topoi}。
\end{lemma}

\begin{proof}
本证明中用 $\mathcal{O}_X$ 表示小 \'etale 网站 $X_\etale$ 的结构层，
$\mathcal{O}_Y$ 亦然。记 $Y = \Spec(B)$ 且 $X = \Spec(A)$。由于
$B = \Gamma(Y_\etale, \mathcal{O}_Y)$、
$A = \Gamma(X_\etale, \mathcal{O}_X)$，可知 $g^\sharp$ 诱导环同态
$\varphi : B \to A$。令 $f = \Spec(\varphi) : X \to Y$ 为相应的仿射概形态射。
我们将证明这个 $f$ 满足要求。

\medskip\noindent
设 $V \to Y$ 是仿射概形，并且在 $Y$ 上是 \'etale 的。因此可以写成
$V = \Spec(C)$，其中 $C$ 是 \'etale $B$-代数。可以写成
$$
C = B[x_1, \ldots, x_n]/(P_1, \ldots, P_n)
$$
其中 $P_i$ 是多项式，并且 $\Delta = \det(\partial P_i/ \partial x_j)$ 在 $C$ 中可逆；
例如见 Algebra，引理 \ref{algebra-lemma-etale-standard-smooth}。
若 $T$ 是 $Y$ 上的概形，则一个 $T$-值点（取值于 $V$）由 $n$ 个截面
（属于 $\Gamma(T, \mathcal{O}_T)$）给出，它们满足多项式方程
$P_1 = 0, \ldots, P_n = 0$。换言之，层 $h_V$（在 $Y_\etale$ 上）是下列两个映射的等化子：
$$
\xymatrix{
\prod\nolimits_{i = 1, \ldots, n} \mathcal{O}_Y
\ar@<1ex>[r]^a \ar@<-1ex>[r]_b &
\prod\nolimits_{j = 1, \ldots, n} \mathcal{O}_Y
}
$$
其中 $b(h_1, \ldots, h_n) = 0$，而
$a(h_1, \ldots, h_n) =
(P_1(h_1, \ldots, h_n), \ldots, P_n(h_1, \ldots, h_n))$。
由于 $g^{-1}$ 正合，可知下列实线交换图的顶行也是等化子图：
$$
\xymatrix{
g^{-1}h_V \ar[r] \ar@{..>}[d] &
\prod\nolimits_{i = 1, \ldots, n} g^{-1}\mathcal{O}_Y
\ar@<1ex>[r]^{g^{-1}a} \ar@<-1ex>[r]_{g^{-1}b} \ar[d]^{\prod g^\sharp} &
\prod\nolimits_{j = 1, \ldots, n} g^{-1}\mathcal{O}_Y \ar[d]^{\prod g^\sharp}\\
h_{X \times_Y V} \ar[r] &
\prod\nolimits_{i = 1, \ldots, n} \mathcal{O}_X
\ar@<1ex>[r]^{a'} \ar@<-1ex>[r]_{b'} &
\prod\nolimits_{j = 1, \ldots, n} \mathcal{O}_X  \\
}
$$
这里 $b'$ 是零映射，而 $a'$ 是由诸像
$P'_i = \varphi(P_i) \in A[x_1, \ldots, x_n]$ 按同一规则定义的映射：
$a'(h_1, \ldots, h_n) =
(P'_1(h_1, \ldots, h_n), \ldots, P'_n(h_1, \ldots, h_n))$；这就是定义 $a$ 时所用的规则。
图的交换性来自整体截面上的等式 $\varphi = g^\sharp$。底行同样是等化子图；
理由与前面完全相同，因为 $X \times_Y V$ 是仿射概形
$\Spec(A \otimes_B C)$，并且
$A \otimes_B C = A[x_1, \ldots, x_n]/(P'_1, \ldots, P'_n)$。
因此，得到唯一的虚线箭头 $g^{-1}h_V \to h_{X \times_Y V}$，使之嵌入该图。

\medskip\noindent
我们断言层映射 $g^{-1}h_V \to h_{X \times_Y V}$ 是同构。由于 $X$ 的小
\'etale 网站有足够多的点（定理 \ref{etale-cohomology-theorem-exactness-stalks}），只需在茎上证明。
因此，设 $\overline{x}$ 是 $X$ 的几何点，并以 $p$ 表示 $X$ 的小 \'etale 拓扑斯中
相应的点。置 $q = g \circ p$。这是 $Y$ 的小 \'etale 拓扑斯的一个点。
由引理 \ref{etale-cohomology-lemma-points-small-etale-site}，可知 $q$ 对应于几何点
$\overline{y}$（它位于 $Y$ 上）。考虑茎映射
$$
(g^\sharp)_p :
(\mathcal{O}_Y)_{\overline{y}} =
\mathcal{O}_{Y, q} =
(g^{-1}\mathcal{O}_Y)_p
\longrightarrow
\mathcal{O}_{X, p} =
(\mathcal{O}_X)_{\overline{x}}
$$
由于 $(g, g^\sharp)$ 是{\it 局部}环化拓扑斯的态射，$(g^\sharp)_p$ 是严格 Hensel
局部环之间的局部环同态。对上面的大交换图取局部化，并应用 Algebra，引理
\ref{algebra-lemma-strictly-henselian-solutions}，可知
$(g^{-1}h_V)_p \to (h_{X \times_Y V})_p$ 是所需的同构。

\medskip\noindent
我们断言同构 $g^{-1}h_V \to h_{X \times_Y V}$ 是函子性的。具体地，设
$V_1 \to V_2$ 是仿射概形间的态射，且这些概形在 $Y$ 上是 \'etale 的。写成
$V_i = \Spec(C_i)$，其中
$$
C_i = B[x_{i, 1}, \ldots, x_{i, n_i}]/(P_{i, 1}, \ldots, P_{i, n_i})
$$
态射 $V_1 \to V_2$ 由 $B$-代数同态 $C_2 \to C_1$ 给出，后者又由一些多项式
$Q_j \in B[x_{1, 1}, \ldots, x_{1, n_1}]$（其中 $j = 1, \ldots, n_2$）给出。
于是容易证明下列层的图交换：
$$
\xymatrix{
h_{V_1} \ar[d] \ar[r] & \prod_{i = 1, \ldots, n_1} \mathcal{O}_Y
\ar[d]^{Q_1, \ldots, Q_{n_2}}\\
h_{V_2} \ar[r] & \prod_{i = 1, \ldots, n_2} \mathcal{O}_Y
}
$$
将它拉回到 $X_\etale$，得到下列实线交换图：
$$
\xymatrix{
g^{-1}h_{V_1} \ar@{..>}[dd] \ar[rrd] \ar[r] &
\prod_{i = 1, \ldots, n_1} g^{-1}\mathcal{O}_Y
\ar[dd]^{g^\sharp}
\ar[rrd]^{Q_1, \ldots, Q_{n_2}} \\
& & g^{-1}h_{V_2} \ar@{..>}[dd] \ar[r] &
\prod_{i = 1, \ldots, n_2} g^{-1}\mathcal{O}_Y
\ar[dd]^{g^\sharp} \\
h_{X \times_Y V_1} \ar[r] \ar[rrd] &
\prod\nolimits_{i = 1, \ldots, n_1} \mathcal{O}_X
\ar[rrd]^{Q'_1, \ldots, Q'_{n_2}} \\
& & h_{X \times_Y V_2} \ar[r] &
\prod\nolimits_{i = 1, \ldots, n_2} \mathcal{O}_X
}
$$
其中 $Q'_j \in A[x_{1, 1}, \ldots, x_{1, n_1}]$ 是 $Q_j$ 经 $\varphi$ 得到的像。
由于虚线箭头存在、使两个方块交换，而且水平箭头是单射，所以整个图交换。
这证明了函子性（也证明了 $g^{-1}h_V \to h_{X \times_Y V}$ 的构造不依赖于
表示的选择，尽管严格说来我们不需要证明这一点）。

\medskip\noindent
此时可以证明 $f_{small, *} \cong g_*$。事实上，设 $\mathcal{F}$ 是 $X_\etale$
上的层。对每个仿射的 $V \in \Ob(X_\etale)$，有
\begin{align*}
(g_*\mathcal{F})(V)
& =
\Mor_{\Sh(Y_\etale)}(h_V, g_*\mathcal{F}) \\
& =
\Mor_{\Sh(X_\etale)}(g^{-1}h_V, \mathcal{F}) \\
& =
\Mor_{\Sh(X_\etale)}(h_{X \times_Y V}, \mathcal{F}) \\
& =
\mathcal{F}(X \times_Y V) \\
& =
f_{small, *}\mathcal{F}(V)
\end{align*}
其中第三个等式使用上面构造的同构 $g^{-1}h_V \cong h_{X \times_Y V}$。
这些同构关于 $\mathcal{F}$ 显然是函子性的，关于 $V$ 也具有函子性，因为同构
$g^{-1}h_V \cong h_{X \times_Y V}$ 本身是函子性的。$Y_\etale$ 上的任意层都由其
到仿射概形子范畴的限制所决定（Topologies，引理
\ref{topologies-lemma-alternative}），所以得到所需的函子同构
$f_{small, *} \cong g_*$。

\medskip\noindent
最后，需要检验：借助上面的同构 $f_{small, *} \cong g_*$，映射
$f_{small}^\sharp$ 与 $g^\sharp$ 相同。根据构造，这对 $\mathcal{O}_Y$ 的整体截面，
即对 $B$ 的元素已经成立。只需对仿射概形 $V$ 上的截面检验，其中它在
$Y$ 上是 \'etale 的（再次使用 Topologies，引理
\ref{topologies-lemma-alternative}）。像前面一样写
$V = \Spec(C)$、$C = B[x_i]/(P_j)$，只需检验坐标函数 $x_i$ 映到
$\mathcal{O}_X$ 在 $X \times_Y V$ 上的同一截面。这恰好就是下图交换的含义：
$$
\xymatrix{
g^{-1}h_V \ar[r] \ar@{..>}[d] &
\prod\nolimits_{i = 1, \ldots, n} g^{-1}\mathcal{O}_Y
\ar[d]^{\prod g^\sharp} \\
h_{X \times_Y V} \ar[r] &
\prod\nolimits_{i = 1, \ldots, n} \mathcal{O}_X
}
$$
因此引理得证。
\end{proof}

\noindent
下面给出一般概形的版本。

\begin{theorem}
\label{etale-cohomology-theorem-fully-faithful}
设 $X$、$Y$ 是概形。设
$$
(g, g^\#) :
(\Sh(X_\etale), \mathcal{O}_X)
\longrightarrow
(\Sh(Y_\etale), \mathcal{O}_Y)
$$
是局部环化拓扑斯的态射。则存在唯一的概形态射 $f : X \to Y$，使得
$(g, g^\#)$ 同构于 $(f_{small}, f_{small}^\sharp)$。换言之，构造
$$
\Sch \longrightarrow \textit{Locally ringed topoi},
\quad
X \longrightarrow (X_\etale, \mathcal{O}_X)
$$
是全忠实的（右侧的态射取到 $2$-同构为止）。
\end{theorem}

\begin{proof}
可以仔细调整引理 \ref{etale-cohomology-lemma-morphism-ringed-etale-topoi-affines} 的证明论证，
使之适用于整体情形，从而证明本定理。不过，我们希望说明如何利用引理
\ref{etale-cohomology-lemma-2-morphism} 的结论所提供的刚性，把前一引理的结果粘合成整体态射。
遗憾的是，这有些繁琐。

\medskip\noindent
先在 $Y$ 仿射时证明存在性。此时选取仿射开覆盖 $X = \bigcup U_i$。对每个 $i$，
包含态射 $j_i : U_i \to X$ 诱导局部环化拓扑斯的态射
$(j_{i, small}, j_{i, small}^\sharp) :
(\Sh(U_{i, \etale}), \mathcal{O}_{U_i})
\to
(\Sh(X_\etale), \mathcal{O}_X)$
，这是引理 \ref{etale-cohomology-lemma-morphism-locally-ringed} 的结论。把它与 $(g, g^\sharp)$
复合，得到局部环化拓扑斯的态射
$$
(g, g^\sharp) \circ (j_{i, small}, j_{i, small}^\sharp) :
(\Sh(U_{i, \etale}), \mathcal{O}_{U_i})
\to
(\Sh(Y_\etale), \mathcal{O}_Y)
$$
；见 Modules on Sites，引理
\ref{sites-modules-lemma-composition-morphisms-locally-ringed-topoi}。
由引理 \ref{etale-cohomology-lemma-morphism-ringed-etale-topoi-affines}，存在唯一的概形态射
$f_i : U_i \to Y$ 和一个 $2$-同构
$$
t_i :
(f_{i, small}, f_{i, small}^\sharp)
\longrightarrow
(g, g^\sharp) \circ (j_{i, small}, j_{i, small}^\sharp).
$$
置 $U_{i, i'} = U_i \cap U_{i'}$，并以 $j_{i, i'} : U_{i, i'} \to U_i$ 表示包含态射。
由于有
$j_i \circ j_{i, i'} = j_{i'} \circ j_{i', i}$
，可知
\begin{align*}
(g, g^\sharp) \circ
(j_{i, small}, j_{i, small}^\sharp) \circ
(j_{i, i', small}, j_{i, i', small}^\sharp)
= \\
(g, g^\sharp) \circ
(j_{i', small}, j_{i', small}^\sharp) \circ
(j_{i', i, small}, j_{i', i, small}^\sharp)
\end{align*}
因此由唯一性（见引理 \ref{etale-cohomology-lemma-faithful}），得到
$f_i \circ j_{i, i'} = f_{i'} \circ j_{i', i}$。换言之，概形态射
$f_i = f \circ j_i$ 是某个整体概形态射 $f : X \to Y$ 的限制。考虑下列
$2$-同构的图（为简化记号，我们省略分量 ${}^\sharp$）：
$$
\xymatrix{
g \circ j_{i, small} \circ j_{i, i', small}
\ar[rr]^{t_i \star \text{id}_{j_{i, i', small}}}
\ar@{=}[d] & &
f_{small} \circ j_{i, small} \circ j_{i, i', small} \ar@{=}[d] \\
g \circ j_{i', small} \circ j_{i', i, small}
\ar[rr]^{t_{i'} \star \text{id}_{j_{i', i, small}}} & &
f_{small} \circ j_{i', small} \circ j_{i', i, small}
}
$$
记号 $\star$ 表示水平复合；一般情形见 Categories，定义
\ref{categories-definition-2-category}，本情形见 Sites，第
\ref{sites-section-2-category} 节。由引理 \ref{etale-cohomology-lemma-2-morphism}，该图交换。
因此，对任意层 $\mathcal{G}$（在 $Y_\etale$ 上），同构
$t_i : f_{small}^{-1}\mathcal{G}|_{U_i} \to g^{-1}\mathcal{G}|_{U_i}$
在 $U_{i, i'}$ 上相同，从而得到整体同构
$t : f_{small}^{-1}\mathcal{G} \to g^{-1}\mathcal{G}$。显然，该同构关于
$\mathcal{G}$ 是函子性的，并且与映射 $f_{small}^\sharp$ 和 $g^\sharp$ 相容
（因为它在局部与这些映射相容）。这证明了 $Y$ 仿射时的定理。

\medskip\noindent
在一般情形，设 $V \subset Y$ 是仿射开集。那么 $h_V$ 是终层 $*$ 的子层
（该终层位于 $Y_\etale$ 上）。由于 $g$ 正合，可知 $g^{-1}h_V$ 是
$X_\etale$ 上终层的子层。因此，由引理 \ref{etale-cohomology-lemma-support-subsheaf-final}，
存在开子概形 $W \subset X$，使得 $g^{-1}h_V = h_W$。由 Modules on Sites，
引理 \ref{sites-modules-lemma-localize-morphism-locally-ringed-topoi}，存在下列
局部环化拓扑斯态射的交换图：
$$
\xymatrix{
(\Sh(W_\etale), \mathcal{O}_W) \ar[r] \ar[d]_{g'} &
(\Sh(X_\etale), \mathcal{O}_X) \ar[d]^g \\
(\Sh(V_\etale), \mathcal{O}_V) \ar[r] &
(\Sh(Y_\etale), \mathcal{O}_Y)
}
$$
其中水平箭头是局部化态射（由包含态射 $V \to Y$ 和 $W \to X$ 诱导），
而 $g'$ 由 $g$ 诱导。由上一段的结论，得到概形态射 $f' : W \to V$ 和一个
$2$-同构
$t : (f'_{small}, (f'_{small})^\sharp) \to (g', (g')^\sharp)$.
与前面完全一样，这些态射 $f'$（当仿射开集 $V \subset Y$ 变化时）由唯一性
在交叠处相同，故得到态射 $f : X \to Y$。此外，这些 $2$-同构 $t$ 由引理
\ref{etale-cohomology-lemma-2-morphism} 再次可知在交叠处相容，因而得到整体 $2$-同构
$(f_{small}, (f_{small})^\sharp) \to (g, (g)^\sharp)$.
这正是所需结论。略去了一些细节。
\end{proof}










\section{推前与拉回}
\label{etale-cohomology-section-monomorphisms}

\noindent
设 $f : X \to Y$ 是概形态射。下面列出我们将要考虑的条件：
\begin{enumerate}
\item[(A)] 对每个 \'etale 态射 $U \to X$ 和 $u \in U$，存在 \'etale 态射
$V \to Y$、不交并分解 $X \times_Y V = W \amalg W'$，以及态射
$h : W \to U$（在 $X$ 上），使得 $u$ 属于 $h$ 的像。
\item[(B)] 对每个 \'etale 态射 $V \to Y$ 和每个 \'etale 覆盖
$\{U_i \to X \times_Y V\}$，存在 \'etale 覆盖 $\{V_j \to V\}$，使得对每个
$j$ 都有 $X \times_Y V_j = \coprod W_{ij}$，其中 $W_{ij} \to X \times_Y V$
经 $U_i \to X \times_Y V$ 分解（对某个 $i$）。
\item[(C)] 对每个 \'etale 态射 $U \to X$，存在 \'etale 态射 $V \to Y$ 和满态射
$X \times_Y V \to U$（在 $X$ 上）。
\end{enumerate}
这些性质中的每一个都可以用函子 $f_{small, *}$ 的行为来表述。我们将在随后几节
阐明这一点。



\section{性质 (A)}
\label{etale-cohomology-section-A}

\noindent
性质 (A) 的定义见第 \ref{etale-cohomology-section-monomorphisms} 节。

\begin{lemma}
\label{etale-cohomology-lemma-property-A-implies}
设 $f : X \to Y$ 是概形态射。假设 (A) 成立。
\begin{enumerate}
\item
$f_{small, *} :
\textit{Ab}(X_\etale)
\to
\textit{Ab}(Y_\etale)$
 保持对单射和满射的反映，
\item 映射 $f_{small}^{-1}f_{small, *}\mathcal{F} \to \mathcal{F}$
对任意 Abel 层 $\mathcal{F}$（在 $X_\etale$ 上）都是满射，
\item
$f_{small, *} :
\textit{Ab}(X_\etale)
\to
\textit{Ab}(Y_\etale)$
是忠实的。
\end{enumerate}
\end{lemma}

\begin{proof}
设 $\mathcal{F}$ 是 Abel 层（在 $X_\etale$ 上）。设 $U$ 是 $X_\etale$ 的对象。
根据假设，可以找到覆盖 $\{W_i \to U\}$（它位于 $X_\etale$ 中），使每个 $W_i$
都是 $X \times_Y V_i$ 的开闭子概形，其中 $V_i$ 是 $Y_\etale$ 的某个对象。
层条件给出
$$
\mathcal{F}(U) \subset \prod \mathcal{F}(W_i)
$$
，且 $\mathcal{F}(W_i)$ 是
$\mathcal{F}(X \times_Y V_i) = f_{small, *}\mathcal{F}(V_i)$ 的直和项。
因此，$f_{small, *}$ 显然反映单射。

\medskip\noindent
接着，设 $a : \mathcal{G} \to \mathcal{F}$ 是 Abel 层的映射，且
$f_{small, *}a$ 是满射。设 $s \in \mathcal{F}(U)$，其中 $U$ 如上。
用上面的 $W_i$、$V_i$，可知只需证明 $s|_{W_i}$ 在 \'etale 局部是
$\mathcal{G}$ 的某个截面在 $a$ 下的像。由于 $\mathcal{F}(W_i)$ 是
$\mathcal{F}(X \times_Y V_i)$ 的直和项，只需证明：对任意
$V \in \Ob(Y_\etale)$，任意元素 $s \in \mathcal{F}(X \times_Y V)$
在 $X \times_Y V$ 的 \'etale 局部是 $\mathcal{G}$ 的某个截面在 $a$ 下的像。
由于
$\mathcal{F}(X \times_Y V) = f_{small, *}\mathcal{F}(V)$
，根据假设存在覆盖 $\{V_j \to V\}$，使得 $s$ 是
$s_j \in f_{small, *}\mathcal{G}(V_j) = \mathcal{G}(X \times_Y V_j)$ 的像。
这证明了 $f_{small, *}$ 反映满射。

\medskip\noindent
(2)、(3) 由 (1) 形式地得出；见 Modules on Sites，引理
\ref{sites-modules-lemma-reflect-surjections}。
\end{proof}

\begin{lemma}
\label{etale-cohomology-lemma-locally-quasi-finite-A}
设 $f : X \to Y$ 是分离且局部拟有限的概形态射。则上面的性质 (A) 成立。
\end{lemma}

\begin{proof}
设 $U \to X$ 是 \'etale 态射，且 $u \in U$。几何断言 (A) 直接归结为
$U$ 与 $Y$ 都是仿射概形的情形。以 $x \in X$ 和 $y \in Y$ 表示 $u$ 的像。
由于 $X \to Y$ 局部拟有限，而 $U \to X$ 局部拟有限（见 Morphisms，引理
\ref{morphisms-lemma-etale-locally-quasi-finite}），可知 $U \to Y$ 局部拟有限
（见 Morphisms，引理 \ref{morphisms-lemma-composition-quasi-finite}）。此外，
$X \to Y$ 与 $U \to Y$ 都分离。因此 More on Morphisms，引理
\ref{more-morphisms-lemma-etale-splits-off-quasi-finite-part-technical-variant}
可应用于这两个态射。这意味着可以选择 \'etale 邻域
$(V, v) \to (Y, y)$，使得
$$
X \times_Y V = W \amalg R, \quad
U \times_Y V = W' \amalg R'
$$
以及点 $w \in W$、$w' \in W'$，使得
\begin{enumerate}
\item $W$、$R$ 在 $X \times_Y V$ 中开且闭；
\item $W'$、$R'$ 在 $U \times_Y V$ 中开且闭；
\item $W \to V$ 与 $W' \to V$ 有限；
\item $w$、$w'$ 映到 $v$；
\item $\kappa(v) \subset \kappa(w)$ 和 $\kappa(v) \subset \kappa(w')$
是纯不可分扩张；
\item $W$ 或 $W'$ 中没有其他点映到 $v$。
\end{enumerate}
有如下交换图：
$$
\xymatrix{
U \ar[d] & U \times_Y V \ar[l] \ar[d] & W' \amalg R' \ar[d] \ar[l] \\
X \ar[d] & X \times_Y V \ar[l] \ar[d] & W \amalg R \ar[l] \\
Y & V \ar[l]
}
$$
缩小 $V$ 后，可以假设 $W'$ 映入 $W$：只需去掉 $R$ 的逆像（它在 $W'$ 中）的像；
这是闭集（因为 $W' \to V$ 有限），且不含 $v$。于是 $W' \to W$ 有限，
因为 $W \to V$ 与 $W' \to V$ 都有限。因此 $W' \to W$ 是有限 \'etale 态射，
并且在 $w$ 上的纤维中恰有一个点满足 $\kappa(w) = \kappa(w')$。故
$W' \to W$ 在某个开邻域 $W^\circ$（它包含 $w$）上是同构；见 \'Etale Morphisms，
引理 \ref{etale-lemma-finite-etale-one-point}。由于 $W \to V$ 有限，
$W \setminus W^\circ$ 的像是闭子集 $T$（在 $V$ 中），且不含 $v$。因此，把
$V$ 替换成 $V \setminus T$ 后，可以假设 $W' \to W$ 是同构。此时分解
$X \times_Y V = W \amalg R$ 和态射 $W \to U$ 正是所需，证明完成。
\end{proof}

\begin{lemma}
\label{etale-cohomology-lemma-integral-A}
设 $f : X \to Y$ 是整概形态射。则性质 (A) 成立。
\end{lemma}

\begin{proof}
设 $U \to X$ 是 \'etale 态射，且 $u \in U$ 是一点。需要找到 \'etale 态射
$V \to Y$、不交并分解 $X \times_Y V = W \amalg W'$，以及 $X$-态射
$W \to U$，使得 $u$ 属于其像。可以缩小 $U$ 和 $Y$，并假设 $U$ 与 $Y$
仿射。此时 $X$ 也仿射，因为整态射按定义是仿射的。写成
$Y = \Spec(A)$、$X = \Spec(B)$ 和 $U = \Spec(C)$。那么 $A \to B$ 是整环同态，
而 $B \to C$ 是 \'etale 环同态。由 Algebra，引理 \ref{algebra-lemma-etale}，
可以找到有限 $A$-子代数 $B' \subset B$ 和 \'etale 环同态 $B' \to C'$，使得
$C = B \otimes_{B'} C'$。因此，问题归结为 \'etale 态射
$U' = \Spec(C') \to X' = \Spec(B')$，它位于有限态射 $X' \to Y$ 上。在这种情形，结论由引理
\ref{etale-cohomology-lemma-locally-quasi-finite-A} 得出。
\end{proof}

\begin{lemma}
\label{etale-cohomology-lemma-when-push-pull-surjective}
设 $f : X \to Y$ 是概形态射。以
$f_{small} :
\Sh(X_\etale)
\to
\Sh(Y_\etale)$
表示相应的小 \'etale 拓扑斯态射。假设下列条件至少有一个成立：
\begin{enumerate}
\item $f$ 是整态射；
\item $f$ 分离且局部拟有限。
\end{enumerate}
则函子
$f_{small, *} :
\textit{Ab}(X_\etale)
\to
\textit{Ab}(Y_\etale)$
具有下列性质：
\begin{enumerate}
\item 映射
$f_{small}^{-1}f_{small, *}\mathcal{F} \to \mathcal{F}$
总是满射；
\item $f_{small, *}$ 是忠实的；
\item $f_{small, *}$ 反映单射和满射。
\end{enumerate}
\end{lemma}

\begin{proof}
合用引理 \ref{etale-cohomology-lemma-locally-quasi-finite-A}、
\ref{etale-cohomology-lemma-integral-A} 和 \ref{etale-cohomology-lemma-property-A-implies}。
\end{proof}



\section{性质 (B)}
\label{etale-cohomology-section-B}

\noindent
性质 (B) 的定义见第 \ref{etale-cohomology-section-monomorphisms} 节。

\begin{lemma}
\label{etale-cohomology-lemma-property-B-implies}
设 $f : X \to Y$ 是概形态射。假设 (B) 成立。则函子
$f_{small, *} :
\Sh(X_\etale)
\to
\Sh(Y_\etale)$
把满射变为满射。
\end{lemma}

\begin{proof}
这由 Sites，引理 \ref{sites-lemma-weaker} 得出。
\end{proof}

\begin{lemma}
\label{etale-cohomology-lemma-simplify-B}
设 $f : X \to Y$ 是概形态射。假设
\begin{enumerate}
\item $V \to Y$ 是概形的 \'etale 态射；
\item $\{U_i \to X \times_Y V\}$ 是 \'etale 覆盖；
\item $v \in V$ 是一点。
\end{enumerate}
假设对任意这样的数据，都存在 \'etale 邻域 $(V', v') \to (V, v)$、不交并分解
$X \times_Y V' = \coprod W'_i$，以及态射 $W'_i \to U_i$
（在 $X \times_Y V$ 上）。则性质 (B) 成立。
\end{lemma}

\begin{proof}
略。
\end{proof}

\begin{lemma}
\label{etale-cohomology-lemma-finite-B}
设 $f : X \to Y$ 是有限概形态射。则性质 (B) 成立。
\end{lemma}

\begin{proof}
考虑 \'etale 态射 $V \to Y$、\'etale 覆盖 $\{U_i \to X \times_Y V\}$，以及
$v \in V$。需要找到引理 \ref{etale-cohomology-lemma-simplify-B} 中的 $V' \to V$、分解和映射。
可以缩小 $V$ 和 $Y$，因而可以假设 $V$ 与 $Y$ 仿射。由于 $X$ 在 $Y$ 上有限，
这也蕴含 $X$ 仿射。在证明过程中，可以（有限多次）把 $(V, v)$ 替换为
\'etale 邻域 $(V', v')$，并相应地把覆盖 $\{U_i \to X \times_Y V\}$ 替换为
$\{V' \times_V U_i \to X \times_Y V'\}$。

\medskip\noindent
由于 $X \times_Y V \to V$ 有限，存在有限多个（两两不同的）点
$x_1, \ldots, x_n \in X \times_Y V$ 映到 $v$。把 More on Morphisms，引理
\ref{more-morphisms-lemma-etale-splits-off-quasi-finite-part-technical-variant}
应用于 $X \times_Y V \to V$ 及诸点 $x_1, \ldots, x_n$（它们位于 $v$ 上），得到
\'etale 邻域 $(V', v') \to (V, v)$，使得
$$
X \times_Y V' = R \amalg \coprod T_a
$$
，其中 $T_a \to V'$ 有限，并且恰有一个点 $p_a$ 位于 $v'$ 上；此外
$\kappa(v') \subset \kappa(p_a)$ 是纯不可分扩张，而 $R \to V'$ 在 $v'$ 上的纤维为空。
由于 $X \to Y$ 有限，$R \to V'$ 也有限。因此，缩小 $V'$ 后可以假设
$R = \emptyset$。于是可以假设
$X \times_Y V = X_1 \amalg \ldots \amalg X_n$，其中恰有一个点
$x_l \in X_l$ 位于 $v$ 上，并且 $\kappa(v) \subset \kappa(x_l)$ 是纯不可分扩张。
注意，这个性质在加细 \'etale 邻域 $(V, v)$ 时保持不变。

\medskip\noindent
对每个 $l$，选择一个 $i_l$ 和一点 $u_l \in U_{i_l}$，它映到 $x_l$。现在，把性质 (A)
应用于有限态射 $X \times_Y V \to V$、\'etale 态射
$U_{i_l} \to X \times_Y V$ 以及诸点 $u_l$。引理 \ref{etale-cohomology-lemma-integral-A}
保证可以这样做。由此得到 \'etale 邻域 $(V', v') \to (V, v)$ 和分解
$$
X \times_Y V' = W_l \amalg R_l
$$
，以及 $X$-态射 $a_l : W_l \to U_{i_l}$，其像包含 $u_{i_l}$。图示如下：
$$
\xymatrix{
& & & U_{i_l} \ar[d] & \\
W_l \ar[rrru] \ar[r] & W_l \amalg R_l \ar@{=}[r] &
X \times_Y V' \ar[r] \ar[d] &
X \times_Y V \ar[r] \ar[d] & X \ar[d] \\
& & V' \ar[r] & V \ar[r] & Y
}
$$
把 $(V, v)$ 替换为 $(V', v')$ 后，可知每个 $x_l$ 都包含在开闭邻域 $W_l$ 中，
使得包含态射 $W_l \to X \times_Y V$ 经 $U_i \to X \times_Y V$ 分解
（对某个 $i$）。把 $W_l$ 替换为 $W_l \cap X_l$ 后，可知这些开闭集彼此不交，
而且 $\{x_1, \ldots, x_n\} \subset W_1 \cup \ldots \cup W_n$。
由于 $X \times_Y V \to V$ 有限，可以缩小 $V$ 并假设
$X \times_Y V = W_1 \amalg \ldots \amalg W_n$，正是所需。
\end{proof}

\begin{lemma}
\label{etale-cohomology-lemma-integral-B}
设 $f : X \to Y$ 是整概形态射。则性质 (B) 成立。
\end{lemma}

\begin{proof}
考虑 \'etale 态射 $V \to Y$、\'etale 覆盖 $\{U_i \to X \times_Y V\}$，以及
$v \in V$。需要找到引理 \ref{etale-cohomology-lemma-simplify-B} 中的 $V' \to V$、分解和映射。
可以缩小 $V$ 与 $Y$，因而可以假设 $V$ 与 $Y$ 仿射。由于 $X$ 在 $Y$ 上整，
这也蕴含 $X$ 与 $X \times_Y V$ 仿射。可以加细覆盖
$\{U_i \to X \times_Y V\}$，因而可以假设
$\{U_i \to X \times_Y V\}_{i = 1, \ldots, n}$ 是标准 \'etale 覆盖。写成
$Y = \Spec(A)$、$X = \Spec(B)$、$V = \Spec(C)$ 和 $U_i = \Spec(B_i)$。
那么 $A \to B$ 是整环同态，而 $B \otimes_A C \to B_i$ 是 \'etale 环同态。
由 Algebra，引理 \ref{algebra-lemma-etale}，可以找到有限 $A$-子代数
$B' \subset B$ 和 \'etale 环同态 $B' \otimes_A C \to B'_i$
（其中 $i = 1, \ldots, n$），使得 $B_i = B \otimes_{B'} B'_i$。
因此，问题归结为 \'etale 覆盖
$\{\Spec(B'_i) \to X' \times_Y V\}_{i = 1, \ldots, n}$，其中
$X' = \Spec(B')$ 在 $Y$ 上有限。在这种情形，结论由引理 \ref{etale-cohomology-lemma-finite-B} 得出。
\end{proof}

\begin{lemma}
\label{etale-cohomology-lemma-what-integral}
设 $f : X \to Y$ 是概形态射。假设 $f$ 是整态射（例如有限态射）。则
\begin{enumerate}
\item $f_{small, *}$ 把满射变为满射（对集合层和 Abel 层都是如此）；
\item $f_{small}^{-1}f_{small, *}\mathcal{F} \to \mathcal{F}$
对任意 Abel 层 $\mathcal{F}$（在 $X_\etale$ 上）都是满射；
\item
$f_{small, *} :
\textit{Ab}(X_\etale)
\to
\textit{Ab}(Y_\etale)$
是忠实的，并反映单射和满射；
\item
$f_{small, *} :
\textit{Ab}(X_\etale)
\to
\textit{Ab}(Y_\etale)$
正合。
\end{enumerate}
\end{lemma}

\begin{proof}
(2)、(3) 已在引理 \ref{etale-cohomology-lemma-when-push-pull-surjective} 中看到。
(1) 由引理 \ref{etale-cohomology-lemma-integral-B} 和 \ref{etale-cohomology-lemma-property-B-implies} 得出。
(4) 是 (1) 的推论；见 Modules on Sites，引理
\ref{sites-modules-lemma-exactness}。
\end{proof}








\section{性质 (C)}
\label{etale-cohomology-section-C}

\noindent
性质 (C) 的定义见第 \ref{etale-cohomology-section-monomorphisms} 节。

\begin{lemma}
\label{etale-cohomology-lemma-property-C-implies}
设 $f : X \to Y$ 是概形态射。假设 (C) 成立。则函子
$f_{small, *} :
\Sh(X_\etale)
\to
\Sh(Y_\etale)$
反映单射和满射。
\end{lemma}

\begin{proof}
这由 Sites，引理 \ref{sites-lemma-cover-from-below} 得出。略去如下验证：
性质 (C) 蕴含函子
$Y_\etale \to X_\etale$, $V \mapsto X \times_Y V$
满足 Sites，引理 \ref{sites-lemma-cover-from-below} 的假设。
\end{proof}

\begin{remark}
\label{etale-cohomology-remark-property-C-strong}
若 $f : X \to Y$ 是开浸入，则性质 (C) 成立。事实上，若
$U \in \Ob(X_\etale)$，则也可以把 $U$ 看作 $Y_\etale$ 的对象，并且
$U \times_Y X = U$。因此，性质 (C) 并不蕴含 $f_{small, *}$ 正合，
因为开浸入一般不具有该正合性。
\end{remark}

\begin{lemma}
\label{etale-cohomology-lemma-property-C-closed-implies}
设 $f : X \to Y$ 是概形态射。假设对任意 \'etale 态射 $V \to Y$，都有
\begin{enumerate}
\item $X \times_Y V \to V$ 具有性质 (C)；
\item $X \times_Y V \to V$ 是闭映射。
\end{enumerate}
则函子
$Y_\etale \to X_\etale$, $V \mapsto X \times_Y V$
几乎余连续；见 Sites，定义 \ref{sites-definition-almost-cocontinuous}。
\end{lemma}

\begin{proof}
设 $V \to Y$ 是 $Y_\etale$ 的对象，并设
$\{U_i \to X \times_Y V\}_{i \in I}$ 是 $X_\etale$ 的覆盖。由假设 (1)，
对每个 $i$，可以找到 \'etale 态射 $h_i : V_i \to V$ 和满态射
$X \times_Y V_i \to U_i$（在 $X \times_Y V$ 上）。注意
$\bigcup h_i(V_i) \subset V$ 是包含闭集
$Z = \Im(X \times_Y V \to V)$ 的开集。令
$h_0 : V_0 = V \setminus Z \to V$ 为开浸入。显然，
$\{V_i \to V\}_{i \in I \cup \{0\}}$ 是 \'etale 覆盖，使得对每个
$i \in I \cup \{0\}$，或者 $V_i \times_Y X = \emptyset$（即 $i = 0$），
或者 $V_i \times_Y X \to V \times_Y X$ 经 $U_i \to X \times_Y V$ 分解
（若 $i \not = 0$）。因此函子 $Y_\etale \to X_\etale$ 几乎余连续。
\end{proof}

\begin{lemma}
\label{etale-cohomology-lemma-integral-homeo-onto-image-C}
设 $f : X \to Y$ 是整概形态射，它给出 $X$ 与 $Y$ 的一个闭子集之间的同胚。
则性质 (C) 成立。
\end{lemma}

\begin{proof}
设 $g : U \to X$ 是 \'etale 态射。需要找到一个对象 $V \to Y$（属于
$Y_\etale$）和满态射 $X \times_Y V \to U$（在 $X$ 上）。假设对每个
$u \in U$，都能找到一个对象 $V_u \to Y$（属于 $Y_\etale$）和态射
$h_u : X \times_Y V_u \to U$（在 $X$ 上），且 $u \in \Im(h_u)$。
那么可以取 $V = \coprod V_u$ 和 $h = \coprod h_u$，从而得证。因此，给定一点
$u \in U$，我们来寻找如上的一对 $(V_u, h_u)$。为此可以缩小 $U$ 并假设
$U$ 仿射。在这种情形，$g : U \to X$ 局部拟有限。设
$g^{-1}(g(\{u\})) = \{u, u_2, \ldots, u_n\}$。由于不存在特殊化
$u_i \leadsto u$，可以用一个仿射邻域替换 $U$，使得
$g^{-1}(g(\{u\})) = \{u\}$。

\medskip\noindent
像 $g(U) \subset X$ 是开集，故 $f(g(U))$ 在 $Y$ 中局部闭。选择开集
$V \subset Y$，使得 $f(g(U)) = f(X) \cap V$。于是 $g$ 经
$X \times_Y V$ 分解，并且所得的 $\{U \to X \times_Y V\}$ 是 \'etale 覆盖。
由于 $f$ 具有性质 (B)，见引理 \ref{etale-cohomology-lemma-integral-B}，存在 \'etale 覆盖
$\{V_j \to V\}$，使得 $X \times_Y V_j \to X \times_Y V$ 经 $U$ 分解。
这蕴含 $V' = \coprod V_j$ 在 $Y$ 上是 \'etale 的，并且存在态射
$h : X \times_Y V' \to U$，其像满射到 $g(U)$。由于 $u$ 是其纤维中唯一的点，
它必属于 $h$ 的像，从而得证。
\end{proof}

\noindent
我们建议读者把下面的引理看作通往整且普遍单射态射的 $f_{small, *}$ 之最终真相的
旅途中的驿站\footnote{驿站是人们在长途旅行中停下来进食和休息的地方。}。

\begin{lemma}
\label{etale-cohomology-lemma-integral-universally-injective}
设 $f : X \to Y$ 是概形态射。假设 $f$ 普遍单射且整（例如闭浸入）。则
\begin{enumerate}
\item
$f_{small, *} :
\Sh(X_\etale)
\to
\Sh(Y_\etale)$
反映单射和满射；
\item
$f_{small, *} :
\Sh(X_\etale)
\to
\Sh(Y_\etale)$
与推出和余等化子可交换（更一般地，与有限连通余极限可交换）；
\item $f_{small, *}$ 把满射变为满射（对集合层和 Abel 层都是如此）；
\item 映射
$f_{small}^{-1}f_{small, *}\mathcal{F} \to \mathcal{F}$
对任意层 $\mathcal{F}$（集合层或 Abel 群层，在 $X_\etale$ 上）都是满射；
\item 函子 $f_{small, *}$ 是忠实的（对集合层和 Abel 层都是如此）；
\item
$f_{small, *} :
\textit{Ab}(X_\etale)
\to
\textit{Ab}(Y_\etale)$
正合；
\item 函子 $Y_\etale \to X_\etale$、$V \mapsto X \times_Y V$ 几乎余连续。
\end{enumerate}
\end{lemma}

\begin{proof}
由引理 \ref{etale-cohomology-lemma-integral-A}、\ref{etale-cohomology-lemma-integral-B} 和
\ref{etale-cohomology-lemma-integral-homeo-onto-image-C}，可知态射 $f$ 具有性质 (A)、(B) 和 (C)。
此外，由引理 \ref{etale-cohomology-lemma-property-C-closed-implies}，可知函子
$Y_\etale \to X_\etale$ 几乎余连续。现在有
\begin{enumerate}
\item 由引理 \ref{etale-cohomology-lemma-property-C-implies}，性质 (C) 蕴含 (1)；
\item 由 Sites，引理 \ref{sites-lemma-morphism-of-sites-almost-cocontinuous}，
几乎余连续蕴含 (2)；
\item 由引理 \ref{etale-cohomology-lemma-property-B-implies}，性质 (B) 蕴含 (3)。
\end{enumerate}
性质 (4)、(5)、(6) 从前三项形式地得出；见 Sites，引理
\ref{sites-lemma-exactness-properties} 和 Modules on Sites，引理
\ref{sites-modules-lemma-exactness}。性质 (7) 已在上面看到。
\end{proof}




\section{小 \'etale 网站的拓扑不变性}
\label{etale-cohomology-section-topological-invariance}

\noindent
下一个定理表明，小 \'etale 网站在如下意义下是拓扑不变量：若
$f : X \to Y$ 是作为普遍同胚的概形态射，则网站
$X_\etale \cong Y_\etale$。这改进了 \'Etale Morphisms，定理
\ref{etale-theorem-remarkable-equivalence} 的结果。我们先对态射证明该结论，
然后陈述范畴层面的结论。

\begin{theorem}
\label{etale-cohomology-theorem-etale-topological}
设 $X$ 和 $Y$ 是基概形 $S$ 上的两个概形。设 $S' \to S$ 是普遍同胚。
以 $X'$（相应地以 $Y'$）表示到 $S'$ 的基变换。若 $X$ 在 $S$ 上是 \'etale 的，
则映射
$$
\Mor_S(Y, X) \longrightarrow \Mor_{S'}(Y', X')
$$
是双射。
\end{theorem}

\begin{proof}
沿 $Y \to S$ 作基变换后，可以假设 $Y = S$。因此，可以并且确实假设
$X$ 与 $Y$ 都在 $S$ 上是 \'etale 的。换言之，定理断言：基变换函子从
$S$ 上 \'etale 概形的范畴到 $S'$ 上 \'etale 概形的范畴是全忠实函子。

\medskip\noindent
考虑遗忘函子
\begin{equation}
\label{etale-cohomology-equation-descent-etale-forget}
\begin{matrix}
\text{关于 }S'/S\text{ 的下降数据 }(X', \varphi') \\
\text{，其中 }X'\text{ 在 }S'\text{ 上 \'etale}
\end{matrix}
\longrightarrow
\text{在 }S'\text{ 上 \'etale 的概形 }X'
\end{equation}
我们断言这个函子是等价。另一方面，函子
\begin{equation}
\label{etale-cohomology-equation-descent-etale}
\text{在 }S\text{ 上 \'etale 的概形 }X \longrightarrow
\begin{matrix}
\text{关于 }S'/S\text{ 的下降数据 }(X', \varphi') \\
\text{，其中 }X'\text{ 在 }S'\text{ 上 \'etale}
\end{matrix}
\end{equation}
由 \'Etale Morphisms，引理 \ref{etale-lemma-fully-faithful-cases} 是全忠实的。
因此，该断言蕴含本定理。

\medskip\noindent
证明该断言。回忆，普遍同胚等价于整、普遍单射且满的态射；见 Morphisms，引理
\ref{morphisms-lemma-universal-homeomorphism}。特别地，由 Morphisms，引理
\ref{morphisms-lemma-universally-injective}，对角态射
$\Delta : S' \to S' \times_S S'$ 是加厚。因此，由 \'Etale Morphisms，定理
\ref{etale-theorem-etale-topological}，可知给定 \'etale 态射 $X' \to S'$，
存在唯一同构
$$
\varphi' : X' \times_S S' \to S' \times_S X'
$$
，它是 $S' \times_S S'$ 上 \'etale 概形的同构，并且沿 $\Delta$ 拉回后成为
$\text{id} : X' \to X'$（在 $S'$ 上）。由于
$S' \to S' \times_S S' \times_S S'$ 也是加厚（它是双射且为闭浸入），
可知 $(X', \varphi')$ 是相对于 $S'/S$ 的下降数据。$\varphi'$ 构造的典范性表明，
它与 $S'$ 上 \'etale 概形之间的态射相容。换言之，得到函子
(\ref{etale-cohomology-equation-descent-etale-forget}) 的拟逆 $X' \mapsto (X', \varphi')$。
这证明了断言，并完成了定理的证明。
\end{proof}

\begin{theorem}
\label{etale-cohomology-theorem-topological-invariance}
\begin{reference}
\cite[IV Theorem 18.1.2]{EGA}
\end{reference}
设 $f : X \to Y$ 是概形态射。假设 $f$ 整、普遍单射且满
（即 $f$ 是普遍同胚；见 Morphisms，引理
\ref{morphisms-lemma-universal-homeomorphism}）。函子
$$
V \longmapsto V_X = X \times_Y V
$$
给出范畴等价
$$
\{
\text{在 }Y\text{ 上 \'etale 的概形 }V
\}
\leftrightarrow
\{
\text{在 }X\text{ 上 \'etale 的概形 }U
\}
$$
\end{theorem}

\noindent
我们给出两个证明。第一个证明使用拟紧、分离、\'etale 态射相对于满整态射的下降
有效性。第二个证明使用本章前面讨论的性质 (A)、(B)、(C)。

\begin{proof}[第一个证明]
由定理 \ref{etale-cohomology-theorem-etale-topological}，可知该函子全忠实。还需证明它本质满。
设 $U \to X$ 是概形的 \'etale 态射。

\medskip\noindent
假设 $U$ 与 $Y$ 仿射时结论成立。在这种情形，选择仿射开覆盖
$U = \bigcup U_i$，使每个 $U_i$ 映入 $Y$ 的一个仿射开集。根据假设
（仿射情形），可以找到 \'etale 态射 $V_i \to Y$，使得
$X \times_Y V_i \cong U_i$（作为 $X$ 上的概形）。令 $V_{i, i'} \subset V_i$ 是开子概形，
其底拓扑空间对应于 $U_i \cap U_{i'}$。由于有同构
$$
X \times_Y V_{i, i'} \cong U_i \cap U_{i'} \cong X \times_Y V_{i', i}
$$
（作为 $X$ 上的概形），由全忠实性得到同构
$\theta_{i, i'} : V_{i, i'} \to V_{i', i}$（它是 $Y$ 上概形的同构）。略去这些同构满足 Schemes，第
\ref{schemes-section-glueing-schemes} 节的余圈条件的验证。应用 Schemes，引理
\ref{schemes-lemma-glue-schemes}，得到概形 $V \to Y$；它是把概形 $V_i$ 沿同一
$\theta_{i, i'}$ 粘合得到的。根据构造，$V \to Y$ 显然是 \'etale 的，且
$X \times_Y V \cong U$。

\medskip\noindent
因此，只需在 $U$ 与 $Y$ 仿射时证明该引理。回忆，在定理
\ref{etale-cohomology-theorem-etale-topological} 的证明中，我们已证明 $U$ 带有唯一下降数据
$(U, \varphi)$（相对于 $X/Y$）。由 \'Etale Morphisms，命题
\ref{etale-proposition-effective}（它适用是因为通过归约到仿射情形，
$U \to X$ 不仅是 \'etale 的，而且拟紧且分离），存在 \'etale 态射
$V \to Y$，使得 $X \times_Y V \cong U$，证明完成。
\end{proof}

\begin{proof}[第二个证明]
由定理 \ref{etale-cohomology-theorem-etale-topological}，可知该函子全忠实。还需证明它本质满。
设 $U \to X$ 是概形的 \'etale 态射。

\medskip\noindent
假设 $U$ 与 $Y$ 仿射时结论成立。在这种情形，选择仿射开覆盖
$U = \bigcup U_i$，使每个 $U_i$ 映入 $Y$ 的一个仿射开集。根据假设
（仿射情形），可以找到 \'etale 态射 $V_i \to Y$，使得
$X \times_Y V_i \cong U_i$（作为 $X$ 上的概形）。令
$V_{i, i'} \subset V_i$ 是开子概形，其底拓扑空间对应于
$U_i \cap U_{i'}$。由于有同构
$$
X \times_Y V_{i, i'} \cong U_i \cap U_{i'} \cong X \times_Y V_{i', i}
$$
（作为 $X$ 上的概形），由全忠实性得到同构
$\theta_{i, i'} : V_{i, i'} \to V_{i', i}$（它是 $Y$ 上概形的同构）。
略去这些同构满足 Schemes，第 \ref{schemes-section-glueing-schemes} 节的余圈条件
的验证。应用 Schemes，引理 \ref{schemes-lemma-glue-schemes}，得到概形
$V \to Y$；它是把概形 $V_i$ 沿同一 $\theta_{i, i'}$ 粘合得到的。根据构造，
$V \to Y$ 显然是 \'etale 的，且 $X \times_Y V \cong U$。

\medskip\noindent
因此，只需在 $X$ 与 $Y$ 仿射时证明函子
\begin{equation}
\label{etale-cohomology-equation-affine-etale}
\{
\text{在 }Y\text{ 上 \'etale 的仿射概形 }V
\}
\leftrightarrow
\{
\text{在 }X\text{ 上 \'etale 的仿射概形 }U
\}
\end{equation}
本质满。

\medskip\noindent
设 $U \to X$ 是在 $X$ 上 \'etale 的仿射概形。需要找到 \'etale（且仿射）的
$V \to Y$，使得 $X \times_Y V$ 同构于 $U$（在 $X$ 上）。注意，仿射概形间的
\'etale 态射具有普遍有界的纤维；见 Morphisms，引理
\ref{morphisms-lemma-etale-locally-quasi-finite} 和
\ref{morphisms-lemma-locally-quasi-finite-qc-source-universally-bounded}。
因此，可以对整数 $n$ 作归纳，它限制 $U \to X$ 之纤维的次数。\'etale 态射情形下
这个整数的描述见 Morphisms，引理
\ref{morphisms-lemma-etale-universally-bounded}。若 $n = 1$，则
$U \to X$ 是开浸入（见 \'Etale Morphisms，定理
\ref{etale-theorem-etale-radicial-open}），结论显然。假设 $n > 1$。

\medskip\noindent
由引理 \ref{etale-cohomology-lemma-integral-homeo-onto-image-C}，存在概形的 \'etale 态射
$W \to Y$ 和满态射 $W_X \to U$（在 $X$ 上）。由于 $U$ 拟紧，
可以用 $W$ 的有限多个仿射开集的不交并替换 $W$，因而也可以假设 $W$ 仿射。
有如下图：
$$
\xymatrix{
U \ar[d] & U \times_Y W \ar[l] \ar[d] & W_X \amalg R \ar@{=}[l]\\
X \ar[d] & W_X \ar[l] \ar[d] \\
Y & W \ar[l]
}
$$
这个不交并分解之所以出现，是因为按构造，仿射概形的 \'etale 态射
$U \times_Y W \to W_X$ 有截面。现在可知，态射 $R \to X \times_Y W$ 是仿射概形的
\'etale 态射，其纤维次数普遍不超过 $n - 1$。因此，由归纳假设，存在 \'etale 概形
$V' \to W$，使得 $R \cong W_X \times_W V'$。取 $V'' = W \amalg V'$，得到概形
$V''$（它在 $W$ 上是 \'etale 的），其到 $W_X$ 的基变换同构于 $U \times_Y W$
（在 $X \times_Y W$ 上）。

\medskip\noindent
此时可以使用下降找到 $V$（在 $Y$ 上），使其到 $X$ 的基变换同构于 $U$
（在 $X$ 上）。
具体地，由与普遍同胚
$X \times_Y (W \times_Y W) \to (W \times_Y W)$ 对应的函子
(\ref{etale-cohomology-equation-affine-etale}) 的全忠实性，存在唯一同构
$\varphi : V'' \times_Y W \to W \times_Y V''$，它到
$X \times_Y (W \times_Y W)$ 的基变换是 $U \times_Y W$ 在
$X \times_Y W$ 上的典范下降数据。特别地，$\varphi$ 满足余圈条件。因此，由
Descent，引理 \ref{descent-lemma-affine}，可知 $\varphi$ 是有效的（回忆上面的概形
全都仿射）。于是得到 $V \to Y$ 和同构 $V'' \cong W \times_Y V$，使得
$W \times_Y V/W/Y$ 上的典范下降数据与 $\varphi$ 相同。注意，由 Descent，引理
\ref{descent-lemma-descending-property-etale}，$V \to Y$ 是 \'etale 的。此外，
存在同构 $V_X \cong U$，它来自下列同构的下降：
$$
V_X \times_X W_X =
X \times_Y V \times_Y W =
(X \times_Y W) \times_W (W \times_Y V) \cong
W_X  \times_W V'' \cong U \times_Y W
$$
；这个同构由构造给出。略去一些细节。
\end{proof}

\begin{remark}
\label{etale-cohomology-remark-affine-inside-equivalence}
在定理 \ref{etale-cohomology-theorem-topological-invariance} 的情形下，$V \mapsto V_X$ 还诱导
两类对象之间的等价：一类是 \'etale 态射 $V \to Y$ 且 $V$ 仿射，另一类是
\'etale 态射 $U \to X$ 且 $U$ 仿射。例如，这由 Limits，命题
\ref{limits-proposition-affine} 得出。
\end{remark}

\begin{proposition}[\'Etale 上同调的拓扑不变性]
\label{etale-cohomology-proposition-topological-invariance}
设 $X_0 \to X$ 是概形的普遍同胚（例如由幂零理想层定义的闭浸入）。则
\begin{enumerate}
\item \'etale 网站 $X_\etale$ 与 $(X_0)_\etale$ 同构；
\item \'etale 拓扑斯 $\Sh(X_\etale)$ 与 $\Sh((X_0)_\etale)$ 等价；
\item $H^q_\etale(X, \mathcal{F}) = H^q_\etale(X_0, \mathcal{F}|_{X_0})$
对所有 $q$ 以及任意 Abel 层 $\mathcal{F}$（在 $X_\etale$ 上）成立。
\end{enumerate}
\end{proposition}

\begin{proof}
范畴等价 $X_\etale \to (X_0)_\etale$ 由定理
\ref{etale-cohomology-theorem-topological-invariance} 给出。略去在此等价下 \'etale 覆盖相互对应的
证明。因此 (1) 成立，而 (2)、(3) 从 (1) 形式地得出。
\end{proof}






\section{闭浸入与推前}
\label{etale-cohomology-section-closed-immersions}

\noindent
在以正确的一般性陈述并证明命题
\ref{etale-cohomology-proposition-closed-immersion-pushforward} 之前，我们先对闭浸入简要陈述并证明它。
事实上，前面的一些论证在闭浸入情形下更容易理解，所以在此以简化形式重复它们。

\medskip\noindent
本节余下部分中，$i : Z \to X$ 是闭浸入。函子
$$
\Sch/X \longrightarrow \Sch/Z, \quad
U \longmapsto U_Z = Z \times_X U
$$
如上记为 $U \mapsto U_Z$。由于闭浸入在任意基变换下保持，概形 $U_Z$ 是
$U$ 的闭子概形。

\begin{lemma}
\label{etale-cohomology-lemma-closed-immersion-almost-full}
设 $i : Z \to X$ 是概形的闭浸入。设 $U, U'$ 是 $X$ 上的 \'etale 概形。
设 $h : U_Z \to U'_Z$ 是 $Z$ 上的态射。则存在下图
$$
\xymatrix{
U & W \ar[l]_a \ar[r]^b & U'
}
$$
，它位于 $X_\etale$ 中，使得 $a_Z : W_Z \to U_Z$ 是同构，且
$h = b_Z \circ (a_Z)^{-1}$。
\end{lemma}

\begin{proof}
考虑概形 $M = U \times_X U'$。$\Gamma_h \subset M_Z$（即 $h$ 的图像）是开集。
例如，这是因为 $\Gamma_h$ 是 \'etale 态射
$\text{pr}_{1, Z} : M_Z \to U_Z$ 的某个截面的像；见 \'Etale Morphisms，命题
\ref{etale-proposition-properties-sections}。因此，存在开子概形 $W \subset M$，
它与闭子集 $M_Z$ 的交是 $\Gamma_h$。置 $a = \text{pr}_1|_W$ 和
$b = \text{pr}_2|_W$。
\end{proof}

\begin{lemma}
\label{etale-cohomology-lemma-closed-immersion-almost-essentially-surjective}
设 $i : Z \to X$ 是概形的闭浸入。设 $V \to Z$ 是概形的 \'etale 态射。
存在 \'etale 态射 $U_i \to X$ 和态射 $U_{i, Z} \to V$，使得
$\{U_{i, Z} \to V\}$ 是 $V$ 的 Zariski 覆盖。
\end{lemma}

\begin{proof}
由于只需找到 $V$ 的 Zariski 覆盖，其成员是形如 $U_Z$ 的概形，且 $U$ 在
$X$ 上是 \'etale 的，因此可以在 $X$ 和 $V$ 上作 Zariski 局部化。因此可以假设
$X$ 与 $V$ 仿射。在仿射情形，这就是 Algebra，引理
\ref{algebra-lemma-lift-etale}。
\end{proof}

\noindent
若 $\overline{x} : \Spec(k) \to X$ 是 $X$ 的几何点，则或者 $\overline{x}$
（唯一地）经闭子概形 $Z$ 分解，或者 $Z_{\overline{x}} = \emptyset$。若
$\overline{x}$ 经 $Z$ 分解，就称 $\overline{x}$ 为 $Z$ 的几何点（它确实如此），
并用记号``$\overline{x} \in Z$''表示这一点。

\begin{lemma}
\label{etale-cohomology-lemma-stalk-pushforward-closed-immersion}
设 $i : Z \to X$ 是概形的闭浸入。设 $\mathcal{G}$ 是 $Z_\etale$ 上的集合层。
设 $\overline{x}$ 是 $X$ 的几何点。则
$$
(i_{small, *}\mathcal{G})_{\overline{x}} =
\left\{
\begin{matrix}
* & \text{若} & \overline{x} \not \in Z \\
\mathcal{G}_{\overline{x}} & \text{若} & \overline{x} \in Z
\end{matrix}
\right.
$$
其中 $*$ 表示单点集。
\end{lemma}

\begin{proof}
令 $U = X \setminus Z$。注意，$i_{small, *}\mathcal{G}|_{U_\etale} = *$ 是
$U$ 上 \'etale 层范畴中的终对象，即给每个在 $U$ 上 \'etale 的概形赋予单点集的层。
这说明了 $(i_{small, *}\mathcal{G})_{\overline{x}}$ 的取值
（当 $\overline{x} \not \in Z$ 时）。

\medskip\noindent
接着，假设 $\overline{x} \in Z$。注意
$$
(i_{small, *}\mathcal{G})_{\overline{x}}
=
\colim_{(U, \overline{u})} \mathcal{G}(U_Z)
$$
，另一方面
$$
\mathcal{G}_{\overline{x}}
=
\colim_{(V, \overline{v})} \mathcal{G}(V).
$$
令 $\mathcal{C}_1 = \{(U, \overline{u})\}^{opp}$ 为 $\overline{x}$ 在 $X$ 中的
\'etale 邻域范畴的反范畴，并令
$\mathcal{C}_2 = \{(V, \overline{v})\}^{opp}$ 为 $\overline{x}$ 在 $Z$ 中的
\'etale 邻域范畴的反范畴。典范映射
$$
\mathcal{G}_{\overline{x}}
\longrightarrow
(i_{small, *}\mathcal{G})_{\overline{x}}
$$
对应于函子 $F : \mathcal{C}_1 \to \mathcal{C}_2$，其中
$F(U, \overline{u}) = (U_Z, \overline{x})$。现在，引理
\ref{etale-cohomology-lemma-closed-immersion-almost-essentially-surjective} 和
\ref{etale-cohomology-lemma-closed-immersion-almost-full} 蕴含 $\mathcal{C}_1$ 在
$\mathcal{C}_2$ 中余终；见 Categories，定义 \ref{categories-definition-cofinal}。
因此，上面的箭头是同构；见 Categories，引理 \ref{categories-lemma-cofinal}。
\end{proof}

\begin{proposition}
\label{etale-cohomology-proposition-closed-immersion-pushforward}
设 $i : Z \to X$ 是概形的闭浸入。
\begin{enumerate}
\item 函子
$$
i_{small, *} :
\Sh(Z_\etale)
\longrightarrow
\Sh(X_\etale)
$$
是全忠实的，其本质像由如下集合层 $\mathcal{F}$ 构成：它在 $X_\etale$ 上，
且到 $X \setminus Z$ 的限制同构于 $*$；
\item 函子
$$
i_{small, *} :
\textit{Ab}(Z_\etale)
\longrightarrow
\textit{Ab}(X_\etale)
$$
是全忠实的，其本质像由 $X_\etale$ 上且支集包含于 $Z$ 的 Abel 层构成。
\end{enumerate}
在两种情形下，$i_{small}^{-1}$ 都是函子 $i_{small, *}$ 的左逆。
\end{proposition}

\begin{proof}
先讨论集合层的情形。对任意层 $\mathcal{G}$（定义于 $Z$ 上），态射
$i_{small}^{-1}i_{small, *}\mathcal{G} \to \mathcal{G}$
由引理 \ref{etale-cohomology-lemma-stalk-pushforward-closed-immersion}（以及定理
\ref{etale-cohomology-theorem-exactness-stalks}）是同构。这形式地蕴含 $i_{small, *}$ 全忠实；见
Sites，引理 \ref{sites-lemma-exactness-properties}。显然，
$i_{small, *}\mathcal{G}|_{U_\etale} \cong *$，其中 $U = X \setminus Z$。
反之，设 $\mathcal{F}$ 是 $X$ 上的集合层，且
$\mathcal{F}|_{U_\etale} \cong *$。考虑伴随映射
$$
\mathcal{F} \longrightarrow i_{small, *}i_{small}^{-1}\mathcal{F}
$$
合用引理 \ref{etale-cohomology-lemma-stalk-pushforward-closed-immersion} 和
\ref{etale-cohomology-lemma-stalk-pullback}，可知它是同构。这完成了 (1) 的证明。
(2) 的证明相同。
\end{proof}





\section{整普遍单射态射}
\label{etale-cohomology-section-integral-universally-injective}

\noindent
下面给出命题 \ref{etale-cohomology-proposition-closed-immersion-pushforward} 的一般版本。

\begin{proposition}
\label{etale-cohomology-proposition-integral-universally-injective-pushforward}
设 $f : X \to Y$ 是整且普遍单射的概形态射。
\begin{enumerate}
\item 函子
$$
f_{small, *} :
\Sh(X_\etale)
\longrightarrow
\Sh(Y_\etale)
$$
是全忠实的，其本质像由集合层 $\mathcal{F}$ 构成：它定义于 $Y_\etale$ 上，
且到 $Y \setminus f(X)$ 的限制同构于 $*$；
\item 函子
$$
f_{small, *} :
\textit{Ab}(X_\etale)
\longrightarrow
\textit{Ab}(Y_\etale)
$$
是全忠实的，其本质像由 $Y_\etale$ 上且支集包含于 $f(X)$ 的 Abel 层构成。
\end{enumerate}
在两种情形下，$f_{small}^{-1}$ 都是函子 $f_{small, *}$ 的左逆。
\end{proposition}

\begin{proof}
可以把 $f$ 分解为
$$
\xymatrix{
X \ar[r]^h & Z \ar[r]^i & Y
}
$$
，其中 $h$ 整、普遍单射且满，而 $i : Z \to Y$ 是闭浸入。对 $i$ 应用命题
\ref{etale-cohomology-proposition-closed-immersion-pushforward}，再对 $h$ 应用定理
\ref{etale-cohomology-theorem-topological-invariance}。
\end{proof}









\section{大网站与推前}
\label{etale-cohomology-section-big}

\noindent
本节对若干类概形态射证明关于 $f_{big, *}$ 的一些技术性结果。

\begin{lemma}
\label{etale-cohomology-lemma-monomorphism-big-push-pull}
设 $\tau \in \{Zariski, \etale, smooth, syntomic, fppf\}$。
设 $f : X \to Y$ 是概形的单态射。则典范映射
$f_{big}^{-1}f_{big, *}\mathcal{F} \to \mathcal{F}$ 是同构；这里
$\mathcal{F}$ 是定义于 $(\Sch/X)_\tau$ 上的任意层。
\end{lemma}

\begin{proof}
在这种情形下，函子 $(\Sch/X)_\tau \to (\Sch/Y)_\tau$ 连续、余连续且全忠实。
因此，结论由 Sites，引理 \ref{sites-lemma-back-and-forth} 得到。
\end{proof}

\begin{remark}
\label{etale-cohomology-remark-push-pull-shriek}
在引理 \ref{etale-cohomology-lemma-monomorphism-big-push-pull} 的情形下，典范映射
$\mathcal{F} \to f_{big}^{-1}f_{big!}\mathcal{F}$ 是同构；这里
$\mathcal{F}$ 是定义于 $(\Sch/X)_\tau$ 上的任意集合层。证明相同。
这对 Abel 群层也成立。不过要注意，Abel 群层的函子 $f_{big!}$ 定义于
Modules on Sites，第 \ref{sites-modules-section-exactness-lower-shriek} 节，
且一般不同于集合层上的 $f_{big!}$。Abel 群层的结论由
Modules on Sites，引理 \ref{sites-modules-lemma-back-and-forth} 得到。
\end{remark}

\begin{lemma}
\label{etale-cohomology-lemma-closed-immersion-cover-from-below}
设 $f : X \to Y$ 是概形的闭浸入。设 $U \to X$ 是 syntomic（分别为光滑、
\'etale）态射。则存在 syntomic（分别为光滑、\'etale）态射 $V_i \to Y$ 以及态射
$V_i \times_Y X \to U$，使得 $\{V_i \times_Y X \to U\}$ 是 $U$ 的 Zariski 覆盖。
\end{lemma}

\begin{proof}
先证明 $\tau = syntomic$ 时的引理。问题在 $U$ 上是局部的。因此，可以假设
$U$ 是映入 $Y$ 的某个仿射开子概形的仿射概形。于是，归结为证明下列情形：
$Y = \Spec(A)$、$X = \Spec(A/I)$ 且 $U = \Spec(\overline{B})$，其中
$A/I \to \overline{B}$ 是 syntomic 环同态。由 Algebra，引理
\ref{algebra-lemma-lift-syntomic}，可以找到元素
$\overline{g}_i \in \overline{B}$，使得 $\overline{B}_{\overline{g}_i} = A_i/IA_i$
对某些 syntomic 环同态 $A \to A_i$ 成立。这证明了 syntomic 情形。
光滑情形的证明相同，只是改用 Algebra，引理 \ref{algebra-lemma-lift-smooth}。
\'etale 情形使用 Algebra，引理 \ref{algebra-lemma-lift-etale}。
\end{proof}

\begin{lemma}
\label{etale-cohomology-lemma-prepare-closed-immersion-almost-cocontinuous}
设 $f : X \to Y$ 是概形的闭浸入。设 $\{U_i \to X\}$ 是 syntomic（分别为光滑、
\'etale）覆盖。则存在 syntomic（分别为光滑、\'etale）覆盖 $\{V_j \to Y\}$，
使得对每个 $j$，或者 $V_j \times_Y X = \emptyset$，或者态射
$V_j \times_Y X \to X$ 通过 $U_i$ 分解，其中 $i$ 可以依赖于所给覆盖成员。
\end{lemma}

\begin{proof}
对每个 $i$，可以选取 syntomic（分别为光滑、\'etale）态射
$g_{ij} : V_{ij} \to Y$ 以及态射 $V_{ij} \times_Y X \to U_i$（在 $X$ 上），
使得 $\{V_{ij} \times_Y X \to U_i\}$ 是 Zariski 覆盖；见引理
\ref{etale-cohomology-lemma-closed-immersion-cover-from-below}。这特别蕴含
$\bigcup_{ij} g_{ij}(V_{ij})$ 包含闭子集 $f(X)$。因此，syntomic（分别为光滑、
\'etale）映射族 $g_{ij}$ 与开浸入 $Y \setminus f(X) \to Y$ 一起构成所需的
$Y$ 的 syntomic（分别为光滑、\'etale）覆盖。
\end{proof}

\begin{lemma}
\label{etale-cohomology-lemma-closed-immersion-almost-cocontinuous}
设 $f : X \to Y$ 是概形的闭浸入。设 $\tau \in \{syntomic, smooth, \etale\}$。
函子 $V \mapsto X \times_Y V$ 给出几乎余连续函子
（见 Sites，定义 \ref{sites-definition-almost-cocontinuous}）
$(\Sch/Y)_\tau \to (\Sch/X)_\tau$；这是大 $\tau$ 网站之间的函子。
\end{lemma}

\begin{proof}
需要证明如下事实：给定态射 $V \to Y$ 和任意 syntomic（分别为光滑、\'etale）
覆盖 $\{U_i \to X \times_Y V\}$，存在光滑（分别为光滑、\'etale）覆盖
$\{V_j \to V\}$，使得对每个 $j$，或者 $X \times_Y V_j$ 为空，或者
$X \times_Y V_j \to Z \times_Y V$ 通过某个 $U_i$ 分解。这可由对闭浸入
$X \times_Y V \to V$ 应用上面的引理
\ref{etale-cohomology-lemma-prepare-closed-immersion-almost-cocontinuous} 得到。
\end{proof}

\begin{lemma}
\label{etale-cohomology-lemma-closed-immersion-pushforward-exact}
设 $f : X \to Y$ 是概形的闭浸入。设 $\tau \in \{syntomic, smooth, \etale\}$。
\begin{enumerate}
\item 推前
$f_{big, *} :
\Sh((\Sch/X)_\tau)
\to
\Sh((\Sch/Y)_\tau)$
与余等化子及推出交换。
\item 推前
$f_{big, *} :
\textit{Ab}((\Sch/X)_\tau)
\to
\textit{Ab}((\Sch/Y)_\tau)$
是正合的。
\end{enumerate}
\end{lemma}

\begin{proof}
这由 Sites，引理 \ref{sites-lemma-morphism-of-sites-almost-cocontinuous}，
Modules on Sites，引理
\ref{sites-modules-lemma-morphism-ringed-sites-almost-cocontinuous}，以及上面的引理
\ref{etale-cohomology-lemma-closed-immersion-almost-cocontinuous} 得到。
\end{proof}

\begin{remark}
\label{etale-cohomology-remark-fppf-closed-immersion-not-closed}
引理 \ref{etale-cohomology-lemma-closed-immersion-pushforward-exact} 中缺少 $\tau = fppf$ 的情形。
原因是，给定环 $A$、理想 $I$ 以及忠实平坦有限表示的环同态
$A/I \to \overline{B}$，没有理由认为可以找到{\it 任意}平坦有限表示环同态
$A \to B$，它满足 $B/IB \not = 0$，并使得 $A/I \to B/IB$ 通过
$\overline{B}$ 分解。
因此，引理 \ref{etale-cohomology-lemma-closed-immersion-almost-cocontinuous} 的证明不适用于
fppf 拓扑。事实上，
$f_{big, *} : \textit{Ab}((\Sch/X)_{fppf})
\to \textit{Ab}((\Sch/Y)_{fppf})$
在 $f$ 是闭浸入时为正合的说法很可能是假的。如果知道一个例子，请发送电子邮件至
\href{mailto:stacks.project@gmail.com}{stacks.project@gmail.com}。
\end{remark}












\section{大网站下叹号函子的正合性}
\label{etale-cohomology-section-exactness-lower-shriek}

\noindent
这里只需下列技术性结果。注意，一般而言，函子 $f_{big!}$ 与紧支集上同调毫无关系。

\begin{lemma}
\label{etale-cohomology-lemma-exactness-lower-shriek}
设 $\tau \in \{Zariski, \etale, smooth, syntomic, fppf\}$。
设 $f : X \to Y$ 是概形态射。令
$$
f_{big} :
\Sh((\Sch/X)_\tau)
\longrightarrow
\Sh((\Sch/Y)_\tau)
$$
为 Topologies，引理 \ref{topologies-lemma-morphism-big}、
\ref{topologies-lemma-morphism-big-etale},
\ref{topologies-lemma-morphism-big-smooth},
\ref{topologies-lemma-morphism-big-syntomic} 或
\ref{topologies-lemma-morphism-big-fppf} 中相应的拓扑斯态射。
\begin{enumerate}
\item 函子
$f_{big}^{-1} : \textit{Ab}((\Sch/Y)_\tau) \to \textit{Ab}((\Sch/X)_\tau)$
有左伴随
$$
f_{big!} : \textit{Ab}((\Sch/X)_\tau) \to \textit{Ab}((\Sch/Y)_\tau)
$$
，并且后者正合。
\item 函子
$f_{big}^* :
\textit{Mod}((\Sch/Y)_\tau, \mathcal{O})
\to
\textit{Mod}((\Sch/X)_\tau, \mathcal{O})$
有左伴随
$$
f_{big!} :
\textit{Mod}((\Sch/X)_\tau, \mathcal{O})
\to
\textit{Mod}((\Sch/Y)_\tau, \mathcal{O})
$$
，并且后者正合。
\end{enumerate}
此外，这两个函子 $f_{big!}$ 在底层 Abel 群层上相同。
\end{lemma}

\begin{proof}
回忆，$f_{big}$ 是由连续且余连续的函子
$u : (\Sch/X)_\tau \to (\Sch/Y)_\tau$、$U/X \mapsto U/Y$ 所联系的拓扑斯态射。
此外，有 $f_{big}^{-1}\mathcal{O} = \mathcal{O}$。因此，$f_{big!}$ 的存在性分别由
Modules on Sites，引理 \ref{sites-modules-lemma-g-shriek-adjoint} 和
Modules on Sites，引理 \ref{sites-modules-lemma-lower-shriek-modules} 得到。
注意，若 $U$ 是 $(\Sch/X)_\tau$ 的对象，则函子 $u$ 诱导范畴等价
$$
u' :
(\Sch/X)_\tau/U
\longrightarrow
(\Sch/Y)_\tau/U
$$
，因为箭头两边都等于 $(\Sch/U)_\tau$。因此，$f_{big!}$ 在底层 Abel 层上相同由
Modules on Sites，注 \ref{sites-modules-remark-when-shriek-equal} 中的讨论得到。
$f_{big!}$ 的正合性由 Modules on Sites，引理
\ref{sites-modules-lemma-exactness-lower-shriek} 得到，因为上面的函子 $u$
与纤维积和等化子交换。
\end{proof}

\noindent
下面证明一个技术性引理，稍后比较代数栈所联系的不同网站上的模层时将会用到它。

\begin{lemma}
\label{etale-cohomology-lemma-compare-structure-sheaves}
设 $X$ 是概形。设 $\tau \in \{Zariski, \etale, smooth, syntomic, fppf\}$。
设 $\mathcal{C}_1 \subset \mathcal{C}_2 \subset (\Sch/X)_\tau$ 是满足下列性质的
全子范畴：
\begin{enumerate}
\item 对对象 $U/X$（它属于 $\mathcal{C}_t$），
\begin{enumerate}
\item 若 $\{U_i \to U\}$ 是 $(\Sch/X)_\tau$ 的覆盖，则 $U_i/X$ 是
$\mathcal{C}_t$ 的对象；
\item $U \times \mathbf{A}^1/X$ 是 $\mathcal{C}_t$ 的对象。
\end{enumerate}
\item $X/X$ 是 $\mathcal{C}_t$ 的对象。
\end{enumerate}
赋予 $\mathcal{C}_t$ 网站结构；其覆盖恰为那些覆盖 $\{U_i \to U\}$，它们是
$(\Sch/X)_\tau$ 的覆盖并满足 $U \in \Ob(\mathcal{C}_t)$。则
\begin{enumerate}
\item[(a)] 函子 $\mathcal{C}_1 \to \mathcal{C}_2$ 全忠实、连续且余连续。
\end{enumerate}
以 $g : \Sh(\mathcal{C}_1) \to \Sh(\mathcal{C}_2)$ 表示相应的拓扑斯态射。
以 $\mathcal{O}_t$ 表示 $\mathcal{O}$ 到 $\mathcal{C}_t$ 的限制。以 $g_!$ 表示
Modules on Sites，定义 \ref{sites-modules-definition-g-shriek} 中的函子。
\begin{enumerate}
\item[(b)] 典范映射 $g_!\mathcal{O}_1 \to \mathcal{O}_2$ 是同构。
\end{enumerate}
\end{lemma}

\begin{proof}
断言 (a) 由定义立即可得。在本证明中，所有概形都是 $X$ 上的概形，所有概形态射
都是 $X$ 上的概形态射。注意，$g^{-1}$ 由限制给出，所以对对象 $U$
（它属于 $\mathcal{C}_1$），有 $\mathcal{O}_1(U) = \mathcal{O}_2(U) = \mathcal{O}(U)$。
回忆，$g_!\mathcal{O}_1$ 是预层 $g_{p!}\mathcal{O}_1$ 所联系的层；该预层把
对象 $V$（在 $\mathcal{C}_2$ 中）映到群
$$
\colim_{V \to U} \mathcal{O}(U)
$$
，其中 $U$ 遍历 $\mathcal{C}_1$ 的对象，余极限取在 Abel 群范畴中。下面将反复使用
如下事实：若
$$
V \to U \to U'
$$
是满足 $U, U' \in \Ob(\mathcal{C}_1)$ 的态射，且
$f' \in \mathcal{O}(U')$ 限制为 $f \in \mathcal{O}(U)$，则
$(V \to U, f)$ 与 $(V \to U', f')$ 定义余极限的同一个元素。此外，
$g_!\mathcal{O}_1 \to \mathcal{O}_2$ 只是把元素 $(V \to U, f)$ 映到 $f$ 在
$V$ 上的拉回。

\medskip\noindent
满性。设 $V$ 是概形，并设 $h \in \mathcal{O}(V)$。于是得到态射
$V \to X \times \mathbf{A}^1$；它由 $h$ 和结构态射 $V \to X$ 诱导。写成
$\mathbf{A}^1 = \Spec(\mathbf{Z}[x])$，可见元素
$x \in \mathcal{O}(X \times \mathbf{A}^1)$ 拉回到 $h$。由假设 (1)(b) 和 (2)，
$X \times \mathbf{A}^1$ 是 $\mathcal{C}_1$ 的对象，于是得到所需的满性。

\medskip\noindent
单性。设 $V$ 是概形。设
$s = \sum_{i = 1, \ldots, n} (V \to U_i, f_i)$ 是上面所示余极限的一个元素。
对任意 $i$，可以使用态射 $f_i : U_i \to X \times \mathbf{A}^1$，看出
$(V \to U_i, f_i)$ 与 $(f_i : V \to X \times \mathbf{A}^1, x)$ 定义余极限的
同一个元素。然后考虑
$$
f_1 \times \ldots \times f_n : V \to X \times \mathbf{A}^n
$$
，可见 $s$ 在余极限中等价于
$$
\sum\nolimits_{i = 1, \ldots, n}
(f_1 \times \ldots \times f_n : V \to X \times \mathbf{A}^n, x_i) =
(f_1 \times \ldots \times f_n : V \to X \times \mathbf{A}^n,
x_1 + \ldots + x_n)
$$
现在，若 $x_1 + \ldots + x_n$ 在 $V$ 上限制为零，则可见
$f_1 \times \ldots \times f_n$ 通过
$X \times \mathbf{A}^{n - 1} = V(x_1 + \ldots + x_n)$ 分解。因此，$s$
在余极限中等价于零。
\end{proof}











%9.29.09
\section{\'Etale 上同调}
\label{etale-cohomology-section-etale-cohomology}

\noindent
以下各节将证明一些关于 \'etale 上同调的基本结果。下面举出一个对拓扑空间上同调
成立、并且对 \'etale 上同调也成立的例子。

\begin{lemma}[\'Etale 上同调的 Mayer--Vietoris 序列]
\label{etale-cohomology-lemma-mayer-vietoris}
设 $X$ 是概形。假设 $X = U \cup V$ 是两个开子概形之并。对任意 Abel 层
$\mathcal{F}$（定义于 $X_\etale$ 上），存在上同调长正合序列
$$
\begin{matrix}
0 \to
H^0_\etale(X, \mathcal{F}) \to
H^0_\etale(U, \mathcal{F}) \oplus H^0_\etale(V, \mathcal{F}) \to
H^0_\etale(U \cap V, \mathcal{F}) \phantom{\to \ldots} \\
\phantom{0} \to H^1_\etale(X, \mathcal{F}) \to
H^1_\etale(U, \mathcal{F}) \oplus H^1_\etale(V, \mathcal{F}) \to
H^1_\etale(U \cap V, \mathcal{F}) \to \ldots
\end{matrix}
$$
这个长正合序列关于 $\mathcal{F}$ 是函子性的。
\end{lemma}

\begin{proof}
注意，若 $\mathcal{I}$ 是内射 Abel 层，则
$$
0 \to \mathcal{I}(X) \to \mathcal{I}(U) \oplus \mathcal{I}(V) \to
\mathcal{I}(U \cap V) \to 0
$$
正合。由于 $\mathcal{I}$ 是层，首项和中间项处的正合性成立。右端的正合性成立，
因为由 Cohomology on Sites，引理
\ref{sites-cohomology-lemma-restriction-along-monomorphism-surjective}，
$\mathcal{I}(U) \to \mathcal{I}(U \cap V)$ 是满射。另一种证明方法是指出映射
$\mathcal{I}(U) \oplus \mathcal{I}(V) \to
\mathcal{I}(U \cap V)$ 的余核就是
$\mathcal{I}$ 关于覆盖 $X = U \cup V$ 的第一 {\v C}ech 上同调群；它由引理
\ref{etale-cohomology-lemma-hom-injective} 和 \ref{etale-cohomology-lemma-forget-injectives} 消失。因此，若
$\mathcal{F} \to \mathcal{I}^\bullet$ 是内射解消，则
$$
0 \to \mathcal{I}^\bullet(X) \to
\mathcal{I}^\bullet(U) \oplus \mathcal{I}^\bullet(V) \to
\mathcal{I}^\bullet(U \cap V) \to 0
$$
是复形的短正合序列，而相联系的上同调长正合序列正是引理陈述中的序列。
\end{proof}

\begin{lemma}[相对 Mayer--Vietoris 序列]
\label{etale-cohomology-lemma-relative-mayer-vietoris}
设 $f : X \to Y$ 是概形态射。假设 $X = U \cup V$ 是两个开子概形之并。记
$a = f|_U : U \to Y$、$b = f|_V : V \to Y$ 以及
$c = f|_{U \cap V} : U \cap V \to Y$。对每个 Abel 层 $\mathcal{F}$
（定义于 $X_\etale$ 上），存在长正合序列
$$
0 \to
f_*\mathcal{F} \to
a_*(\mathcal{F}|_U) \oplus b_*(\mathcal{F}|_V) \to
c_*(\mathcal{F}|_{U \cap V}) \to
R^1f_*\mathcal{F} \to \ldots
$$
（定义于 $Y_\etale$ 上）。这个长正合序列关于 $\mathcal{F}$ 是函子性的。
\end{lemma}

\begin{proof}
设 $\mathcal{F} \to \mathcal{I}^\bullet$ 是 $\mathcal{F}$ 在 $X_\etale$ 上的
内射解消。我们断言得到复形的短正合序列
$$
0 \to
f_*\mathcal{I}^\bullet \to
a_*\mathcal{I}^\bullet|_U \oplus b_*\mathcal{I}^\bullet|_V \to
c_*\mathcal{I}^\bullet|_{U \cap V} \to
0.
$$
事实上，对任意 $W$（属于 $Y_\etale$）以及任意 $n \geq 0$，在 $W$ 上相应的
截面群序列
$$
0 \to
\mathcal{I}^n(W \times_Y X) \to
\mathcal{I}^n(W \times_Y U)
\oplus \mathcal{I}^n(W \times_Y V) \to
\mathcal{I}^n(W \times_Y (U \cap V)) \to
0
$$
已在引理 \ref{etale-cohomology-lemma-mayer-vietoris} 的证明中证明为短正合。取上同调层，并使用
$\mathcal{I}^\bullet|_U$ 是 $\mathcal{F}|_U$ 的内射解消这一事实，即得引理；
$\mathcal{I}^\bullet|_V$、$\mathcal{I}^\bullet|_{U \cap V}$ 的情形相同。
\end{proof}







\section{余极限}
\label{etale-cohomology-section-colimit}

\noindent
回忆，若 $(\mathcal{F}_i, \varphi_{ii'})$ 是网站 $\mathcal{C}$ 上的层图，则它在层范畴中
的余极限是预层 $U \mapsto \colim_i \mathcal{F}_i(U)$ 的层化；见 Sites，引理
\ref{sites-lemma-colimit-sheaves}。若系统是定向的，$U$ 是 $\mathcal{C}$ 的拟紧对象，
并且它有由拟紧对象组成的余终覆盖系，则
$\mathcal{F}(U) = \colim \mathcal{F}_i(U)$；见 Sites，引理
\ref{sites-lemma-directed-colimits-sections}。关于 Abel 层余极限之高阶上同调群的结果，
见 Cohomology on Sites，引理
\ref{sites-cohomology-lemma-colim-works-over-collection}。

\medskip\noindent
在 Cohomology on Sites，引理 \ref{sites-cohomology-lemma-colimit} 中，我们把这个结果
推广到网站逆系统上的层系统。下面给出概形逆系统上的 \'etale 层系统这一相应概念。

\begin{definition}
\label{etale-cohomology-definition-inverse-system-sheaves}
设 $I$ 是预序集。设 $(X_i, f_{i'i})$ 是以 $I$ 为指标的概形逆系统。
一个{\it 层系统 $(\mathcal{F}_i, \varphi_{i'i})$}（定义于 $(X_i, f_{i'i})$ 上）
由下列数据给出：
\begin{enumerate}
\item 一个层 $\mathcal{F}_i$（定义于 $(X_i)_\etale$ 上），对所有 $i \in I$；
\item 对 $i' \geq i$，层的映射
$\varphi_{i'i} : f_{i'i}^{-1}\mathcal{F}_i \to \mathcal{F}_{i'}$
（定义于 $(X_{i'})_\etale$ 上）。
\end{enumerate}
这些数据满足等式
$\varphi_{i''i} = \varphi_{i''i'} \circ f_{i'' i'}^{-1}\varphi_{i'i}$，
只要 $i'' \geq i' \geq i$。
\end{definition}

\noindent
在定义 \ref{etale-cohomology-definition-inverse-system-sheaves} 的情形下，假设 $I$ 是定向集，并且转移态射
$f_{i'i}$ 仿射。设 $X = \lim X_i$ 是概形范畴中的极限；见 Limits，第
\ref{limits-section-limits} 节。以 $f_i : X \to X_i$ 表示投影态射，并考虑映射
$$
f_i^{-1}\mathcal{F}_i = f_{i'}^{-1}f_{i'i}^{-1}\mathcal{F}_i
\xrightarrow{f_{i'}^{-1}\varphi_{i'i}}
f_{i'}^{-1}\mathcal{F}_{i'}
$$
这使 $f_i^{-1}\mathcal{F}_i$ 成为 $X_\etale$ 上以 $I$ 为指标的层系统
（验证这一点是个很好的练习）。我们经常希望知道是否存在同构
$$
H^q_\etale(X, \colim f_i^{-1}\mathcal{F}_i) =
\colim H^q_\etale(X_i, \mathcal{F}_i)
$$
后面会证明：若 $X_i$ 对所有 $i$ 都拟紧且拟分离，则此式成立；见定理
\ref{etale-cohomology-theorem-colimit}。

\begin{lemma}
\label{etale-cohomology-lemma-colimit-affine-sites}
设 $I$ 是定向集。设 $(X_i, f_{i'i})$ 是以 $I$ 为指标、具有仿射转移态射的概形逆系统。
设 $X = \lim_{i \in I} X_i$。采用 Topologies，引理
\ref{topologies-lemma-alternative} 中的记号，有
$$
X_{affine, \etale} = \colim (X_i)_{affine, \etale}
$$
；这里的网站等式取 Sites，引理 \ref{sites-lemma-colimit-sites} 的意义。
\end{lemma}

\begin{proof}
先在 $X$ 与 $X_i$ 都拟紧且拟分离（对所有 $i$）时证明此式；所有我们关心的情形
都是如此。在这种情形下，$X_{affine, \etale}$ 的任意对象，以及相应地
$(X_i)_{affine, \etale}$ 的任意对象，都在 $X$ 上有限表示。此外，$X$ 上有限表示概形
的范畴是 $X_i$ 上有限表示概形范畴的余极限；见 Limits，引理
\ref{limits-lemma-descend-finite-presentation}。由 Limits，引理
\ref{limits-lemma-limit-affine} 和 \ref{limits-lemma-descend-etale}，对 $X$ 上
\'etale 的仿射对象所成子范畴，同一结论也成立。最后，若 $\{U^j \to U\}$ 是
$X_{affine, \etale}$ 的覆盖，并且 $U_i^j \to U_i$ 是 $X_i$ 上 \'etale 的仿射概形
间态射，其到 $X$ 的基变换为 $U^j \to U$，则
$\{U^j_i \to U_i\}$ 到某个 $X_{i'}$ 的基变换是覆盖（当 $i'$ 足够大）；见 Limits，引理
\ref{limits-lemma-descend-surjective}。

\medskip\noindent
在一般情形下，设 $U$ 是 $X_{affine, \etale}$ 的对象。则 $U \to X$ 是 \'etale
且分离的（因为 $U$ 分离），但一般不拟紧。不过，$U \to X$ 局部有限表示。因此，
由 Limits，引理 \ref{limits-lemma-descend-finite-presentation-variant}，存在某个 $i$、
拟紧拟分离概形 $U_i$ 以及局部有限表示态射 $U_i \to X_i$，其到 $X$ 的基变换
为 $U \to X$。于是 $U = \lim_{i' \geq i} U_{i'}$，其中
$U_{i'} = U_i \times_{X_i} X_{i'}$。增大 $i$ 后，可以假设 $U_i$ 仿射；见
Limits，引理 \ref{limits-lemma-limit-affine}。为验证 $U_i \to X_i$ 对充分大的 $i$
是 \'etale 的，选取有限仿射开覆盖
$U_i = U_{i, 1} \cup \ldots \cup U_{i, m}$，使得
$U_{i, j} \to U_i \to X_i$ 映入仿射开集 $W_{i, j} \subset X_i$。应用 Limits，引理
\ref{limits-lemma-descend-etale}，可知 $U_{i, j} \to W_{i, j}$ 是 \'etale 的
（必要时再次增大 $i$）。由此可见函子
$\colim (X_i)_{affine, \etale} \to X_{affine, \etale}$ 本质满。全忠实性直接由
已经使用过的 Limits，引理
\ref{limits-lemma-descend-finite-presentation-variant} 得到。关于覆盖的断言与证明第一段
完全同样地证明。
\end{proof}

\noindent
利用上面的结果，得到下列关于余极限和上同调的一般结论。

\begin{theorem}
\label{etale-cohomology-theorem-colimit}
设 $X = \lim_{i \in I} X_i$ 是具有仿射转移态射
$f_{i'i} : X_{i'} \to X_i$ 的概形定向系统的极限。假设 $X_i$ 对所有 $i \in I$
都拟紧且拟分离。设 $(\mathcal{F}_i, \varphi_{i'i})$ 是 $(X_i, f_{i'i})$ 上的
Abel 层系统。以 $f_i : X \to X_i$ 表示投影，并置
$\mathcal{F} = \colim f_i^{-1}\mathcal{F}_i$。则
$$
\colim_{i\in I} H_\etale^p(X_i, \mathcal{F}_i) = H_\etale^p(X, \mathcal{F}).
$$
对所有 $p \geq 0$ 成立。
\end{theorem}

\begin{proof}
由 Topologies，引理 \ref{topologies-lemma-alternative}，可以计算
$\mathcal{F}$ 在 $X_{affine, \etale}$ 上的上同调。因此，结论由引理
\ref{etale-cohomology-lemma-colimit-affine-sites} 和 Cohomology on Sites，引理
\ref{sites-cohomology-lemma-colimit} 合并得到。
\end{proof}

\noindent
下面两个结果是上述定理的特殊情形。

\begin{lemma}
\label{etale-cohomology-lemma-colimit}
设 $X$ 是拟紧拟分离概形。设 $I$ 是定向集。设
$(\mathcal{F}_i, \varphi_{ij})$ 是 $X_\etale$ 上以 $I$ 为指标的 Abel 层系统。则
$$
\colim_{i\in I} H_\etale^p(X, \mathcal{F}_i) = H_\etale^p(X,
\colim_{i\in I} \mathcal{F}_i).
$$
\end{lemma}

\begin{proof}
这是定理 \ref{etale-cohomology-theorem-colimit} 的特殊情形。再概述一个直接证明。同时对所有 $X$
证明结论，并对 $p$ 作归纳。
\begin{enumerate}
\item
对任意拟紧拟分离概形 $X$ 以及任意 \'etale 覆盖 $\mathcal{U}$（它覆盖 $X$），证明存在加细
$\mathcal{V} = \{V_j \to X\}_{j\in J}$，其中 $J$ 有限，每个 $V_j$ 拟紧且拟分离，
并且所有 $V_{j_0} \times_X \ldots \times_X V_{j_p}$ 也拟紧且拟分离。
\item
利用上一步和层范畴中余极限的定义，证明定理对 $p = 0$ 以及所有 $X$ 成立。
\item
利用上同调的局部性（引理 \ref{etale-cohomology-lemma-locality-cohomology}）、从 {\v C}ech 到上同调的
谱序列（定理 \ref{etale-cohomology-theorem-cech-ss}），以及归纳假设适用于上述情形中的所有
$V_{j_0} \times_X \ldots \times_X V_{j_p}$ 这一事实，证明归纳步骤 $p \to p + 1$。
\end{enumerate}
\end{proof}

\begin{lemma}
\label{etale-cohomology-lemma-directed-colimit-cohomology}
设 $A$ 是环，$(I, \leq)$ 是定向集，$(B_i, \varphi_{ij})$ 是 $A$-代数系统。
置 $B = \colim_{i\in I} B_i$。设 $X \to \Spec(A)$ 是拟紧拟分离的概形态射。
设 $\mathcal{F}$ 是 $X_\etale$ 上的 Abel 层。记
$Y_i = X \times_{\Spec(A)} \Spec(B_i)$、
$Y = X \times_{\Spec(A)} \Spec(B)$、
$\mathcal{G}_i = (Y_i \to X)^{-1}\mathcal{F}$ 以及
$\mathcal{G} = (Y \to X)^{-1}\mathcal{F}$。则
$$
H_\etale^p(Y, \mathcal{G}) =
\colim_{i\in I} H_\etale^p (Y_i, \mathcal{G}_i).
$$
\end{lemma}

\begin{proof}
这是定理 \ref{etale-cohomology-theorem-colimit} 的特殊情形。再如下概述一个直接证明。
\begin{enumerate}
\item 给定 \'etale 态射 $V \to Y$，其中 $V$ 拟紧且拟分离，则存在 $i\in I$ 和
$V_i \to Y_i$，使得 $V = V_i \times_{Y_i} Y$。若所考虑的概形全都仿射，
这对应于下列代数断言：若 $B = \colim B_i$ 且 $B \to C$ 是 \'etale 的，
则存在 $i \in I$ 以及 \'etale 同态 $B_i \to C_i$，使得
$C \cong B \otimes_{B_i} C_i$。这在 Algebra，引理 \ref{algebra-lemma-etale} 中证明。
\item 在 (1) 的情形下，证明
$\mathcal{G}(V) = \colim_{i' \geq i} \mathcal{G}_{i'}(V_{i'})$
，其中 $V_{i'}$ 是 $V_i$ 到 $Y_{i'}$ 的基变换。
\item 由 (1) 可见，对每个 \'etale 覆盖
$\mathcal{V} = \{V_j \to Y\}_{j\in J}$，其中 $J$ 有限且诸 $V_j$ 拟紧拟分离，
存在 $i \in I$ 以及 \'etale 覆盖
$\mathcal{V}_i = \{V_{ij} \to Y_i\}_{j \in J}$，使得
$\mathcal{V} \cong \mathcal{V}_i \times_{Y_i} Y$。
\item 证明 (2) 和 (3) 蕴含
$$
\check H^*(\mathcal{V}, \mathcal{G})=
\colim_{i\in I} \check H^*(\mathcal{V}_i, \mathcal{G}_i).
$$
\item 巧妙地使用从 {\v C}ech 到上同调的谱序列（定理 \ref{etale-cohomology-theorem-cech-ss}）。
\end{enumerate}
\end{proof}

\begin{lemma}
\label{etale-cohomology-lemma-higher-direct-images}
设 $f: X\to Y$ 是概形态射，并且 $\mathcal{F}\in
\textit{Ab}(X_\etale)$。则 $R^pf_*\mathcal{F}$ 是下列预层所联系的层：
$$
(V \to Y) \longmapsto H_\etale^p(X \times_Y V, \mathcal{F}|_{X \times_Y V}).
$$
更一般地，对 $K \in D(X_\etale)$，$R^pf_*K$ 是下列预层所联系的层：
$$
(V \to Y) \longmapsto H_\etale^p(X \times_Y V, K|_{X \times_Y V}).
$$
\end{lemma}

\begin{proof}
这个引理对拓扑空间也成立，并且该情形的证明相同。层的情形见 Cohomology on Sites，
引理 \ref{sites-cohomology-lemma-higher-direct-images}；Abel 层复形的情形见
Cohomology on Sites，引理
\ref{sites-cohomology-lemma-sheafification-cohomology}。
\end{proof}

\begin{lemma}
\label{etale-cohomology-lemma-relative-colimit}
设 $S$ 是概形。设 $X = \lim_{i \in I} X_i$ 是 $S$ 上概形定向系统的极限，
其转移态射 $f_{i'i} : X_{i'} \to X_i$ 仿射。假设结构态射
$g_i : X_i \to S$ 与 $g : X \to S$ 拟紧且拟分离。设
$(\mathcal{F}_i, \varphi_{i'i})$ 是 $(X_i, f_{i'i})$ 上的 Abel 层系统。
以 $f_i : X \to X_i$ 表示投影，并置
$\mathcal{F} = \colim f_i^{-1}\mathcal{F}_i$。则
$$
\colim_{i\in I} R^p g_{i, *} \mathcal{F}_i = R^p g_* \mathcal{F}
$$
对所有 $p \geq 0$ 成立。
\end{lemma}

\begin{proof}
回忆（引理 \ref{etale-cohomology-lemma-higher-direct-images}），$R^p g_{i, *} \mathcal{F}_i$ 是预层
$U \mapsto H^p_\etale(U \times_S X_i, \mathcal{F}_i)$ 所联系的层；
$R^pg_*\mathcal{F}$ 也类似。此外，层系统的余极限是预层层次上余极限的层化。
注意，$S_\etale$ 的每个对象都有由拟紧拟分离对象（例如仿射概形）组成的覆盖。
此外，若 $U$ 是拟紧拟分离对象，则有
$$
\colim H^p_\etale(U \times_S X_i, \mathcal{F}_i) =
H^p_\etale(U \times_S X, \mathcal{F})
$$
；这由定理 \ref{etale-cohomology-theorem-colimit} 得到。因此引理成立。
\end{proof}

\begin{lemma}
\label{etale-cohomology-lemma-relative-colimit-general}
设 $I$ 是定向集。设 $g_i : X_i \to S_i$ 是以 $I$ 为指标的概形态射逆系统。
假设 $g_i$ 拟紧且拟分离，并且对 $i' \geq i$，转移态射
$f_{i'i} : X_{i'} \to X_i$ 与 $h_{i'i} : S_{i'} \to S_i$ 仿射。设
$g : X \to S$ 是态射 $g_i$ 的极限；见 Limits，第 \ref{limits-section-limits} 节。
以 $f_i : X \to X_i$ 和 $h_i : S \to S_i$ 表示投影。设
$(\mathcal{F}_i, \varphi_{i'i})$ 是 $(X_i, f_{i'i})$ 上的层系统。置
$\mathcal{F} = \colim f_i^{-1}\mathcal{F}_i$。则
$$
R^p g_* \mathcal{F} =
\colim_{i \in I} h_i^{-1}R^p g_{i, *} \mathcal{F}_i
$$
对所有 $p \geq 0$ 成立。
\end{lemma}

\begin{proof}
引理中的映射如何构造？对 $i' \geq i$，有交换图
$$
\xymatrix{
X \ar[r]_{f_{i'}} \ar[d]_g &
X_{i'} \ar[r]_{f_{i'i}} \ar[d]_{g_{i'}} &
X_i \ar[d]^{g_i} \\
S \ar[r]^{h_{i'}} &
S_{i'} \ar[r]^{h_{i'i}} &
S_i
}
$$
把基变换映射
$h_{i'i}^{-1}Rg_{i, *}\mathcal{F}_i \to Rg_{i', *}f_{i'i}^{-1}\mathcal{F}_i$
（Cohomology on Sites，引理
\ref{sites-cohomology-lemma-base-change-map-flat-case} 或
注 \ref{sites-cohomology-remark-base-change}）
与映射 $Rg_{i', *}\varphi_{i'i}$ 合并，得到
$\psi_{i'i} : h_{i' i}^{-1} R^p g_{i, *} \mathcal{F}_i \to
R^pg_{i', *} \mathcal{F}_{i'}$。类似地，利用图中的左方块，得到映射
$\psi_i : h_i^{-1}R^pg_{i, *}\mathcal{F}_i \to R^pg_*\mathcal{F}$.
映射 $h_{i'}^{-1}\psi_{i'i}$ 与 $\psi_i$ 就是引理陈述中使用的映射。为使其有意义，
必须验证
$\psi_{i''i} = \psi_{i''i'} \circ h_{i''i'}^{-1}\psi_{i'i}$ 且
$\psi_{i'} \circ h_{i'}^{-1}\psi_{i'i} = \psi_i$；这由 Cohomology on Sites，注
\ref{sites-cohomology-remark-compose-base-change-horizontal} 得到。

\medskip\noindent
证明该等式。第一种证明使用维数移动\footnote{也可以用这种方法构造引理中的映射。}。
对任意 $U$（它在 $X$ 上仿射且 \'etale），由定理 \ref{etale-cohomology-theorem-colimit} 有
$$
g_*\mathcal{F}(U) =
H^0(U \times_S X, \mathcal{F}) =
\colim H^0(U_i \times_{S_i} X_i, \mathcal{F}_i) =
\colim g_{i, *}\mathcal{F}_i(U_i)
$$
，其中余极限取遍充分大的 $i$，从而存在某个 $i$ 和某个 $U_i$
（它在 $S_i$ 上仿射且 \'etale），其基变换为 $U$（在 $S$ 上；见引理
\ref{etale-cohomology-lemma-colimit-affine-sites}）。
由 Sites，引理 \ref{sites-lemma-colimit}，右边等于
$(\colim h_i^{-1}g_{i, *}\mathcal{F}_i)(U)$。这证明了 $p = 0$ 时的引理。
若 $(\mathcal{G}_i, \varphi_{i'i})$ 是满足
$\mathcal{G} = \colim f_i^{-1}\mathcal{G}_i$ 的系统，使得 $\mathcal{G}_i$ 是
$X_i$ 上的内射 Abel 层（对所有 $i$），则对任意 $U$（它在 $X$ 上仿射且 \'etale），由定理
\ref{etale-cohomology-theorem-colimit} 有
$$
H^p(U \times_S X, \mathcal{G}) =
\colim H^p(U_i \times_{S_i} X_i, \mathcal{G}_i) = 0
$$
，对 $p > 0$ 成立（余极限与前相同）。因此 $R^pg_*\mathcal{G} = 0$，
并且对这样的系统得到 $p > 0$ 时的结论。一般而言，可以选取系统的短正合序列
$$
0 \to (\mathcal{F}_i, \varphi_{i'i}) \to
(\mathcal{G}_i, \varphi_{i'i}) \to
(\mathcal{Q}_i, \varphi_{i'i}) \to 0
$$
，其中 $(\mathcal{G}_i, \varphi_{i'i})$ 如上；见 Cohomology on Sites，引理
\ref{sites-cohomology-lemma-colim-sites-injective}。由归纳，引理对 $p - 1$ 成立；
由上文，对 $p$ 和 $(\mathcal{G}_i, \varphi_{i'i})$ 有消失。因此，上同调长正合序列
给出 $p$ 和 $(\mathcal{F}_i, \varphi_{i'i})$ 的结论。

\medskip\noindent
第二种证明。回忆，$S_{affine, \etale} = \colim (S_i)_{affine, \etale}$；见引理
\ref{etale-cohomology-lemma-colimit-affine-sites}。因此，若 $U$ 是 $S_{affine, \etale}$ 的对象，
则可写成 $U = U_i \times_{S_i} S$，其中有某个 $i$ 以及某个 $U_i$
（它属于 $(S_i)_{affine, \etale}$），并且
$$
(\colim_{i \in I} h_i^{-1}R^p g_{i, *} \mathcal{F}_i)(U) =
\colim_{i' \geq i} (R^p g_{i', *}\mathcal{F}_{i'})(U_i \times_{S_i} S_{i'})
$$
；这由 Sites，引理 \ref{sites-lemma-colimit} 以及上面所述系统中转移映射的构造得到。
由于 $R^pg_{i', *}\mathcal{F}_{i'}$ 是预层
$U_{i'} \mapsto H^p(U_{i'} \times_{S_{i'}} X_{i'}, \mathcal{F}_{i'})$ 所联系的层，
而 $R^pg_*\mathcal{F}$ 是预层 $U \mapsto H^p(U \times_S X, \mathcal{F})$
所联系的层（引理 \ref{etale-cohomology-lemma-higher-direct-images}），得到典范交换图
$$
\xymatrix{
\colim_{i' \geq i}
H^p(U_i \times_{S_i} X_{i'}, \mathcal{F}_{i'}) \ar[r] \ar[d] &
\colim_{i' \geq i}
(R^p g_{i', *}\mathcal{F}_{i'})(U_i \times_{S_i} S_{i'}) \ar[d] \\
H^p(U \times_S X, \mathcal{F}) \ar[r] &
R^pg_*\mathcal{F}(U)
}
$$
注意，由定理 \ref{etale-cohomology-theorem-colimit}，左侧竖直箭头是同构。我们要证明右侧竖直箭头
是同构。不过，已经知道该箭头的源和靶都是 $S_{affine, \etale}$ 上的层。因此，只需
证明：(1) 靶中的元素局部来自源中的元素；(2) 源中映到靶中零的元素局部消失。
(1) 立即由上文和下方水平箭头来自预层映射这一事实得到；该预层映射在层化后成为同构。
对于 (2)，设
$\xi \in  \colim_{i' \geq i}
(R^p g_{i', *}\mathcal{F}_{i'})(U_i \times_{S_i} S_{i'})$ 属于核。选取
$i' \geq i$ 以及元素
$\xi_{i'} \in (R^p g_{i', *}\mathcal{F}_{i'})(U_i \times_{S_i} S_{i'})$，
它表示 $\xi$。
选取标准 \'etale 覆盖
$\{U_{i', k} \to U_i \times_{S_i} S_{i'}\}_{k = 1, \ldots, m}$，
使得 $\xi_{i'}|_{U_{i', k}}$ 来自
$\xi_{i', k} \in H^p(U_{i', k} \times_{S_{i'}} X_{i'}, \mathcal{F}_{i'})$。
由于只需证明 $\xi$ 局部消失，可以替换 $U$，使用 \'etale 覆盖
$\{U_{i', k} \times_{S_{i'}} S \to U = U_i \times_{S_i} S\}$ 的成员。
替换后，可见 $\xi$ 是某个元素 $\xi'$ 的像；该元素属于上图中的群
$\colim_{i' \geq i} H^p(U_i \times_{S_i} X_{i'}, \mathcal{F}_{i'})$。
由于 $\xi'$ 在 $R^pg_*\mathcal{F}(U)$ 中映到零，可以再作一次替换，
并假设 $\xi'$ 在 $H^p(U \times_S X, \mathcal{F})$ 中映到零。不过，由于左侧竖直箭头
是同构，于是得到 $\xi' = 0$，从而如愿有 $\xi = 0$。
\end{proof}

\begin{lemma}
\label{etale-cohomology-lemma-linus-hamann}
设 $X = \lim_{i \in I} X_i$ 是具有仿射转移态射 $f_{i'i}$ 和投影态射
$f_i : X \to X_i$ 的概形定向极限。设 $\mathcal{F}$ 是 $X_\etale$ 上的层。则
\begin{enumerate}
\item 存在典范映射
$\varphi_{i'i} : f_{i'i}^{-1}f_{i, *}\mathcal{F} \to f_{i', *}\mathcal{F}$
，使得 $(f_{i, *}\mathcal{F}, \varphi_{i'i})$ 是定义
\ref{etale-cohomology-definition-inverse-system-sheaves} 意义下 $(X_i, f_{i'i})$ 上的层系统；
\item $\mathcal{F} = \colim f_i^{-1}f_{i, *}\mathcal{F}$.
\end{enumerate}
\end{lemma}

\begin{proof}
通过 Topologies，引理 \ref{topologies-lemma-alternative} 和引理
\ref{etale-cohomology-lemma-colimit-affine-sites}，这是 Sites，引理
\ref{sites-lemma-colimit-push-pull} 的特殊情形。
\end{proof}

\begin{lemma}
\label{etale-cohomology-lemma-compute-strangely}
设 $I$ 是定向集。设 $g_i : X_i \to S_i$ 是以 $I$ 为指标的概形态射逆系统。
假设 $g_i$ 拟紧且拟分离，并且对 $i' \geq i$，转移态射
$X_{i'} \to X_i$ 和 $S_{i'} \to S_i$ 仿射。设 $g : X \to S$ 是态射 $g_i$ 的极限；
见 Limits，第 \ref{limits-section-limits} 节。以 $f_i : X \to X_i$ 和
$h_i : S \to S_i$ 表示投影。设 $\mathcal{F}$ 是 $X$ 上的 Abel 层。则有
$$
R^pg_*\mathcal{F} = \colim_{i \in I} h_i^{-1}R^pg_{i, *}(f_{i, *}\mathcal{F})
$$
\end{lemma}

\begin{proof}
这是引理 \ref{etale-cohomology-lemma-relative-colimit-general} 和
\ref{etale-cohomology-lemma-linus-hamann} 的形式组合。
\end{proof}











\section{余极限与复形}
\label{etale-cohomology-section-colimit-complexes}

\noindent
本节讨论如何在各种情形下对复形系统取上同调，继续第
\ref{etale-cohomology-section-colimit} 节中关于层的讨论。
除非绝对必要，我们强烈建议读者不要阅读本节。

\begin{lemma}
\label{etale-cohomology-lemma-colimit-variant-complexes}
设 $X = \lim_{i \in I} X_i$ 是概形定向系统的极限，其仿射转移态射为
$f_{i'i} : X_{i'} \to X_i$。假设概形 $X_i$ 对所有 $i \in I$
均拟紧且拟分离。设 $\mathcal{F}_i^\bullet$ 是
$X_{i, \etale}$ 上的 Abel 层复形。设
$\varphi_{i'i} : f_{i'i}^{-1}\mathcal{F}_i^\bullet \to
\mathcal{F}_{i'}^\bullet$ 是 $X_{i, \etale}$ 上的复形映射，
并且等式
$\varphi_{i''i} = \varphi_{i''i'} \circ f_{i'' i'}^{-1}\varphi_{i'i}$
在每当 $i'' \geq i' \geq i$ 时成立。假设存在整数 $a$，使得
$\mathcal{F}_i^n = 0$ 对 $n < a$ 以及所有 $i \in I$ 成立。则
$$
H^p_\etale(X, \colim f_i^{-1}\mathcal{F}_i^\bullet) =
\colim H^p_\etale(X_i, \mathcal{F}^\bullet_i)
$$
，其中 $f_i : X \to X_i$ 是投影。
\end{lemma}

\begin{proof}
这是定理 \ref{etale-cohomology-theorem-colimit} 的推论。置
$\mathcal{F}^\bullet = \colim f_i^{-1}\mathcal{F}_i^\bullet$。
该定理告诉我们
$$
\colim_{i \in I} H_\etale^p(X_i, \mathcal{F}_i^n) =
H_\etale^p(X, \mathcal{F}^n)
$$
对所有 $n, p \in \mathbf{Z}$ 成立。使用谱序列
$$
E_{1, i}^{s, t} = H_\etale^t(X_i, \mathcal{F}_i^s) \Rightarrow
H_\etale^{s + t}(X_i, \mathcal{F}_i^\bullet)
$$
以及
$$
E_1^{s, t} = H_\etale^t(X, \mathcal{F}^s) \Rightarrow
H_\etale^{s + t}(X, \mathcal{F}^\bullet)
$$
；见 Derived Categories，引理 \ref{derived-lemma-two-ss-complex-functor}。
由于 $\mathcal{F}_i^n = 0$ 对 $n < a$ 成立（其中 $a$ 与 $i$ 无关），
可见只有固定的有限多个项 $E_{1, i}^{s, t}$（与 $i$ 无关）和
$E_1^{s, t}$ 对
$H^q_\etale(X_i, \mathcal{F}_i^\bullet)$ 与
$H^q_\etale(X, \mathcal{F}^\bullet)$ 有贡献，并且
$E_1^{s, t} = \colim E_{i, i}^{s, t}$。
这就给出了所需结论；略去一些细节。
（也可使用复形的“笨截断”给出另一论证，从而避免使用谱序列。）
\end{proof}

\begin{lemma}
\label{etale-cohomology-lemma-direct-sum-bounded-below-cohomology}
设 $X$ 是拟紧且拟分离概形。设
$K_i \in D(X_\etale)$，$i \in I$，是一族对象。
假设给定 $a \in \mathbf{Z}$，使得 $H^n(K_i) = 0$ 对
$n < a$ 和 $i \in I$ 成立。则
$R\Gamma(X, \bigoplus_i K_i) =
\bigoplus_i R\Gamma(X, K_i)$。
\end{lemma}

\begin{proof}
必须证明
$H^p(X, \bigoplus_i K_i) =
\bigoplus_i H^p(X, K_i)$ 对所有 $p \in \mathbf{Z}$ 成立。
选取复形 $\mathcal{F}_i^\bullet$ 来表示 $K_i$，使得
$\mathcal{F}_i^n = 0$ 对 $n < a$ 成立。
由 Injectives，引理 \ref{injectives-lemma-derived-products}，
复形 $\mathcal{F}_i^\bullet$ 的直和表示对象 $\bigoplus K_i$。
由于 $\bigoplus \mathcal{F}^\bullet$ 是有限直和的滤过余极限，
结论由引理 \ref{etale-cohomology-lemma-colimit-variant-complexes} 得到。
\end{proof}

\begin{lemma}
\label{etale-cohomology-lemma-colimit-complexes}
设 $S$ 是概形。设 $X = \lim_{i \in I} X_i$ 是 $S$ 上概形定向系统的极限，
其仿射转移态射为 $f_{i'i} : X_{i'} \to X_i$。假设概形 $X_i$ 对所有
$i \in I$ 均拟紧且拟分离。设
$K \in D^+(S_\etale)$。则
$$
\colim_{i \in I} H_\etale^p(X_i, K|_{X_i}) = H_\etale^p(X, K|_X).
$$
对所有 $p \in \mathbf{Z}$ 成立，其中 $K|_{X_i}$ 与 $K|_X$
分别是 $K$ 到 $X_i$ 与 $X$ 的拉回。
\end{lemma}

\begin{proof}
可以将 $K$ 表示为向下有界复形 $\mathcal{G}^\bullet$；
该复形由 $S_\etale$ 上的 Abel 层组成。设 $\mathcal{G}^n = 0$ 对
$n < a$ 成立。分别以 $\mathcal{F}^\bullet_i$ 和
$\mathcal{F}^\bullet$ 表示这个复形到 $X_i$ 和 $X$ 的拉回。
这些复形分别表示对象 $K|_{X_i}$ 和 $K|_X$，而且逐项有
$\mathcal{F}^\bullet = \colim f_i^{-1}\mathcal{F}_i^\bullet$。
因此，引理由引理 \ref{etale-cohomology-lemma-colimit-variant-complexes} 得到。
\end{proof}

\begin{lemma}
\label{etale-cohomology-lemma-relative-colimit-general-complexes}
令 $I$、$g_i : X_i \to S_i$、$g : X \to S$、$f_i$、$g_i$、$h_i$
如引理 \ref{etale-cohomology-lemma-relative-colimit-general} 中所述。
设 $0 \in I$ 且 $K_0 \in D^+(X_{0, \etale})$。
对 $i \geq 0$，以 $K_i$ 表示 $K_0$ 到 $X_i$ 的拉回；
以 $K$ 表示 $K_0$ 到 $X$ 的拉回。则
$$
R^pg_*K = \colim_{i \geq 0} h_i^{-1}R^pg_{i, *}K_i
$$
对所有 $p \in \mathbf{Z}$ 成立。
\end{lemma}

\begin{proof}
固定整数 $p_0 \in \mathbf{Z}$。
设 $a$ 是整数，使得 $H^j(K_0) = 0$ 对 $j < a$ 成立。
我们将证明该公式对所有 $p \leq p_0$ 成立，并对 $a$ 作下降归纳。
若 $a > p_0$，则由平凡消失可见，对 $p \leq p_0$，公式两边均为零；
见 Derived Categories，引理 \ref{derived-lemma-negative-vanishing}。
假设 $a \leq p_0$。考虑可区别三角
$$
H^a(K_0)[-a] \to K_0 \to \tau_{\geq a + 1}K_0
$$
将这个可区别三角拉回到 $X_i$ 和 $X$，得到关于 $K_i$ 和 $K$
的相容可区别三角。对 $p \leq p_0$，考虑交换图
$$
\xymatrix{
\colim_{i \geq 0} h_i^{-1}R^{p - 1}g_{i, *}(\tau_{\geq a + 1}K_i)
\ar[r]_-\alpha \ar[d] &
R^{p - 1}g_*(\tau_{\geq a + 1}K) \ar[d] \\
\colim_{i \geq 0} h_i^{-1}R^pg_{i, *}(H^a(K_i)[-a])
\ar[r]_-\beta \ar[d] &
R^pg_*(H^a(K)[-a]) \ar[d] \\
\colim_{i \geq 0} h_i^{-1}R^pg_{i, *}K_i
\ar[r]_-\gamma \ar[d] &
R^pg_*K \ar[d] \\
\colim_{i \geq 0} R^pg_{i, *}\tau_{\geq a + 1}K_i
\ar[r]_-\delta \ar[d] &
R^pg_*\tau_{\geq a + 1}K \ar[d] \\
\colim_{i \geq 0} R^{p + 1}g_{i, *}(H^a(K_i)[-a])
\ar[r]^-\epsilon &
R^{p + 1}g_*(H^a(K)[-a])
}
$$
，其各列正合。由引理 \ref{etale-cohomology-lemma-relative-colimit-general}，
箭头 $\beta$ 和 $\epsilon$ 是同构。
由归纳假设，箭头 $\alpha$ 和 $\delta$ 是同构。
因此 $\gamma$ 是同构，正合所求。
\end{proof}

\begin{lemma}
\label{etale-cohomology-lemma-relative-colimit-general-really-complexes}
令 $I$、$g_i : X_i \to S_i$、$g : X \to S$、$f_{ii'}$、$f_i$、$g_i$、$h_i$
如引理 \ref{etale-cohomology-lemma-relative-colimit-general} 中所述。
设 $\mathcal{F}_i^\bullet$ 是 $X_{i, \etale}$ 上的 Abel 层复形。
设 $\varphi_{i'i} : f_{i'i}^{-1}\mathcal{F}_i^\bullet \to
\mathcal{F}_{i'}^\bullet$ 是 $X_{i, \etale}$ 上的复形映射，
并且等式
$\varphi_{i''i} = \varphi_{i''i'} \circ f_{i'' i'}^{-1}\varphi_{i'i}$
在每当 $i'' \geq i' \geq i$ 时成立。假设存在整数 $a$，使得
$\mathcal{F}_i^n = 0$ 对 $n < a$ 以及所有 $i \in I$ 成立。则
$$
R^pg_*(\colim f_i^{-1}\mathcal{F}_i^\bullet) =
\colim_{i \geq 0} h_i^{-1}R^pg_{i, *}\mathcal{F}_i^\bullet
$$
对所有 $p \in \mathbf{Z}$ 成立。
\end{lemma}

\begin{proof}
这是引理 \ref{etale-cohomology-lemma-relative-colimit-general} 的推论。置
$\mathcal{F}^\bullet = \colim f_i^{-1}\mathcal{F}_i^\bullet$。
该引理告诉我们
$$
\colim_{i \in I} h_i^{-1}R^pg_{i, *}\mathcal{F}_i^n = R^pg_*\mathcal{F}^n
$$
对所有 $n, p \in \mathbf{Z}$ 成立。使用谱序列
$$
E_{1, i}^{s, t} = R^tg_{i, *}\mathcal{F}_i^s \Rightarrow
R^{s + t}g_{i, *}\mathcal{F}_i^\bullet
$$
以及
$$
E_1^{s, t} = R^tg_*\mathcal{F}^s \Rightarrow
R^{s + t}g_*\mathcal{F}^\bullet
$$
；见 Derived Categories，引理 \ref{derived-lemma-two-ss-complex-functor}。
由于 $\mathcal{F}_i^n = 0$ 对 $n < a$ 成立（其中 $a$ 与 $i$ 无关），
可见只有固定的有限多个项 $E_{1, i}^{s, t}$（与 $i$ 无关）和
$E_1^{s, t}$ 有贡献，并且
$E_1^{s, t} = \colim E_{i, i}^{s, t}$。
这就给出了所需结论；略去一些细节。
（也可使用复形的“笨截断”给出另一论证，从而避免使用谱序列。）
\end{proof}

\begin{lemma}
\label{etale-cohomology-lemma-direct-sum-bounded-below-pushforward}
设 $f : X \to Y$ 是拟紧且拟分离的概形态射。设
$K_i \in D(X_\etale)$，$i \in I$，是一族对象。
假设给定 $a \in \mathbf{Z}$，使得 $H^n(K_i) = 0$ 对
$n < a$ 和 $i \in I$ 成立。则
$Rf_*(\bigoplus_i K_i) = \bigoplus_i Rf_*K_i$。
\end{lemma}

\begin{proof}
必须证明
$R^pf_*(\bigoplus_i K_i) =
\bigoplus_i R^pf_*K_i$ 对所有 $p \in \mathbf{Z}$ 成立。
选取复形 $\mathcal{F}_i^\bullet$ 来表示 $K_i$，使得
$\mathcal{F}_i^n = 0$ 对 $n < a$ 成立。
由 Injectives，引理 \ref{injectives-lemma-derived-products}，
复形 $\mathcal{F}_i^\bullet$ 的直和表示对象 $\bigoplus K_i$。
由于 $\bigoplus \mathcal{F}^\bullet$ 是有限直和的滤过余极限，
结论由引理 \ref{etale-cohomology-lemma-relative-colimit-general-really-complexes} 得到。
\end{proof}














\section{高阶直接像的茎}
\label{etale-cohomology-section-stalks-direct-image}

\noindent
高阶直接像的茎常可按如下方式计算。

\begin{theorem}
\label{etale-cohomology-theorem-higher-direct-images}
设 $f: X \to S$ 是拟紧且拟分离的概形态射，
$\mathcal{F}$ 是 $X_\etale$ 上的 Abel 层，并且 $\overline{s}$ 是
$S$ 的一个位于 $s \in S$ 上方的几何点。则
$$
\left(R^nf_* \mathcal{F}\right)_{\overline{s}} =
H_\etale^n( X \times_S \Spec(\mathcal{O}_{S, \overline{s}}^{sh}),
p^{-1}\mathcal{F})
$$
，其中 $p : X \times_S \Spec(\mathcal{O}_{S, \overline{s}}^{sh}) \to X$
是投影。对 $K \in D^+(X_\etale)$ 和 $n \in \mathbf{Z}$，有
$$
\left(R^nf_*K\right)_{\overline{s}} =
H_\etale^n(X \times_S \Spec(\mathcal{O}_{S, \overline{s}}^{sh}), p^{-1}K)
$$
事实上，有
$$
\left(Rf_*K\right)_{\overline{s}}
=
R\Gamma_\etale(X \times_S \Spec(\mathcal{O}_{S, \overline{s}}^{sh}), p^{-1}K)
$$
作为 $D^+(\textit{Ab})$ 中的等式。
\end{theorem}

\begin{proof}
设 $\mathcal{I}$ 是 $\overline{s}$ 在 $S$ 上的 \'etale 邻域所成的范畴。
由引理 \ref{etale-cohomology-lemma-higher-direct-images}，有
$$
(R^nf_*\mathcal{F})_{\overline{s}} =
\colim_{(V, \overline{v}) \in \mathcal{I}^{opp}}
H_\etale^n(X \times_S V, \mathcal{F}|_{X \times_S V}).
$$
可以将 $\mathcal{I}$ 替换为由 $\overline{s}$ 的仿射 \'etale 邻域
组成的始子范畴。注意
$$
\Spec(\mathcal{O}_{S, \overline{s}}^{sh}) =
\lim_{(V, \overline{v}) \in \mathcal{I}} V
$$
；这是引理 \ref{etale-cohomology-lemma-describe-etale-local-ring} 以及
Limits，引理
\ref{limits-lemma-directed-inverse-system-affine-schemes-has-limit}
的结论。由于纤维积与极限交换，还得到
$$
X \times_S \Spec(\mathcal{O}_{S, \overline{s}}^{sh}) =
\lim_{(V, \overline{v}) \in \mathcal{I}} X \times_S V
$$
结论由引理 \ref{etale-cohomology-lemma-directed-colimit-cohomology} 得到。
对于第二种形式，使用相同论证，但以引理
\ref{etale-cohomology-lemma-colimit-complexes} 代替引理
\ref{etale-cohomology-lemma-directed-colimit-cohomology}。

\medskip\noindent
为证明最后一个陈述，只需构造映射
$\left(Rf_*K\right)_{\overline{s}} \to
R\Gamma_\etale(X \times_S \Spec(\mathcal{O}_{S, \overline{s}}^{sh}), p^{-1}K)$，
它属于 $D^+(\textit{Ab})$，并在上述每个次数 $n$
的上同调群上实现同构（对所有 $n$）。
为此，选取向下有界复形 $\mathcal{J}^\bullet$，它由
$X_\etale$ 上的内射 Abel 层组成，并用它表示 $K$。复形
$f_*\mathcal{J}^\bullet$ 表示 $Rf_*K$。因此，复形
$$
(f_*\mathcal{J}^\bullet)_{\overline{s}} =
\colim_{(V, \overline{v}) \in \mathcal{I}^{opp}}
(f_*\mathcal{J}^\bullet)(V)
$$
表示 $(Rf_*K)_{\overline{s}}$。对每个 $V$，有映射
$$
(f_*\mathcal{J}^\bullet)(V) =
\Gamma(X \times_S V, \mathcal{J}^\bullet)
\longrightarrow
\Gamma(X \times_S \Spec(\mathcal{O}_{S, \overline{s}}^{sh}),
p^{-1}\mathcal{J}^\bullet)
$$
，而靶复形表示
$R\Gamma_\etale(X \times_S \Spec(\mathcal{O}_{S, \overline{s}}^{sh}), p^{-1}K)$，
这是 $D^+(\textit{Ab})$ 中的对象。
对这些映射取余极限，即得结论。
\end{proof}

\begin{remark}
\label{etale-cohomology-remark-stalk-fibre}
设 $f : X \to S$ 是概形态射。
设 $K \in D(X_\etale)$。
设 $\overline{s}$ 是 $S$ 的几何点。
总有典范映射
$$
(Rf_*K)_{\overline{s}}
\longrightarrow
R\Gamma(X \times_S \Spec(\mathcal{O}_{S, \overline{s}}^{sh}), p^{-1}K)
\longrightarrow
R\Gamma(X_{\overline{s}}, K|_{X_{\overline{s}}})
$$
，其中 $p : X \times_S \Spec(\mathcal{O}_{S, \overline{s}}^{sh}) \to X$
是投影。具体而言，考虑交换图
$$
\xymatrix{
X_{\overline{s}} \ar[r] \ar[d]^{f_{\overline{s}}} &
X \times_S \Spec(\mathcal{O}_{S, \overline{s}}^{sh}) \ar[r]_-p \ar[d]^{f'} &
X \ar[d]^f \\
\overline{s} \ar[r]^-i &
\Spec(\mathcal{O}_{S, \overline{s}}^{sh}) \ar[r]^-j &
S
}
$$
对于该图中的两个方块，有基变换映射
$$
i^{-1}Rf'_*(p^{-1}K) \to Rf_{\overline{s}, *}(K|_{X_{\overline{s}}})
\quad\text{且}\quad
j^{-1}Rf_*K \to Rf'_*(p^{-1}K)
$$
（Cohomology on Sites，注
\ref{sites-cohomology-remark-base-change}）。
取整体截面，即得所需映射。
由 Cohomology on Sites，注
\ref{sites-cohomology-remark-compose-base-change-horizontal}，
这两个映射的复合就是通常的（基变换）映射
$(Rf_*K)_{\overline{s}} \to R\Gamma(X_{\overline{s}}, K|_{X_{\overline{s}}})$。
\end{remark}






\section{Leray 谱序列}
\label{etale-cohomology-section-leray}

\begin{lemma}
\label{etale-cohomology-lemma-prepare-leray}
设 $f: X \to Y$ 是态射，而 $\mathcal{I}$ 是
$\textit{Ab}(X_\etale)$ 的内射对象。设 $V \in \Ob(Y_\etale)$。则
\begin{enumerate}
\item 对任意覆盖 $\mathcal{V} = \{V_j\to V\}_{j \in J}$，有
$\check H^p(\mathcal{V}, f_*\mathcal{I}) = 0$ 对所有 $p > 0$ 成立；
\item $f_*\mathcal{I}$ 关于函子 $\Gamma(V, -)$ 是无环的；并且
\item 若 $g : Y \to Z$，则 $f_*\mathcal{I}$ 关于 $g_*$ 是无环的。
\end{enumerate}
\end{lemma}

\begin{proof}
注意
$\check{\mathcal{C}}^\bullet(\mathcal{V}, f_*\mathcal{I}) =
\check{\mathcal{C}}^\bullet(\mathcal{V} \times_Y X, \mathcal{I})$，
由引理 \ref{etale-cohomology-lemma-hom-injective}，它的高阶上同调群消失。
这证明了 (1)。第二个陈述如下得到：若一个层对任意覆盖的高阶
{\v C}ech 上同调群均消失，则它的高阶上同调群也消失。
用 {\v C}ech 到上同调的谱序列证明这一点是一道很精彩的练习；
细节以及更精确、更一般的陈述见 Cohomology on Sites，引理
\ref{sites-cohomology-lemma-cech-vanish-collection}。
(3) 是 (2) 以及引理 \ref{etale-cohomology-lemma-higher-direct-images} 中
$R^pg_*$ 的描述的推论。
\end{proof}

\noindent
应用 Grothendieck 谱序列的形式体系，得到如下结论。

\begin{proposition}[Leray 谱序列]
\label{etale-cohomology-proposition-leray}
设 $f: X \to Y$ 是概形态射，而 $\mathcal{F}$ 是 $X$ 上的 \'etale 层。
则存在谱序列
$$
E_2^{p, q} = H_\etale^p(Y, R^qf_*\mathcal{F}) \Rightarrow
H_\etale^{p+q}(X, \mathcal{F}).
$$
\end{proposition}

\begin{proof}
见引理 \ref{etale-cohomology-lemma-prepare-leray} 以及 Derived Categories，第
\ref{derived-section-composition-right-derived-functors} 节。
\end{proof}









\section{有限态射高阶直接像的消失}
\label{etale-cohomology-section-vanishing-finite-morphism}

\noindent
下一个目标是证明，概形有限态射的高阶直接像消失。

\begin{lemma}
\label{etale-cohomology-lemma-vanishing-etale-cohomology-strictly-henselian}
设 $R$ 是严格 Hensel 局部环。置 $S = \Spec(R)$，并设
$\overline{s}$ 是它的闭点。则整体截面函子
$\Gamma(S, -) : \textit{Ab}(S_\etale) \to \textit{Ab}$
是正合的。事实上，有
$\Gamma(S, \mathcal{F}) = \mathcal{F}_{\overline{s}}$，
对任意集合层 $\mathcal{F}$ 成立。特别地，
$$
\forall p\geq 1, \quad H_\etale^p(S, \mathcal{F})=0
$$
对所有 $\mathcal{F}\in \textit{Ab}(S_\etale)$ 成立。
\end{lemma}

\begin{proof}
若证明了 $\Gamma(S, \mathcal{F}) = \mathcal{F}_{\overline{s}}$，
则由于茎函子正合，$\Gamma(S, -)$ 也正合。
设 $(U, \overline{u})$ 是 $\overline{s}$ 的一个 \'etale 邻域。
选取仿射开邻域 $\Spec(A)$；它是 $\overline{u}$ 在 $U$ 中的邻域。
于是 $R \to A$ 是 \'etale 的，并且
$\kappa(\overline{s}) = \kappa(\overline{u})$。
由定理 \ref{etale-cohomology-theorem-henselian}，可见存在同构
$A \cong R \times A'$；它是 $R$-代数同构，并与到
$\kappa(\overline{s}) = \kappa(\overline{u})$ 的映射相容。因此得到截面
$$
\xymatrix{
\Spec(A) \ar[r] & U \ar[d]\\
& S \ar[ul]
}
$$
于是，在 $\overline{s}$ 的 \'etale 邻域系统中，恒等映射
$(S, \overline{s}) \to (S, \overline{s})$ 是余终的。
因此 $\Gamma(S, \mathcal{F}) = \mathcal{F}_{\overline{s}}$。
引理的最后一个陈述是因为正合函子的高阶导出函子为零；见
Derived Categories，引理 \ref{derived-lemma-right-derived-exact-functor}。
\end{proof}

\begin{proposition}
\label{etale-cohomology-proposition-finite-higher-direct-image-zero}
设 $f : X \to Y$ 是概形的有限态射。
\begin{enumerate}
\item 对任意几何点 $\overline{y} : \Spec(k) \to Y$，有
$$
(f_*\mathcal{F})_{\overline{y}} =
\prod\nolimits_{\overline{x} : \Spec(k) \to X,\ f(\overline{x}) =
\overline{y}} \mathcal{F}_{\overline{x}}.
$$
对 $\mathcal{F}$ 属于 $\Sh(X_\etale)$ 的情形成立，并且
$$
(f_*\mathcal{F})_{\overline{y}} =
\bigoplus\nolimits_{\overline{x} : \Spec(k) \to X,\ f(\overline{x}) =
\overline{y}} \mathcal{F}_{\overline{x}}.
$$
对 $\mathcal{F}$ 属于 $\textit{Ab}(X_\etale)$ 的情形成立。
\item 对任意 $q \geq 1$，有 $R^q f_*\mathcal{F} = 0$，
其中 $\mathcal{F}$ 属于 $\textit{Ab}(X_\etale)$。
\end{enumerate}
\end{proposition}

\begin{proof}
以 $X_{\overline{y}}^{sh}$ 表示纤维积
$X \times_Y \Spec(\mathcal{O}_{Y, \overline{y}}^{sh})$。
由定理 \ref{etale-cohomology-theorem-higher-direct-images}，
$R^qf_*\mathcal{F}$ 在 $\overline{y}$ 处的茎由
$H_\etale^q(X_{\overline{y}}^{sh}, \mathcal{F})$ 计算。
由于 $f$ 有限，$X_{\bar y}^{sh}$ 在
$\Spec(\mathcal{O}_{Y, \overline{y}}^{sh})$ 上有限，因而有
$X_{\bar y}^{sh} = \Spec(A)$，其中某个环 $A$ 在
$\mathcal{O}_{Y, \bar y}^{sh}$ 上有限。
由于后一个环严格 Hensel，引理 \ref{etale-cohomology-lemma-finite-over-henselian}
说明 $A$ 是 Hensel 局部环的有限乘积
$A = A_1 \times \ldots \times A_r$。由于
$\mathcal{O}_{Y, \overline{y}}^{sh}$ 的剩余域可分闭，每个
$A_i$ 也具有该性质。因此 $A_i$ 严格 Hensel。
这说明 $X_{\overline{y}}^{sh} = \coprod_{i = 1}^r \Spec(A_i)$。
引理 \ref{etale-cohomology-lemma-vanishing-etale-cohomology-strictly-henselian}
中的消失给出 $(R^qf_*\mathcal{F})_{\overline{y}} = 0$ 对
$q > 0$ 成立，再由定理 \ref{etale-cohomology-theorem-exactness-stalks} 得到 (2)。
(1) 由引理 \ref{etale-cohomology-lemma-vanishing-etale-cohomology-strictly-henselian}
中的相应陈述得到。
\end{proof}

\begin{lemma}
\label{etale-cohomology-lemma-finite-pushforward-commutes-with-base-change}
考虑笛卡尔方块
$$
\xymatrix{
X' \ar[r]_{g'} \ar[d]_{f'} & X \ar[d]^f \\
Y' \ar[r]^g & Y
}
$$
，其中 $f$ 是有限态射。
对任意集合层 $\mathcal{F}$（定义在 $X_\etale$ 上），有
$f'_*(g')^{-1}\mathcal{F} = g^{-1}f_*\mathcal{F}$。
\end{lemma}

\begin{proof}
在极大的一般性下，存在拉回映射
$g^{-1}f_*\mathcal{F} \to f'_*(g')^{-1}\mathcal{F}$；见
Sites，第 \ref{sites-section-pullback} 节。
只需在茎上验证（定理 \ref{etale-cohomology-theorem-exactness-stalks}）。
设 $\overline{y}' : \Spec(k) \to Y'$ 是几何点。有
\begin{align*}
(f'_*(g')^{-1}\mathcal{F})_{\overline{y}'}
& =
\prod\nolimits_{\overline{x}' : \Spec(k) \to X',\ f' \circ \overline{x}' =
\overline{y}'}
((g')^{-1}\mathcal{F})_{\overline{x}'} \\
& =
\prod\nolimits_{\overline{x}' : \Spec(k) \to X',\ f' \circ \overline{x}' =
\overline{y}'} \mathcal{F}_{g' \circ \overline{x}'} \\
& =
\prod\nolimits_{\overline{x} : \Spec(k) \to X,\ f \circ \overline{x} =
g \circ \overline{y}'} \mathcal{F}_{\overline{x}} \\
& =
(f_*\mathcal{F})_{g \circ \overline{y}'} \\
& =
(g^{-1}f_*\mathcal{F})_{\overline{y}'}
\end{align*}
第一个等式由命题 \ref{etale-cohomology-proposition-finite-higher-direct-image-zero} 得到。
第二个等式由引理 \ref{etale-cohomology-lemma-stalk-pullback} 得到。
第三个等式成立，是因为该图是笛卡尔方块，故映射
$$
\{\overline{x}' : \Spec(k) \to X',\ f' \circ \overline{x}' =
\overline{y}'\}
\longrightarrow
\{\overline{x} : \Spec(k) \to X,\ f \circ \overline{x} =
g \circ \overline{y}'\}
$$
把 $\overline{x}'$ 映到 $g' \circ \overline{x}'$，且是双射。
第四个等式由命题 \ref{etale-cohomology-proposition-finite-higher-direct-image-zero} 得到。
第五个等式由引理 \ref{etale-cohomology-lemma-stalk-pullback} 得到。
\end{proof}

\begin{lemma}
\label{etale-cohomology-lemma-integral-pushforward-commutes-with-base-change}
考虑笛卡尔方块
$$
\xymatrix{
X' \ar[r]_{g'} \ar[d]_{f'} & X \ar[d]^f \\
Y' \ar[r]^g & Y
}
$$
，其中 $f$ 是整态射。
对任意集合层 $\mathcal{F}$（定义在 $X_\etale$ 上），有
$f'_*(g')^{-1}\mathcal{F} = g^{-1}f_*\mathcal{F}$。
\end{lemma}

\begin{proof}
问题在 $Y$ 上是局部的，因而可以假设 $Y$ 仿射。
于是可写成 $X = \lim X_i$，其中 $f_i : X_i \to Y$ 有限
（在仿射情形这很容易，不过可参见
Limits，引理 \ref{limits-lemma-integral-limit-finite-and-finite-presentation}）。
以 $p_{i'i} : X_{i'} \to X_i$ 表示转移态射，并以
$p_i : X \to X_i$ 表示投影。置
$\mathcal{F}_i = p_{i, *}\mathcal{F}$，则由引理
\ref{etale-cohomology-lemma-linus-hamann} 得到系统
$(\mathcal{F}_i, \varphi_{i'i})$，它满足
$\mathcal{F} = \colim p_i^{-1}\mathcal{F}_i$。
由引理 \ref{etale-cohomology-lemma-relative-colimit} 得到
$f_*\mathcal{F} = \colim f_{i, *}\mathcal{F}_i$。
置 $X'_i = Y' \times_Y X_i$，其投影为 $f'_i$ 与 $g'_i$。
由于极限与极限交换，有 $X' = \lim X'_i$。
以 $p'_i : X' \to X'_i$ 表示投影。有
\begin{align*}
g^{-1}f_*\mathcal{F}
& =
g^{-1} \colim f_{i, *}\mathcal{F}_i \\
& =
\colim g^{-1}f_{i, *}\mathcal{F}_i \\
& =
\colim f'_{i, *}(g'_i)^{-1}\mathcal{F}_i \\
& =
f'_*(\colim (p'_i)^{-1}(g'_i)^{-1}\mathcal{F}_i) \\
& =
f'_*(\colim (g')^{-1}p_i^{-1}\mathcal{F}_i) \\
& =
f'_*(g')^{-1} \colim p_i^{-1}\mathcal{F}_i \\
& =
f'_*(g')^{-1}\mathcal{F}
\end{align*}
正合所求。第一个等式见上文。
第二个等式使用拉回与余极限交换。
第三个等式使用有限情形；见引理
\ref{etale-cohomology-lemma-finite-pushforward-commutes-with-base-change}。
第四个等式使用引理 \ref{etale-cohomology-lemma-relative-colimit}。
第五个等式使用 $g'_i \circ p'_i = p_i \circ g'$。
第六个等式使用拉回与余极限交换。
第七个等式使用 $\mathcal{F} = \colim p_i^{-1}\mathcal{F}_i$。
\end{proof}

\noindent
下一个引理是处理 \'etale 层与满有限态射的一种上同调下降。
证明正常基变换定理后，我们将大幅推广这一结果。

\begin{lemma}
\label{etale-cohomology-lemma-cohomological-descent-finite}
设 $f : X \to Y$ 是概形的满有限态射。
置 $f_n : X_n \to Y$ 为 $(n + 1)$ 重纤维积；该纤维积由
$X$ 在 $Y$ 上取得。
对 $\mathcal{F} \in \textit{Ab}(Y_\etale)$，置
$\mathcal{F}_n = f_{n, *}f_n^{-1}\mathcal{F}$。存在正合序列
$$
0 \to \mathcal{F} \to \mathcal{F}_0 \to \mathcal{F}_1 \to
\mathcal{F}_2 \to \ldots
$$
于 $Y_\etale$ 上。此外，存在谱序列
$$
E_1^{p, q} = H^q_\etale(X_p, f_p^{-1}\mathcal{F})
$$
收敛到 $H^{p + q}(Y_\etale, \mathcal{F})$。
该谱序列关于 $\mathcal{F}$ 函子性。
\end{lemma}

\begin{proof}
若证明了引理的第一个陈述，则得到以
$E_1^{p, q} = H^q_\etale(Y, \mathcal{F}_p)$ 为首项并收敛到
$H^{p + q}(Y_\etale, \mathcal{F})$ 的谱序列；见
Derived Categories，引理 \ref{derived-lemma-two-ss-complex-functor}。
另一方面，由于
$R^if_{p, *}f_p^{-1}\mathcal{F} = 0$ 对 $i > 0$ 成立
（命题 \ref{etale-cohomology-proposition-finite-higher-direct-image-zero}），得到
$$
H^q_\etale(X_p, f_p^{-1}\mathcal{F}) =
H^q_\etale(Y, f_{p, *}f_p^{-1} \mathcal{F}) =
H^q_\etale(Y, \mathcal{F}_p)
$$
；这是命题 \ref{etale-cohomology-proposition-leray} 的结论，从而得到引理中的谱序列。

\medskip\noindent
为证明引理的第一个陈述，注意 $X_n$ 构成 $Y$ 上的单纯概形；
见 Simplicial，例 \ref{simplicial-example-fibre-products-simplicial-object}。
还要注意，对每个投影 $d_j : X_{n + 1} \to X_n$，存在映射
$d_j^{-1} f_n^{-1}\mathcal{F} \to f_{n + 1}^{-1}\mathcal{F}$。
这些映射诱导映射
$$
\delta_j : \mathcal{F}_n \to \mathcal{F}_{n + 1}
$$
，其中 $j = 0, \ldots, n + 1$。我们用这些映射的交错和来定义微分
$\mathcal{F}_n \to \mathcal{F}_{n + 1}$。
类似地，有典范增广 $\mathcal{F} \to \mathcal{F}_0$；
它就是典范映射 $\mathcal{F} \to f_*f^{-1}\mathcal{F}$。
要验证这个层序列是正合复形，只需在几何点的茎上验证
（定理 \ref{etale-cohomology-theorem-exactness-stalks}）。
因此，设 $\overline{y} : \Spec(k) \to Y$ 是几何点。置
$E = \{\overline{x} : \Spec(k) \to X \mid f(\overline{x}) = \overline{y}\}$。
于是 $E$ 是有限非空集合，并且
$$
(\mathcal{F}_n)_{\overline{y}} =
\bigoplus\nolimits_{e \in E^{n + 1}} \mathcal{F}_{\overline{y}}
$$
；这是命题 \ref{etale-cohomology-proposition-finite-higher-direct-image-zero}
和引理 \ref{etale-cohomology-lemma-stalk-pullback} 的结论。
因此，必须验证：给定 Abel 群 $M$，序列
$$
0 \to M \to \bigoplus\nolimits_{e \in E} M \to
\bigoplus\nolimits_{e \in E^2} M \to \ldots
$$
正合。这里第一个映射是对角映射，而映射
$\bigoplus_{e \in E^{n + 1}} M  \to \bigoplus_{e \in E^{n + 2}} M$
是由 $(n + 2)$ 个投影 $E^{n + 2} \to E^{n + 1}$ 所诱导映射的交错和。
这可以直接证明，也可将 Simplicial，引理
\ref{simplicial-lemma-fibre-products-simplicial-object-w-section}
应用于映射 $E \to \{*\}$ 得到。
\end{proof}

\begin{remark}
\label{etale-cohomology-remark-cohomological-descent-finite}
在引理 \ref{etale-cohomology-lemma-cohomological-descent-finite} 的情形中，
若 $\mathcal{G}$ 是 $Y_\etale$ 上的集合层，则
$$
\Gamma(Y, \mathcal{G}) =
\text{等化子}(
\xymatrix{
\Gamma(X_0, f_0^{-1}\mathcal{G})
\ar@<1ex>[r] \ar@<-1ex>[r] &
\Gamma(X_1, f_1^{-1}\mathcal{G})
}
)
$$
证明完全相同：说明层 $\mathcal{G}$ 是两个映射
$f_{0, *}f_0^{-1}\mathcal{G} \to f_{1, *}f_1^{-1}\mathcal{G}$
的等化子即可。
\end{remark}










%10.06.09
\section{茎上的 Galois 作用}
\label{etale-cohomology-section-galois-action-stalks}

\noindent
本节定义如下作用：点 $s$（其所在概形为 $S$）的剩余域之绝对 Galois 群，
作用于任意位于 $s$ 上方的几何点处的茎函子。

\medskip\noindent
茎上的 Galois 作用。
设 $S$ 是概形。
设 $\overline{s}$ 是 $S$ 的几何点。
设 $\sigma \in \text{Aut}(\kappa(\overline{s})/\kappa(s))$。
如下定义 $\sigma$ 在茎 $\mathcal{F}_{\overline{s}}$
上的作用，其中 $\mathcal{F}$ 是层：
\begin{equation}
\label{etale-cohomology-equation-galois-action}
\begin{matrix}
\mathcal{F}_{\overline{s}} &
\longrightarrow &
\mathcal{F}_{\overline{s}} \\
(U, \overline{u}, t) &
\longmapsto &
(U, \overline{u} \circ \Spec(\sigma), t).
\end{matrix}
\end{equation}
这里使用定义 \ref{etale-cohomology-definition-stalk} 后讨论中以三元组描述茎元素的方法。
这是左作用，因为若
$\sigma_i \in \text{Aut}(\kappa(\overline{s})/\kappa(s))$，则
\begin{align*}
\sigma_1 \cdot (\sigma_2 \cdot (U, \overline{u}, t))
& =
\sigma_1 \cdot (U, \overline{u} \circ \Spec(\sigma_2), t) \\
& =
(U, \overline{u} \circ \Spec(\sigma_2) \circ \Spec(\sigma_1), t) \\
& =
(U, \overline{u} \circ \Spec(\sigma_1 \circ \sigma_2), t) \\
& =
(\sigma_1 \circ \sigma_2) \cdot (U, \overline{u}, t)
\end{align*}
显然，这个作用关于层 $\mathcal{F}$ 是函子性的。
也可以直接引用注 \ref{etale-cohomology-remark-map-stalks} 来定义这个作用。

\begin{definition}
\label{etale-cohomology-definition-algebraic-geometric-point}
设 $S$ 是概形。
设 $\overline{s}$ 是位于点 $s$（属于 $S$）上方的几何点。
以
$\kappa(s) \subset \kappa(s)^{sep} \subset \kappa(\overline{s})$
表示 $\kappa(s)$ 在代数闭域 $\kappa(\overline{s})$ 中的可分代数闭包。
\begin{enumerate}
\item 在此情形，$\kappa(s)$ 的{\it 绝对 Galois 群}是
$\text{Gal}(\kappa(s)^{sep}/\kappa(s))$。有时记作
$\text{Gal}_{\kappa(s)}$。
\item 若几何点 $\overline{s}$ 所对应的扩张
$\kappa(s) \subset \kappa(\overline{s})$ 是 $\kappa(s)$ 的代数闭包，
则称它是{\it 代数的}。
\end{enumerate}
\end{definition}

\begin{example}
\label{etale-cohomology-example-stupid}
几何点
$\Spec(\mathbf{C}) \to \Spec(\mathbf{Q})$
不是代数的。
\end{example}

\noindent
令
$\kappa(s) \subset \kappa(s)^{sep} \subset \kappa(\overline{s})$
如定义中所述。注意，由于 $\kappa(\overline{s})$ 代数闭，映射
$$
\text{Aut}(\kappa(\overline{s})/\kappa(s))
\longrightarrow
\text{Gal}(\kappa(s)^{sep}/\kappa(s)) = \text{Gal}_{\kappa(s)}
$$
是满射。设 $(U, \overline{u})$ 是 $\overline{s}$ 的
\'etale 邻域，并设 $\overline{u}$ 位于点 $u$（属于 $U$）上方。
由于 $U \to S$ 是 \'etale 的，剩余域扩张
$\kappa(u)/\kappa(s)$ 有限可分。
由此得到
\begin{enumerate}
\item 若 $\sigma \in \text{Aut}(\kappa(\overline{s})/\kappa(s)^{sep})$，
则 $\sigma$ 在 $\mathcal{F}_{\overline{s}}$ 上的作用平凡。
\item 更精确地，
$\text{Aut}(\kappa(\overline{s})/\kappa(s))$ 的作用决定绝对 Galois 群
$\text{Gal}_{\kappa(s)}$ 在 $\mathcal{F}_{\overline{s}}$
上的作用，并且由后者决定。
\item 给定 $(U, \overline{u}, t)$，它表示元素 $\xi$，而该元素属于
$\mathcal{F}_{\overline{s}}$；群
$\text{Gal}(\kappa(s)^{sep}/K)$ 的任意元素作用平凡，其中
$\kappa(s) \subset K \subset \kappa(s)^{sep}$ 是
$\overline{u}^\sharp : \kappa(u) \to \kappa(\overline{s})$ 的像。
\end{enumerate}
综上，$\mathcal{F}_{\overline{s}}$ 成为
$\text{Gal}_{\kappa(s)}$-集合（见 Fundamental Groups，定义
\ref{pione-definition-G-set-continuous}）。
因此，可以把茎函子看作函子
$$
\Sh(S_\etale) \longrightarrow
\text{Gal}_{\kappa(s)}\textit{-Sets},
\quad
\mathcal{F} \longmapsto \mathcal{F}_{\overline{s}}
$$
；此后通常都这样看待茎函子。

\begin{theorem}
\label{etale-cohomology-theorem-equivalence-sheaves-point}
设 $S = \Spec(K)$，其中 $K$ 是域。
设 $\overline{s}$ 是 $S$ 的几何点。
以 $G = \text{Gal}_{\kappa(s)}$ 表示绝对 Galois 群。
取茎诱导范畴等价
$$
\Sh(S_\etale) \longrightarrow G\textit{-Sets},
\quad
\mathcal{F} \longmapsto \mathcal{F}_{\overline{s}}.
$$
\end{theorem}

\begin{proof}
构造这个函子的逆。在 Fundamental Groups，引理
\ref{pione-lemma-sheaves-point} 中已经看到：给定 $G$-集合 $M$，
存在 \'etale 态射
$X \to \Spec(K)$，
使得 $\Mor_K(\Spec(K^{sep}), X)$ 同构于 $M$，且这是 $G$-集合的同构。
考虑层 $\mathcal{F}$，它定义在 $\Spec(K)_\etale$ 上，并由规则
$U \mapsto \Mor_K(U, X)$ 定义。由于 \'etale 拓扑是次典范的，
这确实是层。于是可见
$\mathcal{F}_{\overline{s}} = \Mor_K(\Spec(K^{sep}), X) = M$
作为 $G$-集合成立（略去细节）。这给出了该函子的逆，结论得证。
\end{proof}

\begin{remark}
\label{etale-cohomology-remark-every-sheaf-representable}
定理 \ref{etale-cohomology-theorem-equivalence-sheaves-point} 与 Fundamental Groups，引理
\ref{pione-lemma-sheaves-point} 的结论也可以表述为：
$\Spec(K)_\etale$ 上的每个层都可由概形 $X$ 表示，且它
在 $\Spec(K)$ 上是 \'etale 的。
这并不意味着每个层在 Sites，定义
\ref{sites-definition-representable-sheaf} 的意义下都可表示。
原因是，在构造 $\Spec(K)_\etale$ 时，我们选取了一个充分大的、在
$\Spec(K)$ 上 \'etale 的概形集合，而
$\Spec(K)_\etale$ 上的层构成真类。
\end{remark}

\begin{lemma}
\label{etale-cohomology-lemma-global-sections-point}
假设与记号同定理 \ref{etale-cohomology-theorem-equivalence-sheaves-point}。
存在函子性双射
$$
\Gamma(S, \mathcal{F}) = (\mathcal{F}_{\overline{s}})^G
$$
\end{lemma}

\begin{proof}
可以使用形式论证和定理 \ref{etale-cohomology-theorem-equivalence-sheaves-point}
的结果来证明。给定层 $\mathcal{F}$，它对应于
$G$-集合 $M = \mathcal{F}_{\overline{s}}$，有
\begin{eqnarray*}
\Gamma(S, \mathcal{F}) & = &
\Mor_{\Sh(S_\etale)}(h_{\Spec(K)},  \mathcal{F})
\\
& = & \Mor_{G\textit{-Sets}}(\{*\}, M) \\
& = & M^G
\end{eqnarray*}
这里第一个等同由 Sites，第 \ref{sites-section-presheaves} 节和
第 \ref{sites-section-representable-sheaves} 节解释；
第二个由定理 \ref{etale-cohomology-theorem-equivalence-sheaves-point} 得到；
第三个显然。还给出一个直接证明\footnote{献给仍心存疑虑的读者。}。

\medskip\noindent
设 $t \in \Gamma(S, \mathcal{F})$ 是整体截面。
则三元组 $(S, \overline{s}, t)$ 定义
$\mathcal{F}_{\overline{s}}$ 的一个元素；它显然在 $G$ 的作用下不变。
反之，设 $(U, \overline{u}, t)$ 定义
$\mathcal{F}_{\overline{s}}$ 的一个不变元素。
可以缩小 $U$，并假设 $U = \Spec(L)$，其中
$K$ 有一个有限可分域扩张；见命题 \ref{etale-cohomology-proposition-etale-morphisms}。
在这种情形，映射 $\mathcal{F}(U) \to \mathcal{F}_{\overline{s}}$
是单射，因为对 \'etale 邻域的任意态射
$(U', \overline{u}') \to (U, \overline{u})$，限制映射
$\mathcal{F}(U) \to \mathcal{F}(U')$ 是单射；这是因为
$U' \to U$ 是 $S_\etale$ 的覆盖。
稍微扩大 $L$ 后，可以假设 $K \subset L$ 是有限 Galois 扩张。
此时使用
$$
\Spec(L) \times_{\Spec(K)} \Spec(L)
=
\coprod\nolimits_{\sigma \in \text{Gal}(L/K)} \Spec(L)
$$
，其中映射 $\Spec(L) \to \Spec(L \otimes_K L)$
来自环映射 $a \otimes b \mapsto a\sigma(b)$。因此可见，
$(U, \overline{u}, t)$ 在整个 $G$ 下不变这一条件说明
$t \in \mathcal{F}(\Spec(L))$ 通过任一投影作限制后，
都映到
$\mathcal{F}(\Spec(L) \times_{\Spec(K)} \Spec(L))$
的同一个元素（这里使用了上述单射性；略去细节）。
于是，层 $\mathcal{F}$ 关于 \'etale 覆盖
$\{\Spec(L) \to \Spec(K)\}$ 的层条件发挥作用，得到
$t$ 来自 $\mathcal{F}$ 在 $\Spec(K)$ 上的唯一截面。
\end{proof}

\begin{remark}
\label{etale-cohomology-remark-stalk-pullback}
设 $S$ 是概形，并设 $\overline{s} : \Spec(k) \to S$
是 $S$ 的几何点。按定义，这意味着 $k$ 代数闭。
特别地，$k$ 的绝对 Galois 群平凡。因此，由定理
\ref{etale-cohomology-theorem-equivalence-sheaves-point}，
$\Spec(k)_\etale$ 上的层范畴等价于集合范畴。
该等价由取 $\Spec(k)$ 上的截面给出。
这最终给出了茎函子的另一种定义。具体而言，函子
$$
\Sh(S_\etale) \longrightarrow \textit{Sets}, \quad
\mathcal{F} \longmapsto \mathcal{F}_{\overline{s}}
$$
同构于函子
$$
\Sh(S_\etale)
\longrightarrow
\Sh(\Spec(k)_\etale) = \textit{Sets},
\quad
\mathcal{F} \longmapsto \overline{s}^*\mathcal{F}
$$
要严格证明这一点，可以使用引理 \ref{etale-cohomology-lemma-stalk-pullback} 的 (3)，
其中 $f = \overline{s}$。此外，由此，引理
\ref{etale-cohomology-lemma-stalk-pullback} 的 (3) 的一般情形
也可从拉回的函子性得到。
\end{remark}






\section{群上同调}
\label{etale-cohomology-section-group-cohomology}

\noindent
下文若写 $H^i(G, M)$，其含义是：$G$ 是拓扑群，
$M$ 是具有连续 $G$-作用的离散 $G$-模，而 $H^i(G, -)$ 是
第 $i$ 个右导出函子，定义在范畴 $\text{Mod}_G$ 上；
该范畴由这样的 $G$-模组成；
见定义 \ref{etale-cohomology-definition-G-module-continuous} 和
\ref{etale-cohomology-definition-galois-cohomology}。这包括抽象群 $G$ 的情形，
其含义就是把 $G$ 看成具有离散拓扑的拓扑群。

\medskip\noindent
当模具有非离散拓扑时，将使用记号
$H^i_{cont}(G, M)$ 表示 \cite{Tate} 中引入的连续上同调群；
见第 \ref{etale-cohomology-section-continuous-group-cohomology} 节。

\begin{definition}
\label{etale-cohomology-definition-G-module-continuous}
设 $G$ 是拓扑群。
\begin{enumerate}
\item {\it $G$-模}（有时称为{\it 离散 $G$-模}）是 Abel 群
$M$，它带有由群同态给出的左作用 $a : G \times M \to M$，
并且 $a$ 连续（此时给 $M$ 配备离散拓扑）。
\item {\it $G$-模的态射} $f : M \to N$ 是
$G$-等变同态，其源为 $M$、靶为 $N$。
\item $G$-模的范畴记作 $\text{Mod}_G$。
\end{enumerate}
设 $R$ 是环。
\begin{enumerate}
\item {\it $R\text{-}G$-模}是 $R$-模 $M$，它带有左作用
$a : G \times M \to M$，该作用由 $R$-线性映射给出；并且
$a$ 连续（此时给 $M$ 配备离散拓扑）。
\item {\it $R\text{-}G$-模的态射} $f : M \to N$ 是
$G$-等变 $R$-模映射，其源为 $M$、靶为 $N$。
\item $R\text{-}G$-模的范畴记作 $\text{Mod}_{R, G}$。
\end{enumerate}
\end{definition}

\noindent
条件 $a : G \times M \to M$ 连续，等价于任意 $x \in M$
的稳定子在 $G$ 中开。若 $G$ 是抽象群，则在给 $G$ 配备离散拓扑后，
这就对应于带 $G$-作用的 Abel 群这一概念。注意
$\text{Mod}_{\mathbf{Z}, G} = \text{Mod}_G$。

\medskip\noindent
范畴 $\text{Mod}_G$ 有足够多内射对象；见 Injectives，引理
\ref{injectives-lemma-G-modules}。考虑左正合函子
$$
\text{Mod}_G \longrightarrow \textit{Ab},
\quad
M \longmapsto M^G =
\{x \in M \mid g \cdot x = x\ \forall g \in G\}
$$
有时记 $M^G = H^0(G, M)$，有时写作
$M^G = \Gamma_G(M)$。该函子有全右导出函子
$R\Gamma_G(M)$，并且其第 $i$ 个右导出函子为
$R^i\Gamma_G(M) = H^i(G, M)$，对任意 $i \geq 0$ 成立。

\medskip\noindent
同一构造也适用于
$H^0(G, -) : \text{Mod}_{R, G} \to \text{Mod}_R$。
引理 \ref{etale-cohomology-lemma-modules-abelian} 将说明，这与底层 $G$-模的上同调一致。

\begin{definition}
\label{etale-cohomology-definition-galois-cohomology}
设 $G$ 是拓扑群。设 $M$ 是具有连续 $G$-作用的离散 $G$-模。
换言之，$M$ 是定义 \ref{etale-cohomology-definition-G-module-continuous}
中引入的范畴 $\text{Mod}_G$ 的对象。
\begin{enumerate}
\item 右导出函子 $H^i(G, M)$（它导自 $H^0(G, M)$，并定义在范畴
$\text{Mod}_G$ 上）称为 $M$ 的{\it 连续群上同调群}。
\item 若 $G$ 是配备离散拓扑的抽象群，则
$H^i(G, M)$ 称为 $M$ 的{\it 群上同调群}。
\item 若 $G$ 是 Galois 群，则群 $H^i(G, M)$ 称为
$M$ 的{\it Galois 上同调群}。
\item 若 $G$ 是域 $K$ 的绝对 Galois 群，则群
$H^i(G, M)$ 有时称为{\it $K$ 的、以 $M$ 为系数的 Galois 上同调群}。
在此情形，有时写
$H^i(K, M)$ 来代替 $H^i(G, M)$。
\end{enumerate}
\end{definition}

\begin{lemma}
\label{etale-cohomology-lemma-modules-abelian}
设 $G$ 是拓扑群。设 $R$ 是环。
对每个 $i \geq 0$，图
$$
\xymatrix{
\text{Mod}_{R, G} \ar[rr]_{H^i(G, -)} \ar[d] & &
\text{Mod}_R \ar[d] \\
\text{Mod}_G \ar[rr]^{H^i(G, -)} & &
\textit{Ab}
}
$$
交换，其中竖直箭头是遗忘函子。
\end{lemma}

\begin{proof}
以 $F : \text{Mod}_{R, G} \to \text{Mod}_G$ 表示遗忘函子。
则 $F$ 有左伴随
$H : \text{Mod}_G \to \text{Mod}_{R, G}$，由
$H(M) = M \otimes_\mathbf{Z} R$ 给出。注意，
$\text{Mod}_G$ 的每个对象都是形如
$\mathbf{Z}[G/U]$ 的模的直和之商，其中 $U \subset G$ 是开子群。
这里 $\mathbf{Z}[G/U]$ 表示 $G$-模；它由有限
$\mathbf{Z}$-线性组合组成，所组合的是右 $U$ 同余类（位于 $G$ 中），
并带有左 $G$-作用。
因此，$\text{Mod}_G$ 中每个向上有界复形都拟同构于
$\text{Mod}_G$ 中的一个向上有界复形，后者的底层各项是平坦
$\mathbf{Z}$-模
（Derived Categories，引理 \ref{derived-lemma-subcategory-left-resolution}）。
于是显然，$LH$ 存在于 $D^-(\text{Mod}_G)$ 上，并可在任意复形上逐项
计算 $H$，只要其各项是平坦 $\mathbf{Z}$-模；这由 Derived Categories，引理
\ref{derived-lemma-subcategory-left-acyclics} 和命题
\ref{derived-proposition-enough-acyclics} 得到。
由 Derived Categories，引理
\ref{derived-lemma-pre-derived-adjoint-functors} 可知
$$
\text{Ext}^i(\mathbf{Z}, F(M)) = \text{Ext}^i(R, M)
$$
对 $M$（它属于 $\text{Mod}_{R, G}$）成立。
注意，有
$H^0(G, -) = \Hom(\mathbf{Z}, -)$ 于 $\text{Mod}_G$ 上，其中
$\mathbf{Z}$ 表示具有平凡作用的 $G$-模。因此，有
$H^i(G, -) = \text{Ext}^i(\mathbf{Z}, -)$ 于 $\text{Mod}_G$ 上。
类似地，有 $H^i(G, -) = \text{Ext}^i(R, -)$ 于
$\text{Mod}_{R, G}$ 上。综合以上即得引理。
\end{proof}

\begin{lemma}
\label{etale-cohomology-lemma-ext-modules-hom}
设 $G$ 是拓扑群。设 $R$ 是环。
设 $M$、$N$ 是 $R\text{-}G$-模。若 $M$ 作为 $R$-模有限投射，则
$\text{Ext}^i(M, N) = H^i(G, M^\vee \otimes_R N)$
（记号见证明）。
\end{lemma}

\begin{proof}
模 $M^\vee = \Hom_R(M, R)$ 带有 $G$ 的逆步作用。
具体而言，有 $(g \cdot \lambda)(m) = \lambda(g^{-1} \cdot m)$，
其中 $g \in G$、$\lambda \in M^\vee$、$m \in M$。
$G$ 在 $M^\vee \otimes_R N$ 上的作用是对角作用，即
$g \cdot (\lambda \otimes n) = g \cdot \lambda \otimes g \cdot n$。
注意，对第三个 $R\text{-}G$-模 $E$，有
$\Hom(E, M^\vee \otimes_R N) = \Hom(M \otimes_R E, N)$。
事实上，由 Algebra，引理 \ref{algebra-lemma-hom-from-tensor-product}
和 \ref{algebra-lemma-evaluation-map-iso-finite-projective}，
这在 $R$-模层面成立；而 $G$-作用的定义正是为了使它对
$R\text{-}G$-模仍然成立。因此，$M^\vee \otimes_R N$
是内射 $R\text{-}G$-模，只要 $N$ 是内射 $R\text{-}G$-模。
所以，若 $N \to N^\bullet$ 是内射分解，则
$M^\vee \otimes_R N \to M^\vee \otimes_R N^\bullet$
也是内射分解。于是
$$
\Hom(M, N^\bullet) = \Hom(R, M^\vee \otimes_R N^\bullet) =
(M^\vee \otimes_R N^\bullet)^G
$$
由于左边计算 $\text{Ext}^i(M, N)$，右边计算
$H^i(G, M^\vee \otimes_R N)$，证明完成。
\end{proof}

\begin{lemma}
\label{etale-cohomology-lemma-finite-dim-group-cohomology}
设 $G$ 是拓扑群。设 $k$ 是域。
设 $V$ 是 $k\text{-}G$-模。
若 $G$ 拓扑有限生成且
$\dim_k(V) < \infty$，则 $\dim_k H^1(G, V) < \infty$。
\end{lemma}

\begin{proof}
设 $g_1, \ldots, g_r \in G$ 是拓扑生成 $G$ 的元素；
也就是说，由 $g_1, \ldots, g_r$ 生成的子群稠密。
由引理 \ref{etale-cohomology-lemma-ext-modules-hom} 可见，
$H^1(G, V)$ 是 $k$-向量空间，其元素是扩张
$$
0 \to V \to E \to k \to 0
$$
；这些扩张属于 $k\text{-}G$-模。选取 $e \in E$，它映到 $1 \in k$。
写成
$$
g_i \cdot e = v_i + e
$$
，其中某个 $v_i \in V$。这是可能的，因为 $g_i \cdot 1 = 1$。
断言元素列表 $v_1, \ldots, v_r \in V$ 决定扩张 $E$ 的同构类。
一旦证明这一点，引理即得，因为这意味着我们的 Ext 向量空间同构于
$k$-向量空间 $V^{\oplus r}$ 的一个子商；略去一些细节。
由于 $E$ 是定义 \ref{etale-cohomology-definition-G-module-continuous} 中定义的范畴的对象，
知道存在开子群 $U$，使得 $u \cdot e = e$
对所有 $u \in U$ 成立。
现任取 $g \in G$。则 $gU$ 包含一个字 $w$，它由元素
$g_1, \ldots, g_r$ 组成。
设 $gu = w$。由于元素 $w \cdot e$ 由
$v_1, \ldots, v_r$ 决定，可见
$g \cdot e = (gu) \cdot e = w \cdot e$ 也由其决定。
\end{proof}

\begin{lemma}
\label{etale-cohomology-lemma-profinite-group-cohomology-torsion}
设 $G$ 是预有限拓扑群。则
\begin{enumerate}
\item $H^i(G, M)$ 是挠群，只要 $i > 0$，且任取
$G$-模 $M$；并且
\item $H^i(G, M) = 0$，若 $M$ 是 $\mathbf{Q}$-向量空间。
\end{enumerate}
\end{lemma}

\begin{proof}
证明 (1)。由维数移动，只需证明 $H^1(G, M)$ 是挠群，
其中任取 $G$-模 $M$。
选取正合序列 $0 \to M \to I \to N \to 0$，其中 $I$
是 $G$-模范畴的内射对象。
则 $H^1(G, M)$ 的任意元素都是某个元素
$y \in N^G$ 的像。选取 $x \in I$，它映到 $y$。
稳定子 $U \subset G$（关于 $x$）是开的，因此具有有限指数 $r$。
设 $g_1, \ldots, g_r \in G$ 是 $G/U$ 的一组代表元。
于是 $\sum g_i(x)$ 是 $I$ 的不变元素，映到 $ry$。
因此 $r$ 消去最初选取的 $H^1(G, M)$ 的元素。
(2) 得证，因为此时 $H^i(G, M)$ 既是 $\mathbf{Q}$-向量空间又是挠群。
\end{proof}





\section{Tate 连续上同调}
\label{etale-cohomology-section-continuous-group-cohomology}

\noindent
Tate 连续上同调（\cite{Tate}）由连续非齐次上链复形定义。
当 $M$ 是带有连续 $G$-作用的任意拓扑 Abel 群时，可以作此定义。
具体而言，考虑复形
$$
C^\bullet_{cont}(G, M) :
M \to \text{Maps}_{cont}(G, M) \to
\text{Maps}_{cont}(G \times G, M) \to \ldots
$$
其中对 $n \geq 1$，边界映射由下式定义：
\begin{align*}
\text{d}(f)(g_1, \ldots, g_{n + 1})
& = g_1(f(g_2, \ldots, g_{n + 1})) \\
&
+ \sum\nolimits_{j = 1, \ldots, n}
(-1)^jf(g_1, \ldots, g_jg_{j + 1}, \ldots, g_{n + 1}) \\
&
+ (-1)^{n + 1}f(g_1, \ldots, g_n)
\end{align*}
而当 $n = 0$ 时，它把 $m \in M$ 映到映射
$g \mapsto g(m) - m$。定义
$$
H^i_{cont}(G, M) = H^i(C^\bullet_{cont}(G, M))
$$
由于该复形的各项涉及从 $G$ 以及 $G$ 的自身乘积到拓扑模
$M$ 的连续映射，并不清楚它是否把拓扑模的短正合序列
变为上同调长正合序列。另一个困难是拓扑 Abel 群范畴并非 Abel 范畴！

\medskip\noindent
不过，离散 $G$-模的短正合序列确实给出连续上链复形的短正合序列，
从而给出连续上同调群 $H^i_{cont}(G, -)$ 的上同调长正合序列。
因此，在定义 \ref{etale-cohomology-definition-G-module-continuous} 的范畴
$\text{Mod}_G$ 上，函子 $H^i_{cont}(G, M)$ 构成上同调
$\delta$-函子，其定义见 Homology，第
\ref{homology-section-cohomological-delta-functor} 节。
由于定义 \ref{etale-cohomology-definition-galois-cohomology} 的上同调
$H^i(G, M)$ 是泛 $\delta$-函子
（Derived Categories，引理 \ref{derived-lemma-right-derived-delta-functor}），
得到典范映射
$$
H^i(G, M) \longrightarrow H^i_{cont}(G, M)
$$
，其中 $M \in \text{Mod}_G$。已知当 $G$ 是抽象群
（即 $G$ 具有离散拓扑），或者当 $G$ 是预有限群时，
这些映射是同构（此处插入未来的参考文献）。
若你知道某个例子，说明对拓扑群 $G$ 和
$M \in \Ob(\text{Mod}_G)$，此映射不是同构，请发送电子邮件至
\href{mailto:stacks.project@gmail.com}{stacks.project@gmail.com}。




\section{点的上同调}
\label{etale-cohomology-section-cohomology-point}

\noindent
作为前几节讨论的推论，得到域的谱的 \'etale 上同调与 Galois 上同调的等价。

\begin{lemma}
\label{etale-cohomology-lemma-equivalence-abelian-sheaves-point}
设 $S = \Spec(K)$，其中 $K$ 是域。
设 $\overline{s}$ 是 $S$ 的几何点。
以 $G = \text{Gal}_{\kappa(s)}$ 表示绝对 Galois 群。
茎函子诱导范畴等价
$$
\textit{Ab}(S_\etale) \longrightarrow \text{Mod}_G,
\quad
\mathcal{F} \longmapsto \mathcal{F}_{\overline{s}}.
$$
\end{lemma}

\begin{proof}
在定理 \ref{etale-cohomology-theorem-equivalence-sheaves-point} 中，
已经看到集合层与 $G$-集合之间的等价。
本引理由此形式地得到，因为 Abel 层无非是带有交换群律的集合层，
而 $G$-模无非是带有交换群律的 $G$-集合。
\end{proof}

\begin{lemma}
\label{etale-cohomology-lemma-compare-cohomology-point}
记号与假设同引理 \ref{etale-cohomology-lemma-equivalence-abelian-sheaves-point}。
设 $\mathcal{F}$ 是 $\Spec(K)_\etale$ 上的 Abel 层，
对应于 $G$-模 $M$。则
\begin{enumerate}
\item 在 $D(\textit{Ab})$ 中有典范同构
$R\Gamma(S, \mathcal{F}) = R\Gamma_G(M)$；
\item $H_\etale^0(S, \mathcal{F}) = M^G$；并且
\item $H_\etale^q(S, \mathcal{F}) = H^q(G, M)$。
\end{enumerate}
\end{lemma}

\begin{proof}
结合引理 \ref{etale-cohomology-lemma-equivalence-abelian-sheaves-point}
与引理 \ref{etale-cohomology-lemma-global-sections-point}。
\end{proof}

\begin{example}
\label{etale-cohomology-example-sheaves-point}
$\Spec(K)_\etale$ 上的层。
设 $G = \text{Gal}(K^{sep}/K)$ 是 $K$ 的绝对 Galois 群。
\begin{enumerate}
\item 常值层 $\underline{\mathbf{Z}/n\mathbf{Z}}$ 对应于模
$\mathbf{Z}/n\mathbf{Z}$，其 $G$-作用平凡；
\item 层 $\mathbf{G}_m|_{\Spec(K)_\etale}$ 对应于
$(K^{sep})^*$，带有其 $G$-作用；
\item 层 $\mathbf{G}_a|_{\Spec(K^{sep})}$ 对应于
$(K^{sep}, +)$，带有其 $G$-作用；并且
\item 层 $\mu_n|_{\Spec(K^{sep})}$ 对应于
$\mu_n(K^{sep})$，带有其 $G$-作用。
\end{enumerate}
由注 \ref{etale-cohomology-remark-special-case-fpqc-cohomology-quasi-coherent}
和定理 \ref{etale-cohomology-theorem-picard-group}，有如下上同调群等同：
\begin{align*}
H_\etale^0(S_\etale, \mathbf{G}_m) & =
\Gamma(S, \mathcal{O}_S^*) \\
H_\etale^1(S_\etale, \mathbf{G}_m) & =
H_{Zar}^1(S, \mathcal{O}_S^*) = \Pic(S) \\
H_\etale^i(S_\etale, \mathbf{G}_a) & =
H_{Zar}^i(S, \mathcal{O}_S)
\end{align*}
此外，对任意拟凝聚层 $\mathcal{F}$（定义在 $S_\etale$ 上），有
$$
H^i(S_\etale, \mathcal{F}) = H_{Zar}^i(S, \mathcal{F}),
$$
见定理 \ref{etale-cohomology-theorem-zariski-fpqc-quasi-coherent}。
特别地，这给出等式序列
$$
0 =
\Pic(\Spec(K)) =
H_\etale^1(\Spec(K)_\etale, \mathbf{G}_m) =
H^1(G, (K^{sep})^*)
$$
，它正是 Hilbert 定理 90。类似地，对 $i \geq 1$，
$$
0 = H^i(\Spec(K), \mathcal{O})
= H_\etale^i(\Spec(K)_\etale, \mathbf{G}_a)
= H^i(G, K^{sep})
$$
其中 $K^{sep}$ 表示作为 Galois 模的 $K^{sep}$，其群律为加法。
这样，可以把迄今所做的工作看作计算 Galois 上同调群的一种复杂方法。
\end{example}

\noindent
下一个结果是一则趣闻，初读时应跳过。

\begin{lemma}
\label{etale-cohomology-lemma-all-modules-quasi-coherent}
设 $R$ 是维数 $0$ 的局部环。设 $S = \Spec(R)$。
则每个 $\mathcal{O}_S$-模（定义在 $S_\etale$ 上）都是拟凝聚的。
\end{lemma}

\begin{proof}
设 $\mathcal{F}$ 是 $\mathcal{O}_S$-模，定义在 $S_\etale$ 上。
必须证明 $\mathcal{F}$ 由 $R$-模
$M = \Gamma(S, \mathcal{F})$ 决定。
更精确地，若 $\pi : X \to S$ 是 \'etale 的，必须证明
$\Gamma(X, \mathcal{F}) = \Gamma(X, \pi^*\widetilde{M})$。

\medskip\noindent
设 $\mathfrak m \subset R$ 是极大理想，并设
$\kappa$ 是剩余域。由 Algebra，引理
\ref{algebra-lemma-local-dimension-zero-henselian}，
局部环 $R$ 是 Hensel 的。若 $X \to S$ 是 \'etale 的，
则由 Morphisms，引理 \ref{morphisms-lemma-etale-over-field}，
$X$ 的底层拓扑空间离散，因而 $X$ 是仿射概形的不交并，
且每个分支只有一个点。此外，若 $X = \Spec(A)$ 仿射且只有一个点，
则由 Algebra，引理 \ref{algebra-lemma-mop-up}，
$R \to A$ 是有限 \'etale 的。
在此情形，必须证明 $\Gamma(X, \mathcal{F}) = M \otimes_R A$。

\medskip\noindent
函子 $A \mapsto A/\mathfrak m A$ 定义如下范畴等价：
有限 \'etale $R$-代数的范畴等价于有限可分 $\kappa$-代数的范畴；
见 Algebra，引理 \ref{algebra-lemma-henselian-cat-finite-etale}。
先考虑 $A/\mathfrak m A$ 是 $\kappa$ 的 Galois 扩张且
Galois 群为 $G$ 的情形。
对每个 $\sigma \in G$，以 $\sigma : A \to A$ 表示
$A$ 在 $R$ 上的相应自同构。
置 $N = \Gamma(X, \mathcal{F})$。
于是 $\Spec(\sigma) : X \to X$ 是 $S$ 上的自同构，
故由其拉回定义映射 $\sigma : N \to N$；
这是 $\sigma$-线性映射：
$\sigma(an) = \sigma(a) \sigma(n)$，其中 $a \in A$ 且 $n \in N$。
将 Galois 下降应用于拟凝聚模
$\widetilde{N}$（定义在 $X$ 上），它配备来自 $\sigma$ 在 $N$ 上作用的同构。
见 Descent，引理 \ref{descent-lemma-galois-descent-more-general}。
该引理说明存在同构 $N = N^G \otimes_R A$。
另一方面，显然 $N^G = M$；这是 $\mathcal{F}$ 的层条件。
因此所需同构成立。

\medskip\noindent
一般情形（$A$ 局部且在 $R$ 上有限 \'etale）如下归结为
Galois 情形。选取有限 \'etale 的 $A \to B$，使得 $B$ 局部，
其剩余域在 $\kappa$ 上 Galois。置
$G = \text{Aut}(B/R) = \text{Gal}(\kappa_B/\kappa)$。
设 $H \subset G$ 是对应于 Galois 扩张
$\kappa_B/\kappa_A$ 的 Galois 群。于是如上可证
$\Gamma(X, \mathcal{F}) = \Gamma(\Spec(B), \mathcal{F})^H$。
将 Galois 扩张的结果使用两次，得到
$$
\Gamma(X, \mathcal{F}) = (M \otimes_R B)^H = M \otimes_R A
$$
正合所求。
\end{proof}








\section{曲线的上同调}
\label{etale-cohomology-section-cohomology-curves}

\noindent
眼下的下一个任务是计算代数闭域上光滑曲线取挠系数的 \'etale 上同调，
特别要证明它在次数至少为 3 时消失。为此，将计算一般点处的上同调，
这归结为某些 Galois 上同调。





\section{布劳尔群}
\label{etale-cohomology-section-brauer-groups}

\noindent
域的布劳尔群用有限维中心单代数定义。
本节回顾布劳尔群的相关事实，其中大部分在 Brauer Groups，第
\ref{brauer-section-introduction} 节中讨论。
其他参考文献见 \cite{SerreCorpsLocaux}、
\cite{SerreGaloisCohomology} 或 \cite{Weil}。

\begin{theorem}
\label{etale-cohomology-theorem-central-simple-algebra}
设 $K$ 是域。对含幺、结合（不一定交换）的 $K$-代数 $A$，
以下条件等价：
\begin{enumerate}
\item $A$ 是有限维中心单 $K$-代数；
\item $A$ 是有限维 $K$-向量空间，$K$ 是 $A$ 的中心，并且
$A$ 没有非平凡双边理想；
\item 存在 $d \geq 1$，使得
$A \otimes_K \bar K \cong \text{Mat}(d \times d, \bar K)$；
\item 存在 $d \geq 1$，使得
$A \otimes_K K^{sep} \cong \text{Mat}(d \times d, K^{sep})$；
\item 存在 $d \geq 1$ 以及有限 Galois 扩张 $K'/K$，使得
$A \otimes_K K' \cong \text{Mat}(d \times d, K')$；
\item 存在 $n \geq 1$ 以及有限维中心斜域 $D$（定义在 $K$ 上），使得
$A \cong \text{Mat}(n \times n, D)$。
\end{enumerate}
整数 $d$ 称为 $A$ 的{\it 次数}。
\end{theorem}

\begin{proof}
这是 Brauer Groups，引理
\ref{brauer-lemma-finite-central-simple-algebra} 的复述。
\end{proof}

\begin{lemma}
\label{etale-cohomology-lemma-brauer-inverse}
设 $A$ 是 $K$ 上的有限维中心单代数。则
$$
\begin{matrix}
A \otimes_K A^{opp} & \longrightarrow & \text{End}_K(A) \\
\ a \otimes a' & \longmapsto & (x \mapsto a x a')
\end{matrix}
$$
是 $K$ 上代数的同构。
\end{lemma}

\begin{proof}
见 Brauer Groups，引理 \ref{brauer-lemma-inverse}。
\end{proof}

\begin{definition}
\label{etale-cohomology-definition-brauer-equivalent}
若两个有限维中心单代数 $A_1$ 与 $A_2$（均定义在 $K$ 上）满足：
存在 $m, n \geq 1$，使得
$\text{Mat}(n \times n, A_1)
\cong \text{Mat}(m \times m, A_2)$，
则称它们{\it 相似}或{\it 等价}。记作 $A_1 \sim A_2$。
\end{definition}

\noindent
由 Brauer Groups，引理 \ref{brauer-lemma-similar}，这是等价关系。

\begin{definition}
\label{etale-cohomology-definition-brauer-group}
设 $K$ 是域。$K$ 的{\it 布劳尔群}是相似类集合
$\text{Br} (K)$，其中各类来自 $K$ 上的有限维中心单代数，
其群律由张量积（在 $K$ 上）诱导。
$A$ 在 $\text{Br}(K)$ 中的类记作 $[A]$。单位元为
$[K] = [\text{Mat}(d \times d, K)]$，其中任意 $d \geq 1$。
\end{definition}

\noindent
前一引理说明逆元存在，且 $-[A] = [A^{opp}]$。
域的布劳尔群总是挠群。
事实上，将看到 $[A]$ 的阶整除 $\deg(A)$，其中 $A$
是任意有限维中心单代数（见引理 \ref{etale-cohomology-lemma-annihilated-by-degree}）。
一般而言，布劳尔群不是有限生成的；例如，非 Archimedes 局部域的布劳尔群是
$\mathbf{Q}/\mathbf{Z}$。
$\mathbf{C}(x, y)$ 的布劳尔群不可数。

\begin{lemma}
\label{etale-cohomology-lemma-central-simple-algebra-pgln}
\begin{slogan}
中心单代数由 PGL 的 Galois 上同调分类。
\end{slogan}
设 $K$ 是域，并设 $K^{sep}$ 是可分代数闭包。
则次数为 $d$、定义在 $K$ 上的中心单代数之同构类集合，与非 Abel 上同调
$H^1(\text{Gal}(K^{sep}/K), \text{PGL}_d(K^{sep}))$
之间存在双射。
\end{lemma}

\begin{proof}[证明概要。]
Skolem--Noether 定理（见 Brauer Groups，定理
\ref{brauer-theorem-skolem-noether}）说明，对任意域 $L$，群
$\text{Aut}_{L\text{-Algebras}}(\text{Mat}_d(L))$
等于 $\text{PGL}_d(L)$。由定理 \ref{etale-cohomology-theorem-central-simple-algebra}，
可见次数为 $d$ 的中心单代数对应于 $K$-代数
$\text{Mat}_d(K)$ 的形式。
结合二者可见，次数为 $d$ 的中心单代数之同构类对应于
$H^1(\text{Gal}(K^{sep}/K), \text{PGL}_d(K^{sep}))$
的元素。关于扭曲的更多细节，例如见 \cite{SilvermanEllipticCurves}。
\end{proof}

\noindent
若 $A$ 是次数为 $d$ 的有限维中心单代数，定义在域 $K$ 上，
以 $\xi_A$ 表示相应的上同调类
$H^1(\text{Gal}(K^{sep}/K), \text{PGL}_d(K^{sep}))$。
考虑短正合序列
$$
1 \to (K^{sep})^* \to \text{GL}_d(K^{sep}) \to \text{PGL}_d(K^{sep}) \to 1,
$$
它给出上同调长正合序列（直到次数 2），其连接映射为
$$
\delta_d :
H ^1(\text{Gal}(K^{sep}/K), \text{PGL}_d(K^{sep}))
\longrightarrow
H^2(\text{Gal}(K^{sep}/K), (K^{sep})^*).
$$
具体如下：若 $\xi$ 是由 1-上闭链 $(g_\sigma)$ 表示的上同调类，
则 $\delta_d(\xi)$ 是 2-上闭链之类
\begin{equation}
\label{etale-cohomology-equation-two-cocycle}
(\sigma, \tau)
\longmapsto
\tilde g_\sigma^{-1} \tilde g_{\sigma \tau} \sigma(\tilde g_\tau^{-1})
\in (K^{sep})^*
\end{equation}
，其中 $\tilde g_\sigma \in \text{GL}_d(K^{sep})$ 是
$g_\sigma$ 的一个提升。
利用这一点，可以将映射
$$
\delta : \text{Br}(K) \longrightarrow H^2(\text{Gal}(K^{sep}/K), (K^{sep})^*),
\quad
[A] \longmapsto \delta_{\deg A} (\xi_A)
$$
具体化如下。假设 $A$ 次数为 $d$，且定义在 $K$ 上。选取同构
$\varphi : \text{Mat}_d(K^{sep}) \to A \otimes_K K^{sep}$。
对 $\sigma \in \text{Gal}(K^{sep}/K)$，选取元素
$\tilde g_\sigma \in \text{GL}_d(K^{sep})$，使得
$\varphi^{-1} \circ \sigma(\varphi)$ 等于映射
$x \mapsto \tilde g_\sigma x \tilde g_\sigma^{-1}$。
$H^2$ 中的类由二上闭链（\ref{etale-cohomology-equation-two-cocycle}）定义。

\begin{theorem}
\label{etale-cohomology-theorem-brauer-delta}
设 $K$ 是具有可分代数闭包 $K^{sep}$ 的域。上面定义的映射
$\delta : \text{Br}(K) \to H^2(\text{Gal}(K^{sep}/K), (K^{sep})^*)$
是群同构。
\end{theorem}

\begin{proof}[证明概要]
为证明 $\delta$ 定义群同态，即
$\delta(A \otimes_K B) = \delta(A) + \delta(B)$，
直接用上闭链计算。

\medskip\noindent
证明 $\delta$ 的单射性。在 Abel 情形（$d = 1$）中，有等同
$$
H^1(\text{Gal}(K^{sep}/K), \text{GL}_d(K^{sep})) =
H_\etale^1(\Spec(K), \text{GL}_d(\mathcal{O}))
$$
后者由 fpqc 下降是平凡的。若这在非 Abel 情形也成立，
便立即推出 $\delta$ 的单射性。（见 \cite{SGA4.5}。）
更确切地说，为证明这一点，可以把 $\delta([A])$ 重新解释为
如下存在性的障碍：存在 $K$-向量空间 $V$，它带左 $A$-模结构，
并满足 $\dim_K V = \deg A$。当 $V$ 存在时，有
$A \cong \text{End}_K(V)$。

\medskip\noindent
为证明满射性，选取上同调类
$\xi \in H^2(\text{Gal}(K^{sep}/K), (K^{sep})^*)$，
则存在有限 Galois 扩张 $K^{sep}/K'/K$，使得 $\xi$ 是某个
$\xi' \in H^2(\text{Gal}(K'|K), (K')^*)$ 的像。
然后写出 $K$ 上的一个显式中心单代数，使用数据 $K', \xi'$。
\end{proof}










\section{概形的布劳尔群}
\label{etale-cohomology-section-brauer-scheme}

\noindent
设 $S$ 为概形。若一个 $\mathcal{O}_S$-代数
$\mathcal{A}$ 在 \'etale 局部是矩阵代数，则称其为 {\it Azumaya 代数}；
换言之，若存在一个 \'etale 覆盖
$\mathcal{U} = \{ \varphi_i : U_i \to S\}_{i \in I}$，使得
$\varphi_i^*\mathcal{A} \cong \text{Mat}_{d_i}(\mathcal{O}_{U_i})$
对某些 $d_i \geq 1$ 成立。对于两个这样的
$\mathcal{A}$ 与 $\mathcal{B}$，若存在有限局部自由
$\mathcal{O}_S$-模 $\mathcal{F}$ 与 $\mathcal{G}$，
它们在每个 $s \in S$ 处的秩都为正，并且
$$
\mathcal{A} \otimes_{\mathcal{O}_S}
\SheafHom_{\mathcal{O}_S}(\mathcal{F}, \mathcal{F})
\cong
\mathcal{B} \otimes_{\mathcal{O}_S}
\SheafHom_{\mathcal{O}_S}(\mathcal{G}, \mathcal{G})
$$
作为 $\mathcal{O}_S$-代数同构，则称它们{\it 等价}。$S$ 的{\it 布劳尔群}
是集合 $\text{Br}(S)$，其元素为 Azumaya
$\mathcal{O}_S$-代数的等价类，
其运算由（在 $\mathcal{O}_S$ 上的）张量积诱导。

\begin{lemma}
\label{etale-cohomology-lemma-end-unique-up-to-invertible}
设 $S$ 为概形。设 $\mathcal{F}$ 与 $\mathcal{G}$ 为正秩的有限局部自由
$\mathcal{O}_S$-模层。若存在同构
$\SheafHom_{\mathcal{O}_S}(\mathcal{F}, \mathcal{F}) \cong
\SheafHom_{\mathcal{O}_S}(\mathcal{G}, \mathcal{G})$，且它是
$\mathcal{O}_S$-代数同构，则存在可逆层
$\mathcal{L}$，它定义在 $S$ 上，并且
$\mathcal{F} \otimes_{\mathcal{O}_S} \mathcal{L} \cong \mathcal{G}$，
并且此同构诱导给定的自同态代数同构。
\end{lemma}

\begin{proof}
固定一个同构
$\SheafHom_{\mathcal{O}_S}(\mathcal{F}, \mathcal{F}) \to
\SheafHom_{\mathcal{O}_S}(\mathcal{G}, \mathcal{G})$。
考虑层 $\mathcal{L} \subset \SheafHom(\mathcal{F}, \mathcal{G})$：
作为 $\mathcal{O}_S$-模，它由所有局部同构
$\varphi : \mathcal{F} \to \mathcal{G}$ 生成，其中由
$\varphi$ 作共轭得到的正是给定的自同态代数同构。
局部计算（归结到 $\mathcal{F}$ 与 $\mathcal{G}$ 有限自由且
$S$ 仿射的情形）表明 $\mathcal{L}$ 可逆。另一项局部计算表明取值映射
$$
\mathcal{F} \otimes_{\mathcal{O}_S} \mathcal{L} \longrightarrow \mathcal{G}
$$
是同构。
\end{proof}

\noindent
下述引理证明中的论证可见
\cite{Saltman-torsion}。

\begin{lemma}
\label{etale-cohomology-lemma-annihilated-by-degree}
\begin{reference}
论证取自 \cite{Saltman-torsion}。
\end{reference}
设 $S$ 为概形。设 $\mathcal{A}$ 为 Azumaya 代数，并且是秩为
$d^2$ 的局部自由代数（相对于 $S$）。则
$\mathcal{A}$ 在 $S$ 的布劳尔群中的类被 $d$ 零化。
\end{lemma}

\begin{proof}
选取 \'etale 覆盖 $\{U_i \to S\}$，并选取同构
$\mathcal{A}|_{U_i} \to \SheafHom(\mathcal{F}_i, \mathcal{F}_i)$，
其中某个局部自由 $\mathcal{O}_{U_i}$-模 $\mathcal{F}_i$
秩为 $d$。（可以假设 $\mathcal{F}_i$ 自由。）考虑复合
$$
p_i : \mathcal{F}_i^{\otimes d} \to
\wedge^d(\mathcal{F}_i) \to \mathcal{F}_i^{\otimes d}
$$
第一个箭头是通常的投影；第二个箭头把 $\mathcal{F}_i$ 的最高外幂
同构到 $\mathcal{F}_i^{\otimes d}$ 的如下截面子模：其中的截面在
$d$ 个字母上的对称群作用下依符号特征变换。于是
$p_i^2 = d! p_i$，并且 $p_i$ 的秩为 $1$。利用给定同构
$\mathcal{A}|_{U_i} \to \SheafHom(\mathcal{F}_i, \mathcal{F}_i)$
以及典范同构
$$
\SheafHom(\mathcal{F}_i, \mathcal{F}_i)^{\otimes d} =
\SheafHom(\mathcal{F}_i^{\otimes d}, \mathcal{F}_i^{\otimes d})
$$
，可把 $p_i$ 看作 $\mathcal{A}^{\otimes d}$ 在
$U_i$ 上的截面。我们断言，有
$p_i|_{U_i \times_S U_j} = p_j|_{U_i \times_S U_j}$，
且这是作为 $\mathcal{A}^{\otimes d}$ 的截面的等式。事实上，应用
引理 \ref{etale-cohomology-lemma-end-unique-up-to-invertible}，得到可逆层
$\mathcal{L}_{ij}$ 以及典范同构
$$
\mathcal{F}_i|_{U_i \times_S U_j} \otimes \mathcal{L}_{ij}
\longrightarrow
\mathcal{F}_j|_{U_i \times_S U_j}.
$$
由此同构可见 $p_i$ 映到 $p_j$。由于 $\mathcal{A}^{\otimes d}$
是 $S_\etale$ 上的层（命题
\ref{etale-cohomology-proposition-quasi-coherent-sheaf-fpqc}），我们得到典范整体截面
$p \in \Gamma(S, \mathcal{A}^{\otimes d})$。局部计算表明
$$
\mathcal{H} =
\Im(\mathcal{A}^{\otimes d} \to \mathcal{A}^{\otimes d}, f \mapsto fp)
$$
是秩为 $d^d$ 的局部自由模，并且由 $\mathcal{A}^{\otimes d}$ 作（左）乘
诱导同构
$\mathcal{A}^{\otimes d} \to \SheafHom(\mathcal{H}, \mathcal{H})$。
换言之，$\mathcal{A}^{\otimes d}$ 是 $S$ 的布劳尔群中的平凡元素，
正如所需。
\end{proof}

\noindent
在此情形下，定理 \ref{etale-cohomology-theorem-brauer-delta} 中同构 $\delta$ 的类似物
是映射
$$
\delta_S: \text{Br}(S) \to H_\etale^2(S, \mathbf{G}_m).
$$
映射 $\delta_S$ 确实是单射。若 $S$ 拟紧或连通，则
$\text{Br}(S)$ 是挠群，因而在此情形下 $\delta_S$ 的像包含在 $S$
的{\it 上同调布劳尔群}
$$
\text{Br}'(S) := H_\etale^2(S, \mathbf{G}_m)_\text{torsion}.
$$
中。因此，若 $S$ 拟紧或连通，就有包含关系 $\text{Br}(S)
\subset \text{Br}'(S)$。它并非总是等式：存在非分离奇异曲面 $S$，
使得 $\text{Br}(S) \subset
\text{Br}'(S)$ 是严格包含。若 $S$ 拟射影，则
$\text{Br}(S) = \text{Br}'(S)$。不过，举例来说，
尚不知这一点对 $\mathbf{C}$ 上的光滑固有簇是否成立。



\section{Artin--Schreier 序列}
\label{etale-cohomology-section-artin-schreier}

\noindent
设 $p$ 为素数。设 $S$ 为特征 $p$ 的概形。
{\it Artin--Schreier} 序列是短正合列
$$
0 \longrightarrow \underline{\mathbf{Z}/p\mathbf{Z}}_S \longrightarrow
\mathbf{G}_{a, S} \xrightarrow{F-1} \mathbf{G}_{a, S} \longrightarrow 0
$$
，其中 $F - 1$ 是映射 $x \mapsto x^p - x$。

\begin{lemma}
\label{etale-cohomology-lemma-vanishing-affine-char-p-p}
设 $p$ 为素数。设 $S$ 为特征 $p$ 的概形。
\begin{enumerate}
\item 若 $S$ 仿射，则
$H_\etale^q(S, \underline{\mathbf{Z}/p\mathbf{Z}}) = 0$ 对所有
$q \geq 2$ 都成立。
\item 若 $S$ 是维数为 $d$ 的拟紧拟分离概形，则
$H_\etale^q(S, \underline{\mathbf{Z}/p\mathbf{Z}}) = 0$
对所有 $q \geq 2 + d$ 都成立。
\end{enumerate}
\end{lemma}

\begin{proof}
回忆结构层的 \'etale 上同调等于它在底拓扑空间上的上同调
（定理 \ref{etale-cohomology-theorem-zariski-fpqc-quasi-coherent}）。
第一条由 Artin--Schreier 正合列以及仿射概形上结构层上同调的消失得到
（《概形上同调》引理
\ref{coherent-lemma-quasi-coherent-affine-cohomology-zero}）。
第二条由同一论证、《上同调》命题
\ref{cohomology-proposition-cohomological-dimension-spectral}
中的消失，以及 $S$ 是谱空间这一事实得到
（《性质》引理
\ref{properties-lemma-quasi-compact-quasi-separated-spectral}）。
\end{proof}

\begin{lemma}
\label{etale-cohomology-lemma-F-1}
设 $k$ 是特征 $p > 0$ 的代数闭域。
设 $V$ 是有限维 $k$-向量空间。设 $F : V \to V$
为 Frobenius 线性映射，即加法映射，并且满足
$F(\lambda v) = \lambda^p F(v)$，其中任取 $\lambda \in k$
与 $v \in V$。则 $F - 1 : V \to V$ 是满射，其核是有限维
$\mathbf{F}_p$-向量空间，维数 $\leq \dim_k(V)$。
\end{lemma}

\begin{proof}
若 $F = 0$，则结论成立。若 $V$ 有一个由
$F$-稳定子向量空间组成的滤过，并且结论对每个分次片都成立，
则它对 $(V, F)$ 成立。结合这两点，可以假设 $F$ 的核为零。

\medskip\noindent
选取基 $v_1, \ldots, v_n$，它们构成 $V$ 的一组基，并写成
$F(v_i) = \sum a_{ij} v_j$。注意，$v = \sum \lambda_i v_i$
属于该核，当且仅当 $\sum \lambda_i^p a_{ij} v_j = 0$。
由于 $k$ 代数闭，这蕴含矩阵 $(a_{ij})$ 可逆。设
$(b_{ij})$ 为其逆矩阵。为了说明 $F - 1$ 为满射，选取
$w = \sum \mu_i v_i \in V$，并尝试求解
$$
(F - 1)(\sum \lambda_iv_i) =
\sum \lambda_i^p a_{ij} v_j - \sum \lambda_j v_j = \sum \mu_j v_j
$$
这等价于
$$
\sum \lambda_j^p v_j - \sum b_{ij} \lambda_i v_j = \sum b_{ij} \mu_i v_j
$$
，换言之，
$$
\lambda_j^p - \sum b_{ij} \lambda_i = \sum b_{ij} \mu_i,
\quad j = 1, \ldots, \dim(V).
$$
代数
$$
A = k[x_1, \ldots, x_n]/
(x_j^p - \sum b_{ij} x_i - \sum b_{ij} \mu_i)
$$
在 $k$ 上标准光滑（《代数》定义
\ref{algebra-definition-standard-smooth}），因为矩阵 $(b_{ij})$
可逆，而 $x_j^p$ 的偏导数都为零。$A$ 在 $k$ 上的一组基是所有单项式
$x_1^{e_1} \ldots x_n^{e_n}$，其中 $e_i < p$，故
$\dim_k(A) = p^n$。由于 $k$ 代数闭，可见 $\Spec(A)$ 恰有
$p^n$ 个点。因此 $F - 1$ 为满射，并且每条纤维都有
$p^n$ 个点；也就是说，$F - 1$ 的核是有 $p^n$ 个元素的群。
\end{proof}

\begin{lemma}
\label{etale-cohomology-lemma-F-2}
在引理 \ref{etale-cohomology-lemma-F-1} 的情形下，设 $k'/k$ 是代数闭域扩张。
令 $V' = V \otimes_k k'$，并令
$F' : V' \to V'$ 为 $F$ 诱导的 Frobenius 线性映射。
则 $\Ker(F - 1)$ 同构地映到 $\Ker(F' - 1)$ 上。
\end{lemma}

\begin{proof}
设 $v_1, \ldots, v_n$ 是 $V$ 的一组基，并设
$F(v_i) = \sum a_{ij}v_j$。则 $v = \sum \lambda_i v_i$
属于 $\Ker(F - 1)$，当且仅当有
$\sum \lambda_i^pa_{ij} = \lambda_j$，其中 $j = 1, \ldots, n$。
由引理 \ref{etale-cohomology-lemma-F-1}，
这些方程截出一个闭子概形 $Z$；它位于
$\mathbf{A}^n_k = \Spec(k[\lambda_1, \ldots, \lambda_n])$ 中，
并且在 $\Spec(k)$ 上有限。
于是 $Z(k) = Z(k')$，结论随之成立。
\end{proof}

\begin{lemma}
\label{etale-cohomology-lemma-top-cohomology-coherent}
设 $X$ 是域 $k$ 上分离的有限型概形。设 $\mathcal{F}$
为 $\mathcal{O}_X$-模的凝聚层。则
$\dim_k H^d(X, \mathcal{F}) < \infty$，其中 $d = \dim(X)$。
\end{lemma}

\begin{proof}
对 $d$ 作归纳。情形 $d = 0$ 成立，因为此时
$X$ 是有限维 $k$-代数 $A$ 的谱
（《簇》引理 \ref{varieties-lemma-algebraic-scheme-dim-0}），
而每个凝聚层 $\mathcal{F}$ 都对应于有限 $A$-模
$M = H^0(X, \mathcal{F})$，它满足 $\dim_k M < \infty$。

\medskip\noindent
假设 $d > 0$，且已对维数 $< d$ 的分离有限型概形证明结论。
概形 $X$ 是诺特的。考虑性质
$\mathcal{P}$；它定义在 $X$ 上的凝聚层上，规则为
$$
\mathcal{P}(\mathcal{F}) \Leftrightarrow
\dim_k H^d(X, \mathcal{F}) < \infty
$$
我们将使用《概形上同调》引理
\ref{coherent-lemma-property-initial} 的结果，证明
$\mathcal{P}$ 对 $X$ 上每个凝聚层都成立。

\medskip\noindent
设
$$
0 \to \mathcal{F}_1 \to \mathcal{F} \to \mathcal{F}_2 \to 0
$$
是 $X$ 上凝聚层的短正合列。考虑上同调长正合列中的一段
$$
H^d(X, \mathcal{F}_1) \to
H^d(X, \mathcal{F}) \to
H^d(X, \mathcal{F}_2)
$$
因此，若 $\mathcal{P}$ 对 $\mathcal{F}_1$ 与 $\mathcal{F}_2$
成立，则它对 $\mathcal{F}$ 也成立。

\medskip\noindent
设 $Z \subset X$ 为整闭子概形。设 $\mathcal{I}$
为 $Z$ 上的凝聚理想层。为完成证明，须说明
$H^d(X, i_*\mathcal{I}) = H^d(Z, \mathcal{I})$ 是有限维的。
若 $\dim(Z) < d$，则结论成立，因为该上同调群为零
（《上同调》命题
\ref{cohomology-proposition-vanishing-Noetherian}）。
这样便归结到下一段讨论的情形。

\medskip\noindent
假设 $X$ 是维数为 $d$ 的簇，并且
$\mathcal{F} = \mathcal{I}$ 是凝聚理想层。在此情形下有短正合列
$$
0 \to \mathcal{I} \to \mathcal{O}_X \to i_*\mathcal{O}_Z \to 0
$$
，其中 $i : Z \to X$ 是由 $\mathcal{I}$ 定义的闭子概形。
由归纳假设可见
$H^{d - 1}(Z, \mathcal{O}_Z) = H^{d - 1}(X, i_*\mathcal{O}_Z)$
是有限维的。因此，只需对结构层证明结论。

\medskip\noindent
可以把 Chow 引理
（《概形上同调》引理 \ref{coherent-lemma-chow-Noetherian}）
应用于态射 $X \to \Spec(k)$。由此得到 Chow 引理陈述中的图
$$
\xymatrix{
X \ar[rd]_g & X' \ar[d]^-{g'} \ar[l]^\pi \ar[r]_i & \mathbf{P}^n_k \ar[dl] \\
& \Spec(k) &
}
$$
再设 $U \subset X$ 是满足 $\pi^{-1}(U) \to U$ 为同构的稠密开子概形。
还可以假设 $X'$ 也是簇，见
《概形上同调》评注 \ref{coherent-remark-chow-Noetherian}。
态射 $i' = (i, \pi) : X' \to \mathbf{P}^n_X$
是闭浸入（同上）。因此
$$
\mathcal{L} = i^*\mathcal{O}_{\mathbf{P}^n_k}(1) \cong
(i')^*\mathcal{O}_{\mathbf{P}^n_X}(1)
$$
是 $\pi$-相对丰沛的（例如由《态射》引理
\ref{morphisms-lemma-characterize-ample-on-finite-type}）。
所以由《概形上同调》引理
\ref{coherent-lemma-kill-by-twisting}，存在 $n \geq 0$，使得
$R^p\pi_*\mathcal{L}^{\otimes n} = 0$ 对所有 $p > 0$ 都成立。
置 $\mathcal{G} = \pi_*\mathcal{L}^{\otimes n}$。
任取非零整体截面 $s$；它是 $\mathcal{L}^{\otimes n}$ 的截面。
由于 $\mathcal{G} = \pi_*\mathcal{L}^{\otimes n}$，截面 $s$
对应于 $\mathcal{G}$ 的一个截面，即一个映射
$\mathcal{O}_X \to \mathcal{G}$。
有 $s|_U \not = 0$，因为 $X'$ 是簇且 $\mathcal{L}$ 可逆，
所以 $\mathcal{O}_X|_U \to \mathcal{G}|_U$ 非零。
由于 $\mathcal{G}|_U = \mathcal{L}^{\otimes n}|_{\pi^{-1}(U)}$
可逆，得到短正合列
$$
0 \to \mathcal{O}_X \to \mathcal{G} \to \mathcal{Q} \to 0
$$
，其中 $\mathcal{Q}$ 凝聚且支撑在 $X$ 的真闭子概形上。
像先前一样使用归纳假设，可见只需证明
$\dim H^d(X, \mathcal{G}) < \infty$。

\medskip\noindent
由 Leray 谱序列
（《上同调》引理 \ref{cohomology-lemma-apply-Leray}）可见
$H^d(X, \mathcal{G}) = H^d(X', \mathcal{L}^{\otimes n})$。
设 $\overline{X}' \subset \mathbf{P}^n_k$ 为 $X'$ 的闭包。
则 $\overline{X}'$ 是维数为 $d$、定义在 $k$ 上的射影簇，而
$X' \subset \overline{X}'$ 是稠密开子概形。
可逆层 $\mathcal{L}^{\otimes n}$ 是
$\mathcal{O}_{\overline{X}'}(n)$ 在 $X'$ 上的限制。由
《上同调》命题
\ref{cohomology-proposition-cohomological-dimension-spectral}，映射
$$
H^d(\overline{X}', \mathcal{O}_{\overline{X}'}(n))
\longrightarrow
H^d(X', \mathcal{L}^{\otimes n})
$$
是满射。左侧上同调群由《概形上同调》引理
\ref{coherent-lemma-coherent-projective} 是有限维的，证明完成。
\end{proof}

\begin{lemma}
\label{etale-cohomology-lemma-vanishing-variety-char-p-p}
设 $X$ 是代数闭域 $k$ 上分离的有限型概形，且该域的特征为
$p > 0$。则
$H_\etale^q(X, \underline{\mathbf{Z}/p\mathbf{Z}}) = 0$
对 $q \geq dim(X) + 1$ 成立。
\end{lemma}

\begin{proof}
设 $d = \dim(X)$。由引理
\ref{etale-cohomology-lemma-vanishing-affine-char-p-p} 中建立的消失，只需说明
$H_\etale^{d + 1}(X, \underline{\mathbf{Z}/p\mathbf{Z}}) = 0$。
由引理 \ref{etale-cohomology-lemma-top-cohomology-coherent} 可见
$H^d(X, \mathcal{O}_X)$ 是有限维 $k$-向量空间。
因此，与 Artin--Schreier 序列相伴的上同调长正合列以
$$
H^d(X, \mathcal{O}_X) \xrightarrow{F - 1}
H^d(X, \mathcal{O}_X) \to H^{d + 1}_\etale(X, \mathbf{Z}/p\mathbf{Z}) \to 0
$$
结尾。由引理 \ref{etale-cohomology-lemma-F-1}，此序列中的映射 $F - 1$ 是满射。
这就证明了该引理。
\end{proof}

\begin{lemma}
\label{etale-cohomology-lemma-finiteness-proper-variety-char-p-p}
设 $X$ 是代数闭域 $k$ 上的固有概形，且该域的特征为 $p > 0$。则
\begin{enumerate}
\item $H_\etale^q(X, \underline{\mathbf{Z}/p\mathbf{Z}})$
是有限 $\mathbf{Z}/p\mathbf{Z}$-模，对所有 $q$ 都如此；并且
\item 映射
$H^q_\etale(X, \underline{\mathbf{Z}/p\mathbf{Z}}) \to
H^q_\etale(X_{k'}, \underline{\mathbf{Z}/p\mathbf{Z}})$
是同构，只要 $k'/k$ 是代数闭域扩张。
\end{enumerate}
\end{lemma}

\begin{proof}
由《概形上同调》引理
\ref{coherent-lemma-proper-over-affine-cohomology-finite}
以及定理 \ref{etale-cohomology-theorem-zariski-fpqc-quasi-coherent} 中的上同调比较，
上同调群
$H^q_\etale(X, \mathbf{G}_a) = H^q(X, \mathcal{O}_X)$
是有限维 $k$-向量空间。因此，由引理 \ref{etale-cohomology-lemma-F-1}，
与 Artin--Schreier 序列相伴的上同调长正合列分解为短正合列
$$
0 \to H_\etale^q(X, \underline{\mathbf{Z}/p\mathbf{Z}}) \to
H^q(X, \mathcal{O}_X) \xrightarrow{F - 1} H^q(X, \mathcal{O}_X) \to 0
$$
而且 $\mathbf{F}_p$-维数（对上同调群
$H_\etale^q(X, \underline{\mathbf{Z}/p\mathbf{Z}})$ 而言）
小于或等于 $k$-维数（对向量空间
$H^q(X, \mathcal{O}_X)$ 而言）。
这证明了第一条。第二条由引理 \ref{etale-cohomology-lemma-F-2} 得到，因为
$H^q(X, \mathcal{O}_X) \otimes_k k' \to H^q(X_{k'}, \mathcal{O}_{X_{k'}})$
由平坦基变换是同构
（《概形上同调》引理
\ref{coherent-lemma-flat-base-change-cohomology}）。
\end{proof}








\section{局部常值层}
\label{etale-cohomology-section-locally-constant}

\noindent
本节是《网站上的模》第
\ref{sites-modules-section-locally-constant} 节针对 \'etale 网站的类似版本。

\begin{definition}
\label{etale-cohomology-definition-finite-locally-constant}
设 $X$ 为概形。
设 $\mathcal{F}$ 为 $X_\etale$ 上的集合层。
\begin{enumerate}
\item 设 $E$ 为集合。称 $\mathcal{F}$ 为{\it 取值为 $E$ 的常值层}，
如果 $\mathcal{F}$ 是预层 $U \mapsto E$ 的层化。
记作 $\underline{E}_X$ 或 $\underline{E}$。
\item 若 $\mathcal{F}$ 同构于第 (1) 项中的层，则称其为{\it 常值层}。
\item 称 $\mathcal{F}$ {\it 局部常值}，如果存在覆盖
$\{U_i \to X\}$，使得 $\mathcal{F}|_{U_i}$ 为常值层。
\item 若 $\mathcal{F}$ 局部常值，而且各常值取值都是有限集合，
则称其{\it 有限局部常值}。
\end{enumerate}
设 $\mathcal{F}$ 为 $X_\etale$ 上的阿贝尔群层。
\begin{enumerate}
\item 设 $A$ 为阿贝尔群。
称 $\mathcal{F}$ 为{\it 取值为 $A$ 的常值层}，
如果 $\mathcal{F}$ 是预层 $U \mapsto A$ 的层化。
记作 $\underline{A}_X$ 或 $\underline{A}$。
\item 若 $\mathcal{F}$ 作为阿贝尔群层同构于第 (1) 项中的层，
则称其为{\it 常值层}。
\item 称 $\mathcal{F}$ {\it 局部常值}，如果存在覆盖
$\{U_i \to X\}$，使得 $\mathcal{F}|_{U_i}$ 为常值层。
\item 若 $\mathcal{F}$ 局部常值，而且各常值取值都是有限阿贝尔群，
则称其{\it 有限局部常值}。
\end{enumerate}
设 $\Lambda$ 为环。设 $\mathcal{F}$ 为 $\Lambda$-模层，
它定义在 $X_\etale$ 上。
\begin{enumerate}
\item 设 $M$ 为 $\Lambda$-模。
称 $\mathcal{F}$ 为{\it 取值为 $M$ 的常值层}，
如果 $\mathcal{F}$ 是预层 $U \mapsto M$ 的层化。
记作 $\underline{M}_X$ 或 $\underline{M}$。
\item 若 $\mathcal{F}$ 作为 $\Lambda$-模层同构于第 (1) 项中的层，
则称其为{\it 常值层}。
\item 称 $\mathcal{F}$ {\it 局部常值}，如果存在覆盖
$\{U_i \to X\}$，使得 $\mathcal{F}|_{U_i}$ 为常值层。
\end{enumerate}
\end{definition}

\begin{lemma}
\label{etale-cohomology-lemma-pullback-locally-constant}
设 $f : X \to Y$ 为概形态射。若 $\mathcal{G}$ 是局部常值的
集合层、阿贝尔群层或 $\Lambda$-模层，并且定义在 $Y_\etale$ 上，
则 $f^{-1}\mathcal{G}$ 在 $X_\etale$ 上也具有同样的性质。
\end{lemma}

\begin{proof}
这对任意拓扑斯态射都成立，见《网站上的模》引理
\ref{sites-modules-lemma-pullback-locally-constant}。
\end{proof}

\begin{lemma}
\label{etale-cohomology-lemma-pushforward-locally-constant}
设 $f : X \to Y$ 为有限 \'etale 概形态射。
若 $\mathcal{F}$ 是（有限）局部常值的集合层、
（有限）局部常值的阿贝尔群层或（有限型）局部常值的
$\Lambda$-模层，并且定义在 $X_\etale$ 上，则
$f_*\mathcal{F}$ 在 $Y_\etale$ 上也具有同样的性质。
\end{lemma}

\begin{proof}
$f_*$ 的构造与 \'etale 局部化可交换。
有限 \'etale 态射局部同构于若干同构的不交并，见
《\'Etale 态射》引理 \ref{etale-lemma-finite-etale-etale-local}。
因此，该引理是说：若 $\mathcal{F}_i$（$i = 1, \ldots, n$）
都是（有限）局部常值的集合层，则
$\prod_{i = 1, \ldots, n} \mathcal{F}_i$ 也是。这是显然的。
阿贝尔群层和模层的情形类似。
\end{proof}

\begin{lemma}
\label{etale-cohomology-lemma-characterize-finite-locally-constant}
设 $X$ 为概形，$\mathcal{F}$ 为 $X_\etale$ 上的集合层。
则下列条件等价：
\begin{enumerate}
\item $\mathcal{F}$ 有限局部常值；并且
\item 有 $\mathcal{F} = h_U$，其中 $U \to X$ 是某个有限 \'etale 态射。
\end{enumerate}
\end{lemma}

\begin{proof}
有限 \'etale 态射局部同构于若干同构的不交并，见
《\'Etale 态射》引理 \ref{etale-lemma-finite-etale-etale-local}。
因此 (2) 蕴含 (1)。反过来，若 $\mathcal{F}$ 有限局部常值，
则存在 \'etale 覆盖 $\{X_i \to X\}$，使得
$\mathcal{F}|_{X_i}$ 可由有限 \'etale 态射 $U_i \to X_i$ 表示。
完全仿照《下降》引理
\ref{descent-lemma-descent-data-sheaves} 的证明进行论证，
可得到概形下降数据 $(U_i, \varphi_{ij})$，它相对于
$\{X_i \to X\}$（细节从略）。例如由《下降》引理
\ref{descent-lemma-affine}，此下降数据是有效的；所得概形态射
$U \to X$ 由《下降》引理
\ref{descent-lemma-descending-property-finite} 和
\ref{descent-lemma-descending-property-etale} 是有限 \'etale 的。
\end{proof}

\begin{lemma}
\label{etale-cohomology-lemma-morphism-locally-constant}
设 $X$ 为概形。
\begin{enumerate}
\item 设 $\varphi : \mathcal{F} \to \mathcal{G}$ 是
$X_\etale$ 上局部常值集合层的映射。
若 $\mathcal{F}$ 有限局部常值，则存在 \'etale 覆盖
$\{U_i \to X\}$，使得 $\varphi|_{U_i}$
是由集合映射所伴随的常值层映射。
\item 设 $\varphi : \mathcal{F} \to \mathcal{G}$ 是
$X_\etale$ 上局部常值阿贝尔群层的映射。
若 $\mathcal{F}$ 有限局部常值，则存在 \'etale 覆盖
$\{U_i \to X\}$，使得 $\varphi|_{U_i}$
是由阿贝尔群映射所伴随的常值阿贝尔群层映射。
\item 设 $\Lambda$ 为环。
设 $\varphi : \mathcal{F} \to \mathcal{G}$ 是
局部常值 $\Lambda$-模层的映射，并且这些层定义在 $X_\etale$ 上。
若 $\mathcal{F}$ 有限型，则存在 \'etale 覆盖
$\{U_i \to X\}$，使得 $\varphi|_{U_i}$
是由 $\Lambda$-模映射所伴随的常值 $\Lambda$-模层映射。
\end{enumerate}
\end{lemma}

\begin{proof}
这对任意网站都成立，见《网站上的模》引理
\ref{sites-modules-lemma-morphism-locally-constant}。
\end{proof}

\begin{lemma}
\label{etale-cohomology-lemma-kernel-finite-locally-constant}
设 $X$ 为概形。
\begin{enumerate}
\item 有限局部常值集合层的范畴在 $\Sh(X_\etale)$ 内
对有限极限与有限余极限封闭。
\item 有限局部常值阿贝尔群层的范畴是
$\textit{Ab}(X_\etale)$ 的弱 Serre 子范畴。
\item 设 $\Lambda$ 为诺特环。有限型局部常值
$\Lambda$-模层（定义在 $X_\etale$ 上）的范畴是
$\textit{Mod}(X_\etale, \Lambda)$ 的弱 Serre 子范畴。
\end{enumerate}
\end{lemma}

\begin{proof}
这对任意网站都成立，见《网站上的模》引理
\ref{sites-modules-lemma-kernel-finite-locally-constant}。
\end{proof}

\begin{lemma}
\label{etale-cohomology-lemma-tensor-product-locally-constant}
设 $X$ 为概形。设 $\Lambda$ 为环。
两个局部常值 $\Lambda$-模层（定义在 $X_\etale$ 上）的张量积仍是
局部常值 $\Lambda$-模层。
\end{lemma}

\begin{proof}
这对任意网站都成立，见《网站上的模》引理
\ref{sites-modules-lemma-tensor-product-locally-constant}。
\end{proof}

\begin{lemma}
\label{etale-cohomology-lemma-connected-locally-constant}
设 $X$ 为连通概形。设 $\Lambda$ 为环，并设
$\mathcal{F}$ 为局部常值 $\Lambda$-模层。
则存在 $\Lambda$-模 $M$ 以及 \'etale 覆盖
$\{U_i \to X\}$，使得 $\mathcal{F}|_{U_i} \cong \underline{M}|_{U_i}$。
\end{lemma}

\begin{proof}
选取 \'etale 覆盖
$\{U_i \to X\}$，使得 $\mathcal{F}|_{U_i}$ 为常值层；设
$\mathcal{F}|_{U_i} \cong \underline{M_i}_{U_i}$。
注意，$U_i \times_X U_j$ 为空，只要 $M_i$ 不同构于 $M_j$。
对每个 $\Lambda$-模 $M$，令 $I_M = \{i \in I \mid M_i \cong M\}$。
由于 \'etale 态射是开映射，可见
$U_M = \bigcup_{i \in I_M} \Im(U_i \to X)$
是 $X$ 的开子集。于是 $X = \coprod U_M$ 是 $X$ 的不交开覆盖。
由于 $X$ 连通，只有一个 $U_M$ 非空，引理得证。
\end{proof}







\section{局部常值层与基本群}
\label{etale-cohomology-section-pione}

\noindent
在某些情形下，可以把局部常值层与概形的基本群联系起来。

\begin{lemma}
\label{etale-cohomology-lemma-locally-constant-on-connected}
设 $X$ 为连通概形。设 $\overline{x}$ 为 $X$ 的几何点。
\begin{enumerate}
\item 存在范畴等价
$$
\left\{
\begin{matrix}
\text{有限局部常值}\\
X_\etale\text{ 上的集合层}
\end{matrix}
\right\}
\longleftrightarrow
\left\{
\begin{matrix}
\text{有限 }\pi_1(X, \overline{x})\text{-集合}
\end{matrix}
\right\}
$$
\item 存在范畴等价
$$
\left\{
\begin{matrix}
\text{有限局部常值}\\
X_\etale\text{ 上的阿贝尔群层}
\end{matrix}
\right\}
\longleftrightarrow
\left\{
\begin{matrix}
\text{有限 }\pi_1(X, \overline{x})\text{-模}
\end{matrix}
\right\}
$$
\item 设 $\Lambda$ 为有限环。存在范畴等价
$$
\left\{
\begin{matrix}
\text{有限型局部常值}\\
X_\etale\text{ 上的 }\Lambda\text{-模层}
\end{matrix}
\right\}
\longleftrightarrow
\left\{
\begin{matrix}
\text{带有相容结构的有限 }\pi_1(X, \overline{x})\text{-模，}\\
\text{其 }\Lambda\text{-模结构与作用可交换}
\end{matrix}
\right\}
$$
\end{enumerate}
\end{lemma}

\begin{proof}
注意，$\pi_1(X, \overline{x})$ 是预有限拓扑群，见
《基本群》定义 \ref{pione-definition-fundamental-group}。
左侧各范畴定义于第 \ref{etale-cohomology-section-locally-constant} 节。
右侧范畴所用的记号，对集合取自《基本群》定义
\ref{pione-definition-G-set-continuous}，对阿贝尔群取自定义
\ref{etale-cohomology-definition-G-module-continuous}。这就解释了这些记号。

\medskip\noindent
断言 (1) 由引理 \ref{etale-cohomology-lemma-characterize-finite-locally-constant}
和《基本群》定理 \ref{pione-theorem-fundamental-group} 得到。
第 (2)、(3) 条由此立即得到，只须赋予底层（层的）集合额外结构。
例如，$X_\etale$ 上有限局部常值的阿贝尔群层，
等同于一个有限局部常值集合层 $\mathcal{F}$，连同一个满足通常公理的映射
$+ : \mathcal{F} \times \mathcal{F} \to \mathcal{F}$。
(1) 中的等价把积映到积，因而把 $+$ 映到相应有限
$\pi_1(X, \overline{x})$-集合上的加法。由于
$\pi_1(X, \overline{x})$-模等同于带相容阿贝尔群结构的
$\pi_1(X, \overline{x})$-集合，便得到 (2)。第 (3) 条完全同理。
\end{proof}

\begin{lemma}
\label{etale-cohomology-lemma-locally-constant-on-connected-geom-unibranch}
设 $X$ 为不可约几何单分支概形。
设 $\overline{x}$ 为 $X$ 的几何点。
设 $\Lambda$ 为环。存在范畴等价
$$
\left\{
\begin{matrix}
\text{有限型局部常值}\\
X_\etale\text{ 上的 }\Lambda\text{-模层}
\end{matrix}
\right\}
\longleftrightarrow
\left\{
\begin{matrix}
\text{有限 }\Lambda\text{-模 }M\text{，带有连续的}\\
\pi_1(X, \overline{x})\text{-作用}
\end{matrix}
\right\}
$$
\end{lemma}

\begin{proof}
引理 \ref{etale-cohomology-lemma-locally-constant-on-connected} 中给出的证明不适用，
因为有限 $\Lambda$-模 $M$ 的底层集合未必有限。

\medskip\noindent
设 $\nu : X^\nu \to X$ 为正规化态射。由《态射》引理
\ref{morphisms-lemma-normalization-geom-unibranch-univ-homeo}，
这是万有同胚。由《基本群》命题
\ref{pione-proposition-universal-homeomorphism}，它诱导同构
$\pi_1(X^\nu, \overline{x}) \to \pi_1(X, \overline{x})$；
由定理 \ref{etale-cohomology-theorem-topological-invariance}，可得有限型局部常值
$\Lambda$-模层的两个范畴之间的等价；这两个范畴分别定义在
$X_\etale$ 与 $X^\nu_\etale$ 上。
这把问题归结到 $X$ 为整正规概形的情形。

\medskip\noindent
假设 $X$ 是整正规概形。设 $\eta \in X$ 为泛点。
设 $\overline{\eta}$ 为位于 $\eta$ 上方的几何点。
由《基本群》命题 \ref{pione-proposition-normal}，有连续满射
$$
\text{Gal}(\kappa(\eta)^{sep}/\kappa(\eta)) = \pi_1(\eta, \overline{\eta})
\longrightarrow
\pi_1(X, \overline{\eta})
$$
，其核描述于《基本群》引理
\ref{pione-lemma-normal-pione-quotient-inertia}。
设 $\mathcal{F}$ 为有限型局部常值的
$\Lambda$-模层，它定义在 $X_\etale$ 上。设 $M = \mathcal{F}_{\overline{\eta}}$
为 $\mathcal{F}$ 在 $\overline{\eta}$ 处的茎。
由第 \ref{etale-cohomology-section-galois-action-stalks} 节，得到
$\text{Gal}(\kappa(\eta)^{sep}/\kappa(\eta))$ 在 $M$ 上的连续作用。
目标是说明此作用经由上面显示的满射分解。
由于 $\mathcal{F}$ 有限型，$M$ 是有限 $\Lambda$-模。
由于 $\mathcal{F}$ 局部常值，对每个 $x \in X$，
$\mathcal{F}$ 到 $\Spec(\mathcal{O}_{X, x}^{sh})$ 的限制都是常值的。
因此，$\text{Gal}(K^{sep}/K_x^{sh})$ 在 $M$ 上的作用平凡
（记号如《基本群》引理
\ref{pione-lemma-normal-pione-quotient-inertia}）。
由此得到所需分解。

\medskip\noindent
反过来，假设有一个有限 $\Lambda$-模 $M$，
带有 $\pi_1(X, \overline{\eta})$ 的连续作用。
我们将构造 $\mathcal{F}$，使得
$M \cong \mathcal{F}_{\overline{\eta}}$ 作为
$\Lambda[\pi_1(X, \overline{\eta})]$-模。
选取生成元 $m_1, \ldots, m_r \in M$。
由于 $\pi_1(X, \overline{\eta})$ 在 $M$ 上的作用连续，
对每个 $i$ 都存在开子群 $N_i$，它是预有限群
$\pi_1(X, \overline{\eta})$ 的开子群，
使得每个 $\gamma \in H_i$ 都固定 $m_i$。因此开子群
$H = \bigcap_{i = 1, \ldots, r} H_i$ 的每个元素都固定
$M$ 的每个元素。缩小 $H$ 后，可以假设 $H$ 是
$\pi_1(X, \overline{\eta})$ 的开正规子群。
置 $G = \pi_1(X, \overline{\eta})/H$。设 $f : Y \to X$
为相应的 Galois 有限 \'etale $G$-覆盖。
可把 $f_*\underline{\mathbf{Z}}$ 看作
$\mathbf{Z}[G]$-模层，它定义在 $X_\etale$ 上。然后只须取
$$
\mathcal{F} =
f_*\underline{\mathbf{Z}} \otimes_{\underline{\mathbf{Z}[G]}} \underline{M}
$$
把 $\mathcal{F}_{\overline{\eta}}$ 的计算留给读者。
我们也略去此构造是上一段构造之逆的验证。
\end{proof}

\begin{remark}
\label{etale-cohomology-remark-functorial-locally-constant-on-connected}
引理 \ref{etale-cohomology-lemma-locally-constant-on-connected} 与
\ref{etale-cohomology-lemma-locally-constant-on-connected-geom-unibranch} 中的等价
与拉回相容。例如，假设 $f : Y \to X$ 是连通概形之间的态射。
设 $\overline{y}$ 为 $Y$ 的几何点，并置
$\overline{x} = f(\overline{y})$。则图
$$
\xymatrix{
Y_\etale\text{ 上有限局部常值的集合层}
\ar[r] &
\text{有限 }\pi_1(Y, \overline{y})\text{-集合} \\
X_\etale\text{ 上有限局部常值的集合层}
\ar[r] \ar[u]_{f^{-1}} &
\text{有限 }\pi_1(X, \overline{x})\text{-集合} \ar[u]
}
$$
可交换，其中右侧竖直箭头来自连续同态
$\pi_1(Y, \overline{y}) \to \pi_1(X, \overline{x})$；它由 $f$ 诱导。
这立即由《基本群》定理
\ref{pione-theorem-fundamental-group} 中的交换图得到。
其余情形有类似结果。
\end{remark}









\section{迹方法}
\label{etale-cohomology-section-trace-method}

\noindent
本节的参考文献是 \cite[Expos\'e IX, \S 5]{SGA4}。
这里的材料将在下文引理 \ref{etale-cohomology-lemma-vanishing-easier} 的证明中使用。

\medskip\noindent
设 $f : Y \to X$ 为概形的 \'etale 态射。有一列伴随函子
$$
f_!, f^{-1}, f_*
$$
，它们位于 $\textit{Ab}(X_\etale)$ 与 $\textit{Ab}(Y_\etale)$ 之间。
函子 $f_!$ 在第 \ref{etale-cohomology-section-extension-by-zero} 节讨论。
伴随映射 $\text{id} \to f_* f^{-1}$ 称为{\it 限制}。
伴随映射 $f_! f^{-1} \to \text{id}$ 通常称为{\it 迹映射}。
若 $f$ 有限 \'etale，则 $f_* = f_!$
（引理 \ref{etale-cohomology-lemma-shriek-equals-star-finite-etale}），因而可把它看作映射
$f_*f^{-1} \to \text{id}$。

\begin{definition}
\label{etale-cohomology-definition-trace-map}
设 $f : Y \to X$ 为有限 \'etale 概形态射。
上面描述并在下面明确给出的映射 $f_* f^{-1} \to \text{id}$
称为{\it 迹}。
\end{definition}

\noindent
设 $f : Y \to X$ 为有限 \'etale 概形态射。迹映射由以下两个性质刻画：
\begin{enumerate}
\item 它与 $X$ 上的 \'etale 局部化可交换；并且
\item 若 $Y = \coprod_{i = 1}^d X$，则迹映射是求和映射
$f_*f^{-1} \mathcal{F} = \mathcal{F}^{\oplus d} \to \mathcal{F}$。
\end{enumerate}
由《\'Etale 态射》引理 \ref{etale-lemma-finite-etale-etale-local}，
每个有限 \'etale 态射 $f : Y \to X$ 在 $X$ 的 \'etale 局部，
都具有 (2) 中给出的形式，其中某个整数 $d \geq 0$。
因此可以用上述刻画定义迹映射；特别地，为构造迹映射，
无须知道 $f_!$ 的存在性以及 $f_!$ 与 $f_*$ 的相等性。
此描述表明，若 $f$ 的常次数为 $d$，则复合
$$
\mathcal{F} \xrightarrow{res}
f_* f^{-1} \mathcal{F} \xrightarrow{trace}
\mathcal{F}
$$
是乘以 $d$。所谓“迹方法”是如下观察：若 $\mathcal{F}$
是 $X_\etale$ 上的阿贝尔群层，且乘以 $d$ 在 $\mathcal{F}$ 上
是同构，则映射
$$
H^n_\etale(X, \mathcal{F}) \longrightarrow H^n_\etale(Y, f^{-1}\mathcal{F})
$$
是单射。事实上，由高阶直像的消失
（命题 \ref{etale-cohomology-proposition-finite-higher-direct-image-zero}）
和 Leray 谱序列（命题 \ref{etale-cohomology-proposition-leray}），有
$$
H^n_\etale(Y, f^{-1}\mathcal{F}) = H^n_\etale(X, f_*f^{-1}\mathcal{F})
$$
。因此可以考虑映射
$$
H^n_\etale(X, \mathcal{F}) \to
H^n_\etale(Y, f^{-1}\mathcal{F})= H^n_\etale(X, f_*f^{-1}\mathcal{F})
\xrightarrow{trace}
H^n_\etale(X, \mathcal{F})
$$
，而复合是同构（在我们对 $\mathcal{F}$ 与 $f$ 的假设下）。
特别地，若 $H_\etale^q(Y, f^{-1}\mathcal{F}) = 0$，则
$H_\etale^q(X, \mathcal{F}) = 0$ 也成立。
确实，乘以 $d$ 在 $H_\etale^q(X, \mathcal{F})$ 上诱导同构，
而此同构经由 $H_\etale^q(Y, f^{-1}\mathcal{F})= 0$ 分解。

\medskip\noindent
这常与下面的引理结合使用。

\begin{lemma}
\label{etale-cohomology-lemma-pullback-filtered}
设 $S$ 为连通概形。设 $\ell$ 为素数。设
$\mathcal{F}$ 为有限型局部常值的
$\mathbf{F}_\ell$-向量空间层，它定义在 $S_\etale$ 上。
则存在有限 \'etale 态射 $f : T \to S$，其次数与
$\ell$ 互素，并且 $f^{-1}\mathcal{F}$
有有限滤过，其逐次商均为
$\underline{\mathbf{Z}/\ell\mathbf{Z}}_T$。
\end{lemma}

\begin{proof}
选取几何点 $\overline{s}$，它属于 $S$。
通过引理 \ref{etale-cohomology-lemma-locally-constant-on-connected} 的等价，
层 $\mathcal{F}$ 对应于有限维
$\mathbf{F}_\ell$-向量空间 $V$，它带有
$\pi_1(S, \overline{s})$ 的连续作用。
设 $G \subset \text{Aut}(V)$ 为给出该作用的同态
$\rho : \pi_1(S, \overline{s}) \to \text{Aut}(V)$ 的像。
注意 $G$ 有限。连续满同态
$\overline{\rho} : \pi_1(S, \overline{s}) \to G$
对应于 Galois 对象 $Y \to S$，它属于 $\textit{F\'Et}_S$，
其自同构群为 $G = \text{Aut}(Y/S)$，见《基本群》第
\ref{pione-section-finite-etale-under-galois} 节。
设 $H \subset G$ 为 $\ell$-Sylow 子群。
我们断言 $T = Y/H \to S$ 符合要求。事实上，设
$\overline{t} \in T$ 为位于 $\overline{s}$ 上方的几何点。
由基本群的函子性，
$\pi_1(T, \overline{t}) \to \pi_1(S, \overline{s})$ 的像
是 $(\overline{\rho})^{-1}(H)$。因此，
$\pi_1(T, \overline{t})$ 在 $V$ 上的作用（对应于
$f^{-1}\mathcal{F}$）经由映射
$\pi_1(T, \overline{t}) \to H$ 给出，见评注
\ref{etale-cohomology-remark-functorial-locally-constant-on-connected}。
由于 $H$ 是有限 $\ell$-群，表示
$\rho|_{\pi_1(T, \overline{t})}$ 的每个不可约成分
都是秩为 $1$ 的平凡表示（这是有限群表示论中的简单引理；
此处插入将来的参考文献）。
通过引理 \ref{etale-cohomology-lemma-locally-constant-on-connected} 的等价，
这意味着 $f^{-1}\mathcal{F}$ 是一列逐次扩张，
其中各常值层的取值为 $\underline{\mathbf{Z}/\ell\mathbf{Z}}_T$。
此外，$T = Y/H \to S$ 的次数与 $\ell$ 互素，
因为它等于 $H$ 在 $G$ 中的指数。
\end{proof}

\begin{lemma}
\label{etale-cohomology-lemma-action-l-group}
设 $\Lambda$ 为诺特环。
设 $\ell$ 为素数且 $n \geq 1$。
设 $H$ 为有限 $\ell$-群。
设 $M$ 为有限 $\Lambda[H]$-模，并被 $\ell^n$ 零化。
则存在有限滤过
$0 = M_0 \subset M_1 \subset \ldots \subset M_t = M$，
它由 $\Lambda[H]$-子模组成，使得 $H$ 在
$M_{i + 1}/M_i$ 上的作用平凡，对所有
$i = 0, \ldots, t - 1$ 都如此。
\end{lemma}

\begin{proof}
从略。提示：证明增广理想 $\mathfrak m$ 是幂零的；它属于非交换环
$\mathbf{Z}/\ell^n\mathbf{Z}[H]$。
\end{proof}

\begin{lemma}
\label{etale-cohomology-lemma-pullback-filtered-modules}
设 $S$ 为不可约几何单分支概形。
设 $\ell$ 为素数且 $n \geq 1$。设 $\Lambda$
为诺特环。设 $\mathcal{F}$ 为有限型局部常值的
$\Lambda$-模层，它定义在 $S_\etale$ 上，并且被 $\ell^n$ 零化。
则存在有限 \'etale 态射 $f : T \to S$，其次数与 $\ell$
互素，并且
$f^{-1}\mathcal{F}$ 有有限滤过，其逐次商形如
$\underline{M}_T$，其中取某个有限 $\Lambda$-模 $M$。
\end{lemma}

\begin{proof}
选取几何点 $\overline{s}$，它属于 $S$。
通过引理 \ref{etale-cohomology-lemma-locally-constant-on-connected-geom-unibranch}
的等价，层 $\mathcal{F}$ 对应于有限 $\Lambda$-模
$M$，它带有 $\pi_1(S, \overline{s})$ 的连续作用。
设 $G \subset \text{Aut}(V)$ 为给出该作用的同态
$\rho : \pi_1(S, \overline{s}) \to \text{Aut}(M)$ 的像。
注意，$G$ 有限，因为 $M$ 是有限 $\Lambda$-模
（见引理 \ref{etale-cohomology-lemma-locally-constant-on-connected-geom-unibranch} 的证明）。
连续满同态
$\overline{\rho} : \pi_1(S, \overline{s}) \to G$
对应于 Galois 对象 $Y \to S$，它属于 $\textit{F\'Et}_S$，
其自同构群为 $G = \text{Aut}(Y/S)$，见《基本群》第
\ref{pione-section-finite-etale-under-galois} 节。
设 $H \subset G$ 为 $\ell$-Sylow 子群。
我们断言 $T = Y/H \to S$ 符合要求。事实上，设
$\overline{t} \in T$ 为位于 $\overline{s}$ 上方的几何点。
由基本群的函子性，
$\pi_1(T, \overline{t}) \to \pi_1(S, \overline{s})$ 的像
是 $(\overline{\rho})^{-1}(H)$。因此，
$\pi_1(T, \overline{t})$ 在 $M$ 上的作用（对应于
$f^{-1}\mathcal{F}$）经由映射
$\pi_1(T, \overline{t}) \to H$ 给出，见评注
\ref{etale-cohomology-remark-functorial-locally-constant-on-connected}。
取引理 \ref{etale-cohomology-lemma-action-l-group} 中的滤过
$0 = M_0 \subset M_1 \subset \ldots \subset M_t = M$。
它诱导滤过
$0 = \mathcal{F}_0 \subset \mathcal{F}_1 \subset \ldots
\subset \mathcal{F}_t = f^{-1}\mathcal{F}$，
其逐次商是取值为 $M_{i + 1}/M_i$ 的常值层。
最后，$T = Y/H \to S$ 的次数与 $\ell$ 互素，
因为它等于 $H$ 在 $G$ 中的指数。
\end{proof}







\section{Galois 上同调}
\label{etale-cohomology-section-galois-cohomology}

\noindent
本节证明一个关于 Galois 上同调的结果
（命题 \ref{etale-cohomology-proposition-serre-galois}），所用工具是 \'etale 上同调以及
第 \ref{etale-cohomology-section-trace-method} 节中的技巧。
这将使我们能够证明域的谱上的高阶 \'etale 上同调群消失。

\begin{lemma}
\label{etale-cohomology-lemma-nonvanishing-inherited}
设 $\ell$ 是素数，$n$ 是整数且 $> 0$。
设 $S$ 是拟紧拟分离概形。
设 $X = \lim_{i \in I} X_i$ 是一个 $S$-概形有向系的极限，
其中每个 $X_i \to S$ 都是常次数的有限 \'etale 态射，且该次数与 $\ell$ 互素。
下列条件等价：
\begin{enumerate}
\item 存在一个 $\ell$-幂挠层 $\mathcal{G}$ 于 $S$ 上，使得
$H_\etale^n(S, \mathcal{G}) \neq 0$；
\item 存在一个 $\ell$-幂挠层 $\mathcal{F}$ 于 $X$ 上，使得
$H_\etale^n(X, \mathcal{F}) \neq 0$。
\end{enumerate}
事实上，给定 $\mathcal{G}$，可以取
$\mathcal{F} = g^{-1}\mathcal{F}$；给定 $\mathcal{F}$，可以取
$\mathcal{G} = g_*\mathcal{F}$。
\end{lemma}

\begin{proof}
以 $g : X \to S$ 和 $g_i : X_i \to S$ 表示结构态射。
固定一个 $\ell$-幂挠层 $\mathcal{G}$ 于 $S$ 上，使得
$H^n_\etale(S, \mathcal{G}) \not = 0$。
由 $\mathcal{G}_i = g_i^{-1}\mathcal{G}$ 给出的系统满足定理
\ref{etale-cohomology-theorem-colimit} 的条件，其余极限层为 $g^{-1}\mathcal{G}$。因此
$$
\colim_{i\in I} H^n_\etale(X_i, g_i^{-1}\mathcal{G}) =
H^n_\etale(X, \mathcal{G})
$$
由于 $g_i$ 是次数与 $\ell$ 互素的有限 \'etale 态射，可以应用“迹方法”，得到映射
$$
H^n_\etale(S, \mathcal{G}) \to H^n_\etale(X_i, g_i^{-1}\mathcal{G})
$$
全都是单射（并且与转移映射相容）。见第 \ref{etale-cohomology-section-trace-method} 节。
因此余极限非零，即 $H^n(X,g^{-1}\mathcal{G}) \neq 0$，
取 $\mathcal{F} = g^{-1}\mathcal{G}$ 即得所需结论。

\medskip\noindent
反之，设给定一个 $\ell$-幂挠层 $\mathcal{F}$ 于 $X$ 上，使得
$H^n_\etale(X, \mathcal{F}) \not = 0$。注意，由于 $g_i$ 是有限态射，
其高阶正像消失（命题 \ref{etale-cohomology-proposition-finite-higher-direct-image-zero}）。
于是应用引理 \ref{etale-cohomology-lemma-relative-colimit}，可知 $g$ 也具有同样的性质。
高阶正像的消失以及 Leray（命题 \ref{etale-cohomology-proposition-leray}）给出
$H^n_\etale(X, \mathcal{F}) = H^n(S, g_*\mathcal{F}) \neq 0$；
取 $\mathcal{G} = g_*\mathcal{F}$ 即得所需结论。
\end{proof}

\begin{lemma}
\label{etale-cohomology-lemma-reduce-to-l-group}
设 $\ell$ 是素数，$n$ 是整数且 $> 0$。设 $K$ 是域，
$G = Gal(K^{sep}/K)$，并设 $H \subset G$ 是极大 pro-$\ell$ 子群，
$L/K$ 是相应的域扩张。则
$H^n_\etale(\Spec(K), \mathcal{F}) = 0$ 对所有 $\ell$-幂挠层
$\mathcal{F}$ 成立，当且仅当
$H^n_\etale(\Spec(L), \underline{\mathbf{Z}/\ell\mathbf{Z}}) = 0$。
\end{lemma}

\begin{proof}
将 $L = \bigcup L_i$ 写成其在 $K$ 上有限子扩张的并。
由 $H$ 的选取，$[L_i : K]$ 与 $\ell$ 互素。因此如引理
\ref{etale-cohomology-lemma-nonvanishing-inherited} 中那样，
$\Spec(L) = \lim_{i \in I} \Spec(L_i)$。
故可将 $K$ 替换为 $L$，并假设 $G$，即 $K$ 的绝对 Galois 群，
是预有限 pro-$\ell$ 群。

\medskip\noindent
假设 $H^n(\Spec(K), \underline{\mathbf{Z}/\ell\mathbf{Z}}) = 0$。
设 $\mathcal{F}$ 是一个 $\ell$-幂挠层，定义在 $\Spec(K)_\etale$ 上。
我们要证明 $H^n_\etale(\Spec(K), \mathcal{F}) = 0$。
由引理 \ref{etale-cohomology-lemma-equivalence-abelian-sheaves-point} 中给出的对应，
层 $\mathcal{F}$ 对应于一个 $\ell$-幂挠 $G$-模 $M$。
由连续性，任意有限个元素 $x_1, \ldots, x_m \in M$ 都被某个开子群 $U$ 固定。
设 $M'$ 是由 $x_1, \ldots, x_m$ 的轨道生成的模。
这是有限交换 $\ell$-群，因为每个 $x_i$ 都被 $\ell$ 的某个幂零化，
并且轨道有限。由于 $M$ 是这些子模 $M'$ 的滤过余极限，
可见 $\mathcal{F}$ 是相应子层 $\mathcal{F}' \subset \mathcal{F}$ 的滤过余极限。
对此余极限应用定理 \ref{etale-cohomology-theorem-colimit}，便归结到
$\mathcal{F}$ 是有限局部常值层的情形。

\medskip\noindent
设 $M$ 是有限交换 $\ell$-群，带有预有限 pro-$\ell$ 群 $G$ 的连续作用。
则存在 $G$-不变滤过
$$
0 = M_0 \subset M_1 \subset \ldots \subset M_r = M
$$
使得 $M_{i + 1}/M_i \cong \mathbf{Z}/\ell \mathbf{Z}$，且其上的
$G$-作用平凡（这是有限群表示论中的一个简单引理；此处以后补引文）。
因此相应的层 $\mathcal{F}$ 有滤过
$$
0 = \mathcal{F}_0 \subset \mathcal{F}_1 \subset \ldots \subset
\mathcal{F}_r = \mathcal{F}
$$
其逐次商同构于 $\underline{\mathbf{Z}/\ell \mathbf{Z}}$。
于是由归纳法和上同调长正合列得到结论。
\end{proof}

\begin{lemma}
\label{etale-cohomology-lemma-reduce-to-l-group-higher}
设 $\ell$ 是素数，$n$ 是整数且 $> 0$。设 $K$ 是域，
$G = Gal(K^{sep}/K)$，并设 $H \subset G$ 是极大 pro-$\ell$ 子群，
$L/K$ 是相应的域扩张。则
$H^q_\etale(\Spec(K),\mathcal{F}) = 0$ 对 $q \geq n$ 以及所有
$\ell$-挠层 $\mathcal{F}$ 成立，当且仅当
$H^n_\etale(\Spec(L), \underline{\mathbf{Z}/\ell\mathbf{Z}}) = 0$。
\end{lemma}

\begin{proof}
正向是显然的，故只需证明反向。对 $q$ 作归纳。
$q = n$ 的情形就是引理 \ref{etale-cohomology-lemma-reduce-to-l-group}。
现设 $\mathcal{F}$ 是一个 $\ell$-幂挠层，定义在 $\Spec(K)$ 上。
设 $f : \Spec(K^{sep}) \rightarrow \Spec(K)$ 是一个几何点的包含。
考虑正合列
$$
0 \rightarrow \mathcal{F} \xrightarrow{res}
f_* f^{-1} \mathcal{F} \rightarrow f_* f^{-1} \mathcal{F}/\mathcal{F} 
\rightarrow 0
$$
注意，$K^{sep}$ 可写成有限可分扩张的滤过余极限。因此 $f$
是一个有限 \'etale 态射有向系的极限。如引理
\ref{etale-cohomology-lemma-nonvanishing-inherited} 的证明所示，可知 $f$ 的高阶正像消失。
因而 $f_* f^{-1} \mathcal{F}$ 的高阶上同调可表示为几何点上的高阶上同调，
后者显然消失。于是，由于这里的一切仍是 $\ell$-挠的，
结合归纳假设与上同调长正合列，即可得到 $q + 1$ 的结论。
\end{proof}

\begin{proposition}
\label{etale-cohomology-proposition-serre-galois}
\begin{reference}
\cite[Chapter II, Section 3, Proposition 5]{SerreGaloisCohomology}
\end{reference}
设 $K$ 是域，其可分代数闭包为 $K^{sep}$。假设对任意有限扩张
$K'$（扩张自 $K$）都有 $\text{Br}(K') = 0$。则
\begin{enumerate}
\item $H^q(\text{Gal}(K^{sep}/K), (K^{sep})^*) = 0$
对所有 $q \geq 1$ 成立；
\item $H^q(\text{Gal}(K^{sep}/K), M) = 0$
对任意挠 $\text{Gal}(K^{sep}/K)$-模 $M$ 和任意 $q \geq 2$ 成立。
\end{enumerate}
\end{proposition}

\begin{proof}
置 $p = \text{char}(K)$。由引理 \ref{etale-cohomology-lemma-compare-cohomology-point}、
定理 \ref{etale-cohomology-theorem-brauer-delta} 和例 \ref{etale-cohomology-example-sheaves-point}，
本命题等价于证明：若
$H^2(\Spec(K'),\mathbf{G}_m|_{\Spec(K')_\etale}) = 0$
对所有有限扩张 $K'/K$ 都成立，则
\begin{itemize}
\item $H^q(\Spec(K),\mathbf{G}_m|_{\Spec(K)_\etale}) = 0$
对所有 $q \geq 1$ 成立；
\item $H^q(\Spec(K),\mathcal{F}) = 0$
对任意挠层 $\mathcal{F}$ 和任意 $q \geq 2$ 成立。
\end{itemize}
先证明第二部分。由于 $\mathcal{F}$ 是挠层，可以利用 $\ell$-准素分解，
以及上同调与余极限（即直和，见定理 \ref{etale-cohomology-theorem-colimit}）的相容性，
将问题归结为证明 $H^q(\Spec(K),\mathcal{F}) = 0$, $q \geq 2$，
其中对所有 $\ell$-幂挠层并对每个素数 $\ell$ 都要成立。
这样便可逐个分析各素数。

\medskip\noindent
假设 $\ell \neq p$。对任意扩张 $K'/K$，考虑 Kummer 序列
（引理 \ref{etale-cohomology-lemma-kummer-sequence}）
$$
0 \to
\mu_{\ell, \Spec{K'}} \to
\mathbf{G}_{m, \Spec{K'}} \xrightarrow{(\cdot)^{\ell}}
\mathbf{G}_{m, \Spec{K'}} \to 0
$$
按假设，$H^q(\Spec{K'},\mathbf{G}_m|_{\Spec(K')_\etale}) = 0$
在 $q = 2$ 时成立；由定理 \ref{etale-cohomology-theorem-picard-group}，在 $q = 1$ 时，
结合 $\Pic(K) = (0)$ 也成立。故由上同调长正合列可知，
$H^2(\Spec{K'}, \mu_\ell) = 0$ 对任意可分扩张 $K'/K$ 都成立。
现设 $H$ 是极大 pro-$\ell$ 子群，位于 $K$ 的绝对 Galois 群中，
$L$ 是相应扩张。可以把 $L$ 写成有限扩张的余极限；
对此余极限应用定理 \ref{etale-cohomology-theorem-colimit}，得到
$H^2(\Spec(L), \mu_\ell) = 0$。
此时 $\mu_\ell$ 必为常值层。否则就意味着存在一个次数与 $\ell$ 互素的
$L$ 的 Galois 扩张，这与 $L$ 的定义不符（该扩张即由添加
$\ell$ 次单位根到 $L$ 而得到）。
因此由引理 \ref{etale-cohomology-lemma-reduce-to-l-group-higher}，得到 $\ell \neq p$ 时的结论。

\medskip\noindent
现假设 $\ell = p$。考虑 Artin--Schrier 正合列
（第 \ref{etale-cohomology-section-artin-schreier} 节）
$$
0 \longrightarrow \underline{\mathbf{Z}/p\mathbf{Z}}_{\Spec{K}} \longrightarrow
\mathbf{G}_{a, \Spec{K}} \xrightarrow{F-1} \mathbf{G}_{a, \Spec{K}} 
\longrightarrow 0
$$
其中 $F - 1$ 是映射 $x \mapsto x^p - x$。由注
\ref{etale-cohomology-remark-special-case-fpqc-cohomology-quasi-coherent} 以及仿射概形结构层的
高阶上同调消失（Cohomology of Schemes，引理
\ref{coherent-lemma-quasi-coherent-affine-cohomology-zero}），
$\mathbf{G}_{a, \Spec{K}}$ 的高阶上同调消失。注意，这可应用于任意特征为
$p$ 的域；特别地，可应用于域扩张 $L$，该扩张由极大 pro-$p$ 子群 $H$
确定。由此得到
$H^n(\Spec{L}, \underline{\mathbf{Z}/p\mathbf{Z}}_{\Spec{L}}) = 0$
对 $n \geq 2$ 成立；
再由引理 \ref{etale-cohomology-lemma-reduce-to-l-group-higher} 即得 $\ell = p$ 时的结论。

\medskip\noindent
为完成证明，还需证明
$H^q(\text{Gal}(K^{sep}/K), (K^{sep})^*) = 0$
对所有 $q \geq 1$ 成立。
置 $G = \text{Gal}(K^{sep}/K)$，并将 $M = (K^{sep})^*$ 视为 $G$-模。
上文已经证明 $H^1(G, M) = 0$ 和 $H^2(G, M) = 0$。考虑正合列
$$
0 \to A \to M \to M \otimes \mathbf{Q} \to B \to 0
$$
，它是 $G$-模的正合列。由上文有 $H^i(G, A) = 0$ 和
$H^i(G, B) = 0$（$i > 1$），因为 $A$ 和 $B$ 是挠 $G$-模。
由引理 \ref{etale-cohomology-lemma-profinite-group-cohomology-torsion}，有
$H^i(G, M \otimes \mathbf{Q}) = 0$（$i > 0$）。不难验证，这蕴含
$H^i(G, M) = 0$ 对 $i \geq 3$ 也成立。
\end{proof}

%10.08.09

\begin{definition}
\label{etale-cohomology-definition-Cr}
称域 $K$ 为 {\it $C_r$ 域}，如果对任意 $0 < d^r < n$ 以及任意
$f \in K[T_1,
\ldots, T_n]$，若它是次数为 $d$ 的齐次多项式，则都存在
$\alpha = (\alpha_1,
\ldots, \alpha_n)$，其中 $\alpha_i \in K$ 不全为零，使得
$f(\alpha) = 0$。
这样的 $\alpha$ 称为 $f$ 的一个{\it 非平凡解}。
\end{definition}

\begin{example}
\label{etale-cohomology-example-algebraically-closed-field-Cr}
代数闭域是 $C_r$ 域。
\end{example}

\noindent
事实上有如下简单引理。

\begin{lemma}
\label{etale-cohomology-lemma-algebraically-closed-find-solutions}
设 $k$ 是代数闭域。设 $f_1, \ldots, f_s \in k[T_1, \ldots, T_n]$
是次数分别为 $d_1, \ldots, d_s$ 的齐次多项式，其中 $d_i
> 0$。
若 $s < n$，则 $f_1 = \ldots = f_s = 0$ 有一个公共非平凡解。
\end{lemma}

\begin{proof}
这由维数理论得出，例如将 Varieties，引理
\ref{varieties-lemma-intersection-in-affine-space} 应用 $s - 1$ 次。
\end{proof}

\noindent
下面的结果计算 $C_1$ 域的布劳尔群。

\begin{theorem}
\label{etale-cohomology-theorem-C1-brauer-group-zero}
设 $K$ 是 $C_1$ 域。则 $\text{Br}(K) = 0$。
\end{theorem}

\begin{proof}
设 $D$ 是中心为 $K$ 的有限维 $K$-除代数。前面已见
$$
D \otimes_K K^{sep} \cong \text{Mat}_d(K^{sep})
$$
在内同构的意义下唯一。因此行列式 $\det :
\text{Mat}_d(K^{sep}) \to K^{sep}$ 是 Galois 不变的，并下降为次数 $d$ 的齐次映射
$$
\det = N_\text{red} : D \longrightarrow K
$$
，称为{\it 约化范数}。由于 $K$ 是 $C_1$ 域，若 $d > 1$，则存在非零
$x \in D$，使得 $N_\text{red}(x) = 0$。这显然蕴含 $x$ 不可逆，产生矛盾。
故 $\text{Br}(K) = 0$。
\end{proof}

\begin{definition}
\label{etale-cohomology-definition-variety}
设 $k$ 是域。{\it 簇}是 $k$ 上有限型的分离整概形。
{\it 曲线}是维数为 $1$ 的簇。
\end{definition}

\begin{theorem}[Tsen's theorem]
\label{etale-cohomology-theorem-tsen}
维数为 $r$ 的、代数闭域 $k$ 上的簇的函数域是 $C_r$ 域。
\end{theorem}

\begin{proof}
对射影空间，可直接证明域 $k(x_1, \ldots, x_r)$ 是 $C_r$ 域（练习）。

\medskip\noindent
一般情形。不失一般性，可以假设 $X$ 是射影的。设
$f \in k(X)[T_1, \ldots, T_n]_d$ 且 $0 < d^r < n$。
设 $f$ 的系数属于 $\Gamma(X, \mathcal{O}_X(H))$，其中
$H \subset X$ 是某个丰沛对象。取
$\mathbf{\alpha} = (\alpha_1, \ldots, \alpha_n)$，其中
$\alpha_i \in \Gamma(X,
\mathcal{O}_X(eH))$。则 $f(\mathbf{\alpha}) \in \Gamma(X,
\mathcal{O}_X((de + 1)H))$。考虑方程组 $f(\mathbf{\alpha})
=0$。由渐近 Riemann--Roch
（Varieties，命题 \ref{varieties-proposition-asymptotic-riemann-roch}），
存在 $c > 0$，使得
\begin{itemize}
\item 变量数为 $n\dim_k \Gamma(X, \mathcal{O}_X(eH)) \sim n e^r c$；
\item 方程数为
$\dim_k \Gamma(X, \mathcal{O}_X((de + 1)H)) \sim (de + 1)^r c$。
\end{itemize}
由于 $n > d^r$，变量多于方程。方程是齐次的，故由引理
\ref{etale-cohomology-lemma-algebraically-closed-find-solutions} 存在一个解。
\end{proof}

\begin{lemma}
\label{etale-cohomology-lemma-curve-brauer-zero}
设 $C$ 是代数闭域 $k$ 上的曲线。则 $C$ 的函数域的布劳尔群为零：
$\text{Br}(k(C)) = 0$。
\end{lemma}

\begin{proof}
这直接由 Tsen 定理（定理 \ref{etale-cohomology-theorem-tsen}）和定理
\ref{etale-cohomology-theorem-C1-brauer-group-zero} 得出。
\end{proof}

\begin{lemma}
\label{etale-cohomology-lemma-cohomology-Gm-function-field-curve}
设 $k$ 是代数闭域，$K/k$ 是超越次数为 1 的域扩张。则对所有
$q \geq 1$，都有 $H_\etale^q(\Spec(K), \mathbf{G}_m) = 0$。
\end{lemma}

\begin{proof}
回顾一下，由引理 \ref{etale-cohomology-lemma-compare-cohomology-point}，
$H_\etale^q(\Spec(K), \mathbf{G}_m) = H^q(\text{Gal}(K^{sep}/K), (K^{sep})^*)$。
因此由命题 \ref{etale-cohomology-proposition-serre-galois}，只需证明：若 $K'/K$ 是有限域扩张，
则 $\text{Br}(K') = 0$。注意，$K' = \colim K''$，其中 $K''$ 遍历
$k$ 的有限生成子扩张，这些子扩张包含在 $K'$ 中且超越次数为 $1$。
又有 $\text{Br}(K') = \colim \text{Br}(K'')$，故问题归结到有限生成域扩张
$K''/k$，它的超越次数为 $1$。这样的域是 $k$ 上某条曲线的函数域，
所以由引理 \ref{etale-cohomology-lemma-curve-brauer-zero}，其布劳尔群平凡。
\end{proof}






\section{乘法群的高阶消失}
\label{etale-cohomology-section-higher-Gm}

\noindent
本节固定一个代数闭域 $k$ 和一条光滑曲线 $X$，后者定义在 $k$ 上。
以 $i_x : x \hookrightarrow X$ 表示 $X$ 的闭点的包含，以
$j : \eta \hookrightarrow X$ 表示泛点的包含。又以 $X_0$ 表示 $X$ 的闭点集。

\begin{theorem}[基本正合列]
\label{etale-cohomology-theorem-fundamental-exact-sequence}
$X$ 上有一个 \'etale 层的短正合列
$$
0 \longrightarrow
\mathbf{G}_{m, X} \longrightarrow
j_* \mathbf{G}_{m, \eta} \longrightarrow
\bigoplus\nolimits_{x \in X_0} {i_x}_* \underline{\mathbf{Z}}
\longrightarrow 0.
$$
\end{theorem}

\begin{proof}
设 $\varphi : U \to X$ 是 \'etale 态射。由 \'etale 态射的性质
（命题 \ref{etale-cohomology-proposition-etale-morphisms}），有 $U = \coprod_i U_i$，
其中每个 $U_i$ 都是一条映到 $X$ 的光滑曲线。$U$ 的上述序列是每个
$U_i$ 的相应序列之积，故只需处理 $U$ 连通、从而不可约的情形。
在此情形，有熟知的正合列
$$
1 \longrightarrow
\Gamma(U, \mathcal{O}_U^*) \longrightarrow
k(U)^* \longrightarrow \bigoplus\nolimits_{y \in U_0} \mathbf{Z}_y.
$$
这等价于序列
$$
0 \longrightarrow \Gamma(U, \mathcal{O}_U^*) \longrightarrow
\Gamma(\eta \times_X U, \mathcal{O}_{\eta \times_X U}^*) \longrightarrow
\bigoplus\nolimits_{x \in X_0}
\Gamma(x \times_X U, \underline{\mathbf{Z}})
$$
展开定义，这正是序列
$$
0 \longrightarrow \mathbf{G}_m(U) \longrightarrow
j_* \mathbf{G}_{m, \eta}(U) \longrightarrow
\left(\bigoplus\nolimits_{x \in X_0} {i_x}_* \underline{\mathbf{Z}}\right)(U).
$$
这定义了基本正合列中的映射，并表明该列除最后一步可能尚未验证外是正合的。
为看出满射性，回顾：若 $U$ 是非奇异曲线，$D$ 是 $U$ 上的除子，
则存在一个 Zariski 开覆盖 $\{U_j \to U\}$ 于 $U$ 上，使得
$D|_{U_j} = \text{div}(f_j)$，其中某个 $f_j \in k(U)^*$ 具有此性质。
\end{proof}

\begin{lemma}
\label{etale-cohomology-lemma-higher-direct-jstar-Gm}
对任意 $q \geq 1$，都有 $R^q j_*\mathbf{G}_{m, \eta} = 0$。
\end{lemma}

\begin{proof}
需证明 $(R^q j_*\mathbf{G}_{m, \eta})_{\bar x} = 0$ 对每个几何点
$\bar x$（它是 $X$ 的点）都成立。

\medskip\noindent
假设 $\bar x$ 位于闭点 $x$ 之上，该点属于 $X$。设 $\Spec(A)$ 是 $x$ 在 $X$
中的一个仿射开邻域，$K$ 是 $A$ 的分式域。则
$$
\Spec(\mathcal{O}^{sh}_{X, \bar x}) \times_X \eta =
\Spec(\mathcal{O}^{sh}_{X, \bar x} \otimes_A K).
$$
环 $\mathcal{O}^{sh}_{X, \bar x} \otimes_A K$ 是离散赋值环
$\mathcal{O}^{sh}_{X, \bar x}$ 的一个局部化，故它要么仍是
$\mathcal{O}^{sh}_{X, \bar x}$，要么是其分式域 $K^{sh}_{\bar x}$。
但由于某个局部一致化元被取逆，它必为后者。因此
$$
(R^q j_*\mathbf{G}_{m, \eta})_{(X, \bar x)} = H_\etale^q(\Spec
K^{sh}_{\bar x}, \mathbf{G}_m).
$$
现回顾 $\mathcal{O}^{sh}_{X, \bar x} =
\colim_{(U, \bar u) \to \bar x} \mathcal{O}(U) = \colim_{A \subset B} B$
，其中 $A \to B$ 是 \'etale 的。因此 $K^{sh}_{\bar x}$ 是
$K = k(X)$ 的代数扩张，可以应用引理
\ref{etale-cohomology-lemma-cohomology-Gm-function-field-curve} 得到消失结论。

\medskip\noindent
假设 $\bar x = \bar \eta$ 位于泛点 $\eta$ 之上，而后者属于 $X$
（事实上这一情形是多余的）。此时
$\mathcal{O}^{sh}_{X, \bar \eta} = \kappa(\eta)^{sep}$，故
\begin{eqnarray*}
(R^q j_*\mathbf{G}_{m, \eta})_{\bar \eta}
& = &
H_\etale^q(\Spec(\kappa(\eta)^{sep}) \times_X \eta, \mathbf{G}_m) \\
& = & H_\etale^q (\Spec(\kappa(\eta)^{sep}), \mathbf{G}_m) \\
& = & 0 \ \ \text{ 对 } q \geq 1
\end{eqnarray*}
，因为相应的 Galois 群是平凡的。
\end{proof}

\begin{lemma}
\label{etale-cohomology-lemma-cohomology-jstar-Gm}
对所有 $p \geq 1$，都有 $H_\etale^p(X, j_*\mathbf{G}_{m, \eta}) = 0$。
\end{lemma}

\begin{proof}
Leray 谱序列为
$$
E_2^{p, q} = H_\etale^p(X, R^qj_*\mathbf{G}_{m, \eta}) \Rightarrow
H_\etale^{p+q}(\eta, \mathbf{G}_{m, \eta}),
$$
由引理 \ref{etale-cohomology-lemma-cohomology-Gm-function-field-curve}，右端在
$p + q \geq 1$ 时消失。由引理 \ref{etale-cohomology-lemma-higher-direct-jstar-Gm}，
左端在 $q \geq 1$ 时消失。故得到所需的消失结论。
\end{proof}

\begin{lemma}
\label{etale-cohomology-lemma-cohomology-istar-Z}
对所有 $q \geq 1$，都有 $H_\etale^q(X, \bigoplus_{x \in X_0} {i_x}_*
\underline{\mathbf{Z}}) = 0$。
\end{lemma}

\begin{proof}
由于 $X$ 拟紧拟分离，上同调与余极限可交换，故只需证明
$H_\etale^q(X, {i_x}_*
\underline{\mathbf{Z}})$ 消失。但闭点的包含
$i_x$ 是有限的，故由命题 \ref{etale-cohomology-proposition-finite-higher-direct-image-zero}，
有 $R^p {i_x}_* \underline{\mathbf{Z}} = 0$ 对所有 $p \geq 1$ 成立。
应用 Leray 谱序列可见
$H_\etale^q(X, {i_x}_* \underline{\mathbf{Z}}) =
H_\etale^q(x, \underline{\mathbf{Z}})$.
最后，由于 $x$ 是一个代数闭域的谱，$x$ 上所有高阶上同调都消失。
\end{proof}

\noindent
由以上一系列引理，得到如下结果。

\begin{theorem}
\label{etale-cohomology-theorem-vanishing-cohomology-Gm-curve}
设 $X$ 是代数闭域上的光滑曲线。则
$$
H_\etale^q(X, \mathbf{G}_m) = 0 \ \ \text{ 对所有 } q \geq 2.
$$
\end{theorem}

\begin{proof}
见上面的讨论。
\end{proof}

\noindent
还得到上同调长正合列
$$
0 \to
H_\etale^0(X, \mathbf{G}_m) \to
H_\etale^0(X, j_*\mathbf{G}_{m\eta}) \to
H_\etale^0(X, \bigoplus {i_x}_*\underline{\mathbf{Z}}) \to
H_\etale^1(X, \mathbf{G}_m) \to 0
$$
不过这就是熟知的序列
$$
0 \to H_{Zar}^0(X, \mathcal{O}_X^*) \to k(X)^* \to \text{Div}(X)
\to \Pic(X) \to 0.
$$





\section{曲线的 Picard 群}
\label{etale-cohomology-section-pic-curves}

\noindent
下一步利用 Kummer 序列推出曲线以有限群为系数的上同调群的一些信息。
为在长正合列中得到消失结论，先回顾关于 Picard 群的一些事实。

\medskip\noindent
设 $X$ 是代数闭域 $k$ 上的光滑射影曲线。以
$g = \dim_k H^1(X, \mathcal{O}_X)$ 表示 $X$ 的亏格。存在短正合列
$$
0 \to \Pic^0(X) \to \Pic(X) \xrightarrow{\deg} \mathbf{Z} \to 0.
$$
交换群 $\Pic^0(X)$ 可与 $\Pic^0(X) = \underline{\Picardfunctor}^0_{X/k}(k)$
等同，即它是 $k$-值点，所对应的阿贝尔簇
$\underline{\Picardfunctor}^0_{X/k}$ 定义在 $k$ 上且维数为 $g$。
因此，若 $n \in k^*$，则
$\Pic^0(X)[n] \cong (\mathbf{Z}/n\mathbf{Z})^{2g}$
作为交换群同构。见 Picard Schemes of Curves，第
\ref{pic-section-picard-curve} 节，以及 Groupoids，第
\ref{groupoids-section-abelian-varieties} 节。这个关键事实，即代数闭域上
光滑射影曲线的 Picard 群中挠元的描述，看来没有初等证明。

\begin{lemma}
\label{etale-cohomology-lemma-cohomology-smooth-projective-curve}
设 $X$ 是亏格为 $g$ 的光滑射影曲线，定义在代数闭域 $k$ 上，并设
$n \geq 1$ 在 $k$ 中可逆。则有典范等同
$$
H_\etale^q(X, \mu_n) =
\left\{
\begin{matrix}
\mu_n(k) & \text{ 若 }q = 0, \\
\Pic^0(X)[n] & \text{ 若 }q = 1, \\
\mathbf{Z}/n\mathbf{Z} & \text{ 若 }q = 2, \\
0 & \text{ 若 }q \geq 3.
\end{matrix}
\right.
$$
由于 $\mu_n \cong \underline{\mathbf{Z}/n\mathbf{Z}}$，这给出
（非典范的）等同
$$
H_\etale^q(X, \underline{\mathbf{Z}/n\mathbf{Z}}) \cong
\left\{
\begin{matrix}
\mathbf{Z}/n\mathbf{Z} & \text{ 若 }q = 0, \\
(\mathbf{Z}/n\mathbf{Z})^{2g} & \text{ 若 }q = 1, \\
\mathbf{Z}/n\mathbf{Z} & \text{ 若 }q = 2, \\
0 & \text{ 若 }q \geq 3.
\end{matrix}
\right.
$$
\end{lemma}

\begin{proof}
定理 \ref{etale-cohomology-theorem-picard-group} 和
\ref{etale-cohomology-theorem-vanishing-cohomology-Gm-curve} 确定了
$\mathbf{G}_m$ 在 $X$ 上的 \'etale 上同调，其表述使用 $X$ 的 Picard 群。
于是 Kummer 序列
$0\to \mu_{n, X} \to \mathbf{G}_{m, X}
\to \mathbf{G}_{m, X}\to 0$（引理 \ref{etale-cohomology-lemma-kummer-sequence}）
给出上同调长正合列
$$
\xymatrix{
0 \ar[r] & \mu_n(k) \ar[r] &
k^* \ar[r]^{(\cdot)^n} &
k^* \ar@(rd, ul)[rdllllr] \\
& H_\etale^1(X, \mu_n) \ar[r] &
\Pic(X) \ar[r]^{(\cdot)^n} &
\Pic(X) \ar@(rd, ul)[rdllllr] \\
& H_\etale^2(X, \mu_n) \ar[r] & 0 \ar[r] & 0 \ldots
}
$$
这里的 $n$ 次幂映射 $k^* \to k^*$ 是满射，因为 $k$ 是代数闭的。
因此需要计算映射 $n : \Pic(X) \to \Pic(X)$ 的核与余核。
考虑如下各行正合的交换图
$$
\xymatrix{
0 \ar[r] &
\Pic^0(X) \ar[r] \ar@{>>}[d]^{(\cdot)^n} &
\Pic(X) \ar[r]_-\deg \ar[d]^{(\cdot)^n} &
\mathbf{Z} \ar[r] \ar@{^{(}->}[d]^n & 0 \\
0 \ar[r] &
\Pic^0(X) \ar[r] &
\Pic(X) \ar[r]^-\deg &
\mathbf{Z} \ar[r] & 0
}
$$
群 $\Pic^0(X)$ 是 $k$-点，属于群概形
$\underline{\Picardfunctor}^0_{X/k}$，
见 Picard Schemes of Curves，引理 \ref{pic-lemma-picard-pieces}。
同一引理还说明 $\underline{\Picardfunctor}^0_{X/k}$ 是一个
$g$ 维阿贝尔簇，定义在 $k$ 上；这里的定义见 Groupoids，定义
\ref{groupoids-definition-abelian-variety}。故由 Groupoids，命题
\ref{groupoids-proposition-review-abelian-varieties}，左边的竖直映射是满射。
应用蛇引理便得到本引理所述的典范等同。

\medskip\noindent
为得到本引理中的非典范等同，需要证明
$n : \Pic^0(X) \to \Pic^0(X)$ 的核同构于
$(\mathbf{Z}/n\mathbf{Z})^{\oplus 2g}$。这也是 Groupoids，命题
\ref{groupoids-proposition-review-abelian-varieties} 的一部分。
\end{proof}

\begin{lemma}
\label{etale-cohomology-lemma-pullback-on-h2-curve}
设 $\pi : X \to Y$ 是代数闭域 $k$ 上光滑射影曲线之间的非常值态射，
并设 $n \geq 1$ 在 $k$ 中可逆。映射
$$
\pi^* : H^2_\etale(Y, \mu_n) \longrightarrow H^2_\etale(X, \mu_n)
$$
由乘以 $\pi$ 的次数给出。
\end{lemma}

\begin{proof}
注意，该陈述有意义，因为在引理
\ref{etale-cohomology-lemma-cohomology-smooth-projective-curve} 中，我们已经把两个上同调群
$ H^2_\etale(Y, \mu_n)$ 和 $H^2_\etale(X, \mu_n)$ 都与
$\mathbf{Z}/n\mathbf{Z}$ 等同。事实上，若 $\mathcal{L}$ 是次数为
$1$ 的线丛，定义在 $Y$ 上，其类为
$[\mathcal{L}] \in H^1_\etale(Y, \mathbf{G}_m)$，
则 $[\mathcal{L}]$ 的余边界是 $H^2_\etale(Y, \mu_n)$ 的生成元。
这里的余边界是与 Kummer 序列相伴的上同调长正合列中的余边界。
因此，本引理的结果由如下事实得出：线丛 $\pi^*\mathcal{L}$ 在
$X$ 上的次数为 $\deg(\pi)$。省略若干细节。
\end{proof}

\begin{lemma}
\label{etale-cohomology-lemma-vanishing-cohomology-mu-smooth-curve}
设 $X$ 是代数闭域 $k$ 上的仿射光滑曲线，且 $n \in k^*$。
设 $X \subset \overline{X}$ 是一个光滑射影紧化
（Varieties，注 \ref{varieties-remark-smooth-projective-compactification}）。
设 $g$ 是 $\overline{X}$ 的亏格，$r$ 是 $\overline{X} \setminus X$
中的点数。则
\begin{enumerate}
\item $H_\etale^0(X, \mu_n) = \mu_n(k)$;
\item $H_\etale^1(X, \mu_n) \cong (\mathbf{Z}/n\mathbf{Z})^{2g+r-1}$；
\item $H_\etale^q(X, \mu_n) = 0$ 对所有 $q \geq 2$ 成立。
\end{enumerate}
\end{lemma}

\begin{proof}
写成 $X = \overline{X} - \{ x_1, \ldots, x_r\}$。则
$\Pic(X) = \Pic(\overline{X})/ R$，其中 $R$ 是由
$\mathcal{O}_{\overline{X}}(x_i)$（$1 \leq i \leq r$）生成的子群。
由于 $r \geq 1$，可见 $\Pic^0(\overline{X}) \to \Pic(X)$ 是满射，
故 $\Pic(X)$ 可除（见引理
\ref{etale-cohomology-lemma-cohomology-smooth-projective-curve} 的证明中的讨论）。
应用 Kummer 序列得到 (1) 和 (3)。对 (2)，使用
\begin{align*}
H_\etale^1(X, \mu_n)
& =
\{(\mathcal L, \alpha)\}/\cong \\
& =
(\{(\bar{\mathcal L}, D, \bar \alpha)\}/\cong)/\tilde{R}
\end{align*}
第一个等号和记号见引理 \ref{etale-cohomology-lemma-describe-h1-mun}。
第二行中的三元组 $(\bar{\mathcal L}, D, \bar \alpha)$ 由以下数据组成：
可逆模 $\bar{\mathcal L}$，定义在 $\overline{X}$ 上且次数为 $0$；
除子 $D$，定义在 $\overline{X}$ 上且次数为 $0$，其支集包含于
$\left\{x_1, \ldots, x_r\right\}$；以及同构
$\bar{\alpha}: \bar{\mathcal L}^{\otimes n} \cong \mathcal{O}_{\bar{X}}(D)$。
三元组的同构按引理 \ref{etale-cohomology-lemma-describe-h1-mun} 之前讨论中的方式定义。
此外，$\tilde{R}$ 是由如下形式的三元组生成的子群：
$(\mathcal{O}_{\overline{X}}(D'), n D', 1^{\otimes n})$
，其中 $D'$ 的支集包含于 $\left\{x_1, \ldots, x_r\right\}$ 且次数为 0。
上述讨论表明，第二个等号由如下映射给出：它把三元组
$(\bar{\mathcal L},\ D,\ \bar \alpha)$ 送到对
$(\bar{\mathcal L}|_X, \bar \alpha|_X)$。于是得到正合列
$$
0 \longrightarrow
H_\etale^1(\overline{X}, \mu_n) \longrightarrow
H_\etale^1(X, \mu_n) \longrightarrow
\bigoplus_{i = 1}^r \mathbf{Z}/n\mathbf{Z}
\xrightarrow{\ \sum\ }
\mathbf{Z}/n\mathbf{Z} \longrightarrow 0
$$
其中中间的映射把三元组 $(\bar{ \mathcal L}, D, \bar \alpha)$ 的类
（这里 $D = \sum_{i = 1}^r a_i (x_i)$）送到 $r$-元组
$(a_i)_{i = 1}^r$。现只需利用引理
\ref{etale-cohomology-lemma-cohomology-smooth-projective-curve} 计算秩。
\end{proof}

\begin{remark}
\label{etale-cohomology-remark-natural-proof}
证明前一个推论的“自然”方法是把 $X$ 从 $\bar
X$ 中切除。
这是可行的，只是我们尚未建立该理论。
\end{remark}

\begin{remark}
\label{etale-cohomology-remark-normalize-H1-Gm}
设 $k$ 是代数闭域。设 $n$ 是与 $k$ 的特征互素的整数。回顾
$$
\mathbf{G}_{m, k} = \mathbf{A}^1_k \setminus \{0\} =
\mathbf{P}^1_k \setminus \{0, \infty\}
$$
我们断言存在典范同构
$$
H^1_\etale(\mathbf{G}_{m, k}, \mu_n) = \mathbf{Z}/n\mathbf{Z}
$$
这是什么意思？这意味着存在元素 $1_k$ 于
$H^1_\etale(\mathbf{G}_{m, k}, \mu_n)$ 中，使得对每个态射
$\Spec(k') \to \Spec(k)$，映射 $\mathbf{G}_{m, k'} \to \mathbf{G}_{m, k}$
在 \'etale 上同调上诱导的拉回映射把 $1_k$ 送到 $1_{k'}$。
（特别地，此元素被 $k$ 的所有自同构固定。）为看出这一点，考虑
$\mu_{n, \mathbf{Z}}$-挠子
$\mathbf{G}_{m, \mathbf{Z}} \to \mathbf{G}_{m, \mathbf{Z}}$，
$x \mapsto x^n$。由挠子与一阶上同调的等同，拉回此挠子便得到我们的
典范元素 $1_k$。反向扭曲可见存在典范等同
$$
H^1_\etale(\mathbf{G}_{m, k}, \mathbf{Z}/n\mathbf{Z}) =
\Hom(\mu_n(k), \mathbf{Z}/n\mathbf{Z}),
$$
即这些同构关于代数闭域之间的映射相容，特别地，关于 $k$ 的自同构相容。
\end{remark}










\section{零延拓}
\label{etale-cohomology-section-extension-by-zero}

\noindent
Modules on Sites，第 \ref{sites-modules-section-localize} 节中的一般材料
使我们可以作如下定义。

\begin{definition}
\label{etale-cohomology-definition-extension-zero}
设 $j : U \to X$ 是概形的 \'etale 态射。
\begin{enumerate}
\item 限制函子
$j^{-1} : \Sh(X_\etale) \to \Sh(U_\etale)$
有左伴随
$j_!^{Sh} : \Sh(U_\etale) \to \Sh(X_\etale)$.
\item 限制函子
$j^{-1} : \textit{Ab}(X_\etale) \to \textit{Ab}(U_\etale)$
有左伴随，记为
$j_! : \textit{Ab}(U_\etale) \to \textit{Ab}(X_\etale)$
，称为{\it 零延拓}。
\item 设 $\Lambda$ 是环。限制函子
$j^{-1} : \textit{Mod}(X_\etale, \Lambda) \to
\textit{Mod}(U_\etale, \Lambda)$
有左伴随，记为
$j_! : \textit{Mod}(U_\etale, \Lambda) \to
\textit{Mod}(X_\etale, \Lambda)$
，称为{\it 零延拓}。
\end{enumerate}
\end{definition}

\noindent
若 $\mathcal{F}$ 是 $X_\etale$ 上的交换群层，则一般有
$j_!\mathcal{F} \not = j_!^{Sh}\mathcal{F}$。另一方面，函子 $j_!$ 对
$\Lambda$-模层的作用与底层交换群层上的 $j_!$ 一致
（Modules on Sites，注 \ref{sites-modules-remark-localize-shriek-equal}）。
函子 $j_!$ 由函子性同构
$$
\Hom_X(j_!\mathcal{F}, \mathcal{G}) = \Hom_U(\mathcal{F}, j^{-1}\mathcal{G})
$$
刻画，其中对所有 $\mathcal{F} \in \textit{Ab}(U_\etale)$ 和
$\mathcal{G} \in \textit{Ab}(X_\etale)$ 均成立。$\Lambda$-模层的情形类似。

\medskip\noindent
为更明确地描述定义 \ref{etale-cohomology-definition-extension-zero} 中的函子，回顾：
$j^{-1}$ 就是沿函子 $U_\etale \to X_\etale$ 的限制。换言之，
$j^{-1}\mathcal{G}(U') = \mathcal{G}(U')$ 对 $U'$ 成立，其中它在
$U$ 上是 \'etale 的。另一方面，对
$\mathcal{F} \in \textit{Ab}(U_\etale)$，考虑预层
\begin{equation}
\label{etale-cohomology-equation-j-p-shriek}
j_{p!}\mathcal{F} : X_\etale \longrightarrow \textit{Ab},
\quad
V \longmapsto \bigoplus\nolimits_{V \to U} \mathcal{F}(V \to U)
\end{equation}
则 $j_!\mathcal{F}$ 是 $j_{p!}\mathcal{F}$ 的层化。
这一点在 Modules on Sites，引理
\ref{sites-modules-lemma-extension-by-zero} 中证明；更一般地，见
Modules on Sites，第 \ref{sites-modules-section-localize} 节和第
\ref{sites-modules-section-exactness-lower-shriek} 节中的讨论。

\begin{exercise}
\label{etale-cohomology-exercise-jshriek-direct}
直接证明：函子 $j_!$（它定义为 (\ref{etale-cohomology-equation-j-p-shriek}) 中函子
$j_{p!}$ 的层化）是 $j^{-1}$ 的左伴随。
\end{exercise}

\begin{proposition}
\label{etale-cohomology-proposition-describe-jshriek}
设 $j : U \to X$ 是概形的 \'etale 态射。设
$\mathcal{F}$ 属于 $\textit{Ab}(U_\etale)$。若
$\overline{x} : \Spec(k) \to X$ 是 $X$ 的几何点，则
$$
(j_!\mathcal{F})_{\overline{x}} =
\bigoplus\nolimits_{\overline{u} : \Spec(k) \to U,\ j(\overline{u}) =
\overline{x}} \mathcal{F}_{\bar{u}}.
$$
特别地，$j_!$ 是正合函子。
\end{proposition}

\begin{proof}
$j_!$ 的正合性是非常一般的，见 Modules on Sites，引理
\ref{sites-modules-lemma-extension-by-zero-exact}。当然，它也可由茎的描述得出。
茎的公式由 Modules on Sites，引理
\ref{sites-modules-lemma-stalk-j-shriek}，以及用几何点描述小 \'etale
网站的点这一结果得出；后者见引理 \ref{etale-cohomology-lemma-points-small-etale-site}。

\medskip\noindent
为以后使用，注意同构
\begin{align*}
(j_!\mathcal{F})_{\overline{x}}
& =
(j_{p!}\mathcal{F})_{\overline{x}} \\
& =
\colim_{(V, \overline{v})} j_{p!}\mathcal{F}(V) \\
& =
\colim_{(V, \overline{v})}
\bigoplus\nolimits_{\varphi : V \to U}
\mathcal{F}(V \xrightarrow{\varphi} U) \\
& \to
\bigoplus\nolimits_{\overline{u} : \Spec(k) \to U,\ j(\overline{u}) =
\overline{x}} \mathcal{F}_{\bar{u}}.
\end{align*}
在 Modules on Sites，引理 \ref{sites-modules-lemma-stalk-j-shriek}
中构造；它把 $(V, \overline{v}, \varphi, s)$ 送到 $s$ 在
$\mathcal{F}$ 于 $\overline{u} = \varphi(\overline{v})$ 处的茎中的类。
\end{proof}

\begin{lemma}
\label{etale-cohomology-lemma-jshriek-open}
设 $j : U \to X$ 是概形的开浸入。对任意交换群层 $\mathcal{F}$ 于
$U_\etale$ 上，伴随映射 $j^{-1}j_*\mathcal{F} \to \mathcal{F}$ 和
$\mathcal{F} \to j^{-1}j_!\mathcal{F}$ 都是同构。事实上，
$j_!\mathcal{F}$ 是 $X_\etale$ 上唯一的交换群层，其到 $U$ 的限制为
$\mathcal{F}$，并且它在 $X \setminus U$ 的几何点处的茎为零。
\end{lemma}

\begin{proof}
我们鼓励读者展开定义证明第一项陈述，但这里仅使用它是非常一般的
Modules on Sites，引理 \ref{sites-modules-lemma-restrict-back} 的特殊情形。
对第二项陈述，注意：若 $\mathcal{G}$ 是 $X_\etale$ 上的交换群层，
其到 $U$ 的限制为 $\mathcal{F}$，则由伴随性得到映射
$j_!\mathcal{F} \to \mathcal{G}$。由命题
\ref{etale-cohomology-proposition-describe-jshriek}，此映射在 $U$ 的几何点的茎上是同构。
因此，若 $\mathcal{G}$ 在 $X \setminus U$ 的几何点处的茎消失，
则由定理 \ref{etale-cohomology-theorem-exactness-stalks}，
$j_!\mathcal{F} \to \mathcal{G}$ 是同构。
\end{proof}

\begin{lemma}[零延拓与基变换可交换]
\label{etale-cohomology-lemma-shriek-base-change}
设 $f: Y \to X$ 是概形态射。设 $j: V \to X$ 是 \'etale 态射。
考虑纤维积
$$
\xymatrix{
V' = Y \times_X V \ar[d]_{f'} \ar[r]_-{j'} & Y \ar[d]^f \\
V \ar[r]^j & X
}
$$
则在交换群层和模层上都有 $j'_! f'^{-1} = f^{-1} j_!$。
\end{lemma}

\begin{proof}
这是因为 $j'_! f'^{-1}$ 是 $f'_* (j')^{-1}$ 的左伴随，而
$f^{-1} j_!$ 是 $j^{-1}f_*$ 的左伴随。又有
$f'_* (j')^{-1} = j^{-1}f_*$，因为 $f_*$ 按构造与 \'etale 局部化可交换。
事实上，本引理在网站态射的框架下非常一般地成立，见
Modules on Sites，引理
\ref{sites-modules-lemma-localize-morphism-ringed-sites}。
\end{proof}

\begin{lemma}
\label{etale-cohomology-lemma-shriek-into-star-separated-etale}
设 $j : U \to X$ 是分离 \'etale 态射。则有函子性单射
$j_!\mathcal{F} \to j_*\mathcal{F}$，定义在交换群层和 $\Lambda$-模层上。
\end{lemma}

\begin{proof}
我们在交换群层的情形证明此结论。对任意交换群层，构造典范映射
$$
j_{p!}\mathcal{F} \to j_*\mathcal{F}
$$
，它是 $X_\etale$ 上交换群预层的映射；这里 $\mathcal{F}$ 定义在
$U_\etale$ 上，$j_{p!}$ 如 (\ref{etale-cohomology-equation-j-p-shriek}) 中所定义。
将此映射层化便得到所需映射 $j_!\mathcal{F} \to j_*\mathcal{F}$。
在 \'etale 态射 $V \to X$ 上计算两端，得到
$$
j_{p!}\mathcal{F}(V) =
\bigoplus\nolimits_{\varphi : V \to U} \mathcal{F}(V \xrightarrow{\varphi} U)
\quad\text{且}\quad
j_*\mathcal{F}(V) = \mathcal{F}(V \times_X U)
$$
对每个 $\varphi$，有既开又闭的浸入
$$
\Gamma_\varphi = (1, \varphi) : V \longrightarrow V \times_X U
$$
，它定义在 $U$ 上。它是开的，因为它是 $U$ 上 \'etale 概形之间的态射；
它是闭的，因为它是一个 $V$ 上分离概形的截面（Schemes，引理
\ref{schemes-lemma-section-immersion}）。因此，对截面
$s_\varphi \in \mathcal{F}(V \xrightarrow{\varphi} U)$，存在唯一截面
$s'_\varphi$ 于 $\mathcal{F}(V \times_X U)$ 中，使其拉回为
$s_\varphi$，所沿映射为 $\Gamma_\varphi$，且其到 $\Gamma_\varphi$
的像之补集的限制为零。

\medskip\noindent
为证明映射是单射，假设上述公式中的元素
$\sum_{i = 1, \ldots, n} s_{\varphi_i}$ 属于 $j_{p!}\mathcal{F}(V)$，
且它在 $j_*\mathcal{F}(V)$ 中的像为零。任务是证明
$\sum_{i = 1, \ldots, n} s_{\varphi_i}$ 到 $V$ 的某个 \'etale 覆盖的
各成员上的限制为零。考察所有两两等化子（它们在 $V$ 中既开又闭），
这些等化子来自态射 $\varphi_i : V \to U$；在 $V$ 上局部工作，可以假设态射
$\Gamma_{\varphi_1}, \ldots, \Gamma_{\varphi_n}$ 的像两两不交。
由于假设 $\sum_{i = 1, \ldots, n} s'_{\varphi_i} = 0$，立即得到
$s'_{\varphi_i} = 0$ 对每个 $i$ 都成立（因为这些截面的支集两两不交），
进而 $s_{\varphi_i} = 0$ 对所有 $i$ 都成立，如所需。
\end{proof}

\begin{lemma}
\label{etale-cohomology-lemma-shriek-equals-star-finite-etale}
设 $j : U \to X$ 是有限 \'etale 态射。则映射
$j_! \to j_*$（见引理 \ref{etale-cohomology-lemma-shriek-into-star-separated-etale}）
在交换群层和 $\Lambda$-模层上都是同构。
\end{lemma}

\begin{proof}
只需验证 $j_!\mathcal{F} \to j_*\mathcal{F}$ 是同构，且可在
$X$ 上 \'etale 局部地验证。因此可以假设
$U \to X$ 是有限个同构的不交并，见 \'Etale Morphisms，引理
\ref{etale-lemma-finite-etale-etale-local}。省略这一情形的证明。
\end{proof}

\begin{lemma}
\label{etale-cohomology-lemma-ses-associated-to-open}
设 $X$ 是概形。设 $Z \subset X$ 是闭子概形，$U \subset X$ 是其补集。
以 $i : Z \to X$ 和 $j : U \to X$ 表示包含态射。对每个交换群层
$\mathcal{F}$ 于 $X_\etale$ 上，有典范短正合列
$$
0 \to j_!j^{-1}\mathcal{F} \to \mathcal{F} \to i_*i^{-1}\mathcal{F} \to 0
$$
于 $X_\etale$ 上。
\end{lemma}

\begin{proof}
这些映射由所涉函子的伴随性质得到。对几何点 $\overline{x}$ 于
$X$ 中，或者 $\overline{x} \in U$；此时左侧映射在茎上是同构，
而 $i_*i^{-1}\mathcal{F}$ 的茎为零。或者 $\overline{x} \in Z$；
此时右侧映射在茎上是同构，而 $j_!j^{-1}\mathcal{F}$ 的茎为零。
这里使用了引理 \ref{etale-cohomology-lemma-stalk-pushforward-closed-immersion} 和命题
\ref{etale-cohomology-proposition-describe-jshriek} 中对茎的描述。
\end{proof}

\begin{lemma}
\label{etale-cohomology-lemma-compatible-shriek-push-finite}
考虑概形的笛卡儿图
$$
\xymatrix{
U \ar[d]_g \ar[r]_{j'} & X \ar[d]^f \\
V \ar[r]^j & Y
}
$$
，其中 $f$ 是有限态射，$j$ 是开浸入。则有
$f_* \circ j'_! = j_! \circ g_*$，作为函子
$\textit{Ab}(U_\etale) \to \textit{Ab}(Y_\etale)$。
\end{lemma}

\begin{proof}
设 $\mathcal{F}$ 是 $\textit{Ab}(U_\etale)$ 的对象。设
$\overline{y}$ 是 $Y$ 的几何点，且不包含在开集 $V$ 中。则
$$
(f_*j'_!\mathcal{F})_{\overline{y}} =
\bigoplus\nolimits_{\overline{x},\ f(\overline{x}) = \overline{y}}
(j'_!\mathcal{F})_{\overline{x}} = 0
$$
由命题 \ref{etale-cohomology-proposition-finite-higher-direct-image-zero}，并且由引理
\ref{etale-cohomology-lemma-jshriek-open}，$j'_!\mathcal{F}$ 在
$\overline{x} \not \in U$ 处的茎为零。另一方面，有
$$
j^{-1}f_*j'_!\mathcal{F} =
g_*(j')^{-1}j'_!\mathcal{F} =
g_*\mathcal{F}
$$
，这由引理 \ref{etale-cohomology-lemma-finite-pushforward-commutes-with-base-change} 和
\ref{etale-cohomology-lemma-jshriek-open} 得出。因此由引理 \ref{etale-cohomology-lemma-jshriek-open}
中对 $j_!$ 的刻画，可知 $f_*j'_!\mathcal{F} = j_!g_*\mathcal{F}$。
省略验证此等同关于 $\mathcal{F}$ 具有函子性。
\end{proof}







\section{可构造层}
\label{etale-cohomology-section-constructible}

\noindent
设 $X$ 为概形。$X$ 的一个{\it 可构造局部闭子概形}，是指一个局部闭子概形
$T \subset X$，使得 $T$ 的底拓扑空间是 $X$ 的可构造子集。
若 $T, T' \subset X$ 是具有同一底拓扑空间的局部闭子概形，则由 \'etale
位点的拓扑不变性（定理 \ref{etale-cohomology-theorem-topological-invariance}），有
$T_\etale \cong T'_\etale$。因此，在以下定义中可以假定所考虑的局部闭
子概形都是既约的。

\begin{definition}
\label{etale-cohomology-definition-constructible}
设 $X$ 为概形。
\begin{enumerate}
\item 若 $X_\etale$ 上的集合层满足如下条件，则称它是{\it 可构造的}：
对每个仿射开集 $U \subset X$，都存在将 $U$ 有限分解为可构造局部闭
子概形的分解 $U = \coprod_i U_i$，使得 $\mathcal{F}|_{U_i}$ 对所有 $i$
都是有限局部常值的。
\item 若 $X_\etale$ 上的阿贝尔群层满足如下条件，则称它是{\it 可构造的}：
对每个仿射开集 $U \subset X$，都存在将 $U$ 有限分解为可构造局部闭
子概形的分解 $U = \coprod_i U_i$，使得 $\mathcal{F}|_{U_i}$ 对所有 $i$
都是有限局部常值的。
\item 设 $\Lambda$ 为诺特环。若一个 $\Lambda$-模层在 $X_\etale$ 上满足如下条件，
则称它是{\it 可构造的}：对每个仿射开集 $U \subset X$，都存在将 $U$
有限分解为可构造局部闭子概形的分解 $U = \coprod_i U_i$，使得
$\mathcal{F}|_{U_i}$ 对所有 $i$ 都是有限型且局部常值的。
\end{enumerate}
\end{definition}

\noindent
这看来是公认的定义。另有一种更便于推广到概形 \'etale 位点以外情形的
定义方式：从
$h_U$, $j_{U!}\mathbf{Z}/n\mathbf{Z}$ 以及 $j_{U!}\underline{\Lambda}$
出发，其中 $U$ 遍历 $X_\etale$ 的所有拟紧拟分离对象，再取
$\Sh(X_\etale)$、$\textit{Ab}(X_\etale)$ 以及
$\textit{Mod}(X_\etale, \underline{\Lambda})$ 中包含这些对象且对有限极限和
有限余极限封闭的最小满子范畴。由引理 \ref{etale-cohomology-lemma-constructible-abelian}
以及引理 \ref{etale-cohomology-lemma-category-constructible-sets},
\ref{etale-cohomology-lemma-category-constructible-abelian} 和
\ref{etale-cohomology-lemma-category-constructible-modules}
可知，当 $X$ 拟紧拟分离时，这种做法给出同一范畴；但一般而言，它们给出的
范畴并不相同。

\medskip\noindent
概形的一个由局部闭子概形组成的不交并分解 $U = \coprod U_i$，称为 $U$ 的一个
{\it 分划}（参见拓扑学，第 \ref{topology-section-stratifications} 节）。

\begin{lemma}
\label{etale-cohomology-lemma-constructible-quasi-compact-quasi-separated}
设 $X$ 为拟紧拟分离概形，$\mathcal{F}$ 为 $X_\etale$ 上的集合层。以下条件等价：
\begin{enumerate}
\item $\mathcal{F}$ 是可构造的；
\item 存在开覆盖 $X = \bigcup U_i$，使得 $\mathcal{F}|_{U_i}$ 是可构造的；
\item 存在由可构造局部闭子概形组成的分划 $X = \bigcup X_i$，使得
$\mathcal{F}|_{X_i}$ 是有限局部常值的。
\end{enumerate}
若 $\Lambda$ 为诺特环，则对阿贝尔层和 $\Lambda$-模层也有类似结论。
\end{lemma}

\begin{proof}
(1) 蕴含 (2) 是显然的。

\medskip\noindent
假设 (2) 成立。对每个 $x \in X$，可以找到某个 $i$ 以及一个仿射开邻域
$V_x \subset U_i$，后者是 $x$ 的邻域。因而可以找到有限仿射开覆盖
$X = \bigcup V_j$，使得对每个
$j$，都有一个由局部闭可构造子集组成的有限分解
$V_j = \coprod V_{j, k}$，并且 $\mathcal{F}|_{V_{j, k}}$ 是有限局部常值的。
由拓扑学，引理
\ref{topology-lemma-quasi-compact-open-immersion-constructible-image}
可知，每个 $V_{j, k}$ 作为 $X$ 的子集都是可构造的。由拓扑学，引理
\ref{topology-lemma-constructible-partition-refined-by-stratification}
可找到一个具有可构造局部闭层的有限分层 $X = \coprod X_l$，使得每个
$V_{j, k}$ 都是若干 $X_l$ 的并。因此 (3) 成立。

\medskip\noindent
假设 (3) 成立。设 $U \subset X$ 为仿射开集。则 $U \cap X_i$ 是 $U$ 的可构造
局部闭子集（例如由性质，引理 \ref{properties-lemma-locally-constructible}），
而 $U = \coprod U \cap X_i$ 是定义 \ref{etale-cohomology-definition-constructible} 中所述的
$U$ 的分划。因此 (1) 成立。
\end{proof}

\begin{lemma}
\label{etale-cohomology-lemma-constructible-constructible}
设 $X$ 为拟紧拟分离概形。设 $\mathcal{F}$ 为集合层、阿贝尔群层或
$\Lambda$-模层（其中 $\Lambda$ 为诺特环），定义在 $X_\etale$ 上。若存在可构造局部闭子概形
$T_i \subset X$，使得 (a) $X = \bigcup T_j$ 且 (b) $\mathcal{F}|_{T_j}$
是可构造的，则 $\mathcal{F}$ 是可构造的。
\end{lemma}

\begin{proof}
首先，可以假定该覆盖是有限的，因为 $X$ 在谱拓扑下是拟紧的
（拓扑学，引理 \ref{topology-lemma-constructible-hausdorff-quasi-compact}
以及性质，引理
\ref{properties-lemma-quasi-compact-quasi-separated-spectral}）。
注意到每个 $T_i$ 本身都是拟紧拟分离概形（因为它在 $X$ 中可构造；细节略去）。
因此，可以找到一个有限分划 $T_i = \coprod T_{i, j}$，其各部分是
$T_i$ 的局部闭可构造子集，并且 $\mathcal{F}|_{T_{i, j}}$ 是有限局部常值的
（引理 \ref{etale-cohomology-lemma-constructible-quasi-compact-quasi-separated}）。由拓扑学，
引理 \ref{topology-lemma-constructible-in-constructible} 可知，$T_{i, j}$ 是
$X$ 的可构造局部闭子概形。于是可以对 $X = \bigcup T_{i, j}$ 应用拓扑学，引理
\ref{topology-lemma-constructible-partition-refined-by-stratification}
以得到所需的 $X$ 的分划。
\end{proof}

\begin{lemma}
\label{etale-cohomology-lemma-constructible-local}
设 $X$ 为概形。集合层、阿贝尔群层或 $\Lambda$-模层（其中 $\Lambda$ 为诺特环）
的可构造性可以在 $X$ 上按 Zariski 局部检验。
\end{lemma}

\begin{proof}
该陈述的含义是：若 $X = \bigcup U_i$ 是开覆盖，且 $\mathcal{F}|_{U_i}$
是可构造的，则 $\mathcal{F}$ 是可构造的。若 $U \subset X$ 为仿射开集，则
$U = \bigcup U \cap U_i$，且 $\mathcal{F}|_{U \cap U_i}$ 是可构造的
（可构造层限制到开集仍可构造是显然的）。由引理
\ref{etale-cohomology-lemma-constructible-quasi-compact-quasi-separated} 可知，
$\mathcal{F}|_U$ 是可构造的，也就是说，$U$ 存在一个合适的分划。
\end{proof}

\begin{lemma}
\label{etale-cohomology-lemma-pullback-constructible}
设 $f : X \to Y$ 为概形态射。若 $\mathcal{F}$ 是可构造集合层、阿贝尔群层或
$\Lambda$-模层（其中 $\Lambda$ 为诺特环），定义在 $Y_\etale$ 上，则
$f^{-1}\mathcal{F}$ 在 $X_\etale$ 上也具有同样的性质。
\end{lemma}

\begin{proof}
由引理 \ref{etale-cohomology-lemma-constructible-local}，可归约到 $X$ 和 $Y$ 都是仿射的情形。
由引理 \ref{etale-cohomology-lemma-constructible-quasi-compact-quasi-separated}，只需找到一个
由可构造局部闭子概形组成的 $X$ 的有限分划，使得 $f^{-1}\mathcal{F}$ 在每一部分
上都是有限局部常值的。为此，只须拉回 $Y$ 的适配于 $\mathcal{F}$ 的分划，并应用
引理 \ref{etale-cohomology-lemma-pullback-locally-constant}。
\end{proof}

\begin{lemma}
\label{etale-cohomology-lemma-constructible-abelian}
设 $X$ 为概形。
\begin{enumerate}
\item 可构造集合层范畴在 $\Sh(X_\etale)$ 中对有限极限和有限余极限封闭。
\item 可构造阿贝尔层范畴是 $\textit{Ab}(X_\etale)$ 的弱 Serre 子范畴。
\item 设 $\Lambda$ 为诺特环。可构造 $\Lambda$-模层在 $X_\etale$ 上的范畴
是 $\textit{Mod}(X_\etale, \Lambda)$ 的弱 Serre 子范畴。
\end{enumerate}
\end{lemma}

\begin{proof}
证明 (3)。我们将使用同调，引理
\ref{homology-lemma-characterize-weak-serre-subcategory} 的判据。假设
$\varphi : \mathcal{F} \to \mathcal{G}$ 是可构造 $\Lambda$-模层之间的态射。
需要证明 $\mathcal{K} = \Ker(\varphi)$ 和 $\mathcal{Q} = \Coker(\varphi)$
都是可构造的。类似地，假设
$0 \to \mathcal{F} \to \mathcal{E} \to \mathcal{G} \to 0$
是 $\Lambda$-模层的短正合列，其中 $\mathcal{F}$、$\mathcal{G}$ 可构造。
需要证明 $\mathcal{E}$ 可构造。在这两种情形中，都可以用某个仿射开覆盖的成员
替换 $X$，因而可以假定 $X$ 是仿射的。进而，可以用 $X$ 的一个有限分划的各个
成员替换 $X$；该分划由可构造局部闭子概形组成，并且 $\mathcal{F}$ 和
$\mathcal{G}$ 在其上都是有限型且局部常值的。于是应用引理
\ref{etale-cohomology-lemma-kernel-finite-locally-constant} 即得结论。

\medskip\noindent
(1) 和 (2) 的证明非常类似，故略去。
\end{proof}

\begin{lemma}
\label{etale-cohomology-lemma-support-constructible}
设 $X$ 为拟紧拟分离概形。
\begin{enumerate}
\item 设 $\mathcal{F} \to \mathcal{G}$ 为 $X_\etale$ 上可构造集合层之间的态射。
则满足如下条件的点 $x \in X$ 所组成的集合，即
$\mathcal{F}_{\overline{x}} \to \mathcal{G}_{\overline{x}}$ 为满射、分别地为单射、
分别地同构于某个给定集合映射的点集，在 $X$ 中可构造。
\item 设 $\mathcal{F}$ 为 $X_\etale$ 上的可构造阿贝尔层。则 $\mathcal{F}$ 的
支集可构造。
\item 设 $\Lambda$ 为诺特环。设 $\mathcal{F}$ 为可构造 $\Lambda$-模层，
定义在 $X_\etale$ 上。则 $\mathcal{F}$ 的支集可构造。
\end{enumerate}
\end{lemma}

\begin{proof}
证明 (1)。设 $X = \coprod X_i$ 是由局部闭可构造子概形组成的 $X$ 的分划，
使得 $\mathcal{F}$ 和 $\mathcal{G}$ 在各部分上都是有限局部常值的（对
$\mathcal{F}$ 和 $\mathcal{G}$ 分别应用引理
\ref{etale-cohomology-lemma-constructible-quasi-compact-quasi-separated}，再取一个共同加细）。
然后，对该态射在每一部分上的限制应用引理
\ref{etale-cohomology-lemma-morphism-locally-constant}。

\medskip\noindent
(2) 和 (3) 的证明略去。
\end{proof}

\noindent
以下引理稍后会非常有用。粗略地说，它表明可构造层范畴具有某种弱
``诺特'' 性质。

\begin{lemma}
\label{etale-cohomology-lemma-colimit-constructible}
设 $X$ 为拟紧拟分离概形。设
$\mathcal{F} = \colim_{i \in I} \mathcal{F}_i$ 是集合层、阿贝尔层或模层的
滤过余极限。
\begin{enumerate}
\item 若 $\mathcal{F}$ 和 $\mathcal{F}_i$ 都是可构造集合层，则 ind-对象
$\mathcal{F}_i$ 本质上为常值，取值为 $\mathcal{F}$。
\item 若 $\mathcal{F}$ 和 $\mathcal{F}_i$ 都是可构造阿贝尔群层，则 ind-对象
$\mathcal{F}_i$ 本质上为常值，取值为 $\mathcal{F}$。
\item 设 $\Lambda$ 为诺特环。若 $\mathcal{F}$ 和 $\mathcal{F}_i$ 都是可构造
$\Lambda$-模层，则 ind-对象 $\mathcal{F}_i$ 本质上为常值，取值为
$\mathcal{F}$。
\end{enumerate}
\end{lemma}

\begin{proof}
证明 (1)。下文不再另行说明地使用如下事实：可构造层的有限极限和有限余极限
仍然可构造（引理 \ref{etale-cohomology-lemma-kernel-finite-locally-constant}）。对每个 $i$，
令 $T_i \subset X$ 为所有点 $x \in X$ 所组成的集合，其中
$\mathcal{F}_{i, \overline{x}} \to \mathcal{F}_{\overline{x}}$ 不是满射。
由于 $\mathcal{F}_i$ 和 $\mathcal{F}$ 可构造，
$T_i$ 是 $X$ 的可构造子集（引理 \ref{etale-cohomology-lemma-support-constructible}）。因为
$\mathcal{F}$ 的茎都是有限的，且
$\mathcal{F} = \colim_{i \in I} \mathcal{F}_i$，所以对每个 $x \in X$，
$x \not \in T_i$ 对充分大的 $i$ 成立。由性质，引理
\ref{properties-lemma-quasi-compact-quasi-separated-spectral}
可知 $X$ 是谱空间；由拓扑学，引理
\ref{topology-lemma-constructible-hausdorff-quasi-compact}，$X$ 上的可构造拓扑
是拟紧的。因此，$T_i = \emptyset$ 对充分大的 $i$ 成立。从而，
$\mathcal{F}_i \to \mathcal{F}$ 对充分大的 $i$ 是满射。现在假设
$\mathcal{F}_i \to \mathcal{F}$ 对所有 $i$ 都是满射。取 $i \in I$。对 $i' \geq i$，
记 $S_{i'} \subset X$ 为如下点 $x$ 组成的集合：
$\Im(\mathcal{F}_{i, \overline{x}} \to \mathcal{F}_{\overline{x}})$
中的元素个数不等于
$\Im(\mathcal{F}_{i, \overline{x}} \to \mathcal{F}_{i', \overline{x}})$
中的元素个数。由于 $\mathcal{F}_i$、$\mathcal{F}_{i'}$ 和 $\mathcal{F}$
可构造，$S_{i'}$ 是 $X$ 的可构造子集（细节略去；提示：使用引理
\ref{etale-cohomology-lemma-support-constructible}）。由于 $\mathcal{F}_i$ 和 $\mathcal{F}$
的茎都是有限的，且 $\mathcal{F} = \colim_{i' \geq i} \mathcal{F}_{i'}$，
所以对每个 $x \in X$，$x \not \in S_{i'}$ 对充分大的 $i'$ 成立。
由与上面相同的论证，可找到充分大的 $i'$，使得 $S_{i'} = \emptyset$。
因此，$\mathcal{F}_i \to \mathcal{F}_{i'}$ 如所需地通过 $\mathcal{F}$ 分解。

\medskip\noindent
证明 (2)。注意到可构造阿贝尔层也是可构造集合层。因此 (2) 由 (1) 得出。

\medskip\noindent
证明 (3)。下文不再另行说明地使用如下事实：可构造 $\Lambda$-模层范畴是阿贝尔
范畴（引理 \ref{etale-cohomology-lemma-kernel-finite-locally-constant}）。对每个 $i$，令
$\mathcal{Q}_i$ 为态射 $\mathcal{F}_i \to \mathcal{F}$ 的余核。令 $T_i$ 为
$\mathcal{Q}_i$ 的支集；它是 $X$ 的可构造子集，因为 $\mathcal{Q}_i$ 可构造
（引理 \ref{etale-cohomology-lemma-support-constructible}）。由于 $\mathcal{F}$ 的茎都是有限
$\Lambda$-模，且 $\mathcal{F} = \colim_{i \in I} \mathcal{F}_i$，所以对每个
$x \in X$，$x \not \in T_i$ 对充分大的 $i$ 成立。由性质，引理
\ref{properties-lemma-quasi-compact-quasi-separated-spectral}
可知 $X$ 是谱空间；由拓扑学，引理
\ref{topology-lemma-constructible-hausdorff-quasi-compact}，$X$ 上的可构造拓扑
是拟紧的。因此，$T_i = \emptyset$ 对充分大的 $i$ 成立。这就证明了第一个断言。
对于第二个断言，现在假设 $\mathcal{F}_i \to \mathcal{F}$ 对所有 $i$ 都是
满射。取 $i \in I$。对 $i' \geq i$，记 $\mathcal{K}_{i'}$ 为
$\Ker(\mathcal{F}_i \to \mathcal{F})$ 在 $\mathcal{F}_{i'}$ 中的像。令
$S_{i'}$ 为 $\mathcal{K}_{i'}$ 的支集；它是 $X$ 的可构造子集，因为
$\mathcal{K}_{i'}$ 可构造。由于 $\Ker(\mathcal{F}_i \to \mathcal{F})$
的茎都是有限 $\Lambda$-模，且
$\mathcal{F} = \colim_{i' \geq i} \mathcal{F}_{i'}$，所以对每个 $x \in X$，
$x \not \in S_{i'}$ 对充分大的 $i'$ 成立。由与上面相同的论证，可找到充分大的
$i'$，使得 $S_{i'} = \emptyset$。因此，
$\mathcal{F}_i \to \mathcal{F}_{i'}$ 如所需地通过 $\mathcal{F}$ 分解。
\end{proof}

\begin{lemma}
\label{etale-cohomology-lemma-tensor-product-constructible}
设 $X$ 为概形，$\Lambda$ 为诺特环。两个可构造 $\Lambda$-模层在
$X_\etale$ 上的张量积仍是可构造 $\Lambda$-模层。
\end{lemma}

\begin{proof}
该问题立即归约到 $X$ 为仿射的情形。由于 $X$ 的任意两个由可构造局部闭层组成的
分划都有同类的共同加细，并且拉回与张量积可交换，所以可归约到引理
\ref{etale-cohomology-lemma-tensor-product-locally-constant}。
\end{proof}

\begin{lemma}
\label{etale-cohomology-lemma-tensor-constructible}
设 $\Lambda \to \Lambda'$ 为诺特环同态。设 $X$ 为概形，$\mathcal{F}$ 为
可构造 $\Lambda$-模层，定义在 $X_\etale$ 上。则
$\mathcal{F} \otimes_{\underline{\Lambda}} \underline{\Lambda'}$
是可构造 $\Lambda'$-模层。
\end{lemma}

\begin{proof}
略。提示：在仿射局部，可以使用同一个分层。
\end{proof}










\section{关于态射的辅助引理}
\label{etale-cohomology-section-stratify-morphisms}

\noindent
下面给出一些在证明可构造层的函子性性质时有用的引理。

\begin{lemma}
\label{etale-cohomology-lemma-etale-stratified-finite}
设 $U \to X$ 为拟紧拟分离概形之间的 \'etale 态射（例如诺特概形之间的
\'etale 态射）。则存在一个由可构造局部闭子概形组成的分划
$X = \coprod_i X_i$，使得 $X_i \times_X U \to X_i$ 对所有 $i$ 都是有限
\'etale 态射。
\end{lemma}

\begin{proof}
若 $U \to X$ 是分离的，则这就是态射进阶，引理
\ref{more-morphisms-lemma-stratify-flat-fp-qf}。一般而言，可以假定 $X$ 是仿射的。
选择一个有限仿射开覆盖 $U = \bigcup U_j$。对所有态射 $U_j \to X$ 和
$U_j \cap U_{j'} \to X$ 应用前一种情形，并选取所得各分划的共同加细
$X = \coprod X_i$。进一步加细该分划后，还可以假定 $X_i$ 是仿射的。
固定 $i$，并令 $V = U \times_X X_i$。态射
$V_j = U_j \times_X X_i \to X_i$ 以及
$V_{jj'} = (U_j \cap U_{j'}) \times_X X_i \to X_i$ 都是有限 \'etale 的。
因此，$V_j$ 和 $V_{jj'}$ 都是仿射概形，且 $V_{jj'} \subset V_j$ 既闭又开
（因为 $V_{jj'} \to X_i$ 是固有的，故可应用态射，引理
\ref{morphisms-lemma-image-proper-scheme-closed}）。于是 $V = \bigcup V_j$ 是分离的，
因为 $\mathcal{O}(V_j) \to \mathcal{O}(V_{jj'})$ 是满射，见概形，引理
\ref{schemes-lemma-characterize-separated}。所以前一种情形适用于 $V \to X_i$，
而且如果需要，可以进一步加细分划（事实上不需要，但我们不必用到这一点）。
\end{proof}

\noindent
在诺特情形中，可以用诺特归纳法和下面这个有趣的引理来证明前一个引理。

\begin{lemma}
\label{etale-cohomology-lemma-generically-finite}
设 $f: X \to Y$ 为拟紧、拟分离且局部有限型的概形态射。若 $\eta$ 是 $Y$
的某个不可约分支的泛点，且 $f^{-1}(\eta)$ 是有限的，则存在开集
$V \subset Y$，它包含 $\eta$，并且 $f^{-1}(V) \to V$ 是有限态射。
\end{lemma}

\begin{proof}
这就是态射，引理 \ref{morphisms-lemma-generically-finite}。
\end{proof}

\noindent
下面引理的陈述还可以略作加强。

\begin{lemma}
\label{etale-cohomology-lemma-decompose-quasi-finite-morphism}
设 $f : Y \to X$ 为仿射概形之间的拟有限且有限表现态射。
\begin{enumerate}
\item 存在仿射概形的满态射 $X' \to X$ 和闭子概形
$Z' \subset Y' = X' \times_X Y$，使得
\begin{enumerate}
\item $Z' \subset Y'$ 是一个加厚；
\item $Z' \to X'$ 是有限 \'etale 态射。
\end{enumerate}
\item 存在一个由局部闭、可构造、仿射的层组成的有限分划
$X = \coprod X_i$，以及有限局部自由满态射 $X'_i \to X_i$，使得
$Y'_i = X'_i \times_X Y \to X'_i$ 的既约化同构于
$\coprod_{j = 1}^{n_i} (X'_i)_{red} \to (X'_i)_{red}$，其中 $n_i$ 为某个整数。
\end{enumerate}
\end{lemma}

\begin{proof}
令 $X' = \coprod X'_i$，可见 (2) 蕴含 (1)。写成 $X = \Spec(A)$ 和
$Y = \Spec(B)$。把 $A$ 写成有限型 $\mathbf{Z}$-代数 $A_i$ 的滤过余极限。
由于 $B$ 是有限表现 $A$-代数，可知存在 $0 \in I$ 以及有限型环同态
$A_0 \to B_0$，使得 $B = \colim B_i$，其中
$B_i = A_i \otimes_{A_0} B_0$，见代数，引理
\ref{algebra-lemma-colimit-category-fp-algebras}。当 $i$ 充分大时，
$A_i \to B_i$ 是拟有限的，见极限，引理
\ref{limits-lemma-descend-quasi-finite}。因此，可归约到 $\mathbf{Z}$ 上有限型代数
的情形，特别地，可归约到诺特情形。（细节略去。）

\medskip\noindent
假设 $X$ 和 $Y$ 都是诺特的。在此情形中，$X$ 的任意局部闭子集都是可构造的。
由引理 \ref{etale-cohomology-lemma-generically-finite} 和诺特归纳法可知，存在有限分划
$X = \coprod X_i$，它是 $X$ 的由局部闭层组成的分划，并使得
$Y \times_X X_i \to X_i$ 是有限的。可以加细该分划，使其各层仿射。因此，
把 $X$ 替换为 $X' = \coprod X_i$ 后，
可以假定 $Y \to X$ 是有限的。

\medskip\noindent
假设 $X$ 和 $Y$ 都是诺特的，且 $Y \to X$ 有限。假设在沿某个满、平坦、
拟有限态射 $U \to X$ 作基变换后可以证明 (2)。于是有分划
$U = \coprod U_i$ 和有限局部自由态射 $U'_i \to U_i$，使得
$U'_i \times_X Y \to U'_i$ 同构于
$\coprod_{j = 1}^{n_i} (U'_i)_{red} \to (U'_i)_{red}$，其中 $n_i$ 为某个整数。
由前一段的论证，可以找到一个具有局部闭仿射层的分划 $X = \coprod X_j$，使得
$X_j \times_X U_i \to X_j$ 对所有 $i, j$ 都是有限的。由态射，引理
\ref{morphisms-lemma-finite-flat}，每个 $X_j \times_X U_i \to X_j$ 都是有限
局部自由的。因此，$X_j \times_X U'_i \to X_j$ 是有限局部自由的（态射，引理
\ref{morphisms-lemma-composition-finite-locally-free}）。由此，
$X = \coprod X_j$ 和 $X_j' = \coprod_i X_j \times_X U'_i$ 给出了
$Y \to X$ 的一个解。因此，只需在作满、平坦、拟有限基变换后证明该结论
（在诺特情形中）。

\medskip\noindent
应用态射，引理 \ref{morphisms-lemma-massage-finite}，可知可以假定 $Y$ 是某个
仿射概形 $Z$ 的闭子概形，而后者在集合意义下是闭子概形的有限并
$Z = \bigcup_{i \in I} Z_i$，其中这些闭子概形同构地映到 $X$。在这种情形下，
我们将找到一个可行的、具有仿射局部闭层的有限分划 $X = \coprod X_j$
（换言之，$X'_j = X_j$）。令 $T_i = Y \cap Z_i$。这是 $X$ 的闭子概形。
由于 $X$ 是诺特的，可以找到一个由仿射局部闭子概形组成的有限分划
$X = \coprod X_j$，使得每个 $X_j \times_X T_i$ 在集合意义下都是若干层
$X_j \times_X Z_i$ 的并。把 $X$ 替换为 $X_j$，可知可以假定
$I = I_1 \amalg I_2$，其中 $Z_i \subset Y$ 对 $i \in I_1$ 成立，而
$Z_i \cap Y = \emptyset$ 对 $i \in I_2$ 成立。把 $Z$ 替换为
$\bigcup_{i \in I_1} Z_i$，可知可以假定 $Y = Z$。
最后，可以再次用上述分划的各成员替换 $X$，使得对每个 $i, i' \subset I$，交
$Z_i \cap Z_{i'}$ 或者为空，或者在集合意义下同时等于 $Z_i$ 和 $Z_{i'}$。
这显然意味着 $Y$ 在集合意义下等于诸 $Z_i$ 的不交并，而这正是所要证明的。
\end{proof}








\section{关于可构造层的进一步结果}
\label{etale-cohomology-section-more-constructible}

\noindent
设 $\Lambda$ 为诺特环，$X$ 为概形。我们经常把 $X_\etale$ 看成环化位点，
其环层为 $\underline{\Lambda}$。在阿贝尔层的情形中，我们经常取
$\Lambda = \mathbf{Z}/n\mathbf{Z}$，其中 $n$ 为某个合适的整数。

\begin{lemma}
\label{etale-cohomology-lemma-jshriek-constructible}
设 $j : U \to X$ 为拟紧拟分离概形之间的 \'etale 态射。
\begin{enumerate}
\item 层 $h_U$ 是可构造集合层。
\item 层 $j_!\underline{M}$ 对有限阿贝尔群 $M$ 是可构造阿贝尔层。
\item 若 $\Lambda$ 为诺特环，且 $M$ 是有限 $\Lambda$-模，则
$j_!\underline{M}$ 是可构造 $\Lambda$-模层，定义在 $X_\etale$ 上。
\end{enumerate}
\end{lemma}

\begin{proof}
由引理 \ref{etale-cohomology-lemma-etale-stratified-finite}，存在分划 $\coprod_i X_i$，使得
$\pi_i : j^{-1}(X_i) \to X_i$ 是有限 \'etale 态射。$h_U$ 在 $X_i$ 上的限制是
$h_{j^{-1}(X_i)}$；由引理 \ref{etale-cohomology-lemma-characterize-finite-locally-constant}，
它是有限局部常值的。对于 (2) 和 (3)，注意到
$$
j_!(\underline{M})|_{X_i} =
\pi_{i!}(\underline{M}) =
\pi_{i*}(\underline{M})
$$
由引理 \ref{etale-cohomology-lemma-shriek-base-change} 和
\ref{etale-cohomology-lemma-shriek-equals-star-finite-etale} 成立。因此，只需对有限 \'etale 态射
$\pi : Y \to X$ 证明该引理。这就是引理
\ref{etale-cohomology-lemma-pushforward-locally-constant}。
\end{proof}

\begin{lemma}
\label{etale-cohomology-lemma-torsion-colimit-constructible}
设 $X$ 为拟紧拟分离概形。
\begin{enumerate}
\item 设 $\mathcal{F}$ 为 $X_\etale$ 上的集合层。则 $\mathcal{F}$ 是可构造
集合层的滤过余极限。
\item 设 $\mathcal{F}$ 为 $X_\etale$ 上的挠阿贝尔层。则 $\mathcal{F}$ 是
可构造阿贝尔层的滤过余极限。
\item 设 $\Lambda$ 为诺特环，$\mathcal{F}$ 为 $\Lambda$-模层，定义在
$X_\etale$ 上。则 $\mathcal{F}$ 是可构造 $\Lambda$-模层的滤过余极限。
\end{enumerate}
\end{lemma}

\begin{proof}
令 $\mathcal{B}$ 为 $X_\etale$ 的拟紧拟分离对象所组成的集合。由位点上的模，
引理 \ref{sites-modules-lemma-module-filtered-colimit-constructibles}，
任意集合层都是如下形式的层的滤过余极限：
$$
\text{Coequalizer}\left(
\xymatrix{
\coprod\nolimits_{j = 1, \ldots, m} h_{V_j}
\ar@<1ex>[r] \ar@<-1ex>[r] &
\coprod\nolimits_{i = 1, \ldots, n} h_{U_i}
}
\right)
$$
其中 $V_j$ 和 $U_i$ 是 $X_\etale$ 的拟紧拟分离对象。由引理
\ref{etale-cohomology-lemma-jshriek-constructible} 和 \ref{etale-cohomology-lemma-constructible-abelian}，
这些余等化子是可构造的。这就证明了 (1)。

\medskip\noindent
设 $\Lambda$ 为诺特环。由位点上的模，引理
\ref{sites-modules-lemma-module-filtered-colimit-constructibles}，
$\Lambda$-模层 $\mathcal{F}$ 是如下形式的模层的滤过余极限：
$$
\Coker\left(
\bigoplus\nolimits_{j = 1, \ldots, m} j_{V_j!}\underline{\Lambda}_{V_j}
\longrightarrow
\bigoplus\nolimits_{i = 1, \ldots, n} j_{U_i!}\underline{\Lambda}_{U_i}
\right)
$$
其中 $V_j$ 和 $U_i$ 是 $X_\etale$ 的拟紧拟分离对象。由引理
\ref{etale-cohomology-lemma-jshriek-constructible} 和 \ref{etale-cohomology-lemma-constructible-abelian}，
这些余核是可构造的。这就证明了 (3)。

\medskip\noindent
证明 (2)。首先写成 $\mathcal{F} = \bigcup \mathcal{F}[n]$，其中
$\mathcal{F}[n]$ 是 $n$-挠子层。于是可以把 $\mathcal{F}[n]$ 看成
$\mathbf{Z}/n\mathbf{Z}$-模层，并应用 (3)。
\end{proof}

\begin{lemma}
\label{etale-cohomology-lemma-check-constructible}
设 $f : X \to Y$ 为拟紧拟分离概形之间的满态射。
\begin{enumerate}
\item 设 $\mathcal{F}$ 为 $Y_\etale$ 上的集合层。则 $\mathcal{F}$ 可构造，
当且仅当 $f^{-1}\mathcal{F}$ 可构造。
\item 设 $\mathcal{F}$ 为 $Y_\etale$ 上的阿贝尔层。则 $\mathcal{F}$ 可构造，
当且仅当 $f^{-1}\mathcal{F}$ 可构造。
\item 设 $\Lambda$ 为诺特环。设 $\mathcal{F}$ 为 $\Lambda$-模层，定义在
$Y_\etale$ 上。则 $\mathcal{F}$ 可构造，当且仅当 $f^{-1}\mathcal{F}$ 可构造。
\end{enumerate}
\end{lemma}

\begin{proof}
一个方向由引理 \ref{etale-cohomology-lemma-pullback-constructible} 得出。反过来，假设
$f^{-1}\mathcal{F}$ 可构造。使用引理
\ref{etale-cohomology-lemma-torsion-colimit-constructible}，把
$\mathcal{F} = \colim \mathcal{F}_i$ 写成可构造层（集合层、阿贝尔群层或模层）的
滤过余极限。由于 $f^{-1}$ 是左伴随，它与余极限可交换（范畴，引理
\ref{categories-lemma-adjoint-exact}），从而
$f^{-1}\mathcal{F} = \colim f^{-1}\mathcal{F}_i$.
由引理 \ref{etale-cohomology-lemma-colimit-constructible} 可知，
$f^{-1}\mathcal{F}_i \to f^{-1}\mathcal{F}$
对所有充分大的 $i$ 都是满射。由于 $f$ 是满射，通过使用引理
\ref{etale-cohomology-lemma-stalk-pullback} 和定理 \ref{etale-cohomology-theorem-exactness-stalks} 考察茎可得，
$\mathcal{F}_i \to \mathcal{F}$ 对所有充分大的 $i$ 都是满射。因此，
$\mathcal{F}$ 是某个可构造层 $\mathcal{G}$ 的商。再把该论证应用于
$\mathcal{G} \times_\mathcal{F} \mathcal{G}$ 或
$\mathcal{G} \to \mathcal{F}$ 的核，并使用 $f^{-1}$ 的正合性以及可构造层
（集合层、阿贝尔群层或模层）范畴在下列范畴内部对有限（余）极限或（余）核封闭：
$\Sh(Y_\etale)$, $\Sh(X_\etale)$, $\textit{Ab}(Y_\etale)$,
$\textit{Ab}(X_\etale)$, $\textit{Mod}(Y_\etale, \Lambda)$ 以及
$\textit{Mod}(X_\etale, \Lambda)$，见引理
\ref{etale-cohomology-lemma-constructible-abelian}，即可得出结论。
\end{proof}

\begin{lemma}
\label{etale-cohomology-lemma-pushforward-constructible}
设 $f : X \to Y$ 为概形的有限 \'etale 态射，$\Lambda$ 为诺特环。若
$\mathcal{F}$ 是可构造集合层、可构造阿贝尔群层或可构造 $\Lambda$-模层，
定义在 $X_\etale$ 上，则 $f_*\mathcal{F}$ 在 $Y_\etale$ 上也具有同样的性质。
\end{lemma}

\begin{proof}
由引理 \ref{etale-cohomology-lemma-constructible-local}，只需在 $Y$ 上按 Zariski 局部检验；
由引理 \ref{etale-cohomology-lemma-check-constructible}，可以用某个 \'etale 覆盖替换 $Y$
（$f_*$ 的构造与 \'etale 局部化可交换）。有限 \'etale 态射在 \'etale 局部
同构于若干同构的不交并，见 \'Etale 态射，引理
\ref{etale-lemma-finite-etale-etale-local}。因此，在集合层的情形中，该引理是说：
若 $\mathcal{F}_i$（$i = 1, \ldots, n$）是可构造集合层，则
$\prod_{i = 1, \ldots, n} \mathcal{F}_i$ 也是可构造的。这是显然的。
对阿贝尔群层和模层也类似。
\end{proof}

\begin{lemma}
\label{etale-cohomology-lemma-category-constructible-sets}
设 $X$ 为拟紧拟分离概形。可构造集合层范畴正是 $\Sh(X_\etale)$ 的如下满子范畴：
其对象为可写成下列余等化子的层 $\mathcal{F}$：
$$
\xymatrix{
\mathcal{F}_1
\ar@<1ex>[r] \ar@<-1ex>[r]
&
\mathcal{F}_0 \ar[r]
&
\mathcal{F}}
$$
其中 $\mathcal{F}_i$（$i = 0, 1$）是形如 $h_U$ 的层的有限余积，而 $U$ 是
$X_\etale$ 的拟紧拟分离对象。
\end{lemma}

\begin{proof}
在引理 \ref{etale-cohomology-lemma-torsion-colimit-constructible} 的证明中，我们已经看到这种形式的
层是可构造的。反过来，假设对每个可构造集合层 $\mathcal{F}$，都能找到一个满射
$\mathcal{F}_0 \to \mathcal{F}$，其中 $\mathcal{F}_0$ 如引理中所述。那么就能找到
所需的满射
$\mathcal{F}_1 \to \mathcal{F}_0 \times_\mathcal{F} \mathcal{F}_0$
，因为由引理 \ref{etale-cohomology-lemma-constructible-abelian}，后一个层是可构造的。

\medskip\noindent
由拓扑学，引理
\ref{topology-lemma-constructible-partition-refined-by-stratification}
，可以选择有限分层 $X = \coprod_{i \in I} X_i$，使得 $\mathcal{F}$ 在每一层
上都是有限局部常值的。我们对 $I$ 的基数作归纳来证明该结论。设 $i \in I$
是 $I$ 的偏序中的极小元。则 $X_i \subset X$ 是闭的。由归纳假设，存在有限多个
拟紧拟分离对象 $U_\alpha$，它们属于 $(X \setminus X_i)_\etale$，并且存在满射
$\coprod h_{U_\alpha} \to \mathcal{F}|_{X \setminus X_i}$。它们确定一个态射
$$
\coprod h_{U_\alpha} \to \mathcal{F}
$$
，它限制到 $X \setminus X_i$ 后是满射。由引理
\ref{etale-cohomology-lemma-characterize-finite-locally-constant}，可知
$\mathcal{F}|_{X_i} = h_V$，其中某个概形 $V$ 在 $X_i$ 上有限 \'etale。设
$\overline{v}$ 为 $V$ 的位于 $\overline{x} \in X_i$ 上方的几何点。可以把
$\overline{v}$ 看成茎 $\mathcal{F}_{\overline{x}} = V_{\overline{x}}$ 的元素。
因此，可以找到一个 \'etale 邻域 $(U, \overline{u})$，其中心为 $\overline{x}$，以及截面
$s \in \mathcal{F}(U)$，其在 $\overline{x}$ 处的茎给出 $\overline{v}$。把 $s$
看成态射 $s : h_U \to \mathcal{F}$，并限制到 $X_i$，得到态射
$s|_{X_i} : U \times_X X_i \to V$；它在 $X_i$ 上，并把 $\overline{u}$ 映到
$\overline{v}$。由于 $V$ 是拟紧的（它在闭子概形 $X_i$ 上有限，而后者属于拟紧概形
$X$），有限多个 $s^{(1)}, \ldots, s^{(m)}$（它们是 $\mathcal{F}$ 在
$U^{(1)}, \ldots, U^{(m)}$ 上的这些截面）
上，将确定一个联合满态射
$$
\coprod s^{(j)}|_{X_i} : \coprod U^{(j)} \times_X X_i \longrightarrow V
$$
于是得到所需的满射
$$
\coprod h_{U_\alpha} \amalg \coprod h_{U^{(j)}} \to \mathcal{F}
$$
。
\end{proof}

\begin{lemma}
\label{etale-cohomology-lemma-category-constructible-modules}
设 $X$ 为拟紧拟分离概形，$\Lambda$ 为诺特环。可构造 $\Lambda$-模层范畴
正是由如下形式的模层组成的范畴：
$$
\Coker\left(
\bigoplus\nolimits_{j = 1, \ldots, m} j_{V_j!}\underline{\Lambda}_{V_j}
\longrightarrow
\bigoplus\nolimits_{i = 1, \ldots, n} j_{U_i!}\underline{\Lambda}_{U_i}
\right)
$$
其中 $V_j$ 和 $U_i$ 是 $X_\etale$ 的拟紧拟分离对象。事实上，甚至可以假定
$U_i$ 和 $V_j$ 都是仿射的。
\end{lemma}

\begin{proof}
在引理 \ref{etale-cohomology-lemma-torsion-colimit-constructible} 的证明中，我们已经看到这种形式的
模层是可构造的。由于可构造模层范畴是阿贝尔范畴（引理
\ref{etale-cohomology-lemma-constructible-abelian}），只需证明：给定可构造模层 $\mathcal{F}$，
存在满射
$$
\bigoplus\nolimits_{i = 1, \ldots, n} j_{U_i!}\underline{\Lambda}_{U_i}
\longrightarrow \mathcal{F}
$$
，其中 $U_i$ 是 $X_\etale$ 的某些仿射对象。由位点上的模，引理
\ref{sites-modules-lemma-module-filtered-colimit-constructibles}
，存在满射
$$
\Psi :
\bigoplus\nolimits_{i \in I} j_{U_i!}\underline{\Lambda}_{U_i}
\longrightarrow
\mathcal{F}
$$
，其中 $U_i$ 是仿射的，且直和取遍一个可能无限的指标集 $I$。对每个有限子集
$I' \subset I$，令
$$
T_{I'} = \text{Supp}(\Coker(
\bigoplus\nolimits_{i \in I'} j_{U_i!}\underline{\Lambda}_{U_i}
\longrightarrow \mathcal{F}))
$$
由可构造层的定义，集合 $T_{I'}$ 是 $X$ 的可构造子集。我们要证明
$T_{I'} = \emptyset$ 对某个 $I'$ 成立。由于每个茎 $\mathcal{F}_{\overline{x}}$ 都是有限型
$\Lambda$-模，且 $\Psi$ 是满射，所以对每个 $x \in X$，存在某个 $I'$，使得
$x \not \in T_{I'}$。换言之，有
$\emptyset = \bigcap_{I' \subset I\text{ 有限}} T_{I'}$。由于
$X$ 由性质，引理
\ref{properties-lemma-quasi-compact-quasi-separated-spectral}
是谱空间，而由拓扑学，引理
\ref{topology-lemma-constructible-hausdorff-quasi-compact}，$X$ 上的可构造拓扑
是拟紧的。因此，$T_{I'} = \emptyset$ 对某个有限的 $I' \subset I$ 成立，
如所需。
\end{proof}

\begin{lemma}
\label{etale-cohomology-lemma-category-constructible-abelian}
设 $X$ 为拟紧拟分离概形。可构造阿贝尔层范畴正是由如下形式的阿贝尔层组成的范畴：
$$
\Coker\left(
\bigoplus\nolimits_{j = 1, \ldots, m}
j_{V_j!}\underline{\mathbf{Z}/m_j\mathbf{Z}}_{V_j}
\longrightarrow
\bigoplus\nolimits_{i = 1, \ldots, n}
j_{U_i!}\underline{\mathbf{Z}/n_i\mathbf{Z}}_{U_i}
\right)
$$
其中 $V_j$ 和 $U_i$ 是 $X_\etale$ 的拟紧拟分离对象，且 $m_j$、$n_i$ 为正整数。
事实上，甚至可以假定 $U_i$ 和 $V_j$ 都是仿射的。
\end{lemma}

\begin{proof}
这由对 $\Lambda = \mathbf{Z}/n\mathbf{Z}$ 应用引理
\ref{etale-cohomology-lemma-category-constructible-modules}，以及如下事实得出：由于 $X$ 拟紧，
每个可构造阿贝尔层都被某个正整数 $n$ 零化（细节略去）。
\end{proof}

\begin{lemma}
\label{etale-cohomology-lemma-constructible-is-compact}
设 $X$ 为拟紧拟分离概形，$\Lambda$ 为诺特环。设 $\mathcal{F}$ 为可构造集合层、
阿贝尔群层或 $\Lambda$-模层，定义在 $X_\etale$ 上。设
$\mathcal{G} = \colim \mathcal{G}_i$ 为集合层、阿贝尔群层或 $\Lambda$-模层的
滤过余极限。则
$$
\Mor(\mathcal{F}, \mathcal{G}) = \colim \Mor(\mathcal{F}, \mathcal{G}_i)
$$
在集合层、阿贝尔群层或 $\Lambda$-模层的范畴中成立，该范畴定义在 $X_\etale$ 上。
\end{lemma}

\begin{proof}
先考虑集合层的情形。由引理 \ref{etale-cohomology-lemma-category-constructible-sets}，只需对
$h_U$ 证明该引理，其中 $U$ 是 $X_\etale$ 的拟紧拟分离对象。回忆
$\Mor(h_U, \mathcal{G}) = \mathcal{G}(U)$。因此，结论由位点，引理
\ref{sites-lemma-directed-colimits-sections} 得出。

\medskip\noindent
在阿贝尔层或模层的情形中，使用引理
\ref{etale-cohomology-lemma-category-constructible-abelian} 和
\ref{etale-cohomology-lemma-category-constructible-modules}，以同样方式得到结论。对于阿贝尔层，
还要补充注意到
$\Mor(j_{U!}\underline{\mathbf{Z}/n\mathbf{Z}}, \mathcal{G})$
等于所有 $n$-挠元素所组成的集合，这些元素属于 $\mathcal{G}(U)$。
\end{proof}

\begin{lemma}
\label{etale-cohomology-lemma-finite-pushforward-constructible}
设 $f : X \to Y$ 为概形的有限且有限表现态射，$\Lambda$ 为诺特环。若
$\mathcal{F}$ 是可构造集合层、阿贝尔群层或 $\Lambda$-模层，定义在 $X_\etale$ 上，
则 $f_*\mathcal{F}$ 也是可构造的。
\end{lemma}

\begin{proof}
由引理 \ref{etale-cohomology-lemma-constructible-local}，只需在 $X$ 和 $Y$ 都是仿射时证明该结论。
由引理 \ref{etale-cohomology-lemma-finite-pushforward-commutes-with-base-change} 和
\ref{etale-cohomology-lemma-check-constructible}，可以基变换到任意满射到 $X$ 的仿射概形。由引理
\ref{etale-cohomology-lemma-decompose-quasi-finite-morphism}，这归约到有限 \'etale 态射的情形
（因为加厚给出 \'etale 拓扑斯、甚至小 \'etale 位点之间的等价，见定理
\ref{etale-cohomology-theorem-topological-invariance}）。有限 \'etale 情形就是引理
\ref{etale-cohomology-lemma-pushforward-constructible}。
\end{proof}

\begin{lemma}
\label{etale-cohomology-lemma-category-is-colimit}
设 $X = \lim_{i \in I} X_i$ 为一个具有仿射转移态射的概形有向系的极限。
假设 $X_i$ 对所有 $i \in I$ 都拟紧拟分离。
\begin{enumerate}
\item $X_\etale$ 上的可构造集合层范畴是各 $(X_i)_\etale$ 上可构造集合层范畴
的余极限。
\item $X_\etale$ 上的可构造阿贝尔层范畴是各 $(X_i)_\etale$ 上可构造阿贝尔层
范畴的余极限。
\item 设 $\Lambda$ 为诺特环。可构造 $\Lambda$-模层在 $X_\etale$ 上的范畴是
可构造 $\Lambda$-模层在各 $(X_i)_\etale$ 上的范畴的余极限。
\end{enumerate}
\end{lemma}

\begin{proof}
证明 (1)。用 $f_i : X \to X_i$ 表示投影态射。证明分为三部分，分别对应于
``忠实''、``全忠实'' 和 ``本质满射''。

\medskip\noindent
忠实性。选取 $0 \in I$，并设 $\mathcal{F}_0$、$\mathcal{G}_0$ 为 $X_0$ 上的
可构造层。假设 $a, b : \mathcal{F}_0 \to \mathcal{G}_0$ 是满足
$f_0^{-1}a = f_0^{-1}b$ 的态射。令 $E \subset X_0$ 为如下点
$x \in X_0$ 所组成的集合：$a_{\overline{x}} = b_{\overline{x}}$。由引理
\ref{etale-cohomology-lemma-support-constructible}，子集 $E \subset X_0$ 可构造。根据假设，
$X \to X_0$ 映入 $E$。由极限，引理
\ref{limits-lemma-limit-contained-in-constructible}，可找到 $i \geq 0$，使得
$X_i \to X_0$ 映入 $E$。因此 $f_{i0}^{-1}a = f_{i0}^{-1}b$。

\medskip\noindent
全忠实性。选取 $0 \in I$，并设 $\mathcal{F}_0$、$\mathcal{G}_0$ 为 $X_0$ 上的
可构造层。假设 $a : f_0^{-1}\mathcal{F}_0 \to f_0^{-1}\mathcal{G}_0$ 是一个态射。
我们断言，存在某个 $i$ 和态射
$a_i : f_{i0}^{-1}\mathcal{F}_0 \to f_{i0}^{-1}\mathcal{G}_0$
，它的拉回即为 $a$，这是 $X$ 上的态射。由引理
\ref{etale-cohomology-lemma-category-constructible-sets}，可以替换 $\mathcal{F}_0$，所用对象是若干
由 $(X_0)_\etale$ 的拟紧拟分离对象表示的层的有限余积。因此，需要证明：
若 $U_0 \to X_0$ 是 $(X_0)_\etale$ 的这样一个
对象，则
$$
f_0^{-1}\mathcal{G}(U) = \colim_{i \geq 0} f_{i0}^{-1}\mathcal{G}(U_i)
$$
，其中 $U = X \times_{X_0} U_0$ 且 $U_i = X_i \times_{X_0} U_0$。这是定理
\ref{etale-cohomology-theorem-colimit} 的特殊情形。

\medskip\noindent
本质满射性。需要证明，每个可构造层 $\mathcal{F}$（定义在 $X$ 上）都同构于
$f_i^{-1}\mathcal{F}$，其中某个可构造层 $\mathcal{F}_i$ 定义在 $X_i$ 上。
应用引理 \ref{etale-cohomology-lemma-category-constructible-sets} 并使用前两段的结论，可知只需对
$h_U$ 证明这一点，其中 $U$ 是 $X_\etale$ 的某个拟紧拟分离对象。在这种情形中，
需要证明 $U$ 是某个在 $X_i$ 上 \'etale 的拟紧拟分离概形的基变换，其中 $i$
是某个指标。这由极限，引理 \ref{limits-lemma-descend-finite-presentation} 和
\ref{limits-lemma-descend-etale} 得出。

\medskip\noindent
证明 (3)。论证与集合层的论证非常类似，只是用引理
\ref{etale-cohomology-lemma-category-constructible-modules} 代替引理
\ref{etale-cohomology-lemma-category-constructible-sets}。细节略去。(2) 由 (3) 得出，因为拟紧概形上的
每个可构造阿贝尔层都是可构造 $\mathbf{Z}/n\mathbf{Z}$-模层，其中 $n$ 为某个整数。
\end{proof}

\begin{lemma}
\label{etale-cohomology-lemma-category-loc-cst-is-colimit}
设 $X = \lim_{i \in I} X_i$ 为一个具有仿射转移态射的概形有向系的极限。
假设 $X_i$ 对所有 $i \in I$ 都拟紧拟分离。
\begin{enumerate}
\item $X_\etale$ 上的有限局部常值层范畴是各 $(X_i)_\etale$ 上有限局部常值层
范畴的余极限。
\item $X_\etale$ 上的有限局部常值阿贝尔层范畴是各 $(X_i)_\etale$ 上有限局部
常值阿贝尔层范畴的余极限。
\item 设 $\Lambda$ 为诺特环。有限型局部常值 $\Lambda$-模层在 $X_\etale$ 上的
范畴是有限型局部常值 $\Lambda$-模层在各 $(X_i)_\etale$ 上的范畴的余极限。
\end{enumerate}
\end{lemma}

\begin{proof}
由引理 \ref{etale-cohomology-lemma-category-is-colimit}，每种情形中的函子都是全忠实的。由同一引理，
为完成情形 (1) 的证明，只需证明如下结论：给定可构造层 $\mathcal{F}_i$，定义在
$X_i$ 上；若其拉回 $\mathcal{F}$ 定义在 $X$ 上，并且是有限局部常值的，则存在
$i' \geq i$，使得拉回 $\mathcal{F}_{i'}$（它来自 $\mathcal{F}_i$，并定义在
$X_{i'}$ 上）
是有限局部常值的。根据假设，存在 \'etale 覆盖
$\mathcal{U} = \{U_j \to X\}_{j \in J}$
，使得有 $\mathcal{F}|_{U_j} \cong \underline{S_j}$，其中 $S_j$ 为某个有限集。
可以假定 $U_j$ 对所有 $j \in J$ 都是仿射的。由于 $X$ 拟紧，可以假定 $J$
有限。由引理 \ref{etale-cohomology-lemma-colimit-affine-sites}，可找到 $i' \geq i$ 和 \'etale 覆盖
$\mathcal{U}_{i'} = \{U_{i', j} \to X_{i'}\}_{j \in J}$，其到 $X$ 的基变换为
$\mathcal{U}$。于是，可构造层 $\mathcal{F}_{i'}|_{U_{i', j}}$ 和
$\underline{S_j}$ 定义在 $(U_{i', j})_\etale$ 上，并且拉回到 $U_j$ 后同构。因此，
增大 $i'$ 后可使 $\mathcal{F}_{i'}|_{U_{i', j}}$ 和 $\underline{S_j}$ 同构。
所以 $\mathcal{F}_{i'}$ 是有限局部常值的。情形 (2) 和 (3) 的证明完全相同。
\end{proof}

\begin{lemma}
\label{etale-cohomology-lemma-irreducible-subsheaf-constant-zero}
设 $X$ 为具有泛点 $\eta$ 的不可约概形。
\begin{enumerate}
\item 设 $S' \subset S$ 为集合的包含。若在
$\underline{S'} \subset \mathcal{G} \subset \underline{S}$
在 $\Sh(X_\etale)$ 中成立，且 $S' = \mathcal{G}_{\overline{\eta}}$，则
$\mathcal{G} = \underline{S'}$。
\item 设 $A' \subset A$ 为阿贝尔群的包含。若在
$\underline{A'} \subset \mathcal{G} \subset \underline{A}$
在 $\textit{Ab}(X_\etale)$ 中成立，且 $A' = \mathcal{G}_{\overline{\eta}}$，则
$\mathcal{G} = \underline{A'}$。
\item 设 $M' \subset M$ 为环 $\Lambda$ 上模的包含。若
$\underline{M'} \subset \mathcal{G} \subset \underline{M}$ 在
$\textit{Mod}(X_\etale, \underline{\Lambda})$ 中成立，且
$M' = \mathcal{G}_{\overline{\eta}}$，则 $\mathcal{G} = \underline{M'}$。
\end{enumerate}
\end{lemma}

\begin{proof}
这是因为，对每个 \'etale 态射 $U \to X$，只要 $U \not = \emptyset$，点
$\eta$ 都在其像中。
\end{proof}

\begin{lemma}
\label{etale-cohomology-lemma-push-constant-sheaf-from-open}
设 $X$ 为函数域为 $K$ 的整正规概形。设 $E$ 为集合。
\begin{enumerate}
\item 设 $g : \Spec(K) \to X$ 为泛点的包含。则
$g_*\underline{E} = \underline{E}$。
\item 设 $j : U \to X$ 为非空开集的包含。则
$j_*\underline{E} = \underline{E}$。
\end{enumerate}
\end{lemma}

\begin{proof}
证明 (1)。设 $x \in X$ 为一点。令
$\mathcal{O}^{sh}_{X, \overline{x}}$
为 $\mathcal{O}_{X, x}$ 的一个严格 Hensel 化。由代数进阶，引理
\ref{more-algebra-lemma-henselization-normal}，可知
$\mathcal{O}^{sh}_{X, \overline{x}}$ 是正规整环。因此
$\Spec(K) \times_X \Spec(\mathcal{O}^{sh}_{X, \overline{x}})$ 不可约。
由此，茎 $(g_*\underline{E})_{\overline{x}}$ 等于 $E$，见定理
\ref{etale-cohomology-theorem-higher-direct-images}。

\medskip\noindent
证明 (2)。由于 $g$ 通过 $j$ 分解，存在态射
$j_*\underline{E} \to g_*\underline{E}$。该态射是单射，因为对于每个概形 $V$
（它在 $X$ 上 \'etale），集合 $\Spec(K) \times_X V$ 在 $U \times_X V$ 中稠密。
另一方面，有态射 $\underline{E} \to j_*\underline{E}$，因而得出结论。
\end{proof}

\begin{lemma}
\label{etale-cohomology-lemma-zero-in-generic-point}
设 $X$ 为拟紧拟分离概形。设 $\eta \in X$ 为 $X$ 的某个不可约分支的泛点。
\begin{enumerate}
\item 设 $\mathcal{F}$ 为 $X_\etale$ 上的挠阿贝尔层，其茎
$\mathcal{F}_{\overline{\eta}}$ 为零。则 $\mathcal{F} = \colim \mathcal{F}_i$
是可构造阿贝尔层 $\mathcal{F}_i$ 的滤过余极限，使得对每个 $i$，
$\mathcal{F}_i$ 的支集都包含在某个不含 $\eta$ 的闭子概形中。
\item 设 $\Lambda$ 为诺特环，$\mathcal{F}$ 为 $\Lambda$-模层，定义在 $X_\etale$ 上，
其茎 $\mathcal{F}_{\overline{\eta}}$ 为零。则
$\mathcal{F} = \colim \mathcal{F}_i$
是可构造 $\Lambda$-模层 $\mathcal{F}_i$ 的滤过余极限，使得对每个 $i$，
$\mathcal{F}_i$ 的支集都包含在某个不含 $\eta$ 的闭子概形中。
\end{enumerate}
\end{lemma}

\begin{proof}
证明 (1)。由引理 \ref{etale-cohomology-lemma-torsion-colimit-constructible}，可以写成
$\mathcal{F} = \colim_{i \in I} \mathcal{F}_i$，其中 $\mathcal{F}_i$ 是可构造
阿贝尔层。选取 $i \in I$。由于根据假设 $\mathcal{F}|_\eta$ 为零，可知存在
$i'(i) \geq i$，使得 $\mathcal{F}_i|_\eta \to \mathcal{F}_{i'(i)}|_\eta$ 为零，
见引理 \ref{etale-cohomology-lemma-colimit-constructible}。于是
$\mathcal{G}_i = \Im(\mathcal{F}_i \to \mathcal{F}_{i'(i)})$ 是可构造阿贝尔层
（引理 \ref{etale-cohomology-lemma-constructible-abelian}），且其在 $\eta$ 处的茎为零。因此，令
$E_i$ 为 $\mathcal{G}_i$ 的支集；它是 $X$ 的不含 $\eta$ 的可构造子集。由于 $\eta$
是 $X$ 的某个不可约分支的泛点，由拓扑学，引理
\ref{topology-lemma-generic-point-in-constructible} 可知
$\eta \not \in Z_i = \overline{E_i}$。定义新的有向集 $I'$，其底集取为 $I$，
并用如下规则定义序：$i_1$ 大于或等于 $i_2$，当且仅当
$i_1 \geq i'(i_2)$。于是，层 $\mathcal{G}_i$ 构成 $I'$ 上的一个系，其余极限为
$\mathcal{F}$，证明完成。

\medskip\noindent
情形 (2) 的证明完全相同，故略去。
\end{proof}








\section{诺特概形上的可构造层}
\label{etale-cohomology-section-noetherian-constructible}

\noindent
若 $X$ 是诺特概形，则任意局部闭子集都是可构造局部闭子集
（拓扑学，引理 \ref{topology-lemma-constructible-Noetherian-space}）。
因此，阿贝尔层 $\mathcal{F}$ 定义在 $X_\etale$ 上时，它是可构造的，当且仅当存在有限分划
$X = \coprod X_i$，使得 $\mathcal{F}|_{X_i}$ 是有限局部常值的。
（按照约定，拓扑空间的分划以局部闭子集为各部分；见拓扑学，
第 \ref{topology-section-stratifications} 节。）换言之，在定义
\ref{etale-cohomology-definition-constructible} 中可以省略形容词“可构造”。事实上，诺特概形上的
可构造层范畴还具有一些附加性质，我们将在本节中逐一列出。

\begin{proposition}
\label{etale-cohomology-proposition-constructible-over-noetherian}
设 $X$ 为诺特概形，$\Lambda$ 为诺特环。
\begin{enumerate}
\item 可构造集合层的任意子层或商层都是可构造的。
\item $X_\etale$ 上的可构造阿贝尔层范畴是
$\textit{Ab}(X_\etale)$ 的（强）Serre 子范畴。特别地，$X_\etale$ 上可构造
阿贝尔层的任意子层和商层都是可构造的。
\item 可构造 $\Lambda$-模层在 $X_\etale$ 上的范畴是
$\textit{Mod}(X_\etale, \Lambda)$ 的（强）Serre 子范畴。特别地，
可构造 $\Lambda$-模层在 $X_\etale$ 上的任意子模层和商模层都是可构造的。
\end{enumerate}
\end{proposition}

\begin{proof}
证明 (1)。设 $\mathcal{G} \subset \mathcal{F}$，其中 $\mathcal{F}$ 是
$X_\etale$ 上的可构造集合层。设 $\eta \in X$ 为 $X$ 的某个不可约分支的泛点。
由诺特归纳法，只需找到一个开邻域 $U$，它包含 $\eta$，且使 $\mathcal{G}|_U$ 局部常值。
为此，可以替换 $X$，代之以 $\eta$ 的一个 \'etale 邻域。因而可以假定
$\mathcal{F}$ 是常值层且 $X$ 不可约。

\medskip\noindent
设 $\mathcal{F} = \underline{S}$，其中 $S$ 为某个有限集。于是
$S' = \mathcal{G}_{\overline{\eta}} \subset S$；记
$S' = \{s_1, \ldots, s_t\}$。选取一个 \'etale 邻域
$(U, \overline{u})$，它是 $\overline{\eta}$ 的邻域；再选取截面
$\sigma_1, \ldots, \sigma_t \in \mathcal{G}(U)$，使其映到 $s_i$，后者视为
$\mathcal{G}_{\overline{\eta}} \subset S$ 中的元素。由于 $\sigma_i$ 映到元素
$s_i \in S' \subset S = \Gamma(X, \mathcal{F})$，故 $\sigma_i$ 在
$U \times_X U$ 上的两个拉回作为 $\mathcal{G}$ 的截面相同。由 $\mathcal{G}$ 的
层条件可知，$\sigma_i$ 来自 $\mathcal{G}$ 在开集
$\Im(U \to X)$ 上的一个截面；这是 $X$ 的开集。缩小 $X$ 后，可以假定
$\underline{S'} \subset \mathcal{G} \subset \underline{S}$。于是由引理
\ref{etale-cohomology-lemma-irreducible-subsheaf-constant-zero} 可知
$\underline{S'} = \mathcal{G}$。

\medskip\noindent
设 $\mathcal{F} \to \mathcal{Q}$ 为满射，其中 $\mathcal{F}$ 是 $X_\etale$ 上的
可构造集合层。令
$\mathcal{G} = \mathcal{F} \times_\mathcal{Q} \mathcal{F}$。
由证明的第一部分，$\mathcal{G}$ 作为 $\mathcal{F} \times \mathcal{F}$ 的子层是
可构造的。这又蕴含 $\mathcal{Q}$ 可构造；见引理
\ref{etale-cohomology-lemma-constructible-abelian}。

\medskip\noindent
证明 (3)。我们已经知道，可构造模层构成一个弱 Serre 子范畴；见引理
\ref{etale-cohomology-lemma-constructible-abelian}。因此，只需证明关于子模层的断言。

\medskip\noindent
设 $\mathcal{G} \subset \mathcal{F}$ 是某个可构造 $\Lambda$-模层的子模层，
该层定义在 $X_\etale$ 上。设 $\eta \in X$ 为 $X$ 的某个不可约分支的泛点。
由诺特归纳法，只需找到一个开邻域 $U$，它包含 $\eta$，且使 $\mathcal{G}|_U$ 局部常值。
为此，可以替换 $X$，代之以 $\eta$ 的一个 \'etale 邻域。因而可以假定
$\mathcal{F}$ 是常值层且 $X$ 不可约。

\medskip\noindent
设 $\mathcal{F} = \underline{M}$，其中所取的有限 $\Lambda$-模记作 $M$。于是
$M' = \mathcal{G}_{\overline{\eta}} \subset M$。选取有限多个元素
$s_1, \ldots, s_t$，使其生成 $M'$ 这个 $\Lambda$-模。（这是可能的，因为
$\Lambda$ 是诺特的且 $M$ 有限。）选取一个 \'etale 邻域
$(U, \overline{u})$，它是 $\overline{\eta}$ 的邻域；再选取截面
$\sigma_1, \ldots, \sigma_t \in \mathcal{G}(U)$，使其映到 $s_i$，后者视为
$\mathcal{G}_{\overline{\eta}} \subset M$ 中的元素。由于 $\sigma_i$ 映到元素
$s_i \in M' \subset M = \Gamma(X, \mathcal{F})$，故 $\sigma_i$ 在
$U \times_X U$ 上的两个拉回作为 $\mathcal{G}$ 的截面相同。由 $\mathcal{G}$ 的
层条件可知，$\sigma_i$ 来自 $\mathcal{G}$ 在开集
$\Im(U \to X)$ 上的一个截面；这是 $X$ 的开集。缩小 $X$ 后，可以假定
$\underline{M'} \subset \mathcal{G} \subset \underline{M}$。于是由引理
\ref{etale-cohomology-lemma-irreducible-subsheaf-constant-zero} 可知
$\underline{M'} = \mathcal{G}$。

\medskip\noindent
证明 (2)。这以通常方式由 (3) 得出。细节略去。
\end{proof}

\noindent
下面的引理说明，$X$ 上可构造层的阿贝尔范畴中的每个对象都是“诺特的”，
即对子对象满足升链条件（a.c.c.）。

\begin{lemma}
\label{etale-cohomology-lemma-constructible-over-noetherian-noetherian}
设 $X$ 为诺特概形，$\Lambda$ 为诺特环。考虑包含关系
$$
\mathcal{F}_1 \subset \mathcal{F}_2 \subset \mathcal{F}_3 \subset \ldots
\subset \mathcal{F}
$$
，它们位于集合层、阿贝尔群层或 $\Lambda$-模层的范畴中。若 $\mathcal{F}$
可构造，则对某个 $n$ 有
$\mathcal{F}_n = \mathcal{F}_{n + 1} = \mathcal{F}_{n + 2} = \ldots$。
\end{lemma}

\begin{proof}
由命题 \ref{etale-cohomology-proposition-constructible-over-noetherian} 可知，
$\mathcal{F}_i$ 和 $\colim \mathcal{F}_i$ 都可构造。于是该引理由引理
\ref{etale-cohomology-lemma-colimit-constructible} 得出。
\end{proof}

\begin{lemma}
\label{etale-cohomology-lemma-constructible-maps-into-constant}
设 $X$ 为诺特概形。
\begin{enumerate}
\item 设 $\mathcal{F}$ 为 $X_\etale$ 上的可构造集合层。存在层的单射
$$
\mathcal{F} \longrightarrow
\prod\nolimits_{i = 1, \ldots, n} f_{i, *}\underline{E_i}
$$
，其中 $f_i : Y_i \to X$ 是有限态射，$E_i$ 是有限集。
\item 设 $\mathcal{F}$ 为 $X_\etale$ 上的可构造阿贝尔层。存在阿贝尔层的单射
$$
\mathcal{F} \longrightarrow
\bigoplus\nolimits_{i = 1, \ldots, n} f_{i, *}\underline{M_i}
$$
，其中 $f_i : Y_i \to X$ 是有限态射，$M_i$ 是有限阿贝尔群。
\item 设 $\Lambda$ 为诺特环，设 $\mathcal{F}$ 为可构造 $\Lambda$-模层，
定义在 $X_\etale$ 上。存在模层的单射
$$
\mathcal{F} \longrightarrow
\bigoplus\nolimits_{i = 1, \ldots, n} f_{i, *}\underline{M_i}
$$
，其中 $f_i : Y_i \to X$ 是有限态射，$M_i$ 是有限 $\Lambda$-模。
\end{enumerate}
此外，可以假定每个 $Y_i$ 都不可约且既约，并映满某个不可约既约闭子概形
$Z_i \subset X$，且 $Y_i \to Z_i$ 在 $Z_i$ 的某个非空开集上有限 \'etale。
\end{lemma}

\begin{proof}
证明 (1)。由于 $\mathcal{F}$ 的子层满足升链条件（引理
\ref{etale-cohomology-lemma-constructible-over-noetherian-noetherian}），只需证明：对每个点
$x \in X$，可以找到态射 $\varphi : \mathcal{F} \to f_*\underline{E}$，其中
$f : Y \to X$ 有限且 $E$ 为有限集，并使
$\varphi_{\overline{x}} : \mathcal{F}_{\overline{x}} \to
(f_*S)_{\overline{x}}$ 为单射。（也可以避开这一论证：选取引理
\ref{etale-cohomology-lemma-constructible-quasi-compact-quasi-separated} 中那样的 $X$ 的一个分划，
并对每个部分的每个不可约分支构造一个 $Y_i \to X$。）令 $Z \subset X$ 为
$\overline{\{x\}}$ 上的诱导既约概形结构（概形，定义
\ref{schemes-definition-reduced-induced-scheme}）。由于 $\mathcal{F}$ 可构造，
存在有限可分扩张 $K/\kappa(x)$，使得 $\mathcal{F}|_{\Spec(K)}$ 是以
$E$ 为值的常值层，其中 $E$ 为某个有限集。令 $Y \to Z$ 为 $Z$ 在
$\Spec(K)$ 中的正规化。由态射，引理
\ref{morphisms-lemma-normal-normalization}，$Y$ 是正规整概形。由于
$K/\kappa(x)$ 是有限扩张，显然 $K$ 是 $Y$ 的函数域。用
$g : \Spec(K) \to Y$ 表示该包含。态射
$\mathcal{F}|_{\Spec(K)} \to \underline{E}$ 与态射
$\mathcal{F}|_Y \to g_*\underline{E} = \underline{E}$ 伴随（引理
\ref{etale-cohomology-lemma-push-constant-sheaf-from-open}）。后者又与态射
$\varphi : \mathcal{F} \to f_*\underline{E}$ 伴随。注意，$\varphi$ 在几何点
$\overline{x}$ 处的茎是单射：可以选取一个提升 $\overline{y} \in Y$，
它是 $\overline{x}$ 的提升，而交换图
$$
\xymatrix{
\mathcal{F}_{\overline{x}} \ar@{=}[r] \ar[d] &
(\mathcal{F}|_Y)_{\overline{y}} \ar@{=}[d] \\
(f_*\underline{E})_{\overline{x}} \ar[r] &
\underline{E}_{\overline{y}}
}
$$
证明了单射性。不过，证明尚未完成，因为态射 $f : Y \to Z$ 是整的，
但一般并非有限的\footnote{若 $X$ 是 Nagata 概形，例如域上的有限型概形，
则 $Y \to Z$ 是有限的。}。

\medskip\noindent
为解决前一段末句所述的问题，写成 $Y = \lim_{i \in I} Y_i$，其中 $Y_i$
不可约、整，且在 $Z$ 上有限。具体而言，把性质，引理
\ref{properties-lemma-integral-algebra-directed-colimit-finite} 应用于
$f_*\mathcal{O}_Y$，把它视为 $\mathcal{O}_Z$-代数层，再应用函子
$\underline{\Spec}_Z$。于是由引理 \ref{etale-cohomology-lemma-relative-colimit} 有
$f_*\underline{E} = \colim f_{i, *}\underline{E}$。由引理
\ref{etale-cohomology-lemma-constructible-is-compact}，态射 $\mathcal{F} \to f_*\underline{E}$
通过 $f_{i, *}\underline{E}$ 分解，其中 $i$ 是某个指标。由于 $Y_i \to Z$ 是整概形之间的
有限态射，且该态射诱导的函数域扩张有限可分，故该态射在 $Z$ 的某个非空开集上
有限 \'etale（使用代数，引理 \ref{algebra-lemma-smooth-at-generic-point}；
细节略去）。这就完成了 (1) 的证明。

\medskip\noindent
(2) 和 (3) 的证明与 (1) 的证明完全相同。
\end{proof}

\noindent
在下面的引理中，我们用一个标准技巧把一个非常一般的陈述归约到诺特情形。

\begin{lemma}
\label{etale-cohomology-lemma-constructible-maps-into-constant-general}
\begin{reference}
\cite[Exposee IX, Proposition 2.14]{SGA4}
\end{reference}
设 $X$ 为拟紧拟分离概形。
\begin{enumerate}
\item 设 $\mathcal{F}$ 为 $X_\etale$ 上的可构造集合层。存在层的单射
$$
\mathcal{F} \longrightarrow
\prod\nolimits_{i = 1, \ldots, n} f_{i, *}\underline{E_i}
$$
，其中 $f_i : Y_i \to X$ 是有限且有限表现的态射，$E_i$ 是有限集。
\item 设 $\mathcal{F}$ 为 $X_\etale$ 上的可构造阿贝尔层。存在阿贝尔层的单射
$$
\mathcal{F} \longrightarrow
\bigoplus\nolimits_{i = 1, \ldots, n} f_{i, *}\underline{M_i}
$$
，其中 $f_i : Y_i \to X$ 是有限且有限表现的态射，$M_i$ 是有限阿贝尔群。
\item 设 $\Lambda$ 为诺特环，设 $\mathcal{F}$ 为可构造 $\Lambda$-模层，
定义在 $X_\etale$ 上。存在模层的单射
$$
\mathcal{F} \longrightarrow
\bigoplus\nolimits_{i = 1, \ldots, n} f_{i, *}\underline{M_i}
$$
，其中 $f_i : Y_i \to X$ 是有限且有限表现的态射，$M_i$ 是有限 $\Lambda$-模。
\end{enumerate}
\end{lemma}

\begin{proof}
我们通过绝对诺特逼近把该引理归约到诺特情形。具体而言，由极限，命题
\ref{limits-proposition-approximate}，可以写成
$X = \lim_{t \in T} X_t$，其中每个 $X_t$ 在 $\Spec(\mathbf{Z})$ 上有限型，
且转移态射为仿射态射。由引理 \ref{etale-cohomology-lemma-category-is-colimit}，可构造层
（集合层、阿贝尔群层或 $\Lambda$-模层）在 $X_\etale$ 上的范畴是各
$X_t$ 上相应范畴的余极限。因此，我们的可构造层 $\mathcal{F}$ 是同类可构造层
$\mathcal{F}_t$ 的拉回；后者定义在 $X_t$ 上，其中 $t$ 是某个指标。随后对诺特情形应用引理
\ref{etale-cohomology-lemma-constructible-maps-into-constant}，得到单射
$$
\mathcal{F}_t \longrightarrow
\prod\nolimits_{i = 1, \ldots, n} f_{i, *}\underline{E_i}
\quad\text{或}\quad
\mathcal{F}_t \longrightarrow
\bigoplus\nolimits_{i = 1, \ldots, n} f_{i, *}\underline{M_i}
$$
，它位于 $X_t$ 上，其中 $f_i : Y_i \to X_t$ 是某些有限态射。由于 $X_t$
是诺特的，态射 $f_i$ 是有限表现的。由于拉回正合，而且 $f_{i, *}$ 的形成与
基变换可交换（引理 \ref{etale-cohomology-lemma-finite-pushforward-commutes-with-base-change}），
结论得证。
\end{proof}

\begin{lemma}
\label{etale-cohomology-lemma-support-in-subset}
设 $X$ 为诺特概形，$E \subset X$ 为在特殊化下封闭的子集。
\begin{enumerate}
\item 设 $\mathcal{F}$ 为 $X_\etale$ 上的挠阿贝尔层，其支集包含于 $E$。则
$\mathcal{F} = \colim \mathcal{F}_i$ 是可构造阿贝尔层 $\mathcal{F}_i$ 的滤过
余极限，且对每个 $i$，$\mathcal{F}_i$ 的支集包含在某个包含于 $E$ 的闭子集中。
\item 设 $\Lambda$ 为诺特环，$\mathcal{F}$ 为 $\Lambda$-模层，定义在 $X_\etale$ 上，
其支集包含于 $E$。则 $\mathcal{F} = \colim \mathcal{F}_i$ 是可构造
$\Lambda$-模层 $\mathcal{F}_i$ 的滤过余极限，且对每个 $i$，$\mathcal{F}_i$
的支集包含在某个包含于 $E$ 的闭子集中。
\end{enumerate}
\end{lemma}

\begin{proof}
证明 (1)。由引理 \ref{etale-cohomology-lemma-torsion-colimit-constructible}，可以写成
$\mathcal{F} = \colim_{i \in I} \mathcal{F}_i$，其中 $\mathcal{F}_i$ 是可构造
阿贝尔层。由命题 \ref{etale-cohomology-proposition-constructible-over-noetherian}，像
$\mathcal{F}'_i \subset \mathcal{F}$（它来自态射
$\mathcal{F}_i \to \mathcal{F}$）是可构造的。因此
$\mathcal{F} = \colim \mathcal{F}'_i$，而 $\mathcal{F}'_i$ 的支集包含在 $E$ 中。
由于 $\mathcal{F}'_i$ 的支集是可构造的（根据可构造层的定义），故其闭包也包含于
$E$；例如见拓扑学，引理
\ref{topology-lemma-constructible-stable-specialization-closed}。

\medskip\noindent
情形 (2) 的证明完全相同，故略去。
\end{proof}








\section{特殊化与 \'etale 层}
\label{etale-cohomology-section-specialization}

\noindent
拓扑图景如下：设 $X$ 为拓扑空间，$x' \leadsto x$ 为 $X$ 中点的一个特殊化。
则 $x$ 的每个开邻域都包含 $x'$。因此，对任意层 $\mathcal{F}$（定义在 $X$ 上），
存在茎的{\it 特殊化映射}
$$
sp : \mathcal{F}_x \longrightarrow \mathcal{F}_{x'}
$$
，它把二元组 $(U, s)$ 在 $\mathcal{F}_x$ 中的等价类映到
二元组 $(U, s)$ 在 $\mathcal{F}_{x'}$ 中的等价类。关于用二元组的等价类描述茎，
见层，第 \ref{sheaves-section-stalks} 节。这个映射当然关于 $\mathcal{F}$ 具有
函子性；换言之，$sp$ 是函子之间的变换。

\medskip\noindent
对 \'etale 拓扑中的层，可以仿照这一构造；见
\cite[Exposee VIII, 7.7, page 397]{SGA4}。为此，设有概形 $S$、几何点
$\overline{s}$（它属于 $S$），以及几何点 $\overline{t}$（它属于
$\Spec(\mathcal{O}^{sh}_{S, \overline{s}})$）。对任意层 $\mathcal{F}$（定义在
$S_\etale$ 上），我们将构造
{\it 特殊化映射}
$$
sp : \mathcal{F}_{\overline{s}} \longrightarrow \mathcal{F}_{\overline{t}}
$$
这里略有语言上的滥用：不应写成 $\mathcal{F}_{\overline{t}}$，而应写成
$\mathcal{F}_{p(\overline{t})}$，其中
$p : \Spec(\mathcal{O}^{sh}_{S, \overline{s}}) \to S$ 是典范态射。回忆
$$
\mathcal{F}_{\overline{s}} = \colim_{(U, \overline{u})} \mathcal{F}(U)
$$
，其中余极限取遍所有 \'etale 邻域 $(U, \overline{u})$，它们是
$(S, \overline{s})$ 的邻域；见第 \ref{etale-cohomology-section-stalks} 节。由于
$\mathcal{O}^{sh}_{S, \overline{s}}$ 是结构层的茎，对每个 \'etale 邻域
$(U, \overline{u})$（它是 $(S, \overline{s})$ 的邻域），都有典范映射
$\mathcal{O}_{U, u} \to \mathcal{O}^{sh}_{S, \overline{s}}$。因此得到唯一分解
$$
\Spec(\mathcal{O}^{sh}_{S, \overline{s}}) \to U \to S
$$
若 $\overline{v}$ 表示 $\overline{t}$ 在 $U$ 中的像，则
$(U, \overline{v})$ 是 $(S, \overline{t})$ 的一个 \'etale 邻域。该构造定义了一个
函子，它从 $(S, \overline{s})$ 的 \'etale 邻域范畴到
$(S, \overline{t})$ 的 \'etale 邻域范畴。因此，可以定义映射
$sp : \mathcal{F}_{\overline{s}} \to \mathcal{F}_{\overline{t}}$
：它把 $(U, \overline{u}, \sigma)$ 的等价类（其中
$\sigma \in \mathcal{F}(U)$）映到 $(U, \overline{v}, \sigma)$ 的等价类。

\medskip\noindent
设 $K \in D(S_\etale)$。对如上的 $\overline{s}$ 和 $\overline{t}$，有
{\it 特殊化映射}
$$
sp : K_{\overline{s}} \longrightarrow K_{\overline{t}}
\quad\text{于}\quad D(\textit{Ab})
$$
具体而言，若 $K$ 由阿贝尔层复形 $\mathcal{F}^\bullet$ 表示，则直接取映射
$$
K_{\overline{s}} = \mathcal{F}^\bullet_{\overline{s}}
\longrightarrow
\mathcal{F}^\bullet_{\overline{t}} = K_{\overline{t}}
$$
，其各项由上面构造的层的特殊化映射给出。由于茎函子正合，这一映射与表示
$K$ 的复形之选择无关（即逐项取复形的茎在导出范畴上有良好定义）。

\medskip\noindent
显然，该构造关于层 $\mathcal{F}$ 具有函子性，该层定义在 $S_\etale$ 上。
若把茎函子看成拓扑斯的态射
$\overline{s}, \overline{t} : \textit{Sets} \to \Sh(S_\etale)$，则可以把
$sp$ 看成如下 $2$-态射：
$$
\xymatrix{
\textit{Sets}
\rrtwocell^{\overline{t}}_{\overline{s}}{\ sp}
&
&
\Sh(S_\etale)
}
$$
。

\begin{remark}[sp 的另一种描述]
\label{etale-cohomology-remark-alternative-sp}
设 $S$、$\overline{s}$ 和 $\overline{t}$ 如上。描述特殊化映射的另一种方式是使用
$$
\mathcal{F}_{\overline{s}} =
\Gamma(\Spec(\mathcal{O}^{sh}_{S, \overline{s}}), p^{-1}\mathcal{F})
\quad\text{且}\quad
\mathcal{F}_{\overline{t}} = \Gamma(\overline{t},
\overline{t}^{-1}p^{-1}\mathcal{F})
$$
。第一个等式由把定理 \ref{etale-cohomology-theorem-higher-direct-images} 应用于
$\text{id}_S : S \to S$ 得出，第二个等式由引理
\ref{etale-cohomology-lemma-stalk-pullback} 得出。于是，可以把 $sp$ 看成映射
$$
sp :
\mathcal{F}_{\overline{s}} =
\Gamma(\Spec(\mathcal{O}^{sh}_{S, \overline{s}}), p^{-1}\mathcal{F})
\xrightarrow{\text{pullback by }\overline{t}}
\Gamma(\overline{t}, \overline{t}^{-1}p^{-1}\mathcal{F}) =
\mathcal{F}_{\overline{t}}
$$
\end{remark}

\begin{remark}[sp 的又一种描述]
\label{etale-cohomology-remark-another-sp}
设 $S$、$\overline{s}$ 和 $\overline{t}$ 如上。还可以使用唯一态射
$$
c :
\Spec(\mathcal{O}^{sh}_{S, \overline{t}})
\longrightarrow
\Spec(\mathcal{O}^{sh}_{S, \overline{s}})
$$
，它位于 $S$ 上，并且与给定态射
$\overline{t} \to \Spec(\mathcal{O}^{sh}_{S, \overline{s}})$ 及态射
$\overline{t} \to \Spec(\mathcal{O}^{sh}_{t, \overline{t}})$ 相容。把代数，引理
\ref{algebra-lemma-map-into-henselian-colimit} 应用于 $\mathcal{O}_{S, s}$、
$\mathcal{O}^{sh}_{S, \overline{t}}$ 和
$\mathcal{O}^{sh}_{S, \overline{s}} \to \kappa(\overline{t})$，即可得到所示箭头的
唯一性与存在性。于是有
$$
sp :
\mathcal{F}_{\overline{s}} =
\Gamma(\Spec(\mathcal{O}^{sh}_{S, \overline{s}}), \mathcal{F})
\xrightarrow{\text{pullback by }c}
\Gamma(\Spec(\mathcal{O}^{sh}_{S, \overline{t}}), \mathcal{F}) =
\mathcal{F}_{\overline{t}}
$$
（采用显然的记号约定）。事实上，该过程也适用于对象 $K$，该对象属于
$D(S_\etale)$：
$K$ 的特殊化映射等于映射
$$
sp :
K_{\overline{s}} =
R\Gamma(\Spec(\mathcal{O}^{sh}_{S, \overline{s}}), K)
\xrightarrow{\text{pullback by }c}
R\Gamma(\Spec(\mathcal{O}^{sh}_{S, \overline{t}}), K) =
K_{\overline{t}}
$$
。这些等号成立，是因为在严格 Hensel 概形
$\Spec(\mathcal{O}^{sh}_{S, \overline{s}})$ 和
$\Spec(\mathcal{O}^{sh}_{S, \overline{t}})$ 上取全局截面是正合的（并且分别等同于
在 $\overline{s}$ 和 $\overline{t}$ 处取茎）；因此，不会出现与 $K$ 可能无界这一
事实有关的微妙问题。
\end{remark}

\begin{remark}[提升特殊化]
\label{etale-cohomology-remark-can-lift}
设 $S$ 为概形，$t \leadsto s$ 为 $S$ 上点的一个特殊化。选取几何点
$\overline{t}$ 和 $\overline{s}$，它们分别位于 $t$ 和 $s$ 上方。由概形，引理
\ref{schemes-lemma-specialize-points}，$t$ 对应于
$\Spec(\mathcal{O}_{S, s})$ 的一个点；又因为
$\mathcal{O}_{S, s} \to \mathcal{O}^{sh}_{S, \overline{s}}$ 忠实平坦，故可找到点
$t' \in \Spec(\mathcal{O}^{sh}_{S, \overline{s}})$，它映到 $t$。由于
$\Spec(\mathcal{O}^{sh}_{S, \overline{s}})$ 是 $S$ 上 \'etale 概形的极限，
$\kappa(t')/\kappa(t)$ 是可分代数扩张（当然通常不是有限的）。由于
$\kappa(\overline{t})$ 是代数闭的，可以选取嵌入
$\kappa(t') \to \kappa(\overline{t})$，它是 $\kappa(t)$ 上扩张之间的嵌入。
这一选择给出交换图
$$
\xymatrix{
\overline{t} \ar[d] \ar[r] &
\Spec(\mathcal{O}^{sh}_{S, \overline{s}}) \ar[d] &
\overline{s} \ar[l] \ar[d] \\
t \ar[r] & S & s \ar[l]
}
$$
，它由点和几何点组成。因此，若 $t \leadsto s$，总可以把 $\overline{t}$“提升”为
$S$ 在 $\overline{s}$ 处的严格 Hensel 化的一个几何点，并得到上述特殊化映射。
\end{remark}

\begin{lemma}
\label{etale-cohomology-lemma-specialization-map-pullback}
设 $g : S' \to S$ 为概形态射，$\mathcal{F}$ 为 $S_\etale$ 上的层。设
$\overline{s}'$ 为 $S'$ 的几何点，$\overline{t}'$ 为
$\Spec(\mathcal{O}^{sh}_{S', \overline{s}'})$ 的几何点。记
$\overline{s} = g(\overline{s}')$、$\overline{t} = h(\overline{t}')$，其中典范态射为
$h : \Spec(\mathcal{O}^{sh}_{S', \overline{s}'}) \to
\Spec(\mathcal{O}^{sh}_{S, \overline{s}})$。对任意层 $\mathcal{F}$（定义在
$S_\etale$ 上），特殊化映射
$$
sp :
(g^{-1}\mathcal{F})_{\overline{s}'}
\longrightarrow
(g^{-1}\mathcal{F})_{\overline{t}'}
$$
在下列等同下等于特殊化映射
$sp : \mathcal{F}_{\overline{s}} \to \mathcal{F}_{\overline{t}}$：
$(g^{-1}\mathcal{F})_{\overline{s}'} = \mathcal{F}_{\overline{s}}$ 以及
$(g^{-1}\mathcal{F})_{\overline{t}'} = \mathcal{F}_{\overline{t}}$。
这些等同来自引理 \ref{etale-cohomology-lemma-stalk-pullback}。
\end{lemma}

\begin{proof}
略。
\end{proof}

\begin{lemma}
\label{etale-cohomology-lemma-characterize-locally-constant}
设 $S$ 为概形，并且 $S$ 的每个拟紧开集都只有有限个不可约分支（例如，$S$ 的
底拓扑空间为诺特空间，或 $S$ 局部诺特）。设 $\mathcal{F}$ 为 $S_\etale$ 上的集合层。
下列条件等价：
\begin{enumerate}
\item $\mathcal{F}$ 是有限局部常值的；
\item $\mathcal{F}$ 的所有茎都是有限集，并且所有特殊化映射
$sp : \mathcal{F}_{\overline{s}} \to \mathcal{F}_{\overline{t}}$
都是双射。
\end{enumerate}
\end{lemma}

\begin{proof}
假设 (2)。设 $\overline{s}$ 为 $S$ 的几何点，位于 $s \in S$ 上方。为证明 (1)，
必须找到一个 \'etale 邻域 $(U, \overline{u})$，它是 $(S, \overline{s})$ 的邻域，并使得
$\mathcal{F}|_U$ 是常值层。可以且确实假定 $S$ 是仿射的。

\medskip\noindent
由于 $\mathcal{F}_{\overline{s}}$ 有限，可以选取 $(U, \overline{u})$、
$n \geq 0$ 以及两两不同的元素
$\sigma_1, \ldots, \sigma_n \in \mathcal{F}(U)$，使得
$\{\sigma_1, \ldots, \sigma_n\} \subset \mathcal{F}(U)$ 双射地映到
$\mathcal{F}_{\overline{s}}$，所用映射是
$\mathcal{F}(U) \to \mathcal{F}_{\overline{s}}$。考虑映射
$$
\varphi : \underline{\{1, \ldots, n\}} \longrightarrow \mathcal{F}|_U
$$
，它定义在 $U_\etale$ 上，并由 $\sigma_1, \ldots, \sigma_n$ 给出。按照构造，
该映射在 $\overline{u}$ 处的茎上是双射。考虑子集
$$
E = \{u' \in U \mid \varphi_{\overline{u}'}\text{ 是双射}\} \subset U
$$
。这里，$\overline{u}'$ 是 $U$ 的位于 $u'$ 上方的任意几何点（由注
\ref{etale-cohomology-remark-map-stalks}，该条件与选择无关）。点 $u \in U$ 是 $\overline{u}$ 的像，并
属于 $E$。由关于 $\mathcal{F}$ 的特殊化映射的假设、注
\ref{etale-cohomology-remark-can-lift} 以及引理 \ref{etale-cohomology-lemma-specialization-map-pullback}，
$E$ 在拓扑空间 $U$ 中对特殊化和一般化封闭。

\medskip\noindent
缩小 $U$ 后，还可以假定 $U$ 仿射。由下降，引理
\ref{descent-lemma-locally-finite-nr-irred-local-fppf}，$U$ 只有有限个不可约分支。
删去不经过 $u$ 的不可约分支后，可以假定 $U$ 的每个不可约分支都经过 $u$。
由于 $U$ 是 Sober 拓扑空间，故 $E = U$；再由定理
\ref{etale-cohomology-theorem-exactness-stalks} 可知 $\varphi$ 是同构。因此 (1) 成立。

\medskip\noindent
省略 (1) 蕴含 (2) 的证明。
\end{proof}

\begin{lemma}
\label{etale-cohomology-lemma-characterize-locally-constant-module}
设 $S$ 为概形，并且 $S$ 的每个拟紧开集都只有有限个不可约分支（例如，$S$ 的
底拓扑空间为诺特空间，或 $S$ 局部诺特）。设 $\Lambda$ 为诺特环。设
$\mathcal{F}$ 为 $\Lambda$-模层，定义在 $S_\etale$ 上。下列条件等价：
\begin{enumerate}
\item $\mathcal{F}$ 是有限型局部常值 $\Lambda$-模层；
\item $\mathcal{F}$ 的所有茎都是有限 $\Lambda$-模，并且所有特殊化映射
$sp : \mathcal{F}_{\overline{s}} \to \mathcal{F}_{\overline{t}}$
都是双射。
\end{enumerate}
\end{lemma}

\begin{proof}
该引理的证明与引理 \ref{etale-cohomology-lemma-characterize-locally-constant} 的证明相同。
假设 (2)。设 $\overline{s}$ 为 $S$ 的几何点，位于 $s \in S$ 上方。为证明 (1)，
必须找到一个 \'etale 邻域 $(U, \overline{u})$，它是 $(S, \overline{s})$ 的邻域，并使得
$\mathcal{F}|_U$ 是常值层。可以且确实假定 $S$ 是仿射的。

\medskip\noindent
由于 $M = \mathcal{F}_{\overline{s}}$ 是有限 $\Lambda$-模，且 $\Lambda$ 诺特，
可以选取一个表现
$$
\Lambda^{\oplus m} \xrightarrow{A} \Lambda^{\oplus n} \to M \to 0
$$
，其中 $A = (a_{ji})$ 是一个系数位于 $\Lambda$ 中的矩阵。可以选取
$(U, \overline{u})$ 以及元素
$\sigma_1, \ldots, \sigma_n \in \mathcal{F}(U)$，使得
$\sum a_{ji}\sigma_i = 0$ 在 $\mathcal{F}(U)$ 中成立，并且 $\sigma_i$ 在
$\mathcal{F}_{\overline{s}} = M$ 中的像，是 $\Lambda^n$ 的标准基元素在上述 $M$ 的表现中的
像。考虑映射
$$
\varphi : \underline{M} \longrightarrow \mathcal{F}|_U
$$
，它定义在 $U_\etale$ 上，并由 $\sigma_1, \ldots, \sigma_n$ 给出。按照构造，
该映射在 $\overline{u}$ 处的茎上是双射。考虑子集
$$
E = \{u' \in U \mid \varphi_{\overline{u}'}\text{ 是双射}\} \subset U
$$
。这里，$\overline{u}'$ 是 $U$ 的位于 $u'$ 上方的任意几何点（由注
\ref{etale-cohomology-remark-map-stalks}，该条件与选择无关）。点 $u \in U$ 是 $\overline{u}$ 的像，并
属于 $E$。由关于 $\mathcal{F}$ 的特殊化映射的假设、注
\ref{etale-cohomology-remark-can-lift} 以及引理 \ref{etale-cohomology-lemma-specialization-map-pullback}，
$E$ 在拓扑空间 $U$ 中对特殊化和一般化封闭。

\medskip\noindent
缩小 $U$ 后，还可以假定 $U$ 仿射。由下降，引理
\ref{descent-lemma-locally-finite-nr-irred-local-fppf}，$U$ 只有有限个不可约分支。
删去不经过 $u$ 的不可约分支后，可以假定 $U$ 的每个不可约分支都经过 $u$。
由于 $U$ 是 Sober 拓扑空间，故 $E = U$；再由定理
\ref{etale-cohomology-theorem-exactness-stalks} 可知 $\varphi$ 是同构。因此 (1) 成立。

\medskip\noindent
省略 (1) 蕴含 (2) 的证明。
\end{proof}

\begin{lemma}
\label{etale-cohomology-lemma-specialization-map-pushforward}
设 $f : X \to S$ 为概形的拟紧拟分离态射。设 $K \in D^+(X_\etale)$。
设 $\overline{s}$ 为 $S$ 的几何点，$\overline{t}$ 为
$\Spec(\mathcal{O}^{sh}_{S, \overline{s}})$ 的几何点。则有交换图
$$
\xymatrix{
(Rf_*K)_{\overline{s}} \ar[r]_{sp} \ar@{=}[d] &
(Rf_*K)_{\overline{t}} \ar@{=}[d] \\
R\Gamma(X \times_S \Spec(\mathcal{O}^{sh}_{S, \overline{s}}), K)
\ar[r] &
R\Gamma(X \times_S \Spec(\mathcal{O}^{sh}_{S, \overline{t}}), K)
}
$$
，其中下方水平箭头来自沿态射 $\text{id}_X \times c$ 的拉回，而
$c : \Spec(\mathcal{O}^{sh}_{S, \overline{t}}) \to
\Spec(\mathcal{O}^{sh}_{S, \overline{s}})$ 是注 \ref{etale-cohomology-remark-another-sp}
中引入的态射。竖直箭头由定理 \ref{etale-cohomology-theorem-higher-direct-images} 给出。
\end{lemma}

\begin{proof}
这立即由注 \ref{etale-cohomology-remark-another-sp} 中对 $sp$ 的描述得出。
\end{proof}

\begin{remark}
\label{etale-cohomology-remark-specialization-map-and-fibres}
设 $f : X \to S$ 为概形态射，$K \in D(X_\etale)$。设 $\overline{s}$ 为 $S$
的几何点，$\overline{t}$ 为 $\Spec(\mathcal{O}^{sh}_{S, \overline{s}})$ 的几何点。
设 $c$ 如注 \ref{etale-cohomology-remark-another-sp} 中所述。总可以作出交换图
$$
\xymatrix{
(Rf_*K)_{\overline{s}} \ar[r] \ar[d]_{sp} &
R\Gamma(X \times_S \Spec(\mathcal{O}_{S, \overline{s}}^{sh}), K)
\ar[r] \ar[d]_{(\text{id}_X \times c)^{-1}} &
R\Gamma(X_{\overline{s}}, K) \\
(Rf_*K)_{\overline{t}} \ar[r] &
R\Gamma(X \times_S \Spec(\mathcal{O}_{S, \overline{t}}^{sh}), K) \ar[r] &
R\Gamma(X_{\overline{t}}, K)
}
$$
，其中水平箭头是注 \ref{etale-cohomology-remark-stalk-fibre} 中的箭头。一般而言，右侧并不存在一个
夹在此图中的竖直映射，它连接 $K$ 在诸纤维上的上同调；即使在引理
\ref{etale-cohomology-lemma-specialization-map-pushforward} 的情形中也是如此。
\end{remark}















\section{具有可构造上同调的复形}
\label{etale-cohomology-section-Dc}

\noindent
设 $\Lambda$ 为环。记 $D(X_\etale, \Lambda)$ 为 $\Lambda$-模层（定义在
$X_\etale$ 上）的导出范畴。记 $D^b(X_\etale, \Lambda)$（分别记
$D^+$、$D^-$）为 $D(X_\etale, \Lambda)$ 中有界（分别为上有界、下有界）
复形所成的充满子范畴。

\begin{definition}
\label{etale-cohomology-definition-c}
设 $X$ 为概形，$\Lambda$ 为诺特环。记
{\it $D_c(X_\etale, \Lambda)$} 为 $D(X_\etale, \Lambda)$ 的充满子范畴，
其对象是上同调层均为可构造 $\Lambda$-模层的复形。
\end{definition}

\noindent
由引理 \ref{etale-cohomology-lemma-constructible-abelian} 及导出范畴，第
\ref{derived-section-triangulated-sub} 节，此定义是合理的。因此，
$D_c(X_\etale, \Lambda)$ 是 $D(X_\etale, \Lambda)$ 的严格充满、饱和三角子范畴。

\begin{lemma}
\label{etale-cohomology-lemma-restrict-and-shriek-from-etale-c}
设 $\Lambda$ 为诺特环。若 $j : U \to X$ 为概形的 \'etale 态射，则
\begin{enumerate}
\item 结论 $K|_U \in D_c(U_\etale, \Lambda)$ 在
$K \in D_c(X_\etale, \Lambda)$ 时成立；
\item 结论 $j_!M \in D_c(X_\etale, \Lambda)$ 在
$M \in D_c(U_\etale, \Lambda)$ 且态射 $j$ 拟紧拟分离时成立。
\end{enumerate}
\end{lemma}

\begin{proof}
第一条是显然的。第二条由 $j_!$ 正合及引理
\ref{etale-cohomology-lemma-jshriek-constructible} 得出。
\end{proof}

\begin{lemma}
\label{etale-cohomology-lemma-pullback-c}
设 $\Lambda$ 为诺特环，$f : X \to Y$ 为概形态射。若
$K \in D_c(Y_\etale, \Lambda)$，则 $Lf^*K \in D_c(X_\etale, \Lambda)$。
\end{lemma}

\begin{proof}
这是因为 $f^{-1} = f^*$ 正合，再由引理
\ref{etale-cohomology-lemma-pullback-constructible} 即得。
\end{proof}

\begin{lemma}
\label{etale-cohomology-lemma-one-constructible}
设 $X$ 为拟紧拟分离概形，$\Lambda$ 为诺特环。设
$K \in D(X_\etale, \Lambda)$、$b \in \mathbf{Z}$，且 $H^b(K)$ 可构造。
则存在层 $\mathcal{F}$，它是有限多个 $j_{U!}\underline{\Lambda}$ 的直和，
其中 $U \in \Ob(X_\etale)$ 为仿射对象；并且存在映射
$\mathcal{F}[-b] \to K$，它位于 $D(X_\etale, \Lambda)$ 中，并诱导满射
$\mathcal{F} \to H^b(K)$。
\end{lemma}

\begin{proof}
将 $K$ 表示为复形 $\mathcal{K}^\bullet$，该复形由 $\Lambda$-模层组成。考虑满射
$$
\Ker(\mathcal{K}^b \to \mathcal{K}^{b + 1})
\longrightarrow
H^b(K)
$$
。由位点上的模，引理
\ref{sites-modules-lemma-module-quotient-direct-sum}
，可以选取满射
$\bigoplus_{i \in I} j_{U_i!} \underline{\Lambda} \to
\Ker(\mathcal{K}^b \to \mathcal{K}^{b + 1})$
，其中 $U_i$ 仿射。对有限子集 $I' \subset I$，记
$\mathcal{H}_{I'} \subset H^b(K)$ 为
$\bigoplus_{i \in I'} j_{U_i!} \underline{\Lambda}$ 的像。由引理
\ref{etale-cohomology-lemma-colimit-constructible}，可知 $\mathcal{H}_{I'} = H^b(K)$ 对某个有限子集
$I' \subset I$ 成立。取
$\mathcal{F} = \bigoplus_{i \in I'} j_{U_i!} \underline{\Lambda}$ 即得结论。
\end{proof}

\begin{lemma}
\label{etale-cohomology-lemma-bounded-above-c}
设 $X$ 为拟紧拟分离概形，$\Lambda$ 为诺特环，并设
$K \in D^-(X_\etale, \Lambda)$。下列条件等价：
\begin{enumerate}
\item $K$ 属于 $D_c(X_\etale, \Lambda)$；
\item $K$ 可由一个上有界复形表示，该复形的每一项都是有限多个
$j_{U!}\underline{\Lambda}$ 的直和，其中 $U \in \Ob(X_\etale)$ 仿射；
\item $K$ 可由一个上有界的平坦可构造 $\Lambda$-模层复形表示。
\end{enumerate}
\end{lemma}

\begin{proof}
(2) 蕴含 (3)、(3) 蕴含 (1) 都是显然的。假设
$K$ 属于 $D_c^-(X_\etale, \Lambda)$。设 $H^i(K) = 0$ 对所有 $i > b$ 成立。
我们对 $a$ 归纳构造复形 $\mathcal{F}^a \to \ldots \to \mathcal{F}^b$，
使得每个 $\mathcal{F}^i$ 都是有限多个 $j_{U!}\underline{\Lambda}$ 的直和，
其中 $U \in \Ob(X_\etale)$ 仿射；还构造映射
$\mathcal{F}^\bullet \to K$，它诱导同构
$H^i(\mathcal{F}^\bullet) \to H^i(K)$（对 $i > a$），并诱导满射
$H^a(\mathcal{F}^\bullet) \to H^a(K)$。当 $a = b$ 时，由引理
\ref{etale-cohomology-lemma-one-constructible} 可以完成此构造。给定这样的数据，选取优三角
$$
\mathcal{F}^\bullet \to K \to L \to \mathcal{F}^\bullet[1]
$$
。于是 $H^i(L) = 0$ 对所有 $i \geq a$ 成立。按引理
\ref{etale-cohomology-lemma-one-constructible}，选取
$\mathcal{F}^{a - 1}[-a +1] \to L$。复合
$\mathcal{F}^{a - 1}[-a +1] \to L \to \mathcal{F}^\bullet$
对应于映射 $\mathcal{F}^{a - 1} \to \mathcal{F}^a$，且它与
$\mathcal{F}^a \to \mathcal{F}^{a + 1}$ 的复合为零。由 TR4 得到映射
$$
(\mathcal{F}^{a - 1} \to \ldots \to \mathcal{F}^b) \to K
$$
，它位于 $D(X_\etale, \Lambda)$ 中。这就完成了归纳步骤及引理的证明。
\end{proof}

\begin{lemma}
\label{etale-cohomology-lemma-tensor-c}
设 $X$ 为概形，$\Lambda$ 为诺特环，并设
$K, L \in D_c^-(X_\etale, \Lambda)$。则
$K \otimes_\Lambda^\mathbf{L} L$ 属于 $D_c^-(X_\etale, \Lambda)$。
\end{lemma}

\begin{proof}
这由引理 \ref{etale-cohomology-lemma-bounded-above-c} 和
\ref{etale-cohomology-lemma-tensor-product-constructible} 得出。
\end{proof}




\section{具有可构造上同调的有限 Tor 维数复形}
\label{etale-cohomology-section-ctf}

\noindent
设 $X$ 为概形，$\Lambda$ 为诺特环。导出范畴 $D_c(X_\etale, \Lambda)$
（定义见第 \ref{etale-cohomology-section-Dc} 节）的一个常用子范畴，是由（局部）具有有限
Tor 维数的对象组成的充满子范畴。下面给出正式定义。

\begin{definition}
\label{etale-cohomology-definition-ctf}
设 $X$ 为概形，$\Lambda$ 为诺特环。记
{\it $D_{ctf}(X_\etale, \Lambda)$} 为 $D_c(X_\etale, \Lambda)$ 的充满子范畴，
其对象局部具有有限 Tor 维数。
\end{definition}

\noindent
这是 $D_c(X_\etale, \Lambda)$ 和 $D(X_\etale, \Lambda)$ 的严格充满、饱和
三角子范畴。按照我们的约定，见位点上的上同调，定义
\ref{sites-cohomology-definition-tor-amplitude}，有
$$
D_{ctf}(X_\etale, \Lambda) \subset D^b_c(X_\etale, \Lambda)
\subset D^b(X_\etale, \Lambda)
$$
，只要 $X$ 拟紧。关于 $D_{ctf}(X_\etale, \Lambda)$ 中的对象，引理
\ref{etale-cohomology-lemma-when-ctf} 给出了一种很好的理解方式。

\begin{remark}
\label{etale-cohomology-remark-different}
从某种意义上说，导出范畴 $D_{ctf}(X_\etale, \Lambda)$ 中的对象，比
$D(\mathcal{O}_X)$ 中的完美对象具有更好的整体性质。具体而言，一个
$\mathcal{O}_X$-模复形可能局部拟同构于有限局部自由
$\mathcal{O}_X$-模的有限复形，却不整体拟同构于局部自由
$\mathcal{O}_X$-模的有界复形。下面的引理说明，在诺特概形上的
$D_{ctf}$ 中不会出现这种现象。
\end{remark}

\begin{lemma}
\label{etale-cohomology-lemma-when-ctf}
设 $\Lambda$ 为诺特环，$X$ 为拟紧拟分离概形，并设
$K \in D(X_\etale, \Lambda)$。下列条件等价：
\begin{enumerate}
\item $K \in D_{ctf}(X_\etale, \Lambda)$；
\item $K$ 可由可构造平坦 $\Lambda$-模层的有限复形表示。
\end{enumerate}
事实上，若 $K$ 的 Tor 振幅落在 $[a, b]$ 中，则可将 $K$ 表示为复形
$\mathcal{F}^a \to \ldots \to \mathcal{F}^b$，其中
$\mathcal{F}^p$ 是可构造平坦 $\Lambda$-模层。
\end{lemma}

\begin{proof}
可构造平坦 $\Lambda$-模层的有限复形具有有限 Tor 维数，这是显然的。
它属于 $D_c(X_\etale, \Lambda)$ 也是显然的。因此 (2) 蕴含 (1)。

\medskip\noindent
假设 (1)。选取 $a, b \in \mathbf{Z}$，使得
$H^i(K \otimes_\Lambda^\mathbf{L} \mathcal{G}) = 0$ if
$i \not \in [a, b]$，并且这对所有 $\Lambda$-模层 $\mathcal{G}$ 都成立。
我们对 $b - a$ 归纳证明最后一个断言。若
$a = b$，则 $K = H^a(K)[-a]$ 是平坦可构造层，结论成立。接着假设
$b > a$。将 $K$ 表示为复形 $\mathcal{K}^\bullet$，其项为 $\Lambda$-模层。
考虑满射
$$
\Ker(\mathcal{K}^b \to \mathcal{K}^{b + 1})
\longrightarrow
H^b(K)
$$
。由引理 \ref{etale-cohomology-lemma-category-constructible-modules}，可找到有限多个仿射概形
$U_i$，它们在 $X$ 上 \'etale；还可找到满射
$\bigoplus j_{U_i!}\underline{\Lambda}_{U_i} \to H^b(K)$。把 $U_i$ 换成标准
\'etale 覆盖 $\{U_{ij} \to U_i\}$ 后，可以假定该满射提升为映射
$\mathcal{F} = \bigoplus j_{U_i!}\underline{\Lambda}_{U_i} \to
\Ker(\mathcal{K}^b \to \mathcal{K}^{b + 1})$.
该映射决定优三角
$$
\mathcal{F}[-b] \to K \to L \to \mathcal{F}[-b + 1]
$$
，它位于 $D(X_\etale, \Lambda)$ 中。由于 $D_{ctf}(X_\etale, \Lambda)$ 是三角
子范畴，$L$ 也属于其中。事实上，$L$ 的 Tor 振幅落在 $[a, b - 1]$ 中，因为
$\mathcal{F}$ 满射到 $H^b(K)$（细节略）。由归纳假设，可以找到有限复形
$\mathcal{F}^a \to \ldots \to \mathcal{F}^{b - 1}$，它由可构造平坦
$\Lambda$-模层组成，并表示 $L$。映射
$L \to \mathcal{F}[-b + 1]$ 对应于映射
$\mathcal{F}^b \to \mathcal{F}$，后者在
$\mathcal{F}^{b - 1} \to \mathcal{F}^b$ 的像上为零。于是由公理 TR3，$K$ 可由复形
$$
\mathcal{F}^a \to \ldots \to \mathcal{F}^{b - 1} \to \mathcal{F}^b
$$
表示，从而完成证明。
\end{proof}

\begin{remark}
\label{etale-cohomology-remark-projective-each-degree}
设 $\Lambda$ 为诺特环，$X$ 为概形。对有界复形 $K^\bullet$，它由可构造平坦
$\Lambda$-模层组成并定义在 $X_\etale$ 上，每个茎
$K^p_{\overline{x}}$ 都是有限投射 $\Lambda$-模。因此，该复形的各个茎都是
$\Lambda$-模的完美复形。
\end{remark}

\begin{lemma}
\label{etale-cohomology-lemma-restrict-and-shriek-from-etale-ctf}
设 $\Lambda$ 为诺特环。若 $j : U \to X$ 为概形的 \'etale 态射，则
\begin{enumerate}
\item 结论 $K|_U \in D_{ctf}(U_\etale, \Lambda)$ 在
$K \in D_{ctf}(X_\etale, \Lambda)$ 时成立；
\item 结论 $j_!M \in D_{ctf}(X_\etale, \Lambda)$ 在
$M \in D_{ctf}(U_\etale, \Lambda)$ 且态射 $j$ 拟紧拟分离时成立。
\end{enumerate}
\end{lemma}

\begin{proof}
证明本引理或许最容易的方法，是先约化到 $X$ 仿射的情形，再应用引理
\ref{etale-cohomology-lemma-when-ctf}，把它转化为关于可构造平坦 $\Lambda$-模层有限复形的断言；
此时结论由引理 \ref{etale-cohomology-lemma-jshriek-constructible} 得出。
\end{proof}

\begin{lemma}
\label{etale-cohomology-lemma-pullback-ctf}
设 $\Lambda$ 为诺特环，$f : X \to Y$ 为概形态射。若
$K \in D_{ctf}(Y_\etale, \Lambda)$，则
$Lf^*K \in D_{ctf}(X_\etale, \Lambda)$。
\end{lemma}

\begin{proof}
应用引理 \ref{etale-cohomology-lemma-when-ctf}，把问题约化为关于可构造平坦
$\Lambda$-模层有限复形的问题。于是，由 $f^{-1} = f^*$ 正合及引理
\ref{etale-cohomology-lemma-pullback-constructible} 即得结论。
\end{proof}

\begin{lemma}
\label{etale-cohomology-lemma-connected-ctf-locally-constant}
设 $X$ 为连通概形，$\Lambda$ 为诺特环。设
$K \in D_{ctf}(X_\etale, \Lambda)$ 具有局部常值上同调层。则存在有限投射
$\Lambda$-模的有限复形 $M^\bullet$ 以及 \'etale 覆盖 $\{U_i \to X\}$，使得
$K|_{U_i} \cong \underline{M^\bullet}|_{U_i}$ 在 $D(U_{i, \etale}, \Lambda)$ 中成立。
\end{lemma}

\begin{proof}
选取 \'etale 覆盖 $\{U_i \to X\}$，使得 $K|_{U_i}$ 为常值复形，比如
$K|_{U_i} \cong \underline{M_i^\bullet}_{U_i}$；这里所用有限复形由有限
$\Lambda$-模组成，记为 $M_i^\bullet$。见位点上的上同调，引理
\ref{sites-cohomology-lemma-locally-constant}.
注意，$U_i \times_X U_j$ 在 $M_i^\bullet$ 与 $M_j^\bullet$ 于
$D(\Lambda)$ 中不同构时为空。对每个 $\Lambda$-模复形 $M^\bullet$，令
$I_{M^\bullet} =
\{i \in I \mid M_i^\bullet \cong M^\bullet\text{ 于 }D(\Lambda)\}$.
由于 \'etale 态射是开映射，故
$U_{M^\bullet} = \bigcup_{i \in I_{M^\bullet}} \Im(U_i \to X)$
是 $X$ 的开子集。于是 $X = \coprod U_{M^\bullet}$ 是 $X$ 的不交开覆盖。
由于 $X$ 连通，只有一个 $U_{M^\bullet}$ 非空。又因为 $K$ 属于
$D_{ctf}(X_\etale, \Lambda)$，可知 $M^\bullet$ 是 $\Lambda$-模的完美复形；
见代数的进一步结果，引理 \ref{more-algebra-lemma-perfect}。因此，可以假定
$M^\bullet$ 是有限投射 $\Lambda$-模的有限复形。
\end{proof}








\section{挠层}
\label{etale-cohomology-section-torsion}

\noindent
本节简要讨论挠阿贝尔层及其 \'etale 上同调。设 $\mathcal{C}$ 为位点。我们已经在
位点上的上同调，引理 \ref{sites-cohomology-lemma-torsion} 中证明：
$D(\mathcal{C})$ 中上同调层均为挠层的任何对象，都可由每一项皆为挠层的复形表示。

\begin{lemma}
\label{etale-cohomology-lemma-torsion-cohomology}
设 $X$ 为拟紧拟分离概形。
\begin{enumerate}
\item 若 $\mathcal{F}$ 是 $X_\etale$ 上的挠阿贝尔层，则
$H^n_\etale(X, \mathcal{F})$ 对所有 $n$ 都是挠阿贝尔群。
\item 若 $K$ 属于 $D^+(X_\etale)$ 且具有挠上同调层，则
$H^n_\etale(X, K)$ 对所有 $n$ 都是挠阿贝尔群。
\end{enumerate}
\end{lemma}

\begin{proof}
为证明 (1)，写成 $\mathcal{F} = \bigcup \mathcal{F}[n]$，其中
$\mathcal{F}[d]$ 是 $d$-挠子层。由引理 \ref{etale-cohomology-lemma-colimit}，有
$H^n_\etale(X, \mathcal{F}) = \colim H^n_\etale(X, \mathcal{F}[d])$.
由于 $H^n_\etale(X, \mathcal{F}[d])$ 被 $d$ 零化，这就证明了 (1)。

\medskip\noindent
为证明 (2)，可以使用谱序列
$E_2^{p, q} = H^p_\etale(X, H^q(K))$，它收敛到 $H^n_\etale(X, K)$
（导出范畴，引理 \ref{derived-lemma-two-ss-complex-functor}），以及对层所得的结论。
\end{proof}

\begin{lemma}
\label{etale-cohomology-lemma-torsion-direct-image}
设 $f : X \to Y$ 为概形的拟紧拟分离态射。
\begin{enumerate}
\item 若 $\mathcal{F}$ 是 $X_\etale$ 上的挠阿贝尔层，则
$R^nf_*\mathcal{F}$ 是 $Y_\etale$ 上的挠阿贝尔层，对所有 $n$ 都成立。
\item 若 $K$ 属于 $D^+(X_\etale)$ 且具有挠上同调层，则
$Rf_*K$ 是 $D^+(Y_\etale)$ 的对象，并且其上同调层均为挠阿贝尔层。
\end{enumerate}
\end{lemma}

\begin{proof}
(1) 的证明。回忆 $R^nf_*\mathcal{F}$ 是与预层
$V \mapsto H^n_\etale(X \times_Y V, \mathcal{F})$ 相伴的层，该预层定义在
$Y_\etale$ 上。见位点上的上同调，引理
\ref{sites-cohomology-lemma-higher-direct-images}。若选取仿射的 $V$，则
$X \times_Y V$ 因 $f$ 拟紧拟分离而拟紧拟分离，故可应用引理
\ref{etale-cohomology-lemma-torsion-cohomology}，看出
$H^n_\etale(X \times_Y V, \mathcal{F})$ 是挠群。

\medskip\noindent
(2) 的证明。回忆 $R^nf_*K$ 是与预层
$V \mapsto H^n_\etale(X \times_Y V, K)$ 相伴的层，该预层定义在
$Y_\etale$ 上。见位点上的上同调，引理
\ref{sites-cohomology-lemma-unbounded-describe-higher-direct-images}。若选取仿射的
$V$，则 $X \times_Y V$ 因 $f$ 拟紧拟分离而拟紧拟分离，故可应用引理
\ref{etale-cohomology-lemma-torsion-cohomology}，看出
$H^n_\etale(X \times_Y V, K)$ 是挠群。
\end{proof}









\section{支撑在闭子概形上的上同调}
\label{etale-cohomology-section-cohomology-support}

\noindent
设 $X$ 是一个概形，且 $Z \subset X$ 是一个闭子概形。
设 $\mathcal{F}$ 是 $X_\etale$ 上的一个阿贝尔层。令
$$
\Gamma_Z(X, \mathcal{F}) =
\{s \in \mathcal{F}(X) \mid \text{Supp}(s) \subset Z\}
$$
为支撑在 $Z$ 上的截面（定义 \ref{etale-cohomology-definition-support}）。
这是一个左正合函子，但一般并不正合。
因此得到导出函子
$$
R\Gamma_Z(X, -) : D(X_\etale) \longrightarrow D(\textit{Ab})
$$
并把支撑在 $Z$ 上的上同调群定义为
$H^q_Z(X, \mathcal{F}) = R^q\Gamma_Z(X, \mathcal{F})$。

\medskip\noindent
设 $\mathcal{I}$ 是 $X_\etale$ 上的一个内射阿贝尔层。
令 $U = X \setminus Z$。
则限制映射 $\mathcal{I}(X) \to \mathcal{I}(U)$ 是满射
（位点上的上同调，引理
\ref{sites-cohomology-lemma-restriction-along-monomorphism-surjective})
，其核为 $\Gamma_Z(X, \mathcal{I})$。立即可知，对
$K \in D(X_\etale)$ 有优三角
$$
R\Gamma_Z(X, K) \to R\Gamma(X, K) \to R\Gamma(U, K) \to R\Gamma_Z(X, K)[1]
$$
，它位于 $D(\textit{Ab})$ 中。由此得到上同调长正合列
$$
\ldots \to H^i_Z(X, K) \to H^i(X, K) \to H^i(U, K) \to
H^{i + 1}_Z(X, K) \to \ldots
$$
，对任意 $K$（作为 $D(X_\etale)$ 中的对象）都成立。

\medskip\noindent
对阿贝尔层 $\mathcal{F}$（定义在 $X_\etale$ 上），可以考虑{\it 支撑在
$Z$ 上的截面子层}，记作 $\mathcal{H}_Z(\mathcal{F})$，其定义规则为
$$
\mathcal{H}_Z(\mathcal{F})(U) =
\{s \in \mathcal{F}(U) \mid \text{Supp}(s) \subset U \times_X Z\}
$$
这里使用定义 \ref{etale-cohomology-definition-support} 中截面的支撑。利用命题
\ref{etale-cohomology-proposition-closed-immersion-pushforward} 的等价，可以把
$\mathcal{H}_Z(\mathcal{F})$ 看作 $Z_\etale$ 上的一个阿贝尔层。
于是得到函子
$$
\textit{Ab}(X_\etale) \longrightarrow \textit{Ab}(Z_\etale),\quad
\mathcal{F} \longmapsto \mathcal{H}_Z(\mathcal{F})
$$
，它是左正合的，但一般并不正合。

\begin{lemma}
\label{etale-cohomology-lemma-sections-with-support-acyclic}
设 $i : Z \to X$ 是概形的一个闭浸入。
设 $\mathcal{I}$ 是 $X_\etale$ 上的一个内射阿贝尔层。
则 $\mathcal{H}_Z(\mathcal{I})$ 是 $Z_\etale$ 上的一个内射阿贝尔层。
\end{lemma}

\begin{proof}
注意，对任意阿贝尔层 $\mathcal{G}$（定义在 $Z_\etale$ 上），有
$$
\Hom_Z(\mathcal{G}, \mathcal{H}_Z(\mathcal{F})) =
\Hom_X(i_*\mathcal{G}, \mathcal{F})
$$
，因为 $i_*\mathcal{G}$ 的任意截面都支撑在 $Z$ 上。由于
$i_*$ 是正合的（第 \ref{etale-cohomology-section-closed-immersions} 节），而
$\mathcal{I}$ 在 $X_\etale$ 上是内射的，所以
$\mathcal{H}_Z(\mathcal{I})$ 在 $Z_\etale$ 上是内射的。
\end{proof}

\noindent
将导出函子记作
$$
R\mathcal{H}_Z : D(X_\etale) \longrightarrow D(Z_\etale)
$$
。置
$\mathcal{H}^q_Z(\mathcal{F}) = R^q\mathcal{H}_Z(\mathcal{F})$，于是
$\mathcal{H}^0_Z(\mathcal{F}) = \mathcal{H}_Z(\mathcal{F})$。
由上面的引理，有 Grothendieck 谱序列
$$
E_2^{p, q} = H^p(Z, \mathcal{H}^q_Z(\mathcal{F}))
\Rightarrow H^{p + q}_Z(X, \mathcal{F})
$$

\begin{lemma}
\label{etale-cohomology-lemma-cohomology-with-support-sheaf-on-support}
设 $i : Z \to X$ 是概形的一个闭浸入。
设 $\mathcal{G}$ 是 $Z_\etale$ 上的一个内射阿贝尔层。
则 $\mathcal{H}^p_Z(i_*\mathcal{G}) = 0$ 对 $p > 0$ 成立。
\end{lemma}

\begin{proof}
这是因为函子 $i_*$ 正合，并且把内射阿贝尔层变为内射阿贝尔层
（位点上的上同调，引理
\ref{sites-cohomology-lemma-pushforward-injective-flat}）。
\end{proof}

\begin{lemma}
\label{etale-cohomology-lemma-cohomology-with-support-triangle}
设 $i : Z \to X$ 是概形的一个闭浸入。
设 $j : U \to X$ 是 $Z$ 的补集的包含映射。
设 $\mathcal{F}$ 是 $X_\etale$ 上的一个阿贝尔层。
有优三角
$$
i_*R\mathcal{H}_Z(\mathcal{F}) \to \mathcal{F} \to Rj_*(\mathcal{F}|_U) \to
i_*R\mathcal{H}_Z(\mathcal{F})[1]
$$
，它位于 $D(X_\etale)$ 中。由此得到正合列
$$
0 \to i_*\mathcal{H}_Z(\mathcal{F}) \to \mathcal{F} \to
j_*(\mathcal{F}|_U) \to i_*\mathcal{H}^1_Z(\mathcal{F}) \to 0
$$
以及同构
$R^pj_*(\mathcal{F}|_U) \cong i_*\mathcal{H}^{p + 1}_Z(\mathcal{F})$
，其中 $p \geq 1$。
\end{lemma}

\begin{proof}
为得到这个优三角，选取内射分解
$\mathcal{F} \to \mathcal{I}^\bullet$。于是由上面的讨论得到复形的短正合列
$$
0 \to
i_*\mathcal{H}_Z(\mathcal{I}^\bullet) \to \mathcal{I}^\bullet
\to j_*(\mathcal{I}^\bullet|_U) \to 0
$$
。因此由导出范畴，第
\ref{derived-section-canonical-delta-functor} 节，得到上述优三角。
\end{proof}

\noindent
设 $X$ 是一个概形，且 $Z \subset X$ 是一个闭子概形。用
$D_Z(X_\etale)$ 表示 $D(X_\etale)$ 的严格充满、饱和三角子范畴，
其对象是上同调层支撑在 $Z$ 上的复形。注意，
$D_Z(X_\etale)$ 仅依赖于 $X$ 的底闭子集。

\begin{lemma}
\label{etale-cohomology-lemma-complexes-with-support-on-closed}
设 $i : Z \to X$ 是概形的一个闭浸入。函子
$Ri_{small, *} = i_{small, *} : D(Z_\etale) \to D(X_\etale)$
诱导等价 $D(Z_\etale) \to D_Z(X_\etale)$，其拟逆为
$$
i_{small}^{-1}|_{D_Z(X_\etale)} = R\mathcal{H}_Z|_{D_Z(X_\etale)}
$$
\end{lemma}

\begin{proof}
回忆 $i_{small}^{-1}$ 与 $i_{small, *}$ 构成一对正合函子的伴随，
并且 $i_{small}^{-1}i_{small, *}$ 同构于阿贝尔层范畴上的恒等函子。
见命题 \ref{etale-cohomology-proposition-closed-immersion-pushforward} 和引理
\ref{etale-cohomology-lemma-stalk-pullback}。因此
$i_{small, *} : D(Z_\etale) \to D_Z(X_\etale)$ 充分忠实，而
$i_{small}^{-1}$ 给出一个左逆。另一方面，设 $K$ 是
$D_Z(X_\etale)$ 的一个对象，并考虑伴随映射
$K \to i_{small, *}i_{small}^{-1}K$。利用 $i_{small, *}$ 和
$i_{small}^{-1}$ 的正合性，它在上同调层上诱导伴随映射
$H^n(K) \to i_{small, *}i_{small}^{-1}H^n(K)$。由于这些上同调层
支撑在 $Z$ 上，可知这些伴随映射都是同构，因而
$D(Z_\etale) \to D_Z(X_\etale)$ 是一个等价。

\medskip\noindent
为完成证明，还须说明 $R\mathcal{H}_Z(K) = i_{small}^{-1}K$，
其中 $K$ 是 $D_Z(X_\etale)$ 的对象。为此可以使用刚刚证明的等式
$K = i_{small, *}i_{small}^{-1}K$。于是可以选取一个 K-内射代表
$\mathcal{I}^\bullet$ 来表示 $i_{small}^{-1}K$。
由于 $i_{small, *}$ 是正合函子 $i_{small}^{-1}$ 的右伴随，复形
$i_{small, *}\mathcal{I}^\bullet$ 是 K-内射的
（导出范畴，引理 \ref{derived-lemma-adjoint-preserve-K-injectives}）。
可见 $R\mathcal{H}_Z(K)$ 由
$\mathcal{H}_Z(i_{small, *}\mathcal{I}^\bullet) = \mathcal{I}^\bullet$
计算，这正是所需结论。
\end{proof}

\begin{lemma}
\label{etale-cohomology-lemma-cohomology-with-support-quasi-coherent}
设 $X$ 是一个概形，$Z \subset X$ 是一个闭子概形。
设 $\mathcal{F}$ 是一个拟凝聚 $\mathcal{O}_X$-模，并用
$\mathcal{F}^a$ 表示 $X$ 的小 \'etale 网站上的伴随拟凝聚层
（命题 \ref{etale-cohomology-proposition-quasi-coherent-sheaf-fpqc}）。则
\begin{enumerate}
\item $H^q_Z(X, \mathcal{F})$ 与 $H^q_Z(X_\etale, \mathcal{F}^a)$ 一致；
\item 若 $Z$ 的补集在 $X$ 中是逆紧致的，则
$i_*\mathcal{H}^q_Z(\mathcal{F}^a)$ 是一个拟凝聚
$\mathcal{O}_X$-模层，并且等于 $(i_*\mathcal{H}^q_Z(\mathcal{F}))^a$。
\end{enumerate}
\end{lemma}

\begin{proof}
设 $j : U \to X$ 是 $Z$ 的补集的包含映射。关于上同调群的陈述 (1)
由带支撑上同调的长正合列以及等式
$H^q(X_\etale, \mathcal{F}^a) = H^q(X, \mathcal{F})$ 和
$H^q(U_\etale, \mathcal{F}^a) = H^q(U, \mathcal{F})$ 得出，见定理
\ref{etale-cohomology-theorem-zariski-fpqc-quasi-coherent}。若
$j : U \to X$ 是拟紧态射，即若 $U \subset X$ 是逆紧致的，则
$R^qj_*$ 把拟凝聚层变为拟凝聚层
（概形的上同调，引理
\ref{coherent-lemma-quasi-coherence-higher-direct-images})
，并且与在 \'etale 网站上取伴随层可交换
（下降，引理 \ref{descent-lemma-higher-direct-images-small-etale}）。
应用引理 \ref{etale-cohomology-lemma-cohomology-with-support-triangle} 即得结论。
\end{proof}







\section{具有严格 Hensel 局部环的概形}
\label{etale-cohomology-section-strictly-henselian-local-rings}

\noindent
本节汇集一些关于局部环严格 Hensel 的概形之 \'etale 上同调的结果。
例如，下面是引理
\ref{etale-cohomology-lemma-vanishing-etale-cohomology-strictly-henselian}
的一个颇有趣的推广。

\begin{lemma}
\label{etale-cohomology-lemma-local-rings-strictly-henselian}
设 $S$ 是一个所有局部环都严格 Hensel 的概形。
则对任意阿贝尔层 $\mathcal{F}$（定义在 $S_\etale$ 上），有
$H^i(S_\etale, \mathcal{F}) = H^i(S_{Zar}, \mathcal{F})$。
\end{lemma}

\begin{proof}
设 $\epsilon : S_\etale \to S_{Zar}$ 是由包含函子给出的网站态射。
Zariski 层 $R^p\epsilon_*\mathcal{F}$ 是预层
$U \mapsto H^p_\etale(U, \mathcal{F})$ 的伴随层。因此，它在
$x \in X$ 处的茎为
$\colim H^p_\etale(U, \mathcal{F}) =
H^p_\etale(\Spec(\mathcal{O}_{X, x}), \mathcal{G}_x)$
，其中 $\mathcal{G}_x$ 表示 $\mathcal{F}$ 到
$\Spec(\mathcal{O}_{X, x})$ 的拉回，见引理
\ref{etale-cohomology-lemma-directed-colimit-cohomology}。因此，由引理
\ref{etale-cohomology-lemma-vanishing-etale-cohomology-strictly-henselian}，
高阶正像 $R^p\epsilon_*\mathcal{F}$ 为零；再由 Leray 谱序列即得结论。
\end{proof}

\begin{lemma}
\label{etale-cohomology-lemma-gabber-criterion}
设 $R$ 是一个所有局部环都严格 Hensel 的环。
设 $\mathcal{F}$ 是 $\Spec(R)_\etale$ 上的一个层。
假设对所有 $f, g \in R$，映射
$$
H^1_\etale(D(f + g), \mathcal{F})
\longrightarrow
H^1_\etale(D(f(f + g)), \mathcal{F}) \oplus
H^1_\etale(D(g(f + g)), \mathcal{F})
$$
的核为零。则 $H^q_\etale(\Spec(R), \mathcal{F}) = 0$ 对 $q > 0$ 成立。
\end{lemma}

\begin{proof}
由引理 \ref{etale-cohomology-lemma-local-rings-strictly-henselian} 可知，
$\mathcal{F}$ 的 \'etale 上同调在 $\Spec(R)$ 的任意开子集上
都与 Zariski 上同调一致。我们对 $i$ 归纳证明如下陈述：若
$h \in R$，则 $H^q_\etale(D(h), \mathcal{F}) = 0$ 对
$1 \leq q \leq i$ 成立。基础情形 $i = 0$ 是平凡的。假设 $i \geq 1$。

\medskip\noindent
取 $\xi \in H^q_\etale(D(h), \mathcal{F})$，其中
$1 \leq q \leq i$ 且 $h \in R$。若 $q < i$，则由归纳已经完成，
故假设 $q = i$。把 $R$ 替换为 $R_h$ 后，可以假设
$\xi \in H^i_\etale(\Spec(R), \mathcal{F})$；这里略去一些细节。
令 $I \subset R$ 为元素 $f \in R$ 所成的集合，这些元素满足
$\xi|_{D(f)} = 0$。由于 $\xi$ 在 Zariski 局部平凡，
对每个素理想 $\mathfrak p$（属于 $R$），存在 $f \in I$ 使得
$f \not \in \mathfrak p$。因此，只要证明 $I$ 是理想，就有
$1 \in I$，从而完成证明。显然，$f \in I$、$r \in R$ 蕴含
$rf \in I$。于是设 $f, g \in I$，并证明 $f + g \in I$。
若 $q = i = 1$，这恰好就是引理的假设！因而得到 $i = 1$ 的情形。
当 $q = i > 1$ 时，注意
$$
D(f + g) = D(f(f + g)) \cup D(g(f + g))
$$
由 Mayer--Vietoris（上同调，引理
\ref{cohomology-lemma-mayer-vietoris}；它适用是因为
$\Spec(R)$ 的开子概形上的 \'etale 上同调等于 Zariski 上同调），
有正合列
$$
\xymatrix{
H^{i - 1}_\etale(D(fg(f + g)), \mathcal{F}) \ar[d] \\
H^i_\etale(D(f + g), \mathcal{F}) \ar[d] \\
H^i_\etale(D(f(f + g)), \mathcal{F}) \oplus
H^i_\etale(D(g(f + g)), \mathcal{F})
}
$$
而第一个群由归纳为零，故结论成立。
\end{proof}

\begin{lemma}
\label{etale-cohomology-lemma-affine-only-closed-points}
设 $S$ 是一个仿射概形，满足：(1) 所有点都是闭点；(2) 所有剩余域
都是可分代数闭的。则对任意阿贝尔层 $\mathcal{F}$（定义在
$S_\etale$ 上），有 $H^i(S_\etale, \mathcal{F}) = 0$ 对 $i > 0$ 成立。
\end{lemma}

\begin{proof}
条件 (1) 蕴含 $S$ 的底拓扑空间是预有限的，见代数，引理
\ref{algebra-lemma-ring-with-only-minimal-primes}。因此，拓扑空间
$S$ 上阿贝尔层的高阶上同调群（即 Zariski 上同调）为零，见上同调，
引理 \ref{cohomology-lemma-vanishing-for-profinite}。由代数，引理
\ref{algebra-lemma-local-dimension-zero-henselian}，其局部环严格 Hensel。
于是由引理 \ref{etale-cohomology-lemma-local-rings-strictly-henselian}，$S$ 的
\'etale 上同调由 Zariski 上同调计算，证明完成。
\end{proof}

\noindent
绝对整闭环的谱是所有局部环都严格 Hensel 的概形的一个例子，见代数进阶，
引理
\ref{more-algebra-lemma-absolutely-integrally-closed-strictly-henselian}。
事实证明，具有可分闭分式域的正规整环具有更强的性质，如下面的引理所述。

\begin{lemma}
\label{etale-cohomology-lemma-normal-scheme-with-alg-closed-function-field}
设 $X$ 是一个函数域可分闭的整正规概形。
\begin{enumerate}
\item 分离 \'etale 态射 $U \to X$ 是开浸入的不交并。
\item $X$ 的所有局部环都严格 Hensel。
\end{enumerate}
\end{lemma}

\begin{proof}
设 $R$ 是一个分式域可分代数闭的正规整环。设
$R \to A$ 是 \'etale 环映射。则 $A \otimes_R K$ 作为
$K$-代数是有限乘积 $\prod_{i = 1, \ldots, n} K$，其各因子都是
$K$。设 $e_i$（$i = 1, \ldots, n$）
是 $A \otimes_R K$ 中相应的幂等元。由于 $A$ 正规
（代数，引理 \ref{algebra-lemma-normal-goes-up}），这些幂等元
$e_i$ 属于 $A$（代数，引理
\ref{algebra-lemma-normal-ring-integrally-closed}）。因此
$A = \prod Ae_i$，可以假设 $A \otimes_R K = K$。由于
$A \subset A \otimes_R K = K$（利用 $R \to A$ 的平坦性以及
$R \subset K$），可知 $A$ 是整环。由同一论证可知
$A \otimes_R A \subset (A \otimes_R A) \otimes_R K = K$。
于是映射 $A \otimes_R A \to A$ 既是单射也是满射。因此，由态射，
引理 \ref{morphisms-lemma-universally-injective} 和 \'Etale 态射，
定理 \ref{etale-theorem-etale-radicial-open}，$R \to A$ 定义一个开浸入。

\medskip\noindent
设 $f : U \to X$ 是一个分离 \'etale 态射。设 $\eta \in X$
为一般点，并令 $f^{-1}(\{\eta\}) = \{\xi_i\}_{i \in I}$。
上一段的结果说明：对任意仿射开集 $U' \subset U$，只要其在
$X$ 中的像包含于一个仿射开集，就有
$U' = \coprod_{i \in I} U'_i$，其中 $U'_i$ 是 $U'$ 中所有
$\xi_i$ 的特殊化点所成的集合。此外，态射 $U'_i \to X$
是开浸入。由此可知，$U_i = \overline{\{\xi_i\}}$ 是 $U$ 的
开闭子概形，并且 $U_i \to X$ 在源上局部是同构。由态射，引理
\ref{morphisms-lemma-distinct-local-rings}，$U_i \to X$ 的分离性蕴含
$U_i \to X$ 是单射，故 $U_i \to X$ 是开浸入，即 (1) 成立。

\medskip\noindent
由 (1) 以及 $\mathcal{O}_{X, x}$ 的严格 Hensel 化可描述为
$\overline{x}$ 处局部环（位于 $X$ 的 \'etale 网站上）这一事实
（引理 \ref{etale-cohomology-lemma-describe-etale-local-ring}），可得 (2)。
也可以直接证明，见基本群，引理
\ref{pione-lemma-normal-local-domain-separablly-closed-fraction-field}。
\end{proof}

\begin{lemma}
\label{etale-cohomology-lemma-Rf-star-zero-normal-with-alg-closed-function-field}
设 $f : X \to Y$ 是概形态射，其中 $X$ 是函数域可分闭的整正规概形。
则 $R^qf_*\underline{M} = 0$ 对 $q > 0$ 以及任意阿贝尔群 $M$ 成立。
\end{lemma}

\begin{proof}
回忆，$R^qf_*\underline{M}$ 是预层
$V \mapsto H^q_\etale(V \times_Y X, M)$ 在 $Y_\etale$ 上的伴随层，见引理
\ref{etale-cohomology-lemma-higher-direct-images}。若 $V$ 仿射，则
$V \times_Y X \to X$ 是分离 \'etale 态射。因此
$V \times_Y X = \coprod U_i$ 是开子概形 $U_i$（属于 $X$）的不交并，
见引理 \ref{etale-cohomology-lemma-normal-scheme-with-alg-closed-function-field}。
由引理 \ref{etale-cohomology-lemma-local-rings-strictly-henselian}，
$H^q_\etale(U_i, M)$ 等于 $H^q_{Zar}(U_i, M)$。由上同调，引理
\ref{cohomology-lemma-irreducible-constant-cohomology-zero}，后者为零。
\end{proof}

\begin{lemma}
\label{etale-cohomology-lemma-closed-of-affine-normal-scheme-with-alg-closed-function-field}
设 $X$ 是函数域可分闭的仿射整正规概形。设
$Z \subset X$ 是一个闭子概形。设 $V \to Z$ 是 \'etale 态射，
且 $V$ 仿射。则 $V$ 是 $Z$ 的有限多个开子概形的不交并。
若 $V \to Z$ 是满射且有限 \'etale，则 $V \to Z$ 有一个截面。
\end{lemma}

\begin{proof}
由代数，引理 \ref{algebra-lemma-lift-etale}，可以把 $V$
提升为仿射概形 $U$，它在 $X$ 上是 \'etale 的。对 $U \to X$ 应用引理
\ref{etale-cohomology-lemma-normal-scheme-with-alg-closed-function-field}，得到第一个陈述。

\medskip\noindent
最后一个陈述是第一个陈述的推论。设
$V = \coprod_{i = 1, \ldots, n} V_i$ 是分解为开闭子概形的一个有限分解，
并且 $V_i \to Z$ 是开浸入。由于 $V \to Z$ 有限，可知
$V_i \to Z$ 也是闭的。设 $U_i \subset Z$ 是其像。于是有开闭子概形分解
$$
Z =
\coprod\nolimits_{(A, B)}
\bigcap\nolimits_{i \in A} U_i \cap
\bigcap\nolimits_{i \in B} U_i^c
$$
，其中不交并遍历所有分解 $\{1, \ldots, n\} = A \amalg B$，
且 $A$ 至少有一个元素。每个层都包含在某个 $U_i$ 中，
于是得到所需截面。
\end{proof}

\begin{lemma}
\label{etale-cohomology-lemma-gabber-for-h1-absolutely-algebraically-closed}
设 $X$ 是函数域可分闭的正规整仿射概形。设
$Z \subset X$ 是一个闭子概形。对任意有限阿贝尔群
$M$，有 $H^1_\etale(Z, \underline{M}) = 0$。
\end{lemma}

\begin{proof}
由位点上的上同调，引理 \ref{sites-cohomology-lemma-torsors-h1}，
$H^1_\etale(Z, \underline{M})$ 的一个元素对应于一个
$\underline{M}$-旋子 $\mathcal{F}$，它定义在 $Z_\etale$ 上。
这样的旋子显然是有限局部常值层。
因此，由引理 \ref{etale-cohomology-lemma-characterize-finite-locally-constant}，
$\mathcal{F}$ 可由一个概形 $V$ 表示，它在 $Z$ 上是有限 \'etale 的。
旋子局部平凡，所以 $V \to Z$ 当然是满射。由引理
\ref{etale-cohomology-lemma-closed-of-affine-normal-scheme-with-alg-closed-function-field}，
$V \to Z$ 有截面，故得结论。
\end{proof}

\begin{lemma}
\label{etale-cohomology-lemma-gabber-for-absolutely-algebraically-closed}
设 $X$ 是函数域可分闭的正规整仿射概形。设
$Z \subset X$ 是一个闭子概形。对任意有限阿贝尔群 $M$，有
$H^q_\etale(Z, \underline{M}) = 0$ 对 $q \geq 1$ 成立。
\end{lemma}

\begin{proof}
写 $X = \Spec(R)$ 且 $Z = \Spec(R')$，于是有环满射
$R \to R'$。由引理
\ref{etale-cohomology-lemma-normal-scheme-with-alg-closed-function-field} 和代数，引理
\ref{algebra-lemma-quotient-strict-henselization}，$R'$ 的所有局部环
都严格 Hensel。此外可知，对任意 $f' \in R'$，都有满射
$R_f \to R'_{f'}$，其中 $f \in R$ 是 $f'$ 的一个提升。
由于 $R_f$ 是分式域可分闭的正规整环，由引理
\ref{etale-cohomology-lemma-gabber-for-h1-absolutely-algebraically-closed} 可知
$H^1_\etale(D(f'), \underline{M}) = 0$。因此可以对
$Z = \Spec(R')$ 应用引理 \ref{etale-cohomology-lemma-gabber-criterion} 得到结论。
\end{proof}

\begin{lemma}
\label{etale-cohomology-lemma-integral-cover-trivial-cohomology}
设 $X$ 是一个仿射概形。
\begin{enumerate}
\item 存在整满射 $X' \to X$，使得对每个闭子概形
$Z' \subset X'$、每个有限阿贝尔群 $M$ 以及每个 $q \geq 1$，都有
$H^q_\etale(Z', \underline{M}) = 0$。
\item 对任意闭子概形 $Z \subset X$、有限阿贝尔群 $M$、
$q \geq 1$ 以及 $\xi \in H^q_\etale(Z, \underline{M})$，存在有限表现的
有限满射 $X' \to X$，使得 $\xi$ 在
$H^q_\etale(X' \times_X Z, \underline{M})$ 中的拉回为零。
\end{enumerate}
\end{lemma}

\begin{proof}
写 $X = \Spec(A)$。写 $A = \mathbf{Z}[x_i]/J$，其中 $J$ 是某个理想。
设 $R$ 是 $\mathbf{Z}[x_i]$ 的分式域的
某个代数闭包中 $\mathbf{Z}[x_i]$ 的整闭包。令
$A' = R/JR$，并置 $X' = \Spec(A')$。由引理
\ref{etale-cohomology-lemma-gabber-for-absolutely-algebraically-closed}，
这给出 (1) 中的一个例子。

\medskip\noindent
证明 (2)。设 $X' \to X$ 是上面找到的整满射。显然，$\xi$ 在
$H^q_\etale(X' \times_X Z, \underline{M})$ 中的像为零。可以把
$X'$ 写成极限 $X' = \lim X'_i$，其中各概形在 $X$ 上有限且有限表现；
在当前仿射情形下这很容易做到，而一般地，这是极限，
引理 \ref{limits-lemma-integral-limit-finite-and-finite-presentation}
的一个特殊情形。由引理 \ref{etale-cohomology-lemma-directed-colimit-cohomology}，
可知 $\xi$ 在 $H^q_\etale(X'_i \times_X Z, \underline{M})$
中的像为零，其中 $i$ 充分大。
\end{proof}








\section{绝对整闭情形的消失}
\label{etale-cohomology-section-aic-vanishing}

\noindent
回忆，若环 $R$ 上的每个首一多项式都在 $R$ 中有根，则称 $R$
是绝对整闭的（代数进阶，定义
\ref{more-algebra-definition-absolutely-integrally-closed}）。
本节证明 $\Spec(R)$ 以有限挠群为系数的 \'etale 上同调在正次数消失
（命题 \ref{etale-cohomology-proposition-aic-vanishing}），从而略微改进先前的引理
\ref{etale-cohomology-lemma-gabber-for-absolutely-algebraically-closed}。
建议读者跳过本节。

\begin{lemma}
\label{etale-cohomology-lemma-find-extension}
设 $A$ 是一个环。设 $a, b \in A$，满足
$aA + bA = A$，且 $a \bmod bA$ 是一个单位根。
则存在单生成扩张 $A \subset B$ 和元素 $y \in B$，
使得 $u = a - by$ 是一个单位。
\end{lemma}

\begin{proof}
设 $a^n \equiv 1 \bmod bA$。特别地，$a^i$ 模
$b^mA$ 对所有 $i, m \geq 1$ 都是单位。我们断言，存在
$a_1, \ldots, a_n \in A$，使得
$$
1 = a^n + a_1 a^{n - 1}b + a_2 a^{n - 2}b^2 + \ldots + a_n b^n
$$
事实上，由于 $1 - a^n \in bA$，可以找到元素 $a_1 \in A$，使得
$1 - a^n - a_1 a^{n - 1} b \in b^2A$；这里使用了
$a^{n - 1}$ 模 $bA$ 为单位这一事实。接着可以找到元素
$a_2 \in A$，使得
$1 - a^n - a_1 a^{n - 1} b - a_2 a^{n - 2} b^2 \in b^3A$。
如此继续。最终找到 $a_1, \ldots, a_{n - 1} \in A$，使得
$1 - (a^n + a_1 a^{n - 1}b + a_2 a^{n - 2}b^2 +
\ldots + a_{n - 1} ab^{n - 1}) \in b^nA$。于是可以找到
$a_n \in A$，使上述等式成立。

\medskip\noindent
对上述 $a_1, \ldots, a_n$，我们断言取
$$
B = A[y]/(y^n + a_1 y^{n - 1} + a_2 y^{n - 2} + \ldots + a_n)
$$
即可。事实上，设 $\mathfrak q \subset B$ 是位于
$\mathfrak p \subset A$ 之上的素理想。为得到矛盾，假设
$u = a - by$ 属于 $\mathfrak q$。若 $b \in \mathfrak p$，
则 $a \not \in \mathfrak p$，因为 $aA + bA = A$，从而
$u$ 不属于 $\mathfrak q$。因此可以假设 $b \not \in \mathfrak p$，
即 $b \not \in \mathfrak q$。这蕴含 $y \bmod \mathfrak q$
等于 $a/b \bmod \mathfrak q$。然而此时得到
$$
0 = y^n + a_1 y^{n - 1} + a_2 y^{n - 2} + \ldots + a_n =
b^{-n}(a^n + a_1 a^{n - 1}b + a_2 a^{n - 2}b^2 + \ldots + a_nb^n) =
b^{-n}
$$
，矛盾。证明完成。
\end{proof}

\noindent
为说明证明，需要引入一些群概形。固定一个素数 $\ell$。令
$$
A = \mathbf{Z}[\zeta] =
\mathbf{Z}[x]/(x^{\ell - 1} + x^{\ell - 2} + \ldots + 1)
$$
换言之，$A$ 是 $\mathbf{Z}$ 的单生成扩张，它由一个本原
$\ell$ 次单位根 $\zeta$ 生成。置
$$
\pi = \zeta - 1
$$
一个略去的计算表明，$\ell$ 可被 $\pi^{\ell - 1}$ 整除；
这里的计算在 $A$ 中进行。
第一个 $A$ 上群概形为
$$
G = \Spec(A[s, \frac{1}{\pi s + 1}])
$$
，其群律由余乘法
$$
\mu :
A[s, \frac{1}{\pi s + 1}]
\longrightarrow
A[s, \frac{1}{\pi s + 1}] \otimes_A A[s, \frac{1}{\pi s + 1}],\quad
s \longmapsto \pi s \otimes s + s \otimes 1 + 1 \otimes s
$$
给出。由此有
$$
\mu(\pi s + 1) = (\pi s + 1) \otimes (\pi s + 1)
$$
，所以确实得到所示的 $A$-代数映射。略去它确实定义群律的验证。
第二个 $A$ 上群概形为
$$
H = \Spec(A[t, \frac{1}{\pi^\ell t + 1}])
$$
，其群律由余乘法
$$
\mu : A[t, \frac{1}{\pi^\ell t + 1}]
\longrightarrow
A[t, \frac{1}{\pi^\ell t + 1}] \otimes_A A[t, \frac{1}{\pi^\ell t + 1}],\quad
t \longmapsto \pi^\ell t \otimes t + t \otimes 1 + 1 \otimes t
$$
给出。与前面相同的验证表明，这定义了一个群律。接着注意多项式
$$
\Phi(s) = \frac{(\pi s + 1)^\ell - 1}{\pi^\ell}
$$
属于 $A[s]$，次数为 $\ell$，并且关于 $s$ 首一。事实上，
$s^i$ 的系数在 $0 < i < \ell$ 时等于
${\ell \choose i}\pi^{i - \ell}$；由于
$\pi^{\ell - 1}$ 整除 $\ell$（在 $A$ 中），这是 $A$ 的元素。于是得到环映射
$$
A[t, \frac{1}{\pi^\ell t + 1}]
\longrightarrow
A[s, \frac{1}{\pi s + 1}],\quad
t \longmapsto \Phi(s)
$$
；读者容易验证它与余乘法相容。因此得到群概形态射
$$
f : G \to H
$$
下面的引理特别说明，该态射忠实平坦（事实上将看到它是有限
\'etale 满射）。

\begin{lemma}
\label{etale-cohomology-lemma-monogenic-one}
有
$$
A[s, \frac{1}{\pi s + 1}] =
\left(A[t, \frac{1}{\pi^\ell t + 1}]\right)[s]/(\Phi(s) - t)
$$
特别地，$G$ 的 Hopf 代数是 $H$ 的 Hopf 代数的单生成扩张。
\end{lemma}

\begin{proof}
这由上面的讨论和 $\Phi(s)$ 的形式得出。特别注意，利用
$\Phi(s) = t$，元素 $\frac{1}{\pi^\ell t + 1}$ 变为元素
$\frac{1}{(\pi s + 1)^\ell}$。
\end{proof}

\noindent
接下来计算 $f$ 的核。由于 $H$ 的单位元
由 $t = 0$ 在 $H$ 中给出，可知 $f$ 的核由 $\Phi(s) = 0$ 给出。现在注意，
下列 $A$-值点
$\sigma_0, \ldots, \sigma_{\ell - 1}$
属于 $G$，其定义为
$$
\sigma_i :
s = \frac{\zeta^i - 1}{\pi} = \frac{\zeta^i - 1}{\zeta - 1} =
\zeta^{i - 1} + \zeta^{i - 2} + \ldots + 1,\quad
i = 0, 1, \ldots, \ell - 1
$$
。它们显然包含在 $\Ker(f)$ 中。此外，它们在 $G \to \Spec(A)$
的{\bf 每个}纤维中都两两不同。读者还可计算出
$\sigma_i +_G \sigma_j = \sigma_{i + j \bmod \ell}$。因此得到群概形的闭浸入
$$
\underline{\mathbf{Z}/\ell \mathbf{Z}}_A \longrightarrow \Ker(f)
$$
，它把 $i$ 送到 $\sigma_i$。然而，按构造，$\Ker(f)$ 在
$\Spec(A)$ 上是次数为 $\ell$ 的有限平坦概形。因此该映射是同构。
总之，得到 $A$ 上群概形的短正合列
\begin{equation}
\label{etale-cohomology-equation-ses}
0 \to
\underline{\mathbf{Z}/\ell \mathbf{Z}}_A
\to G
\to H
\to 0
\end{equation}
。

\begin{lemma}
\label{etale-cohomology-lemma-lift-points-H-to-G}
设 $R$ 是一个绝对整闭的 $A$-代数。则 $G(R) \to H(R)$ 是满射。
\end{lemma}

\begin{proof}
设 $h \in H(R)$ 对应于 $A$-代数映射
$A[t, \frac{1}{\pi^\ell t + 1}] \to R$，它把 $t$ 送到 $a \in A$。
由于 $\Phi(s)$ 首一，可以找到 $b \in A$，使得 $\Phi(b) = a$。
由引理 \ref{etale-cohomology-lemma-monogenic-one}，把 $s$ 送到 $b$，得到唯一的
$A$-代数映射 $A[s, \frac{1}{\pi s + 1}] \to R$，它与上面的映射
$A[t, \frac{1}{\pi^\ell t + 1}] \to R$ 相容。后者又对应于元素
$g \in G(R)$，它映到 $h \in H(R)$。
\end{proof}

\begin{lemma}
\label{etale-cohomology-lemma-interpolate}
设 $R$ 是一个绝对整闭的 $A$-代数。设 $I, J \subset R$
是满足 $I + J = R$ 的理想。存在 $g \in G(R)$，使得
$g \bmod I = \sigma_0$ 且 $g \bmod J = \sigma_1$。
\end{lemma}

\begin{proof}
选取 $x \in I$，使得 $x \equiv 1 \bmod J$。可以并且确实把
$I$ 替换为 $xR$，把 $J$ 替换为 $(x - 1)R$。于是要找
$s \in R$，使得
\begin{enumerate}
\item $1 + \pi s$ 是单位；
\item $s \equiv 0 \bmod xR$；以及
\item $s \equiv 1 \bmod (x - 1)R$。
\end{enumerate}
后两个条件说明，$s = x + x(x - 1)y$，其中某个 $y \in R$。
第一个条件说明
$1 + \pi s = 1 + \pi x + \pi x (x - 1) y$ 需要是 $R$ 的单位。
但注意，$1 + \pi x$ 与 $\pi x (x - 1)$ 生成 $R$ 的单位理想，
而且 $1 + \pi x$ 是 $\ell$ 次方等于 $1$ 的元素（模
$\pi x (x - 1)$）
\footnote{因为 $1 + \pi x$ 同余于 $1$（模 $\pi$），同余于
$1$（模 $x$），并且同余于 $1 + \pi = \zeta$（模 $x - 1$）；
此外，$(\pi) \cap (x) \cap (x - 1) = (\pi x (x - 1))$ 在
$A[x]$ 中成立。}。
由引理 \ref{etale-cohomology-lemma-find-extension} 以及 $R$ 绝对整闭这一事实即得结论。
\end{proof}

\begin{proposition}
\label{etale-cohomology-proposition-aic-vanishing}
设 $R$ 是一个绝对整闭环。设 $M$ 是有限阿贝尔群。
则 $H^i_\etale(\Spec(R), \underline{M}) = 0$ 对 $i > 0$ 成立。
\end{proposition}

\begin{proof}
任意有限阿贝尔群都有一个有限滤过，其子商是素数阶循环群，
故可以假设 $M = \mathbf{Z}/\ell\mathbf{Z}$，其中 $\ell$ 是素数。

\medskip\noindent
注意，$R$ 的所有局部环都严格 Hensel，见代数进阶，引理
\ref{more-algebra-lemma-absolutely-integrally-closed-strictly-henselian}。
此外，由代数进阶，引理
\ref{more-algebra-lemma-absolutely-integrally-closed-quotient-localization}，
$R$ 的任意局部化也绝对整闭。因此，引理
\ref{etale-cohomology-lemma-gabber-criterion} 告诉我们，只需证明映射
$$
H^1_\etale(D(f + g), \mathbf{Z}/\ell\mathbf{Z})
\longrightarrow
H^1_\etale(D(f(f + g)), \mathbf{Z}/\ell\mathbf{Z}) \oplus
H^1_\etale(D(g(f + g)), \mathbf{Z}/\ell\mathbf{Z})
$$
的核对任意 $f, g \in R$ 都为零。把 $R$ 替换为 $R_{f + g}$ 后，
归结为如下断言：给定
$\xi \in H^1_\etale(\Spec(R), \mathbf{Z}/\ell\mathbf{Z})$
以及仿射开覆盖 $\Spec(R) = U \cup V$，若 $\xi|_U$ 和
$\xi|_V$ 平凡，则 $\xi = 0$。

\medskip\noindent
令 $A = \mathbf{Z}[\zeta]$ 如上。由于 $\mathbf{Z} \subset A$
是单生成的，可以找到环映射 $A \to R$。从现在起，把 $R$
看作 $A$-代数，并把 $\Spec(R)$ 看作 $\Spec(A)$ 上的概形。
把短正合列 (\ref{etale-cohomology-equation-ses}) 基变换到 $\Spec(R)$ 并取
\'etale 上同调，得到
$$
G(R) \to H(R) \to
H^1_\etale(\Spec(R), \mathbf{Z}/\ell\mathbf{Z}) \to
H^1_\etale(\Spec(R), G)
$$
请在证明余下部分始终记住这一点。

\medskip\noindent
设 $\tau \in \Gamma(U \cap V, \mathbf{Z}/\ell\mathbf{Z})$ 是一个截面，
它在 Mayer--Vietoris 正合列（引理 \ref{etale-cohomology-lemma-mayer-vietoris}）中的
边界给出 $\xi$。对 $i = 0, 1, \ldots, \ell - 1$，令
$A_i \subset U \cap V$ 为 $\tau$ 取值 $i \bmod \ell$ 的开闭子集。
于是有有限不交并分解
$$
U \cap V = A_0 \amalg \ldots \amalg A_{\ell - 1}
$$
，使得 $\tau$ 在每个 $A_i$ 上常值。对
$i = 0, 1, \ldots, \ell - 1$，用
$\tau_i \in H^0(U \cap V, \mathbf{Z}/\ell\mathbf{Z})$ 表示如下元素：
它取值 $1$ 的地方是 $A_i$，取值 $0$ 的地方是 $A_j$，其中
$j \not = i$。
于是 $\tau$ 是各 $\tau_i$ 的倍数之和\footnote{不计计算错误，
有 $\tau = \sum i \tau_i$。}。因此只需证明对应于 $\tau_i$
的上同调类平凡。这把问题归结到 $\tau$ 只取两个不同值
$1$ 和 $0$ 的情形。

\medskip\noindent
假设 $\tau$ 只取值 $1$ 和 $0$。写作
$$
U \cap V = A \amalg B
$$
其中 $A$ 是 $\tau = 0$ 的点集，$B$ 是
$\tau = 1$ 的点集。于是 $A$ 与 $B$ 是互不相交的闭子集。
用 $\overline{A}$ 和 $\overline{B}$ 表示 $A$ 与 $B$ 在
$\Spec(R)$ 中的闭包。于是得到一个“香蕉形”：即有
$$
\overline{A} \cap \overline{B} = Z_1 \amalg Z_2
$$
其中 $Z_1 \subset U$ 与 $Z_2 \subset V$ 是互不相交的闭子集。
置 $T_1 = \Spec(R) \setminus V$ 和 $T_2 = \Spec(R) \setminus U$。
注意 $Z_1 \subset T_1 \subset U$、$Z_2 \subset T_2 \subset V$，且
$T_1 \cap T_2 = \emptyset$。从拓扑上可写作
$$
\Spec(R) = \overline{A} \cup \overline{B} \cup T_1 \cup T_2
$$
建议画图以便直观理解。
为证明 $\xi$ 为零，我们可以并且确实把 $R$ 替换为其约化
（命题 \ref{etale-cohomology-proposition-topological-invariance}）。
下文把 $A$、$\overline{A}$、$B$、$\overline{B}$、$T_1$、$T_2$
看作 $\Spec(R)$ 的约化闭子概形。接着，在 $Z_1$
和 $Z_2$ 上采用如下概形结构：
$$
Z_1 = \overline{A} \cap (\overline{B} \cup T_1)
\quad\text{且}\quad
Z_2 = \overline{A} \cap (\overline{B} \cup T_2)
$$
（即态射，定义
\ref{morphisms-definition-scheme-theoretic-intersection-union}
中的概形论并与交）。

\medskip\noindent
用 $X$ 表示一个 $G$-旋子；它定义在 $\Spec(R)$ 上，并对应于
$\xi$ 在 $H^1(\Spec(R), G)$ 中的像。若 $X$ 平凡，则 $\xi$
来自一个元素 $h \in H(R)$（见上面的上同调正合列）。
然而，由引理 \ref{etale-cohomology-lemma-lift-points-H-to-G}，元素 $h$
可提升为 $G(R)$ 的一个元素，因而如所愿得到
$\xi = 0$。所以我们的目标是证明 $X$ 平凡。

\medskip\noindent
回忆，嵌入 $\mathbf{Z}/\ell \mathbf{Z} \to G(R)$
把 $i \bmod \ell$ 送到 $\sigma_i \in G(R)$。注意 $\overline{A}$
是一个绝对整闭环的谱（即 $R$ 的一个商）。由
引理 \ref{etale-cohomology-lemma-interpolate}，可以找到 $g \in G(\overline{A})$，使得
$g|_{\overline{A} \cap Z_1} = \sigma_0$ 且
$g|_{\overline{A} \cap Z_2} = \sigma_1$（按概形论意义）。
于是可以定义
\begin{enumerate}
\item $g_1 \in G(U)$：它等于 $g$ 于 $\overline{A} \cap U$，
等于 $\sigma_0$ 于 $\overline{B} \cap U$，并且等于 $\sigma_0$ 于 $T_1$；以及
\item $g_2 \in G(V)$：它等于 $g$ 于 $\overline{A} \cap V$，
等于 $\sigma_1$ 于 $\overline{B} \cap V$，并且等于 $\sigma_1$ 于 $T_2$。
\end{enumerate}
具体而言，为找到 (1) 中的 $g_1$，把截面
$\sigma_0$ 在 $\Omega = (\overline{B} \cup T_1) \cap U$ 上的取值
与截面 $g$ 在 $\Omega' = \overline{A} \cap U$ 上的限制黏合。
注意 $U = \Omega \cup \Omega'$（按概形论意义），
因为 $U$ 是约化的，并且
$\Omega \cap \Omega' = Z_1$（按概形论意义）；这里使用了我们对
$Z_1$ 的选取。因此由
态射，引理 \ref{morphisms-lemma-scheme-theoretic-union}，
$U$ 是 $\Omega$ 与 $\Omega'$ 沿 $Z_1$ 的推出。
所以可以找到 $g_1$。
(2) 中 $g_2$ 的存在性同理。
于是有
$$
\tau = g_2|_{A \cup B} - g_1|_{A \cup B} \quad(\text{addition in group law})
$$
从而看出 $X$ 平凡，证明完成。
\end{proof}




















\section{真基变换的仿射类似}
\label{etale-cohomology-section-gabber-affine-proper}

\noindent
本节讨论 Ofer Gabber 的一个结果，见
\cite{gabber-affine-proper}。Roland Huber 也证明了这一结果，见
\cite{Huber-henselian}。Gabber 证明所需的一部分工作已经在
第 \ref{etale-cohomology-section-strictly-henselian-local-rings} 节中完成。

\begin{lemma}
\label{etale-cohomology-lemma-efface-cohomology-on-closed-by-finite-cover}
设 $X$ 是仿射概形。设 $\mathcal{F}$ 是 $X_\etale$ 上的挠阿贝尔层。
设 $Z \subset X$ 是闭子概形。设
$\xi \in H^q_\etale(Z, \mathcal{F}|_Z)$，其中 $q > 0$。
则存在挠阿贝尔层的单射 $\mathcal{F} \to \mathcal{F}'$
（在 $X_\etale$ 上），使得 $\xi$ 在
$H^q_\etale(Z, \mathcal{F}'|_Z)$ 中的像为零。
\end{lemma}

\begin{proof}
由引理 \ref{etale-cohomology-lemma-torsion-colimit-constructible} 和 \ref{etale-cohomology-lemma-colimit}，
可以找到态射 $\mathcal{G} \to \mathcal{F}$，其中 $\mathcal{G}$
是可构造阿贝尔层，且 $\xi$ 来自某个元素 $\zeta$，该元素属于
$H^q_\etale(Z, \mathcal{G}|_Z)$。假设可以找到挠阿贝尔层的单射
$\mathcal{G} \to \mathcal{G}'$（在 $X_\etale$ 上），使得 $\zeta$ 在
$H^q_\etale(Z, \mathcal{G}'|_Z)$ 中的像为零。
此时可以取 $\mathcal{F}'$ 为推出
$$
\mathcal{F}' = \mathcal{G}' \amalg_{\mathcal{G}} \mathcal{F}
$$
，从而得到引理的结论。（注意，限制到 $Z$ 是正合的，故与有限极限
和余极限交换；此外，它作为推前的左伴随，也与任意余极限交换。）
因此可以假设 $\mathcal{F}$ 可构造。

\medskip\noindent
假设 $\mathcal{F}$ 可构造。由
引理 \ref{etale-cohomology-lemma-constructible-maps-into-constant-general}，
只需在 $\mathcal{F}$ 形如 $f_*\underline{M}$ 时证明结论，
其中 $M$ 是有限阿贝尔群，而 $f : Y \to X$ 是有限呈示的有限态射
（由引理 \ref{etale-cohomology-lemma-finite-pushforward-constructible}，
这种层仍然可构造，但这里不需要这一点）。
由于 $f_*$ 的形成与任意基变换交换
（引理 \ref{etale-cohomology-lemma-finite-pushforward-commutes-with-base-change}），
$f_*\underline{M}$ 限制到 $Z$ 等于 $\underline{M}$ 沿
$Y \times_X Z \to Z$ 的推前。由 Leray 谱序列
（命题 \ref{etale-cohomology-proposition-leray}）
和高阶正像的消失
（命题 \ref{etale-cohomology-proposition-finite-higher-direct-image-zero}），得到
$$
H^q_\etale(Z, f_*\underline{M}|_Z) = H^q_\etale(Y \times_X Z, \underline{M}).
$$
由引理 \ref{etale-cohomology-lemma-integral-cover-trivial-cohomology}，
可以找到有限呈示的有限满态射 $Y' \to Y$，使得 $\xi$ 在
$H^q(Y' \times_X Z, \underline{M})$ 中的像为零。
用 $f' : Y' \to X$ 表示复合 $Y' \to Y \to X$。我们断言映射
$$
f_*\underline{M} \longrightarrow f'_*\underline{M}
$$
是单射；结合上述讨论，这就完成证明。为验证所需的单射性，
可以考察茎。具体地，若 $\overline{x} : \Spec(k) \to X$ 是几何点，则
$$
(f_*\underline{M})_{\overline{x}} =
\bigoplus\nolimits_{f(\overline{y}) = \overline{x}} M
$$
，这由命题 \ref{etale-cohomology-proposition-finite-higher-direct-image-zero} 得到；
另一个层也同样。由于 $Y' \to Y$ 是有限满态射，
它在位于 $\overline{x}$ 上方的几何点之间诱导的映射也满，
故得结论。
\end{proof}

\noindent
上面的引理将处理 Gabber 结果中的高阶上同调群。
下面的引理用于处理整体截面。

\begin{lemma}
\label{etale-cohomology-lemma-gabber-h0}
设 $X$ 是拟紧拟分离概形。设 $i : Z \to X$ 是闭浸入。假设
\begin{enumerate}
\item 对任意层 $\mathcal{F}$（在 $X_{Zar}$ 上），映射
$\Gamma(X, \mathcal{F}) \to \Gamma(Z, i^{-1}\mathcal{F})$
是双射；并且
\item 对任意有限态射 $X' \to X$，假设 (1) 对
$Z \times_X X' \to X'$ 成立。
\end{enumerate}
则对任意层 $\mathcal{F}$（在 $X_\etale$ 上），有
$\Gamma(X, \mathcal{F}) = \Gamma(Z, i^{-1}_{small}\mathcal{F})$。
\end{lemma}

\begin{proof}
设 $\mathcal{F}$ 是 $X_\etale$ 上的层。有典范的（基变换）映射
$$
i^{-1}(\mathcal{F}|_{X_{Zar}})
\longrightarrow
(i_{small}^{-1}\mathcal{F})|_{Z_{Zar}}
$$
，这是 $Z_{Zar}$ 上层的映射。通过考察茎来证明该映射为单射。
左边在 $z \in Z$ 处的茎是 $\mathcal{F}|_{X_{Zar}}$ 在 $z$ 处的茎。
右边的茎是对所有初等 \'etale 邻域
$(U, u) \to (X, z)$ 取余极限，其中 $U \times_X Z \to Z$
在 $z$ 的某个邻域上有截面。由于 \'etale 态射是开映射，
$U \to X$ 的像是开邻域 $U_0$，它是 $z$ 在 $X$ 中的邻域。
映射 $\mathcal{F}(U_0) \to \mathcal{F}(U)$ 是单射；
这是由层 $\mathcal{F}$ 关于 \'etale 覆盖 $U \to U_0$ 的层条件得到的。
对所有 $U$ 和 $U_0$ 取余极限，便得到茎上的单射性。

\medskip\noindent
由此和假设 (1)，映射
$\Gamma(X, \mathcal{F}) \to \Gamma(Z, i^{-1}_{small}\mathcal{F})$
是单射。由 (2)，对所有 $X'$（在 $X$ 上有限）也有同一结论。

\medskip\noindent
在证明满射性之前先引入一些记号。对任意对象 $U$（属于 $X_\etale$）
和 $s \in \mathcal{F}(U)$，用
$i^{-1}s \in (i_{small}^{-1}\mathcal{F})(U \times_X Z)$
表示拉回（更准确地说，这是 $s$ 在位点，引理
\ref{sites-lemma-recover} 的映射与从预层拉回到层拉回的映射之复合下的像）。
这一构造与 $X_\etale$ 中态射的拉回以及（有限）态射
$X' \to X$ 的拉回相容，具体细节留给读者说明。

\medskip\noindent
设 $t \in \Gamma(Z, i^{-1}_{small}\mathcal{F})$。由
$i^{-1}_{small}\mathcal{F}$ 的构造，存在 \'etale 覆盖
$\{V_j \to Z\}_{j \in J}$、\'etale 态射 $U_j \to X$、截面
$s_j \in \mathcal{F}(U_j)$ 以及态射 $V_j \to U_j$（在 $X$ 上），
使得 $t|_{V_j}$ 是 $i^{-1}s_j$ 沿
$V_j \to U_j \times_X Z$ 的拉回。由于 $V_j \to U_j \times_X Z$
是 \'etale 的，其像是开子概形。由于
$U_j \times_X Z \subset U_j$ 是闭的，把 $U_j$ 替换为一个开子概形后，
可以假设 $V_j \to U_j \times_X Z$ 满。于是
$\{V_j \to U_j \times_X Z\}$ 是 \'etale 覆盖，从而
$t|_{U_j \times_X Z} = i^{-1}s_j$。
注意，每个非空闭子概形 $T \subset X$ 都与 $Z$ 相交；例如，把假设 (1)
应用于层 $(T \to X)_*\underline{\mathbf{Z}}$ 即得。
因此 $\coprod U_j \to X$ 满。由态射进阶，引理
\ref{more-morphisms-lemma-there-is-a-scheme-integral-over}
，可以找到有限满态射 $X' \to X$，使得 $X' \to X$ 在 Zariski 局部
经 $\coprod U_j \to X$ 分解。设
$X' = \bigcup_{\alpha \in A} W_\alpha$ 是开覆盖，$j : A \to J$ 是映射，
并且存在态射 $g_\alpha : W_\alpha \to U_{j(\alpha)}$（在 $X$ 上）。
用 $\mathcal{F}'$ 表示 $\mathcal{F}$ 到 $X'_\etale$ 的拉回。
用 $s'_\alpha \in \mathcal{F}'(W_\alpha)$ 表示
$s_{j(\alpha)}$ 沿 $g_\alpha$ 的拉回。
置 $Z' = X' \times_X Z$。用 $i' : Z' \to X'$ 表示 $i$ 的基变换。
于是 $(i'_{small})^{-1}\mathcal{F}' =
(Z' \to Z)_{small}^{-1}i_{small}^{-1}\mathcal{F}$。
在这个等同下，用
$t' \in
\Gamma(Z', (i')^{-1}_{small}\mathcal{F}')$
表示 $t$ 到 $Z'$ 的拉回。
按构造，$t'|_{W_\alpha \times_X Z}$ 等于
$(i')^{-1}s'_\alpha$。因此 $t'$ 是如下子层的截面：
$$
(i')^{-1}(\mathcal{F}'|_{X'_{Zar}})
\subset
(i'_{small})^{-1}\mathcal{F}')|_{Z'_{Zar}}
$$
由假设 (2)，截面 $t'$ 来自截面
$s' \in \Gamma(X', \mathcal{F}'|_{X'_{Zar}}) = \Gamma(X', \mathcal{F}')$。
由第二段证明的单射性，$s'$ 到 $X' \times_X X'$ 的两个拉回相同
（毕竟 $t'$ 到 $Z' \times_Z Z'$ 的拉回具有这一性质）。
因此由注 \ref{etale-cohomology-remark-cohomological-descent-finite}，
$s'$ 来自 $\mathcal{F}$ 在 $X$ 上的一个截面。
\end{proof}

\begin{lemma}
\label{etale-cohomology-lemma-connected-topological}
设 $Z \subset X$ 是拓扑空间 $X$ 的闭子集。假设
\begin{enumerate}
\item $X$ 是谱空间
（拓扑，定义 \ref{topology-definition-spectral-space}）；并且
\item 对 $x \in X$，交 $Z \cap \overline{\{x\}}$
连通（特别地非空）。
\end{enumerate}
若 $Z = Z_1 \amalg Z_2$，且 $Z_i$ 在 $Z$ 中闭，
则存在分解 $X = X_1 \amalg X_2$，其中 $X_i$ 在 $X$ 中闭，
并且 $Z_i = Z \cap X_i$。
\end{lemma}

\begin{proof}
注意 $Z_i$ 拟紧。因此，由拓扑，引理
\ref{topology-lemma-make-spectral-space}，点集
$W_i$（其点特化到 $Z_i$）在可构造拓扑中闭。假设 (2) 蕴含
$X = W_1 \amalg W_2$。
设 $x \in \overline{W_1}$。由拓扑，引理
\ref{topology-lemma-constructible-stable-specialization-closed} 的 (1)，
存在特化 $x_1 \leadsto x$，其中 $x_1 \in W_1$。
于是 $\overline{\{x\}} \subset \overline{\{x_1\}}$，从而
$x \in W_1$。换言之，取 $X_i = W_i$ 即可。
\end{proof}

\begin{lemma}
\label{etale-cohomology-lemma-h0-topological}
设 $Z \subset X$ 是拓扑空间 $X$ 的闭子集。假设
\begin{enumerate}
\item $X$ 是谱空间
（拓扑，定义 \ref{topology-definition-spectral-space}）；并且
\item 对 $x \in X$，交 $Z \cap \overline{\{x\}}$
连通（特别地非空）。
\end{enumerate}
则对任意层 $\mathcal{F}$（在 $X$ 上），有
$\Gamma(X, \mathcal{F}) = \Gamma(Z, \mathcal{F}|_Z)$。
\end{lemma}

\begin{proof}
若 $x \leadsto x'$ 是点的特化，则有典范映射
$\mathcal{F}_{x'} \to \mathcal{F}_x$；它与开集上的截面相容，
并且关于 $\mathcal{F}$ 函子性。由于 $X$ 的每个点都特化到 $Z$ 中的点，
可知 $\Gamma(X, \mathcal{F}) \to \Gamma(Z, \mathcal{F}|_Z)$ 为单射。
困难之处是证明它满。

\medskip\noindent
用 $\mathcal{B}$ 表示 $X$ 的全体拟紧开集。把 $\mathcal{F}$ 写成滤过余极限
$\mathcal{F} = \colim \mathcal{F}_i$，其中每个 $\mathcal{F}_i$
都如模，等式 (\ref{modules-equation-towards-constructible-sets}) 所述。
见模，引理 \ref{modules-lemma-filtered-colimit-constructibles}。
于是有 $\mathcal{F}|_Z = \colim \mathcal{F}_i|_Z$；
这是因为限制到 $Z$ 是左伴随（范畴，引理
\ref{categories-lemma-adjoint-exact} 和层，引理
\ref{sheaves-lemma-f-map}）。
由层，引理 \ref{sheaves-lemma-directed-colimits-sections}，函子
$\Gamma(X, -)$ 与 $\Gamma(Z, -)$ 和滤过余极限交换。
因此可以假设层 $\mathcal{F}$ 如模，等式
(\ref{modules-equation-towards-constructible-sets}) 所述。

\medskip\noindent
假设有嵌入 $\mathcal{F} \subset \mathcal{G}$。则
$$
\Gamma(X, \mathcal{F}) =
\Gamma(Z, \mathcal{F}|_Z) \cap \Gamma(X, \mathcal{G})
$$
，其中交取在 $\Gamma(Z, \mathcal{G}|_Z)$ 内。
这由证明开头的观察得出：通过考察茎，可以检验 $\mathcal{G}$
的整体截面是否属于 $\mathcal{F}$，而 $X$ 的每个点都特化到 $Z$ 中的点。

\medskip\noindent
由模，引理 \ref{modules-lemma-constructible-in-constant}，存在单射
$\mathcal{F} \to \prod (Z_i \to X)_*\underline{S_i}$，
其中乘积有限，$Z_i \subset X$ 闭，且 $S_i$ 有限。
因此只需对层 $(Z_i \to X)_*\underline{S_i}$ 证明满射性。注意
$$
\Gamma(X, (Z_i \to X)_*\underline{S_i}) = \Gamma(Z_i, \underline{S_i})
\quad\text{且}\quad
\Gamma(X, (Z_i \to X)_*\underline{S_i}|_Z) =
\Gamma(Z \cap Z_i, \underline{S_i})
$$
此外，条件 (1) 和 (2) 由 $Z_i$ 继承；对 (2) 这是显然的，
对 (1) 则由拓扑，引理 \ref{topology-lemma-spectral-sub} 得出。
因此只需在（有限）常值层的情形证明引理。
这一情形正是引理 \ref{etale-cohomology-lemma-connected-topological} 的重新表述，
证明完成。
\end{proof}

\begin{example}
\label{etale-cohomology-example-quinard}
若 $X$ 不是谱空间，则引理 \ref{etale-cohomology-lemma-h0-topological} 不成立。
下面给出一个例子。设 $Y$ 是 $T_1$ 拓扑空间，$y \in Y$ 是非开点。
令 $X = Y \amalg \{ x \}$，并赋予如下拓扑：其闭集为
$\emptyset$、$\{y\}$ 以及所有 $F \amalg \{ x \}$，
其中 $F$ 是 $Y$ 的闭子集。于是 $Z = \{x, y\}$ 是 $X$ 的闭子集，
并满足引理 \ref{etale-cohomology-lemma-h0-topological} 的假设 (2)。
但 $X$ 连通而 $Z$ 不连通。因此，对取值于 $\{0, 1\}$ 的常值层
（在 $X$ 上），引理的结论不成立。
\end{example}

\begin{lemma}
\label{etale-cohomology-lemma-h0-henselian-pair}
设 $(A, I)$ 是 Hensel 对。置 $X = \Spec(A)$ 和
$Z = \Spec(A/I)$。对任意层 $\mathcal{F}$（在 $X_\etale$ 上），有
$\Gamma(X, \mathcal{F}) = \Gamma(Z, \mathcal{F}|_Z)$。
\end{lemma}

\begin{proof}
回忆，任意环的谱都是谱空间，见代数，引理
\ref{algebra-lemma-spec-spectral}。由代数进阶，引理
\ref{more-algebra-lemma-irreducible-henselian-pair-connected}
可知，$\overline{\{x\}} \cap Z$ 连通，其中 $x \in X$ 任意。
由引理 \ref{etale-cohomology-lemma-h0-topological}，该断言对 $X_{Zar}$ 上的层成立。
对任意有限态射 $X' \to X$，有 $X' = \Spec(A')$ 和
$Z \times_X X' = \Spec(A'/IA')$，且 $(A', IA')$ 是 Hensel 对，
见代数进阶，引理
\ref{more-algebra-lemma-integral-over-henselian-pair}
；故同一断言对 $(X')_{Zar}$ 上的层成立。
因此应用引理 \ref{etale-cohomology-lemma-gabber-h0} 即得结论。
\end{proof}

\noindent
最后，可以陈述并证明 Gabber 定理。

\begin{theorem}[Gabber]
\label{etale-cohomology-theorem-gabber}
设 $(A, I)$ 是 Hensel 对。置 $X = \Spec(A)$ 和
$Z = \Spec(A/I)$。对任意挠阿贝尔层 $\mathcal{F}$（在 $X_\etale$ 上），有
$H^q_\etale(X, \mathcal{F}) = H^q_\etale(Z, \mathcal{F}|_Z)$。
\end{theorem}

\begin{proof}
由引理 \ref{etale-cohomology-lemma-h0-henselian-pair}，结论对 $q = 0$ 成立。
设 $q \geq 1$。假设结论已在所有次数 $< q$ 中证明。
设 $\mathcal{F}$ 是挠阿贝尔层。令
$\mathcal{F} \to \mathcal{F}'$
为挠阿贝尔层的单射（稍后选取），其余核为 $\mathcal{Q}$，
从而有短正合列
$$
0 \to \mathcal{F} \to \mathcal{F}' \to \mathcal{Q} \to 0
$$
，这是 $X_\etale$ 上挠阿贝尔层的短正合列。它给出 $X$ 和 $Z$ 上
上同调长正合列之间的映射，其中一部分如下：
$$
\xymatrix{
H^{q - 1}_\etale(X, \mathcal{F}') \ar[d] \ar[r] &
H^{q - 1}_\etale(X, \mathcal{Q}) \ar[d] \ar[r] &
H^q_\etale(X, \mathcal{F}) \ar[d] \ar[r] &
H^q_\etale(X, \mathcal{F}') \ar[d] \\
H^{q - 1}_\etale(Z, \mathcal{F}'|_Z) \ar[r] &
H^{q - 1}_\etale(Z, \mathcal{Q}|_Z) \ar[r] &
H^q_\etale(Z, \mathcal{F}|_Z) \ar[r] &
H^q_\etale(Z, \mathcal{F}'|_Z)
}
$$
利用这个行正合的阿贝尔群交换图来完成证明。

\medskip\noindent
先证关于 $\mathcal{F}$ 的单射性。设 $\xi$ 是
$H^q_\etale(X, \mathcal{F})$ 的非零元素。把引理
\ref{etale-cohomology-lemma-efface-cohomology-on-closed-by-finite-cover} 应用于
$Z = X$（！），可以找到 $\mathcal{F} \subset \mathcal{F}'$，
使得 $\xi$ 向右的像为零。于是 $\xi$ 是
$H^{q - 1}_\etale(X, \mathcal{Q})$ 中某个元素的像，而对
$q - 1$ 的双射性说明 $\xi$ 在
$H^q_\etale(Z, \mathcal{F}|_Z)$ 中的像不为零。

\medskip\noindent
再证关于 $\mathcal{F}$ 的满射性。设 $\xi$ 是
$H^q_\etale(Z, \mathcal{F}|_Z)$ 的一个元素。把引理
\ref{etale-cohomology-lemma-efface-cohomology-on-closed-by-finite-cover} 应用于
$Z = Z$，可以找到 $\mathcal{F} \subset \mathcal{F}'$，
使得 $\xi$ 向右的像为零。于是 $\xi$ 是
$H^{q - 1}_\etale(Z, \mathcal{Q}|_Z)$ 中某个元素的像，而对
$q - 1$ 的双射性说明 $\xi$ 属于竖直映射的像。
\end{proof}

\begin{lemma}
\label{etale-cohomology-lemma-vanishing-restriction-injective}
设 $X$ 是具有仿射对角的概形，并可由 $n + 1$ 个仿射开集覆盖。
设 $Z \subset X$ 是闭子概形。设 $\mathcal{A}$ 是 $X_\etale$
上的挠环层，而 $\mathcal{I}$ 是 $\mathcal{A}$-模的内射层
（在 $X_\etale$ 上）。则
$H^q_\etale(Z, \mathcal{I}|_Z) = 0$ 对 $q > n$ 成立。
\end{lemma}

\begin{proof}
对 $n$ 作归纳。若 $n = 0$，则 $X$ 仿射。设
$X = \Spec(A)$ 且 $Z = \Spec(A/I)$。令 $A^h$ 为所有 \'etale
$A$-代数 $B$ 的滤过余极限，其中 $A/I \to B/IB$ 是同构。
于是 $(A^h, IA^h)$ 是 Hensel 对，且 $A/I = A^h/IA^h$；
见代数进阶，引理 \ref{more-algebra-lemma-henselization} 及其证明。
置 $X^h = \Spec(A^h)$。由定理 \ref{etale-cohomology-theorem-gabber}，有
$$
H^q_\etale(Z, \mathcal{I}|_Z) = H^q_\etale(X^h, \mathcal{I}|_{X^h})
$$
由定理 \ref{etale-cohomology-theorem-colimit}，有
$$
H^q_\etale(X^h, \mathcal{I}|_{X^h}) =
\colim_{A \to B} H^q_\etale(\Spec(B), \mathcal{I}|_{\Spec(B)})
$$
，其中余极限取遍上述 $A$-代数 $B$。
由于态射 $\Spec(B) \to \Spec(A)$ 是 \'etale 的，限制
$\mathcal{I}|_{\Spec(B)}$ 是 $\mathcal{A}|_{\Spec(B)}$-模的内射层
（位点上的上同调，引理
\ref{sites-cohomology-lemma-cohomology-of-open}）。
因此右边的上同调群为零，从而得到这一情形的结论。

\medskip\noindent
归纳步骤可用 Mayer--Vietoris 完成。具体地，假设
$X = U \cup V$，其中 $U$ 是 $n$ 个仿射开集之并，而 $V$ 仿射。
利用 $X$ 的对角是仿射的，可知 $U \cap V$ 是 $n$ 个仿射开集之并。
Mayer--Vietoris 给出正合列
$$
H^{q - 1}_\etale(U \cap V \cap Z, \mathcal{I}|_Z) \to
H^q_\etale(Z, \mathcal{I}|_Z) \to
H^q_\etale(U \cap Z, \mathcal{I}|_Z) \oplus
H^q_\etale(V \cap Z, \mathcal{I}|_Z)
$$
，由归纳假设，当 $q > n$ 时得到所需的消失。
\end{proof}





\section{曲线上挠层的上同调}
\label{etale-cohomology-section-vanishing-torsion}

\noindent
本节的目标是证明曲线上挠层上同调的基本有限性与消失结果，
见定理 \ref{etale-cohomology-theorem-vanishing-affine-curves}。
第 \ref{etale-cohomology-section-vanishing-torsion-coefficients} 节将讨论
Noether 环上挠模的可构造层。

\begin{situation}
\label{etale-cohomology-situation-what-to-prove}
这里 $k$ 是代数闭域，$X$ 是维数 $\leq 1$ 的分离有限型概形（在 $k$ 上），
而 $\mathcal{F}$ 是 $X_\etale$ 上的挠阿贝尔层。
\end{situation}

\noindent
在情形 \ref{etale-cohomology-situation-what-to-prove} 中，要证明下列断言：
\begin{enumerate}
\item
\label{etale-cohomology-item-vanishing}
$H^q_\etale(X, \mathcal{F}) = 0$ 对 $q > 2$ 成立；
\item
\label{etale-cohomology-item-vanishing-affine}
$H^q_\etale(X, \mathcal{F}) = 0$ 对 $q > 1$ 成立，若 $X$ 仿射；
\item
\label{etale-cohomology-item-vanishing-p-p}
$H^q_\etale(X, \mathcal{F}) = 0$ 对 $q > 1$ 成立，若
$p = \text{char}(k) > 0$ 且 $\mathcal{F}$ 是 $p$-幂挠的；
\item
\label{etale-cohomology-item-finite-prime-to-p}
$H^q_\etale(X, \mathcal{F})$ 有限，若 $\mathcal{F}$
可构造且是与 $\text{char}(k)$ 互素的挠层；
\item
\label{etale-cohomology-item-finite-proper}
$H^q_\etale(X, \mathcal{F})$ 有限，若 $X$ 真且
$\mathcal{F}$ 可构造；
\item
\label{etale-cohomology-item-base-change-prime-to-p}
$H^q_\etale(X, \mathcal{F}) \to
H^q_\etale(X_{k'}, \mathcal{F}|_{X_{k'}})$ 对代数闭域的任意扩张
$k'/k$ 都是同构，前提是 $\mathcal{F}$ 是与
$\text{char}(k)$ 互素的挠层；
\item
\label{etale-cohomology-item-base-change-proper}
$H^q_\etale(X, \mathcal{F}) \to
H^q_\etale(X_{k'}, \mathcal{F}|_{X_{k'}})$ 对代数闭域的任意扩张
$k'/k$ 都是同构，前提是 $X$ 真；
\item
\label{etale-cohomology-item-surjective}
$H^2_\etale(X, \mathcal{F}) \to H^2_\etale(U, \mathcal{F})$
对所有开子集 $U \subset X$ 都是满射。
\end{enumerate}
给定任意情形 \ref{etale-cohomology-situation-what-to-prove}，若适用于该情形的断言都为真，
就说“断言 (\ref{etale-cohomology-item-vanishing})--(\ref{etale-cohomology-item-surjective}) 成立”。
证明从常系数上同调计算的如下推论开始。

\begin{lemma}
\label{etale-cohomology-lemma-constant-smooth-statements}
在情形 \ref{etale-cohomology-situation-what-to-prove} 中，假设 $X$ 光滑，且
$\mathcal{F} = \underline{\mathbf{Z}/\ell\mathbf{Z}}$，其中 $\ell$ 是素数。
则断言 (\ref{etale-cohomology-item-vanishing})--(\ref{etale-cohomology-item-surjective})
对 $\mathcal{F}$ 成立。
\end{lemma}

\begin{proof}
由于 $X$ 光滑，$X$ 是光滑曲线与点的有限不交并。
点（$k$ 的谱）上任意层的高阶 \'etale 上同调平凡，例如可由
引理 \ref{etale-cohomology-lemma-compare-cohomology-point}
或 \ref{etale-cohomology-lemma-affine-only-closed-points} 得到。
因此可以假设 $X$ 是光滑曲线。

\medskip\noindent
情形 I：$\ell$ 不同于 $k$ 的特征。
这一情形由引理 \ref{etale-cohomology-lemma-cohomology-smooth-projective-curve}
（射影情形）和引理 \ref{etale-cohomology-lemma-vanishing-cohomology-mu-smooth-curve}
（仿射情形）得到。关于上同调和代数闭基域扩张的断言
(\ref{etale-cohomology-item-base-change-prime-to-p}) 来自如下事实：亏格 $g$
与“穿孔”数 $r$ 在从 $k$ 到 $k'$ 时不变。
断言 (\ref{etale-cohomology-item-surjective}) 也成立，因为
$H^2_\etale(U, \mathcal{F})$ 只要 $U \not = X$ 就为零；此时 $U$ 仿射
（簇，引理 \ref{varieties-lemma-proper-minus-point} 和
\ref{varieties-lemma-curve-affine-projective}）。

\medskip\noindent
情形 II：$\ell$ 等于 $k$ 的特征。
消失性由引理 \ref{etale-cohomology-lemma-vanishing-variety-char-p-p} 得到。
断言 (\ref{etale-cohomology-item-finite-proper}) 和 (\ref{etale-cohomology-item-base-change-proper})
由引理 \ref{etale-cohomology-lemma-finiteness-proper-variety-char-p-p} 得到。
\end{proof}

\begin{remark}[代数闭基域扩张下的不变性]
\label{etale-cohomology-remark-invariance}
设 $k$ 是特征 $p > 0$ 的代数闭域。在第
\ref{etale-cohomology-section-artin-schreier} 节中已经看到正合列
$$
k[x] \to k[x] \to
H^1_\etale(\mathbf{A}^1_k, \mathbf{Z}/p\mathbf{Z}) \to 0
$$
，其中第一个箭头把 $f(x)$ 送到 $f^p - f$。余核的一组代表元由多项式
$$
\sum\nolimits_{p \not | n} \lambda_n x^n
$$
组成，其中 $\lambda_n \in k$。（若 $k$ 不是代数闭的，
还必须加入某些常数。）特别地，当 $k'/k$ 是代数闭扩张时，映射
$$
H^1_\etale(\mathbf{A}^1_k, \mathbf{Z}/p\mathbf{Z})
\to
H^1_\etale(\mathbf{A}^1_{k'}, \mathbf{Z}/p\mathbf{Z})
$$
一般不是同构。特别地，\'etale 基本群之间的映射
$\pi_1(\mathbf{A}^1_{k'}) \to \pi_1(\mathbf{A}^1_k)$
（在此插入未来的参考文献）也不是同构。因此，仿射直线的
\'etale 同伦型依赖于代数闭基域。由上面的引理
\ref{etale-cohomology-lemma-constant-smooth-statements} 可知，这一现象只在特征
$p$ 且系数为 $p$-幂挠时发生。
\end{remark}

\begin{lemma}
\label{etale-cohomology-lemma-ses-statements}
设 $k$ 是代数闭域。设 $X$ 是 $k$ 上维数 $\leq 1$ 的分离有限型概形。
设
$0 \to \mathcal{F}_1 \to \mathcal{F} \to \mathcal{F}_2 \to 0$
是 $X$ 上挠阿贝尔层的短正合列。若断言
(\ref{etale-cohomology-item-vanishing})--(\ref{etale-cohomology-item-surjective}) 对
$\mathcal{F}_1$ 和 $\mathcal{F}_2$ 成立，则它们也对 $\mathcal{F}$ 成立。
\end{lemma}

\begin{proof}
这大都直接来自定义和上同调长正合列。还要注意，$\mathcal{F}$
可构造（相应地，是与 $k$ 的特征互素的挠层）当且仅当
$\mathcal{F}_1$ 与 $\mathcal{F}_2$ 都可构造
（相应地，是与 $k$ 的特征互素的挠层）。见命题
\ref{etale-cohomology-proposition-constructible-over-noetherian}。略去若干细节。
\end{proof}

\begin{lemma}
\label{etale-cohomology-lemma-finite-pushforward-statements}
设 $k$ 是代数闭域。设 $f : X \to Y$ 是 $k$ 上维数 $\leq 1$
的分离有限型概形之间的有限态射。设 $\mathcal{F}$ 是 $X$ 上的挠阿贝尔层。
若断言 (\ref{etale-cohomology-item-vanishing})--(\ref{etale-cohomology-item-surjective})
对 $\mathcal{F}$ 成立，则它们也对 $f_*\mathcal{F}$ 成立。
\end{lemma}

\begin{proof}
事实上，有 $H^q_\etale(X, \mathcal{F}) = H^q_\etale(Y, f_*\mathcal{F})$，
这是由 $R^qf_*$ 对 $q > 0$ 的消失
（命题 \ref{etale-cohomology-proposition-finite-higher-direct-image-zero}）和 Leray 谱序列
（位点上的上同调，引理 \ref{sites-cohomology-lemma-apply-Leray}）得到的。
对 (\ref{etale-cohomology-item-surjective})，使用 $f_*$ 的形成与任意基变换交换这一事实
（引理 \ref{etale-cohomology-lemma-finite-pushforward-commutes-with-base-change}）。
\end{proof}

\begin{lemma}
\label{etale-cohomology-lemma-restrict-to-open}
在情形 \ref{etale-cohomology-situation-what-to-prove} 中，假设 $\mathcal{F}$ 可构造。
设 $j : X' \to X$ 是稠密开子概形的包含。则断言
(\ref{etale-cohomology-item-vanishing})--(\ref{etale-cohomology-item-surjective}) 对 $\mathcal{F}$ 成立
当且仅当它们对 $j_!j^{-1}\mathcal{F}$ 成立。
\end{lemma}

\begin{proof}
由于 $X'$ 稠密，$Z = X \setminus X'$ 的维数为 $0$，
故它是有限集 $Z = \{x_1, \ldots, x_n\}$，其元素都是 $k$-有理点。
考虑短正合列
$$
0 \to j_!j^{-1}\mathcal{F} \to \mathcal{F} \to i_*i^{-1}\mathcal{F} \to 0
$$
，见引理 \ref{etale-cohomology-lemma-ses-associated-to-open}。注意
$H^q_\etale(X, i_*i^{-1}\mathcal{F}) = H^q_\etale(Z, i^*\mathcal{F})$.
事实上，$i : Z \to X$ 是闭浸入，因而有限；于是由命题
\ref{etale-cohomology-proposition-finite-higher-direct-image-zero}，
$R^qi_*$ 对 $q > 0$ 消失，故该等式由 Leray 谱序列得到
（位点上的上同调，引理 \ref{sites-cohomology-lemma-apply-Leray}）。
由于 $Z$ 是代数闭域之谱的不交并，可知
$H^q_\etale(Z, i^*\mathcal{F}) = 0$ 对 $q > 0$ 成立，并且
$$
H^0_\etale(Z, i^{-1}\mathcal{F}) =
\bigoplus\nolimits_{i = 1, \ldots, n} \mathcal{F}_{x_i}
$$
是有限的，因为 $\mathcal{F}_{x_i}$ 由 $\mathcal{F}$
可构造这一假设而有限。上同调长正合列给出正合列
$$
0 \to
H^0_\etale(X, j_!j^{-1}\mathcal{F}) \to
H^0_\etale(X, \mathcal{F}) \to
H^0_\etale(Z, i^{-1}\mathcal{F}) \to
H^1_\etale(X, j_!j^{-1}\mathcal{F}) \to
H^1_\etale(X, \mathcal{F}) \to 0
$$
以及同构 $H^q_\etale(X, j_!j^{-1}\mathcal{F}) \to
H^q_\etale(X, \mathcal{F})$，对 $q > 1$ 成立。

\medskip\noindent
此时容易推出，(\ref{etale-cohomology-item-vanishing})--(\ref{etale-cohomology-item-surjective})
中的每个断言对 $\mathcal{F}$ 成立当且仅当它对
$j_!j^{-1}\mathcal{F}$ 成立。补充几点以帮助读者：
(a) 若 $\mathcal{F}$ 是与 $k$ 的特征互素的挠层，则
$j_!j^{-1}\mathcal{F}$ 也是；
(b) 层 $j_!j^{-1}\mathcal{F}$ 可构造；
(c) 有 $H^0_\etale(Z, i^{-1}\mathcal{F}) =
H^0_\etale(Z_{k'}, i^{-1}\mathcal{F}|_{Z_{k'}})$；以及 (d)
若 $U \subset X$ 是开集，则 $U' = U \cap X'$ 在 $U$ 中稠密。
\end{proof}

\begin{lemma}
\label{etale-cohomology-lemma-even-easier}
在情形 \ref{etale-cohomology-situation-what-to-prove} 中，假设 $X$ 光滑。
设 $j : U \to X$ 是开浸入。设 $\ell$ 是素数，并令
$\mathcal{F} = j_!\underline{\mathbf{Z}/\ell\mathbf{Z}}$。
则断言 (\ref{etale-cohomology-item-vanishing})--(\ref{etale-cohomology-item-surjective})
对 $\mathcal{F}$ 成立。
\end{lemma}

\begin{proof}
由于 $X$ 光滑，它是光滑曲线的不交并，故可以假设 $X$
是一条曲线（即不可约）。此时要么 $U = \emptyset$，无需证明；
要么 $U \subset X$ 稠密。在后一情形，引理由
引理 \ref{etale-cohomology-lemma-constant-smooth-statements} 和
\ref{etale-cohomology-lemma-restrict-to-open} 得到。
\end{proof}

\begin{lemma}
\label{etale-cohomology-lemma-somewhat-easier}
在情形 \ref{etale-cohomology-situation-what-to-prove} 中，假设 $X$ 约化。
设 $j : U \to X$ 是开浸入。设 $\ell$ 是素数，并令
$\mathcal{F} = j_!\underline{\mathbf{Z}/\ell\mathbf{Z}}$。
则断言 (\ref{etale-cohomology-item-vanishing})--(\ref{etale-cohomology-item-surjective})
对 $\mathcal{F}$ 成立。
\end{lemma}

\begin{proof}
与引理 \ref{etale-cohomology-lemma-even-easier} 的区别是，这里不假设 $X$ 光滑。
设 $\nu : X^\nu \to X$ 是正规化态射。则 $\nu$ 有限
（簇，引理 \ref{varieties-lemma-normalization-locally-algebraic}），
且 $X^\nu$ 光滑
（簇，引理 \ref{varieties-lemma-regular-point-on-curve}）。
设 $j^\nu : U^\nu \to X^\nu$ 是 $U$ 的逆像。
由引理 \ref{etale-cohomology-lemma-even-easier}，结论对
$j^\nu_!\underline{\mathbf{Z}/\ell\mathbf{Z}}$ 成立。
由引理 \ref{etale-cohomology-lemma-finite-pushforward-statements}，结论对
$\nu_*j^\nu_!\underline{\mathbf{Z}/\ell\mathbf{Z}}$ 成立。
一般而言，
$\nu_*j^\nu_!\underline{\mathbf{Z}/\ell\mathbf{Z}}$ 不等于
$j_!\underline{\mathbf{Z}/\ell\mathbf{Z}}$，
但可以如下绕过这一点。由于 $X$ 约化，态射 $\nu : X^\nu \to X$
在某个稠密开集 $j' : X' \to X$ 上是同构
（簇，引理 \ref{varieties-lemma-normalization-locally-algebraic}）。
在这个开集上有等式
$$
(j')^{-1}(\nu_*j^\nu_!\underline{\mathbf{Z}/\ell\mathbf{Z}})
=
(j')^{-1}(j_!\underline{\mathbf{Z}/\ell\mathbf{Z}})
$$
。对 $j' : X' \to X$ 和上述两个层使用两次
引理 \ref{etale-cohomology-lemma-restrict-to-open} 即得结论。
\end{proof}

\begin{lemma}
\label{etale-cohomology-lemma-vanishing-easier}
在情形 \ref{etale-cohomology-situation-what-to-prove} 中，假设 $X$ 约化。
设 $j : U \to X$ 是开浸入，且 $U$ 连通。设 $\ell$ 是素数。
设 $\mathcal{G}$ 是 $\mathbf{F}_\ell$-向量空间的有限局部常值层
（在 $U$ 上）。
令 $\mathcal{F} = j_!\mathcal{G}$。则断言
(\ref{etale-cohomology-item-vanishing})--(\ref{etale-cohomology-item-surjective}) 对 $\mathcal{F}$ 成立。
\end{lemma}

\begin{proof}
设 $f : V \to U$ 是引理 \ref{etale-cohomology-lemma-pullback-filtered} 中那样的有限
\'etale 态射，其次数与 $\ell$ 互素。第 \ref{etale-cohomology-section-trace-method} 节
的讨论给出映射
$$
\mathcal{G} \to f_*f^{-1}\mathcal{G} \to \mathcal{G}
$$
，其复合是同构。因此只需对
$\mathcal{F} = j_!f_*f^{-1}\mathcal{G}$ 证明引理。
由 Zariski 主定理（态射进阶，引理
\ref{more-morphisms-lemma-quasi-finite-separated-pass-through-finite})
，可以选取图
$$
\xymatrix{
V \ar[r]_{j'} \ar[d]_f & Y \ar[d]^{\overline{f}} \\
U \ar[r]^j & X
}
$$
，其中 $\overline{f} : Y \to X$ 有限，而 $j'$ 是像稠密的开浸入。
可以把 $Y$ 替换为其约化（这不改变 $V$，因为 $V$
作为 $U$ 上的 \'etale 概形是约化的）。
由于 $f$ 有限且 $V$ 在 $Y$ 中稠密，有 $V = U \times_X Y$。
由引理 \ref{etale-cohomology-lemma-compatible-shriek-push-finite}，有
$$
j_!f_*f^{-1}\mathcal{G} = \overline{f}_*j'_!f^{-1}\mathcal{G}
$$
由引理 \ref{etale-cohomology-lemma-finite-pushforward-statements}，只需考虑
$j'_!f^{-1}\mathcal{G}$。利用引理 \ref{etale-cohomology-lemma-pullback-filtered}
给出的滤过、$j'_!$ 的正合性以及引理 \ref{etale-cohomology-lemma-ses-statements}，
归结到
$\mathcal{F} = j'_!\underline{\mathbf{Z}/\ell\mathbf{Z}}$
的情形，这正是引理 \ref{etale-cohomology-lemma-somewhat-easier}。
\end{proof}


%10.20.09

\begin{theorem}
\label{etale-cohomology-theorem-vanishing-affine-curves}
若 $k$ 是代数闭域，$X$ 是维数 $\leq 1$ 的分离有限型概形（在 $k$ 上），
而 $\mathcal{F}$ 是 $X_\etale$ 上的挠阿贝尔层，则
\begin{enumerate}
\item
$H^q_\etale(X, \mathcal{F}) = 0$ 对 $q > 2$ 成立；
\item
$H^q_\etale(X, \mathcal{F}) = 0$ 对 $q > 1$ 成立，若 $X$ 仿射；
\item
$H^q_\etale(X, \mathcal{F}) = 0$ 对 $q > 1$ 成立，若
$p = \text{char}(k) > 0$ 且 $\mathcal{F}$ 是 $p$-幂挠的；
\item
$H^q_\etale(X, \mathcal{F})$ 有限，若 $\mathcal{F}$
可构造且是与 $\text{char}(k)$ 互素的挠层；
\item
$H^q_\etale(X, \mathcal{F})$ 有限，若 $X$ 真且
$\mathcal{F}$ 可构造；
\item
$H^q_\etale(X, \mathcal{F}) \to
H^q_\etale(X_{k'}, \mathcal{F}|_{X_{k'}})$ 对代数闭域的任意扩张
$k'/k$ 都是同构，前提是 $\mathcal{F}$ 是与
$\text{char}(k)$ 互素的挠层；
\item
$H^q_\etale(X, \mathcal{F}) \to
H^q_\etale(X_{k'}, \mathcal{F}|_{X_{k'}})$ 对代数闭域的任意扩张
$k'/k$ 都是同构，前提是 $X$ 真；
\item
$H^2_\etale(X, \mathcal{F}) \to H^2_\etale(U, \mathcal{F})$
对所有开子集 $U \subset X$ 都是满射。
\end{enumerate}
\end{theorem}

\begin{proof}
定理所说的正是：在情形 \ref{etale-cohomology-situation-what-to-prove} 中，
断言 (\ref{etale-cohomology-item-vanishing})--(\ref{etale-cohomology-item-surjective}) 成立。
第一步是把 $X$ 替换为其约化；命题
\ref{etale-cohomology-proposition-topological-invariance} 说明这是允许的。
由引理 \ref{etale-cohomology-lemma-torsion-colimit-constructible}，可以把
$\mathcal{F}$ 写成可构造阿贝尔层的滤过余极限。
取上同调与余极限交换，见引理 \ref{etale-cohomology-lemma-colimit}。
此外，沿 $X_{k'} \to X$ 的拉回作为左伴随与余极限交换。
因此只需对可构造层证明这些断言。

\medskip\noindent
本段中不再另行说明地使用引理 \ref{etale-cohomology-lemma-ses-statements}。写作
$\mathcal{F} = \mathcal{F}_1 \oplus \ldots \oplus \mathcal{F}_r$
，其中 $\mathcal{F}_i$ 对某个素数 $\ell_i$ 是 $\ell_i$-准素的，
便可假设 $\ell^n$ 零化 $\mathcal{F}$，其中 $\ell$ 为某个素数。
现在考虑正合列
$$
0 \to \mathcal{F}[\ell] \to \mathcal{F} \to \mathcal{F}/\mathcal{F}[\ell] \to 0.
$$
由此只需假设 $\mathcal{F}$ 是 $\ell$-挠的。这意味着
$\mathcal{F}$ 是 $\mathbf{F}_\ell$-向量空间的可构造层，
其中 $\ell$ 为某个素数。

\medskip\noindent
按定义，这意味着存在稠密开集 $U \subset X$，使得
$\mathcal{F}|_U$ 是 $\mathbf{F}_\ell$-向量空间的有限局部常值层。
由于 $\dim(X) \leq 1$，缩小 $U$ 后，可以假设
$U = U_1 \amalg \ldots \amalg U_n$ 是不可约概形的不交并
（只需去掉位于 $\geq 2$ 个分支之交中的 $U$ 的闭点）。
由引理 \ref{etale-cohomology-lemma-restrict-to-open}，归结到
$\mathcal{F} = j_!\mathcal{G}$ 的情形，其中 $\mathcal{G}$
是 $\mathbf{F}_\ell$-向量空间的有限局部常值层，定义在 $U$ 上。

\medskip\noindent
由于所选 $U = U_1 \amalg \ldots \amalg U_n$ 中每个 $U_i$ 都不可约，有
$$
j_!\mathcal{G} =
j_{1!}(\mathcal{G}|_{U_1}) \oplus \ldots \oplus
j_{n!}(\mathcal{G}|_{U_n})
$$
，其中 $j_i : U_i \to X$ 是包含态射。
$j_{i!}(\mathcal{G}|_{U_i})$ 的情形由引理
\ref{etale-cohomology-lemma-vanishing-easier} 处理。
\end{proof}

\begin{theorem}
\label{etale-cohomology-theorem-vanishing-curves}
设 $X$ 是维数为 $1$ 的有限型概形（在代数闭域 $k$ 上）。
设 $\mathcal{F}$ 是 $X_\etale$ 上的挠层。则
$$
H_\etale^q(X, \mathcal{F}) = 0, \quad \forall q \geq 3.
$$
若 $X$ 仿射，则还有 $H_\etale^2(X, \mathcal{F}) = 0$。
\end{theorem}

\begin{proof}
若 $X$ 分离，结论立即来自更精确的定理
\ref{etale-cohomology-theorem-vanishing-affine-curves}。
若 $X$ 不分离，选取仿射开覆盖
$X = X_1 \cup \ldots \cup X_n$。对 $n$ 归纳，可以假设消失性在
$U = X_1 \cup \ldots \cup X_{n - 1}$ 上成立。
于是 Mayer--Vietoris（引理 \ref{etale-cohomology-lemma-mayer-vietoris}）给出
$$
H^2_\etale(U, \mathcal{F}) \oplus H^2_\etale(X_n, \mathcal{F}) \to
H^2_\etale(U \cap X_n, \mathcal{F}) \to
H^3_\etale(X, \mathcal{F}) \to 0
$$
然而，$U \cap X_n$ 是仿射概形的开子概形，并且由维数假设它也仿射，
故群 $H^2_\etale(U \cap X_n, \mathcal{F})$ 由定理
\ref{etale-cohomology-theorem-vanishing-affine-curves} 消失。
\end{proof}

\begin{lemma}
\label{etale-cohomology-lemma-base-change-dim-1-separably-closed}
设 $k'/k$ 是可分闭域的扩张。设 $X$ 是 $k$ 上维数
$\leq 1$ 的真概形。设 $\mathcal{F}$ 是 $X$ 上的挠阿贝尔层。
则映射 $H^q_\etale(X, \mathcal{F}) \to
H^q_\etale(X_{k'}, \mathcal{F}|_{X_{k'}})$ 对 $q \geq 0$ 是同构。
\end{lemma}

\begin{proof}
定理 \ref{etale-cohomology-theorem-vanishing-affine-curves} 已证明代数闭域的情形。
给定引理中的 $k \subset k'$，可以选取图
$$
\xymatrix{
k' \ar[r] & \overline{k}' \\
k \ar[u] \ar[r] & \overline{k} \ar[u]
}
$$
，其中 $k \subset \overline{k}$ 和 $k' \subset \overline{k}'$
分别为代数闭包。由于 $k$ 与 $k'$ 可分闭，域扩张
$\overline{k}/k$ 与 $\overline{k}'/k'$
是代数纯不可分的。此时态射
$X_{\overline{k}} \to X$ 与 $X_{\overline{k}'} \to X_{k'}$
是普遍同胚。因此，$\mathcal{F}$ 的上同调可在
$X_{\overline{k}}$ 上计算，而 $\mathcal{F}|_{X_{k'}}$ 的上同调可在
$X_{\overline{k}'}$ 上计算，见命题
\ref{etale-cohomology-proposition-topological-invariance}。
故一般情形由代数闭域的情形推出。
\end{proof}






\section{曲线上挠模的上同调}
\label{etale-cohomology-section-vanishing-torsion-coefficients}

\noindent
本节对 Noether 环上的挠模可构造层重复第
\ref{etale-cohomology-section-vanishing-torsion} 节的论证。先从最有意思的一步开始。

\begin{lemma}
\label{etale-cohomology-lemma-constant-statements-coefficients}
设 $\Lambda$ 是 Noether 环，$M$ 是有限 $\Lambda$-模且被整数
$n > 0$ 零化，$k$ 是代数闭域，而 $X$ 是维数
$\leq 1$ 的分离有限型概形（在 $k$ 上）。则
\begin{enumerate}
\item $H^q_\etale(X, \underline{M})$ 是有限 $\Lambda$-模，
若 $n$ 与 $\text{char}(k)$ 互素；
\item $H^q_\etale(X, \underline{M})$ 是有限 $\Lambda$-模，
若 $X$ 真。
\end{enumerate}
\end{lemma}

\begin{proof}
若 $n = \ell n'$，其中 $\ell$ 是某个素数，则得到短正合列
$0 \to M[\ell] \to M \to M' \to 0$（有限 $\Lambda$-模），
且 $M'$ 被 $n'$ 零化。这给出相应的常值层短正合列，
进而给出上同调模的正合列
$$
H^q_\etale(X, \underline{M[n]}) \to
H^q_\etale(X, \underline{M}) \to
H^q_\etale(X, \underline{M'})
$$
因此，只要能证明 $M$ 被素数零化的情形，就可对 $n$ 归纳而得结论。

\medskip\noindent
设 $\ell$ 是素数，且 $\ell$ 零化 $M$。于是可把 $\Lambda$ 替换为
$\mathbf{F}_\ell$-代数 $\Lambda/\ell \Lambda$。事实上，
$\mathcal{F}$ 作为 $\Lambda$-模层的上同调与
$\mathcal{F}$ 作为 $\Lambda/\ell \Lambda$-模层的上同调相同；
例如可由位点上的上同调，引理
\ref{sites-cohomology-lemma-cohomology-modules-abelian-agree} 得到。

\medskip\noindent
假设 $\ell$ 是素数，并且 $\ell$ 同时零化 $M$ 和 $\Lambda$。
把问题归结到 $M$ 是有限自由 $\Lambda$-模的情形。
事实上，选取短正合列
$$
0 \to N \to \Lambda^{\oplus m} \to M \to 0
$$
这确定了正合列
$$
H^q_\etale(X, \underline{\Lambda^{\oplus m}}) \to
H^q_\etale(X, \underline{M}) \to
H^{q + 1}_\etale(X, \underline{N})
$$
对 $q$ 下降归纳便得到 $M$ 的结论，只要已知
$\Lambda^{\oplus m}$ 的结论。这里使用定理
\ref{etale-cohomology-theorem-vanishing-affine-curves}
所给出的上同调群在次数 $> 2$ 时消失。

\medskip\noindent
设 $\ell$ 是素数，并假设 $\ell$ 零化 $\Lambda$。尚需证明上同调群
$H^q_\etale(X, \underline{\Lambda})$ 是有限 $\Lambda$-模。
这里用一个技巧证明；“正确”的论证使用稍后才会证明的系数定理。
选取基 $\Lambda = \bigoplus_{i \in I} \mathbf{F}_\ell e_i$，
使得 $e_0 = 1$，其中 $0 \in I$。这一基的选取确定层同构
$$
\underline{\Lambda} = \bigoplus \underline{\mathbf{F}_\ell} e_i
$$
（在 $X_\etale$ 上）。因此
$$
H^q_\etale(X, \underline{\Lambda}) =
H^q_\etale(X, \bigoplus \underline{\mathbf{F}_\ell} e_i) =
\bigoplus H^q_\etale(X, \underline{\mathbf{F}_\ell})e_i
$$
，因为由定理 \ref{etale-cohomology-theorem-colimit}
（或引理 \ref{etale-cohomology-lemma-colimit}，或
引理 \ref{etale-cohomology-lemma-direct-sum-bounded-below-cohomology}），
在 $X$ 上取上同调与直和交换。已经知道
$H^q_\etale(X, \underline{\mathbf{F}_\ell})$ 是有限维
$\mathbf{F}_\ell$-向量空间
（由定理 \ref{etale-cohomology-theorem-vanishing-affine-curves}），
故 $H^q_\etale(X, \underline{\Lambda})$ 是同秩的自由
$\Lambda$-模。具体而言，给定一组基
$\xi_1, \ldots, \xi_m$（属于
$H^q_\etale(X, \underline{\mathbf{F}_\ell})$），则
$\xi_1 e_0, \ldots, \xi_m e_0$ 构成一组 $\Lambda$-基
（属于 $H^q_\etale(X, \underline{\Lambda})$）。
\end{proof}

\begin{lemma}
\label{etale-cohomology-lemma-finite-pushforward-coefficients}
设 $\Lambda$ 是 Noether 环，$k$ 是代数闭域，
$f : X \to Y$ 是 $k$ 上维数 $\leq 1$ 的分离有限型概形之间的有限态射，
而 $\mathcal{F}$ 是 $\Lambda$-模层（在 $X_\etale$ 上）。若
$H^q_\etale(X, \mathcal{F})$ 是有限 $\Lambda$-模，则
$H^q_\etale(Y, f_*\mathcal{F})$ 也如此。
\end{lemma}

\begin{proof}
事实上，有 $H^q_\etale(X, \mathcal{F}) = H^q_\etale(Y, f_*\mathcal{F})$，
这是由 $R^qf_*$ 对 $q > 0$ 的消失
（命题 \ref{etale-cohomology-proposition-finite-higher-direct-image-zero}）和 Leray 谱序列
（位点上的上同调，引理 \ref{sites-cohomology-lemma-apply-Leray}）得到的。
\end{proof}

\begin{lemma}
\label{etale-cohomology-lemma-restrict-to-open-coefficients}
设 $\Lambda$ 是 Noether 环，$k$ 是代数闭域，
$X$ 是 $k$ 上维数 $\leq 1$ 的分离有限型概形，
$\mathcal{F}$ 是可构造 $\Lambda$-模层（在 $X_\etale$ 上），
而 $j : X' \to X$ 是稠密开子概形的包含。则
$H^q_\etale(X, \mathcal{F})$ 是有限 $\Lambda$-模，当且仅当
$H^q_\etale(X, j_!j^{-1}\mathcal{F})$ 是有限 $\Lambda$-模。
\end{lemma}

\begin{proof}
由于 $X'$ 稠密，$Z = X \setminus X'$ 的维数为 $0$，
故它是有限集 $Z = \{x_1, \ldots, x_n\}$，其元素都是 $k$-有理点。
考虑短正合列
$$
0 \to j_!j^{-1}\mathcal{F} \to \mathcal{F} \to i_*i^{-1}\mathcal{F} \to 0
$$
，见引理 \ref{etale-cohomology-lemma-ses-associated-to-open}。注意
$H^q_\etale(X, i_*i^{-1}\mathcal{F}) = H^q_\etale(Z, i^*\mathcal{F})$.
事实上，$i : Z \to X$ 是闭浸入，因而有限；于是由命题
\ref{etale-cohomology-proposition-finite-higher-direct-image-zero}，
$R^qi_*$ 对 $q > 0$ 消失，故该等式由 Leray 谱序列得到
（位点上的上同调，引理 \ref{sites-cohomology-lemma-apply-Leray}）。
由于 $Z$ 是代数闭域之谱的不交并，可知
$H^q_\etale(Z, i^*\mathcal{F}) = 0$ 对 $q > 0$ 成立，并且
$$
H^0_\etale(Z, i^{-1}\mathcal{F}) =
\bigoplus\nolimits_{i = 1, \ldots, n} \mathcal{F}_{x_i}
$$
是有限 $\Lambda$-模；$\mathcal{F}_{x_i}$ 的有限性来自
$\mathcal{F}$ 是可构造 $\Lambda$-模层这一假设。
上同调长正合列给出正合列
$$
0 \to
H^0_\etale(X, j_!j^{-1}\mathcal{F}) \to
H^0_\etale(X, \mathcal{F}) \to
H^0_\etale(Z, i^{-1}\mathcal{F}) \to
H^1_\etale(X, j_!j^{-1}\mathcal{F}) \to
H^1_\etale(X, \mathcal{F}) \to 0
$$
以及同构 $H^0_\etale(X, j_!j^{-1}\mathcal{F}) \to
H^0_\etale(X, \mathcal{F})$，对 $q > 1$ 成立。引理由此容易推出。
\end{proof}

\begin{lemma}
\label{etale-cohomology-lemma-somewhat-easier-coefficients}
设 $\Lambda$ 是 Noether 环，$M$ 是有限 $\Lambda$-模且被整数
$n > 0$ 零化，$k$ 是代数闭域，$X$ 是维数
$\leq 1$ 的分离有限型概形（在 $k$ 上），而
$j : U \to X$ 是开浸入。则
\begin{enumerate}
\item $H^q_\etale(X, j_!\underline{M})$ 是有限 $\Lambda$-模，
若 $n$ 与 $\text{char}(k)$ 互素；
\item $H^q_\etale(X, j_!\underline{M})$ 是有限 $\Lambda$-模，
若 $X$ 真。
\end{enumerate}
\end{lemma}

\begin{proof}
由于 $\dim(X) \leq 1$，存在开集 $V \subset X$，它与 $U$ 不交，
使得 $X' = U \cup V$ 是 $X$ 中的稠密开集（略去细节）。
若 $j' : X' \to X$ 表示包含态射，则
$j_!\underline{M}$ 是 $j'_!\underline{M}$ 的直和项。
因此只需证明 $U$ 在 $X$ 中开且稠密的情形。这一情形由引理
\ref{etale-cohomology-lemma-restrict-to-open-coefficients} 和
\ref{etale-cohomology-lemma-constant-statements-coefficients} 得到。
\end{proof}

\begin{lemma}
\label{etale-cohomology-lemma-ses-statements-coefficients}
设 $\Lambda$ 是 Noether 环，$k$ 是代数闭域，
$X$ 是 $k$ 上维数 $\leq 1$ 的分离有限型概形，而
$0 \to \mathcal{F}_1 \to \mathcal{F} \to \mathcal{F}_2 \to 0$
是 $\Lambda$-模层的短正合列（在 $X_\etale$ 上）。若
$H^q_\etale(X, \mathcal{F}_i)$（$i = 1, 2$）是有限
$\Lambda$-模，则 $H^q_\etale(X, \mathcal{F})$ 是有限
$\Lambda$-模。
\end{lemma}

\begin{proof}
这直接来自上同调长正合列。
\end{proof}

\begin{lemma}
\label{etale-cohomology-lemma-vanishing-easier-coefficients}
设 $\Lambda$ 是 Noether 环，$k$ 是代数闭域，
$X$ 是维数 $\leq 1$ 的分离有限型概形（在 $k$ 上），
$j : U \to X$ 是开浸入且 $U$ 连通，$\ell$ 是素数，
$n > 0$，而 $\mathcal{G}$ 是有限型局部常值
$\Lambda$-模层（在 $U_\etale$ 上），并被 $\ell^n$ 零化。则
\begin{enumerate}
\item $H^q_\etale(X, j_!\mathcal{G})$ 是有限 $\Lambda$-模，
若 $\ell$ 与 $\text{char}(k)$ 互素；
\item $H^q_\etale(X, j_!\mathcal{G})$ 是有限 $\Lambda$-模，
若 $X$ 真。
\end{enumerate}
\end{lemma}

\begin{proof}
设 $f : V \to U$ 是引理 \ref{etale-cohomology-lemma-pullback-filtered-modules}
中那样的有限 \'etale 态射，其次数与 $\ell$ 互素。
第 \ref{etale-cohomology-section-trace-method} 节的讨论给出映射
$$
\mathcal{G} \to f_*f^{-1}\mathcal{G} \to \mathcal{G}
$$
，其复合是同构。因此只需证明
$H^q_\etale(X, j_!f_*f^{-1}\mathcal{G})$ 的有限性。
由 Zariski 主定理（态射进阶，引理
\ref{more-morphisms-lemma-quasi-finite-separated-pass-through-finite}），
可以选取图
$$
\xymatrix{
V \ar[r]_{j'} \ar[d]_f & Y \ar[d]^{\overline{f}} \\
U \ar[r]^j & X
}
$$
，其中 $\overline{f} : Y \to X$ 有限，而 $j'$ 是像稠密的开浸入。
由于 $f$ 有限且 $V$ 在 $Y$ 中稠密，有 $V = U \times_X Y$。
由引理 \ref{etale-cohomology-lemma-compatible-shriek-push-finite}，有
$$
j_!f_*f^{-1}\mathcal{G} = \overline{f}_*j'_!f^{-1}\mathcal{G}
$$
由引理 \ref{etale-cohomology-lemma-finite-pushforward-coefficients}，只需考虑
$j'_!f^{-1}\mathcal{G}$。利用引理
\ref{etale-cohomology-lemma-pullback-filtered-modules} 给出的滤过、
$j'_!$ 的正合性以及引理 \ref{etale-cohomology-lemma-ses-statements-coefficients}，
归结到 $\mathcal{F} = j'_!\underline{M}$ 的情形，其中所用的是有限
$\Lambda$-模 $M$；这正是引理
\ref{etale-cohomology-lemma-somewhat-easier-coefficients}。
\end{proof}

\begin{theorem}
\label{etale-cohomology-theorem-vanishing-affine-curves-coefficients}
设 $\Lambda$ 是 Noether 环，$k$ 是代数闭域，
$X$ 是维数 $\leq 1$ 的分离有限型概形（在 $k$ 上），而
$\mathcal{F}$ 是可构造且为挠的 $\Lambda$-模层
（在 $X_\etale$ 上）。则
\begin{enumerate}
\item
\label{etale-cohomology-item-finite-prime-to-p-coefficients}
$H^q_\etale(X, \mathcal{F})$ 是有限 $\Lambda$-模，
若 $\mathcal{F}$ 是与 $\text{char}(k)$ 互素的挠层；
\item
\label{etale-cohomology-item-finite-proper-coefficients}
$H^q_\etale(X, \mathcal{F})$ 是有限 $\Lambda$-模，若 $X$ 真。
\end{enumerate}
\end{theorem}

\begin{proof}
不再另行说明。写作
$\mathcal{F} = \mathcal{F}_1 \oplus \ldots \oplus \mathcal{F}_r$，
其中 $\mathcal{F}_i$ 被 $\ell_i^{n_i}$ 零化，这里的
$\ell_i$ 是素数，而 $n_i > 0$ 是整数。由引理
\ref{etale-cohomology-lemma-ses-statements-coefficients}，只需对
$\mathcal{F}_i$ 证明定理。因此可以并且确实假设
$\ell^n$ 零化 $\mathcal{F}$，其中 $\ell$ 是素数而 $n > 0$ 是整数。

\medskip\noindent
由于 $\mathcal{F}$ 作为 $\Lambda$-模层是可构造的，存在稠密开集
$U \subset X$，使得 $\mathcal{F}|_U$ 是有限型局部常值
$\Lambda$-模层。由于 $\dim(X) \leq 1$，缩小 $U$ 后，可以假设
$U = U_1 \amalg \ldots \amalg U_n$ 是不可约概形的不交并
（只需去掉位于 $\geq 2$ 个分支之交中的 $U$ 的闭点）。
由引理 \ref{etale-cohomology-lemma-restrict-to-open-coefficients}，归结到
$\mathcal{F} = j_!\mathcal{G}$ 的情形，其中 $\mathcal{G}$
是有限型局部常值 $\Lambda$-模层，定义在 $U$ 上
（且被 $\ell^n$ 零化）。

\medskip\noindent
由于所选 $U = U_1 \amalg \ldots \amalg U_n$ 中每个 $U_i$
都不可约，有
$$
j_!\mathcal{G} =
j_{1!}(\mathcal{G}|_{U_1}) \oplus \ldots \oplus
j_{n!}(\mathcal{G}|_{U_n})
$$
，其中 $j_i : U_i \to X$ 是包含态射。
$j_{i!}(\mathcal{G}|_{U_i})$ 的情形由引理
\ref{etale-cohomology-lemma-vanishing-easier-coefficients} 处理。
\end{proof}








\section{真概形的第一上同调}
\label{etale-cohomology-section-finite-etale-over-proper}

\noindent
在基本群，第 \ref{pione-section-finite-etale-over-proper} 节中，
我们已经在某种意义下看到：取 $R^1f_*\underline{G}$ 与基变换交换，
若 $f : X \to Y$ 是真态射，而 $G$ 是有限群（不必交换）。
本节推出这些结果的一个有用推论。

\begin{lemma}
\label{etale-cohomology-lemma-proper-over-henselian-and-h1}
设 $A$ 是 Hensel 局部环。设 $X$ 是 $A$ 上的真概形，
其闭纤维为 $X_0$。设 $M$ 是有限阿贝尔群。则
$H^1_\etale(X, \underline{M}) = H^1_\etale(X_0, \underline{M})$。
\end{lemma}

\begin{proof}
由位点上的上同调，引理 \ref{sites-cohomology-lemma-torsors-h1}，
$H^1_\etale(X, \underline{M})$ 的一个元素对应于一个
$\underline{M}$-挠子 $\mathcal{F}$（在 $X_\etale$ 上）。
这种挠子显然是有限局部常值层。因此，由引理
\ref{etale-cohomology-lemma-characterize-finite-locally-constant}，
$\mathcal{F}$ 可由概形 $V$ 表示，而该概形在 $X$ 上有限
\'etale。反过来，若概形 $V$ 在 $X$ 上有限 \'etale，并带有
$M$-作用，使它成为 $M$-挠子（在 $X$ 上），则它给出一个上同调类。
$X_0$ 上的上同调类与 $X_0$ 上有限 \'etale 挠子之间也有同样的对应。
因此，引理是基本群，引理
\ref{pione-lemma-finite-etale-on-proper-over-henselian}
中范畴等价的推论。
\end{proof}

\noindent
下列技术引理是真基变换定理证明中的关键材料。
这一论证逐字适用于 $A$ 上任意特殊纤维维数 $\leq 1$ 的真概形；
但事实上，该结论是真基变换定理的推论，而在其证明中我们只需要
这里这个特殊版本。

\begin{lemma}
\label{etale-cohomology-lemma-efface-cohomology-on-fibre-by-finite-cover}
设 $A$ 是 Hensel 局部环。令 $X = \mathbf{P}^1_A$。
设 $X_0 \subset X$ 是闭纤维。设 $\ell$ 是素数。设
$\mathcal{I}$ 是 $\mathbf{Z}/\ell\mathbf{Z}$-模的内射层
（在 $X_\etale$ 上）。则
$H^q_\etale(X_0, \mathcal{I}|_{X_0}) = 0$ 对 $q > 0$ 成立。
\end{lemma}

\begin{proof}
注意 $X$ 是可由 $2$ 个仿射开集覆盖的分离概形。因此对 $q > 1$，
结论来自 Gabber 的真基变换定理的仿射版本，见引理
\ref{etale-cohomology-lemma-vanishing-restriction-injective}。故可假设 $q = 1$。令
$\xi \in H^1_\etale(X_0, \mathcal{I}|_{X_0})$。
目标：证明 $\xi$ 是 $0$。
由引理 \ref{etale-cohomology-lemma-torsion-colimit-constructible} 和
\ref{etale-cohomology-lemma-colimit}，可以找到映射 $\mathcal{F} \to \mathcal{I}$，
其中 $\mathcal{F}$ 是可构造的
$\mathbf{Z}/\ell\mathbf{Z}$-模层，而 $\xi$ 来自某个元素 $\zeta$
（属于 $H^1_\etale(X_0, \mathcal{F}|_{X_0})$）。
假设有层的单射 $\mathcal{F} \to \mathcal{F}'$，这些层是
$\mathbf{Z}/\ell\mathbf{Z}$-模层（在 $X_\etale$ 上）。
由于 $\mathcal{I}$ 内射，可以把给定映射
$\mathcal{F} \to \mathcal{I}$ 延拓为映射
$\mathcal{F}' \to \mathcal{I}$。在这种情况下，可以把
$\mathcal{F}$ 替换为 $\mathcal{F}'$，并把 $\zeta$ 替换为
$\zeta$ 在 $H^1_\etale(X_0, \mathcal{F}'|_{X_0})$ 中的像。
此外，若 $\mathcal{F} = \mathcal{F}_1 \oplus \mathcal{F}_2$ 是直和，
则可以把 $\mathcal{F}$ 替换为 $\mathcal{F}_i$，并把 $\zeta$
替换为 $\zeta$ 在 $H^1_\etale(X_0, \mathcal{F}_i|_{X_0})$ 中的像。

\medskip\noindent
由引理 \ref{etale-cohomology-lemma-constructible-maps-into-constant-general}
及以上说明，可以假设 $\mathcal{F}$ 形如 $f_*\underline{M}$，
其中 $M$ 是有限 $\mathbf{Z}/\ell\mathbf{Z}$-模，而
$f : Y \to X$ 是有限表现的有限态射
（由引理 \ref{etale-cohomology-lemma-finite-pushforward-constructible}，
这种层仍然可构造，但我们不会用到这一点）。
由于 $f_*$ 的形成与任意基变换交换
（引理 \ref{etale-cohomology-lemma-finite-pushforward-commutes-with-base-change}），
$f_*\underline{M}$ 在 $X_0$ 上的限制等于
$\underline{M}$ 沿特殊纤维的诱导态射 $Y_0 \to X_0$ 的推前。
由 Leray 谱序列（命题 \ref{etale-cohomology-proposition-leray}）和高阶直像的消失
（命题 \ref{etale-cohomology-proposition-finite-higher-direct-image-zero}），得到
$$
H^1_\etale(X_0, f_*\underline{M}|_{X_0}) = H^1_\etale(Y_0, \underline{M}).
$$
由于 $Y \to \Spec(A)$ 真，可用引理
\ref{etale-cohomology-lemma-proper-over-henselian-and-h1} 看出
$H^1_\etale(Y_0, \underline{M})$ 等于
$H^1_\etale(Y, \underline{M})$。因此，上同调类 $\zeta$
可提升为上同调类
$$
\tilde\zeta \in H^1_\etale(Y, \underline{M}) = H^1_\etale(X, f_*\underline{M})
$$
然而，$\tilde \zeta$ 在 $H^1_\etale(X, \mathcal{I})$ 中的像为零，
因为 $\mathcal{I}$ 内射；再由下图的交换性
$$
\xymatrix{
H^1_\etale(X, f_*\underline{M}) \ar[r] \ar[d] &
H^1_\etale(X, \mathcal{I}) \ar[d] \\
H^1_\etale(X_0, (f_*\underline{M})|_{X_0}) \ar[r] &
H^1_\etale(X_0, \mathcal{I}|_{X_0})
}
$$
可知，$\xi$（$\zeta$ 的像）也为零。
\end{proof}











\section{基变换的预备知识}
\label{etale-cohomology-section-base-change-preliminaries}

\noindent
若你只关心光滑基变换定理或真基变换定理，可以直接跳到相应各节。
本节及随后几节考虑概形的交换图
$$
\xymatrix{
X \ar[d]_f & Y \ar[l]^h \ar[d]^e \\
S & T \ar[l]_g
}
$$
；通常假设此图笛卡尔，即 $Y = X \times_S T$。
上述交换图诱导小 \'etale 位点的交换图
$$
\xymatrix{
X_\etale \ar[d]_{f_{small}} & Y_\etale \ar[d]^{e_{small}} \ar[l]^{h_{small}} \\
S_\etale & T_\etale \ar[l]_{g_{small}}
}
$$
。采用记号
$$
f^{-1} = f_{small}^{-1}, \quad
g_* = g_{small, *}, \quad
e^{-1} = e_{small}^{-1}, \text{ 且}\quad
h_* = h_{small, *}.
$$
由位点，第 \ref{sites-section-pullback} 节，得到基变换（或拉回）映射
$$
f^{-1}g_*\mathcal{F}
\longrightarrow
h_*e^{-1}\mathcal{F}
$$
，其中 $\mathcal{F}$ 是 $T_\etale$ 上的层。若 $\mathcal{F}$
是 $T_\etale$ 上的阿贝尔层，则得到导出基变换映射
$$
f^{-1}Rg_*\mathcal{F}
\longrightarrow
Rh_*e^{-1}\mathcal{F}
$$
，见位点上的上同调，引理
\ref{sites-cohomology-lemma-base-change-map-flat-case}。
最后，若 $K$ 是 $D(T_\etale)$ 的任意对象，则有基变换映射
$$
f^{-1}Rg_*K
\longrightarrow
Rh_*e^{-1}K
$$
，见位点上的上同调，评注
\ref{sites-cohomology-remark-base-change}。

\begin{lemma}
\label{etale-cohomology-lemma-base-change-local}
考虑概形的笛卡尔图
$$
\xymatrix{
X \ar[d]_f & Y \ar[l]^h \ar[d]^e \\
S & T \ar[l]_g
}
$$
设 $\{U_i \to X\}$ 是 \'etale 覆盖，使得 $U_i \to S$
可分解为 $U_i \to V_i \to S$，其中 $V_i \to S$ 为 \'etale，
并考虑笛卡尔图
$$
\xymatrix{
U_i \ar[d]_{f_i} & U_i \times_X Y \ar[l]^{h_i} \ar[d]^{e_i} \\
V_i & V_i \times_S T \ar[l]_{g_i}
}
$$
设 $\mathcal{F}$ 是 $T_\etale$ 上的层。设 $K$ 属于 $D(T_\etale)$。
令 $K_i = K|_{V_i \times_S T}$，并令
$\mathcal{F}_i = \mathcal{F}|_{V_i \times_S T}$。
\begin{enumerate}
\item 若 $f_i^{-1}g_{i, *}\mathcal{F}_i = h_{i, *}e_i^{-1}\mathcal{F}_i$
对所有 $i$ 成立，则 $f^{-1}g_*\mathcal{F} = h_*e^{-1}\mathcal{F}$。
\item 若 $f_i^{-1}Rg_{i, *}K_i = Rh_{i, *}e_i^{-1}K_i$
对所有 $i$ 成立，则 $f^{-1}Rg_*K = Rh_*e^{-1}K$。
\item 若 $\mathcal{F}$ 是阿贝尔层，并且
$f_i^{-1}R^qg_{i, *}\mathcal{F}_i = R^qh_{i, *}e_i^{-1}\mathcal{F}_i$
对所有 $i$ 成立，则
$f^{-1}R^qg_*\mathcal{F} = R^qh_*e^{-1}\mathcal{F}$。
\end{enumerate}
\end{lemma}

\begin{proof}
证明 (1)。首先注意
$$
(f^{-1}g_*\mathcal{F})|_{U_i} =
f_i^{-1}(g_*\mathcal{F}|_{V_i}) =
f_i^{-1}g_{i, *}\mathcal{F}_i
$$
第一个等式成立是因为 $U_i \to X \to S$ 等于
$U_i \to V_i \to S$；第二个等式成立是因为由位点，引理
\ref{sites-lemma-localize-morphism-strong}，有
$g_*\mathcal{F}|_{V_i} = g_{i, *}\mathcal{F}_i$。类似地有
$$
(h_*e^{-1}\mathcal{F})|_{U_i} =
h_{i, *}(e^{-1}\mathcal{F}|_{U_i \times_X Y}) =
h_{i, *}e_i^{-1}\mathcal{F}_i
$$
因此，若基变换映射
$f_i^{-1}g_{i, *}\mathcal{F}_i \to h_{i, *}e_i^{-1}\mathcal{F}_i$
对所有 $i$ 都是同构，则基变换映射
$f^{-1}g_*\mathcal{F} \to h_*e^{-1}\mathcal{F}$
在 $U_i$ 上的限制对所有 $i$ 都是同构。由于 $\{U_i \to X\}$
是 \'etale 覆盖，可知它本身是同构。

\medskip\noindent
对其余两个断言，只需把对位点，引理
\ref{sites-lemma-localize-morphism-strong} 的引用替换为对
位点上的上同调，引理
\ref{sites-cohomology-lemma-restrict-direct-image-open} 的引用。
\end{proof}

\begin{lemma}
\label{etale-cohomology-lemma-base-change-compose}
考虑概形笛卡尔图组成的塔
$$
\xymatrix{
W \ar[d]_i & Z \ar[l]^j \ar[d]^k \\
X \ar[d]_f & Y \ar[l]^h \ar[d]^e \\
S & T \ar[l]_g
}
$$
设 $K$ 属于 $D(T_\etale)$。若
$$
f^{-1}Rg_*K \to Rh_*e^{-1}K
\quad\text{且}\quad
i^{-1}Rh_*e^{-1}K \to Rj_*k^{-1}e^{-1}K
$$
是同构，则
$(f \circ i)^{-1}Rg_*K \to Rj_*(e \circ k)^{-1}K$
是同构。类似地，若 $\mathcal{F}$ 是 $T_\etale$ 上的阿贝尔层，并且
$$
f^{-1}R^qg_*\mathcal{F} \to R^qh_*e^{-1}\mathcal{F}
\quad\text{且}\quad
i^{-1}R^qh_*e^{-1}\mathcal{F} \to R^qj_*k^{-1}e^{-1}\mathcal{F}
$$
是同构，则
$(f \circ i)^{-1}R^qg_*\mathcal{F} \to R^qj_*(e \circ k)^{-1}\mathcal{F}$
是同构。
\end{lemma}

\begin{proof}
这是形式的，只要验证这些基变换映射的复合就是外矩形的基变换映射；
见位点上的上同调，评注
\ref{sites-cohomology-remark-compose-base-change-horizontal}。
\end{proof}

\begin{lemma}
\label{etale-cohomology-lemma-base-change-Rf-star-colim}
设 $I$ 是有向集。考虑概形笛卡尔图的逆系统
$$
\xymatrix{
X_i \ar[d]_{f_i} & Y_i \ar[l]^{h_i} \ar[d]^{e_i} \\
S_i & T_i \ar[l]_{g_i}
}
$$
，其转移态射仿射，且 $g_i$ 拟紧拟分离。令 $X = \lim X_i$、
$S = \lim S_i$、$T = \lim T_i$ 和 $Y = \lim Y_i$，得到笛卡尔图
$$
\xymatrix{
X \ar[d]_f & Y \ar[l]^h \ar[d]^e \\
S & T \ar[l]_g
}
$$
设 $(\mathcal{F}_i, \varphi_{i'i})$ 是定义
\ref{etale-cohomology-definition-inverse-system-sheaves} 中那样的层系统（在
$(T_i)$ 上）。令 $\mathcal{F} = \colim p_i^{-1}\mathcal{F}_i$
（在 $T$ 上），其中 $p_i : T \to T_i$ 是投影。则有
\begin{enumerate}
\item 若 $f_i^{-1}g_{i, *}\mathcal{F}_i = h_{i, *}e_i^{-1}\mathcal{F}_i$
对所有 $i$ 成立，则
$f^{-1}g_*\mathcal{F} = h_*e^{-1}\mathcal{F}$。
\item 若 $\mathcal{F}_i$ 对所有 $i$ 都是阿贝尔层，并且
$f_i^{-1}R^qg_{i, *}\mathcal{F}_i = R^qh_{i, *}e_i^{-1}\mathcal{F}_i$
对所有 $i$ 成立，则
$f^{-1}R^qg_*\mathcal{F} = R^qh_*e^{-1}\mathcal{F}$。
\end{enumerate}
\end{lemma}

\begin{proof}
证明 (2)，略去 (1) 的证明。以下不再另行说明地使用：
层的拉回作为左伴随与余极限交换。注意 $h_i$ 作为 $g_i$ 的基变换
是拟紧拟分离的。以 $q_i : Y \to Y_i$ 表示投影，注意
$e^{-1}\mathcal{F} = \colim e^{-1}p_i^{-1}\mathcal{F}_i =
\colim q_i^{-1}e_i^{-1}\mathcal{F}_i$.
由引理 \ref{etale-cohomology-lemma-relative-colimit-general}，这给出
$$
R^qh_*e^{-1}\mathcal{F} = \colim r_i^{-1}R^qh_{i, *}e_i^{-1}\mathcal{F}_i
$$
，其中 $r_i : X \to X_i$ 是投影。类似地，有
$$
f^{-1}Rg_*\mathcal{F} =
f^{-1}\colim s_i^{-1}R^qg_{i, *}\mathcal{F}_i =
\colim r_i^{-1}f_i^{-1}R^qg_{i, *}\mathcal{F}_i
$$
，其中 $s_i : S \to S_i$ 是投影。引理得证。
\end{proof}

\begin{lemma}
\label{etale-cohomology-lemma-base-change-Rf-star-colim-complexes}
设 $I$、$X_i$、$Y_i$、$S_i$、$T_i$、$f_i$、$h_i$、$e_i$、$g_i$、
$X$、$Y$、$S$、$T$、$f$、$h$、$e$、$g$ 如引理
\ref{etale-cohomology-lemma-base-change-Rf-star-colim} 的陈述中所述。
设 $0 \in I$，并设 $K_0 \in D^+(T_{0, \etale})$。
对 $i \in I$、$i \geq 0$，以 $K_i$ 表示 $K_0$ 到 $T_i$ 的拉回。
以 $K$ 表示 $K_0$ 到 $T$ 的拉回。若
$f_i^{-1}Rg_{i, *}K_i = Rh_{i, *}e_i^{-1}K_i$
对所有 $i \geq 0$ 成立，则 $f^{-1}Rg_*K = Rh_*e^{-1}K$。
\end{lemma}

\begin{proof}
只需证明基变换映射
$f^{-1}Rg_*K \to Rh_*e^{-1}K$ 在上同调层上诱导同构。
换言之，必须证明 $f^{-1}R^pg_*K \to R^ph_*e^{-1}K$
对所有 $p \in \mathbf{Z}$ 都是同构，前提是已知
$f_i^{-1}R^pg_{i, *}K_i \to R^ph_{i, *}e_i^{-1}K_i$
对所有 $i \geq 0$ 和 $p \in \mathbf{Z}$ 都是同构。
此时可以完全照搬引理
\ref{etale-cohomology-lemma-base-change-Rf-star-colim} 的证明，只需把对引理
\ref{etale-cohomology-lemma-relative-colimit-general} 的引用替换为对引理
\ref{etale-cohomology-lemma-relative-colimit-general-complexes} 的引用。
\end{proof}

\begin{lemma}
\label{etale-cohomology-lemma-base-change-f-star-general-stalks}
考虑概形的笛卡尔图
$$
\xymatrix{
X \ar[d]_f & Y \ar[l]^h \ar[d]^e \\
S & T \ar[l]_g
}
$$
，其中 $g : T \to S$ 拟紧拟分离。设 $\mathcal{F}$
是 $T_\etale$ 上的阿贝尔层。设 $q \geq 0$。下列条件等价：
\begin{enumerate}
\item 对每个几何点 $\overline{x}$（属于 $X$），若其像为
$\overline{s} = f(\overline{x})$，则
$$
H^q(\Spec(\mathcal{O}^{sh}_{X, \overline{x}}) \times_S T, \mathcal{F})
=
H^q(\Spec(\mathcal{O}^{sh}_{S, \overline{s}}) \times_S T, \mathcal{F})
$$
\item $f^{-1}R^qg_*\mathcal{F} \to R^qh_*e^{-1}\mathcal{F}$
是同构。
\end{enumerate}
\end{lemma}

\begin{proof}
由于 $Y = X \times_S T$，有
$\Spec(\mathcal{O}^{sh}_{X, \overline{x}}) \times_X Y =
\Spec(\mathcal{O}^{sh}_{X, \overline{x}}) \times_S T$。
因此，由定理 \ref{etale-cohomology-theorem-higher-direct-images}
（以及引理 \ref{etale-cohomology-lemma-stalk-pullback}），(1) 中的映射就是
(2) 中映射在 $\overline{x}$ 处的茎映射。于是结论由定理
\ref{etale-cohomology-theorem-exactness-stalks} 得到。
\end{proof}

\begin{lemma}
\label{etale-cohomology-lemma-check-stalks-better}
设 $f : X \to S$ 是概形态射。设 $\overline{x}$ 是 $X$ 的几何点，
其像为 $\overline{s}$（在 $S$ 中）。设
$\Spec(K) \to \Spec(\mathcal{O}^{sh}_{S, \overline{s}})$ 是态射，
其中 $K$ 是可分闭域。设 $\mathcal{F}$ 是
$\Spec(K)_\etale$ 上的阿贝尔层。设 $q \geq 0$。下列条件等价：
\begin{enumerate}
\item
$H^q(\Spec(\mathcal{O}^{sh}_{X, \overline{x}}) \times_S \Spec(K), \mathcal{F}) =
H^q(\Spec(\mathcal{O}^{sh}_{S, \overline{s}}) \times_S \Spec(K), \mathcal{F})$
\item
$H^q(\Spec(\mathcal{O}^{sh}_{X, \overline{x}})
\times_{\Spec(\mathcal{O}^{sh}_{S, \overline{s}})} \Spec(K), \mathcal{F}) =
H^q(\Spec(K), \mathcal{F})$
\end{enumerate}
\end{lemma}

\begin{proof}
注意 $\Spec(K) \times_S \Spec(\mathcal{O}^{sh}_{S, \overline{s}})$
是 $K$ 上 \'etale 代数的滤过余极限之谱。由于 $K$ 可分闭，
每个 \'etale $K$-代数都是若干个 $K$ 的有限直积。因此可写成
$$
\Spec(K) \times_S \Spec(\mathcal{O}^{sh}_{S, \overline{s}}) =
\lim_{i \in I} \coprod\nolimits_{a \in A_i} \Spec(K)
$$
这个余滤过极限的每一项都是 $\Spec(K)$ 之不交并，
由有限集 $A_i$ 标号。注意 $A_i$ 非空，因为已给定
$\Spec(K) \to \Spec(\mathcal{O}^{sh}_{S, \overline{s}})$。于是
\begin{align*}
\Spec(\mathcal{O}^{sh}_{X, \overline{x}}) \times_S \Spec(K)
& =
\Spec(\mathcal{O}^{sh}_{X, \overline{x}})
\times_{\Spec(\mathcal{O}^{sh}_{S, \overline{s}})}
\left(
\Spec(\mathcal{O}^{sh}_{S, \overline{s}})
\times_S \Spec(K)\right) \\
& =
\lim_{i \in I} \coprod\nolimits_{a \in A_i}
\Spec(\mathcal{O}^{sh}_{X, \overline{x}})
\times_{\Spec(\mathcal{O}^{sh}_{S, \overline{s}})} \Spec(K)
\end{align*}
由于在当前情形取上同调与概形的极限交换
（定理 \ref{etale-cohomology-theorem-colimit}），结论随即成立。
\end{proof}













\section{正像的基变换}
\label{etale-cohomology-section-base-change-f-star}

\noindent
本节属于预备内容，初读时可略过。
本节讨论哪些态射 $f : X \to S$ 具有如下性质：对所有（集合）层都有
$f^{-1}g_* = h_*e^{-1}$，而且这对每个笛卡尔图
$$
\xymatrix{
X \ar[d]_f & Y \ar[l]^h \ar[d]^e \\
S & T \ar[l]_g
}
$$
都成立，只要 $g$ 拟紧拟分离。

\begin{lemma}
\label{etale-cohomology-lemma-base-change-f-star-general}
考虑概形的笛卡尔图
$$
\xymatrix{
X \ar[d]_f & Y \ar[l]^h \ar[d]^e \\
S & T \ar[l]_g
}
$$
。假设 $f$ 平坦，并且每个对象 $U$（属于 $X_\etale$）都有覆盖
$\{U_i \to U\}$，使得 $U_i \to S$ 可分解为
$U_i \to V_i \to S$，其中 $V_i \to S$ 为 \'etale，
而 $U_i \to V_i$ 拟紧且有几何连通纤维。
则对任意集合层 $\mathcal{F}$（在 $T_\etale$ 上），有
$f^{-1}g_*\mathcal{F} = h_*e^{-1}\mathcal{F}$。
\end{lemma}

\begin{proof}
设 $U \to X$ 是 \'etale 态射，且 $U \to S$ 可分解为
$U \to V \to S$，其中 $V \to S$ 为 \'etale，而
$U \to V$ 拟紧且有几何连通纤维。注意 $U \to V$ 平坦
（见平坦性的进一步结果，引理 \ref{flat-lemma-etale-flat-up-down}）。
我们断言
\begin{align*}
f^{-1}g_*\mathcal{F}(U)
& = g_*\mathcal{F}(V) \\
& = \mathcal{F}(V \times_S T) \\
& = e^{-1}\mathcal{F}(U \times_X Y) \\
& = h_*e^{-1}\mathcal{F}(U)
\end{align*}
。事实上，把 $U$ 看作 $X_\etale$ 的对象，把 $V$ 看作
$S_\etale$ 的对象，由引理 \ref{etale-cohomology-lemma-sections-upstairs}
得到第一个等式\footnote{严格地说，还用到了
$f^{-1}g_*\mathcal{F}$ 在 $U_\etale$ 上的限制，是经
$U \to V$ 拉回 $g_*\mathcal{F}$ 在 $V_\etale$ 上的限制所得；
见位点，引理 \ref{sites-lemma-localize-morphism-strong}。}。
把 $V \times_S T$ 看作 $T_\etale$ 的对象，由 $g_*$ 的定义
得到第二个等式。注意
$U \times_X Y = U \times_S T$（因为 $Y = X \times_S T$），
故 $U \times_X Y \to V \times_S T$ 是 $U \to V$ 的基变换，
从而有几何连通纤维。把 $U \times_X Y$ 看作 $Y_\etale$ 的对象，
与上述相同，由引理 \ref{etale-cohomology-lemma-sections-upstairs} 得到第三个等式。
最后，第四个等式来自 $h_*$ 的定义。

\medskip\noindent
由假设，$X_\etale$ 的每个对象都有一个 \'etale 覆盖，
上一段的论证适用于该覆盖，故引理成立。
\end{proof}

\begin{lemma}
\label{etale-cohomology-lemma-fppf-reduced-fibres-base-change-f-star}
考虑概形的笛卡尔图
$$
\xymatrix{
X \ar[d]_f & Y \ar[l]^h \ar[d]^e \\
S & T \ar[l]_g
}
$$
，其中 $f$ 平坦、局部有限表现，且有几何既约纤维。
则有 $f^{-1}g_*\mathcal{F} = h_*e^{-1}\mathcal{F}$，
其中 $\mathcal{F}$ 是 $T_\etale$ 上的任意层。
\end{lemma}

\begin{proof}
结合引理 \ref{etale-cohomology-lemma-base-change-f-star-general} 与
态射的进一步结果，引理
\ref{more-morphisms-lemma-cover-by-geometrically-connected} 即得。
\end{proof}

\begin{lemma}
\label{etale-cohomology-lemma-base-change-f-star-field}
考虑概形的笛卡尔图
$$
\xymatrix{
X \ar[d]_f & Y \ar[l]^h \ar[d]^e \\
S & T \ar[l]_g
}
$$
。假设 $S$ 是可分闭域之谱。
则有 $f^{-1}g_*\mathcal{F} = h_*e^{-1}\mathcal{F}$，
其中 $\mathcal{F}$ 是 $T_\etale$ 上的任意层。
\end{lemma}

\begin{proof}
可以在 $X$ 上局部论证，故可假设 $X$ 仿射。
于是可将 $X$ 写成 $S$ 上有限型仿射概形的余滤过极限。
由引理 \ref{etale-cohomology-lemma-base-change-Rf-star-colim}，可假设 $X$
在 $S$ 上有限型。此时引理 \ref{etale-cohomology-lemma-base-change-f-star-general}
适用，因为可分闭域上的任意有限型概形都是有限多个连通且
几何连通概形之不交并（见簇，引理
\ref{varieties-lemma-separably-closed-field-connected-components}）。
\end{proof}

\begin{lemma}
\label{etale-cohomology-lemma-base-change-f-star-valuation}
考虑概形的笛卡尔图
$$
\xymatrix{
X \ar[d]_f & Y \ar[l]^h \ar[d]^e \\
S & T \ar[l]_g
}
$$
。假设
\begin{enumerate}
\item $f$ 平坦且开；
\item $S$ 的剩余域都是可分代数闭的；
\item 给定 \'etale 态射 $U \to X$，其中 $U$ 为仿射概形，
可将 $U$ 写成 $X$ 的有限多个开子概形之不交并
（例如，当 $X$ 是函数域可分闭的正规整概形时）；
\item 任一纤维 $X_s$（属于 $f$）的每个非空开子集都连通
（例如，当 $X_s$ 不可约或为空时）。
\end{enumerate}
则对任意集合层 $\mathcal{F}$（在 $T_\etale$ 上），有
$f^{-1}g_*\mathcal{F} = h_*e^{-1}\mathcal{F}$。
\end{lemma}

\begin{proof}
略。提示：这些假设几乎直接蕴含引理
\ref{etale-cohomology-lemma-base-change-f-star-general} 的条件。
第 (3) 项括号内的例子来自引理
\ref{etale-cohomology-lemma-normal-scheme-with-alg-closed-function-field}。
\end{proof}

\noindent
下面的引理其实不属于这里，但似乎也没有更合适的位置。

\begin{lemma}
\label{etale-cohomology-lemma-fppf-reduced-fibres-pullback-products}
设 $f : X \to S$ 是平坦、局部有限表现且有几何既约纤维的概形态射。
则 $f^{-1} : \Sh(S_\etale) \to \Sh(X_\etale)$ 与积交换。
\end{lemma}

\begin{proof}
设 $I$ 是集合，并设 $\mathcal{G}_i$ 是 $S_\etale$ 上的层
（$i \in I$）。设 $U \to X$ 是 \'etale 态射，使得
$U \to S$ 可分解为 $U \to V \to S$，其中 $V \to S$ 为
\'etale，而 $U \to V$ 平坦、有限表现且有几何连通纤维。于是
\begin{align*}
f^{-1}(\prod \mathcal{G}_i)(U)
& =
(\prod \mathcal{G}_i)(V) \\
& =
\prod \mathcal{G}_i(V) \\
& =
\prod f^{-1}\mathcal{G}_i(U) \\
& =
(\prod f^{-1}\mathcal{G}_i)(U)
\end{align*}
，其中第一和第三个等式使用了引理
\ref{etale-cohomology-lemma-sections-upstairs}（还使用了如下事实：
$f^{-1}\mathcal{G}$ 在 $U_\etale$ 上的限制，是经
$U \to V$ 拉回 $\mathcal{G}$ 在 $V_\etale$ 上的限制所得；
见位点，引理 \ref{sites-lemma-localize-morphism-strong}）。
由态射的进一步结果，引理
\ref{more-morphisms-lemma-cover-by-geometrically-connected}，
每个对象 $U$（属于 $X_\etale$）都有 \'etale 覆盖
$\{U_i \to U\}$，使上一段的讨论适用于 $U_i$。引理得证。
\end{proof}

\begin{lemma}
\label{etale-cohomology-lemma-base-change-f-star}
设 $f : X \to S$ 是平坦的概形态射，并且对每个几何点
$\overline{x}$（属于 $X$），映射
$$
\mathcal{O}_{S, f(\overline{x})}^{sh}
\longrightarrow
\mathcal{O}_{X, \overline{x}}^{sh}
$$
都有几何连通纤维。则对每个概形的笛卡尔图
$$
\xymatrix{
X \ar[d]_f & Y \ar[l]^h \ar[d]^e \\
S & T \ar[l]_g
}
$$
，若 $g$ 拟紧拟分离，则有
$f^{-1}g_*\mathcal{F} = h_*e^{-1}\mathcal{F}$
，其中 $\mathcal{F}$ 是 $T_\etale$ 上的任意集合层。
\end{lemma}

\begin{proof}
只需在茎上检验等式；见定理 \ref{etale-cohomology-theorem-exactness-stalks}。
由定理 \ref{etale-cohomology-theorem-higher-direct-images}，有
$$
(h_*e^{-1}\mathcal{F})_{\overline{x}} =
\Gamma(\Spec(\mathcal{O}_{X, \overline{x}}^{sh}) \times_X Y, e^{-1}\mathcal{F})
$$
；类似地有
$$
(f^{-1}g_*^{-1}\mathcal{F})_{\overline{x}} =
(g_*^{-1}\mathcal{F})_{f(\overline{x})} =
\Gamma(\Spec(\mathcal{O}_{S, f(\overline{x})}^{sh}) \times_S T, \mathcal{F})
$$
。把引理 \ref{etale-cohomology-lemma-sections-upstairs} 应用于态射
$$
\Spec(\mathcal{O}_{X, \overline{x}}^{sh}) \times_X Y
\longrightarrow
\Spec(\mathcal{O}_{S, f(\overline{x})}^{sh}) \times_S T
$$
，可知这些集合相等；该态射是
$\Spec(\mathcal{O}_{X, \overline{x}}^{sh}) \to
\Spec(\mathcal{O}_{S, f(\overline{x})}^{sh})$
的基变换，因为 $Y = X \times_S T$。
\end{proof}







\section{高阶正像的基变换}
\label{etale-cohomology-section-base-change}

\noindent
本节是第 \ref{etale-cohomology-section-base-change-f-star} 节关于高阶正像的对应版本。
本节属于预备内容，初次阅读时应当略过。

\begin{remark}
\label{etale-cohomology-remark-base-change-holds}
设 $f : X \to S$ 是概形态射。设 $n$ 是整数。
我们称 $BC(f, n, q_0)$ 成立，如果对于每个交换图
$$
\xymatrix{
X \ar[d]_f & X' \ar[l] \ar[d]_{f'} & Y \ar[l]^h \ar[d]^e \\
S & S' \ar[l] & T \ar[l]_g
}
$$
其中 $X' = X \times_S S'$、$Y = X' \times_{S'} T$，且
$g$ 拟紧且拟分离，并且对于每个 Abel 层 $\mathcal{F}$，若它在
$T_\etale$ 上且被 $n$ 零化，则基变换映射
$$
(f')^{-1}R^qg_*\mathcal{F}
\longrightarrow
R^qh_*e^{-1}\mathcal{F}
$$
对 $q \leq q_0$ 都是同构。
\end{remark}

\begin{lemma}
\label{etale-cohomology-lemma-base-change-q-injective}
设 $f : X \to S$ 与 $n$ 如评注 \ref{etale-cohomology-remark-base-change-holds} 中所述，
并假设对某个 $q \geq 1$ 有 $BC(f, n, q - 1)$。则对于每个交换图
$$
\xymatrix{
X \ar[d]_f & X' \ar[l] \ar[d]_{f'} & Y \ar[l]^h \ar[d]^e \\
S & S' \ar[l] & T \ar[l]_g
}
$$
其中 $X' = X \times_S S'$、$Y = X' \times_{S'} T$，且
$g$ 拟紧且拟分离，并且对于每个 Abel 层 $\mathcal{F}$，若它在
$T_\etale$ 上且被 $n$ 零化，则有
\begin{enumerate}
\item 基变换映射
$(f')^{-1}R^qg_*\mathcal{F}\to R^qh_*e^{-1}\mathcal{F}$
是单射；
\item 如果 $\mathcal{F} \subset \mathcal{G}$，其中 $\mathcal{G}$
是 $T_\etale$ 上被 $n$ 零化的层，则
$$
\Coker\left(
(f')^{-1}R^qg_*\mathcal{F}\to R^qh_*e^{-1}\mathcal{F}
\right)
\subset
\Coker\left(
(f')^{-1}R^qg_*\mathcal{G}\to R^qh_*e^{-1}\mathcal{G}
\right)
$$
\item 如果 (2) 中的层 $\mathcal{G}$ 是 $\mathbf{Z}/n\mathbf{Z}$-模的内射层，则
$$
\Coker\left((f')^{-1}R^qg_*\mathcal{F}\to R^qh_*e^{-1}\mathcal{F} \right)
\subset R^qh_*e^{-1}\mathcal{G}
$$
\end{enumerate}
\end{lemma}

\begin{proof}
取短正合列
$0 \to \mathcal{F} \to \mathcal{I} \to \mathcal{Q} \to 0$
，其中 $\mathcal{I}$ 是 $\mathbf{Z}/n\mathbf{Z}$-模的内射层。
考虑由此诱导的图
$$
\xymatrix{
(f')^{-1}R^{q - 1}g_*\mathcal{I} \ar[d]_{\cong} \ar[r] &
(f')^{-1}R^{q - 1}g_*\mathcal{Q} \ar[d]_{\cong} \ar[r] &
(f')^{-1}R^qg_*\mathcal{F} \ar[d] \ar[r] &
0 \ar[d] \\
R^{q - 1}h_*e^{-1}\mathcal{I} \ar[r] &
R^{q - 1}h_*e^{-1}\mathcal{Q} \ar[r] &
R^qh_*e^{-1}\mathcal{F} \ar[r] &
R^qh_*e^{-1}\mathcal{I}
}
$$
，其两行正合。右上角为零，这是因为 $\mathcal{I}$ 是内射的。由
$BC(f, n, q - 1)$，左侧两个竖直箭头都是同构。因此 (1) 成立。
上图还表明
$$
\Coker\left(
(f')^{-1}R^qg_*\mathcal{F}\to R^qh_*e^{-1}\mathcal{F}
\right)
\subset
R^qh_*e^{-1}\mathcal{I}
$$
，故 (3) 成立。为证明 (2)，取
$\mathcal{F} \subset \mathcal{G} \subset \mathcal{I}$ 即可。
\end{proof}

\begin{lemma}
\label{etale-cohomology-lemma-base-change-q-integral-top}
设 $f : X \to S$ 与 $n$ 如评注 \ref{etale-cohomology-remark-base-change-holds} 中所述，
并假设对某个 $q \geq 1$ 有 $BC(f, n, q - 1)$。考虑交换图
$$
\vcenter{
\xymatrix{
X \ar[d]_f &
X' \ar[d]_{f'} \ar[l] &
Y \ar[l]^h \ar[d]^e &
Y' \ar[l]^{\pi'} \ar[d]^{e'} \\
S &
S' \ar[l] &
T \ar[l]_g &
T' \ar[l]_\pi
}
}
\quad\text{且}\quad
\vcenter{
\xymatrix{
X' \ar[d]_{f'} & & Y' \ar[ll]^{h' = h \circ \pi'} \ar[d]^{e'} \\
S' & & T' \ar[ll]_{g' = g \circ \pi}
}
}
$$
其中所有方块都是笛卡尔的，$g$ 拟紧且拟分离，而 $\pi$ 是整满射。
设 $\mathcal{F}$ 是 $T_\etale$ 上被 $n$ 零化的 Abel 层，并令
$\mathcal{F}' = \pi^{-1}\mathcal{F}$。如果基变换映射
$$
(f')^{-1}R^qg'_*\mathcal{F}' \longrightarrow R^qh'_*(e')^{-1}\mathcal{F}'
$$
是同构，则基变换映射
$(f')^{-1}R^qg_*\mathcal{F} \to R^qh_*e^{-1}\mathcal{F}$
也是同构。
\end{lemma}

\begin{proof}
注意到 $\mathcal{F} \to \pi_*\pi^{-1}\mathcal{F}'$ 是单射，
因为 $\pi$ 是满射（可在茎上检验）。因此，由引理
\ref{etale-cohomology-lemma-base-change-q-injective}，只需证明基变换映射
$$
(f')^{-1}R^qg_*\pi_*\mathcal{F}'
\longrightarrow
R^qh_*e^{-1}\pi_*\mathcal{F}'
$$
是同构。由假设，这一点成立，因为
$R^qg_*\pi_*\mathcal{F}' = R^qg'_*\mathcal{F}'$，
$e^{-1}\pi_*\mathcal{F}' =\pi'_*(e')^{-1}\mathcal{F}'$，并且
$R^qh_*\pi'_*(e')^{-1}\mathcal{F}' = R^qh'_*(e')^{-1}\mathcal{F}'$。
这些等式来自引理
\ref{etale-cohomology-lemma-integral-pushforward-commutes-with-base-change}、
\ref{etale-cohomology-lemma-what-integral} 以及相对 Leray 谱序列
（Cohomology on Sites，引理 \ref{sites-cohomology-lemma-relative-Leray}）。
\end{proof}

\begin{lemma}
\label{etale-cohomology-lemma-base-change-q-integral-bottom}
设 $f : X \to S$ 与 $n$ 如评注 \ref{etale-cohomology-remark-base-change-holds} 中所述，
并假设对某个 $q \geq 1$ 有 $BC(f, n, q - 1)$。考虑交换图
$$
\vcenter{
\xymatrix{
X \ar[d]_f &
X' \ar[d]_{f'} \ar[l] &
X'' \ar[l]^{\pi'} \ar[d]_{f''} &
Y \ar[l]^{h'} \ar[d]^e \\
S &
S' \ar[l] &
S'' \ar[l]_\pi &
T \ar[l]_{g'}
}
}
\quad\text{且}\quad
\vcenter{
\xymatrix{
X' \ar[d]_{f'} & & Y \ar[ll]^{h = h' \circ \pi'} \ar[d]^e \\
S' & & T \ar[ll]_{g = g' \circ \pi}
}
}
$$
其中所有方块都是笛卡尔的，$g'$ 拟紧且拟分离，而 $\pi$ 是整态射。
设 $\mathcal{F}$ 是 $T_\etale$ 上被 $n$ 零化的 Abel 层。如果基变换映射
$$
(f')^{-1}R^qg_*\mathcal{F} \longrightarrow R^qh_*e^{-1}\mathcal{F}
$$
是同构，则基变换映射
$(f'')^{-1}R^qg'_*\mathcal{F} \to R^qh'_*e^{-1}\mathcal{F}$
也是同构。
\end{lemma}

\begin{proof}
由于 $\pi$ 与 $\pi'$ 都是整态射，有 $R\pi_* = \pi_*$ 和
$R\pi'_* = \pi'_*$，见引理 \ref{etale-cohomology-lemma-what-integral}。
还有 $(f')^{-1}\pi_* = \pi'_*(f'')^{-1}$。因此
$\pi'_*(f'')^{-1}R^qg'_*\mathcal{F} = (f')^{-1}R^qg_*\mathcal{F}$
以及
$\pi'_*R^qh'_*e^{-1}\mathcal{F} = R^qh_*e^{-1}\mathcal{F}$.
所以，假设意味着对所考察的映射应用函子 $\pi'_*$ 后得到同构。
于是由引理 \ref{etale-cohomology-lemma-what-integral}，该映射本身是同构。
\end{proof}

\begin{lemma}
\label{etale-cohomology-lemma-formal-argument}
设 $T$ 是拟紧且拟分离的概形。设 $P$ 是 $T$ 上拟紧拟分离概形的一条性质。
假设
\begin{enumerate}
\item 如果 $T'' \to T'$ 是 $T$ 上拟紧拟分离概形之间的加厚，则
$P(T'')$ 成立当且仅当 $P(T')$ 成立。
\item 如果 $T' = \lim T_i$ 是 $T$ 上拟紧拟分离概形构成的逆系统之极限，
其转移态射为仿射态射，并且 $P(T_i)$ 对所有 $i$ 都成立，则 $P(T')$ 成立。
\item 如果 $Z \subset T'$ 是闭子概形，其补集 $V \subset T'$ 拟紧，
并且 $P(T')$ 成立，则 $P(V)$ 或 $P(Z)$ 至少有一个成立。
\end{enumerate}
那么，由 $P(T)$ 可推出：$P(\Spec(K))$ 成立，其中存在态射
$\Spec(K) \to T$，且 $K$ 是域。
\end{lemma}

\begin{proof}
考虑集合 $\mathfrak T$，其元素是闭子概形 $T' \subset T$，且满足 $P(T')$。
由假设 (2)，这个集合有极小元，记为 $T'$。由假设 (1)，$T'$ 是既约的。
设 $\eta \in T'$ 是 $T'$ 的某个不可约分支的泛点。则
$\eta = \Spec(K)$，其中 $K$ 是某个域，并且 $\eta = \lim V$，其中极限取遍
仿射开子概形 $V \subset T'$；这些开子概形都包含 $\eta$。由假设 (3) 以及 $T'$ 的极小性，
$P(V)$ 对所有这些 $V$ 都成立。因此由 (2)，$P(\eta)$ 成立，证明完毕。
\end{proof}

\begin{lemma}
\label{etale-cohomology-lemma-base-change-does-not-hold-pre}
设 $f : X \to S$ 与 $n$ 如评注 \ref{etale-cohomology-remark-base-change-holds} 中所述，
并假设对某个 $q \geq 1$，$BC(f, n, q - 1)$ 成立而
$BC(f, n, q)$ 不成立。那么存在交换图
$$
\xymatrix{
X \ar[d]_f & X' \ar[d]_{f'} \ar[l] & Y \ar[l]^h \ar[d]^e \\
S & S' \ar[l] & \Spec(K) \ar[l]_g
}
$$
，其中 $X' = X \times_S S'$、$Y = X' \times_{S'} \Spec(K)$，
$K$ 是域，并且存在 Abel 层 $\mathcal{F}$，它在 $\Spec(K)$ 上且被 $n$ 零化，使得
$(f')^{-1}R^qg_*\mathcal{F} \to R^qh_*e^{-1}\mathcal{F}$
不是同构。
\end{lemma}

\begin{proof}
取交换图
$$
\xymatrix{
X \ar[d]_f & X' \ar[l] \ar[d]_{f'} & Y \ar[l]^h \ar[d]^e \\
S & S' \ar[l] & T \ar[l]_g
}
$$
，其中 $X' = X \times_S S'$、$Y = X' \times_{S'} T$，
$g$ 拟紧且拟分离；再取 Abel 层 $\mathcal{F}$，它在 $T_\etale$ 上且被
$n$ 零化，使基变换映射
$(f')^{-1}R^qg_*\mathcal{F} \to R^qh_*e^{-1}\mathcal{F}$
不是同构。当然，可以并且确实用 $S'$ 的某个仿射开集替换 $S'$；
这使得 $T$ 拟紧且拟分离。由引理 \ref{etale-cohomology-lemma-base-change-q-injective}，映射
$(f')^{-1}R^qg_*\mathcal{F} \to R^qh_*e^{-1}\mathcal{F}$
是单射。取一个几何点 $\overline{x}$，它是 $X'$ 的几何点；再取元素
$\xi$，它属于 $(R^qh_*q^{-1}\mathcal{F})_{\overline{x}}$，
使它不属于下列映射的像：
$((f')^{-1}R^qg_*\mathcal{F})_{\overline{x}} \to
(R^qh_*e^{-1}\mathcal{F})_{\overline{x}}$.

\medskip\noindent
考虑态射 $\pi : T' \to T$，其中 $T'$ 拟紧且拟分离，并记
$\mathcal{F}' = \pi^{-1}\mathcal{F}$。记
$\pi' : Y' = Y \times_T T' \to Y$ 为 $\pi$ 的基变换，
$e' : Y' \to T'$ 为 $e$ 的基变换。图示为
$$
\vcenter{
\xymatrix{
X' \ar[d]_{f'} & Y \ar[l]^h \ar[d]^e & Y' \ar[l]^{\pi'} \ar[d]^{e'} \\
S' & T \ar[l]_g & T' \ar[l]_\pi
}
}
\quad\text{且}\quad
\vcenter{
\xymatrix{
X' \ar[d]_{f'} & & Y' \ar[ll]^{h' = h \circ \pi'} \ar[d]^{e'} \\
S' & & T' \ar[ll]_{g' = g \circ \pi}
}
}
$$
利用拉回映射，得到典范交换图
$$
\xymatrix{
(f')^{-1}R^qg_*\mathcal{F} \ar[r] \ar[d] &
(f')^{-1}R^qg'_*\mathcal{F}' \ar[d] \\
R^qh_*e^{-1}\mathcal{F} \ar[r] &
R^qh'_*(e')^{-1}\mathcal{F}'
}
$$
，这是 $X'$ 上 Abel 层的交换图。令 $P(T')$ 表示下列性质：
\begin{itemize}
\item 像 $\xi'$ 是 $\xi$ 在
$(Rh'_*(e')^{-1}\mathcal{F}')_{\overline{x}}$ 中的像，并且它
不属于下列映射的像：
$(f^{-1}R^qg'_*\mathcal{F}')_{\overline{x}} \to
(R^qh'_*(e')^{-1}\mathcal{F}')_{\overline{x}}$.
\end{itemize}
我们断言，性质 $P$ 满足引理 \ref{etale-cohomology-lemma-formal-argument} 的假设
(1)、(2)、(3)，从而即可证明本引理。

\medskip\noindent
性质 $P$ 满足引理 \ref{etale-cohomology-lemma-formal-argument} 的条件 (1)，
因为概形与该概形的加厚具有相同的 \'etale 拓扑。见命题
\ref{etale-cohomology-proposition-topological-invariance}。

\medskip\noindent
设 $I$ 是有向集，而 $T_i$ 是以 $I$ 为指标的逆系统，
其对象是 $T$ 上拟紧拟分离的概形，转移态射为仿射态射。
令 $T' = \lim T_i$。分别用 $\mathcal{F}'$ 与 $\mathcal{F}_i$
表示 $\mathcal{F}$ 到 $T'$ 与 $T_i$ 的拉回。考虑图
$$
\vcenter{
\xymatrix{
X \ar[d]_{f'} & Y \ar[l]^h \ar[d]^e & Y_i \ar[l]^{\pi_i'} \ar[d]^{e_i} \\
S & T \ar[l]_g & T_i \ar[l]_{\pi_i}
}
}
\quad\text{且}\quad
\vcenter{
\xymatrix{
X \ar[d]_{f'} & & Y_i \ar[ll]^{h_i = h \circ \pi_i'} \ar[d]^{e_i} \\
S & & T_i \ar[ll]_{g_i = g \circ \pi_i}
}
}
$$
，记号与上一段相同。显然，$\mathcal{F}'$（定义在 $T'$ 上）是各
$\mathcal{F}_i$ 到 $T'$ 的拉回之余极限，而 $(e')^{-1}\mathcal{F}'$
是各 $e_i^{-1}\mathcal{F}_i$ 到 $Y'$ 的拉回之余极限。
由引理 \ref{etale-cohomology-lemma-relative-colimit-general}，有
$$
R^qh'_*(e')^{-1}\mathcal{F}' = \colim R^qh_{i, *}e_i^{-1}\mathcal{F}_i
\quad\text{且}\quad
(f')^{-1}R^qg'_*\mathcal{F}' = \colim (f')^{-1}R^qg_{i, *}\mathcal{F}_i
$$
由此可见，如果 $P(T_i)$ 对所有 $i$ 都成立，则 $P(T')$ 成立。
所以性质 $P$ 满足引理 \ref{etale-cohomology-lemma-formal-argument} 的条件 (2)。

\medskip\noindent
最有意思的是引理 \ref{etale-cohomology-lemma-formal-argument} 的条件 (3)。
设 $T'$ 是 $T$ 上拟紧拟分离的概形，并且 $P(T')$ 成立。
设 $Z \subset T'$ 是闭子概形，其补集 $V \subset T'$ 拟紧。考虑图
$$
\xymatrix{
Y' \times_{T'} Z \ar[d]_{e_Z} \ar[r]_{i'} &
Y' \ar[d]_{e'} &
Y' \times_{T'} V \ar[l]^{j'} \ar[d]^{e_V} \\
Z \ar[r]^i &
T' &
V \ar[l]_j
}
$$
取单射 $j^{-1}\mathcal{F}' \to \mathcal{J}$，其中 $\mathcal{J}$
是 $\mathbf{Z}/n\mathbf{Z}$-模的内射层，定义在 $V$ 上。在茎上考察可知，映射
$$
\mathcal{F}' \to \mathcal{G} = j_*\mathcal{J} \oplus i_*i^{-1}\mathcal{F}'
$$
是单射。因此，$\xi'$ 映到下列群中的一个非零元素：
\begin{align*}
& \Coker\left(
((f')^{-1}R^qg'_*\mathcal{G})_{\overline{x}}
\to
(R^qh'_*(e')^{-1}\mathcal{G})_{\overline{x}}
\right) = \\
&
\Coker\left(
((f')^{-1}R^qg'_*j_*\mathcal{J})_{\overline{x}}
\to
(R^qh'_*(e')^{-1}j_*\mathcal{J})_{\overline{x}}
\right) \oplus \\
& \Coker\left(
((f')^{-1}R^qg'_*i_*i^{-1}\mathcal{F}')_{\overline{x}}
\to
(R^qh'_*(e')^{-1}i_*i^{-1}\mathcal{F}')_{\overline{x}}
\right)
\end{align*}
这里使用了引理 \ref{etale-cohomology-lemma-base-change-q-injective} 的 (2)。
如果 $\xi'$ 在第二个直和项中的像非零，则利用
$$
(f')^{-1}R^qg'_*i_*i^{-1}\mathcal{F}' =
(f')^{-1}R^q(g' \circ i)_*i^{-1}\mathcal{F}'
$$
（因为由命题 \ref{etale-cohomology-proposition-finite-higher-direct-image-zero}，
$Ri_* = i_*$）以及
$$
R^qh'_*(e')^{-1}i_*i^{-1}\mathcal{F} =
R^qh'_*i'_*e_Z^{-1}i^{-1}\mathcal{F} =
R^q(h' \circ i')_*e_Z^{-1}i^{-1}\mathcal{F}'
$$
（第一个等式由引理 \ref{etale-cohomology-lemma-finite-pushforward-commutes-with-base-change}，
第二个等式则来自命题 \ref{etale-cohomology-proposition-finite-higher-direct-image-zero}
所给出的 $Ri'_* = i'_*$），可知 $P(Z)$ 成立。
最后，假设 $\xi'$ 在第一个直和项中的像非零。有
$$
(e')^{-1}j_*\mathcal{J} = j'_*e_V^{-1}\mathcal{J}
\quad\text{且}\quad
R^aj'_*e_V^{-1}\mathcal{J} = 0, \quad a = 1, \ldots, q - 1
$$
；这是把 $BC(f, n, q - 1)$ 应用于图
$$
\xymatrix{
X \ar[d]_f & Y' \ar[l] \ar[d]_{e'} & Y \ar[l]^{j'} \ar[d]^{e_V} \\
S & T' \ar[l] & V \ar[l]_j
}
$$
并使用 $\mathcal{J}$ 的内射性得到的。
由 $h' \circ j'$ 的相对 Leray 谱序列
（Cohomology on Sites，引理 \ref{sites-cohomology-lemma-relative-Leray}），可知
$$
R^qh'_*(e')^{-1}j_*\mathcal{J} =
R^qh'_*j'_*e_V^{-1}\mathcal{J}
\longrightarrow
R^q(h' \circ j')_* e_V^{-1}\mathcal{J}
$$
是单射。因此，$\xi$ 映到
$(R^q(h' \circ j')_* e_V^{-1}\mathcal{J})_{\overline{x}}$
中的一个非零元素。把引理 \ref{etale-cohomology-lemma-base-change-q-injective} 的 (3)
应用于单射 $j^{-1}\mathcal{F}' \to \mathcal{J}$，即得 $P(V)$ 成立。
\end{proof}

\begin{lemma}
\label{etale-cohomology-lemma-base-change-does-not-hold}
设 $f : X \to S$ 与 $n$ 如评注 \ref{etale-cohomology-remark-base-change-holds} 中所述，
并假设对某个 $q \geq 1$，$BC(f, n, q - 1)$ 成立而
$BC(f, n, q)$ 不成立。那么存在交换图
$$
\xymatrix{
X \ar[d]_f & X' \ar[d] \ar[l] & Y \ar[l]^h \ar[d] \\
S & S' \ar[l] & \Spec(K) \ar[l]
}
$$
，其中两个方块都是笛卡尔的，并且
\begin{enumerate}
\item $S'$ 是仿射、整且正规的概形，其函数域代数闭；
\item $K$ 代数闭，并且 $\Spec(K) \to S'$ 是支配态射
（换言之，$K$ 是 $S'$ 的函数域的一个扩域）。
\end{enumerate}
此外，存在整数 $d | n$，使得 $R^qh_*(\mathbf{Z}/d\mathbf{Z})$ 非零。
\end{lemma}

\noindent
反过来，引理中 $R^qh_*(\mathbf{Z}/d\mathbf{Z})$ 的非消失蕴含
$BC(f, n, q)$ 不成立，因为引理
\ref{etale-cohomology-lemma-Rf-star-zero-normal-with-alg-closed-function-field} 表明
$R^q(\Spec(K) \to S')_*\mathbf{Z}/d\mathbf{Z} = 0$。

\begin{proof}
首先取引理 \ref{etale-cohomology-lemma-base-change-does-not-hold-pre} 中的图与
$\mathcal{F}$。可以并且确实假设 $S'$ 是仿射的
（这是显然的；若有疑问，可参见该引理的证明）。
由引理 \ref{etale-cohomology-lemma-base-change-q-integral-top}，可以假设 $K$ 代数闭。
此时 $\mathcal{F}$ 对应于一个 $\mathbf{Z}/n\mathbf{Z}$-模。
这样的模是若干 $\mathbf{Z}/d\mathbf{Z}$ 的直和，其中 $d | n$ 可变，
因此可以假设 $\mathcal{F}$ 是常值层，其值为 $\mathbf{Z}/d\mathbf{Z}$。
由引理 \ref{etale-cohomology-lemma-base-change-q-integral-bottom}，可以把 $S'$ 替换为
$S'$ 在 $\Spec(K)$ 中的正规化，证明完毕。
\end{proof}










\section{光滑基变换}
\label{etale-cohomology-section-smooth-base-change}

\noindent
本节证明光滑基变换定理。

\begin{lemma}
\label{etale-cohomology-lemma-smooth-base-change-fields}
设 $K/k$ 是域扩张。设 $X$ 是 $k$ 上的光滑仿射曲线，
并有有理点 $x \in X(k)$。设 $\mathcal{F}$ 是 $\Spec(K)$ 上的 Abel 层，
它被整数 $n$ 零化，而该整数在 $k$ 中可逆。设 $q > 0$，且
$$
\xi \in H^q(X_K, (X_K \to \Spec(K))^{-1}\mathcal{F})
$$
则存在
\begin{enumerate}
\item 有限扩张 $K'/K$ 与 $k'/k$，其中 $k' \subset K'$；
\item 有限 \'etale Galois 覆盖 $Z \to X_{k'}$，其 Galois 群为 $G$。
\end{enumerate}
使得 $G$ 的阶整除 $n$ 的某个幂，$Z \to X_{k'}$ 在 $x_{k'}$ 上分裂，
并且 $\xi$ 在 $H^q(Z_{K'}, (Z_{K'} \to \Spec(K))^{-1}\mathcal{F})$ 中变为零。
\end{lemma}

\begin{proof}
当 $q > 1$ 时，已知 $\xi$ 在
$H^q(X_{\overline{K}}, (X_{\overline{K}} \to \Spec(K))^{-1}\mathcal{F})$
中变为零（定理 \ref{etale-cohomology-theorem-vanishing-affine-curves}）。
由引理 \ref{etale-cohomology-lemma-directed-colimit-cohomology}，这意味着存在有限扩张
$K'/K$，使得 $\xi$ 在
$H^q(X_{K'}, (X_{K'} \to \Spec(K))^{-1}\mathcal{F})$ 中变为零。
因此在这种情形下可以取 $k' = k$ 以及 $Z = X$。

\medskip\noindent
假设 $q = 1$。回顾，$\mathcal{F}$ 对应于离散模 $M$，它带有连续
$\text{Gal}_K$-作用，见引理
\ref{etale-cohomology-lemma-equivalence-abelian-sheaves-point}。由于 $M$ 是 $n$-挠的，
它是有限 $\text{Gal}_K$-稳定子群的并。因此可归约到 $M$ 是被 $n$
零化的有限 Abel 群之情形，见引理 \ref{etale-cohomology-lemma-colimit}。以一个有限扩张替换
$K$ 后，可以假设 $\text{Gal}_K$ 在 $M$ 上的作用平凡。
所以可以假设 $\mathcal{F} = \underline{M}$ 是常值层，其值为有限 Abel 群
$M$，且该群被 $n$ 零化。

\medskip\noindent
可以把 $M$ 写成循环群的直和。任意两个有限 \'etale Galois 覆盖，
若其 Galois 群的阶在 $k$ 中可逆，就都可被第三个 Galois 群阶也在
$k$ 中可逆的覆盖控制（Fundamental Groups，第
\ref{pione-section-finite-etale-under-galois} 节）。因此只需在
$M = \mathbf{Z}/d\mathbf{Z}$ 且 $d | n$ 时证明本引理。

\medskip\noindent
假设 $M = \mathbf{Z}/d\mathbf{Z}$，其中 $d | n$。
在这种情形下，$\overline{\xi} = \xi|_{X_{\overline{K}}}$ 是下列群的元素：
$$
H^1(X_{\overline{k}}, \mathbf{Z}/d\mathbf{Z}) =
H^1(X_{\overline{K}}, \mathbf{Z}/d\mathbf{Z})
$$
见定理 \ref{etale-cohomology-theorem-vanishing-affine-curves}。这个群分类
$\mathbf{Z}/d\mathbf{Z}$-挠子，见 Cohomology on Sites，引理
\ref{sites-cohomology-lemma-torsors-h1}。与 $\overline{\xi}$ 对应的挠子
（看作 $X_{\overline{k}, \etale}$ 上的层）又给出有限 \'etale 态射
$T \to X_{\overline{k}}$，它带有 $\mathbf{Z}/d\mathbf{Z}$ 的作用，且该作用
在 $T$ 位于 $x_{\overline{k}}$ 上方的纤维上可迁，见引理
\ref{etale-cohomology-lemma-characterize-finite-locally-constant}。取一个连通分支
$T' \subset T$（如果 $\overline{\xi}$ 的阶为 $d$，则 $T$ 已经连通）。
于是 $T' \to X_{\overline{k}}$ 是有限 \'etale Galois 覆盖，其 Galois 群是子群
$G \subset \mathbf{Z}/d\mathbf{Z}$（略去一个小细节）。此外，元素
$\overline{\xi}$ 在映射 $H^1(X_{\overline{k}}, \mathbf{Z}/d\mathbf{Z}) \to
H^1(T', \mathbf{Z}/d\mathbf{Z})$ 下映为零，因为这是 $T$ 的定义性质之一。

\medskip\noindent
接下来使用极限论证，选取有限扩张 $k'/k$，它包含于 $\overline{k}$，使得
$T' \to X_{\overline{k}}$ 下降为有限 \'etale Galois 覆盖
$Z \to X_{k'}$，其 Galois 群为 $G$。见 Limits，引理
\ref{limits-lemma-descend-finite-presentation}、
\ref{limits-lemma-descend-finite-finite-presentation} 和
\ref{limits-lemma-descend-etale}。增大 $k'$ 后，可以假设 $Z$ 在 $x_{k'}$
上分裂。按照构造，$\xi$ 在
$H^1(Z_{\overline{K}}, \mathbf{Z}/d\mathbf{Z})$ 中的像为零。
因此由引理 \ref{etale-cohomology-lemma-directed-colimit-cohomology}，可找到有限子扩张
$\overline{K}/K'/K$，它包含 $k'$，使得 $\xi$ 在 $H^1(Z_{K'}, \mathbf{Z}/d\mathbf{Z})$
中变为零，证明完毕。
\end{proof}

\begin{theorem}[光滑基变换]
\label{etale-cohomology-theorem-smooth-base-change}
考虑概形的笛卡尔图
$$
\xymatrix{
X \ar[d]_f & Y \ar[l]^h \ar[d]^e \\
S & T \ar[l]_g
}
$$
，其中 $f$ 光滑，而 $g$ 拟紧且拟分离。那么
$$
f^{-1}R^qg_*\mathcal{F} = R^qh_*e^{-1}\mathcal{F}
$$
对任意 $q$ 以及任意 Abel 层 $\mathcal{F}$ 都成立；这里该层定义在
$T_\etale$ 上，并且它在几何点处的所有茎都是挠群，其阶在 $S$ 上可逆。
\end{theorem}

\begin{proof}[光滑基变换的第一种证明]
这个证明很长，但比下面给出的第二种证明更直接（使用的一般理论较少）。

\medskip\noindent
本定理在 $X_\etale$ 上是局部的。更准确地，假设有 \'etale 态射
$U \to X$，使得 $U \to S$ 分解为 $U \to V \to S$，其中
$V \to S$ 是 \'etale 的。于是可以考虑笛卡尔方块
$$
\xymatrix{
U \ar[d]_{f'} & U \times_X Y \ar[l]^{h'} \ar[d]^{e'} \\
V & V \times_S T \ar[l]_{g'}
}
$$
并令 $\mathcal{F}' = \mathcal{F}|_{V \times_S T}$。于是
$f^{-1}R^qg_*\mathcal{F}|_U = (f')^{-1}R^qg'_*\mathcal{F}'$，且
$R^qh_*e^{-1}\mathcal{F}|_U = R^qh'_*(e')^{-1}\mathcal{F}'$；
这是网站态射与局部化相容的结果，见 Sites，引理
\ref{sites-lemma-localize-morphism-strong} 以及 Cohomology on Sites，引理
\ref{sites-cohomology-lemma-restrict-direct-image-open}。因此，只需构造
$X$ 的一个 \'etale 覆盖，其成员为 $U \to X$，并构造如上的分解
$U \to V \to S$，使定理对带有 $f'$、$h'$、$g'$、$e'$ 的图成立。

\medskip\noindent
由光滑态射的局部结构，见 Morphisms，引理
\ref{morphisms-lemma-smooth-etale-over-affine-space}，可以假设 $X$ 与
$S$ 都是仿射的，并且 $X \to S$ 经过 \'etale 态射
$X \to \mathbf{A}^d_S$ 分解。如果有一叠笛卡尔图
$$
\xymatrix{
W \ar[d]_i & Z \ar[l]^j \ar[d]^k \\
X \ar[d]_f & Y \ar[l]^h \ar[d]^e \\
S & T \ar[l]_g
}
$$
，并且定理对下面和上面的方块成立，那么它也对外侧矩形成立；这是形式的。
把 $X \to S$ 写成复合
$$
X \to \mathbf{A}^{d - 1}_S \to \mathbf{A}^{d - 2}_S \to \ldots \to
\mathbf{A}^1_S \to S
$$
，可知只需在 $X$ 与 $S$ 都仿射且 $X \to S$ 的相对维数为 $1$ 时证明定理。

\medskip\noindent
对每个 $n \geq 1$，若它在 $S$ 上可逆，则令 $\mathcal{F}[n]$ 为
$\mathcal{F}$ 的被 $n$ 零化的截面所成的子层。由假设，
$\mathcal{F} = \colim \mathcal{F}[n]$，这是因为 $\mathcal{F}$ 的茎具有所述性质。函子 $e^{-1}$ 与 $f^{-1}$
作为左伴随与余极限交换。由引理 \ref{etale-cohomology-lemma-relative-colimit}，函子
$R^qh_*$ 与 $R^qg_*$ 和滤过余极限交换。因此只需对 $\mathcal{F}[n]$
证明定理。从现在起固定整数 $n$，使用 $\mathbf{Z}/n\mathbf{Z}$-模层，
并假设 $S$ 是 $\Spec(\mathbf{Z}[1/n])$ 上的概形。

\medskip\noindent
接下来归约到 $T$ 仿射的情形。由于 $g$ 拟紧且拟分离而 $S$ 仿射，
概形 $T$ 拟紧且拟分离。因此可以使用 Cohomology of Schemes，引理
\ref{coherent-lemma-induction-principle} 的归纳原理。故只需证明：如果
$T = W \cup W'$ 是开覆盖，并且定理对下列方块成立
$$
\xymatrix{
X \ar[d] & e^{-1}(W) \ar[l]^i \ar[d] \\
S & W \ar[l]_a
}
\quad
\xymatrix{
X \ar[d] & e^{-1}(W') \ar[l]^j \ar[d] \\
S & W' \ar[l]_b
}
\quad
\xymatrix{
X \ar[d] & e^{-1}(W \cap W') \ar[l]^-k \ar[d] \\
S & W \cap W' \ar[l]_c
}
$$
，则定理对原图成立。为此考虑图
$$
\xymatrix{
f^{-1}R^{q - 1}c_*\mathcal{F}|_{W \cap W'} \ar[d]^{\cong} \ar[r] &
f^{-1}R^qg_*\mathcal{F} \ar[d] \ar[r] &
f^{-1}R^qa_*\mathcal{F}|_W \oplus
f^{-1}R^qb_*\mathcal{F}|_{W'} \ar[d]_{\cong} \\
R^qk_*e^{-1}\mathcal{F}|_{e^{-1}(W \cap W')} \ar[r] &
R^qh_*e^{-1}\mathcal{F} \ar[r] &
R^qi_*e^{-1}\mathcal{F}|_{e^{-1}(W)} \oplus
R^qj_*e^{-1}\mathcal{F}|_{e^{-1}(W')}
}
$$
，其两行是引理 \ref{etale-cohomology-lemma-relative-mayer-vietoris} 的长正合列。
于是 $5$-引理给出所需结论。

\medskip\noindent
综上，可以假设 $S$、$X$、$T$ 与 $Y$ 都仿射，$\mathcal{F}$ 是
$n$-挠的，$X \to S$ 光滑且相对维数为 $1$，并且 $S$ 是
$\mathbf{Z}[1/n]$ 上的概形。对 $q$ 作归纳来证明定理。初始情形
$q = 0$ 由引理 \ref{etale-cohomology-lemma-fppf-reduced-fibres-base-change-f-star} 处理。
假设 $q > 0$，并且定理对所有更小次数都成立。取短正合列
$0 \to \mathcal{F} \to \mathcal{I} \to \mathcal{Q} \to 0$
，其中 $\mathcal{I}$ 是 $\mathbf{Z}/n\mathbf{Z}$-模的内射层。考虑诱导图
$$
\xymatrix{
f^{-1}R^{q - 1}g_*\mathcal{I} \ar[d]_{\cong} \ar[r] &
f^{-1}R^{q - 1}g_*\mathcal{Q} \ar[d]_{\cong} \ar[r] &
f^{-1}R^qg_*\mathcal{F} \ar[d] \ar[r] &
0 \ar[d] \\
R^{q - 1}h_*e^{-1}\mathcal{I} \ar[r] &
R^{q - 1}h_*e^{-1}\mathcal{Q} \ar[r] &
R^qh_*e^{-1}\mathcal{F} \ar[r] &
R^qh_*e^{-1}\mathcal{I}
}
$$
，其两行正合。右上角为零，因为 $\mathcal{I}$ 是内射的。
由归纳假设，左侧两个竖直箭头都是同构。因此只需证明
$R^qh_*e^{-1}\mathcal{I} = 0$。

\medskip\noindent
写成 $S = \Spec(A)$ 与 $T = \Spec(B)$，并设态射 $T \to S$
由环映射 $A \to B$ 给出。可以把
$A \to B = \colim_{i \in I} (A_i \to B_i)$ 写成
$\mathbf{Z}[1/n]$ 上有限型环映射的滤过余极限（见 Algebra，引理
\ref{algebra-lemma-limit-no-condition}）。对 $i \in I$，令
$S_i = \Spec(A_i)$ 和 $T_i = \Spec(B_i)$。当 $i$ 足够大时，可以找到
光滑态射 $X_i \to S_i$，其相对维数为 $1$，并使得
$X = X_i \times_{S_i} S$，见 Limits，引理
\ref{limits-lemma-descend-finite-presentation}、
\ref{limits-lemma-descend-smooth} 和
\ref{limits-lemma-descend-dimension-d}。令 $Y_i = X_i \times_{S_i} T_i$，得到方块
$$
\xymatrix{
X_i \ar[d]_{f_i} & Y_i \ar[l]^{h_i} \ar[d]^{e_i} \\
S_i & T_i \ar[l]_{g_i}
}
$$
注意，$\mathcal{I}_i = (T \to T_i)_*\mathcal{I}$ 是
$\mathbf{Z}/n\mathbf{Z}$-模的内射层，定义在 $T_i$ 上，见 Cohomology on Sites，引理
\ref{sites-cohomology-lemma-pushforward-injective-flat}。由引理
\ref{etale-cohomology-lemma-linus-hamann}，有
$\mathcal{I} = \colim (T \to T_i)^{-1}\mathcal{I}_i$。沿 $e$ 拉回得到
$e^{-1}\mathcal{I} = \colim (Y \to Y_i)^{-1}e_i^{-1}\mathcal{I}_i$。
把引理 \ref{etale-cohomology-lemma-relative-colimit-general} 应用于态射系统
$Y_i \to X_i$（其极限为 $Y \to X$），得到
$$
R^qh_*e^{-1}\mathcal{I} =
\colim (X \to X_i)^{-1} R^qh_{i, *} e_i^{-1}\mathcal{I}_i
$$
这把问题归约到 $T$ 与 $S$ 都是 $\mathbf{Z}[1/n]$ 上有限型仿射概形的情形。

\medskip\noindent
综上，有整数 $q \geq 1$，使定理在次数 $< q$ 时成立；概形 $S$ 与 $T$
都是 $\mathbf{Z}[1/n]$ 上有限型的仿射概形；$X \to S$ 光滑且相对维数为
$1$，其中 $X$ 仿射；$\mathcal{I}$ 是 $\mathbf{Z}/n\mathbf{Z}$-模的内射层；
需要证明 $R^qh_*e^{-1}\mathcal{I} = 0$。这将通过对 $\dim(T)$ 作归纳来完成。

\medskip\noindent
初始情形是 $T = \emptyset$，即 $\dim(T) < 0$。若不喜欢这种约定，
也可以把 $\dim(T) = 0$ 作为初始情形。此时 $T \to S$ 是有限的
（事实上 $T \to \Spec(\mathbf{Z}[1/n])$ 也是有限的，因为目标是 Jacobson 的；
细节略去），所以 $h$ 也有限，从而其高阶正像消失（见下文的参考文献）。

\medskip\noindent
假设 $\dim(T) = d \geq 0$，并且对所有 $T$ 维数更低的情形已知结论。
取 $U$，它在 $X$ 上仿射且 \'etale；再取截面 $\xi$，它是
$R^qh_*q^{-1}\mathcal{I}$ 在 $U$ 上的截面。需要证明 $\xi$ 为零。当然，
可以替换 $X$，所用概形为 $U$（并相应地替换 $Y$，所用概形为
$U \times_X Y$），并假设
$\xi \in H^0(X, R^qh_*e^{-1}\mathcal{I})$。此外，由于
$R^qh_*e^{-1}\mathcal{I}$ 是层，只需证明 $\xi$ 局部为零，这里的局部性取在 $X$ 上。
因此，可以用一个 \'etale 覆盖的成员替换 $X$。特别地，利用引理
\ref{etale-cohomology-lemma-higher-direct-images}，可以假设 $\xi$ 是某个元素
$\tilde \xi \in H^q(Y, e^{-1}\mathcal{I})$ 的像。用 $\tilde \xi$ 表述，
任务是证明 $\tilde \xi$ 在 $H^q(U_i \times_X Y, e^{-1}\mathcal{I})$
中变为零，其中 $\{U_i \to X\}$ 是某个 \'etale 覆盖。

\medskip\noindent
由 More on Morphisms，引理 \ref{more-morphisms-lemma-cover-smooth-by-special}，
可以假设 $X \to S$ 分解为 $X \to V \to S$，其中 $V \to S$ 是
\'etale 的，而 $X \to V$ 是仿射概形之间相对维数为 $1$ 的光滑态射，
它有截面且纤维几何连通。注意
$\dim(V \times_S T) \leq \dim(T) = d$，例如可用 More on Algebra，引理
\ref{more-algebra-lemma-dimension-etale-extension}。因此，可以替换 $S$，
所用概形为 $V$；并替换 $T$，所用概形为 $V \times_S T$
（与证明第一段中的讨论完全相同）。
所以可以假设 $X \to S$ 光滑且相对维数为 $1$，纤维几何连通，并有截面
$\sigma : S \to X$。

\medskip\noindent
设 $\pi : T' \to T$ 是有限满射。下面将用到
$\dim(T') \leq \dim(T) = d$，见 Algebra，引理
\ref{algebra-lemma-integral-dim-up}。选取单射 $\pi^{-1}\mathcal{I} \to \mathcal{I}'$，
其中目标是 $\mathbf{Z}/n\mathbf{Z}$-模的内射层。于是
$\mathcal{I} \to \pi_*\mathcal{I}'$ 是单射，因而有分裂（因为
$\mathcal{I}$ 是 $\mathbf{Z}/n\mathbf{Z}$-模的内射层）。记
$\pi' : Y' = Y \times_T T' \to Y$ 为 $\pi$ 的基变换，
$e' : Y' \to T'$ 为 $e$ 的基变换。图示为
$$
\xymatrix{
X \ar[d]_f & Y \ar[l]^h \ar[d]^e & Y' \ar[l]^{\pi'} \ar[d]^{e'} \\
S & T \ar[l]_g & T' \ar[l]_\pi
}
$$
由命题 \ref{etale-cohomology-proposition-finite-higher-direct-image-zero} 和引理
\ref{etale-cohomology-lemma-finite-pushforward-commutes-with-base-change}，有
$R\pi'_*(e')^{-1}\mathcal{I}' = e^{-1}\pi_*\mathcal{I}'$.
因此由 Leray 谱序列（Cohomology on Sites，引理
\ref{sites-cohomology-lemma-Leray}），有
$$
H^q(Y', (e')^{-1}\mathcal{I}') =
H^q(Y, e^{-1}\pi_*\mathcal{I}') \supset H^q(Y, e^{-1}\mathcal{I})
$$
，而且在沿任意 \'etale 态射 $U \to X$ 作基变换后这仍然成立。
因此可以替换 $T$，所用概形为 $T'$；替换 $\mathcal{I}$，所用层为
$\mathcal{I}'$；并用 $\tilde \xi$ 在
$H^q(Y', (e')^{-1}\mathcal{I}')$ 中的像替换它。

\medskip\noindent
假设有分解 $T \to S' \to S$，其中 $\pi : S' \to S$ 有限。
令 $X' = S' \times_S X$，可以考虑诱导图
$$
\xymatrix{
X \ar[d]_f & X' \ar[l]^{\pi'} \ar[d]^{f'} & Y \ar[l]^{h'} \ar[d]^e \\
S & S' \ar[l]_\pi & T \ar[l]_g
}
$$
由于 $\pi'$ 的高阶正像消失，由 Leray 谱序列可知
$R^qh_*e^{-1}\mathcal{I} = \pi'_*R^qh'_*e^{-1}\mathcal{I}$。因此
$H^0(X, R^qh_*e^{-1}\mathcal{I}) = H^0(X', R^qh'_*e^{-1}\mathcal{I})$。
所以 $\xi$ 为零当且仅当 $R^qh'_*e^{-1}\mathcal{I}$ 的相应截面为零\footnote{也可用另一种方式理解这一步。也就是说，需要证明存在 \'etale 覆盖
$\{U_i \to X\}$，使得 $\tilde \xi$ 在
$H^q(U_i \times_X Y, e^{-1}\mathcal{I})$ 中变为零。然而，如果能证明存在
\'etale 覆盖 $\{U'_j \to X'\}$，使得 $\tilde \xi$ 在
$H^q(U'_i \times_{X'} Y, e^{-1}\mathcal{I})$ 中变为零，那么由
$X' \to X$ 的性质 (B)（引理 \ref{etale-cohomology-lemma-finite-B}），存在 \'etale 覆盖
$\{U_i \to X\}$，使 $U_i \times_X X'$ 是若干 $X'$ 上概形的不交并，
其中每个概形都经过 $U'_j$ 分解（对某个 $j$）。于是如所需，$\tilde \xi$ 在
$H^q(U_i \times_X Y, e^{-1}\mathcal{I})$ 中变为零。}。
因此可以替换 $S$，所用概形为 $S'$；并替换 $X$，所用概形为 $X'$。注意
$\sigma : S \to X$ 基变换为 $\sigma' : S' \to X'$；所以替换后，
$X \to S$ 仍然有截面 $\sigma$，且纤维几何连通。

\medskip\noindent
我们将使用 $S$ 与 $T$ 都是 Nagata 概形这一事实，见 Algebra，命题
\ref{algebra-proposition-ubiquity-nagata}；这保证各种正规化都是有限的，
见 Morphisms，引理 \ref{morphisms-lemma-nagata-normalization-finite} 和
\ref{morphisms-lemma-nagata-normalization}。特别地，可以先用正规化替换
$T$，再替换 $S$，所用概形为 $S$ 在 $T$ 中的正规化。于是 $T \to S$ 是整正规概形
之间支配态射的不交并，见 Morphisms，引理
\ref{morphisms-lemma-normal-normalization}。显然可以逐个连通分支论证，
故可假设 $T \to S$ 是整正规概形之间的支配态射。

\medskip\noindent
设 $s \in S$ 与 $t \in T$ 为泛点。由引理
\ref{etale-cohomology-lemma-smooth-base-change-fields}，存在有限域扩张 $K/\kappa(t)$ 和
$k/\kappa(s)$，其中 $k$ 包含于 $K$；还存在有限 \'etale Galois 覆盖
$Z \to X_k$，其 Galois 群为 $G$，群阶整除 $n$ 的某个幂，且该覆盖在
$\sigma(\Spec(k))$ 上分裂，使得 $\tilde \xi$ 在
$H^q(Z_K, e^{-1}\mathcal{I}|_{Z_K})$ 中的像为零。令 $T' \to T$ 为
$T$ 在 $\Spec(K)$ 中的正规化，并令 $S' \to S$ 为 $S$ 在
$\Spec(k)$ 中的正规化。于是得到交换图
$$
\xymatrix{
S' \ar[d] & T' \ar[l] \ar[d] \\
S & T \ar[l]
}
$$
，其竖直箭头有限。由上述论证，可以并且确实替换 $S$ 与 $T$，
所用概形分别为 $S'$ 与 $T'$（并相应地替换 $X$，所用概形为
$X \times_S S'$；替换 $Y$，所用概形为 $Y \times_T T'$）。替换后，
得到有限 \'etale Galois 覆盖 $Z \to X_s$；它是 $X \to S$ 的泛纤维上的覆盖，
其 Galois 群 $G$ 的阶整除 $n$ 的某个幂，
该覆盖在 $\sigma(s)$ 上分裂，并且 $\tilde \xi$ 在
$H^q(Z_t, (Z_t \to Y)^{-1}e^{-1}\mathcal{I})$ 中映为零。这里
$Z_t = Z \times_S t = Z \times_s t = Z \times_{X_s} Y_t$。由于 $n$ 在
$S$ 上可逆，由 Fundamental Groups，引理 \ref{pione-lemma-extend-covering}，
可以找到有限 \'etale 态射 $U \to X$，其到 $X_s$ 的限制为 $Z$。

\medskip\noindent
此时替换 $X$，所用概形为 $U$；并替换 $Y$，所用概形为 $U \times_X Y$。替换后，
$X \to S$ 的纤维可能不再几何连通（仍有截面，但不再使用它）；
所得好处是，替换后 $\tilde \xi$ 在 $H^q(Y_t, e^{-1}\mathcal{I})$ 中映为零，
即 $\tilde \xi$ 在 $Y \to T$ 的泛纤维上的限制为零。

\medskip\noindent
回顾，$t$ 是 $T$ 的函数域之谱；也就是说，作为概形，$t$ 是 $T$ 的
非空仿射开子概形之极限。由引理 \ref{etale-cohomology-lemma-directed-colimit-cohomology}，
存在非空开子概形 $V \subset T$，使得 $\tilde \xi$ 在
$H^q(Y \times_T V, e^{-1}\mathcal{I}|_{Y \times_T V})$ 中映为零。

\medskip\noindent
记 $Z = T \setminus V$。考虑图
$$
\xymatrix{
Y \times_T Z \ar[d]_{e_Z} \ar[r]_{i'} &
Y \ar[d]_e &
Y \times_T V \ar[l]^{j'} \ar[d]^{e_V} \\
Z \ar[r]^i &
T &
V \ar[l]_j
}
$$
取单射 $i^{-1}\mathcal{I} \to \mathcal{I}'$，其目标为
$\mathbf{Z}/n\mathbf{Z}$-模的内射层，定义在 $Z$ 上。在茎上考察可知，映射
$$
\mathcal{I} \to j_*\mathcal{I}|_V \oplus i_*\mathcal{I}'
$$
是单射，并且由于 $\mathcal{I}$ 是 $\mathbf{Z}/n\mathbf{Z}$-模的内射层，
该单射分裂。因此只需证明 $\tilde \xi$ 在
$$
H^q(Y, e^{-1}j_*\mathcal{I}|_V) \oplus
H^q(Y, e^{-1}i_*\mathcal{I}')
$$
中映为零，至少在用一个 \'etale 覆盖的成员替换 $X$ 后如此。注意
$$
e^{-1}j_*\mathcal{I}|_V = j'_*e_V^{-1}\mathcal{I}|_V,\quad
e^{-1}i_*\mathcal{I}' = i'_*e_Z^{-1}\mathcal{I}'
$$
由关于 $q$ 的归纳假设，有
$$
R^aj'_*e_V^{-1}\mathcal{I}|_V = 0, \quad a = 1, \ldots, q - 1
$$
由 $j'$ 的 Leray 谱序列及上述消失，可知
$$
H^q(Y, j'_*(e_V^{-1}\mathcal{I}|_V))
\longrightarrow
H^q(Y \times_T V, e_V^{-1}\mathcal{I}_V) =
H^q(Y \times_T V, e^{-1}\mathcal{I}|_{Y \times_T V})
$$
是单射。因此，$\tilde \xi$ 在上述第一个直和项中的像为零，因为已知
$\tilde \xi$ 在 $H^q(Y \times_T V, e^{-1}\mathcal{I}|_{Y \times_T V})$
中消失。由于 $\dim(Z) < \dim(T) = d$，由归纳，$\tilde \xi$
在第二个直和项
$$
H^q(Y, e^{-1}i_*\mathcal{I}') =
H^q(Y, i'_*e_Z^{-1}\mathcal{I}') =
H^q(Y \times_T Z, e_Z^{-1}\mathcal{I}')
$$
中的像在用一个合适的 \'etale 覆盖的成员替换 $X$ 后消失。
这就完成了光滑基变换定理的证明。
\end{proof}

\begin{proof}[光滑基变换的第二种证明]
这个证明与较长的第一种证明相同；它较短，仅仅因为我们已把所用论证
分别整理成了若干引理。

\medskip\noindent
$q = 0$ 的情形是引理
\ref{etale-cohomology-lemma-fppf-reduced-fibres-base-change-f-star}。因此可以假设
$q > 0$，并且结论对所有更小次数都成立。

\medskip\noindent
对每个 $n \geq 1$，若它在 $S$ 上可逆，则令 $\mathcal{F}[n]$ 为
$\mathcal{F}$ 的被 $n$ 零化的截面所成的子层。由假设，
$\mathcal{F} = \colim \mathcal{F}[n]$，这是因为 $\mathcal{F}$ 的茎具有所述性质。函子 $e^{-1}$ 与 $f^{-1}$
作为左伴随与余极限交换。由引理 \ref{etale-cohomology-lemma-relative-colimit}，函子
$R^qh_*$ 与 $R^qg_*$ 和滤过余极限交换。因此只需对 $\mathcal{F}[n]$
证明定理。从现在起固定整数 $n$，并假设它在 $S$ 上可逆；使用
$\mathbf{Z}/n\mathbf{Z}$-模层。

\medskip\noindent
由引理 \ref{etale-cohomology-lemma-base-change-local}，问题在 $X$ 与 $S$ 上都是 \'etale
局部的。由光滑态射的局部结构，见 Morphisms，引理
\ref{morphisms-lemma-smooth-etale-over-affine-space}，可以假设 $X$ 与 $S$
都仿射，并且 $X \to S$ 经过 \'etale 态射
$X \to \mathbf{A}^d_S$ 分解。把 $X \to S$ 写成复合
$$
X \to \mathbf{A}^{d - 1}_S \to \mathbf{A}^{d - 2}_S \to \ldots \to
\mathbf{A}^1_S \to S
$$
，由引理 \ref{etale-cohomology-lemma-base-change-compose} 可知，只需在 $X$ 与 $S$
都仿射且 $X \to S$ 的相对维数为 $1$ 时证明定理。

\medskip\noindent
由引理 \ref{etale-cohomology-lemma-base-change-does-not-hold}，只需证明
$R^qh_*\mathbf{Z}/d\mathbf{Z} = 0$ 对 $d | n$ 成立，只要有笛卡尔图
$$
\xymatrix{
X \ar[d] & Y \ar[d] \ar[l]^h \\
S & \Spec(K) \ar[l]
}
$$
，其中 $X \to S$ 仿射、光滑且相对维数为 $1$，$S$ 是正规整环
$A$ 的谱，其分式域 $L$ 代数闭，而 $K/L$ 是代数闭域的扩张。

\medskip\noindent
回顾，$R^qh_*\mathbf{Z}/d\mathbf{Z}$ 是下列预层的相伴层：
$$
U \longmapsto
H^q(U \times_X Y, \mathbf{Z}/d\mathbf{Z}) =
H^q(U \times_S \Spec(K), \mathbf{Z}/d\mathbf{Z})
$$
；它定义在 $X_\etale$ 上（引理 \ref{etale-cohomology-lemma-higher-direct-images}）。
因此只需证明：给定 $U$ 以及
$\xi \in H^q(U \times_S \Spec(K), \mathbf{Z}/d\mathbf{Z})$，存在 \'etale 覆盖
$\{U_i \to U\}$，使得 $\xi$ 在
$H^q(U_i \times_S \Spec(K), \mathbf{Z}/d\mathbf{Z})$ 中变为零。

\medskip\noindent
当然，可以取 $U$ 仿射。此时 $U \times_S \Spec(K)$ 是 $K$ 上的
（光滑）仿射曲线，因而由定理 \ref{etale-cohomology-theorem-vanishing-affine-curves}，
当 $q > 1$ 时有消失。

\medskip\noindent
最后一种情形：$q = 1$。可以像 More on Morphisms，引理
\ref{more-morphisms-lemma-cover-smooth-by-special} 中那样，用一个 \'etale
覆盖的成员替换 $U$。于是 $U \to S$ 分解为 $U \to V \to S$，其中
$U \to V$ 的纤维几何连通，$U$、$V$ 仿射，$V \to S$ 是 \'etale 的，
并且有截面 $\sigma : V \to U$。由引理
\ref{etale-cohomology-lemma-normal-scheme-with-alg-closed-function-field}，$V$ 同构于 $S$
的（有限）多个（仿射）开子概形之不交并。显然，可以用其中一个替换
$S$，并替换 $X$，所用的是 $U$ 的相应分支。因此可以假设 $X \to S$ 的纤维
几何连通，有截面 $\sigma$，并且
$\xi \in H^1(Y, \mathbf{Z}/d\mathbf{Z})$。由于 $K$ 与 $L$ 都代数闭，有
$$
H^1(X_L, \mathbf{Z}/d\mathbf{Z}) =
H^1(Y, \mathbf{Z}/d\mathbf{Z})
$$
，见引理 \ref{etale-cohomology-lemma-base-change-dim-1-separably-closed}。因此存在有限
\'etale Galois 覆盖 $Z \to X_L$，其 Galois 群为
$G \subset \mathbf{Z}/d\mathbf{Z}$，并且该覆盖零化 $\xi$。
这可以从引理 \ref{etale-cohomology-lemma-smooth-base-change-fields} 的陈述或证明看出，
也可直接使用 $\xi$ 对应于一个 $\mathbf{Z}/d\mathbf{Z}$-挠子，
它定义在 $X_L$ 上。最后，由 Fundamental Groups，引理
\ref{pione-lemma-extend-covering-general}，得到（必然满的）有限 \'etale 态射
$X' \to X$，其到 $X_L$ 的限制是 $Z \to X_L$。由于 $\xi$ 在
$X'_K$ 中变为零，证明完毕。
\end{proof}

\noindent
下面这个光滑基变换定理的直接推论，才是实践中经常使用的形式。

\begin{lemma}
\label{etale-cohomology-lemma-smooth-base-change-general}
设 $S$ 是概形。设 $S' = \lim S_i$ 是概形 $S_i$ 构成的有向逆极限，
其中各概形在 $S$ 上光滑，转移态射仿射。设 $f : X \to S$
拟紧且拟分离，并作纤维方块
$$
\xymatrix{
X' \ar[d]_{f'} \ar[r]_{g'} & X \ar[d]^f \\
S' \ar[r]^g & S
}
$$
那么
$$
g^{-1}Rf_*E = R(f')_*(g')^{-1}E
$$
对任意 $E \in D^+(X_\etale)$ 都成立，只要其上同调层 $H^q(E)$ 的茎
都是挠群，且其阶在 $S$ 上可逆。
\end{lemma}

\begin{proof}
考虑谱序列
$$
E_2^{p, q} = R^pf_*H^q(E)
\quad\text{且}\quad
{E'}_2^{p, q} = R^pf'_*H^q((g')^{-1}E) = R^pf'_*(g')^{-1}H^q(E)
$$
，它们收敛到 $R^nf_*E$ 与 $R^nf'_*(g')^{-1}E$。
这些谱序列在 Derived Categories，引理
\ref{derived-lemma-two-ss-complex-functor} 中构造。把光滑基变换定理
（定理 \ref{etale-cohomology-theorem-smooth-base-change}）与引理
\ref{etale-cohomology-lemma-base-change-Rf-star-colim} 结合，可知
$$
g^{-1}R^pf_*H^q(E) = R^p(f')_*(g')^{-1}H^q(E)
$$
综合以上各点即得本引理。
\end{proof}












\section{光滑基变换的应用}
\label{etale-cohomology-section-applications-smooth-base-change}

\noindent
本节讨论光滑基变换定理的一些颇为直接的推论。

\begin{lemma}
\label{etale-cohomology-lemma-base-change-field-extension}
设 $L/K$ 是域扩张。设 $g : T \to S$ 是 $K$ 上概形之间的拟紧拟分离态射。
记 $g_L : T_L \to S_L$ 为 $g$ 到 $\Spec(L)$ 的基变。
设 $E \in D^+(T_\etale)$ 的上同调层的茎都是挠群，且其阶在 $K$ 中可逆。
设 $E_L$ 为 $E$ 到 $(T_L)_\etale$ 的拉回。则
$Rg_{L, *}E_L$ 是 $Rg_*E$ 到 $S_L$ 的拉回。
\end{lemma}

\begin{proof}
若 $L/K$ 可分，则 $L$ 是光滑 $K$-代数的滤过余极限，见
Algebra, Lemma \ref{algebra-lemma-colimit-syntomic}。
因此在这种情形下，本引理立即由
Lemma \ref{etale-cohomology-lemma-smooth-base-change-general} 得出。
一般情形下，令 $K'$ 和 $L'$ 分别为 $K$ 和 $L$ 的完美闭包
（Algebra, Definition \ref{algebra-definition-perfection}）。
态射 $\Spec(K') \to \Spec(K)$ 和 $\Spec(L') \to \Spec(L)$
都是泛同胚，因为 $K'/K$ 和 $L'/L$ 都是纯不可分扩张
（见 Algebra, Lemma \ref{algebra-lemma-p-ring-map}）。
由 \'etale 上同调的拓扑不变性，我们于是有
$(T_{K'})_\etale = T_\etale$、
$(S_{K'})_\etale = S_\etale$、
$(T_{L'})_\etale = (T_L)\etale$ 以及
$(S_{L'})_\etale = (S_L)_\etale$，见
Proposition \ref{etale-cohomology-proposition-topological-invariance}。
这将本引理归结为可分域扩张 $L'/K'$ 的情形
（由完美域的定义可知，见 Algebra, Definition \ref{algebra-definition-perfect}）。
\end{proof}

\begin{lemma}
\label{etale-cohomology-lemma-smooth-base-change-separably-closed}
设 $K/k$ 是可分闭域的扩张。设 $X$ 是 $k$ 上的拟紧拟分离概形。
设 $E \in D^+(X_\etale)$ 的上同调层的茎都是挠群，且其阶在 $k$ 中可逆。则
\begin{enumerate}
\item 映射 $H^q_\etale(X, E) \to H^q_\etale(X_K, E|_{X_K})$ 是同构；并且
\item $E \to R(X_K \to X)_*E|_{X_K}$ 是同构。
\end{enumerate}
\end{lemma}

\begin{proof}
证明 (1)。首先，令 $\overline{k}$ 和 $\overline{K}$
分别为 $k$ 和 $K$ 的代数闭包。态射
$\Spec(\overline{k}) \to \Spec(k)$ 和
$\Spec(\overline{K}) \to \Spec(K)$ 都是泛同胚，因为
$\overline{k}/k$ 和 $\overline{K}/K$ 都是纯不可分扩张
（见 Algebra, Lemma \ref{algebra-lemma-p-ring-map}）。
因此，由 \'etale 上同调的拓扑不变性，
$H^q_\etale(X, \mathcal{F}) =
H^q_\etale(X_{\overline{k}}, \mathcal{F}_{X_{\overline{k}}})$，见
Proposition \ref{etale-cohomology-proposition-topological-invariance}。
对 $X_K$ 和 $X_{\overline{K}}$ 也同样成立。
因此可以假设 $k$ 和 $K$ 都是代数闭域。
在这种情形下，$K$ 是光滑 $k$-代数的极限，见
Algebra, Lemma \ref{algebra-lemma-colimit-syntomic}。
由此，本引理是 Theorem \ref{etale-cohomology-theorem-smooth-base-change}
按 Lemma \ref{etale-cohomology-lemma-smooth-base-change-general} 中的形式所述的一个特例。

\medskip\noindent
证明 (2)。取 $U$ 为 $X_\etale$ 中任意拟紧拟分离对象。上述结论表明，
映射 $E \to R(X_K \to X)_*E|_{X_K}$ 的限制诱导上同调上的同构。
由于 $X_\etale$ 的每个对象都有一个由此类 $U$ 构成的 \'etale 覆盖，
这就证明了所需结论。
\end{proof}

\begin{lemma}
\label{etale-cohomology-lemma-base-change-does-not-hold-post}
设 $f : X \to S$ 和 $n$ 如 Remark \ref{etale-cohomology-remark-base-change-holds} 中所述。
假设 $n$ 在 $S$ 上可逆，并且对某个 $q \geq 1$，
$BC(f, n, q - 1)$ 成立，而 $BC(f, n, q)$ 不成立。
则存在交换图
$$
\xymatrix{
X \ar[d]_f & X' \ar[d] \ar[l] & Y \ar[l]^h \ar[d] \\
S & S' \ar[l] & \Spec(K) \ar[l]
}
$$
其中两个方块都是笛卡尔的，$S'$ 是仿射、整且正规的，
其函数域 $K$ 为代数闭域；并且存在整数 $d | n$，使得
$R^qh_*(\mathbf{Z}/d\mathbf{Z})$ 非零。
\end{lemma}

\begin{proof}
首先按照 Lemma \ref{etale-cohomology-lemma-base-change-does-not-hold}
选取一个图和 $\mathcal{F}$。
可以并且确实假设 $S'$ 是仿射的（这是显然的；如有疑问，
可参见该引理的证明）。令 $K'$ 为 $S'$ 的函数域，并令
$Y' = X' \times_{S'} \Spec(K')$，得到交换图
$$
\xymatrix{
X \ar[d]_f &
X' \ar[d] \ar[l] &
Y' \ar[l]^{h'} \ar[d] &
Y \ar[l] \ar[d] \\
S &
S' \ar[l] &
\Spec(K') \ar[l] &
\Spec(K) \ar[l]
}
$$
由 Lemma \ref{etale-cohomology-lemma-smooth-base-change-separably-closed}，
总正像 $R(Y \to Y')_*\mathbf{Z}/d\mathbf{Z}$
同构于 $\mathbf{Z}/d\mathbf{Z}$，此同构位于 $D(Y'_\etale)$ 中；这里用到了
$n$ 在 $S$ 上可逆这一条件。因此，由相对 Leray 谱序列可得
$Rh'_*\mathbf{Z}/d\mathbf{Z} = Rh_*\mathbf{Z}/d\mathbf{Z}$。
证毕。
\end{proof}










\section{真基变换定理}
\label{etale-cohomology-section-proper-base-change}

\noindent
本节陈述并证明真基变换定理。
我们的方法大致遵循 \cite[XII, Theorem 5.1]{SGA4} 中的证明，
并使用 Gabber（源自仿射情形）的思想稍稍简化论证。

\begin{lemma}
\label{etale-cohomology-lemma-zariski-h0-proper-over-henselian-pair}
设 $(A, I)$ 是 Hensel 对。设 $f : X \to \Spec(A)$ 是概形之间的真态射。
令 $Z = X \times_{\Spec(A)} \Spec(A/I)$。对任意层 $\mathcal{F}$，
若它定义在与 $X$ 相伴的拓扑空间上，则
$\Gamma(X, \mathcal{F}) = \Gamma(Z, \mathcal{F}|_Z)$。
\end{lemma}

\begin{proof}
我们用 Lemma \ref{etale-cohomology-lemma-h0-topological} 来证明这一点。首先，由 Properties, Lemma
\ref{properties-lemma-quasi-compact-quasi-separated-spectral}，
$X$ 的底拓扑空间是谱空间。设 $Y \subset X$ 是不可约闭子概形。
为完成证明，我们说明 $Y \cap Z = Y \times_{\Spec(A)} \Spec(A/I)$ 连通。
把 $X$ 替换为 $Y$ 后，可以假设 $X$ 不可约，而需要证明 $Z$ 连通。
设 $X \to \Spec(B) \to \Spec(A)$ 为 $f$ 的 Stein 分解
（More on Morphisms, Theorem
\ref{more-morphisms-theorem-stein-factorization-general}）。
此时 $A \to B$ 是整的，并且 $(B, IB)$ 是 Hensel 对
（More on Algebra, Lemma \ref{more-algebra-lemma-integral-over-henselian-pair}）。
因此可以假设 $X \to \Spec(A)$ 的纤维几何连通。另一方面，
$T \subset \Spec(A)$ 是 $f$ 的像，因而不可约；又因 $X$ 在 $A$ 上真，该像是闭的。
故由 More on Algebra, Lemma
\ref{more-algebra-lemma-irreducible-henselian-pair-connected}，
$T \cap V(I)$ 连通。现在
$Y \times_{\Spec(A)} \Spec(A/I) \to T \cap V(I)$
是具有连通纤维的满闭映射。结论于是由 Topology, Lemma
\ref{topology-lemma-connected-fibres-quotient-topology-connected-components} 得出。
\end{proof}

\begin{lemma}
\label{etale-cohomology-lemma-h0-proper-over-henselian-pair}
设 $(A, I)$ 是 Hensel 对。设 $f : X \to \Spec(A)$ 是概形之间的真态射。
令 $i : Z \to X$ 为 $X \times_{\Spec(A)} \Spec(A/I)$ 到 $X$ 的闭浸入。
对任意层 $\mathcal{F}$，若它定义在 $X_\etale$ 上，则
$\Gamma(X, \mathcal{F}) = \Gamma(Z, i_{small}^{-1}\mathcal{F})$。
\end{lemma}

\begin{proof}
这由 Lemma \ref{etale-cohomology-lemma-gabber-h0}、
\ref{etale-cohomology-lemma-zariski-h0-proper-over-henselian-pair}，
以及在 $X$ 上有限的任意概形在 $\Spec(A)$ 上都是真这一事实得出。
\end{proof}

\begin{lemma}
\label{etale-cohomology-lemma-h0-proper-over-henselian-local}
设 $A$ 是 Hensel 局部环。设 $f : X \to \Spec(A)$ 是概形之间的真态射。
令 $X_0 \subset X$ 为 $f$ 在闭点上方的纤维。对任意层 $\mathcal{F}$，
若它定义在 $X_\etale$ 上，则
$\Gamma(X, \mathcal{F}) = \Gamma(X_0, \mathcal{F}|_{X_0})$。
\end{lemma}

\begin{proof}
这是 Lemma \ref{etale-cohomology-lemma-h0-proper-over-henselian-pair} 的特例。
\end{proof}

\noindent
设 $f : X \to S$ 是概形之间的态射，并设
$\overline{s} : \Spec(k) \to S$ 是几何点。$f$ 在 $\overline{s}$ 处的纤维是概形
$X_{\overline{s}} = \Spec(k) \times_{\overline{s}, S} X$，
视为 $\Spec(k)$ 上的概形。若 $\mathcal{F}$ 是 $X_\etale$ 上的层，则记
$\mathcal{F}_{\overline{s}} = p_{small}^{-1}\mathcal{F}$
为 $\mathcal{F}$ 到 $(X_{\overline{s}})_\etale$ 的拉回。
下文将考虑集合
$$
\Gamma(X_{\overline{s}}, \mathcal{F}_{\overline{s}})
$$
令 $s \in S$ 为 $\overline{s}$ 的像点。令 $\kappa(s)^{sep}$
为 Definition \ref{etale-cohomology-definition-algebraic-geometric-point} 中所述的
$\kappa(s)$ 在 $k$ 中的可分代数闭包。
由 Lemma \ref{etale-cohomology-lemma-sections-base-field-extension}，拉回给出双射
$$
\Gamma(X_{\kappa(s)^{sep}},
p_{sep}^{-1} \mathcal{F})
\longrightarrow
\Gamma(X_{\overline{s}}, \mathcal{F}_{\overline{s}})
$$
其中 $p_{sep} : X_{\kappa(s)^{sep}} = \Spec(\kappa(s)^{sep}) \times_S X \to X$
是投影。

\begin{lemma}
\label{etale-cohomology-lemma-proper-pushforward-stalk}
设 $f : X \to S$ 是概形之间的真态射，并设
$\overline{s} \to S$ 是几何点。对任意层 $\mathcal{F}$，
若它定义在 $X_\etale$ 上，则典范映射
$$
(f_*\mathcal{F})_{\overline{s}} \longrightarrow
\Gamma(X_{\overline{s}}, \mathcal{F}_{\overline{s}})
$$
是双射。
\end{lemma}

\begin{proof}
由 Theorem \ref{etale-cohomology-theorem-higher-direct-images}（用于集合层）有
$$
(f_*\mathcal{F})_{\overline{s}} =
\Gamma(X \times_S \Spec(\mathcal{O}_{S, \overline{s}}^{sh}),
p_{small}^{-1}\mathcal{F})
$$
其中 $p : X \times_S \Spec(\mathcal{O}_{S, \overline{s}}^{sh}) \to X$
是投影。严格 Hensel 局部环 $\mathcal{O}_{S, \overline{s}}^{sh}$
的剩余域是 $\kappa(s)^{sep}$，故本引理由上述讨论和
Lemma \ref{etale-cohomology-lemma-h0-proper-over-henselian-local} 得出。
\end{proof}

\begin{lemma}
\label{etale-cohomology-lemma-proper-base-change-f-star}
设 $f : X \to Y$ 是概形之间的真态射，$g : Y' \to Y$
是概形之间的态射。置 $X' = Y' \times_Y X$，其投影为
$f' : X' \to Y'$ 和 $g' : X' \to X$。设 $\mathcal{F}$ 是
$X_\etale$ 上的任意层。则 $g^{-1}f_*\mathcal{F} = f'_*(g')^{-1}\mathcal{F}$。
\end{lemma}

\begin{proof}
存在典范映射 $g^{-1}f_*\mathcal{F} \to f'_*(g')^{-1}\mathcal{F}$。
具体地说，它是下述映射的伴随：
$$
f_*\mathcal{F} \longrightarrow
g_*f'_*(g')^{-1}\mathcal{F} = f_*g'_*(g')^{-1}\mathcal{F}
$$
它由 $f_*$ 作用于典范映射
$\mathcal{F} \to g'_*(g')^{-1}\mathcal{F}$ 得到。要检验该映射为同构，
可以计算其在茎上的作用。设 $y' : \Spec(k) \to Y'$ 是几何点，
其像 $y$ 位于 $Y$ 中。由 Lemma \ref{etale-cohomology-lemma-proper-pushforward-stalk}，
两个茎分别为 $\Gamma(X'_{y'}, \mathcal{F}_{y'})$ 和
$\Gamma(X_y, \mathcal{F}_y)$。这里，层 $\mathcal{F}_y$ 和
$\mathcal{F}_{y'}$ 分别是 $\mathcal{F}$ 沿投影 $X_y \to X$
和 $X'_{y'} \to X$ 的拉回。由于作为概形有 $X_y = X'_{y'}$，
可见这些茎相同。我们略去此同构与上述映射相容的验证。
\end{proof}

\noindent
现在开始讨论真基变换定理。为此先引入一些记号。考虑交换图
\begin{equation}
\label{etale-cohomology-equation-base-change-diagram}
\vcenter{
\xymatrix{
X' \ar[r]_{g'} \ar[d]_{f'} & X \ar[d]^f \\
Y' \ar[r]^g & Y
}
}
\end{equation}
其中各箭头都是概形态射。于是得到位点的交换图
$$
\xymatrix{
X'_\etale \ar[r]_{g'_{small}} \ar[d]_{f'_{small}} &
X_\etale \ar[d]^{f_{small}} \\
Y'_\etale \ar[r]^{g_{small}} &
Y_\etale
}
$$
对任意对象 $E$，若它属于 $D(X_\etale)$，则得到典范基变换映射
\begin{equation}
\label{etale-cohomology-equation-base-change}
g_{small}^{-1}Rf_{small, *}E \longrightarrow Rf'_{small, *}(g'_{small})^{-1}E
\end{equation}
它位于 $D(Y'_\etale)$ 中。见 Cohomology on Sites, Remark
\ref{sites-cohomology-remark-base-change}，其中以常值层
$\mathbf{Z}$ 作为环层。此公式中通常省略下标 ${}_{small}$。
例如，若 $E = \mathcal{F}[0]$，其中 $\mathcal{F}$ 是
$X_\etale$ 上的阿贝尔层，则基变换映射为
\begin{equation}
\label{etale-cohomology-equation-base-change-sheaf}
g^{-1}Rf_*\mathcal{F} \longrightarrow Rf'_*(g')^{-1}\mathcal{F}
\end{equation}
它位于 $D(Y'_\etale)$ 中。

\medskip\noindent
在上述一般性下，映射 (\ref{etale-cohomology-equation-base-change}) 不可能总是同构。
我们的目标是证明：如果图 (\ref{etale-cohomology-equation-base-change-diagram}) 是笛卡尔的，
$f : X \to Y$ 为真态射，$E$ 的上同调层都是挠层，且 $E$ 下有界，
那么该映射是同构。为研究这个问题，引入如下术语。
称{\it 对 $f : X \to Y$，上同调与基变换可交换}，如果
(\ref{etale-cohomology-equation-base-change-sheaf}) 对每个图
(\ref{etale-cohomology-equation-base-change-diagram})（其中 $X' = Y' \times_Y X$）
以及每个挠阿贝尔层 $\mathcal{F}$ 都是同构。

\begin{lemma}
\label{etale-cohomology-lemma-proper-base-change-in-terms-of-injectives}
设 $f : X \to Y$ 是概形之间的真态射。下列条件等价：
\begin{enumerate}
\item 对 $f$，上同调与基变换可交换（见上文）；
\item 对每个素数 $\ell$、每个内射
$\mathbf{Z}/\ell\mathbf{Z}$-模层 $\mathcal{I}$（定义在 $X_\etale$ 上），以及每个满足
$X' = Y' \times_Y X$ 的图 (\ref{etale-cohomology-equation-base-change-diagram})，
层 $R^qf'_*(g')^{-1}\mathcal{I}$ 在 $q > 0$ 时为零。
\end{enumerate}
\end{lemma}

\begin{proof}
显然 (1) 蕴含 (2)。反过来，假设 (2)，并设 $\mathcal{F}$ 是
$X_\etale$ 上的挠阿贝尔层。设 $Y' \to Y$ 是概形之间的态射，并令
$X' = Y' \times_Y X$，其投影 $g' : X' \to X$ 和
$f' : X' \to Y'$ 如图 (\ref{etale-cohomology-equation-base-change-diagram}) 所示。
我们要证明层的映射
$$
g^{-1}R^qf_*\mathcal{F} \longrightarrow R^qf'_*(g')^{-1}\mathcal{F}
$$
对所有 $q \geq 0$ 都是同构。

\medskip\noindent
对每个 $n \geq 1$，令 $\mathcal{F}[n]$ 为 $\mathcal{F}$ 中由
$n$ 零化的截面构成的子层。于是
$\mathcal{F} = \colim \mathcal{F}[n]$.
函子 $g^{-1}$ 和 $(g')^{-1}$ 与任意余极限可交换（因为它们是左伴随）。
由 Lemma \ref{etale-cohomology-lemma-relative-colimit}，沿 $f$ 或 $f'$ 取高阶正像与滤过余极限可交换。
因此有
$$
g^{-1}R^qf_*\mathcal{F} = \colim g^{-1}R^qf_*\mathcal{F}[n]
\quad\text{且}\quad
R^qf'_*(g')^{-1}\mathcal{F} =
\colim R^qf'_*(g')^{-1}\mathcal{F}[n]
$$
故只需对 $\mathcal{F}$ 被正整数 $n$ 零化的情形证明结论。

\medskip\noindent
若 $n = \ell n'$，其中 $\ell$ 是某个素数，则得到短正合列
$$
0 \to \mathcal{F}[\ell] \to \mathcal{F} \to
\mathcal{F}/\mathcal{F}[\ell] \to 0
$$
注意 $\mathcal{F}/\mathcal{F}[\ell]$ 被 $n'$ 零化。此外，若结论对
$\mathcal{F}[\ell]$ 和 $\mathcal{F}/\mathcal{F}[\ell]$ 都成立，
则由高阶正像的长正合列（以及 $5$ 引理），结论成立。
这样便归结为 $\mathcal{F}$ 被某个素数 $\ell$ 零化的情形。

\medskip\noindent
假设 $\mathcal{F}$ 被素数 $\ell$ 零化。选取内射分解
$\mathcal{F} \to \mathcal{I}^\bullet$，它位于
$D(X_\etale, \mathbf{Z}/\ell\mathbf{Z})$ 中。应用条件 (2) 和 Leray 非循环引理
（Derived Categories, Lemma \ref{derived-lemma-leray-acyclicity}），可知
$$
f'_*(g')^{-1}\mathcal{I}^\bullet
$$
计算 $Rf'_*(g')^{-1}\mathcal{F}$。再应用
Lemma \ref{etale-cohomology-lemma-proper-base-change-f-star} 即得结论。
\end{proof}

\begin{lemma}
\label{etale-cohomology-lemma-sandwich}
设 $f : X \to Y$ 和 $g : Y \to Z$ 是概形之间的真态射。假设
\begin{enumerate}
\item 对 $f$，上同调与基变换可交换；
\item 对 $g \circ f$，上同调与基变换可交换；并且
\item $f$ 是满射。
\end{enumerate}
则对 $g$，上同调与基变换可交换。
\end{lemma}

\begin{proof}
我们将使用 Lemma \ref{etale-cohomology-lemma-proper-base-change-in-terms-of-injectives}
中的等价性而不再特别说明。设 $\ell$ 是素数。
设 $\mathcal{I}$ 是一个内射层，它是 $\mathbf{Z}/\ell\mathbf{Z}$-模层，
并定义在 $Y_\etale$ 上。选取层的单射
$f^{-1}\mathcal{I} \to \mathcal{J}$，其中 $\mathcal{J}$ 是一个内射层，
它是 $\mathbf{Z}/\ell\mathbf{Z}$-模层，并定义在 $X_\etale$ 上。
由于 $f$ 是满射，映射 $\mathcal{I} \to f_*\mathcal{J}$ 是单射
（在几何点处看茎即可）。由于 $\mathcal{I}$ 内射，可知
$\mathcal{I}$ 是 $f_*\mathcal{J}$ 的直和项。因此只需对
$f_*\mathcal{J}$ 证明所需的消失。

\medskip\noindent
设 $Z' \to Z$ 是概形之间的态射，并置
$Y' = Z' \times_Z Y$ 和 $X' = Z' \times_Z X = Y' \times_ Y X$。
记 $a : X' \to X$、$b : Y' \to Y$ 和 $c : Z' \to Z$ 为投影；
$f' : X' \to Y'$ 和 $g' : Y' \to Z'$ 也同样记。由
Lemma \ref{etale-cohomology-lemma-proper-base-change-f-star} 有
$b^{-1}f_*\mathcal{J} = f'_*a^{-1}\mathcal{J}$.
另一方面，已知 $R^qf'_*a^{-1}\mathcal{J}$ 和
$R^q(g' \circ f')_*a^{-1}\mathcal{J}$ 在 $q > 0$ 时为零。
利用谱序列
(Cohomology on Sites，引理 \ref{sites-cohomology-lemma-relative-Leray})
$$
R^pg'_* R^qf'_* a^{-1}\mathcal{J} \Rightarrow
R^{p + q}(g' \circ f')_* a^{-1}\mathcal{J}
$$
可得
$ R^pg'_*(b^{-1}f_*\mathcal{J}) = R^pg'_*(f'_*a^{-1}\mathcal{J}) = 0$
在 $p > 0$ 时成立，这正是所需结论。
\end{proof}

\begin{lemma}
\label{etale-cohomology-lemma-composition}
设 $f : X \to Y$ 和 $g : Y \to Z$ 是概形之间的真态射。假设
\begin{enumerate}
\item 对 $f$，上同调与基变换可交换；并且
\item 对 $g$，上同调与基变换可交换。
\end{enumerate}
则对 $g \circ f$，上同调与基变换可交换。
\end{lemma}

\begin{proof}
我们将使用 Lemma \ref{etale-cohomology-lemma-proper-base-change-in-terms-of-injectives}
中的等价性而不再特别说明。设 $\ell$ 是素数。
设 $\mathcal{I}$ 是一个内射层，它是 $\mathbf{Z}/\ell\mathbf{Z}$-模层，
并定义在 $X_\etale$ 上。于是 $f_*\mathcal{I}$ 是一个内射层，
它是 $\mathbf{Z}/\ell\mathbf{Z}$-模层，并定义在 $Y_\etale$ 上
（Cohomology on Sites, Lemma
\ref{sites-cohomology-lemma-pushforward-injective-flat}）。
结论可由此形式地得出，但我们也把论证写明。

\medskip\noindent
设 $Z' \to Z$ 是概形之间的态射，并置
$Y' = Z' \times_Z Y$ 和 $X' = Z' \times_Z X = Y' \times_ Y X$。
记 $a : X' \to X$、$b : Y' \to Y$ 和 $c : Z' \to Z$ 为投影；
$f' : X' \to Y'$ 和 $g' : Y' \to Z'$ 也同样记。由
Lemma \ref{etale-cohomology-lemma-proper-base-change-f-star} 有
$b^{-1}f_*\mathcal{I} = f'_*a^{-1}\mathcal{I}$.
另一方面，已知 $R^qf'_*a^{-1}\mathcal{I}$ 和
$R^q(g')_*b^{-1}f_*\mathcal{I}$ 在 $q > 0$ 时为零。
利用谱序列
(Cohomology on Sites，引理 \ref{sites-cohomology-lemma-relative-Leray})
$$
R^pg'_* R^qf'_* a^{-1}\mathcal{I} \Rightarrow
R^{p + q}(g' \circ f')_* a^{-1}\mathcal{I}
$$
可得 $R^p(g' \circ f')_*a^{-1}\mathcal{I} = 0$ 在
$p > 0$ 时成立，这正是所需结论。
\end{proof}

\begin{lemma}
\label{etale-cohomology-lemma-finite}
\begin{slogan}
\'Etale 上同调的真基变换对有限态射成立。
\end{slogan}
设 $f : X \to Y$ 是概形之间的有限态射。
则对 $f$，上同调与基变换可交换。
\end{lemma}

\begin{proof}
注意有限态射是真态射，见
Morphisms, Lemma \ref{morphisms-lemma-finite-proper}。
此外，有限态射的基变换仍然有限，见
Morphisms, Lemma \ref{morphisms-lemma-base-change-finite}。
因此，将 Lemma \ref{etale-cohomology-lemma-proper-base-change-in-terms-of-injectives}
与 Proposition \ref{etale-cohomology-proposition-finite-higher-direct-image-zero}
结合便得结论。
\end{proof}

\begin{lemma}
\label{etale-cohomology-lemma-reduce-to-P1}
要证明每个概形真态射的上同调都与基变换可交换，
只需证明该结论对态射 $\mathbf{P}^1_S \to S$ 成立，其中 $S$ 是任意概形。
\end{lemma}

\begin{proof}
设 $f : X \to Y$ 是概形之间的真态射。设 $Y = \bigcup Y_i$
是仿射开覆盖，并置 $X_i = f^{-1}(Y_i)$。如果能够证明对
$X_i \to Y_i$ 上同调与基变换可交换，那么对 $f$ 也可交换。
这是因为高阶正像的形成与基底上的 Zariski（甚至 \'etale）局部化可交换，见
Lemma \ref{etale-cohomology-lemma-higher-direct-images}。因此可以假设 $Y$ 仿射。

\medskip\noindent
设 $Y$ 是仿射概形，并设 $X \to Y$ 是真态射。由 Chow 引理，存在交换图
$$
\xymatrix{
X \ar[rd] & X' \ar[d] \ar[l]^\pi \ar[r] & \mathbf{P}^n_Y \ar[dl] \\
& Y &
}
$$
其中 $X' \to \mathbf{P}^n_Y$ 是浸入，而 $\pi : X' \to X$
是真满射，见 Limits, Lemma \ref{limits-lemma-chow-finite-type}。
由于 $X \to Y$ 真，可知 $X' \to Y$ 真
（Morphisms, Lemma \ref{morphisms-lemma-composition-proper}）。
故 $X' \to \mathbf{P}^n_Y$ 是闭浸入
（Morphisms, Lemma \ref{morphisms-lemma-image-proper-scheme-closed}）。
从而 $X' \to X \times_Y \mathbf{P}^n_Y = \mathbf{P}^n_X$
是闭浸入（因为它是像闭的浸入）。

\medskip\noindent
由 Lemma \ref{etale-cohomology-lemma-sandwich}，只需证明对 $\pi$ 和 $X' \to Y$，
上同调与基变换可交换。这两个态射都可分解为一个闭浸入，
再接投影 $\mathbf{P}^n_S \to S$（其中 $S$ 为某个概形）。
由 Lemma \ref{etale-cohomology-lemma-finite}，结论对闭浸入成立（因为闭浸入有限）。
由 Lemma \ref{etale-cohomology-lemma-composition}，只需对投影
$\mathbf{P}^n_S \to S$ 证明结论。

\medskip\noindent
对每个 $n \geq 1$，存在有限满态射
$$
\mathbf{P}^1_S \times_S \ldots \times_S \mathbf{P}^1_S
\longrightarrow
\mathbf{P}^n_S
$$
它在坐标上由下式给出：
$$
((x_1 : y_1), (x_2 : y_2), \ldots, (x_n : y_n))
\longmapsto
(F_0 : \ldots : F_n)
$$
其中 $F_0, \ldots, F_n$ 是变量 $x_1, \ldots, y_n$
的整系数多项式，满足
$$
\prod (x_i t + y_i) = F_0 t^n + F_1 t^{n - 1} + \ldots + F_n
$$
再次应用 Lemmas \ref{etale-cohomology-lemma-sandwich}, \ref{etale-cohomology-lemma-finite}, and
\ref{etale-cohomology-lemma-composition}，即知本引理成立。
\end{proof}

\begin{theorem}[真基变换]
\label{etale-cohomology-theorem-proper-base-change}
设 $f : X \to Y$ 是概形之间的真态射，$g : Y' \to Y$
是概形之间的态射。置 $X' = Y' \times_Y X$，并考虑笛卡尔图
$$
\xymatrix{
X' \ar[r]_{g'} \ar[d]_{f'} & X \ar[d]^f \\
Y' \ar[r]^g & Y
}
$$
设 $\mathcal{F}$ 是 $X_\etale$ 上的挠阿贝尔层。则基变换映射
$$
g^{-1}Rf_*\mathcal{F} \longrightarrow Rf'_*(g')^{-1}\mathcal{F}
$$
是同构。
\end{theorem}

\begin{proof}
按照上面引入的术语，这就是说，对概形之间的每个真态射，
上同调与基变换可交换。由 Lemma \ref{etale-cohomology-lemma-reduce-to-P1}，
只需证明上同调对态射 $\mathbf{P}^1_S \to S$ 与基变换可交换，
其中 $S$ 是任意概形。

\medskip\noindent
设 $S$ 是严格 Hensel 局部环的谱，其闭点为 $s$。置
$X = \mathbf{P}^1_S$ 和 $X_0 = X_s = \mathbf{P}^1_s$。
设 $\mathcal{F}$ 是一个 $\mathbf{Z}/\ell\mathbf{Z}$-模层，
定义在 $X_\etale$ 上。证明的关键是
$$
H^q_\etale(X, \mathcal{F}) = H^q_\etale(X_0, \mathcal{F}|_{X_0}).
$$
具体地，选取分解 $\mathcal{F} \to \mathcal{I}^\bullet$，
其各项是内射 $\mathbf{Z}/\ell\mathbf{Z}$-模层。于是
$\mathcal{I}^\bullet|_{X_0}$ 是 $\mathcal{F}|_{X_0}$ 的一个分解，
其各项对右函子 $H^0_\etale(X_0, -)$ 非循环，见
Lemma \ref{etale-cohomology-lemma-efface-cohomology-on-fibre-by-finite-cover}。
Leray 非循环引理说明，右端由复形
$H^0_\etale(X_0, \mathcal{I}^\bullet|_{X_0})$ 计算；由
Lemma \ref{etale-cohomology-lemma-h0-proper-over-henselian-local}，它等于
$H^0_\etale(X, \mathcal{I}^\bullet)$。该复形计算左端。

\medskip\noindent
现在设 $S$ 是一般概形，$\mathcal{F}$ 是一个
$\mathbf{Z}/\ell\mathbf{Z}$-模层，定义在 $X_\etale$ 上。
设 $\overline{s} : \Spec(k) \to S$ 是 $S$ 的几何点，
它位于 $s \in S$ 上方。于是
$$
(R^qf_*\mathcal{F})_{\overline{s}} =
H^q_\etale(\mathbf{P}^1_{\mathcal{O}_{S, \overline{s}}^{sh}},
\mathcal{F}|_{\mathbf{P}^1_{\mathcal{O}_{S, \overline{s}}^{sh}}}) =
H^q_\etale(\mathbf{P}^1_{\kappa(s)^{sep}},
\mathcal{F}|_{\mathbf{P}^1_{\kappa(s)^{sep}}})
$$
其中 $\kappa(s)^{sep}$ 是 $\mathcal{O}_{S, \overline{s}}^{sh}$ 的剩余域，
即 $\kappa(s)$ 在 $k$ 中的可分代数闭包。
第一个等号由 Theorem \ref{etale-cohomology-theorem-higher-direct-images} 得出，
第二个等号由上一段的显示公式得出。

\medskip\noindent
最后，考虑概形之间的任意态射 $g : T \to S$，其中
$S$ 和 $\mathcal{F}$ 如上。令 $f' : \mathbf{P}^1_T \to T$ 为投影，
并令 $g' : \mathbf{P}^1_T \to \mathbf{P}^1_S$ 为 $g$ 诱导的态射。
考虑基变换映射
$$
g^{-1}R^qf_*\mathcal{F}
\longrightarrow
R^qf'_*(g')^{-1}\mathcal{F}
$$
设 $\overline{t}$ 是 $T$ 的几何点，其像为
$\overline{s} = g(\overline{t})$。由上述讨论，该映射在
$\overline{t}$ 处的茎上的映射为
$$
H^q_\etale(\mathbf{P}^1_{\kappa(s)^{sep}},
\mathcal{F}|_{\mathbf{P}^1_{\kappa(s)^{sep}}})
\longrightarrow
H^q_\etale(\mathbf{P}^1_{\kappa(t)^{sep}},
\mathcal{F}|_{\mathbf{P}^1_{\kappa(t)^{sep}}})
$$
由于 $\kappa(s)^{sep} \subset \kappa(t)^{sep}$，由
Lemma \ref{etale-cohomology-lemma-base-change-dim-1-separably-closed}，该映射是同构。

\medskip\noindent
这证明了对于 $\mathbf{P}^1_S \to S$ 和
$\mathbf{Z}/\ell\mathbf{Z}$-模层，上同调与基变换可交换。
特别地，对内射 $\mathbf{Z}/\ell\mathbf{Z}$-模层，
任意基变换后的高阶正像都为零。换言之，
引理 \ref{etale-cohomology-lemma-proper-base-change-in-terms-of-injectives}
的条件 (2) 成立，证明完成。
\end{proof}

\begin{lemma}
\label{etale-cohomology-lemma-proper-base-change}
设 $f : X \to Y$ 是概形之间的真态射，$g : Y' \to Y$
是概形之间的态射。置 $X' = Y' \times_Y X$，并记
$f' : X' \to Y'$ 和 $g' : X' \to X$ 为投影。
设 $E \in D^+(X_\etale)$ 的上同调层都是挠层。
则基变换映射 (\ref{etale-cohomology-equation-base-change})
$g^{-1}Rf_*E \to Rf'_*(g')^{-1}E$ 是同构。
\end{lemma}

\begin{proof}
这是利用下列谱序列从真基变换定理
（Theorem \ref{etale-cohomology-theorem-proper-base-change}）直接得到的推论：
$$
E_2^{p, q} = R^pf_*H^q(E)
\quad\text{且}\quad
{E'}_2^{p, q} = R^pf'_*(g')^{-1}H^q(E)
$$
它们分别收敛于 $R^nf_*E$ 和 $R^nf'_*(g')^{-1}E$。
这些谱序列在 Derived Categories, Lemma
\ref{derived-lemma-two-ss-complex-functor} 中构造。略去若干细节。
\end{proof}

\begin{lemma}
\label{etale-cohomology-lemma-proper-base-change-stalk}
设 $f : X \to Y$ 是概形之间的真态射，并设 $\overline{y} \to Y$
是几何点。
\begin{enumerate}
\item 对挠阿贝尔层 $\mathcal{F}$，若它定义在 $X_\etale$ 上，则有
$(R^nf_*\mathcal{F})_{\overline{y}} =
H^n_\etale(X_{\overline{y}}, \mathcal{F}_{\overline{y}})$.
\item 对上同调层均为挠层的 $E \in D^+(X_\etale)$，有
$(R^nf_*E)_{\overline{y}} =
H^n_\etale(X_{\overline{y}}, E|_{X_{\overline{y}}})$.
\end{enumerate}
\end{lemma}

\begin{proof}
在陈述中，$\mathcal{F}_{\overline{y}}$ 表示 $\mathcal{F}$
到概形论纤维 $X_{\overline{y}} = \overline{y} \times_Y X$ 的拉回。
由于沿 $\overline{y} \to Y$ 拉回给出 $\mathcal{F}$ 的茎，
本引理的第一条是 Theorem \ref{etale-cohomology-theorem-proper-base-change} 的特例；
第二条是 Lemma \ref{etale-cohomology-lemma-proper-base-change} 的特例。
\end{proof}






\section{真基变换的应用}
\label{etale-cohomology-section-applications-proper-base-change}

\noindent
本节讨论真基变换定理的一些颇为直接的推论。

\begin{lemma}
\label{etale-cohomology-lemma-base-change-separably-closed}
设 $K/k$ 是可分闭域的扩张。设 $X$ 是 $k$ 上的真概形。
设 $\mathcal{F}$ 是 $X_\etale$ 上的挠阿贝尔层。则映射
$H^q_\etale(X, \mathcal{F}) \to
H^q_\etale(X_K, \mathcal{F}|_{X_K})$ 对 $q \geq 0$ 是同构。
\end{lemma}

\begin{proof}
考察茎可知，这是 Theorem \ref{etale-cohomology-theorem-proper-base-change} 的特例。
\end{proof}

\begin{lemma}
\label{etale-cohomology-lemma-cohomological-dimension-proper}
设 $f : X \to Y$ 是概形之间的真态射，其所有纤维的维数均为 $\leq n$。
则对任意挠阿贝尔层 $\mathcal{F}$，若它定义在 $X_\etale$ 上，
都有 $R^qf_*\mathcal{F} = 0$，只要 $q > 2n$。
\end{lemma}

\begin{proof}
我们对 $n$ 归纳，同时为所有真态射证明此结论。

\medskip\noindent
若 $n = 0$，则 $f$ 是有限态射
(More on Morphisms，引理 \ref{more-morphisms-lemma-characterize-finite})
，而结论由 Proposition \ref{etale-cohomology-proposition-finite-higher-direct-image-zero} 成立。

\medskip\noindent
若 $n > 0$，则由 Lemma \ref{etale-cohomology-lemma-proper-base-change-stalk}，
只需证明 $H^i_\etale(X, \mathcal{F}) = 0$，其中 $i > 2n$，
$X$ 是真概形，满足 $\dim(X) \leq n$，且定义在代数闭域 $k$ 上，
而 $\mathcal{F}$ 是 $X$ 上的挠阿贝尔层。

\medskip\noindent
若 $n = 1$，结论由 Theorem \ref{etale-cohomology-theorem-vanishing-curves} 得出。
假设 $n > 1$。由 Proposition \ref{etale-cohomology-proposition-topological-invariance}，
可以用约化概形代替 $X$。令 $\nu : X^\nu \to X$ 为正规化。
这是满的双有理有限态射
（见 Varieties, Lemma \ref{varieties-lemma-normalization-locally-algebraic}），
因而在稠密开子集 $U \subset X$ 上是同构
（Morphisms, Lemma \ref{morphisms-lemma-birational-birational}）。
于是 $c : \mathcal{F} \to \nu_*\nu^{-1}\mathcal{F}$ 是单射
（因为 $\nu$ 满），并且在 $U$ 上是同构。记 $i : Z \to X$
为 $U$ 的补集的包含映射。由于 $U$ 在 $X$ 中稠密，有
$\dim(Z) < \dim(X) = n$。由
Proposition \ref{etale-cohomology-proposition-closed-immersion-pushforward}，有
$\Coker(c) = i_*\mathcal{G}$，其中 $\mathcal{G}$ 是
$Z_\etale$ 上的某个挠阿贝尔层。于是
$H^q_\etale(X, \Coker(c)) = H^q_\etale(Z, \mathcal{G})$
（由 Proposition \ref{etale-cohomology-proposition-finite-higher-direct-image-zero}
和 Leray 谱序列）；归纳假设说明 $c$ 的余核只在次数
$\leq 2(n - 1)$ 上有非零上同调。因此只需对
$\nu_*\nu^{-1}\mathcal{F}$ 证明结论。由于 $\nu$ 有限，
这归结为证明 $H^i_\etale(X^\nu, \nu^{-1}\mathcal{F})$
在 $i > 2n$ 时为零。下一段处理这种情形。

\medskip\noindent
假设 $X$ 是 $k$ 上维数为 $n$ 的整正规真概形。
选取非常值有理函数 $f$；它定义在 $X$ 上。子概形
$X' \subset X \times \mathbf{P}^1_k$ 是 $f$ 的图，并位于下图中：
$$
X \xleftarrow{b} X' \xrightarrow{f} \mathbf{P}^1_k
$$
注意 $b$ 在开子概形 $U \subset X$ 上是同构；其补集是闭子概形
$Z \subset X$，其余维为 $\geq 2$。具体地，$U$ 是 $f$ 的定义域，
它包含所有余维为 $1$ 且位于 $X$ 中的点，见
Morphisms，引理 \ref{morphisms-lemma-rational-map-from-reduced-to-separated}
和 \ref{morphisms-lemma-extend-across}
（并结合 Serre 正规性判据，见
Properties, Lemma \ref{properties-lemma-criterion-normal}）。
此外，$b$ 的纤维维数为 $\leq 1$（因为它们是 $\mathbf{P}^1$ 的闭子概形）。
因此，由归纳假设，$R^ib_*b^{-1}\mathcal{F}$ 仅当
$i \in \{0, 1, 2\}$ 时才可能非零。选取可区别三角
$$
\mathcal{F} \to Rb_*b^{-1}\mathcal{F} \to Q \to \mathcal{F}[1]
$$
像前面一样，利用 $\mathcal{F} \to b_*b^{-1}\mathcal{F}$ 是单射这一事实，
并结合刚才的结论，可知 $Q$ 仅在次数 $0, 1, 2$ 上有非零上同调层，
且这些层支撑在 $Z$ 上。此外，由 Lemma \ref{etale-cohomology-lemma-torsion-direct-image}，
这些上同调层都是挠层。利用归纳假设和
Derived Categories，引理 \ref{derived-lemma-two-ss-complex-functor}
的谱序列，可知 $H^i(X, Q)$ 在
$i > 2 + 2\dim(Z) \leq 2 + 2(n - 2) = 2n - 2$ 时为零。
因此只需证明 $H^i(X', b^{-1}\mathcal{F}) = 0$ 在
$i > 2n$ 时成立。此时使用态射
$$
f : X' \to \mathbf{P}^1_k
$$
其纤维维数为 $< n$。因此，由归纳假设，
$R^if_*b^{-1}\mathcal{F} = 0$ 在 $i > 2(n - 1)$ 时成立。
最后应用 Leray 谱序列
$$
H^i(\mathbf{P}^1_k, R^jf_*b^{-1}\mathcal{F})
\Rightarrow
H^{i + j}(X', b^{-1}\mathcal{F})
$$
以及事实 $\dim(\mathbf{P}^1_k) = 1$ 即得结论。
\end{proof}

\noindent
使用模 $n$ 系数时，可以对无界复形作真基变换。

\begin{lemma}
\label{etale-cohomology-lemma-proper-base-change-mod-n}
设 $f : X \to Y$ 是概形之间的真态射，$g : Y' \to Y$
是概形之间的态射。置 $X' = Y' \times_Y X$，并记
$f' : X' \to Y'$ 和 $g' : X' \to X$ 为投影。
设 $n \geq 1$ 是整数。设 $E \in D(X_\etale, \mathbf{Z}/n\mathbf{Z})$。
则基变换映射 (\ref{etale-cohomology-equation-base-change})
$g^{-1}Rf_*E \to Rf'_*(g')^{-1}E$ 是同构。
\end{lemma}

\begin{proof}
只需在 $Y$ 和 $Y'$ 都拟紧时证明。由 Morphisms, Lemma
\ref{morphisms-lemma-morphism-finite-type-bounded-dimension}，
可知 $f : X \to Y$ 和 $f' : X' \to Y'$ 的纤维维数有界。
因此 Lemma \ref{etale-cohomology-lemma-cohomological-dimension-proper} 说明
$$
f_* :
\textit{Mod}(X_\etale, \mathbf{Z}/n\mathbf{Z})
\longrightarrow
\textit{Mod}(Y_\etale, \mathbf{Z}/n\mathbf{Z})
$$
以及
$$
f'_* :
\textit{Mod}(X'_\etale, \mathbf{Z}/n\mathbf{Z})
\longrightarrow
\textit{Mod}(Y'_\etale, \mathbf{Z}/n\mathbf{Z})
$$
在 Derived Categories, Lemma
\ref{derived-lemma-unbounded-right-derived} 的意义下具有有限上同调维数。
选取 K-内射复形 $\mathcal{I}^\bullet$，它由
$\mathbf{Z}/n\mathbf{Z}$-模层组成；其每一项 $\mathcal{I}^n$ 都是内射
$\mathbf{Z}/n\mathbf{Z}$-模层，并且它表示 $E$。见 Injectives, Theorem
\ref{injectives-theorem-K-injective-embedding-grothendieck}。
由通常的真基变换定理，有
$R^qf'_*(g')^{-1}\mathcal{I}^n = 0$，只要 $q > 0$，见
Theorem \ref{etale-cohomology-theorem-proper-base-change}。因此，由
Derived Categories, Lemma \ref{derived-lemma-unbounded-right-derived}，
可以计算 $Rf'_*(g')^{-1}E$，所用复形为
$f'_*(g')^{-1}\mathcal{I}^\bullet$。再次应用通常的真基变换定理，
可知它等于 $g^{-1}f_*\mathcal{I}^\bullet$，这正是所需结论。
\end{proof}

\begin{lemma}
\label{etale-cohomology-lemma-pull-out-constant}
设 $X$ 是拟紧拟分离概形。设 $E \in D^+(X_\etale)$ 且
$K \in D^+(\mathbf{Z})$。则
$$
R\Gamma(X, E \otimes_\mathbf{Z}^\mathbf{L} \underline{K}) =
R\Gamma(X, E) \otimes_\mathbf{Z}^\mathbf{L} K
$$
\end{lemma}

\begin{proof}
假设 $H^i(E) = 0$ 对 $i \geq a$ 成立，且 $H^j(K) = 0$
对 $j \geq b$ 成立。可以表示 $K$，所用的是下有界复形
$K^\bullet$，其各项为无挠 $\mathbf{Z}$-模。（选取 K-平坦复形
$L^\bullet$ 表示 $K$，
再取 $K^\bullet = \tau_{\geq b - 1}L^\bullet$。这是可行的，因为
$\mathbf{Z}$ 的整体维数为 $1$。见 More on Algebra, Lemma
\ref{more-algebra-lemma-last-one-flat}。）可以表示 $E$，所用的是下有界复形
$\mathcal{E}^\bullet$。于是
$E \otimes_\mathbf{Z}^\mathbf{L} \underline{K}$ 由下列复形表示：
$$
\text{Tot}(\mathcal{E}^\bullet \otimes_\mathbf{Z} \underline{K}^\bullet)
$$
利用可区别三角
$$
\sigma_{\geq -b + n + 1}K^\bullet
\to K^\bullet \to
\sigma_{\leq -b + n}K^\bullet
$$
以及显然的消失
$$
H^n(X,
\text{Tot}(\mathcal{E}^\bullet \otimes_\mathbf{Z}
\sigma_{\geq -a + n + 1}\underline{K}^\bullet) = 0
$$
和
$$
H^n(R\Gamma(X, E)
\otimes_\mathbf{Z}^\mathbf{L} \sigma_{\geq -a + n + 1}K^\bullet) = 0
$$
可归结为 $K^\bullet$ 是由平坦 $\mathbf{Z}$-模组成的有界复形的情形。
重复此论证，又归结为 $K^\bullet$ 等于位于某个次数上的单个平坦
$\mathbf{Z}$-模的情形。接着，对 $\mathcal{E}^\bullet$ 使用笨截断，
以完全相同的方式归结为 $\mathcal{E}^\bullet$ 是位于某个次数上的
单个阿贝尔层的情形。因此只需证明
$$
H^n(X, \mathcal{E} \otimes_\mathbf{Z} \underline{M}) =
H^n(X, \mathcal{E}) \otimes_\mathbf{Z} M
$$
其中 $M$ 是平坦 $\mathbf{Z}$-模，而 $\mathcal{E}$ 是 $X$ 上的阿贝尔层。
在这种情形下，把 $M$ 写成有限自由 $\mathbf{Z}$-模的滤过余极限
（Lazard 定理，见 Algebra, Theorem \ref{algebra-theorem-lazard}）。
由 Theorem \ref{etale-cohomology-theorem-colimit}，这归结为有限自由
$\mathbf{Z}$-模 $M$ 的情形，此时结论显然成立。
\end{proof}

\begin{lemma}
\label{etale-cohomology-lemma-projection-formula-proper}
设 $f : X \to Y$ 是概形之间的真态射。
设 $E \in D^+(X_\etale)$ 的上同调层都是挠层。
设 $K \in D^+(Y_\etale)$。则
$$
Rf_*E \otimes_\mathbf{Z}^\mathbf{L} K =
Rf_*(E \otimes_\mathbf{Z}^\mathbf{L} f^{-1}K)
$$
该等式成立于 $D^+(Y_\etale)$ 中。
\end{lemma}

\begin{proof}
由 Cohomology on Sites, Section
\ref{sites-cohomology-section-projection-formula}，存在从左端到右端的典范映射。
我们在茎上检验该等式。回顾，导出张量积的计算与拉回可交换。见
Cohomology on Sites，引理
\ref{sites-cohomology-lemma-pullback-tensor-product}。因此有
$$
(E \otimes_\mathbf{Z}^\mathbf{L} f^{-1}K)_{\overline{x}}
=
E_{\overline{x}} \otimes_\mathbf{Z}^\mathbf{L} K_{\overline{y}}
$$
其中 $\overline{y}$ 是 $\overline{x}$ 在 $Y$ 中的像。
由于 $\mathbf{Z}$ 的整体维数为 $1$，可知该复形在次数
$< - 1 + a + b$ 上的上同调为零，只要
$H^i(E) = 0$ 对 $i \geq a$ 成立且 $H^j(K) = 0$
对 $j \geq b$ 成立。此外，由于 $H^i(E)$ 对每个 $i$
都是挠阿贝尔层，复形
$E \otimes_\mathbf{Z}^\mathbf{L} K$ 的上同调层也是如此。事实上有
$$
(E \otimes_\mathbf{Z}^\mathbf{L} f^{-1}K)
\otimes_{\mathbf{Z}}^\mathbf{L} \mathbf{Q} =
(E \otimes_\mathbf{Z}^\mathbf{L} \mathbf{Q})
\otimes_{\mathbf{Q}}^\mathbf{L}
(f^{-1}K \otimes_{\mathbf{Z}}^\mathbf{L} \mathbf{Q})
$$
它在导出范畴中为零。这样，Lemma \ref{etale-cohomology-lemma-proper-base-change-stalk}
可应用于两端，从而只需证明
$$
R\Gamma(X_{\overline{y}},
E|_{X_{\overline{y}}} \otimes_\mathbf{Z}^\mathbf{L}
(X_{\overline{y}} \to \overline{y})^{-1}K_{\overline{y}}) =
R\Gamma(X_{\overline{y}},
E|_{X_{\overline{y}}}) \otimes_\mathbf{Z}^\mathbf{L} K_{\overline{y}}
$$
这已在 Lemma \ref{etale-cohomology-lemma-pull-out-constant} 中证明。
\end{proof}








\section{局部无环性}
\label{etale-cohomology-section-local-acyclicity}

\noindent
本节从光滑基变换定理推出光滑态射的局部无环性。在 SGA 4
或 SGA 4.5 中，作者先证明光滑态射的一个局部无环性版本，
再由此推出光滑基变换定理。

\medskip\noindent
我们将采用 Deligne \cite[Definition 2.12, page 242]{SGA4.5}
给出的局部无环性表述。设 $f : X \to S$ 是概形态射。
设 $\overline{x}$ 是 $X$ 的几何点，其像
$\overline{s} = f(\overline{x})$ 是 $S$ 的点。设 $\overline{t}$ 是
$\Spec(\mathcal{O}^{sh}_{S, \overline{s}})$ 的几何点。
我们得到交换图
$$
\xymatrix{
F_{\overline{x}, \overline{t}} =
\overline{t} \times_{\Spec(\mathcal{O}^{sh}_{S, \overline{s}})}
\Spec(\mathcal{O}^{sh}_{X, \overline{x}}) \ar[r] \ar[d] &
\Spec(\mathcal{O}^{sh}_{X, \overline{x}}) \ar[r] \ar[d] &
X \ar[d] \\
\overline{t} \ar[r] &
\Spec(\mathcal{O}^{sh}_{S, \overline{s}}) \ar[r] &
S
}
$$
概形 $F_{\overline{x}, \overline{t}}$ 称为
{\it $f$ 在 $\overline{x}$ 处的一个消失圈簇}。
设 $K$ 是 $D(X_\etale)$ 的对象。对任意概形态射
$g : Y\to X$，我们以 $R\Gamma(Y, K)$ 记
$R\Gamma(Y_\etale, g_{small}^{-1}K)$。由于
$\mathcal{O}^{sh}_{X, \overline{x}}$ 是严格 Hensel 的，故有
$K_{\overline{x}} = R\Gamma(\Spec(\mathcal{O}^{sh}_{X, \overline{x}}), K)$.
因此得到典范映射
\begin{equation}
\label{etale-cohomology-equation-alpha-K}
\alpha_{K, \overline{x}, \overline{t}} :
K_{\overline{x}}
\longrightarrow
R\Gamma(F_{\overline{x}, \overline{t}}, K)
\end{equation}
它由沿
$F_{\overline{x}, \overline{t}} \to \Spec(\mathcal{O}^{sh}_{X, \overline{x}})$
拉回上同调得到。

\begin{definition}
\label{etale-cohomology-definition-locally-acyclic}
\begin{reference}
\cite[Definition 2.12, page 242]{SGA4.5} 及
\cite[Definition (1.3), page 54]{SGA4.5}
\end{reference}
设 $f : X \to S$ 是概形态射。
设 $K$ 是 $D(X_\etale)$ 的对象。
\begin{enumerate}
\item
设 $\overline{x}$ 是 $X$ 的几何点，其像为
$\overline{s} = f(\overline{x})$。
称 $f$ {\it 在 $\overline{x}$ 处相对于 $K$ 局部无环}，如果对每个几何点
$\overline{t}$（它是
$\Spec(\mathcal{O}^{sh}_{S, \overline{s}})$ 的几何点），映射 (\ref{etale-cohomology-equation-alpha-K}) 都是同构\footnote{我们不要求
$\overline{t}$ 是
$\Spec(\mathcal{O}^{sh}_{S, \overline{s}})$ 的代数几何点。通常利用
引理 \ref{etale-cohomology-lemma-smooth-base-change-separably-closed}
可约化到这种情形。}。
\item 称 $f$ {\it 相对于 $K$ 局部无环}，如果 $f$ 在 $\overline{x}$ 处
相对于 $K$ 局部无环，其中 $\overline{x}$ 取遍 $X$ 的所有几何点。
\item 称 $f$ {\it 相对于 $K$ 普遍局部无环}，如果对任意概形态射
$S' \to S$，基变换 $f' : X' \to S'$ 都相对于 $K$ 在 $X'$ 上的拉回局部无环。
\item 称 $f$ {\it 局部无环}，如果令 $\overline{x}$ 取遍 $X$ 的所有几何点，
并令 $n$ 取遍与 $\kappa(\overline{x})$ 的特征互素的所有整数，
态射 $f$ 都在 $\overline{x}$ 处相对于取值为
$\mathbf{Z}/n\mathbf{Z}$ 的常值层局部无环。
\item 称 $f$ {\it 普遍局部无环}，如果对任意概形态射
$S' \to S$，基变换 $f' : X' \to S'$ 都局部无环。
\end{enumerate}
\end{definition}

\noindent
设 $M$ 是阿贝尔群。那么 $f : X \to S$ 相对于常值层
$\underline{M}$ 的局部无环性归结为要求
$$
H^q(F_{\overline{x}, \overline{t}}, \underline{M}) =
\left\{
\begin{matrix}
M & \text{若} & q = 0 \\
0 & \text{若} & q \not = 0
\end{matrix}
\right.
$$
对任意几何点 $\overline{x}$（它是 $X$ 的几何点）以及任意几何点
$\overline{t}$（它是 $\Spec(\mathcal{O}^{sh}_{S, f(\overline{x})})$ 的几何点）
都成立。由此可见，局部无环等价于几何纤维
$F_{\overline{x}, \overline{t}}$ 的高阶上同调群消失；这些纤维来自
$\overline{x}$ 和 $\overline{s}$ 处严格 Hensel 化之间的映射。

\begin{proposition}
\label{etale-cohomology-proposition-smooth-locally-acyclic}
设 $f : X \to S$ 是光滑概形态射。则 $f$ 普遍局部无环。
\end{proposition}

\begin{proof}
由于光滑态射的基变换仍光滑，只须证明光滑态射局部无环。
设 $\overline{x}$ 是 $X$ 的几何点，其像为
$\overline{s} = f(\overline{x})$。设
$\overline{t}$ 是
$\Spec(\mathcal{O}^{sh}_{S, f(\overline{x})})$ 的几何点。
由于我们要证明的是环同态
$\mathcal{O}^{sh}_{S, \overline{s}} \to \mathcal{O}^{sh}_{X, \overline{x}}$
的一个性质（见 Definition \ref{etale-cohomology-definition-locally-acyclic} 后的讨论），
所以可且确实把 $f : X \to S$ 替换为基变换
$X \times_S \Spec(\mathcal{O}^{sh}_{S, \overline{s}}) \to
\Spec(\mathcal{O}^{sh}_{S, \overline{s}})$.
因此可且确实假设 $S$ 是严格 Hensel 局部环的谱，并且
$\overline{s}$ 位于 $S$ 的闭点之上。

\medskip\noindent
我们将把 Lemma \ref{etale-cohomology-lemma-base-change-f-star-general-stalks}
应用于图
$$
\xymatrix{
X \ar[d]_f & X_{\overline{t}} \ar[l]^h \ar[d]^e \\
S & \overline{t} \ar[l]_g
}
$$
以及层 $\mathcal{F} = \underline{M}$，其中 $M = \mathbf{Z}/n\mathbf{Z}$，
$n$ 是与 $\overline{x}$ 的剩余域特征互素的某个整数。由光滑基变换可知，映射
$f^{-1}R^qg_*\mathcal{F} \to R^qh_*e^{-1}\mathcal{F}$
是同构，见 Theorem \ref{etale-cohomology-theorem-smooth-base-change}（由 $n$ 的选择，
关于挠性的假设成立）。因此 Lemma \ref{etale-cohomology-lemma-base-change-f-star-general-stalks}
给出下式中间的等号：
$$
H^q(F_{\overline{x}, \overline{t}}, \underline{M}) =
H^q(\Spec(\mathcal{O}^{sh}_{X, \overline{x}}) \times_S \overline{t},
\underline{M})
=
H^q(\Spec(\mathcal{O}^{sh}_{S, \overline{s}}) \times_S \overline{t},
\underline{M}) =
H^q(\overline{t}, \underline{M})
$$
外侧的两个等号用到了
$S = \Spec(\mathcal{O}^{sh}_{S, \overline{s}})$。
由于 $\overline{t}$ 是可分闭域的谱，故有
$$
H^q(F_{\overline{x}, \overline{t}}, \underline{M}) =
\left\{
\begin{matrix}
M & \text{若} & q = 0 \\
0 & \text{若} & q \not = 0
\end{matrix}
\right.
$$
这正是所要证明的结论（见 Definition \ref{etale-cohomology-definition-locally-acyclic}
后的讨论）。
\end{proof}

\begin{lemma}
\label{etale-cohomology-lemma-locally-acyclic-locally-constant}
设 $f : X \to S$ 是概形态射。设 $\mathcal{F}$ 是 $X_\etale$ 上的
局部常值阿贝尔层，使得每当 $\overline{x}$ 是 $X$ 的几何点时，阿贝尔群
$\mathcal{F}_{\overline{x}}$ 都是挠群，且每个元素的阶均与
$\overline{x}$ 的剩余域特征互素。如果 $f$ 局部无环，则 $f$
相对于 $\mathcal{F}$ 局部无环。
\end{lemma}

\begin{proof}
事实上，设 $\overline{x}$ 是 $X$ 的几何点。由于 $\mathcal{F}$ 局部常值，
故 $\mathcal{F}$ 在 $\Spec(\mathcal{O}^{sh}_{X, \overline{x}})$ 上的限制
同构于常值层 $\underline{M}$，其中 $M = \mathcal{F}_{\overline{x}}$。
由假设，可把 $M = \colim M_i$ 写成有限阿贝尔群 $M_i$ 的滤过余极限，
并且它们的阶均与 $\overline{x}$ 的剩余域特征互素。取一个几何点
$\overline{t}$（它是 $\Spec(\mathcal{O}^{sh}_{S, f(\overline{x})})$ 的几何点）。
由于 $F_{\overline{x}, \overline{t}}$ 是仿射的，
由 Lemma \ref{etale-cohomology-lemma-colimit} 有
$$
H^q(F_{\overline{x}, \overline{t}}, \underline{M}) =
\colim H^q(F_{\overline{x}, \overline{t}}, \underline{M_i})
$$
对每个 $i$，可把 $M_i = \bigoplus \mathbf{Z}/n_{i, j}\mathbf{Z}$
写成有限直和，其中整数 $n_{i, j}$ 均与 $\overline{x}$ 的剩余域特征互素。
由于 $f$ 局部无环，故
$$
H^q(F_{\overline{x}, \overline{t}},
\underline{\mathbf{Z}/n_{i, j}\mathbf{Z}}) =
\left\{
\begin{matrix}
\mathbf{Z}/n_{i, j}\mathbf{Z} & \text{若} & q = 0 \\
0 & \text{若} & q \not = 0
\end{matrix}
\right.
$$
见 Definition \ref{etale-cohomology-definition-locally-acyclic} 后的讨论。
取直和与余极限可得
$$
H^q(F_{\overline{x}, \overline{t}}, \underline{M}) =
\left\{
\begin{matrix}
M & \text{若} & q = 0 \\
0 & \text{若} & q \not = 0
\end{matrix}
\right.
$$
结论得证。
\end{proof}

\begin{lemma}
\label{etale-cohomology-lemma-locally-acyclic-quasi-finite-base-change}
设
$$
\xymatrix{
X' \ar[r]_{g'} \ar[d]_{f'} & X \ar[d]^f \\
S' \ar[r]^g & S
}
$$
是概形的笛卡儿图。设 $K$ 是 $D(X_\etale)$ 的对象。
设 $\overline{x}'$ 是 $X'$ 的几何点，其像 $\overline{x}$ 是 $X$ 的几何点。如果
\begin{enumerate}
\item $f$ 在 $\overline{x}$ 处相对于 $K$ 局部无环；且
\item $g$ 局部拟有限，或者 $S' = \lim S_i$ 是在 $S$ 上局部拟有限、
过渡态射为仿射态射的概形的有向逆极限，或者 $g : S' \to S$ 是整态射，
\end{enumerate}
则 $f'$ 在 $\overline{x}'$ 处相对于 $(g')^{-1}K$ 局部无环。
\end{lemma}

\begin{proof}
以 $\overline{s}'$ 和 $\overline{s}$ 分别表示 $\overline{x}'$ 和
$\overline{x}$ 在 $S'$ 和 $S$ 中的像。设 $\overline{t}'$ 是
$\Spec(\mathcal{O}^{sh}_{S', \overline{s}'})$ 的几何点，并以
$\overline{t}$ 表示它在 $\Spec(\mathcal{O}^{sh}_{S, \overline{s}})$ 中的像。
由 Algebra, Lemma
\ref{algebra-lemma-base-change-strict-henselization-quasi-finite}
以及对 $g$ 的假设，下列同态是同构：
$$
\mathcal{O}^{sh}_{X, \overline{x}}
\otimes_{\mathcal{O}^{sh}_{S, \overline{s}}}
\mathcal{O}^{sh}_{S', \overline{s}'}
\longrightarrow
\mathcal{O}^{sh}_{X', \overline{x}'}
$$
由于按照约定
$\kappa(\overline{t}) = \kappa(\overline{t}')$，可得
$$
F_{\overline{x}', \overline{t}'} =
\Spec\left(
\mathcal{O}^{sh}_{X', \overline{x}'}
\otimes_{\mathcal{O}^{sh}_{S', \overline{s}'}}
\kappa(\overline{t}')\right) =
\Spec\left(
\mathcal{O}^{sh}_{X, \overline{x}}
\otimes_{\mathcal{O}^{sh}_{S, \overline{s}}}
\kappa(\overline{t})\right) =
F_{\overline{x}, \overline{t}}
$$
换言之，$f'$ 在 $\overline{x}'$ 处的消失圈簇是
$f$ 在 $\overline{x}$ 处的消失圈簇的实例。本引理由此和定义立即得到。
\end{proof}










\section{反特殊化映射}
\label{etale-cohomology-section-cospecialization}

\noindent
设 $f : X \to S$ 是概形态射。设 $\overline{x}$ 是 $X$ 的几何点，
其像 $\overline{s} = f(\overline{x})$ 是 $S$ 的点。设 $\overline{t}$ 是
$\Spec(\mathcal{O}^{sh}_{S, \overline{s}})$ 的几何点。
设 $K \in D(X_\etale)$。对任意概形态射 $g : Y \to X$，
我们以 $K|_Y$ 记 $g_{small}^{-1}K$，并以 $R\Gamma(Y, K)$
记 $R\Gamma(Y_\etale, g_{small}^{-1}K)$。我们断言，如果
\begin{enumerate}
\item $K$ 下有界，即 $K \in D^+(X_\etale)$；
\item $f$ 相对于 $K$ 局部无环，
\end{enumerate}
那么存在一个{\it 反特殊化映射}
$$
cosp :
R\Gamma(X_{\overline{t}}, K)
\longrightarrow
R\Gamma(X_{\overline{s}}, K)
$$
它与 Section \ref{etale-cohomology-section-specialization}，尤其是
Remark \ref{etale-cohomology-remark-specialization-map-and-fibres} 中考察的特殊化映射密切相关。

\medskip\noindent
为构造该映射，考察态射
$$
X_{\overline{t}}
\xrightarrow{h}
X \times_S \Spec(\mathcal{O}^{sh}_{S, \overline{s}})
\xleftarrow{i}
X_{\overline{s}}
$$
$h^{-1}$ 与 $Rh_*$ 之间伴随的单位给出映射
$$
\beta_{K, \overline{s}, \overline{t}} :
K|_{X \times_S \Spec(\mathcal{O}^{sh}_{S, \overline{s}})}
\longrightarrow
Rh_*(K|_{X_{\overline{t}}})
$$
它位于 $D((X \times_S \Spec(\mathcal{O}^{sh}_{S, \overline{s}}))_\etale)$ 中。
下面的 Lemma \ref{etale-cohomology-lemma-beta-is-isomorphism} 表明，在上述假设下，拉回
$i^{-1}\beta_{K, \overline{s}, \overline{t}}$ 是同构。
因此可把反特殊化映射定义为复合
\begin{align*}
R\Gamma(X_{\overline{t}}, K)
& =
R\Gamma(X \times_S \Spec(\mathcal{O}^{sh}_{S, \overline{s}}),
Rh_*(K|_{X_{\overline{t}}})) \\
&
\xrightarrow{i^{-1}}
R\Gamma(X_{\overline{s}}, i^{-1}Rh_*(K|_{X_{\overline{t}}})) \\
&
\xrightarrow{(i^{-1}\beta_{K, \overline{s}, \overline{t}})^{-1}}
R\Gamma(X_{\overline{s}},
i^{-1}(K|_{X \times_S \Spec(\mathcal{O}^{sh}_{S, \overline{s}})})) \\
& =
R\Gamma(X_{\overline{s}}, K)
\end{align*}

\begin{lemma}
\label{etale-cohomology-lemma-beta-is-isomorphism}
映射 $i^{-1}\beta_{K, \overline{s}, \overline{t}}$ 是同构。
\end{lemma}

\begin{proof}
映射 $h$、$i$、$\beta_{K, \overline{s}, \overline{t}}$ 的构造只依赖于
$X$ 和 $K$ 到 $\Spec(\mathcal{O}^{sh}_{S, \overline{s}})$ 的基变换。
因此可且确实假设 $S$ 是闭点为 $\overline{s}$ 的严格 Hensel 概形。
注意，$f$ 相对于 $K$ 的局部无环性在此基变换下保持（例如可用
Lemma \ref{etale-cohomology-lemma-locally-acyclic-quasi-finite-base-change}，也可在这个非常特殊的情形中
直接比较严格 Hensel 环）。

\medskip\noindent
设 $\overline{x}$ 是 $X_{\overline{s}}$ 的几何点。等价地，设
$\overline{x}$ 是在 $f$ 下的像为 $\overline{s}$ 的几何点。下面计算
$i^{-1}\beta_{K, \overline{s}, \overline{t}}$ 在 $\overline{x}$ 处的茎。
首先有
$$
(i^{-1}\beta_{K, \overline{s}, \overline{t}})_{\overline{x}} =
(\beta_{K, \overline{s}, \overline{t}})_{\overline{x}}
$$
因为拉回保持茎，见 Lemma \ref{etale-cohomology-lemma-stalk-pullback}。
由于处在 $S = \Spec(\mathcal{O}^{sh}_{S, \overline{s}})$ 的情形，
故 $h : X_{\overline{t}} \to X$ 满足
$X_{\overline{t}} \times_X \Spec(\mathcal{O}^{sh}_{X, \overline{x}})
= F_{\overline{x}, \overline{t}}$。因此
$$
(\beta_{K, \overline{s}, \overline{t}})_{\overline{x}} :
K_{\overline{x}}
\longrightarrow
Rh_*(K|_{X_{\overline{t}}})_{\overline{x}} =
R\Gamma(F_{\overline{x}, \overline{t}}, K)
$$
其中等号来自 Theorem \ref{etale-cohomology-theorem-higher-direct-images}。
于是映射 $(\beta_{K, \overline{s}, \overline{t}})_{\overline{x}}$
正是 Definition \ref{etale-cohomology-definition-locally-acyclic} 中使用的映射
$\alpha_{K, \overline{x}, \overline{t}}$。由 Theorem \ref{etale-cohomology-theorem-exactness-stalks}，
可在茎上检验一个映射是否为同构，故结论成立。
\end{proof}

\noindent
反特殊化映射一旦存在，便意在成为特殊化映射的逆。

\begin{lemma}
\label{etale-cohomology-lemma-specialization-cospecialization}
在上述情形中，如果还假设 $f$ 拟紧且拟分离，则图
$$
\xymatrix{
(Rf_*K)_{\overline{s}} \ar[r] \ar[d]_{sp} &
R\Gamma(X_{\overline{s}}, K) \\
(Rf_*K)_{\overline{t}} \ar[r] &
R\Gamma(X_{\overline{t}}, K) \ar[u]_{cosp}
}
$$
交换。
\end{lemma}

\begin{proof}
与 Lemma \ref{etale-cohomology-lemma-beta-is-isomorphism} 的证明相同，
可把 $S$ 替换为 $\Spec(\mathcal{O}^{sh}_{S, \overline{s}})$。
这时映射简化为 $h : X_{\overline{t}} \to X$、
$i : X_{\overline{s}} \to X$ 以及
$\beta_{K, \overline{s}, \overline{t}} : K \to Rh_*(K|_{X_{\overline{t}}})$。
由 Theorem \ref{etale-cohomology-theorem-higher-direct-images}，
$(Rf_*K)_{\overline{s}} = R\Gamma(X, K)$，故 $sp$ 与基变换映射
$(Rf_*K)_{\overline{t}} \to R\Gamma(X_{\overline{t}}, K)$
的复合就是沿 $h$ 拉回上同调。这等于映射
$$
R\Gamma(X, K) \xrightarrow{\beta_{K, \overline{s}, \overline{t}}}
R\Gamma(X, Rh_*(K|_{X_{\overline{t}}})) =
R\Gamma(X_{\overline{t}}, K)
$$
而映射 $cosp$ 先把此陈列公式中的等号 $=$ 反向使用，
再沿 $i$ 拉回，最后应用 $i^{-1}\beta_{K, \overline{s}, \overline{t}}$ 的逆。
因此得到所需的交换性。
\end{proof}

\begin{lemma}
\label{etale-cohomology-lemma-sp-isom-proper-torsion-loc-ac}
设 $f : X \to S$ 是概形态射。设 $K \in D(X_\etale)$。假设
\begin{enumerate}
\item $K$ 下有界，即 $K \in D^+(X_\etale)$；
\item $f$ 相对于 $K$ 局部无环；
\item $f$ 是固有态射；且
\item $K$ 的上同调层是挠层。
\end{enumerate}
则对每个几何点 $\overline{s}$（它是 $S$ 的几何点）以及每个几何点
$\overline{t}$（它是 $\Spec(\mathcal{O}^{sh}_{S, \overline{s}})$ 的几何点），
特殊化映射
$sp : (Rf_*K)_{\overline{s}} \to (Rf_*K)_{\overline{t}}$
与反特殊化映射
$cosp : R\Gamma(X_{\overline{t}}, K) \to R\Gamma(X_{\overline{s}}, K)$
都是同构。
\end{lemma}

\begin{proof}
由固有基变换定理（采用 Lemma \ref{etale-cohomology-lemma-proper-base-change-stalk} 的形式），有
$(Rf_*K)_{\overline{s}} = R\Gamma(X_{\overline{s}}, K)$，
对 $\overline{t}$ 也类似。“正确”的证明应说明
Lemma \ref{etale-cohomology-lemma-specialization-cospecialization} 中的论证表明，
此时 $sp$ 与 $cosp$ 是互逆同构。这里改为直接证明 $cosp$ 是同构。
由上面的讨论，$cosp$ 是同构当且仅当沿 $i$ 的拉回
$$
R\Gamma(X \times_S \Spec(\mathcal{O}^{sh}_{S, \overline{s}}),
Rh_*(K|_{X_{\overline{t}}}))
\longrightarrow
R\Gamma(X_{\overline{s}}, i^{-1}Rh_*(K|_{X_{\overline{t}}}))
$$
在 $D^+(\textit{Ab})$ 中是同构。这由固有态射
$f' : X \times_S \Spec(\mathcal{O}^{sh}_{S, \overline{s}}) \to
\Spec(\mathcal{O}^{sh}_{S, \overline{s}})$
沿态射 $\overline{s} \to \Spec(\mathcal{O}^{sh}_{S, \overline{s}})$
对复形 $K' = Rh_*(K|_{X_{\overline{t}}})$ 应用固有基变换定理得到。
由 Lemma \ref{etale-cohomology-lemma-torsion-direct-image}，复形 $K'$ 下有界且有挠上同调层。
由于 $\Spec(\mathcal{O}^{sh}_{S, \overline{s}})$ 严格 Hensel，
且 $\overline{s}$ 位于闭点之上，故上述箭头的源等于
$(Rf'_*K')_{\overline{s}}$，靶等于
$R\Gamma(X_{\overline{s}}, K')$，并且由已经使用过的
Lemma \ref{etale-cohomology-lemma-proper-base-change-stalk}，上述映射是同构。
于是 Lemma \ref{etale-cohomology-lemma-specialization-cospecialization} 的图中四个箭头有三个是同构，
结论随之成立。
\end{proof}

\begin{lemma}
\label{etale-cohomology-lemma-sp-isom-proper-loc-cst-torsion}
设 $f : X \to S$ 是概形态射。设 $\mathcal{F}$ 是 $X_\etale$ 上的阿贝尔层。假设
\begin{enumerate}
\item $f$ 光滑且固有；
\item $\mathcal{F}$ 局部常值；且
\item $\mathcal{F}_{\overline{x}}$ 是挠群，且每个元素的阶均与
$\overline{x}$ 的剩余特征互素，其中 $\overline{x}$ 取遍 $X$ 的所有几何点。
\end{enumerate}
则对每个几何点 $\overline{s}$（它是 $S$ 的几何点）以及每个几何点
$\overline{t}$（它是 $\Spec(\mathcal{O}^{sh}_{S, \overline{s}})$ 的几何点），
特殊化映射
$sp : (Rf_*\mathcal{F})_{\overline{s}} \to (Rf_*\mathcal{F})_{\overline{t}}$
是同构。
\end{lemma}

\begin{proof}
这由 Lemmas \ref{etale-cohomology-lemma-sp-isom-proper-torsion-loc-ac}、
\ref{etale-cohomology-lemma-locally-acyclic-locally-constant} 以及
Proposition \ref{etale-cohomology-proposition-smooth-locally-acyclic} 得到。
\end{proof}

















\section{上同调维数}
\label{etale-cohomology-section-cd}

\noindent
利用 Section \ref{etale-cohomology-section-vanishing-torsion} 中的结果，
再加上对固有基变换定理的一次看似无害的应用（见
Proposition \ref{etale-cohomology-proposition-cd-affine} 的证明），可以推出概形和域的
上同调维数的一些界。

\begin{definition}
\label{etale-cohomology-definition-cd}
设 $X$ 是拟紧拟分离概形。{\it $X$ 的上同调维数}是集合中最小的元素
$$
\text{cd}(X) \in \{0, 1, 2, \ldots\} \cup \{\infty\}
$$
，它满足：对任意阿贝尔挠层 $\mathcal{F}$（定义在 $X_\etale$ 上），
都有 $H^i_\etale(X, \mathcal{F}) = 0$，只要 $i > \text{cd}(X)$。
如果 $X = \Spec(A)$，有时也称它为 $A$ 的上同调维数。
\end{definition}

\noindent
如果概形具有特征 $p$，则对 $p$-幂挠层的上同调消失通常可以得到更尖锐的界。
我们将在别处讨论这一点（此处插入将来的参考文献）。

\begin{lemma}
\label{etale-cohomology-lemma-cd-limit}
设 $X = \lim X_i$ 是一族拟紧拟分离概形的有向逆极限，且过渡态射均为仿射态射。
则 $\text{cd}(X) \leq \max \text{cd}(X_i)$。
\end{lemma}

\begin{proof}
以 $f_i : X \to X_i$ 表示投影。设 $\mathcal{F}$ 是 $X_\etale$ 上的阿贝尔挠层。
由 Lemma \ref{etale-cohomology-lemma-linus-hamann}，有
$\mathcal{F} = \lim f_i^{-1}f_{i, *}\mathcal{F}$。
因此由 Theorem \ref{etale-cohomology-theorem-colimit}，
$H^q_\etale(X, \mathcal{F}) =
\colim H^q_\etale(X_i, f_{i, *}\mathcal{F})$。引理得证。
\end{proof}

\begin{lemma}
\label{etale-cohomology-lemma-cd-curve-over-field}
设 $K$ 是域。设 $X$ 是 $1$ 维仿射概形，且在 $K$ 上有限型。则
$\text{cd}(X) \leq 1 + \text{cd}(K)$。
\end{lemma}

\begin{proof}
设 $\mathcal{F}$ 是 $X_\etale$ 上的阿贝尔挠层。考察态射
$f : X \to \Spec(K)$ 的 Leray 谱序列，得到
$$
E_2^{p, q} = H^p(\Spec(K), R^qf_*\mathcal{F})
$$
，它收敛到 $H^{p + q}_\etale(X, \mathcal{F})$。
$R^qf_*\mathcal{F}$ 在几何点
$\Spec(\overline{K}) \to \Spec(K)$ 处的茎，是 $\mathcal{F}$ 到
$X_{\overline{K}}$ 的拉回之上同调。因此由
Theorem \ref{etale-cohomology-theorem-vanishing-affine-curves}，它在次数 $\geq 2$ 中消失。
\end{proof}

\begin{lemma}
\label{etale-cohomology-lemma-cd-field-extension}
设 $L/K$ 是域扩张。则
$\text{cd}(L) \leq \text{cd}(K) + \text{trdeg}_K(L)$。
\end{lemma}

\begin{proof}
如果 $\text{trdeg}_K(L) = \infty$，结论显然。否则可找到扩张序列
$L= L_r/L_{r - 1}/ \ldots /L_1/L_0 = K$，使得
$\text{trdeg}_{L_i}(L_{i + 1}) = 1$ 且 $r = \text{trdeg}_K(L)$。
因此只须证明 $r = 1$ 的情形。此时可把 $L = \colim A_i$
写成其有限型 $K$-子代数的滤过余极限。由 Lemma \ref{etale-cohomology-lemma-cd-limit}，
只须证明 $\text{cd}(A_i) \leq 1 + \text{cd}(K)$，这由
Lemma \ref{etale-cohomology-lemma-cd-curve-over-field} 得到。
\end{proof}

\begin{lemma}
\label{etale-cohomology-lemma-strictly-henselian}
设 $K$ 是域。设 $X$ 是 $K$ 上的有限型概形。设 $x \in X$。
置 $a = \text{trdeg}_K(\kappa(x))$ 且 $d = \dim_x(X)$。则存在映射
$$
K(t_1, \ldots, t_a)^{sep} \longrightarrow \mathcal{O}_{X, x}^{sh}
$$
使得
\begin{enumerate}
\item $\mathcal{O}_{X, x}^{sh}$ 的剩余域是
$K(t_1, \ldots, t_a)^{sep}$ 的纯不可分扩张；
\item $\mathcal{O}_{X, x}^{sh}$ 是一族有限型
$K(t_1, \ldots, t_a)^{sep}$-代数的滤过余极限，且这些代数的维数均为
$\leq d - a$。
\end{enumerate}
\end{lemma}

\begin{proof}
可假设 $X$ 是仿射的。由 Noether 正规化，在必要时再次缩小 $X$ 后，
可选取有限态射 $\pi : X \to \mathbf{A}^d_K$，见
Algebra, Lemma \ref{algebra-lemma-Noether-normalization-at-point}。
由于 $\kappa(x)$ 是 $\pi(x)$ 的剩余域的有限扩张，该剩余域的超越次数也为
$a$（相对于 $K$）。因此可找到有限态射
$\pi' : \mathbf{A}^d_K \to \mathbf{A}^d_K$，使得 $\pi'(\pi(x))$
对应于线性子空间 $\mathbf{A}^a_K \subset \mathbf{A}^d_K$ 的一般点；
该子空间由令最后 $d - a$ 个坐标为零给出。因此复合
$$
X \xrightarrow{\pi' \circ \pi} \mathbf{A}^d_K \xrightarrow{p} \mathbf{A}^a_K
$$
由 $\pi' \circ \pi$ 和投影 $p$ 组成，后者投到前 $a$ 个坐标；它把 $x$
映到一般点 $\eta \in \mathbf{A}^a_K$。所诱导的 \'etale 局部环映射
$$
K(t_1, \ldots, t_a)^{sep} =
\mathcal{O}_{\mathbf{A}^a_k, \eta}^{sh}
\longrightarrow \mathcal{O}_{X, x}^{sh}
$$
满足 (1)，因为 $\mathcal{O}_{X, x}^{sh}$ 的剩余域显然是可分闭域
$K(t_1, \ldots, t_a)^{sep}$ 的代数扩张。另一方面，如果 $X = \Spec(B)$，
则 $\mathcal{O}_{X, x}^{sh} = \colim B_j$ 是 \'etale $B$-代数 $B_j$ 的滤过余极限。
注意，$B_j$ 在 $K[t_1, \ldots, t_d]$ 上拟有限，因为 $B$ 在
$K[t_1, \ldots, t_d]$ 上有限。类似地，可把
$K(t_1, \ldots, t_a)^{sep} = \colim A_i$ 写成 \'etale
$K[t_1, \ldots, t_a]$-代数的滤过余极限。对每个 $i$，可找到 $j$，使得
$A_i \to K(t_1, \ldots, t_a)^{sep} \to \mathcal{O}_{X, x}^{sh}$
分解经过映射 $\psi_{i, j} : A_i \to B_j$。于是 $B_j$ 在
$A_i[t_{a + 1}, \ldots, t_d]$ 上拟有限。因此
$$
B_{i, j} = B_j \otimes_{\psi_{i, j}, A_i} K(t_1, \ldots, t_a)^{sep}
$$
的维数为 $\leq d - a$，因为它在
$K(t_1, \ldots, t_a)^{sep}[t_{a + 1}, \ldots, t_d]$ 上拟有限。
现在 (2) 得证，因为 $\mathcal{O}_{X, x}^{sh}$ 是滤过余极限\footnote{设
$R$ 是环。设 $A = \colim_{i \in I} A_i$ 是有限呈现
$R$-代数的滤过余极限。设 $B = \colim_{j \in J} B_j$ 是
$R$-代数的滤过余极限。设 $A \to B$ 是 $R$-代数映射。假设对所有
$i \in I$，存在 $j \in J$ 以及 $R$-代数映射
$\psi_{i, j} : A_i \to B_j$。若
$(i', j', \psi_{i', j'}) \geq (i, j, \psi_{i, j})$，意思是
$i' \geq i$、$j' \geq j$，并且 $\psi_{i, j}$ 与
$\psi_{i', j'}$ 相容，则这些三元组的集合构成有向集，并且
$B = \colim B_j \otimes_{\psi_{i, j} A_i} A$。}，其组成代数为
$B_{i, j}$。略去若干细节。
\end{proof}

\begin{proposition}
\label{etale-cohomology-proposition-cd-affine}
设 $K$ 是域。设 $X$ 是 $K$ 上有限型的仿射概形。则有
$\text{cd}(X) \leq \dim(X) + \text{cd}(K)$。
\end{proposition}

\begin{proof}
对 $\dim(X)$ 作归纳。设 $\mathcal{F}$ 是 $X_\etale$ 上的阿贝尔挠层。

\medskip\noindent
情形 $\dim(X) = 0$。此时结构态射 $f : X \to \Spec(K)$ 是有限的。
因此 $R^if_*\mathcal{F} = 0$，只要 $i > 0$，见
Proposition \ref{etale-cohomology-proposition-finite-higher-direct-image-zero}。
因此 $H^i_\etale(X, \mathcal{F}) = H^i_\etale(\Spec(K), f_*\mathcal{F})$；
这是由 $f$ 的 Leray 谱序列（Cohomology on Sites, Lemma
\ref{sites-cohomology-lemma-Leray}）得到的，
结论显然。

\medskip\noindent
情形 $\dim(X) = 1$。这就是 Lemma \ref{etale-cohomology-lemma-cd-curve-over-field}。

\medskip\noindent
假设 $d = \dim(X) > 1$，并且命题对域上维数 $< d$ 的有限型仿射概形成立。
由 Noether 正规化（例如见 Varieties, Lemma
\ref{varieties-lemma-noether-normalization-affine}），存在有限态射
$f : X \to \mathbf{A}^d_K$。回顾由
Proposition \ref{etale-cohomology-proposition-finite-higher-direct-image-zero}，
$R^if_*\mathcal{F} = 0$，只要 $i > 0$。由 $f$ 的 Leray 谱序列（Cohomology on Sites,
Lemma \ref{sites-cohomology-lemma-Leray}），只须对 $\pi_*\mathcal{F}$
在 $\mathbf{A}^d_K$ 上证明结论。

\medskip\noindent
插曲 I。设 $j : X \to Y$ 是光滑 $d$ 维簇的开浸入，这些簇定义在
$K$ 上（不要求仿射），其补集是某个有效 Cartier 除子 $D$ 的支集。
层 $R^qj_*\mathcal{F}$（其中 $q > 0$）支撑在 $D$ 上。我们断言
$(R^qj_*\mathcal{F})_{\overline{y}} = 0$，只要
$a = \text{trdeg}_K(\kappa(y)) > d - q$。事实上，由
Theorem \ref{etale-cohomology-theorem-higher-direct-images} 有
$$
(R^qj_*\mathcal{F})_{\overline{y}} =
H^q(\Spec(\mathcal{O}_{Y, y}^{sh}) \times_Y X, \mathcal{F})
$$
选取 $f \in \mathfrak m_y \subset \mathcal{O}_{Y, y}$ 作为 $D$ 的局部方程。
于是
$$
\Spec(\mathcal{O}_{Y, y}^{sh}) \times_Y X =
\Spec(\mathcal{O}_{Y, y}^{sh}[1/f])
$$
利用 Lemma \ref{etale-cohomology-lemma-strictly-henselian} 得到嵌入
$$
K(t_1, \ldots, t_a)^{sep}(x) =
K(t_1, \ldots, t_a)^{sep}[x]_{(x)}[1/x]
\longrightarrow
\mathcal{O}_{Y, y}^{sh}[1/f]
$$
由于在 $K$ 上，$\mathcal{O}_{Y, y}^{sh}$ 的分式域的超越次数为 $d$，
可见 $\mathcal{O}_{Y, y}^{sh}[1/f]$ 是一族有限型
$(d - a - 1)$ 维代数的滤过余极限；这些代数定义在域
$K(t_1, \ldots, t_a)^{sep}(x)$ 上，而该域由
Lemma \ref{etale-cohomology-lemma-cd-field-extension} 具有上同调维数 $1$。
因此由归纳假设和 Lemma \ref{etale-cohomology-lemma-cd-limit} 得到所需消失。

\medskip\noindent
插曲 II。设 $Z$ 是定义在 $K$ 上的光滑簇，其维数为 $d - 1$。设
$E_a \subset Z$ 是由满足条件的点 $z \in Z$ 构成的集合，该条件为
$\text{trdeg}_K(\kappa(z)) \leq a$。注意 $E_a$ 对特殊化封闭，见
Varieties, Lemma \ref{varieties-lemma-dimension-locally-algebraic}。
设 $\mathcal{G}$ 是 $Z$ 上的阿贝尔挠层，其支集包含在 $E_a$ 中。
我们断言 $H^b_\etale(Z, \mathcal{G}) = 0$，只要
$b > a + \text{cd}(K)$。事实上，可写
$\mathcal{G} = \colim \mathcal{G}_i$，其中 $\mathcal{G}_i$ 是支撑在闭子概形
$Z_i$ 上的阿贝尔挠层，该闭子概形包含在 $E_a$ 中，见
Lemma \ref{etale-cohomology-lemma-support-in-subset}。于是归纳假设给出
$\mathcal{G}_i$ 的所需消失\footnote{这里先用 Proposition
\ref{etale-cohomology-proposition-closed-immersion-pushforward}，把 $\mathcal{G}_i$ 写成
$Z_i$ 上某层的正像；归纳假设给出该层在 $Z_i$ 上的消失，而
$Z_i \to Z$ 的 Leray 谱序列给出 $\mathcal{G}_i$ 的消失。}。
最后由 Theorem \ref{etale-cohomology-theorem-colimit} 得到结论。

\medskip\noindent
考察交换图
$$
\xymatrix{
\mathbf{A}^d_K \ar[rd]_f \ar[rr]_-j & &
\mathbf{P}^1_K \times_K \mathbf{A}^{d - 1}_K \ar[ld]^g \\
& \mathbf{A}^{d - 1}_K
}
$$
注意 $j$ 是光滑 $d$ 维簇的开浸入，其补集为有效 Cartier 除子 $D$。
因此可以使用插曲 I 中得到的结果。下面考察相对 Leray 谱序列
$$
E_2^{p, q} = R^pg_*R^qj_*\mathcal{F} \Rightarrow R^{p + q}f_*\mathcal{F}
$$
层 $R^qj_*\mathcal{F}$ 在 $q > 0$ 时支撑在 $D$ 上；又因
$g|_D : D \to \mathbf{A}^{d - 1}_K$ 是同构，故
$R^pg_*R^qj_*\mathcal{F} = 0$，只要 $p > 0$ 且 $q > 0$。此外，
$R^qj_*\mathcal{F} = 0$，只要
$q > d$。另一方面，$g$ 是相对维数为 $1$ 的固有态射。
因此由 Lemma \ref{etale-cohomology-lemma-cohomological-dimension-proper}，
$R^pg_*j_*\mathcal{F} = 0$，只要 $p > 2$。于是谱序列的 $E_2$-页如下：
$$
\begin{matrix}
g_*R^dj_*\mathcal{F} & 0 & 0 \\
\ldots & \ldots & \ldots \\
g_*R^2j_*\mathcal{F} & 0 & 0 \\
g_*R^1j_*\mathcal{F} & 0 & 0 \\
g_*j_*\mathcal{F} & R^1g_*j_*\mathcal{F} & R^2g_*j_*\mathcal{F}
\end{matrix}
$$
因此 $R^qf_*\mathcal{F} = g_*R^qj_*\mathcal{F}$，只要 $q > 2$。
由插曲 I，$R^qf_*\mathcal{F}$ 在 $q > 2$ 时的支集包含在
$\mathbf{A}^{d - 1}_K$ 的如下点集之中：其剩余域的超越次数为
$\leq d - q$。由插曲 II，
$$
H^p(\mathbf{A}^{d - 1}_K, R^qf_*\mathcal{F}) = 0
\text{ 当 }p > d - q + \text{cd}(K)\text{ 且 }q > 2
$$
另一方面，由 Theorem \ref{etale-cohomology-theorem-higher-direct-images} 有
$R^2f_*\mathcal{F}_{\overline{\eta}} =
H^2(\mathbf{A}^1_{\overline{\eta}}, \mathcal{F}) = 0$
（由维数 $1$ 的情形消失），其中 $\eta$ 是 $\mathbf{A}^{d - 1}_K$ 的一般点。
因此再次由插曲 II，
$$
H^p(\mathbf{A}^{d - 1}_K, R^2f_*\mathcal{F}) = 0
\text{ 对 }p > d - 2 + \text{cd}(K)
$$
最后，由归纳假设有
$$
H^p(\mathbf{A}^{d - 1}_K, R^qf_*\mathcal{F}) = 0
\text{ 当 }p > d - 1 + \text{cd}(K)\text{ 且 }q = 0, 1
$$
。把以上结论与 Leray 谱序列
$H^p(\mathbf{A}^{d - 1}_K, R^qf_*\mathcal{F}) \Rightarrow
H^{p + q}(\mathbf{A}^d_K, \mathcal{F})$ 结合起来即得结论。
\end{proof}

\begin{lemma}
\label{etale-cohomology-lemma-interlude-II}
设 $K$ 是域。设 $X$ 是在 $K$ 上有限型的仿射概形。设
$E_a \subset X$ 是由点 $x \in X$ 构成的集合，这些点满足
$\text{trdeg}_K(\kappa(x)) \leq a$。设 $\mathcal{F}$ 是定义在
$X_\etale$ 上的阿贝尔挠层，其支集包含在 $E_a$ 中。则
$H^b_\etale(X, \mathcal{F}) = 0$，只要 $b > a + \text{cd}(K)$。
\end{lemma}

\begin{proof}
可写 $\mathcal{F} = \colim \mathcal{F}_i$，其中 $\mathcal{F}_i$ 是支撑在
闭子概形 $Z_i$ 上的阿贝尔挠层，该闭子概形包含在 $E_a$ 中，见
Lemma \ref{etale-cohomology-lemma-support-in-subset}。于是 Proposition
\ref{etale-cohomology-proposition-cd-affine} 给出 $\mathcal{F}_i$ 的所需消失。
略去细节；提示：先用 Proposition \ref{etale-cohomology-proposition-closed-immersion-pushforward}
把 $\mathcal{F}_i$ 写成 $Z_i$ 上某层的正像，再使用该层在 $Z_i$ 上的消失，
并用 $Z_i \to Z$ 的 Leray 谱序列得到 $\mathcal{F}_i$ 的消失。
最后由 Theorem \ref{etale-cohomology-theorem-colimit} 得到结论。
\end{proof}

\begin{lemma}
\label{etale-cohomology-lemma-interlude-I}
设 $f : X \to Y$ 是定义在域 $K$ 上的有限型概形之间的仿射态射。设
$E_a(X)$ 是由点 $x \in X$ 构成的集合，这些点满足
$\text{trdeg}_K(\kappa(x)) \leq a$。
设 $\mathcal{F}$ 是 $X_\etale$ 上的阿贝尔挠层，其支集包含在 $E_a$ 中。
则 $R^qf_*\mathcal{F}$ 的支集包含在 $E_{a - q}(Y)$ 中。
\end{lemma}

\begin{proof}
问题在 $Y$ 上是局部的，故可假设 $Y$ 仿射。于是 $X$ 也仿射，并可选取图
$$
\xymatrix{
X \ar[d]_f \ar[r]_i & \mathbf{A}^{n + m}_K \ar[d]^{\text{pr}} \\
Y \ar[r]^j & \mathbf{A}^n_K
}
$$
，其中水平箭头是闭浸入，右侧竖直箭头是投影（略去细节）。
有 $j_*R^qf_*\mathcal{F} = R^q\text{pr}_*i_*\mathcal{F}$；
这是由 $i$ 和 $j$ 的高阶正像消失得到的，见
Proposition \ref{etale-cohomology-proposition-finite-higher-direct-image-zero}。
此外，命题中对 $j_*$ 的茎的描述表明，只须证明
$j_*R^qf_*\mathcal{F}$ 的消失。因此可假设 $f$ 是投影态射
$\text{pr} : \mathbf{A}^{n + m}_K \to \mathbf{A}^n_K$，并设
$\mathcal{F}$ 是 $\mathbf{A}^{n + m}_K$ 上满足本引理假设的阿贝尔挠层。

\medskip\noindent
设 $y$ 是 $\mathbf{A}^n_K$ 的点。由 Theorem
\ref{etale-cohomology-theorem-higher-direct-images} 有
$$
(R^q\text{pr}_*\mathcal{F})_{\overline{y}} =
H^q(\mathbf{A}^{n + m}_K \times_{A^n_K} \Spec(\mathcal{O}_{Y, y}^{sh}),
\mathcal{F}) =
H^q(\mathbf{A}^m_{\mathcal{O}_{Y, y}^{sh}}, \mathcal{F})
$$
记 $b = \text{trdeg}_K(\kappa(y))$。由 Lemma \ref{etale-cohomology-lemma-strictly-henselian}
得到嵌入
$$
L = K(t_1, \ldots, t_b)^{sep} \longrightarrow \mathcal{O}_{Y, y}^{sh}
$$
把 $\mathcal{O}_{Y, y}^{sh} = \colim B_i$ 写成有限型 $L$-子代数
$B_i \subset \mathcal{O}_{Y, y}^{sh}$ 的滤过余极限；这些子代数包含正则函数环
$K[T_1, \ldots, T_n]$，即 $\mathbf{A}^n_K$ 上的正则函数环。于是
$$
\mathbf{A}^m_{\mathcal{O}_{Y, y}^{sh}} =
\lim \mathbf{A}^m_{B_i}
$$
如果 $z \in \mathbf{A}^m_{B_i}$ 是 $\mathcal{F}$ 的支集中的点，
并且 $x$ 是 $z$ 在 $\mathbf{A}^{m + n}_K$ 中的像，则有
$\text{trdeg}_K(\kappa(x)) \leq a$；这是由引理中对 $\mathcal{F}$ 的假设得到的。
由于 $\mathcal{O}_{Y, y}^{sh}$ 是 $K[T_1, \ldots, T_n]$ 上 \'etale 代数的
滤过余极限，并且 $B_i \subset \mathcal{O}_{Y, y}^{sh}$，可见
$\kappa(z)/\kappa(x)$ 是代数扩张（略去若干细节）。于是
$\text{trdeg}_K(\kappa(z)) \leq a$，从而
$\text{trdeg}_L(\kappa(z)) \leq a - b$。由 Lemma \ref{etale-cohomology-lemma-interlude-II}，
$$
H^q(\mathbf{A}^m_{B_i}, \mathcal{F}) = 0\text{ 对 }q > a - b
$$
因此由 Theorem \ref{etale-cohomology-theorem-colimit}，
$(Rf_*\mathcal{F})_{\overline{y}} = 0$，只要 $q > a - b$，这正是所需结论。
\end{proof}








\section{有限上同调维数}
\label{etale-cohomology-section-finite-cd}

\noindent
我们继续 Section \ref{etale-cohomology-section-cd} 中开始的讨论。

\begin{definition}
\label{etale-cohomology-definition-cd-f}
设 $f : X \to Y$ 是概形之间的拟紧拟分离态射。
{\it $f$ 的上同调维数}是集合中最小的元素
$$
\text{cd}(f) \in \{0, 1, 2, \ldots\} \cup \{\infty\}
$$
，它满足：对任意阿贝尔挠层 $\mathcal{F}$（定义在 $X_\etale$ 上），
都有 $R^if_*\mathcal{F} = 0$，只要 $i > \text{cd}(f)$。
\end{definition}

\begin{lemma}
\label{etale-cohomology-lemma-finite-cd}
设 $K$ 是域。
\begin{enumerate}
\item 如果 $f : X \to Y$ 是 $K$ 上有限型概形之间的态射，
则 $\text{cd}(f) < \infty$。
\item 如果 $\text{cd}(K) < \infty$，则 $\text{cd}(X) < \infty$；
这里 $X$ 是 $K$ 上任意有限型概形。
\end{enumerate}
\end{lemma}

\begin{proof}
证明 (1)。可以假设 $Y$ 是仿射的。我们将使用
Cohomology of Schemes，引理 \ref{coherent-lemma-induction-principle}
的归纳原理。如果 $X$ 也是仿射的，则结论由
Lemma \ref{etale-cohomology-lemma-interlude-I} 成立。因此只须证明：若
$X = U \cup V$，且结论对 $U \to Y$、$V \to Y$ 和
$U \cap V \to Y$ 成立，则它对 $f$ 也成立。这由相对
Mayer--Vietoris 序列得到，见 Lemma \ref{etale-cohomology-lemma-relative-mayer-vietoris}。

\medskip\noindent
证明 (2)。我们将使用
Cohomology of Schemes，引理 \ref{coherent-lemma-induction-principle}
的归纳原理。如果 $X$ 是仿射的，则结论由
Proposition \ref{etale-cohomology-proposition-cd-affine} 成立。因此只须证明：若
$X = U \cup V$，且结论对 $U$、$V$ 和 $U \cap V $ 成立，
则它对 $X$ 也成立。这由 Mayer--Vietoris 序列得到，见
Lemma \ref{etale-cohomology-lemma-mayer-vietoris}。
\end{proof}

\begin{lemma}
\label{etale-cohomology-lemma-finite-cd-mod-n-direct-sums}
上同调与直和。设 $n \geq 1$ 是整数。
\begin{enumerate}
\item 设 $f : X \to Y$ 是拟紧拟分离的概形态射，且
$\text{cd}(f) < \infty$。则函子
$$
Rf_* :
D(X_\etale, \mathbf{Z}/n\mathbf{Z})
\longrightarrow
D(Y_\etale, \mathbf{Z}/n\mathbf{Z})
$$
与直和可交换。
\item 设 $X$ 是拟紧拟分离概形，且
$\text{cd}(X) < \infty$。则函子
$$
R\Gamma(X, -) :
D(X_\etale, \mathbf{Z}/n\mathbf{Z})
\longrightarrow
D(\mathbf{Z}/n\mathbf{Z})
$$
与直和可交换。
\end{enumerate}
\end{lemma}

\begin{proof}
证明 (1)。由于 $\text{cd}(f) < \infty$，可知
$$
f_* :
\textit{Mod}(X_\etale, \mathbf{Z}/n\mathbf{Z})
\longrightarrow
\textit{Mod}(Y_\etale, \mathbf{Z}/n\mathbf{Z})
$$
在 Derived Categories, Lemma
\ref{derived-lemma-unbounded-right-derived} 的意义下具有有限上同调维数。
设 $I$ 是集合，并对 $i \in I$ 令 $E_i$ 为
$D(X_\etale, \mathbf{Z}/n\mathbf{Z})$ 的对象。
选取 K-内射复形 $\mathcal{I}_i^\bullet$；它由
$\mathbf{Z}/n\mathbf{Z}$-模组成，每一项
$\mathcal{I}_i^n$ 都是 $\mathbf{Z}/n\mathbf{Z}$-模的内射层，
而且该复形表示 $E_i$。见 Injectives, Theorem
\ref{injectives-theorem-K-injective-embedding-grothendieck}。
于是 $\bigoplus E_i$ 由复形
$\bigoplus \mathcal{I}_i^\bullet$（逐项直和）表示，见
Injectives, Lemma \ref{injectives-lemma-derived-products}。
由 Lemma \ref{etale-cohomology-lemma-relative-colimit} 有
$$
R^qf_*(\bigoplus \mathcal{I}_i^n) =
\bigoplus R^qf_*(\mathcal{I}_i^n) = 0
$$
，只要 $q > 0$，且对任意 $n$ 都成立。因此由
Derived Categories, Lemma \ref{derived-lemma-unbounded-right-derived}，
可知，可以用复形计算 $Rf_*(\bigoplus E_i)$：
$$
f_*(\bigoplus \mathcal{I}_i^\bullet) =
\bigoplus f_*(\mathcal{I}_i^\bullet)
$$
（等式仍由 Lemma \ref{etale-cohomology-lemma-relative-colimit} 得到）。根据已经使用过的
Injectives, Lemma \ref{injectives-lemma-derived-products}，该复形表示
$\bigoplus Rf_*E_i$。

\medskip\noindent
证明 (2)。它与 (1) 的证明完全相同，故略去。
\end{proof}

\begin{lemma}
\label{etale-cohomology-lemma-proper-mod-n-direct-sums}
设 $f : X \to Y$ 是概形之间的固有态射。设 $n \geq 1$ 是整数。则函子
$$
Rf_* :
D(X_\etale, \mathbf{Z}/n\mathbf{Z})
\longrightarrow
D(Y_\etale, \mathbf{Z}/n\mathbf{Z})
$$
与直和可交换。
\end{lemma}

\begin{proof}
只须证明 $Y$ 拟紧的情形。由
Morphisms，引理 \ref{morphisms-lemma-morphism-finite-type-bounded-dimension}
可知 $f : X \to Y$ 的纤维维数有界。因此
Lemma \ref{etale-cohomology-lemma-cohomological-dimension-proper} 蕴含
$\text{cd}(f) < \infty$。于是结论由
Lemma \ref{etale-cohomology-lemma-finite-cd-mod-n-direct-sums} 得到。
\end{proof}

\begin{lemma}
\label{etale-cohomology-lemma-pull-out-constant-mod-n}
设 $X$ 是拟紧拟分离概形，且 $\text{cd}(X) < \infty$。
设 $\Lambda$ 是挠环。设 $E \in D(X_\etale, \Lambda)$ 且
$K \in D(\Lambda)$。则
$$
R\Gamma(X, E \otimes_\Lambda^\mathbf{L} \underline{K}) =
R\Gamma(X, E) \otimes_\Lambda^\mathbf{L} K
$$
\end{lemma}

\begin{proof}
由 Cohomology on Sites, Section
\ref{sites-cohomology-section-projection-formula}，存在从右向左的典范映射。
令 $T(K)$ 表示引理的陈述对 $K \in D(\Lambda)$ 成立这一性质。
我们将验证 More on Algebra, Remark
\ref{more-algebra-remark-P-resolution} 的条件 (1)、(2) 和 (3) 对
$T$ 成立，从而得到结论。
性质 (1) 成立，因为等式两边都与直和可交换，见
Lemma \ref{etale-cohomology-lemma-finite-cd-mod-n-direct-sums}。
性质 (2) 成立，因为我们在比较三角范畴之间的正合函子，并且可以使用
Derived Categories, Lemma
\ref{derived-lemma-third-isomorphism-triangle}。
性质 (3) 是说，引理对 $K = \Lambda[k]$ 成立，其中移位
$k \in \mathbf{Z}$ 任意；这是显然的。
\end{proof}

\begin{lemma}
\label{etale-cohomology-lemma-projection-formula-proper-mod-n}
设 $f : X \to Y$ 是概形之间的固有态射。设 $\Lambda$ 是挠环。
设 $E \in D(X_\etale, \Lambda)$ 且
$K \in D(Y_\etale, \Lambda)$。则
$$
Rf_*E \otimes_\Lambda^\mathbf{L} K =
Rf_*(E \otimes_\Lambda^\mathbf{L} f^{-1}K)
$$
在 $D(Y_\etale, \Lambda)$ 中成立。
\end{lemma}

\begin{proof}
由 Cohomology on Sites, Section
\ref{sites-cohomology-section-projection-formula}，存在从左向右的典范映射。
我们在 $\overline{y}$ 处的茎上检验这个等式。
由固有基变换（采用 Lemma \ref{etale-cohomology-lemma-proper-base-change-mod-n} 中
$Y' = \overline{y}$ 的形式），问题约化到 $Y$ 是代数闭域之谱的情形。
这已在 Lemma \ref{etale-cohomology-lemma-pull-out-constant-mod-n} 中证明；其中由
Lemma \ref{etale-cohomology-lemma-cohomological-dimension-proper} 使用了
$\text{cd}(X) < \infty$。
\end{proof}






\section{\'etale 上同调中的 K\"unneth 公式}
\label{etale-cohomology-section-kunneth}

\noindent
我们先在其中一个因子固有的情形证明 K\"unneth 公式，再用该公式证明
开浸入的一项基变换性质。由此又得到“沿指向域之谱的态射作基变换”
（类似于光滑基变换）。最后再用它得到更一般的 K\"unneth 公式。

\begin{remark}
\label{etale-cohomology-remark-define-kunneth-map}
考察概形范畴中的 Cartesian 图：
$$
\xymatrix{
X \times_S Y \ar[d]_p \ar[r]_q \ar[rd]_c & Y \ar[d]^g \\
X \ar[r]^f & S
}
$$
设 $\Lambda$ 是环，并设 $E \in D(X_\etale, \Lambda)$ 且
$K \in D(Y_\etale, \Lambda)$。则存在典范映射
$$
Rf_*E \otimes_\Lambda^\mathbf{L} Rg_*K
\longrightarrow
Rc_*(p^{-1}E \otimes_\Lambda^\mathbf{L} q^{-1}K)
$$
例如，可以用典范映射 $Rf_*E \to Rc_*p^{-1}E$、
$Rg_*K \to Rc_*q^{-1}K$，以及 Cohomology on Sites,
Remark \ref{sites-cohomology-remark-cup-product} 中定义的相对杯积来定义它。
也可以取下列映射的伴随：
$$
c^{-1}(Rf_*E \otimes_\Lambda^\mathbf{L} Rg_*K)
=
p^{-1}f^{-1}Rf_*E \otimes_\Lambda^\mathbf{L} q^{-1} g^{-1}Rg_*K
\to
p^{-1}E \otimes_\Lambda^\mathbf{L} q^{-1}K
$$
这里使用伴随映射 $f^{-1}Rf_*E \to E$ 和
$g^{-1}Rg_*K \to K$。
\end{remark}


\begin{lemma}
\label{etale-cohomology-lemma-kunneth-one-proper}
设 $k$ 是可分闭域。设 $X$ 是 $k$ 上的固有概形。
设 $Y$ 是 $k$ 上的拟紧拟分离概形。
\begin{enumerate}
\item 如果 $E \in D^+(X_\etale)$ 的上同调层为挠层，且
$K \in D^+(Y_\etale)$，则
$$
R\Gamma(X \times_{\Spec(k)} Y,
\text{pr}_1^{-1}E
\otimes_\mathbf{Z}^\mathbf{L}
\text{pr}_2^{-1}K
)
=
R\Gamma(X, E)
\otimes_\mathbf{Z}^\mathbf{L}
R\Gamma(Y, K)
$$
\item 如果 $n \geq 1$ 是整数，$Y$ 在 $k$ 上有限型，
$E \in D(X_\etale, \mathbf{Z}/n\mathbf{Z})$，且
$K \in D(Y_\etale, \mathbf{Z}/n\mathbf{Z})$，则
$$
R\Gamma(X \times_{\Spec(k)} Y,
\text{pr}_1^{-1}E
\otimes_{\mathbf{Z}/n\mathbf{Z}}^\mathbf{L}
\text{pr}_2^{-1}K
)
=
R\Gamma(X, E)
\otimes_{\mathbf{Z}/n\mathbf{Z}}^\mathbf{L}
R\Gamma(Y, K)
$$
\end{enumerate}
\end{lemma}

\begin{proof}
证明 (1)。由 Lemma \ref{etale-cohomology-lemma-projection-formula-proper} 有
$$
R\text{pr}_{2, *}(
\text{pr}_1^{-1}E
\otimes_\mathbf{Z}^\mathbf{L}
\text{pr}_2^{-1}K) =
R\text{pr}_{2, *}(\text{pr}_1^{-1}E)
\otimes_\mathbf{Z}^\mathbf{L}
K
$$
由固有基变换（采用 Lemma \ref{etale-cohomology-lemma-proper-base-change} 的形式），
它等于对象
$$
\underline{R\Gamma(X, E)}
\otimes_\mathbf{Z}^\mathbf{L}
K
$$
，该对象属于 $D(Y_\etale)$。对此对象取 $R\Gamma(Y, -)$，
再由 $\text{pr}_2$ 的 Leray 谱序列，便得到 (1) 中等式的左边。
于是结论由 Lemma \ref{etale-cohomology-lemma-pull-out-constant} 得到。

\medskip\noindent
证明 (2)。它与 (1) 的证明完全相同，只是要使用
Lemmas \ref{etale-cohomology-lemma-projection-formula-proper-mod-n}、
\ref{etale-cohomology-lemma-proper-base-change-mod-n} 和
\ref{etale-cohomology-lemma-pull-out-constant-mod-n}，以及由
Lemma \ref{etale-cohomology-lemma-finite-cd} 得到的 $\text{cd}(Y) < \infty$。
\end{proof}

\begin{lemma}
\label{etale-cohomology-lemma-supported-in-closed-points}
设 $K$ 是可分闭域。设 $X$ 是 $K$ 上的有限型概形。设
$\mathcal{F}$ 是 $X_\etale$ 上的阿贝尔层，其支集包含在 $X$ 的闭点集内。
则 $H^q(X, \mathcal{F}) = 0$，只要 $q > 0$；并且
$\mathcal{F}$ 由全局截面生成。
\end{lemma}

\begin{proof}
（如果 $\mathcal{F}$ 是挠层，则消失性立即由
Lemma \ref{etale-cohomology-lemma-interlude-II} 得到。）
由 Lemma \ref{etale-cohomology-lemma-support-in-subset}，可把 $\mathcal{F}$
写成可构层 $\mathcal{F}_i$ 的滤过余极限；这些是
$\mathbf{Z}$-模层，其支集 $Z_i \subset X$ 是闭点的有限集。
由 Proposition \ref{etale-cohomology-proposition-closed-immersion-pushforward}，
这种层具有形式 $(Z_i \to X)_*\mathcal{G}_i$，其中
$\mathcal{G}_i$ 是 $Z_i$ 上的层。由于 $K$ 可分闭，概形
$Z_i$ 是若干可分闭域之谱的有限不交并。回顾
$H^q(Z_i, \mathcal{G}_i) = H^q(X, \mathcal{F}_i)$；这是由
$Z_i \to X$ 的 Leray 谱序列以及该态射的高阶正像消失得到的
（Proposition \ref{etale-cohomology-proposition-finite-higher-direct-image-zero}）。
由 Lemmas \ref{etale-cohomology-lemma-equivalence-abelian-sheaves-point} 和
\ref{etale-cohomology-lemma-compare-cohomology-point}，可知
$H^q(Z_i, \mathcal{G}_i)$ 在 $q > 0$ 时为零，而且
$H^0(Z_i, \mathcal{G}_i)$ 生成 $\mathcal{G}_i$。
因此 $H^q(X, \mathcal{F}_i)$ 在 $q > 0$ 时消失，并且
$\mathcal{F}_i$ 由全局截面生成。由 Theorem \ref{etale-cohomology-theorem-colimit} 可知
$H^q(X, \mathcal{F}) = 0$，只要 $q > 0$。
现在证明完成，因为阿贝尔群层中由全局截面生成的层的滤过余极限仍由
全局截面生成（略去细节）。
\end{proof}

\begin{lemma}
\label{etale-cohomology-lemma-vanishing-closed-points}
设 $K$ 是可分闭域。设 $X$ 是 $K$ 上的有限型概形。设
$Q \in D(X_\etale)$。假设 $Q_{\overline{x}}$ 只有在 $x$ 是
$X$ 的闭点时才可能非零。则
$$
Q = 0 \Leftrightarrow H^i(X, Q) = 0 \text{ 对所有 }i
$$
\end{lemma}

\begin{proof}
从左向右的蕴含是平凡的。因此只须证明反向蕴含。

\medskip\noindent
先假设 $Q$ 下有界；这种情形足以处理几乎所有应用。如果 $Q$ 非零，
就考察最小的 $i$，使得上同调层 $H^i(Q)$ 非零。
由 Lemma \ref{etale-cohomology-lemma-supported-in-closed-points} 有
$H^i(X, Q) = H^0(X, H^i(Q)) \not = 0$，结论随即得到。

\medskip\noindent
一般情形。令 $\mathcal{B} \subset \Ob(X_\etale)$ 为拟紧对象的集合。
由 Lemma \ref{etale-cohomology-lemma-supported-in-closed-points}，Cohomology on Sites,
Lemma \ref{sites-cohomology-lemma-cohomology-over-U-trivial} 的假设成立。
因此有 $H^q(U, Q) = H^0(U, H^q(Q))$，对所有
$U \in \mathcal{B}$ 都成立。特别地，这对 $U = X$ 成立。
于是结论由 Lemma \ref{etale-cohomology-lemma-supported-in-closed-points} 得到，因为
$Q$ 在 $D(X_\etale)$ 中为零，当且仅当 $H^q(Q)$ 对所有 $q$ 都为零。
\end{proof}

\begin{lemma}
\label{etale-cohomology-lemma-kunneth-localize-on-X}
设 $K$ 是域。设 $j : U \to X$ 是 $K$ 上有限型概形之间的开浸入。
设 $Y$ 是 $K$ 上的有限型概形。考察图
$$
\xymatrix{
Y \times_{\Spec(K)} X \ar[d]_q  &
Y \times_{\Spec(K)} U \ar[l]^h \ar[d]^p \\
X & U \ar[l]_j
}
$$
则基变换映射 $q^{-1}Rj_*\mathcal{F} \to Rh_*p^{-1}\mathcal{F}$
是同构，其中 $\mathcal{F}$ 是 $U_\etale$ 上的阿贝尔层，
其各茎为阶在 $K$ 中可逆的挠群。
\end{lemma}

\begin{proof}
写成 $\mathcal{F} = \colim \mathcal{F}[n]$，其中余极限取遍
$K$ 中可逆整数构成的乘法系。由于在当前情形上同调与滤过余极限可交换
（精确参考见 Lemma \ref{etale-cohomology-lemma-base-change-Rf-star-colim}），
只须对 $\mathcal{F}[n]$ 证明引理。因此可以假设 $\mathcal{F}$ 是
$\mathbf{Z}/n\mathbf{Z}$-模层，其中某个 $n$ 在 $K$ 中可逆
（证明最后才会用到这一点）。在证明中，用 $X \times_K Y$ 简记
$\Spec(K)$ 上的纤维积。
对 $\dim(X) + \dim(Y)$ 作归纳来证明引理。如果 $\dim(X) \leq 0$，
引理是平凡的，因为此时 $U$ 是 $X$ 的既开又闭子概形。
选取一点 $z \in X \times_K Y$。我们将证明在 $\overline{z}$ 处的茎映射是同构。

\medskip\noindent
假设 $z \mapsto x \in X$，并假设 $\text{trdeg}_K(\kappa(x)) > 0$。
置 $X' = \Spec(\mathcal{O}_{X, x}^{sh})$，并以
$U' \subset X'$ 表示 $U$ 的逆像。考察基变换
$$
\xymatrix{
Y \times_K X' \ar[d]_{q'}  &
Y \times_K U' \ar[l]^{h'} \ar[d]^{p'} \\
X' & U' \ar[l]_{j'}
}
$$
，它由 $X' \to X$ 对原图作基变换得到。注意 $X' \to X$ 是
\'etale 态射的滤过余极限。由采用 Lemma
\ref{etale-cohomology-lemma-smooth-base-change-general} 形式的光滑基变换，
$q^{-1}Rj_*\mathcal{F} \to Rh_*p^{-1}\mathcal{F}$ 在 $X'$ 上、
再到 $Y \times_K X'$ 的拉回是映射
$(q')^{-1}Rj'_*\mathcal{F}' \to Rj'_*(p')^{-1}\mathcal{F}'$，其中
$\mathcal{F}'$ 是 $\mathcal{F}$ 到 $U'$ 的拉回。
（在这一步使用本质上平凡的 \'etale 基变换已经足够。）
因此只须证明 $(q')^{-1}Rj'_*\mathcal{F}' \to Rj'_*(p')^{-1}\mathcal{F}'$
是同构，便可证明原映射在 $\overline{z}$ 处的茎上是同构。
由 Lemma \ref{etale-cohomology-lemma-strictly-henselian}，存在可分闭域 $L/K$，使得
$X' = \lim X_i$，其中 $X_i$ 是 $L$ 上有限型的仿射概形，且
$\dim(X_i) < \dim(X)$。当 $i$ 充分大时，存在开子概形
$U_i \subset X_i$，其限制为 $U'$（在 $X'$ 中）。
可以对图应用归纳假设：
$$
\vcenter{
\xymatrix{
Y \times_K X_i \ar[d]_{q_i}  &
Y \times_K U_i \ar[l]^{h_i} \ar[d]^{p_i} \\
X_i & U_i \ar[l]_{j_i}
}
}
\quad\text{equal to}\quad
\vcenter{
\xymatrix{
Y_L \times_L X_i \ar[d]_{q_i}  &
Y_L \times_L U_i \ar[l]^{h_i} \ar[d]^{p_i} \\
X_i & U_i \ar[l]_{j_i}
}
}
$$
这里底域为 $L$，并取 $\mathcal{F}$ 到这些图的拉回。
由 Lemma \ref{etale-cohomology-lemma-base-change-Rf-star-colim} 可知，映射
$(q')^{-1}Rj'_*\mathcal{F}' \to Rj'_*(p')^{-1}\mathcal{F}$ 是同构。

\medskip\noindent
假设 $z \mapsto y \in Y$，并假设 $\text{trdeg}_K(\kappa(y)) > 0$。
令 $Y' = \Spec(\mathcal{O}_{X, x}^{sh})$。
由 Lemma \ref{etale-cohomology-lemma-strictly-henselian}，存在可分闭域 $L/K$，使得
$Y' = \lim Y_i$，其中 $Y_i$ 是 $L$ 上有限型的仿射概形，且
$\dim(Y_i) < \dim(Y)$。特别地，$Y'$ 是 $L$ 上的概形。
用下标 $L$ 表示从 $K$ 上概形到 $L$ 上概形的基变换。考察交换图
$$
\vcenter{
\xymatrix{
Y' \times_K X \ar[d]_f  &
Y' \times_K U \ar[l]^{h'} \ar[d]^{f'} \\
Y \times_K X \ar[d]_q  &
Y \times_K U \ar[l]^h \ar[d]^p \\
X & U \ar[l]_j
}
}
\quad\text{且}\quad
\vcenter{
\xymatrix{
Y' \times_L X_L \ar[d]_{q'}  &
Y' \times_L U_L \ar[l]^{h'} \ar[d]^{p'} \\
X_L \ar[d] &
U_L \ar[l]^{j_L} \ar[d] \\
X & U \ar[l]_j
}
}
$$
并注意，左图与右图的顶行相同，底行也相同。
由光滑基变换可知
$f^{-1}Rh_*p^{-1}\mathcal{F} = Rh'_*(f')^{-1}p^{-1}\mathcal{F}$
（与上一段类似）。由 $\Spec(L) \to \Spec(K)$ 的光滑基变换
（Lemma \ref{etale-cohomology-lemma-base-change-field-extension}），可知
$Rj_{L, *}\mathcal{F}_L$ 是 $Rj_*\mathcal{F}$ 到 $X_L$ 的拉回。
合并这两个观察可知，只须证明右图上方方形的基变换映射是同构，
便可证明原映射在 $\overline{z}$ 处的茎上是同构\footnote{这里使用：
基变换映射的“竖直复合”仍是基变换映射，见 Cohomology on Sites,
Remark \ref{sites-cohomology-remark-compose-base-change}。}。
然后利用 $Y' = \lim Y_i$，完全照上一段论证，可见归纳假设迫使
$Y' \times_K X$ 上的映射成为同构。

\medskip\noindent
因此，任何使 $\dim(X) + \dim(Y)$ 最小的反例，都只会使映射
$q^{-1}Rj_*\mathcal{F} \to Rh_*p^{-1}\mathcal{F}$
在 $X \times_K Y$ 的闭点之茎上不是同构（因为一点
$z$（属于 $X \times_K Y$）是闭点，当且仅当 $z$ 在 $X$ 与 $Y$
中的像都是闭点）。由于只须局部证明同构，可以假设 $X$ 和 $Y$ 仿射。
但这时可选取稠密开浸入 $Y \to Y'$，使 $Y'$ 射影。
（选取闭浸入 $Y \to \mathbf{A}^n_K$，并令 $Y'$ 为 $Y$ 在
$\mathbf{P}^n_K$ 中的概形论闭包。）于是 $\dim(Y') = \dim(Y)$，
从而得到一个“最小”反例，其中 $Y$ 在 $K$ 上射影。
下一段将说明这是不可能的。

\medskip\noindent
考察引理陈述中的图，并假设 $q^{-1}Rj_*\mathcal{F} \to Rh_*p^{-1}\mathcal{F}$
在 $X \times_K Y$ 的所有非闭点处都是同构，且 $Y$ 固有。
该映射到 $(X \times_K Y)_{K^{sep}}$ 的限制，就是把引理之图基变换到
$K^{sep}$ 后所得的相应映射。因此可以并且确实假设 $K$ 可分闭。
选取正规三角
$$
q^{-1}Rj_*\mathcal{F} \to Rh_*p^{-1}\mathcal{F} \to Q \to
(q^{-1}Rj_*\mathcal{F})[1]
$$
，它位于 $D((X \times_K Y)_\etale)$ 中。由于 $Q$ 支撑在闭点上，
由 Lemma \ref{etale-cohomology-lemma-vanishing-closed-points}，只须证明
$H^i(X \times_K Y, Q) = 0$ 对所有 $i$ 成立。
因此只须证明
$q^{-1}Rj_*\mathcal{F} \to Rh_*p^{-1}\mathcal{F}$
在上同调上诱导同构。回顾 $\mathcal{F}$ 被 $n$ 零化，且该整数
在 $K$ 中可逆。由 Lemma \ref{etale-cohomology-lemma-kunneth-one-proper} 的 K\"unneth 公式有
\begin{align*}
R\Gamma(X \times_K Y, q^{-1}Rj_*\mathcal{F})
&=
R\Gamma(X, Rj_*\mathcal{F})
\otimes_{\mathbf{Z}/n\mathbf{Z}}^\mathbf{L}
R\Gamma(Y, \mathbf{Z}/n\mathbf{Z}) \\
& =
R\Gamma(U, \mathcal{F})
\otimes_{\mathbf{Z}/n\mathbf{Z}}^\mathbf{L}
R\Gamma(Y, \mathbf{Z}/n\mathbf{Z})
\end{align*}
并且
$$
R\Gamma(X \times_K Y, Rh_*p^{-1}\mathcal{F}) =
R\Gamma(U \times_K Y, p^{-1}\mathcal{F}) =
R\Gamma(U, \mathcal{F})
\otimes_{\mathbf{Z}/n\mathbf{Z}}^\mathbf{L}
R\Gamma(Y, \mathbf{Z}/n\mathbf{Z})
$$
证明至此完成。
\end{proof}

\begin{lemma}
\label{etale-cohomology-lemma-punctual-base-change}
设 $K$ 是域。对任意交换图
$$
\xymatrix{
X \ar[d] & X' \ar[l] \ar[d]_{f'} & Y \ar[l]^h \ar[d]^e \\
\Spec(K) & S' \ar[l] & T \ar[l]_g
}
$$
，其中各概形定义在 $K$ 上，
$X' = X \times_{\Spec(K)} S'$ 且 $Y = X' \times_{S'} T$，
$g$ 拟紧拟分离；再对任意阿贝尔层 $\mathcal{F}$（定义在
$T_\etale$ 上），若其各茎是阶在 $K$ 中可逆的挠群，则基变换映射
$$
(f')^{-1}Rg_*\mathcal{F}
\longrightarrow
Rh_*e^{-1}\mathcal{F}
$$
是同构。
\end{lemma}

\begin{proof}
问题在 $X$ 上是局部的，故可假设 $X$ 仿射。由 Limits, Lemma
\ref{limits-lemma-relative-approximation}，可把
$X = \lim X_i$ 写成余滤过极限，过渡态射仿射，而其中的概形
$X_i$ 在 $K$ 上有限型。记
$X'_i =  X_i \times_{\Spec(K)} S'$ 且 $Y_i = X'_i \times_{S'} T$。
由 Lemma \ref{etale-cohomology-lemma-base-change-Rf-star-colim}，只须对顶点为
$X_i, Y_i, S_i, T_i$ 的方形证明断言。因此可以假设 $X$ 在 $K$ 上有限型。
类似地，可写 $\mathcal{F} = \colim \mathcal{F}[n]$，其中余极限取遍
$K$ 中可逆整数构成的乘法系。上面使用的同一引理把问题约化到：
$\mathcal{F}$ 是 $\mathbf{Z}/n\mathbf{Z}$-模层，其中某个
$n$ 在 $K$ 中可逆。

\medskip\noindent
可以把 $K$ 替换为其代数闭包 $\overline{K}$。事实上，依
Lemma \ref{etale-cohomology-lemma-base-change-field-extension}，正像的形成与到
$\overline{K}$ 的基变换可交换（对 $g$ 与 $h$ 都成立）；而且只须在限制到
$X'_{\overline{K}}$ 后证明两者相同。接着，由 \'etale 上同调的拓扑不变性，
可以用既约化代替 $X$，见
Proposition \ref{etale-cohomology-proposition-topological-invariance}。
作此替换后，态射 $X \to \Spec(K)$ 平坦、有限呈现，并具有几何既约纤维；
任意基变换也有同样性质，特别是 $X' \to S'$。因此由
Lemma \ref{etale-cohomology-lemma-fppf-reduced-fibres-base-change-f-star}，
$(f')^{-1}g_*\mathcal{F} \to Rh_*e^{-1}\mathcal{F}$ 是同构。

\medskip\noindent
此时可以应用 Lemma \ref{etale-cohomology-lemma-base-change-does-not-hold-post}，
可见只须证明如下断言：给定交换图
$$
\xymatrix{
X \ar[d]_f & X' \ar[d] \ar[l] & Y \ar[l]^h \ar[d] \\
\Spec(K) & S' \ar[l] & \Spec(L) \ar[l]
}
$$
，两方形均 Cartesian；其中 $S'$ 仿射、整且正规，其函数域
$K$ 代数闭。则 $R^qh_*(\mathbf{Z}/d\mathbf{Z})$ 为零，只要
$q > 0$ 且 $d | n$。注意，该消失性等价于
$$
(f')^{-1}R^q(\Spec(L) \to S')_*\mathbf{Z}/d\mathbf{Z}
\longrightarrow
R^qh_*\mathbf{Z}/d\mathbf{Z}
$$
是同构，因为例如由
Lemma \ref{etale-cohomology-lemma-Rf-star-zero-normal-with-alg-closed-function-field}，
左边为零。

\medskip\noindent
写 $S' = \Spec(B)$，于是 $L$ 是 $B$ 的分式域。把
$B = \bigcup_{i \in I} B_i$ 写成其有限型 $K$-子代数 $B_i$ 的并。
令 $J$ 为二元组 $(i, g)$ 的集合，其中 $i \in I$ 且
$g \in B_i$ 非零；其序为：$(i', g') \geq (i, g)$ 当且仅当
$i' \geq i$，并且 $g$ 在 $(B_{i'})_{g'}$ 中的像可逆。于是
$L = \colim_{(i, g) \in J} (B_i)_g$。
对 $j = (i, g) \in J$，置 $S_j = \Spec(B_i)$ 且
$U_j = \Spec((B_i)_g)$。于是
$$
\vcenter{
\xymatrix{
X' \ar[d] & Y \ar[l]^h \ar[d] \\
S' & \Spec(L) \ar[l]
}
}
\quad\text{是下列对象的余极限}\quad
\vcenter{
\xymatrix{
X \times_K S_j \ar[d] & X \times_K U_j \ar[l]^{h_j} \ar[d] \\
S_j & U_j \ar[l]
}
}
$$
因此可以应用 Lemma \ref{etale-cohomology-lemma-base-change-Rf-star-colim}，
可见只须证明右侧诸图中的基变换成立；这正是
Lemma \ref{etale-cohomology-lemma-kunneth-localize-on-X} 中已经证明的内容。
\end{proof}

\begin{lemma}
\label{etale-cohomology-lemma-punctual-base-change-upgrade}
设 $K$ 是域。设 $n \geq 1$ 在 $K$ 中可逆。考察交换图
$$
\xymatrix{
X \ar[d] & X' \ar[l]^p \ar[d]_{f'} & Y \ar[l]^h \ar[d]^e \\
\Spec(K) & S' \ar[l] & T \ar[l]_g
}
$$
，其中
$X' = X \times_{\Spec(K)} S'$ 且 $Y = X' \times_{S'} T$，
$g$ 拟紧拟分离。典范映射
$$
p^{-1}E \otimes_{\mathbf{Z}/n\mathbf{Z}}^\mathbf{L} (f')^{-1}Rg_*F
\longrightarrow
Rh_*(h^{-1}p^{-1}E \otimes_{\mathbf{Z}/n\mathbf{Z}}^\mathbf{L} e^{-1}F)
$$
是同构，只要 $E$ 属于 $D^+(X_\etale, \mathbf{Z}/n\mathbf{Z})$，
其 Tor 振幅包含在 $[a, \infty]$ 内，其中 $a \in \mathbf{Z}$ 为某个整数；
并且 $F$ 属于 $D^+(T_\etale, \mathbf{Z}/n\mathbf{Z})$。
\end{lemma}

\begin{proof}
本引理把 Lemma \ref{etale-cohomology-lemma-punctual-base-change} 推广到导出范畴中的对象；
断言成立，是因为在 Lemma \ref{etale-cohomology-lemma-punctual-base-change} 中，
概形 $X$ 定义在 $K$ 上且是任意的。我们强烈建议读者跳过这个繁琐的证明
（另一选择是只读最后一段）。

\medskip\noindent
可以把 $E$ 表示为下有界 K-平坦复形 $\mathcal{E}^\bullet$；
该复形由平坦 $\mathbf{Z}/n\mathbf{Z}$-模组成。见 Cohomology on Sites,
Lemma \ref{sites-cohomology-lemma-bounded-below-tor-amplitude}。
选取整数 $b$，使得 $H^i(F) = 0$ 对 $i < b$ 成立。
选取充分大的整数 $N$，并考察短正合列
$$
0 \to \sigma_{\geq N + 1}\mathcal{E}^\bullet \to
\mathcal{E}^\bullet \to
\sigma_{\leq N}\mathcal{E}^\bullet \to 0
$$
，其中使用朴素截断。它产生正规三角
$E'' \to E \to E' \to E''[1]$，该三角位于
$D(X_\etale, \mathbf{Z}/n\mathbf{Z})$ 中。固定 $F$ 后，引理陈述中箭头的两边
都是关于 $E$ 的正合函子。注意
$$
p^{-1}E'' \otimes_{\mathbf{Z}/n\mathbf{Z}}^\mathbf{L} (f')^{-1}Rg_*F
\quad\text{且}\quad
Rh_*(h^{-1}p^{-1}E'' \otimes_{\mathbf{Z}/n\mathbf{Z}}^\mathbf{L} e^{-1}F)
$$
位于次数 $\geq N + b$。因此，如果能对对象 $E'$ 证明引理，
便可知引理在次数 $\leq N + b$ 中成立，从而得到结论。略去若干细节。
于是可以假设 $E$ 由平坦 $\mathbf{Z}/n\mathbf{Z}$-模组成的有界复形表示。
再作一次同类论证，可以假设 $E$ 由单个平坦
$\mathbf{Z}/n\mathbf{Z}$-模 $\mathcal{E}$ 给出。

\medskip\noindent
接着，对变量 $F$ 使用同一论证，把问题约化到 $F$ 由单个
$\mathbf{Z}/n\mathbf{Z}$-模层 $\mathcal{F}$ 给出的情形。
设 $\mathcal{F}$ 被整数 $m | n$ 零化。如果 $\ell$ 是整除
$m$ 的素数，且 $m > \ell$，则可以考察短正合列
$0 \to \mathcal{F}[\ell] \to \mathcal{F} \to
\mathcal{F}/\mathcal{F}[\ell] \to 0$
，从而约化到更小的 $m$。最终，问题约化到 $\mathcal{F}$ 被某个素数
$\ell$ 零化的情形，该素数整除 $n$。这时注意
$$
p^{-1}\mathcal{E}
\otimes_{\mathbf{Z}/n\mathbf{Z}}^\mathbf{L}
(f')^{-1}Rg_*\mathcal{F}
=
p^{-1}(\mathcal{E}/\ell \mathcal{E})
\otimes_{\mathbf{F}_\ell}^\mathbf{L}
(f')^{-1}Rg_*\mathcal{F}
$$
，这是由 $\mathcal{E}$ 的平坦性得到的。另一项同理。
因此问题约化到处理 $\mathbf{F}_\ell$-向量空间层的情形，下面讨论它。

\medskip\noindent
假设 $\ell$ 是在 $K$ 中可逆的素数。假设 $\mathcal{E}$、
$\mathcal{F}$ 是 $\mathbf{F}_\ell$-向量空间层，分别定义在
$X_\etale$ 与 $T_\etale$ 上。我们要证明
$$
p^{-1}\mathcal{E} \otimes_{\mathbf{F}_\ell} (f')^{-1}R^qg_*\mathcal{F}
\longrightarrow
R^qh_*(h^{-1}p^{-1}\mathcal{E} \otimes_{\mathbf{F}_\ell} e^{-1}\mathcal{F})
$$
对每个 $q \geq 0$ 都是同构。这个问题在 $X$ 上是局部的，故可假设
$X$ 仿射。可把 $\mathcal{E}$ 写成可构
$\mathbf{F}_\ell$-向量空间层的滤过余极限，这些层定义在 $X_\etale$ 上，见
Lemma \ref{etale-cohomology-lemma-torsion-colimit-constructible}。
由于张量积与滤过余极限可交换，高阶正像也如此
（Lemma \ref{etale-cohomology-lemma-relative-colimit}），可以假设 $\mathcal{E}$ 是
可构 $\mathbf{F}_\ell$-向量空间层，定义在 $X_\etale$ 上。
于是可以选取整数 $m$ 和有限且有限呈现的态射
$\pi_i : X_i \to X$，其中 $i = 1, \ldots, m$，使得存在单射
$$
\mathcal{E} \to
\bigoplus\nolimits_{i = 1, \ldots, m}
\pi_{i, *}\mathbf{F}_\ell
$$
，见 Lemma \ref{etale-cohomology-lemma-constructible-maps-into-constant-general}。
注意该直和也是可构层
（Lemma \ref{etale-cohomology-lemma-finite-pushforward-constructible}），
故其余核也可构（Lemma \ref{etale-cohomology-lemma-constructible-abelian}）。
由维数移位，即对 $q$ 作归纳，在可构
$\mathbf{F}_\ell$-向量空间层（定义在 $X_\etale$ 上）的范畴中，只须对层
$\pi_{i, *}\mathbf{F}_\ell$ 证明结论
（略去细节；提示：先证明 $q = 0$ 时的单射性，并让
$\mathcal{E}$ 遍历所有可构层）。
为了证明这种情形，把引理中的图扩充为
$$
\xymatrix{
X_i \ar[d]^{\pi_i} &
X'_i \ar[l]^{p_i} \ar[d]^{\pi'_i} &
Y_i \ar[l]^{h_i} \ar[d]^{\rho_i} \\
X \ar[d] & X' \ar[l]^p \ar[d]_{f'} & Y \ar[l]^h \ar[d]^e \\
\Spec(K) & S' \ar[l] & T \ar[l]_g
}
$$
，其中所有方形都 Cartesian。在以下等式中，将使用
$R\pi_{i, *} = \pi_{i, *}$，以及 $\pi'_i$、$\rho_i$ 的类似等式；
还将使用 Leray 谱序列、Lemma
\ref{etale-cohomology-lemma-finite-pushforward-commutes-with-base-change}，并对
$\pi_i$、$\pi'_i$、$\rho_i$ 使用
引理 \ref{etale-cohomology-lemma-projection-formula-proper-mod-n}
（对有限态射而言，该引理几乎是平凡的）。由此可见
\begin{align*}
p^{-1}\pi_{i, *}\mathbf{F}_\ell
\otimes_{\mathbf{F}_\ell} (f')^{-1}R^qg_*\mathcal{F}
& =
\pi'_{i, *}\mathbf{F}_\ell
\otimes_{\mathbf{F}_\ell} (f')^{-1}R^qg_*\mathcal{F} \\
& =
\pi'_{i, *}((\pi'_i)^{-1} (f')^{-1}R^qg_*\mathcal{F})
\end{align*}
类似地，有
\begin{align*}
R^qh_*(h^{-1}p^{-1} \pi_{i, *}\mathbf{F}_\ell
\otimes_{\mathbf{F}_\ell} e^{-1}\mathcal{F})
& =
R^qh_*(\rho_{i, *}\mathbf{F}_\ell
\otimes_{\mathbf{F}_\ell} e^{-1}\mathcal{F}) \\
& =
R^qh_*(\rho_i^{-1}e^{-1}\mathcal{F}) \\
& =
\pi'_{i, *}R^qh_{i, *} \rho_i^{-1}e^{-1}\mathcal{F})
\end{align*}
由于
$R^qh_{i, *} \rho_i^{-1}e^{-1}\mathcal{F} = 
(\pi'_i)^{-1} (f')^{-1}R^qg_*\mathcal{F}$，这由
Lemma \ref{etale-cohomology-lemma-punctual-base-change} 得到，故结论得证。
\end{proof}

\begin{lemma}
\label{etale-cohomology-lemma-punctual-base-change-upgrade-unbounded}
设 $K$ 是域。设 $n \geq 1$ 在 $K$ 中可逆。考察交换图
$$
\xymatrix{
X \ar[d] & X' \ar[l]^p \ar[d]_{f'} & Y \ar[l]^h \ar[d]^e \\
\Spec(K) & S' \ar[l] & T \ar[l]_g
}
$$
，其中各概形在 $K$ 上有限型，
$X' = X \times_{\Spec(K)} S'$ 且 $Y = X' \times_{S'} T$。
典范映射
$$
p^{-1}E \otimes_{\mathbf{Z}/n\mathbf{Z}}^\mathbf{L} (f')^{-1}Rg_*F
\longrightarrow
Rh_*(h^{-1}p^{-1}E \otimes_{\mathbf{Z}/n\mathbf{Z}}^\mathbf{L} e^{-1}F)
$$
是同构，其中 $E$ 属于 $D(X_\etale, \mathbf{Z}/n\mathbf{Z})$，
$F$ 属于 $D(T_\etale, \mathbf{Z}/n\mathbf{Z})$。
\end{lemma}

\begin{proof}
我们将利用这些函子与直和可交换这一事实，把问题约化到
Lemma \ref{etale-cohomology-lemma-punctual-base-change-upgrade}。建议读者跳过证明。
回顾：导出张量积与直和可交换；（导出）拉回也与直和可交换；而
$Rh_*$ 与 $Rg_*$ 同样与直和可交换，见 Lemmas \ref{etale-cohomology-lemma-finite-cd} 和
\ref{etale-cohomology-lemma-finite-cd-mod-n-direct-sums}
（这里用到了各概形在 $K$ 上有限型）。

\medskip\noindent
为完成证明，可作如下论证。先写
$E = \text{hocolim} \tau_{\leq N} E$。由于这些函子与直和可交换，
它们也与同伦余极限可交换。因此只须对上有界的 $E$ 证明引理。
类似地，对 $F$ 可以假设 $F$ 上有界。于是可以把 $E$ 表示为上有界复形
$\mathcal{E}^\bullet$，其项为 $\mathbf{Z}/n\mathbf{Z}$-模层。于是
$$
\mathcal{E}^\bullet = \colim \sigma_{\geq -N}\mathcal{E}^\bullet
$$
（朴素截断）。故可假设 $\mathcal{E}^\bullet$ 是
$\mathbf{Z}/n\mathbf{Z}$-模层组成的有界复形。对 $F$，选取由平坦（！）
$\mathbf{Z}/n\mathbf{Z}$-模层组成的上有界复形。
这样便约化到 $F$ 由平坦 $\mathbf{Z}/n\mathbf{Z}$-模层组成的有界复形表示
这一情形。此时应用 Lemma \ref{etale-cohomology-lemma-punctual-base-change-upgrade} 即得结论。
\end{proof}

\begin{lemma}
\label{etale-cohomology-lemma-kunneth}
设 $k$ 是可分闭域。设 $X$ 与 $Y$ 是 $k$ 上的有限型概形。
设 $n \geq 1$ 是在 $k$ 中可逆的整数。则对
$E \in D(X_\etale, \mathbf{Z}/n\mathbf{Z})$ 与
$K \in D(Y_\etale, \mathbf{Z}/n\mathbf{Z})$ 有
$$
R\Gamma(X \times_{\Spec(k)} Y,
\text{pr}_1^{-1}E
\otimes_{\mathbf{Z}/n\mathbf{Z}}^\mathbf{L}
\text{pr}_2^{-1}K
)
=
R\Gamma(X, E)
\otimes_{\mathbf{Z}/n\mathbf{Z}}^\mathbf{L}
R\Gamma(Y, K)
$$
\end{lemma}

\begin{proof}
由 Lemma \ref{etale-cohomology-lemma-punctual-base-change-upgrade-unbounded} 有
$$
R\text{pr}_{1, *}(
\text{pr}_1^{-1}E
\otimes_{\mathbf{Z}/n\mathbf{Z}}^\mathbf{L}
\text{pr}_2^{-1}K) =
E \otimes_{\mathbf{Z}/n\mathbf{Z}}^\mathbf{L}
\underline{R\Gamma(Y, K)}
$$
结论由 Lemma \ref{etale-cohomology-lemma-pull-out-constant-mod-n} 得到；
可以使用该引理，是因为由 Lemma \ref{etale-cohomology-lemma-finite-cd} 有
$\text{cd}(X) < \infty$。
\end{proof}














\section{混沌拓扑与 Zariski 拓扑的比较}
\label{etale-cohomology-section-compare-chaotic-Zariski}

\noindent
在构造仿射概形的结构层时，我们先构造它在仿射开子集上的取值，
然后再延拓到所有开子集。对于构造概形 $X$ 上的阿贝尔群复形，
类似的构造往往也很有用。回顾 $X_{affine, Zar}$ 表示 $X$ 的
仿射开子集所成的范畴，其拓扑由标准 Zariski 覆盖给出，见
Topologies, Definition \ref{topologies-definition-big-small-Zariski}。
提醒读者，$X_{affine, Zar}$ 的拓扑斯就是 $X$ 的小 Zariski 拓扑斯，
见 Topologies, Lemma \ref{topologies-lemma-alternative-zariski}。
本节用 $X_{affine}$ 表示同一底范畴，但赋予混沌拓扑；也就是说，
层与预层相同。于是得到位点态射
$$
\epsilon : X_{affine, Zar} \longrightarrow X_{affine}
$$
，如 Cohomology on Sites, Section \ref{sites-cohomology-section-compare}
中所述。

\begin{lemma}
\label{etale-cohomology-lemma-check-zar}
在上述情形中，设 $K$ 是 $D^+(X_{affine})$ 的对象。则 $K$ 属于
（全忠实）函子
$R\epsilon_* ; D(X_{affine, Zar}) \to D(X_{affine})$ 的本质像，
当且仅当下列两个条件成立：
\begin{enumerate}
\item $R\Gamma(\emptyset, K)$ 在 $D(\textit{Ab})$ 中为零；
\item 如果 $U = V \cup W$，其中 $U, V, W \subset X$ 是仿射开子集，
而 $V, W \subset U$ 是标准开子集
（Algebra, Definition \ref{algebra-definition-Zariski-topology}），
则 Cohomology on Sites, Lemma
\ref{sites-cohomology-lemma-c-square} 中的映射
$c^K_{U, V, W, V \cap W}$ 是拟同构。
\end{enumerate}
\end{lemma}

\begin{proof}
（由 Cohomology on Sites, Section
\ref{sites-cohomology-section-compare} 中的讨论，函子
$R\epsilon_*$ 是全忠实的。）
除去一个与空集有关的障碍，结论来自非常一般的 Cohomology on Sites,
Lemma \ref{sites-cohomology-lemma-descent-squares}；该引理的假设由
Schemes, Lemma \ref{schemes-lemma-sheaf-on-affines} 和
Cohomology on Sites，引理
\ref{sites-cohomology-lemma-descent-squares-helper} 保证。

\medskip\noindent
为绕开这一障碍，记 $X_{affine, almost-chaotic}$ 为如下位点：
其中空对象 $\emptyset$ 有两个覆盖，即
$\{\emptyset \to \emptyset\}$ 与空覆盖
（有关讨论见 Sites, Example \ref{sites-example-site-topological}）。
于是有位点态射
$$
X_{affine, Zar} \to X_{affine, almost-chaotic} \to X_{affine}
$$
上述论证适用于第一个箭头。然后留给读者验证：对象 $K$ 若属于
$D^+(X_{affine})$，则它属于（全忠实）函子
$D(X_{affine, almost-chaotic}) \to D(X_{affine})$ 的本质像，
当且仅当 $R\Gamma(\emptyset, K)$ 在 $D(\textit{Ab})$ 中为零。
\end{proof}









\section{大拓扑斯与小拓扑斯的比较}
\label{etale-cohomology-section-compare}

\noindent
设 $S$ 是概形。在 Topologies, Lemma
\ref{topologies-lemma-at-the-bottom-etale} 中，我们引入了比较态射
$\pi_S : (\Sch/S)_\etale \to S_\etale$ 和
$i_S : \Sh(S_\etale) \to \Sh((\Sch/S)_\etale)$
，满足 $\pi_S \circ i_S = \text{id}$ 且
$\pi_{S, *} = i_S^{-1}$。更一般地，如果 $f : T \to S$ 是
$(\Sch/S)_\etale$ 的对象，则存在态射
$i_f : \Sh(T_\etale) \to \Sh((\Sch/S)_\etale)$，使得
$f_{small} = \pi_S \circ i_f$，见 Topologies, Lemmas
\ref{topologies-lemma-put-in-T-etale} 和
\ref{topologies-lemma-morphism-big-small-etale}。在 Descent, Remark
\ref{descent-remark-change-topologies-ringed} 中，我们把这些态射延拓为
环化网站的态射
$$
\pi_S : ((\Sch/S)_\etale, \mathcal{O}) \to (S_\etale, \mathcal{O}_S)
$$
以及环化拓扑斯的态射
$$
i_S : (\Sh(S_\etale), \mathcal{O}_S) \to (\Sh((\Sch/S)_\etale), \mathcal{O})
$$
和
$$
i_f : (\Sh(T_\etale), \mathcal{O}_T) \to (\Sh((\Sch/S)_\etale, \mathcal{O}))
$$
。注意，限制 $i_S^{-1} = \pi_{S, *}$（见 Topologies, Definition
\ref{topologies-definition-restriction-small-etale}）把
$\mathcal{O}$ 变为 $\mathcal{O}_S$。类似地，$i_f^{-1}$ 把
$\mathcal{O}$ 变为 $\mathcal{O}_T$。见 Descent, Remark
\ref{descent-remark-change-topologies-ringed}。因此有
$i_S^*\mathcal{F} = i_S^{-1}\mathcal{F}$ 和
$i_f^*\mathcal{F} = i_f^{-1}\mathcal{F}$，这对任意
$\mathcal{O}$-模 $\mathcal{F}$（定义在 $(\Sch/S)_\etale$ 上）成立。
特别地，$i_S^*$ 与
$i_f^*$ 都是正合函子。函子 $i_S^*$ 常记作
$\mathcal{F} \mapsto \mathcal{F}|_{S_\etale}$（这与 Topologies,
Definition \ref{topologies-definition-restriction-small-etale}
中的记号并不冲突）。

\begin{lemma}
\label{etale-cohomology-lemma-compare-injectives}
设 $S$ 是概形。设 $T$ 是 $(\Sch/S)_\etale$ 的对象。
\begin{enumerate}
\item 如果 $\mathcal{I}$ 在 $\textit{Ab}((\Sch/S)_\etale)$ 中是内射对象，则
\begin{enumerate}
\item $i_f^{-1}\mathcal{I}$ 在 $\textit{Ab}(T_\etale)$ 中是内射对象；
\item $\mathcal{I}|_{S_\etale}$ 在 $\textit{Ab}(S_\etale)$ 中是内射对象。
\end{enumerate}
\item 如果 $\mathcal{I}^\bullet$ 是
$\textit{Ab}((\Sch/S)_\etale)$ 中的 K-内射复形，则
\begin{enumerate}
\item $i_f^{-1}\mathcal{I}^\bullet$ 是
$\textit{Ab}(T_\etale)$ 中的 K-内射复形；
\item $\mathcal{I}^\bullet|_{S_\etale}$ 是
$\textit{Ab}(S_\etale)$ 中的 K-内射复形。
\end{enumerate}
\end{enumerate}
对模而言，相应断言并不成立。
\end{lemma}

\begin{proof}
(1)(b) 与 (2)(b) 形式地来自如下事实：限制函子
$\pi_{S, *} = i_S^{-1}$ 是正合函子 $\pi_S^{-1}$ 的右伴随。
见 Homology, Lemma \ref{homology-lemma-adjoint-preserve-injectives} 和
Derived Categories, Lemma
\ref{derived-lemma-adjoint-preserve-K-injectives}。

\medskip\noindent
(1)(a) 与 (2)(a) 有两种证法。第一种证法：可以利用
$i_f^{-1}$ 是正合函子 $i_{f, !}$ 的右伴随。对集合层，该函子构造于
Topologies, Lemma \ref{topologies-lemma-put-in-T-etale}；对阿贝尔层，
构造于 Modules on Sites, Lemma
\ref{sites-modules-lemma-g-shriek-adjoint}。Modules on Sites, Lemma
\ref{sites-modules-lemma-exactness-lower-shriek} 证明了它是正合的。
第二种证法：由 Topologies, Lemma
\ref{topologies-lemma-morphism-big-small-etale} 可知
$i_f = i_T \circ f_{big}$。由于 $f_{big}$ 是局部化，沿它的拉回保持
内射对象与 K-内射对象，见 Cohomology on Sites, Lemmas
\ref{sites-cohomology-lemma-cohomology-of-open} 和
\ref{sites-cohomology-lemma-restrict-K-injective-to-open}。
然后对函子 $i_T^{-1}$ 应用已经证明的 (1)(b) 与 (2)(b)，即可得出结论。

\medskip\noindent
令 $S = \Spec(\mathbf{Z})$，并考察映射
$2 : \mathcal{O}_S \to \mathcal{O}_S$。这是 $\mathcal{O}_S$-模的单射，
定义在 $S_\etale$ 上。然而，拉回
$\pi_S^*(2) : \mathcal{O} \to \mathcal{O}$ 不是单射；在
$\Spec(\mathbf{F}_2)$ 上取值即可看出这一点。现在选取单射
$\alpha : \mathcal{O} \to \mathcal{I}$，其中 $\mathcal{O}$-模
$\mathcal{I}$ 是内射的，并定义在 $(\Sch/S)_\etale$ 上。
考察图
$$
\xymatrix{
\mathcal{O}_S \ar[d]_2 \ar[rr]_{\alpha|_{S_\etale}} & &
\mathcal{I}|_{S_\etale} \\
\mathcal{O}_S \ar@{..>}[rru]
}
$$
虚线箭头不可能存在于 $\mathcal{O}_S$-模范畴中，因为由伴随性，
它将意味着单射 $\alpha$ 通过非单射 $\pi_S^*(2)$ 分解，这是不可能的。
因此 $\mathcal{I}|_{S_\etale}$ 不是内射 $\mathcal{O}_S$-模。
\end{proof}

\noindent
设 $f : T \to S$ 是概形态射。Topologies, Lemma
\ref{topologies-lemma-morphism-big-small-etale} (3) 中的交换图诱导出
下列环化网站交换图：
$$
\xymatrix{
(T_\etale, \mathcal{O}_T) \ar[d]_{f_{small}} &
((\Sch/T)_\etale, \mathcal{O}) \ar[d]^{f_{big}} \ar[l]^{\pi_T} \\
(S_\etale, \mathcal{O}_S) &
((\Sch/S)_\etale, \mathcal{O}) \ar[l]_{\pi_S}
}
$$
把 $f_{small}^\sharp$、$f_{big}^\sharp$、$\pi_S^\sharp$ 和
$\pi_T^\sharp$ 的定义写出来，就很容易看出这一点。特别地，这意味着
\begin{equation}
\label{etale-cohomology-equation-compare-big-small}
(f_{big, *}\mathcal{F})|_{S_\etale} =
f_{small, *}(\mathcal{F}|_{T_\etale})
\end{equation}
对任意层 $\mathcal{F}$（定义在 $(\Sch/T)_\etale$ 上）都成立；如果
$\mathcal{F}$ 是 $\mathcal{O}$-模层，则
(\ref{etale-cohomology-equation-compare-big-small}) 是 $\mathcal{O}_S$-模的同构，
定义在 $S_\etale$ 上。

\begin{lemma}
\label{etale-cohomology-lemma-compare-higher-direct-image}
设 $f : T \to S$ 是概形态射。
\begin{enumerate}
\item 对 $K$（属于 $D((\Sch/T)_\etale)$），有
$
(Rf_{big, *}K)|_{S_\etale} = Rf_{small, *}(K|_{T_\etale})
$
，该等式位于 $D(S_\etale)$ 中。
\item 对 $K$（属于 $D((\Sch/T)_\etale, \mathcal{O})$），有
$
(Rf_{big, *}K)|_{S_\etale} = Rf_{small, *}(K|_{T_\etale})
$
，该等式位于 $D(\textit{Mod}(S_\etale, \mathcal{O}_S))$ 中。
\end{enumerate}
更一般地，设 $g : S' \to S$ 是 $(\Sch/S)_\etale$ 的对象。考察纤维积
$$
\xymatrix{
T' \ar[r]_{g'} \ar[d]_{f'} & T \ar[d]^f \\
S' \ar[r]^g & S
}
$$
则
\begin{enumerate}
\item[(3)] 对 $K$（属于 $D((\Sch/T)_\etale)$），有
$i_g^{-1}(Rf_{big, *}K) = Rf'_{small, *}(i_{g'}^{-1}K)$
，该等式位于 $D(S'_\etale)$ 中。
\item[(4)] 对 $K$（属于 $D((\Sch/T)_\etale, \mathcal{O})$），有
$i_g^*(Rf_{big, *}K) = Rf'_{small, *}(i_{g'}^*K)$
，该等式位于 $D(\textit{Mod}(S'_\etale, \mathcal{O}_{S'}))$ 中。
\item[(5)] 对 $K$（属于 $D((\Sch/T)_\etale)$），有
$g_{big}^{-1}(Rf_{big, *}K) = Rf'_{big, *}((g'_{big})^{-1}K)$
，该等式位于 $D((\Sch/S')_\etale)$ 中。
\item[(6)] 对 $K$（属于 $D((\Sch/T)_\etale, \mathcal{O})$），有
$g_{big}^*(Rf_{big, *}K) = Rf'_{big, *}((g'_{big})^*K)$
，该等式位于 $D(\textit{Mod}(S'_\etale, \mathcal{O}_{S'}))$ 中。
\end{enumerate}
\end{lemma}

\begin{proof}
选取表示 $K$ 的阿贝尔层 K-内射复形后，(1) 由
Lemma \ref{etale-cohomology-lemma-compare-injectives} 和
(\ref{etale-cohomology-equation-compare-big-small}) 得到。

\medskip\noindent
选取表示 $K$ 的阿贝尔层 K-内射复形后，(3) 由
Lemma \ref{etale-cohomology-lemma-compare-injectives} 和 Topologies, Lemma
\ref{topologies-lemma-morphism-big-small-cartesian-diagram-etale} 得到。

\medskip\noindent
(5) 就是 Cohomology on Sites, Lemma
\ref{sites-cohomology-lemma-localize-cartesian-square}.

\medskip\noindent
(6) 就是 Cohomology on Sites, Lemma
\ref{sites-cohomology-lemma-localize-cartesian-square-modules}.

\medskip\noindent
(2) 可证如下。上面已经看到，作为环化网站的态射，有
$\pi_S \circ f_{big} = f_{small} \circ \pi_T$。因此由
Cohomology on Sites，引理
\ref{sites-cohomology-lemma-derived-pushforward-composition} 得到
$R\pi_{S, *} \circ Rf_{big, *} = Rf_{small, *} \circ R\pi_{T, *}$
。由于限制函子 $\pi_{S, *}$ 与 $\pi_{T, *}$ 都正合，结论随即得到。

\medskip\noindent
(4) 由 (6) 以及把 (2) 应用于 $f' : T' \to S'$ 得到。
\end{proof}

\noindent
设 $S$ 是概形，并设 $\mathcal{H}$ 是 $(\Sch/S)_\etale$ 上的阿贝尔层。
回顾，$H^n_\etale(U, \mathcal{H})$ 表示 $\mathcal{H}$ 在对象
$U$ 上的上同调，其中该对象属于 $(\Sch/S)_\etale$。

\begin{lemma}
\label{etale-cohomology-lemma-compare-cohomology}
设 $f : T \to S$ 是概形态射。则
\begin{enumerate}
\item 对 $K$（属于 $D(S_\etale)$），有
$H^n_\etale(S, \pi_S^{-1}K) = H^n(S_\etale, K)$。
\item 对 $K$（属于 $D(S_\etale, \mathcal{O}_S)$），有
$H^n_\etale(S, L\pi_S^*K) = H^n(S_\etale, K)$。
\item 对 $K$（属于 $D(S_\etale)$），有
$H^n_\etale(T, \pi_S^{-1}K) = H^n(T_\etale, f_{small}^{-1}K)$。
\item 对 $K$（属于 $D(S_\etale, \mathcal{O}_S)$），有
$H^n_\etale(T, L\pi_S^*K) = H^n(T_\etale, Lf_{small}^*K)$。
\item 对 $M$（属于 $D((\Sch/S)_\etale)$），有
$H^n_\etale(T, M) = H^n(T_\etale, i_f^{-1}M)$。
\item 对 $M$（属于 $D((\Sch/S)_\etale, \mathcal{O})$），有
$H^n_\etale(T, M) = H^n(T_\etale, i_f^*M)$。
\end{enumerate}
\end{lemma}

\begin{proof}
为证明 (5)，用阿贝尔层 K-内射复形表示 $M$，应用
Lemma \ref{etale-cohomology-lemma-compare-injectives}，再展开定义。由于
$i_f^{-1}\pi_S^{-1} = f_{small}^{-1}$，(3) 由此得到。(1) 是 (3) 的特例。

\medskip\noindent
(6) 由非常一般的 Cohomology on Sites, Lemma
\ref{sites-cohomology-lemma-pullback-same-cohomology} 得到。又因为
$Lf_{small}^* = i_f^* \circ L\pi_S^*$，所以 (4) 也随之得到。
(2) 是 (4) 的特例。
\end{proof}

\begin{lemma}
\label{etale-cohomology-lemma-cohomological-descent-etale}
设 $S$ 是概形。对 $K \in D(S_\etale)$，映射
$$
K \longrightarrow R\pi_{S, *}\pi_S^{-1}K
$$
是同构。
\end{lemma}

\begin{proof}
这是因为 $\pi_S^{-1}$ 与 $\pi_{S, *} = i_S^{-1}$ 都是正合函子，
而复合 $\pi_{S, *} \circ \pi_S^{-1}$ 是恒等函子。
\end{proof}

\begin{lemma}
\label{etale-cohomology-lemma-compare-higher-direct-image-proper}
设 $f : T \to S$ 是概形的固有态射。则有
\begin{enumerate}
\item $\pi_S^{-1} \circ f_{small, *} = f_{big, *} \circ \pi_T^{-1}$
，这是函子 $\Sh(T_\etale) \to \Sh((\Sch/S)_\etale)$ 的等式；
\item $\pi_S^{-1}Rf_{small, *}K = Rf_{big, *}\pi_T^{-1}K$
，其中 $K$ 属于 $D^+(T_\etale)$，且其上同调层为挠层；
\item $\pi_S^{-1}Rf_{small, *}K = Rf_{big, *}\pi_T^{-1}K$
，其中 $K$ 属于 $D(T_\etale, \mathbf{Z}/n\mathbf{Z})$；
\item $\pi_S^{-1}Rf_{small, *}K = Rf_{big, *}\pi_T^{-1}K$
，这对所有 $K$（属于 $D(T_\etale)$）都成立，只要 $f$ 有限。
\end{enumerate}
\end{lemma}

\begin{proof}
证明 (1)。设 $\mathcal{F}$ 是 $T_\etale$ 上的层。设
$g : S' \to S$ 是 $(\Sch/S)_\etale$ 的对象。考察纤维积
$$
\xymatrix{
T' \ar[r]_{f'} \ar[d]_{g'} & S' \ar[d]^g \\
T \ar[r]^f & S
}
$$
则有
$$
(f_{big, *}\pi_T^{-1}\mathcal{F})(S') =
(\pi_T^{-1}\mathcal{F})(T') =
((g'_{small})^{-1}\mathcal{F})(T')  =
(f'_{small, *}(g'_{small})^{-1}\mathcal{F})(S')
$$
，其中第二个等式来自 Lemma \ref{etale-cohomology-lemma-describe-pullback}。另一方面，
$$
(\pi_S^{-1}f_{small, *}\mathcal{F})(S') =
(g_{small}^{-1}f_{small, *}\mathcal{F})(S')
$$
，这里仍使用了 Lemma \ref{etale-cohomology-lemma-describe-pullback}。因此，由集合层的
固有基变换（Lemma \ref{etale-cohomology-lemma-proper-base-change-f-star}），
这两个集合典范同构。该同构与限制映射相容，并定义同构
$\pi_S^{-1}f_{small, *}\mathcal{F} = f_{big, *}\pi_T^{-1}\mathcal{F}$。
于是得到函子的同构
$\pi_S^{-1} \circ f_{small, *} = f_{big, *} \circ \pi_T^{-1}$。

\medskip\noindent
证明 (2)。存在典范基变换映射
$\pi_S^{-1}Rf_{small, *}K \to Rf_{big, *}\pi_T^{-1}K$，它对任意
$K$（属于 $D(T_\etale)$）都有定义，见 Cohomology on Sites, Remark
\ref{sites-cohomology-remark-base-change}。为证明它是同构，只须证明：
基变换映射沿
$i_g : \Sh(S'_\etale) \to \Sh((\Sch/S)_\etale)$ 的拉回是同构，
这对任意对象 $g : S' \to S$（属于 $(\Sch/S)_\etale$）都成立。
令 $T', g', f'$ 如上一段所述。基变换映射的拉回为
\begin{align*}
g_{small}^{-1}Rf_{small, *}K
& =
i_g^{-1}\pi_S^{-1}Rf_{small, *}K \\
& \to
i_g^{-1}Rf_{big, *}\pi_T^{-1}K \\
& =
Rf'_{small, *}(i_{g'}^{-1}\pi_T^{-1}K) \\
& =
Rf'_{small, *}((g'_{small})^{-1}K)
\end{align*}
；这里使用了 $\pi_S \circ i_g = g_{small}$、
$\pi_T \circ i_{g'} = g'_{small}$，以及
Lemma \ref{etale-cohomology-lemma-compare-higher-direct-image}。如果 $K$ 下有界且
$K$ 的上同调层为挠层，则由固有基变换定理
（Lemma \ref{etale-cohomology-lemma-proper-base-change}），该映射是同构。

\medskip\noindent
(3) 的证明与 (2) 完全相同，只须用 Lemma
\ref{etale-cohomology-lemma-proper-base-change-mod-n} 代替 Lemma
\ref{etale-cohomology-lemma-proper-base-change}。

\medskip\noindent
证明 (4)。如果 $f$ 有限，则函子 $f_{small, *}$ 与 $f_{big, *}$
都是正合的。对 $f_{small}$，这由 Proposition
\ref{etale-cohomology-proposition-finite-higher-direct-image-zero} 得到。由于任意基变换
$f'$（来自 $f$）也有限，Lemma
\ref{etale-cohomology-lemma-compare-higher-direct-image} (3) 说明 $f_{big, *}$
也正合（因为高阶导出函子为零）。所以这种情形由 (1) 得到。
\end{proof}




\section{fppf 拓扑与 \'etale 拓扑的比较}
\label{etale-cohomology-section-fppf-etale}

\noindent
本节的一个范本，是 Cohomology on Sites, Section
\ref{sites-cohomology-section-cohomology-LC} 中讨论的局部紧拓扑空间上通常拓扑与
qc 拓扑的比较。我们先回顾 Topologies, Sections
\ref{topologies-section-change-topologies} 和
\ref{topologies-section-etale} 中的一些材料。

\medskip\noindent
设 $S$ 是概形，并设 $(\Sch/S)_{fppf}$ 是一个 fppf 网站。
在同一个底层范畴上还有第二个拓扑，即 \'etale 拓扑，因而还有第二个网站
$(\Sch/S)_\etale$。恒等函子
$(\Sch/S)_\etale \to (\Sch/S)_{fppf}$ 是连续的，并定义网站态射
$$
\epsilon_S : (\Sch/S)_{fppf} \longrightarrow (\Sch/S)_\etale
$$
。见 Cohomology on Sites, Section \ref{sites-cohomology-section-compare}。
请注意，$\epsilon_{S, *}$ 在底层预层上是恒等函子，而
$\epsilon_S^{-1}$ 把一个 \'etale 层变为它的 fppf 层化。令
$S_\etale$ 为小 \'etale 网站。存在网站态射
$$
\pi_S : (\Sch/S)_\etale \longrightarrow S_\etale
$$
，它由连续函子
$S_\etale \to (\Sch/S)_\etale$, $U \mapsto U$ 给出。
事实上，$S_\etale$ 有纤维积和终对象，上述函子与它们可交换，因而可以应用
Sites, Proposition \ref{sites-proposition-get-morphism}。

\begin{lemma}
\label{etale-cohomology-lemma-describe-pullback-pi-fppf}
沿用上述记号。设 $\mathcal{F}$ 是 $S_\etale$ 上的层。规则
$$
(\Sch/S)_{fppf} \longrightarrow \textit{Sets},\quad
(f : X \to S) \longmapsto \Gamma(X, f_{small}^{-1}\mathcal{F})
$$
给出一个层，因而当然也给出 $(\Sch/S)_\etale$ 上的层。
事实上，这个层等于 $\pi_S^{-1}\mathcal{F}$（在
$(\Sch/S)_\etale$ 上），而等于
$\epsilon_S^{-1}\pi_S^{-1}\mathcal{F}$（在 $(\Sch/S)_{fppf}$ 上）。
\end{lemma}

\begin{proof}
关于 \'etale 拓扑的断言就是 Lemma \ref{etale-cohomology-lemma-describe-pullback} 的内容。
为完成证明，只须说明 $\pi_S^{-1}\mathcal{F}$ 对 fppf 拓扑是层。
这也已在 Lemma \ref{etale-cohomology-lemma-describe-pullback} 中证明。
\end{proof}

\noindent
在 Lemma \ref{etale-cohomology-lemma-describe-pullback-pi-fppf} 的情形中，
$\epsilon_S$ 与 $\pi_S$ 的复合以及上述等式确定网站态射
$$
a_S : (\Sch/S)_{fppf} \longrightarrow S_\etale
$$

\begin{lemma}
\label{etale-cohomology-lemma-push-pull-fppf-etale}
沿用上述记号。设 $f : X \to Y$ 是 $(\Sch/S)_{fppf}$ 中的态射。
则有拓扑斯的交换图
$$
\xymatrix{
\Sh((\Sch/X)_{fppf}) \ar[rr]_{f_{big, fppf}} \ar[d]_{\epsilon_X} & &
\Sh((\Sch/Y)_{fppf}) \ar[d]^{\epsilon_Y} \\
\Sh((\Sch/X)_\etale) \ar[rr]^{f_{big, \etale}} & &
\Sh((\Sch/Y)_\etale)
}
$$
以及
$$
\xymatrix{
\Sh((\Sch/X)_{fppf}) \ar[rr]_{f_{big, fppf}} \ar[d]_{a_X} & &
\Sh((\Sch/Y)_{fppf}) \ar[d]^{a_Y} \\
\Sh(X_\etale) \ar[rr]^{f_{small}} & &
\Sh(Y_\etale)
}
$$
，其中 $a_X = \pi_X \circ \epsilon_X$ 且 $a_Y = \pi_X \circ \epsilon_X$。
\end{lemma}

\begin{proof}
这些图的交换性来自 Topologies, Section
\ref{topologies-section-change-topologies} 中的讨论。
\end{proof}

\begin{lemma}
\label{etale-cohomology-lemma-proper-push-pull-fppf-etale}
在 Lemma \ref{etale-cohomology-lemma-push-pull-fppf-etale} 中，如果 $f$ 是固有的，则有
$a_Y^{-1} \circ f_{small, *} = f_{big, fppf, *} \circ a_X^{-1}$.
\end{lemma}

\begin{proof}
可以重复 Lemma \ref{etale-cohomology-lemma-compare-higher-direct-image-proper} (1) 的证明；
这里我们改为由该结论推出本结果。由于 $\epsilon_{Y, *}$ 在底层预层上是
恒等函子，它反映同构。Lemma \ref{etale-cohomology-lemma-describe-pullback-pi-fppf}
中的描述说明 $\epsilon_{Y, *} \circ a_Y^{-1} = \pi_Y^{-1}$，
对 $X$ 也类似。为证明典范映射
$a_Y^{-1}f_{small, *}\mathcal{F} \to f_{big, fppf, *}a_X^{-1}\mathcal{F}$
是同构，只须证明
\begin{align*}
\pi_Y^{-1}f_{small, *}\mathcal{F}
& =
\epsilon_{Y, *}a_Y^{-1}f_{small, *}\mathcal{F} \\
& \to 
\epsilon_{Y, *}f_{big, fppf, *}a_X^{-1}\mathcal{F} \\
& =
f_{big, \etale, *} \epsilon_{X, *}a_X^{-1}\mathcal{F} \\
& =
f_{big, \etale, *}\pi_X^{-1}\mathcal{F}
\end{align*}
是同构。这就是 Lemma \ref{etale-cohomology-lemma-compare-higher-direct-image-proper} (1)。
\end{proof}

\begin{lemma}
\label{etale-cohomology-lemma-descent-sheaf-fppf-etale}
在 Lemma \ref{etale-cohomology-lemma-push-pull-fppf-etale} 中，假设 $f$ 平坦、局部有限表现且满射。
则函子
$$
\Sh(Y_\etale) \longrightarrow
\left\{
(\mathcal{G}, \mathcal{H}, \alpha)
\middle|
\begin{matrix}
\mathcal{G} \in \Sh(X_\etale),\ \mathcal{H} \in \Sh((\Sch/Y)_{fppf}), \\
\alpha : a_X^{-1}\mathcal{G} \to f_{big, fppf}^{-1}\mathcal{H}
\text{ 为同构}
\end{matrix}
\right\}
$$
把 $\mathcal{F}$ 送到
$(f_{small}^{-1}\mathcal{F}, a_Y^{-1}\mathcal{F}, can)$，它是一个等价。
\end{lemma}

\begin{proof}
函子 $a_X^{-1}$ 完全忠实（由 Lemma
\ref{etale-cohomology-lemma-describe-pullback-pi-fppf}，有
$a_{X, *}a_X^{-1} = \text{id}$）。因此遗忘函子
$(\mathcal{G}, \mathcal{H}, \alpha) \mapsto \mathcal{H}$ 把三元组范畴
等同于 $\Sh((\Sch/Y)_{fppf})$ 的一个全子范畴。此外，函子
$a_Y^{-1}$ 完全忠实，故引理中的函子也完全忠实。

\medskip\noindent
假设有一个 \'etale 覆盖 $\{Y_i \to Y\}$。令
$f_i : X_i \to Y_i$ 为 $f$ 的基变换，并记
$f_{ij} = f_i \times f_j : X_i \times_X X_j  \to Y_i \times_Y Y_j$。
我们断言：如果引理对 $f_i$ 和 $f_{ij}$ 都成立（对所有 $i, j$），
那么它对 $f$ 也成立。为说明这一点，注意给定的 \'etale 覆盖在四个网站
$Y_\etale, X_\etale, (\Sch/Y)_{fppf}, (\Sch/X)_{fppf}$ 中都确定了
终对象的一个 \'etale 覆盖。因此在四种情形的每一种中，层范畴都等价于
这个覆盖的黏合数据范畴（Sites, Lemma
\ref{sites-lemma-mapping-property-glue}）。随后，一个庞大的范畴交换图完成
该断言的证明；我们略去细节。该断言表明，可以在 $Y$ 上 \'etale 局部地工作。

\medskip\noindent
注意，$\{X \to Y\}$ 是一个 fppf 覆盖。在 $Y$ 上 \'etale 局部地工作，
可以假设存在态射 $s : X' \to X$，使得复合
$f' = f \circ s : X' \to Y$ 满射、有限且局部自由，见
More on Morphisms，引理
\ref{more-morphisms-lemma-dominate-fppf-etale-locally}。
我们断言：如果引理对 $f'$ 成立，那么它对 $f$ 也成立。
事实上，给定一个三元组 $(\mathcal{G}, \mathcal{H}, \alpha)$（针对 $f$），
沿 $s$ 拉回便得到三元组
$(s_{small}^{-1}\mathcal{G}, \mathcal{H}, s_{big, fppf}^{-1}\alpha)$（针对 $f'$）。
这个三元组的一个解给出层 $\mathcal{F}$（定义在 $Y_\etale$ 上），满足
$a_Y^{-1}\mathcal{F} = \mathcal{H}$。由证明的第一段，这意味着该三元组
属于本质像。于是归约到下一段所述的情形。

\medskip\noindent
假设 $f$ 满射、有限且局部自由。设
$(\mathcal{G}, \mathcal{H}, \alpha)$ 是一个三元组。在此情形考虑三元组
$$
(\mathcal{G}_1, \mathcal{H}_1, \alpha_1) =
(f_{small}^{-1}f_{small, *}\mathcal{G},
f_{big, fppf, *}f_{big, fppf}^{-1}\mathcal{H}, \alpha_1)
$$
，其中 $\alpha_1$ 来自如下等同：
\begin{align*}
a_X^{-1}f_{small}^{-1}f_{small, *}\mathcal{G}
& =
f_{big, fppf}^{-1}a_Y^{-1}f_{small, *}\mathcal{G} \\
& =
f_{big, fppf}^{-1}f_{big, fppf, *}a_X^{-1}\mathcal{G} \\
& \to
f_{big, fppf}^{-1}f_{big, fppf, *}f_{big, fppf}^{-1}\mathcal{H}
\end{align*}
；其中第三个等式是 Lemma \ref{etale-cohomology-lemma-proper-push-pull-fppf-etale}，
箭头由 $\alpha$ 给出。这个三元组属于我们函子的像，因为
$\mathcal{F}_1 = f_{small, *}\mathcal{F}$ 是一个解（为此再次使用
Lemma \ref{etale-cohomology-lemma-proper-push-pull-fppf-etale}；细节从略）。存在三元组的典范映射
$$
(\mathcal{G}, \mathcal{H}, \alpha)
\to
(\mathcal{G}_1, \mathcal{H}_1, \alpha_1)
$$
，它在第二个分量上使用单位
$\text{id} \to f_{big, fppf, *}f_{big, fppf}^{-1}$（由证明的第一段，
只须在第二个分量上指定态射）。由于 $\{f : X \to Y\}$ 是 fppf 覆盖，
映射 $\mathcal{H} \to \mathcal{H}_1$ 是单射（细节从略）。置
$$
\mathcal{G}_2 = \mathcal{G}_1 \amalg_\mathcal{G} \mathcal{G}_1\quad
\mathcal{H}_2 = \mathcal{H}_1 \amalg_\mathcal{H} \mathcal{H}_1
$$
，并令 $\alpha_2$ 为诱导的同构（拉回函子是正合的，因此这是有意义的）。
于是 $\mathcal{H}$ 是两个映射
$\mathcal{H}_1 \to \mathcal{H}_2$ 的等化子。对三元组
$(\mathcal{G}_2, \mathcal{H}_2, \alpha_2)$ 重复上述论证，可得三元组的单态射
$$
(\mathcal{G}_2, \mathcal{H}_2, \alpha_2)
\to
(\mathcal{G}_3, \mathcal{H}_3, \alpha_3)
$$
，使得最后这个三元组属于我们函子的像。设它对应于
$\mathcal{F}_3$（属于 $\Sh(Y_\etale)$）。由完全忠实性，得到两个映射
$\mathcal{F}_1 \to \mathcal{F}_3$；令 $\mathcal{F}$ 为这两个映射的等化子。
由所涉及拉回函子的正合性，便有
$a_Y^{-1}\mathcal{F} = \mathcal{H}$，如所要求。
\end{proof}

\begin{lemma}
\label{etale-cohomology-lemma-compare-fppf-etale}
考虑比较态射
$\epsilon : (\Sch/S)_{fppf} \to (\Sch/S)_\etale$。
令 $\mathcal{P}$ 表示概形的有限态射类。对
$X$（属于 $(\Sch/S)_\etale$），记
$\mathcal{A}'_X \subset \textit{Ab}((\Sch/X)_\etale)$
为由形如 $\pi_X^{-1}\mathcal{F}$ 的层构成的全子范畴，其中
$\mathcal{F}$ 属于 $\textit{Ab}(X_\etale)$。则 Cohomology on Sites 中的性质
(\ref{sites-cohomology-item-base-change-P}),
(\ref{sites-cohomology-item-restriction-A}),
(\ref{sites-cohomology-item-A-sheaf}),
(\ref{sites-cohomology-item-A-and-P})，以及
(\ref{sites-cohomology-item-refine-tau-by-P})
（见 Cohomology on Sites, Situation
\ref{sites-cohomology-situation-compare}）都成立。
\end{lemma}

\begin{proof}
我们先验证 Homology, Lemma
\ref{homology-lemma-characterize-weak-serre-subcategory} 的条件
(1)、(2)、(3)、(4)，从而说明
$\mathcal{A}'_X \subset \textit{Ab}((\Sch/X)_\etale)$ 是弱 Serre 子范畴。
(1)、(2)、(3) 是立即的，因为例如由 Lemma
\ref{etale-cohomology-lemma-cohomological-descent-etale}，$\pi_X^{-1}$ 正合且完全忠实。
如果
$0 \to \pi_X^{-1}\mathcal{F} \to \mathcal{G} \to \pi_X^{-1}\mathcal{F}' \to 0$
是 $\textit{Ab}((\Sch/X)_\etale)$ 中的短正合列，则由 Lemma
\ref{etale-cohomology-lemma-cohomological-descent-etale}，
$0 \to \mathcal{F} \to \pi_{X, *}\mathcal{G} \to \mathcal{F}' \to 0$
正合。因此 $\mathcal{G} = \pi_X^{-1}\pi_{X, *}\mathcal{G}$ 属于
$\mathcal{A}'_X$，这就验证了最后一个条件。

\medskip\noindent
Cohomology on Sites, Property (\ref{sites-cohomology-item-base-change-P})
成立，因为概形有纤维积，而且概形的有限态射经基变换后仍是有限态射；
见 Morphisms, Lemma \ref{morphisms-lemma-base-change-finite}。

\medskip\noindent
Cohomology on Sites, Property (\ref{sites-cohomology-item-restriction-A})
来自 Topologies, Lemma
\ref{topologies-lemma-morphism-big-small-etale} 中的交换图 (3)。

\medskip\noindent
Cohomology on Sites, Property (\ref{sites-cohomology-item-A-sheaf})
就是 Lemma \ref{etale-cohomology-lemma-describe-pullback-pi-fppf}。

\medskip\noindent
Cohomology on Sites, Property (\ref{sites-cohomology-item-A-and-P})
由 Lemma \ref{etale-cohomology-lemma-compare-higher-direct-image-proper} (4) 成立。

\medskip\noindent
Cohomology on Sites, Property (\ref{sites-cohomology-item-refine-tau-by-P})
由 More on Morphisms, Lemma
\ref{more-morphisms-lemma-dominate-fppf-etale-locally} 推出。
\end{proof}

\begin{lemma}
\label{etale-cohomology-lemma-V-C-all-n-etale-fppf}
沿用上述记号。
\begin{enumerate}
\item 对 $X \in \Ob((\Sch/S)_{fppf})$ 和阿贝尔层
$\mathcal{F}$（定义在 $X_\etale$ 上），有
$\epsilon_{X, *}a_X^{-1}\mathcal{F} = \pi_X^{-1}\mathcal{F}$
，并且 $R^i\epsilon_{X, *}(a_X^{-1}\mathcal{F}) = 0$ 对 $i > 0$ 成立。
\item 对有限态射 $f : X \to Y$（属于 $(\Sch/S)_{fppf}$）和
阿贝尔层 $\mathcal{F}$（定义在 $X$ 上），有
$a_Y^{-1}(R^if_{small, *}\mathcal{F}) =
R^if_{big, fppf, *}(a_X^{-1}\mathcal{F})$
，对所有 $i$ 成立。
\item 对概形 $X$ 和 $K$（属于 $D^+(X_\etale)$），映射
$\pi_X^{-1}K \to R\epsilon_{X, *}(a_X^{-1}K)$ 是同构。
\item 对概形的有限态射 $f : X \to Y$ 和 $K$（属于
$D^+(X_\etale)$），有
$a_Y^{-1}(Rf_{small, *}K) = Rf_{big, fppf, *}(a_X^{-1}K)$.
\item 对概形的固有态射 $f : X \to Y$ 和 $K$（属于
$D^+(X_\etale)$，且上同调层为挠层），有
$a_Y^{-1}(Rf_{small, *}K) = Rf_{big, fppf, *}(a_X^{-1}K)$.
\end{enumerate}
\end{lemma}

\begin{proof}
由 Lemma \ref{etale-cohomology-lemma-compare-fppf-etale}，Cohomology on Sites, Section
\ref{sites-cohomology-section-compare-general} 中的各引理都适用于当前情形。
为翻译这些结果，注意 Cohomology on Sites, Lemma
\ref{sites-cohomology-lemma-A} 中的范畴 $\mathcal{A}_X$，正是
$a_X^{-1} : \textit{Ab}(X_\etale) \to \textit{Ab}((\Sch/X)_{fppf})$
的本质像。

\medskip\noindent
(1) 等价于 $(V_n)$ 对所有 $n$ 都成立；这由 Cohomology on Sites, Lemma
\ref{sites-cohomology-lemma-V-C-all-n-general} 成立。

\medskip\noindent
(2) 来自对 Cohomology on Sites, Lemma
\ref{sites-cohomology-lemma-V-implies-C-general} 的结论应用
$\epsilon_Y^{-1}$。

\medskip\noindent
(3) 来自 Cohomology on Sites, Lemma
\ref{sites-cohomology-lemma-V-C-all-n-general} (1)，因为
$\pi_X^{-1}K$ 属于 $D^+_{\mathcal{A}'_X}((\Sch/X)_\etale)$，并且
$a_X^{-1} = \epsilon_X^{-1} \circ a_X^{-1}$。

\medskip\noindent
出于同样的理由，(4) 来自 Cohomology on Sites, Lemma
\ref{sites-cohomology-lemma-V-C-all-n-general} (2)。

\medskip\noindent
证明 (5)。我们使用
\begin{align*}
R\epsilon_{Y, *}Rf_{big, fppf, *}a_X^{-1}K
& =
Rf_{big, \etale, *}R\epsilon_{X, *}a_X^{-1}K \\
& =
Rf_{big, \etale, *}\pi_X^{-1}K \\
& =
\pi_Y^{-1}Rf_{small, *}K \\
& =
R\epsilon_{Y, *} a_Y^{-1}Rf_{small, *}K
\end{align*}
。第一个等式来自 Lemma \ref{etale-cohomology-lemma-push-pull-fppf-etale} 中的交换图以及
Cohomology on Sites，引理
\ref{sites-cohomology-lemma-derived-pushforward-composition}。
第二个等式是 (3)。第三个等式是 Lemma
\ref{etale-cohomology-lemma-compare-higher-direct-image-proper} (2)。第四个等式再次使用 (3)。
因此基变换映射
$a_Y^{-1}(Rf_{small, *}K) \to Rf_{big, fppf, *}(a_X^{-1}K)$
诱导同构
$$
R\epsilon_{Y, *}a_Y^{-1}Rf_{small, *}K \to
R\epsilon_{Y, *}Rf_{big, fppf, *}a_X^{-1}K
$$
。以下观察完成证明：设映射 $\alpha : a_Y^{-1}L \to M$ 中
$L$ 属于 $D^+(Y_\etale)$，$M$ 属于 $D^+((\Sch/Y)_{fppf})$，并且
$R\epsilon_{Y, *}\alpha$ 是同构，则该映射是同构。事实上，
我们对 $i$ 归纳证明 $H^i(\alpha)$ 是同构。对所有充分小的 $i$，
这成立。若它对 $i \leq i_0$ 成立，那么可知
$R^j\epsilon_{Y, *}H^i(M) = 0$（其中 $j > 0$ 且 $i \leq i_0$）；
这是由 (1) 得出的，因为在这个范围内
$H^i(M) = a_Y^{-1}H^i(L)$。于是由谱序列论证，
$\epsilon_{Y, *}H^{i_0 + 1}(M) = H^{i_0 + 1}(R\epsilon_{Y, *}M)$。
因此
$\epsilon_{Y, *}H^{i_0 + 1}(M) = \pi_Y^{-1}H^{i_0 + 1}(L) =
\epsilon_{Y, *}a_Y^{-1}H^{i_0 + 1}(L)$。
这说明 $H^{i_0 + 1}(\alpha)$ 是同构（因为 $\epsilon_{Y, *}$
在底层预层上是恒等函子，所以它反映同构），如所要求。
\end{proof}

\begin{lemma}
\label{etale-cohomology-lemma-cohomological-descent-etale-fppf}
设 $X$ 是概形。对 $K \in D^+(X_\etale)$，映射
$$
K \longrightarrow Ra_{X, *}a_X^{-1}K
$$
是同构，其中
$a_X : \Sh((\Sch/X)_{fppf}) \to \Sh(X_\etale)$ 如上所述。
\end{lemma}

\begin{proof}
我们先把断言归约到 $K$ 由单个阿贝尔层给出的情形。也就是说，取
$K$ 的一个下有界复形表示 $\mathcal{F}^\bullet$。由单个层的情形可知
$\mathcal{F}^n = a_{X, *} a_X^{-1} \mathcal{F}^n$
，而且层 $R^qa_{X, *}a_X^{-1}\mathcal{F}^n$ 对 $q > 0$ 为零。
对 $a_X^{-1}\mathcal{F}^\bullet$ 和函子 $a_{X, *}$ 应用 Leray 无环性引理
（Derived Categories, Lemma \ref{derived-lemma-leray-acyclicity}），
便得到结论。以下假设 $K = \mathcal{F}$。

\medskip\noindent
由 Lemma \ref{etale-cohomology-lemma-describe-pullback-pi-fppf}，有
$a_{X, *}a_X^{-1}\mathcal{F} = \mathcal{F}$。因此只须证明
$R^qa_{X, *}a_X^{-1}\mathcal{F} = 0$ 对 $q > 0$ 成立。为此可以使用
$a_X = \epsilon_X \circ \pi_X$ 和 Leray 谱序列
（Cohomology on Sites, Lemma \ref{sites-cohomology-lemma-relative-Leray}）。
由 Lemma \ref{etale-cohomology-lemma-V-C-all-n-etale-fppf}，有
$R^i\epsilon_{X, *}(a_X^{-1}\mathcal{F}) = 0$（当 $i > 0$），并且
$\epsilon_{X, *}a_X^{-1}\mathcal{F} = \pi_X^{-1}\mathcal{F}$。
由 Lemma \ref{etale-cohomology-lemma-cohomological-descent-etale}，有
$R^j\pi_{X, *}(\pi_X^{-1}\mathcal{F}) = 0$（当 $j > 0$）。证明完毕。
\end{proof}

\begin{lemma}
\label{etale-cohomology-lemma-compare-cohomology-etale-fppf}
设 $X$ 是概形，并设
$a_X : \Sh((\Sch/X)_{fppf}) \to \Sh(X_\etale)$ 如上。则：
\begin{enumerate}
\item $H^q(X_\etale, \mathcal{F}) = H^q_{fppf}(X, a_X^{-1}\mathcal{F})$
，其中 $\mathcal{F}$ 是 $X_\etale$ 上的阿贝尔层；
\item $H^q(X_\etale, K) = H^q_{fppf}(X, a_X^{-1}K)$，其中
$K \in D^+(X_\etale)$。
\end{enumerate}
例如，如果 $A$ 是阿贝尔群，则
$H^q_\etale(X, \underline{A}) = H^q_{fppf}(X, \underline{A})$.
\end{lemma}

\begin{proof}
由 Lemma \ref{etale-cohomology-lemma-cohomological-descent-etale-fppf} 和
Cohomology on Sites, Remark \ref{sites-cohomology-remark-before-Leray}
即得此结论。
\end{proof}







\section{fppf 拓扑与 \'etale 拓扑的比较：模}
\label{etale-cohomology-section-fppf-etale-modules}

\noindent
我们继续 Section \ref{etale-cohomology-section-fppf-etale} 中的讨论，但在本节简要讨论模层的情形。

\medskip\noindent
设 $S$ 是概形。Section \ref{etale-cohomology-section-fppf-etale} 中引入的网站态射
$\epsilon_S$、$\pi_S$ 及其复合 $a_S$，都有向环化网站态射的自然提升。
第一个记作
$$
\epsilon_S :
((\Sch/S)_{fppf}, \mathcal{O})
\longrightarrow
((\Sch/S)_\etale, \mathcal{O})
$$
。注意，我们可以对结构层使用相同的符号，因为这些层确有相同的底层预层。
第二个是
$$
\pi_S :
((\Sch/S)_\etale, \mathcal{O})
\longrightarrow
(S_\etale, \mathcal{O}_S)
$$
。第三个是态射
$$
a_S :
((\Sch/S)_{fppf}, \mathcal{O})
\longrightarrow
(S_\etale, \mathcal{O}_S)
$$
。我们已经知道，概形 $S$ 上的拟凝聚模范畴与
$(S_\etale, \mathcal{O}_S)$ 上的拟凝聚模范畴相同，见
Descent, Proposition \ref{descent-proposition-equivalence-quasi-coherent}。
由于我们关心的是陈述 \'etale 上同调与 fppf 上同调之间的比较，
在本节余下部分，我们将通过小 \'etale 网站来理解拟凝聚层。
下面回顾关于这些网站上拟凝聚模的已有结果。

\begin{lemma}
\label{etale-cohomology-lemma-review-quasi-coherent}
设 $S$ 是概形。设 $\mathcal{F}$ 是拟凝聚 $\mathcal{O}_S$-模
（定义在 $S_\etale$ 上）。
\begin{enumerate}
\item 规则
$$
\mathcal{F}^a : (\Sch/S)_\etale \longrightarrow \textit{Ab},\quad
(f : T \to S) \longmapsto \Gamma(T, f_{small}^*\mathcal{F})
$$
对 fppf 覆盖满足层条件，因而当然也对 \'etale 覆盖满足层条件；
\item $\mathcal{F}^a = \pi_S^*\mathcal{F}$ 在 $(\Sch/S)_\etale$ 上成立；
\item $\mathcal{F}^a = a_S^*\mathcal{F}$ 在 $(\Sch/S)_{fppf}$ 上成立；
\item 规则 $\mathcal{F} \mapsto \mathcal{F}^a$ 定义拟凝聚
$\mathcal{O}_S$-模与
$((\Sch/S)_\etale, \mathcal{O})$ 上拟凝聚模之间的等价；
\item 规则 $\mathcal{F} \mapsto \mathcal{F}^a$ 定义拟凝聚
$\mathcal{O}_S$-模与
$((\Sch/S)_{fppf}, \mathcal{O})$ 上拟凝聚模之间的等价；
\item 有 $\epsilon_{S, *}a_S^*\mathcal{F} = \pi_S^*\mathcal{F}$
和 $a_{S, *}a_S^*\mathcal{F} = \mathcal{F}$；
\item 有 $R^i\epsilon_{S, *}(a_S^*\mathcal{F}) = 0$
和 $R^ia_{S, *}(a_S^*\mathcal{F}) = 0$，其中 $i > 0$。
\end{enumerate}
\end{lemma}

\begin{proof}
我们强烈建议读者根据 Descent, Sections
\ref{descent-section-quasi-coherent-sheaves}、
\ref{descent-section-quasi-coherent-cohomology} 和
\ref{descent-section-quasi-coherent-sheaves-bis} 中的材料，自己找出这些结果的证明。

\medskip\noindent
我们先说明，本引理中的记号 $\mathcal{F}^a$ 为何与我们先前在 Sections
\ref{etale-cohomology-section-quasi-coherent}、\ref{etale-cohomology-section-cohomology-quasi-coherent}
以及 Descent, Section \ref{descent-section-quasi-coherent-sheaves}
中使用的记号相容。事实上，由 Descent, Proposition
\ref{descent-proposition-equivalence-quasi-coherent} 可知，存在一个拟凝聚模
$\mathcal{F}_0$，定义在概形 $S$ 上（换言之，小 Zariski 网站上），
使得 $\mathcal{F}$ 是下述规则的限制：
$$
\mathcal{F}_0^a : (\Sch/S)_\etale \longrightarrow \textit{Ab},\quad
(f : T \to S) \longmapsto \Gamma(T, f^*\mathcal{F})
$$
到子范畴 $S_\etale \subset (\Sch/S)_\etale$；这里 $f^*$ 表示概形上
模层的通常拉回。由于 $\mathcal{F}_0^a$ 是沿环化网站态射
$$
((\Sch/S)_\etale, \mathcal{O}) \longrightarrow (S_{Zar}, \mathcal{O}_{S_{Zar}})
$$
的拉回（Descent, Remark
\ref{descent-remark-change-topologies-ringed-sites}），由拉回的复合立即得到
$\mathcal{F}^a = \mathcal{F}_0^a$。Descent, Lemma
\ref{descent-lemma-sheaf-condition-holds} 于是证明，即使对 fpqc 覆盖也满足
层条件（另见 Proposition \ref{etale-cohomology-proposition-quasi-coherent-sheaf-fpqc}）。
再由 Descent, Remark \ref{descent-remark-change-topologies-ringed-sites}
得到 (2) 和 (3)，由 Descent, Proposition
\ref{descent-proposition-equivalence-quasi-coherent} 得到 (4) 和 (5)
（另见元结果 Theorem \ref{etale-cohomology-theorem-quasi-coherent}）。

\medskip\noindent
(6) 立即来自层 $\mathcal{F}^a = \pi_S^*\mathcal{F} = a_S^*\mathcal{F}$
的描述。

\medskip\noindent
对任意阿贝尔层 $\mathcal{H}$（定义在 $(\Sch/S)_{fppf}$ 上），高阶直接像
$R^p\epsilon_{S, *}\mathcal{H}$ 是预层
$U \mapsto H^p_{fppf}(U, \mathcal{H})$ 所关联的层（定义在
$(\Sch/S)_\etale$ 上）。见 Cohomology on Sites,
Lemma \ref{sites-cohomology-lemma-higher-direct-images}。因此，为证明
$R^p\epsilon_{S, *}a_S^*\mathcal{F} = R^p\epsilon_{S, *}\mathcal{F}^a = 0$
对 $p > 0$ 成立，只须证明任意概形 $U$（在 $S$ 上）都有 \'etale 覆盖
$\{U_i \to U\}_{i \in I}$，使得
$H^p_{fppf}(U_i, \mathcal{F}^a) = 0$ 对 $p > 0$ 成立。
如果取仿射开覆盖，则所需消灭性来自与通常上同调的比较
（Descent, Proposition \ref{descent-proposition-same-cohomology-quasi-coherent}
或 Theorem \ref{etale-cohomology-theorem-zariski-fpqc-quasi-coherent}），以及 Cohomology of
Schemes, Lemma \ref{coherent-lemma-quasi-coherent-affine-cohomology-zero}
所给出的仿射概形上拟凝聚层上同调的消灭性。

\medskip\noindent
为证明 $R^pa_{S, *}a_S^{-1}\mathcal{F} = R^pa_{S, *}\mathcal{F}^a = 0$
对 $p > 0$ 成立，完全同样地论证即可。证明完毕。
\end{proof}

\begin{lemma}
\label{etale-cohomology-lemma-cohomological-descent-etale-fppf-modules}
设 $S$ 是概形。对 $\mathcal{F}$（一个拟凝聚 $\mathcal{O}_S$-模，定义在
$S_\etale$ 上），映射
$$
\pi_S^*\mathcal{F} \longrightarrow R\epsilon_{S, *}(a_S^*\mathcal{F})
\quad\text{且}\quad
\mathcal{F} \longrightarrow Ra_{S, *}(a_S^*\mathcal{F})
$$
都是同构，其中
$a_S : \Sh((\Sch/S)_{fppf}) \to \Sh(S_\etale)$ 如上所述。
\end{lemma}

\begin{proof}
这立即来自 Lemma \ref{etale-cohomology-lemma-review-quasi-coherent} 的 (6) 和 (7)。
\end{proof}

\begin{lemma}
\label{etale-cohomology-lemma-cohomological-descent-complex-modules}
设 $S = \Spec(A)$ 是仿射概形。设 $M^\bullet$ 是 $A$-模复形。
考虑预层复形 $\mathcal{F}^\bullet$，其为 $\mathcal{O}$-模，定义在
$(\textit{Aff}/S)_{fppf}$ 上，并由规则
$$
(U/S) = (\Spec(B)/\Spec(A)) \longmapsto M^\bullet \otimes_A B
$$
给出。这是一个模复形，并且典范映射
$$
M^\bullet \longrightarrow
R\Gamma((\textit{Aff}/S)_{fppf}, \mathcal{F}^\bullet)
$$
是拟同构。
\end{lemma}

\begin{proof}
每个 $\mathcal{F}^n$ 都是模层，因为它等于 Lemma
\ref{etale-cohomology-lemma-review-quasi-coherent} 中的模
$\mathcal{G}^n = (\widetilde{M}^n)^a$ 在
$(\textit{Aff}/S)_{fppf} \subset (\Sch/S)_{fppf}$ 上的限制。
由于这个包含定义环化拓扑斯的等价（Topologies, Lemma
\ref{topologies-lemma-affine-big-site-fppf}），有
$$
R\Gamma((\textit{Aff}/S)_{fppf}, \mathcal{F}^\bullet) =
R\Gamma((\Sch/S)_{fppf}, \mathcal{G}^\bullet)
$$
。注意，例如由 Derived Categories of Schemes, Lemma
\ref{perfect-lemma-affine-compare-bounded}，有
$M^\bullet = R\Gamma(S, \widetilde{M}^\bullet)$。因此要证明比较映射
$$
R\Gamma(S, \widetilde{M}^\bullet)
\longrightarrow
R\Gamma((\Sch/S)_{fppf}, (\widetilde{M}^\bullet)^a)
$$
是同构。如果 $M^\bullet$ 下有界，则这由 Descent, Proposition
\ref{descent-proposition-same-cohomology-quasi-coherent} 和 Derived
Categories, Lemma \ref{derived-lemma-two-ss-complex-functor} 的第一个谱序列得到。
对一般情形，写成 $M^\bullet = \lim M_n^\bullet$，其中
$M_n^\bullet = \tau_{\geq -n}M^\bullet$。于是系统 $M_n^p$ 最终恒等于
$M^p$。我们断言
$$
(\widetilde{M}^\bullet)^a = R\lim (\widetilde{M}_n^\bullet)^a
$$
。也就是说，只须证明从左到右的自然映射在任意仿射对象
$U = \Spec(B)$（属于 $(\Sch/S)_{fppf}$）上诱导上同调同构。
对 $i \in \mathbf{Z}$ 和 $n > |i|$，有
$$
H^i(U, (\widetilde{M}_n^\bullet)^a) =
H^i(\tau_{\geq -n}M^\bullet \otimes_A B) =
H^i(M^\bullet \otimes_A B)
$$
。第一个等式由上面处理过的下有界情形成立。因此由 Cohomology on Sites,
Lemma \ref{sites-cohomology-lemma-RGamma-commutes-with-Rlim} 可知该断言成立。
最后得到
\begin{align*}
R\Gamma((\Sch/S)_{fppf}, (\widetilde{M}^\bullet)^a)
& =
R\Gamma((\Sch/S)_{fppf}, R\lim (\widetilde{M}_n^\bullet)^a) \\
& =
R\lim R\Gamma((\Sch/S)_{fppf}, (\widetilde{M}_n^\bullet)^a) \\
& =
R\lim M_n^\bullet \\
& =
M^\bullet
\end{align*}
，如所要求。第二个等式成立是因为 $R\lim$ 与 $R\Gamma$ 可交换，见
Cohomology on Sites，引理
\ref{sites-cohomology-lemma-RGamma-commutes-with-Rlim}.
\end{proof}










\section{ph 拓扑与 \'etale 拓扑的比较}
\label{etale-cohomology-section-ph-etale}

\noindent
本节的一个范本，是 Cohomology on Sites, Section
\ref{sites-cohomology-section-cohomology-LC} 中讨论的局部紧拓扑空间上通常拓扑与
qc 拓扑的比较。我们先回顾 Topologies, Sections
\ref{topologies-section-change-topologies} 和
\ref{topologies-section-etale} 中的一些材料。

\medskip\noindent
设 $S$ 是概形，并设 $(\Sch/S)_{ph}$ 是一个 ph 网站。
在同一个底层范畴上还有第二个拓扑，即 \'etale 拓扑，因而还有第二个网站
$(\Sch/S)_\etale$。恒等函子
$(\Sch/S)_\etale \to (\Sch/S)_{ph}$ 是连续的
（由 More on Morphisms, Lemma \ref{more-morphisms-lemma-fppf-ph} 和
Topologies，引理
\ref{topologies-lemma-zariski-etale-smooth-syntomic-fppf}），并定义网站态射
$$
\epsilon_S : (\Sch/S)_{ph} \longrightarrow (\Sch/S)_\etale
$$
。见 Cohomology on Sites, Section \ref{sites-cohomology-section-compare}。
请注意，$\epsilon_{S, *}$ 在底层预层上是恒等函子，而
$\epsilon_S^{-1}$ 把一个 \'etale 层变为它的 ph 层化。
令 $S_\etale$ 为小 \'etale 网站。存在网站态射
$$
\pi_S : (\Sch/S)_\etale \longrightarrow S_\etale
$$
，它由连续函子
$S_\etale \to (\Sch/S)_\etale$, $U \mapsto U$ 给出。
事实上，$S_\etale$ 有纤维积和终对象，上述函子与它们可交换，因而可以应用
Sites, Proposition \ref{sites-proposition-get-morphism}。

\begin{lemma}
\label{etale-cohomology-lemma-describe-pullback-pi-ph}
沿用上述记号。设 $\mathcal{F}$ 是 $S_\etale$ 上的层。规则
$$
(\Sch/S)_{ph} \longrightarrow \textit{Sets},\quad
(f : X \to S) \longmapsto \Gamma(X, f_{small}^{-1}\mathcal{F})
$$
给出一个层，因而当然也给出 $(\Sch/S)_\etale$ 上的层。
事实上，这个层等于 $\pi_S^{-1}\mathcal{F}$（在
$(\Sch/S)_\etale$ 上），而等于
$\epsilon_S^{-1}\pi_S^{-1}\mathcal{F}$（在 $(\Sch/S)_{ph}$ 上）。
\end{lemma}

\begin{proof}
关于 \'etale 拓扑的断言就是 Lemma \ref{etale-cohomology-lemma-describe-pullback} 的内容。
为完成证明，只须说明 $\pi_S^{-1}\mathcal{F}$ 对 ph 拓扑是层。
由 Topologies, Lemma \ref{topologies-lemma-characterize-sheaf}，
只须说明：给定固有满射态射 $V \to U$（其为 $S$ 上概形之间的态射），
有等化子图
$$
\xymatrix{
(\pi_S^{-1}\mathcal{F})(U) \ar[r] &
(\pi_S^{-1}\mathcal{F})(V) \ar@<1ex>[r] \ar@<-1ex>[r] &
(\pi_S^{-1}\mathcal{F})(V \times_U V)
}
$$
。置 $\mathcal{G} = \pi_S^{-1}\mathcal{F}|_{U_\etale}$。考虑交换图
$$
\xymatrix{
V \times_U V \ar[r] \ar[rd]_g \ar[d] & V \ar[d]^f \\
V \ar[r]^f & U
}
$$
。我们有
$$
(\pi_S^{-1}\mathcal{F})(V) = \Gamma(V, f^{-1}\mathcal{G}) =
\Gamma(U, f_*f^{-1}\mathcal{G})
$$
；这里用 $f_*$ 和 $f^{-1}$ 表示小 \'etale 网站之间的函子性。其次，有
$$
(\pi_S^{-1}\mathcal{F})(V \times_U V) =
\Gamma(V \times_U V, g^{-1}\mathcal{G}) =
\Gamma(U, g_*g^{-1}\mathcal{G})
$$
。等化子图中的两个映射来自两个映射
$$
f_*f^{-1}\mathcal{G} \longrightarrow g_*g^{-1}\mathcal{G}
$$
。因此只须证明 $\mathcal{G}$ 是这两个层映射的等化子。
设 $\overline{u}$ 是 $U$ 的几何点。置
$\Omega = \mathcal{G}_{\overline{u}}$。由 Lemma
\ref{etale-cohomology-lemma-proper-pushforward-stalk} 在 $\overline{u}$ 处取茎，得到两个映射
$$
H^0(V_{\overline{u}}, \underline{\Omega}) \longrightarrow
H^0((V \times_U V)_{\overline{u}}, \underline{\Omega}) =
H^0(V_{\overline{u}} \times_{\overline{u}} V_{\overline{u}},
\underline{\Omega})
$$
；其中 $\underline{\Omega}$ 表示取值为 $\Omega$ 的常值层。当然，
这些映射是沿投影映射的拉回。显然，由沿到第一因子的投影拉回得到的截面，
在第一投影的纤维上为常值；而由沿到第一因子的投影拉回得到的截面，
在第一投影的纤维上为常值。这些拉回映射之像的交中的截面，在整个
$V_{\overline{u}} \times_{\overline{u}} V_{\overline{u}}$ 上为常值；
也就是说，它们来自 $\Omega$ 中的元素，如所要求。
\end{proof}

\noindent
在 Lemma \ref{etale-cohomology-lemma-describe-pullback-pi-ph} 的情形中，
$\epsilon_S$ 与 $\pi_S$ 的复合以及上述等式确定网站态射
$$
a_S : (\Sch/S)_{ph} \longrightarrow S_\etale
$$

\begin{lemma}
\label{etale-cohomology-lemma-push-pull-ph-etale}
沿用上述记号。设 $f : X \to Y$ 是 $(\Sch/S)_{ph}$ 中的态射。
则有拓扑斯的交换图
$$
\xymatrix{
\Sh((\Sch/X)_{ph}) \ar[rr]_{f_{big, ph}} \ar[d]_{\epsilon_X} & &
\Sh((\Sch/Y)_{ph}) \ar[d]^{\epsilon_Y} \\
\Sh((\Sch/X)_\etale) \ar[rr]^{f_{big, \etale}} & &
\Sh((\Sch/Y)_\etale)
}
$$
以及
$$
\xymatrix{
\Sh((\Sch/X)_{ph}) \ar[rr]_{f_{big, ph}} \ar[d]_{a_X} & &
\Sh((\Sch/Y)_{ph}) \ar[d]^{a_Y} \\
\Sh(X_\etale) \ar[rr]^{f_{small}} & &
\Sh(Y_\etale)
}
$$
，其中 $a_X = \pi_X \circ \epsilon_X$ 且 $a_Y = \pi_X \circ \epsilon_X$。
\end{lemma}

\begin{proof}
这些图的交换性来自 Topologies, Section
\ref{topologies-section-change-topologies} 中的讨论。
\end{proof}

\begin{lemma}
\label{etale-cohomology-lemma-proper-push-pull-ph-etale}
在 Lemma \ref{etale-cohomology-lemma-push-pull-ph-etale} 中，如果 $f$ 是固有的，则有
$a_Y^{-1} \circ f_{small, *} = f_{big, ph, *} \circ a_X^{-1}$.
\end{lemma}

\begin{proof}
可以重复 Lemma \ref{etale-cohomology-lemma-compare-higher-direct-image-proper} (1) 的证明；
这里我们改为由该结论推出本结果。由于 $\epsilon_{Y, *}$ 在底层预层上是
恒等函子，它反映同构。Lemma \ref{etale-cohomology-lemma-describe-pullback-pi-ph}
中的描述说明 $\epsilon_{Y, *} \circ a_Y^{-1} = \pi_Y^{-1}$，
对 $X$ 也类似。为证明典范映射
$a_Y^{-1}f_{small, *}\mathcal{F} \to f_{big, ph, *}a_X^{-1}\mathcal{F}$
是同构，只须证明
\begin{align*}
\pi_Y^{-1}f_{small, *}\mathcal{F}
& =
\epsilon_{Y, *}a_Y^{-1}f_{small, *}\mathcal{F} \\
& \to 
\epsilon_{Y, *}f_{big, ph, *}a_X^{-1}\mathcal{F} \\
& =
f_{big, \etale, *} \epsilon_{X, *}a_X^{-1}\mathcal{F} \\
& =
f_{big, \etale, *}\pi_X^{-1}\mathcal{F}
\end{align*}
是同构。这就是 Lemma \ref{etale-cohomology-lemma-compare-higher-direct-image-proper} (1)。
\end{proof}

\begin{lemma}
\label{etale-cohomology-lemma-compare-ph-etale}
考虑比较态射
$\epsilon : (\Sch/S)_{ph} \to (\Sch/S)_\etale$。
令 $\mathcal{P}$ 表示概形的固有态射类。对
$X$（属于 $(\Sch/S)_\etale$），记
$\mathcal{A}'_X \subset \textit{Ab}((\Sch/X)_\etale)$
为由形如 $\pi_X^{-1}\mathcal{F}$ 的层构成的全子范畴，其中
$\mathcal{F}$ 是 $X_\etale$ 上的挠阿贝尔层。则 Cohomology on Sites 中的性质
(\ref{sites-cohomology-item-base-change-P}),
(\ref{sites-cohomology-item-restriction-A}),
(\ref{sites-cohomology-item-A-sheaf}),
(\ref{sites-cohomology-item-A-and-P})，以及
(\ref{sites-cohomology-item-refine-tau-by-P})
（见 Cohomology on Sites, Situation
\ref{sites-cohomology-situation-compare}）都成立。
\end{lemma}

\begin{proof}
我们先验证 Homology, Lemma
\ref{homology-lemma-characterize-weak-serre-subcategory} 的条件
(1)、(2)、(3)、(4)，从而说明
$\mathcal{A}'_X \subset \textit{Ab}((\Sch/X)_\etale)$ 是弱 Serre 子范畴。
(1)、(2)、(3) 是立即的，因为例如由 Lemma
\ref{etale-cohomology-lemma-cohomological-descent-etale}，$\pi_X^{-1}$ 正合且完全忠实。
如果
$0 \to \pi_X^{-1}\mathcal{F} \to \mathcal{G} \to \pi_X^{-1}\mathcal{F}' \to 0$
是 $\textit{Ab}((\Sch/X)_\etale)$ 中的短正合列，则由 Lemma
\ref{etale-cohomology-lemma-cohomological-descent-etale}，
$0 \to \mathcal{F} \to \pi_{X, *}\mathcal{G} \to \mathcal{F}' \to 0$
正合。特别地，$\pi_{X, *}\mathcal{G}$ 是 $X_\etale$ 上的挠阿贝尔层。
因此 $\mathcal{G} = \pi_X^{-1}\pi_{X, *}\mathcal{G}$ 属于
$\mathcal{A}'_X$，这就验证了最后一个条件。

\medskip\noindent
Cohomology on Sites, Property (\ref{sites-cohomology-item-base-change-P})
成立，因为概形有纤维积，而且概形的固有态射经基变换后仍是固有态射；
见 Morphisms, Lemma \ref{morphisms-lemma-base-change-proper}。

\medskip\noindent
Cohomology on Sites, Property (\ref{sites-cohomology-item-restriction-A})
来自 Topologies, Lemma
\ref{topologies-lemma-morphism-big-small-etale} 中的交换图 (3)。

\medskip\noindent
Cohomology on Sites, Property (\ref{sites-cohomology-item-A-sheaf})
就是 Lemma \ref{etale-cohomology-lemma-describe-pullback-pi-ph}。

\medskip\noindent
Cohomology on Sites, Property (\ref{sites-cohomology-item-A-and-P})
由 Lemma \ref{etale-cohomology-lemma-compare-higher-direct-image-proper} (2) 以及如下事实成立：
$R^if_{small}\mathcal{F}$ 是挠层，只要 $\mathcal{F}$ 是挠阿贝尔层
（定义在 $X_\etale$ 上）；见 Lemma
\ref{etale-cohomology-lemma-torsion-direct-image}。

\medskip\noindent
Cohomology on Sites, Property (\ref{sites-cohomology-item-refine-tau-by-P})
来自 More on Morphisms, Lemma
\ref{more-morphisms-lemma-dominate-fppf-etale-locally}，并结合以下事实：
有限态射是固有的，而满射固有态射定义 ph 覆盖；见
Topologies, Lemma \ref{topologies-lemma-surjective-proper-ph}。
\end{proof}

\begin{lemma}
\label{etale-cohomology-lemma-V-C-all-n-etale-ph}
沿用上述记号。
\begin{enumerate}
\item 对 $X \in \Ob((\Sch/S)_{ph})$ 和挠阿贝尔层 $\mathcal{F}$
（定义在 $X_\etale$ 上），有
$\epsilon_{X, *}a_X^{-1}\mathcal{F} = \pi_X^{-1}\mathcal{F}$
，并且 $R^i\epsilon_{X, *}(a_X^{-1}\mathcal{F}) = 0$ 对 $i > 0$ 成立。
\item 对固有态射 $f : X \to Y$（属于 $(\Sch/S)_{ph}$）和挠阿贝尔层
$\mathcal{F}$（定义在 $X$ 上），有
$a_Y^{-1}(R^if_{small, *}\mathcal{F}) =
R^if_{big, ph, *}(a_X^{-1}\mathcal{F})$
，对所有 $i$ 成立。
\item 对概形 $X$ 和 $K$（属于 $D^+(X_\etale)$，且上同调层为挠层），映射
$\pi_X^{-1}K \to R\epsilon_{X, *}(a_X^{-1}K)$ 是同构。
\item 对概形的固有态射 $f : X \to Y$ 和 $K$（属于
$D^+(X_\etale)$，且上同调层为挠层），有
$a_Y^{-1}(Rf_{small, *}K) = Rf_{big, ph, *}(a_X^{-1}K)$.
\end{enumerate}
\end{lemma}

\begin{proof}
由 Lemma \ref{etale-cohomology-lemma-compare-ph-etale}，Cohomology on Sites, Section
\ref{sites-cohomology-section-compare-general} 中的各引理都适用于当前情形。
为翻译这些结果，注意 Cohomology on Sites, Lemma
\ref{sites-cohomology-lemma-A} 中的范畴 $\mathcal{A}_X$，是
$\textit{Ab}((\Sch/X)_{ph})$ 中由形如 $a_X^{-1}\mathcal{F}$ 的层构成的
全子范畴，其中 $\mathcal{F}$ 是 $X_\etale$ 上的挠阿贝尔层。

\medskip\noindent
(1) 等价于 $(V_n)$ 对所有 $n$ 都成立；这由 Cohomology on Sites, Lemma
\ref{sites-cohomology-lemma-V-C-all-n-general} 成立。

\medskip\noindent
(2) 来自对 Cohomology on Sites, Lemma
\ref{sites-cohomology-lemma-V-implies-C-general} 的结论应用
$\epsilon_Y^{-1}$。

\medskip\noindent
(3) 来自 Cohomology on Sites, Lemma
\ref{sites-cohomology-lemma-V-C-all-n-general} (1)，因为
$\pi_X^{-1}K$ 属于 $D^+_{\mathcal{A}'_X}((\Sch/X)_\etale)$，并且
$a_X^{-1} = \epsilon_X^{-1} \circ a_X^{-1}$。

\medskip\noindent
出于同样的理由，(4) 来自 Cohomology on Sites, Lemma
\ref{sites-cohomology-lemma-V-C-all-n-general} (2)。
\end{proof}

\begin{lemma}
\label{etale-cohomology-lemma-cohomological-descent-etale-ph}
设 $X$ 是概形。对 $K \in D^+(X_\etale)$（其上同调层为挠层），映射
$$
K \longrightarrow Ra_{X, *}a_X^{-1}K
$$
是同构，其中 $a_X : \Sh((\Sch/X)_{ph}) \to \Sh(X_\etale)$ 如上所述。
\end{lemma}

\begin{proof}
我们先把断言归约到 $K$ 由单个阿贝尔层给出的情形。也就是说，把 $K$
表示为挠阿贝尔层的下有界复形 $\mathcal{F}^\bullet$。由 Cohomology on Sites, Lemma
\ref{sites-cohomology-lemma-torsion}，这是可能的。由单个层的情形可知
$\mathcal{F}^n = a_{X, *} a_X^{-1} \mathcal{F}^n$
，而且层 $R^qa_{X, *}a_X^{-1}\mathcal{F}^n$ 对 $q > 0$ 为零。
对 $a_X^{-1}\mathcal{F}^\bullet$ 和函子 $a_{X, *}$ 应用 Leray 无环性引理
（Derived Categories, Lemma \ref{derived-lemma-leray-acyclicity}），
便得到结论。以下假设 $K = \mathcal{F}$，其中 $\mathcal{F}$ 是挠阿贝尔层。

\medskip\noindent
由 Lemma \ref{etale-cohomology-lemma-describe-pullback-pi-ph}，有
$a_{X, *}a_X^{-1}\mathcal{F} = \mathcal{F}$。因此只须证明
$R^qa_{X, *}a_X^{-1}\mathcal{F} = 0$ 对 $q > 0$ 成立。
为此可以使用 $a_X = \epsilon_X \circ \pi_X$ 和 Leray 谱序列
（Cohomology on Sites, Lemma \ref{sites-cohomology-lemma-relative-Leray}）。
由 Lemma \ref{etale-cohomology-lemma-V-C-all-n-etale-ph}，有
$R^i\epsilon_{X, *}(a_X^{-1}\mathcal{F}) = 0$（当 $i > 0$），并且
$\epsilon_{X, *}a_X^{-1}\mathcal{F} = \pi_X^{-1}\mathcal{F}$。
由 Lemma \ref{etale-cohomology-lemma-cohomological-descent-etale}，有
$R^j\pi_{X, *}(\pi_X^{-1}\mathcal{F}) = 0$（当 $j > 0$）。证明完毕。
\end{proof}

\begin{lemma}
\label{etale-cohomology-lemma-compare-cohomology-etale-ph}
设 $X$ 是概形，并设
$a_X : \Sh((\Sch/X)_{ph}) \to \Sh(X_\etale)$ 如上。则：
\begin{enumerate}
\item $H^q(X_\etale, \mathcal{F}) = H^q_{ph}(X, a_X^{-1}\mathcal{F})$
，其中 $\mathcal{F}$ 是 $X_\etale$ 上的挠阿贝尔层；
\item $H^q(X_\etale, K) = H^q_{ph}(X, a_X^{-1}K)$
，其中 $K \in D^+(X_\etale)$ 且其上同调层为挠层。
\end{enumerate}
例如，如果 $A$ 是挠阿贝尔群，则
$H^q_\etale(X, \underline{A}) = H^q_{ph}(X, \underline{A})$.
\end{lemma}

\begin{proof}
由 Lemma \ref{etale-cohomology-lemma-cohomological-descent-etale-ph} 和
Cohomology on Sites, Remark \ref{sites-cohomology-remark-before-Leray}
即得此结论。
\end{proof}












\section{h 拓扑与 \'etale 拓扑的比较}
\label{etale-cohomology-section-h-etale}

\noindent
本节的一个范本，是 Cohomology on Sites, Section
\ref{sites-cohomology-section-cohomology-LC} 中讨论的局部紧拓扑空间上通常拓扑与
qc 拓扑的比较。此外，本节与比较 ph 拓扑和 \'etale 拓扑的那一节几乎逐字相同。
我们先回顾 Topologies, Sections
\ref{topologies-section-change-topologies}、
\ref{topologies-section-etale} 以及 More on Flatness, Section
\ref{flat-section-h} 中的一些材料。

\medskip\noindent
设 $S$ 是概形，并设 $(\Sch/S)_h$ 是一个 h 网站。
在同一个底层范畴上还有第二个拓扑，即 \'etale 拓扑，因而还有第二个网站
$(\Sch/S)_\etale$。恒等函子
$(\Sch/S)_\etale \to (\Sch/S)_h$ 是连续的
（由 More on Flatness, Lemma \ref{flat-lemma-zariski-h} 和
Topologies，引理
\ref{topologies-lemma-zariski-etale-smooth-syntomic-fppf}），并定义网站态射
$$
\epsilon_S : (\Sch/S)_h \longrightarrow (\Sch/S)_\etale
$$
。见 Cohomology on Sites, Section \ref{sites-cohomology-section-compare}。
请注意，$\epsilon_{S, *}$ 在底层预层上是恒等函子，而
$\epsilon_S^{-1}$ 把一个 \'etale 层变为它的 h 层化。
令 $S_\etale$ 为小 \'etale 网站。存在网站态射
$$
\pi_S : (\Sch/S)_\etale \longrightarrow S_\etale
$$
，它由连续函子
$S_\etale \to (\Sch/S)_\etale$, $U \mapsto U$ 给出。
事实上，$S_\etale$ 有纤维积和终对象，上述函子与它们可交换，因而可以应用
Sites, Proposition \ref{sites-proposition-get-morphism}。

\begin{lemma}
\label{etale-cohomology-lemma-describe-pullback-pi-h}
沿用上述记号。设 $\mathcal{F}$ 是 $S_\etale$ 上的层。规则
$$
(\Sch/S)_h \longrightarrow \textit{Sets},\quad
(f : X \to S) \longmapsto \Gamma(X, f_{small}^{-1}\mathcal{F})
$$
给出一个层，因而当然也给出 $(\Sch/S)_\etale$ 上的层。
事实上，这个层等于 $\pi_S^{-1}\mathcal{F}$（在
$(\Sch/S)_\etale$ 上），而等于
$\epsilon_S^{-1}\pi_S^{-1}\mathcal{F}$（在 $(\Sch/S)_h$ 上）。
\end{lemma}

\begin{proof}
关于 \'etale 拓扑的断言就是 Lemma \ref{etale-cohomology-lemma-describe-pullback} 的内容。
为完成证明，只须说明 $\pi_S^{-1}\mathcal{F}$ 对 h 拓扑是层。
然而，我们已在 Lemma \ref{etale-cohomology-lemma-describe-pullback-pi-ph} 中证明，
$\pi_S^{-1}\mathcal{F}$ 即使对更强的 ph 拓扑也是层。
\end{proof}

\noindent
在 Lemma \ref{etale-cohomology-lemma-describe-pullback-pi-h} 的情形中，
$\epsilon_S$ 与 $\pi_S$ 的复合以及上述等式确定网站态射
$$
a_S : (\Sch/S)_h \longrightarrow S_\etale
$$

\begin{lemma}
\label{etale-cohomology-lemma-push-pull-h-etale}
沿用上述记号。设 $f : X \to Y$ 是 $(\Sch/S)_h$ 中的态射。
则有拓扑斯的交换图
$$
\xymatrix{
\Sh((\Sch/X)_h) \ar[rr]_{f_{big, h}} \ar[d]_{\epsilon_X} & &
\Sh((\Sch/Y)_h) \ar[d]^{\epsilon_Y} \\
\Sh((\Sch/X)_\etale) \ar[rr]^{f_{big, \etale}} & &
\Sh((\Sch/Y)_\etale)
}
$$
以及
$$
\xymatrix{
\Sh((\Sch/X)_h) \ar[rr]_{f_{big, h}} \ar[d]_{a_X} & &
\Sh((\Sch/Y)_h) \ar[d]^{a_Y} \\
\Sh(X_\etale) \ar[rr]^{f_{small}} & &
\Sh(Y_\etale)
}
$$
，其中 $a_X = \pi_X \circ \epsilon_X$ 且 $a_Y = \pi_X \circ \epsilon_X$。
\end{lemma}

\begin{proof}
这些图的交换性与 Topologies, Section
\ref{topologies-section-change-topologies} 中所述类似。
\end{proof}

\begin{lemma}
\label{etale-cohomology-lemma-proper-push-pull-h-etale}
在 Lemma \ref{etale-cohomology-lemma-push-pull-h-etale} 中，如果 $f$ 是固有的，则有
$a_Y^{-1} \circ f_{small, *} = f_{big, h, *} \circ a_X^{-1}$.
\end{lemma}

\begin{proof}
可以重复 Lemma \ref{etale-cohomology-lemma-compare-higher-direct-image-proper} (1) 的证明；
这里我们改为由该结论推出本结果。由于 $\epsilon_{Y, *}$ 在底层预层上是
恒等函子，它反映同构。Lemma \ref{etale-cohomology-lemma-describe-pullback-pi-h}
中的描述说明 $\epsilon_{Y, *} \circ a_Y^{-1} = \pi_Y^{-1}$，
对 $X$ 也类似。为证明典范映射
$a_Y^{-1}f_{small, *}\mathcal{F} \to f_{big, h, *}a_X^{-1}\mathcal{F}$
是同构，只须证明
\begin{align*}
\pi_Y^{-1}f_{small, *}\mathcal{F}
& =
\epsilon_{Y, *}a_Y^{-1}f_{small, *}\mathcal{F} \\
& \to 
\epsilon_{Y, *}f_{big, h, *}a_X^{-1}\mathcal{F} \\
& =
f_{big, \etale, *} \epsilon_{X, *}a_X^{-1}\mathcal{F} \\
& =
f_{big, \etale, *}\pi_X^{-1}\mathcal{F}
\end{align*}
是同构。这就是 Lemma \ref{etale-cohomology-lemma-compare-higher-direct-image-proper} (1)。
\end{proof}

\begin{lemma}
\label{etale-cohomology-lemma-compare-h-etale}
考虑比较态射 $\epsilon : (\Sch/S)_h \to (\Sch/S)_\etale$。
令 $\mathcal{P}$ 表示固有态射类。对 $X$（属于 $(\Sch/S)_\etale$），记
$\mathcal{A}'_X \subset \textit{Ab}((\Sch/X)_\etale)$
为由形如 $\pi_X^{-1}\mathcal{F}$ 的层构成的全子范畴，其中
$\mathcal{F}$ 是 $X_\etale$ 上的挠阿贝尔层。则 Cohomology on Sites 中的性质
(\ref{sites-cohomology-item-base-change-P}),
(\ref{sites-cohomology-item-restriction-A}),
(\ref{sites-cohomology-item-A-sheaf}),
(\ref{sites-cohomology-item-A-and-P})，以及
(\ref{sites-cohomology-item-refine-tau-by-P})
（见 Cohomology on Sites, Situation
\ref{sites-cohomology-situation-compare}）都成立。
\end{lemma}

\begin{proof}
我们先验证 Homology, Lemma
\ref{homology-lemma-characterize-weak-serre-subcategory} 的条件
(1)、(2)、(3)、(4)，从而说明
$\mathcal{A}'_X \subset \textit{Ab}((\Sch/X)_\etale)$ 是弱 Serre 子范畴。
(1)、(2)、(3) 是立即的，因为例如由 Lemma
\ref{etale-cohomology-lemma-cohomological-descent-etale}，$\pi_X^{-1}$ 正合且完全忠实。
如果
$0 \to \pi_X^{-1}\mathcal{F} \to \mathcal{G} \to \pi_X^{-1}\mathcal{F}' \to 0$
是 $\textit{Ab}((\Sch/X)_\etale)$ 中的短正合列，则由 Lemma
\ref{etale-cohomology-lemma-cohomological-descent-etale}，
$0 \to \mathcal{F} \to \pi_{X, *}\mathcal{G} \to \mathcal{F}' \to 0$
正合。特别地，$\pi_{X, *}\mathcal{G}$ 是 $X_\etale$ 上的挠阿贝尔层。
因此 $\mathcal{G} = \pi_X^{-1}\pi_{X, *}\mathcal{G}$ 属于
$\mathcal{A}'_X$，这就验证了最后一个条件。

\medskip\noindent
Cohomology on Sites, Property (\ref{sites-cohomology-item-base-change-P})
成立，因为概形有纤维积，概形的固有态射经基变换后仍是固有态射（见
Morphisms, Lemma \ref{morphisms-lemma-base-change-proper}），而有限表现态射经
基变换后仍是有限表现态射（见 Morphisms, Lemma
\ref{morphisms-lemma-base-change-finite-presentation}）。

\medskip\noindent
Cohomology on Sites, Property (\ref{sites-cohomology-item-restriction-A})
来自 Topologies, Lemma
\ref{topologies-lemma-morphism-big-small-etale} 中的交换图 (3)。

\medskip\noindent
Cohomology on Sites, Property (\ref{sites-cohomology-item-A-sheaf})
就是 Lemma \ref{etale-cohomology-lemma-describe-pullback-pi-h}。

\medskip\noindent
Cohomology on Sites, Property (\ref{sites-cohomology-item-A-and-P})
由 Lemma \ref{etale-cohomology-lemma-compare-higher-direct-image-proper} (2) 以及如下事实成立：
$R^if_{small}\mathcal{F}$ 是挠层，只要 $\mathcal{F}$ 是挠阿贝尔层
（定义在 $X_\etale$ 上）；见 Lemma \ref{etale-cohomology-lemma-torsion-direct-image}。

\medskip\noindent
Cohomology on Sites, Property (\ref{sites-cohomology-item-refine-tau-by-P})
由 More on Morphisms, Lemma
\ref{more-morphisms-lemma-dominate-fppf-etale-locally} 推出，并结合以下事实：
满射有限局部自由态射是满射、固有且有限表现的，因而由 More on Flatness,
Lemma \ref{flat-lemma-surjective-proper-finite-presentation-h}，它定义 h 覆盖。
\end{proof}

\begin{lemma}
\label{etale-cohomology-lemma-V-C-all-n-etale-h}
沿用上述记号。
\begin{enumerate}
\item 对 $X \in \Ob((\Sch/S)_{h})$ 和挠阿贝尔层 $\mathcal{F}$
（定义在 $X_\etale$ 上），有
$\epsilon_{X, *}a_X^{-1}\mathcal{F} = \pi_X^{-1}\mathcal{F}$
，并且 $R^i\epsilon_{X, *}(a_X^{-1}\mathcal{F}) = 0$ 对 $i > 0$ 成立。
\item 对固有态射 $f : X \to Y$（属于 $(\Sch/S)_h$）和挠阿贝尔层
$\mathcal{F}$（定义在 $X$ 上），有
$a_Y^{-1}(R^if_{small, *}\mathcal{F}) =
R^if_{big, h, *}(a_X^{-1}\mathcal{F})$
，对所有 $i$ 成立。
\item 对概形 $X$ 和 $K$（属于 $D^+(X_\etale)$，且上同调层为挠层），映射
$\pi_X^{-1}K \to R\epsilon_{X, *}(a_X^{-1}K)$ 是同构。
\item 对概形的固有态射 $f : X \to Y$ 和 $K$（属于
$D^+(X_\etale)$，且上同调层为挠层），有
$a_Y^{-1}(Rf_{small, *}K) = Rf_{big, h, *}(a_X^{-1}K)$.
\end{enumerate}
\end{lemma}

\begin{proof}
由 Lemma \ref{etale-cohomology-lemma-compare-h-etale}，Cohomology on Sites, Section
\ref{sites-cohomology-section-compare-general} 中的各引理都适用于当前情形。
为翻译这些结果，注意 Cohomology on Sites, Lemma
\ref{sites-cohomology-lemma-A} 中的范畴 $\mathcal{A}_X$，是
$\textit{Ab}((\Sch/X)_h)$ 中由形如 $a_X^{-1}\mathcal{F}$ 的层构成的
全子范畴，其中 $\mathcal{F}$ 是 $X_\etale$ 上的挠阿贝尔层。

\medskip\noindent
(1) 等价于 $(V_n)$ 对所有 $n$ 都成立；这由 Cohomology on Sites, Lemma
\ref{sites-cohomology-lemma-V-C-all-n-general} 成立。

\medskip\noindent
(2) 来自对 Cohomology on Sites, Lemma
\ref{sites-cohomology-lemma-V-implies-C-general} 的结论应用
$\epsilon_Y^{-1}$。

\medskip\noindent
(3) 来自 Cohomology on Sites, Lemma
\ref{sites-cohomology-lemma-V-C-all-n-general} (1)，因为
$\pi_X^{-1}K$ 属于 $D^+_{\mathcal{A}'_X}((\Sch/X)_\etale)$，并且
$a_X^{-1} = \epsilon_X^{-1} \circ a_X^{-1}$。

\medskip\noindent
出于同样的理由，(4) 来自 Cohomology on Sites, Lemma
\ref{sites-cohomology-lemma-V-C-all-n-general} (2)。
\end{proof}

\begin{lemma}
\label{etale-cohomology-lemma-cohomological-descent-etale-h}
设 $X$ 是概形。对 $K \in D^+(X_\etale)$（其上同调层为挠层），映射
$$
K \longrightarrow Ra_{X, *}a_X^{-1}K
$$
是同构，其中 $a_X : \Sh((\Sch/X)_h) \to \Sh(X_\etale)$ 如上所述。
\end{lemma}

\begin{proof}
我们先把断言归约到 $K$ 由单个阿贝尔层给出的情形。也就是说，把 $K$
表示为挠阿贝尔层的下有界复形 $\mathcal{F}^\bullet$。由 Cohomology on Sites,
Lemma \ref{sites-cohomology-lemma-torsion}，这是可能的。由单个层的情形可知
$\mathcal{F}^n = a_{X, *} a_X^{-1} \mathcal{F}^n$
，而且层 $R^qa_{X, *}a_X^{-1}\mathcal{F}^n$ 对 $q > 0$ 为零。
对 $a_X^{-1}\mathcal{F}^\bullet$ 和函子 $a_{X, *}$ 应用 Leray 无环性引理
（Derived Categories, Lemma \ref{derived-lemma-leray-acyclicity}），
便得到结论。以下假设 $K = \mathcal{F}$，其中 $\mathcal{F}$ 是挠阿贝尔层。

\medskip\noindent
由 Lemma \ref{etale-cohomology-lemma-describe-pullback-pi-h}，有
$a_{X, *}a_X^{-1}\mathcal{F} = \mathcal{F}$。因此只须证明
$R^qa_{X, *}a_X^{-1}\mathcal{F} = 0$ 对 $q > 0$ 成立。
为此可以使用 $a_X = \epsilon_X \circ \pi_X$ 和 Leray 谱序列
（Cohomology on Sites, Lemma \ref{sites-cohomology-lemma-relative-Leray}）。
由 Lemma \ref{etale-cohomology-lemma-V-C-all-n-etale-h}，有
$R^i\epsilon_{X, *}(a_X^{-1}\mathcal{F}) = 0$（当 $i > 0$），并且
$\epsilon_{X, *}a_X^{-1}\mathcal{F} = \pi_X^{-1}\mathcal{F}$。
由 Lemma \ref{etale-cohomology-lemma-cohomological-descent-etale}，有
$R^j\pi_{X, *}(\pi_X^{-1}\mathcal{F}) = 0$（当 $j > 0$）。证明完毕。
\end{proof}

\begin{lemma}
\label{etale-cohomology-lemma-compare-cohomology-etale-h}
设 $X$ 是概形，并设
$a_X : \Sh((\Sch/X)_h) \to \Sh(X_\etale)$ 如上。则：
\begin{enumerate}
\item $H^q(X_\etale, \mathcal{F}) = H^q_h(X, a_X^{-1}\mathcal{F})$
，其中 $\mathcal{F}$ 是 $X_\etale$ 上的挠阿贝尔层；
\item $H^q(X_\etale, K) = H^q_h(X, a_X^{-1}K)$
，其中 $K \in D^+(X_\etale)$ 且其上同调层为挠层。
\end{enumerate}
例如，如果 $A$ 是挠阿贝尔群，则
$H^q_\etale(X, \underline{A}) = H^q_h(X, \underline{A})$.
\end{lemma}

\begin{proof}
由 Lemma \ref{etale-cohomology-lemma-cohomological-descent-etale-h} 和
Cohomology on Sites, Remark \ref{sites-cohomology-remark-before-Leray}
即得此结论。
\end{proof}











\section{\'etale 层的下降}
\label{etale-cohomology-section-glueing-etale}

\noindent
我们证明 \'etale 层在 fppf 拓扑和 h 拓扑中能够``黏合''，
并证明若干相关结果. 我们已经证明了下列相关结论
\begin{enumerate}
\item 引理 \ref{etale-cohomology-lemma-describe-pullback} 告诉我们，概形 $S$ 的小 \'etale
网站上的一个层确定了 $S$ 的大 \'etale 网站上的一个层，后者对 fpqc 覆盖满足
层条件（因而当然也对 Zariski、\'etale、光滑、合成以及 fppf 覆盖满足层条件），
\item 引理 \ref{etale-cohomology-lemma-describe-pullback-pi-fppf} 是上一点对 fppf 拓扑的重述，
\item 引理 \ref{etale-cohomology-lemma-describe-pullback-pi-ph} 对 ph 拓扑证明了同样的结论，
\item 引理 \ref{etale-cohomology-lemma-describe-pullback-pi-h} 对 h 拓扑证明了同样的结论，
\item 引理 \ref{etale-cohomology-lemma-descent-sheaf-fppf-etale} 是 \'etale 层的 fppf 下降的一个版本，以及
\item 注 \ref{etale-cohomology-remark-cohomological-descent-finite} 告诉我们，对有限满射态射有
\'etale 层的下降（我们将在下文澄清并加强这一点）.
\end{enumerate}
在关于单纯空间的章节中，我们还会证明这方面的一些附加结果；例如见
Simplicial Spaces，第 \ref{spaces-simplicial-section-fppf-hypercovering}
节和第 \ref{spaces-simplicial-section-proper-hypercovering-spaces} 节.

\medskip\noindent
为了方便地表述结果，我们需要一些记号. 设
$\mathcal{U} = \{f_i : X_i \to X\}$
是具有固定靶的概形态射族. 相对于 $\mathcal{U}$ 的 \'etale 层的一个
{\it 下降数据}是一个族
$((\mathcal{F}_i)_{i \in I}, (\varphi_{ij})_{i, j \in I})$，其中
\begin{enumerate}
\item $\mathcal{F}_i$ 属于 $\Sh(X_{i, \etale})$，并且
\item $\varphi_{ij} :
\text{pr}_{0, small}^{-1} \mathcal{F}_i
\longrightarrow
\text{pr}_{1, small}^{-1} \mathcal{F}_j$
是 $\Sh((X_i \times_X X_j)_\etale)$ 中的一个同构.
\end{enumerate}
并且满足{\it 余循环条件}：图表
$$
\xymatrix{
\text{pr}_{0, small}^{-1}\mathcal{F}_i
\ar[dr]_{\text{pr}_{02, small}^{-1}\varphi_{ik}}
\ar[rr]^{\text{pr}_{01, small}^{-1}\varphi_{ij}} & &
\text{pr}_{1, small}^{-1}\mathcal{F}_j
\ar[dl]^{\text{pr}_{12, small}^{-1}\varphi_{jk}} \\
& \text{pr}_{2, small}^{-1}\mathcal{F}_k
}
$$
在 $\Sh((X_i \times_X X_j \times_X X_k)_\etale)$ 中交换.
自然有{\it 下降数据的态射}这一概念，由此得到下降数据范畴.
称下降数据
$((\mathcal{F}_i)_{i \in I}, (\varphi_{ij})_{i, j \in I})$
为{\it 有效的}，如果存在一个 $\mathcal{F}$ 属于 $\Sh(X_\etale)$，并且存在同构
$\varphi_i : f_{i, small}^{-1} \mathcal{F} \to \mathcal{F}_i$
（位于 $\Sh(X_{i, \etale})$ 中），它们与 $\varphi_{ij}$ 相容，也就是说
$$
\varphi_{ij} =
\text{pr}_{1, small}^{-1} (\varphi_j) \circ
\text{pr}_{0, small}^{-1} (\varphi_i^{-1})
$$
也可以如下表述. 给定一个对象 $\mathcal{F}$ 属于 $\Sh(X_\etale)$，我们得到
{\it 典范下降数据}
$(f_{i, small}^{-1}\mathcal{F}_i, c_{ij})$，其中 $c_{ij}$ 是典范同构
$$
c_{ij} : \text{pr}_{0, small}^{-1} f_{i, small}^{-1}\mathcal{F}
\longrightarrow
\text{pr}_{1, small}^{-1} f_{j, small}^{-1}\mathcal{F}
$$
下降数据
$((\mathcal{F}_i)_{i \in I}, (\varphi_{ij})_{i, j \in I})$
有效，当且仅当它同构于某个 $\mathcal{F}$ 所对应的典范下降数据，且该对象属于
$\Sh(X_\etale)$.

\medskip\noindent
如果该族由单个态射 $\{X \to Y\}$ 组成，那么我们把下降数据视为一对
$(\mathcal{F}, \varphi)$，其中 $\mathcal{F}$ 是 $\Sh(X_\etale)$ 的一个对象，
而 $\varphi$ 是同构
$$
\text{pr}_{0, small}^{-1} \mathcal{F}
\longrightarrow
\text{pr}_{1, small}^{-1} \mathcal{F}
$$
，它位于 $\Sh((X \times_Y X)_\etale)$ 中并满足余循环条件：
$$
\xymatrix{
\text{pr}_{0, small}^{-1}\mathcal{F}
\ar[dr]_{\text{pr}_{02, small}^{-1}\varphi}
\ar[rr]^{\text{pr}_{01, small}^{-1}\varphi} & &
\text{pr}_{1, small}^{-1}\mathcal{F}
\ar[dl]^{\text{pr}_{12, small}^{-1}\varphi} \\
& \text{pr}_{2, small}^{-1}\mathcal{F}
}
$$
该图在 $\Sh((X \times_Y X \times_Y X)_\etale)$ 中交换.
与前面完全一样，也有下降数据的态射和有效性的概念.

\medskip\noindent
我们先证明满的整态射的有效下降.

\begin{lemma}
\label{etale-cohomology-lemma-glue-etale-sheaf-section}
设 $f : X \to Y$ 是有截面的概形态射. 则函子
$$
\Sh(Y_\etale)
\longrightarrow
\text{\'etale 层关于 }\{X \to Y\}\text{ 的下降数据}
$$
把 $\mathcal{G}$（它属于 $\Sh(Y_\etale)$）送到典范下降数据；它是范畴等价.
\end{lemma}

\begin{proof}
这是形式的，并且只依赖于拉回函子的函子性. 我们略去细节. 提示：如果
$s : Y \to X$ 是一个截面，那么一个拟逆是把 $(\mathcal{F}, \varphi)$
送到 $s_{small}^{-1}\mathcal{F}$ 的函子.
\end{proof}

\begin{lemma}
\label{etale-cohomology-lemma-glue-etale-sheaf-integral-surjective}
设 $f : X \to Y$ 是概形的满整态射. 函子
$$
\Sh(Y_\etale)
\longrightarrow
\text{\'etale 层关于 }\{X \to Y\}\text{ 的下降数据}
$$
是范畴等价.
\end{lemma}

\begin{proof}
在本证明中，我们省略拉回和正像函子的下标 ${}_{small}$.
记 $X_1 = X \times_Y X$，并记 $f_1 : X_1 \to Y$ 为态射
$f \circ \text{pr}_0 = f \circ \text{pr}_1$.
设 $(\mathcal{F}, \varphi)$ 是 $\{X \to Y\}$ 的一个下降数据.
令 $\mathcal{F}_1 = \text{pr}_0^{-1}\mathcal{F}$.
可把 $\varphi$ 视为定义了一个同构
$\mathcal{F}_1 \to \text{pr}_1^{-1}\mathcal{F}$.
我们断言，把下降数据
$(\mathcal{F}, \varphi)$
送到层
$$
\mathcal{G} =
\text{等化子}\left(
\xymatrix{
f_*\mathcal{F}
\ar@<1ex>[r] \ar@<-1ex>[r] &
f_{1, *}\mathcal{F}_1
}
\right)
$$
的规则是引理陈述中函子的拟逆. 两个箭头中的第一个来自映射
$$
f_*\mathcal{F} \to
f_*\text{pr}_{0, *}\text{pr}_0^{-1}\mathcal{F} =
f_{1, *}\mathcal{F}_1
$$
，而第二个箭头来自映射
$$
f_*\mathcal{F} \to
f_* \text{pr}_{1, *}\text{pr}_1^{-1}\mathcal{F}
\xleftarrow{\varphi}
f_* \text{pr}_{0, *} \text{pr}_0^{-1}\mathcal{F} =
f_{1, *}\mathcal{F}_1
$$
，其中向左的箭头可逆. 为了证明这确实可行，我们必须证明典范映射
$f^{-1}\mathcal{G} \to \mathcal{F}$ 是同构；细节略去. 为此，只需在任意一族
态射 $\Spec(k) \to Y$ 拉回之后核验，其中 $k$ 是代数闭域. 的确，相应的基变换
$X_k \to X$ 是联合满的，并且可以通过考察几何点处的茎来核验
$X_\etale$ 上层之间的一个映射是否为同构；见定理 \ref{etale-cohomology-theorem-exactness-stalks}.
由引理 \ref{etale-cohomology-lemma-integral-pushforward-commutes-with-base-change}，
$\mathcal{G}$ 从下降数据 $(\mathcal{F}, \varphi)$ 的上述构造与任意基变换交换.
因此可以假设 $Y$ 是一个代数闭域的谱（注意基变换保持态射 $f$ 的这些性质；见
Morphisms，引理 \ref{morphisms-lemma-base-change-surjective}
和 \ref{morphisms-lemma-base-change-finite}）.
在这种情形下，态射 $X \to Y$ 有一个截面，故由引理
\ref{etale-cohomology-lemma-glue-etale-sheaf-section} 可知该函子是等价.
另一方面，读者可以证明，该函子是等价当且仅当上述构造是拟逆；细节略去.
这就完成了证明.
\end{proof}

\begin{lemma}
\label{etale-cohomology-lemma-glue-etale-sheaf-proper-surjective}
设 $f : X \to Y$ 是概形的满真态射. 函子
$$
\Sh(Y_\etale)
\longrightarrow
\text{\'etale 层关于 }\{X \to Y\}\text{ 的下降数据}
$$
是范畴等价.
\end{lemma}

\begin{proof}
引理 \ref{etale-cohomology-lemma-glue-etale-sheaf-integral-surjective} 中给出的证明原封不动地适用，
但应当把对引理 \ref{etale-cohomology-lemma-integral-pushforward-commutes-with-base-change}
的引用替换为对引理 \ref{etale-cohomology-lemma-proper-base-change-f-star} 的引用.
\end{proof}

\begin{lemma}
\label{etale-cohomology-lemma-glue-etale-sheaf-check-after-base-change}
设 $f : X \to Y$ 是概形态射. 设 $Z \to Y$ 是概形的满整态射或满真态射.
如果函子
$$
\Sh(Z_\etale)
\longrightarrow
\text{\'etale 层关于 }\{X \times_Y Z \to Z\}\text{ 的下降数据}
$$
和
$$
\Sh((Z \times_Y Z)_\etale)
\longrightarrow
\text{\'etale 层关于 }
\{X \times_Y (Z \times_Y Z) \to Z \times_Y Z\}\text{ 的下降数据}
$$
是范畴等价，那么
$$
\Sh(Y_\etale)
\longrightarrow
\text{\'etale 层关于 }\{X \to Y\}\text{ 的下降数据}
$$
是等价.
\end{lemma}

\begin{proof}
这是定义以及引理 \ref{etale-cohomology-lemma-glue-etale-sheaf-integral-surjective}
和 \ref{etale-cohomology-lemma-glue-etale-sheaf-proper-surjective} 的形式推论. 细节略去.
\end{proof}

\begin{lemma}
\label{etale-cohomology-lemma-glue-etale-sheaf-fppf-cover}
设 $f : X \to Y$ 是满、平且局部有限表现的概形态射. 函子
$$
\Sh(Y_\etale)
\longrightarrow
\text{\'etale 层关于 }\{X \to Y\}\text{ 的下降数据}
$$
是范畴等价.
\end{lemma}

\begin{proof}
与引理 \ref{etale-cohomology-lemma-glue-etale-sheaf-integral-surjective} 的证明完全一样，
我们断言，一个拟逆由把 $(\mathcal{F}, \varphi)$ 送到
$$
\mathcal{G} =
\text{等化子}\left(
\xymatrix{
f_*\mathcal{F}
\ar@<1ex>[r] \ar@<-1ex>[r] &
f_{1, *}\mathcal{F}_1
}
\right)
$$
的函子给出；为了证明这一点，只需证明 $f^{-1}\mathcal{G} \to \mathcal{F}$
是同构. 这一点可局部核验，因此可以并且确实假设 $Y$ 是仿射的. 随后可找到一个
有限满态射 $Z \to Y$，使得存在开覆盖 $Z = \bigcup W_i$，其中
$W_i \to Y$ 经由 $X$ 分解. 见 More on Morphisms，引理
\ref{more-morphisms-lemma-dominate-fppf-finite}.
应用引理 \ref{etale-cohomology-lemma-glue-etale-sheaf-check-after-base-change}，可知只需在把
$Y$ 分别替换为 $Z$ 和 $Z \times_Y Z$、并把 $f$ 替换为其基变换之后证明本引理.
因此可以假设 $f$ 在 Zariski 局部有截面. 当然，利用该问题在 $Y$ 上是局部的，
我们便归结到有一个截面的情形，而这正是引理
\ref{etale-cohomology-lemma-glue-etale-sheaf-section}.
\end{proof}

\begin{lemma}
\label{etale-cohomology-lemma-glue-etale-sheaf-fppf}
设 $\{f_i : X_i \to X\}$ 是概形的一个 fppf 覆盖. 函子
$$
\Sh(X_\etale)
\longrightarrow
\text{\'etale 层关于 }\{f_i : X_i \to X\}\text{ 的下降数据}
$$
是范畴等价.
\end{lemma}

\begin{proof}
对态射 $f : \coprod X_i \to X$，有引理
\ref{etale-cohomology-lemma-glue-etale-sheaf-fppf-cover}. 随后一个形式论证表明，$f$ 的下降数据
与该覆盖的下降数据是一回事；比较 Descent，引理
\ref{descent-lemma-family-is-one}. 细节略去.
\end{proof}

\begin{lemma}
\label{etale-cohomology-lemma-glue-etale-sheaf-modification}
设 $f : X' \to X$ 是概形的真态射. 设 $i : Z \to X$ 是闭浸入.
令 $E = Z \times_X X'$. 图表为
$$
\xymatrix{
E \ar[d]_g \ar[r]_j & X' \ar[d]^f \\
Z \ar[r]^i & X
}
$$
如果 $f$ 在 $X \setminus Z$ 上是同构，那么函子
$$
\Sh(X_\etale)
\longrightarrow
\Sh(X'_\etale) \times_{\Sh(E_\etale)} \Sh(Z_\etale)
$$
是范畴等价.
\end{lemma}

\begin{proof}
我们将使用 Categories，例 \ref{categories-example-2-fibre-product-categories}
中构造的 $2$-纤维积范畴. 该函子把 $\mathcal{F}$ 送到三元组
$(f^{-1}\mathcal{F}, i^{-1}\mathcal{F}, c)$，其中
$c : j^{-1}f^{-1}\mathcal{F} \to g^{-1}i^{-1}\mathcal{F}$ 是典范同构.
我们将构造一个拟逆函子. 设 $(\mathcal{F}', \mathcal{G}, \alpha)$
是箭头右端的一个对象. 我们得到同构
$$
i^{-1}f_*\mathcal{F}' = g_*j^{-1}\mathcal{F}'
\xrightarrow{g_*\alpha}
g_*g^{-1}\mathcal{G}
$$
第一个等式是引理 \ref{etale-cohomology-lemma-proper-base-change-f-star}.
利用它，我们得到映射
$i_*\mathcal{G} \to i_*g_*g^{-1}\mathcal{G}$
和 $f'_*\mathcal{F}' \to i_*g_*g^{-1}\mathcal{G}$. 令
$$
\mathcal{F} = f_*\mathcal{F}' \times_{i_*g_*g^{-1}\mathcal{G}} i_*\mathcal{G}
$$
我们断言，$\mathcal{F}$ 是箭头左端的一个对象，并且它在右端的像同构于我们起初的
三元组. 计算 $\mathcal{F}$ 在几何点 $\overline{x}$ 处的茎，其中该点属于 $X$.

\medskip\noindent
如果 $\overline{x}$ 不在 $Z$ 中，那么一方面，$\overline{x}$ 来自唯一的几何点
$\overline{x}'$，其中该点属于 $X'$，并且
$\mathcal{F}'_{\overline{x}'} = (f_*\mathcal{F}')_{\overline{x}}$
；另一方面，$(i_*\mathcal{G})_{\overline{x}}$ 和
$(i_*g_*g^{-1}\mathcal{G})_{\overline{x}}$ 都是单点集.
所以可见 $\mathcal{F}_{\overline{x}}$ 等于 $\mathcal{F}'_{\overline{x}'}$.

\medskip\noindent
如果 $\overline{x}$ 在 $Z$ 中，也就是说，$\overline{x}$ 是一个几何点
$\overline{z}$ 的像，其中该点属于 $Z$，那么得到
$(i_*\mathcal{G})_{\overline{x}} = \mathcal{G}_{\overline{z}}$
以及
$$
(i_*g_*g^{-1}\mathcal{G})_{\overline{x}} =
(g_*g^{-1}\mathcal{G})_{\overline{z}} =
\Gamma(E_{\overline{z}}, g^{-1}\mathcal{G}|_{E_{\overline{z}}})
$$
（由上面使用的正像真基变换），并且类似地
$$
(f_*\mathcal{F}')_{\overline{x}} =
\Gamma(X'_{\overline{x}}, \mathcal{F}'|_{X'_{\overline{x}}})
$$
因为有等同 $E_{\overline{z}} = X'_{\overline{x}}$，并且 $\alpha$ 定义了层
$\mathcal{F}'|_{X'_{\overline{x}}}$ 与
$g^{-1}\mathcal{G}|_{E_{\overline{z}}}$ 之间的同构，所以得到
$$
\mathcal{F}_{\overline{x}} = \mathcal{G}_{\overline{z}}
$$
.

\medskip\noindent
为了完成证明，注意有典范映射
$i^{-1}\mathcal{F} \to \mathcal{G}$ 和 $f^{-1}\mathcal{F} \to \mathcal{F}'$，
它们与 $\alpha$ 相容，并且在茎上给出我们上面看到的同构.
我们略去这些映射的细致构造.
\end{proof}

\begin{lemma}
\label{etale-cohomology-lemma-h-descent-etale-sheaves}
设 $S$ 是概形. 则群胚化范畴
$$
p : \mathcal{S} \longrightarrow (\Sch/S)_h
$$
在 $U$ 上的纤维范畴是层范畴 $\Sh(U_\etale)$，即 $U$ 的小 \'etale
网站上的层所成的范畴；它是一个群胚叠.
\end{lemma}

\begin{proof}
为证明本引理，我们将核验 More on Flatness，引理
\ref{flat-lemma-refine-check-h-stack} 的条件 (1)、(2) 和 (3).

\medskip\noindent
条件 (1) 成立，因为层可以黏合（并且 Zariski 覆盖是 \'etale 覆盖）.
见 Sites，引理 \ref{sites-lemma-glue-sheaves}.

\medskip\noindent
为看出条件 (2)，设 $f : X \to Y$ 是 $S$ 上的满、平、真且有限表现的态射，
其中 $Y$ 是仿射的. 由引理 \ref{etale-cohomology-lemma-glue-etale-sheaf-fppf-cover}
或引理 \ref{etale-cohomology-lemma-glue-etale-sheaf-proper-surjective}，
我们有 $\{X \to Y\}$ 的下降.

\medskip\noindent
条件 (3) 立即从更一般的引理 \ref{etale-cohomology-lemma-glue-etale-sheaf-modification} 推出.
\end{proof}



















\section{爆破方形与 \'etale 上同调}
\label{etale-cohomology-section-blow-up-square}

\noindent
爆破方形在 More on Flatness，第 \ref{flat-section-blow-up-ph} 节中引入.
利用真基变换定理，可以看出爆破方形有一个 Mayer--Vietoris 型结论.

\begin{lemma}
\label{etale-cohomology-lemma-blow-up-square-cohomological-descent}
设 $X$ 是概形，并设 $Z \subset X$ 是由有限型拟凝聚理想截出的闭子概形.
考虑相应的爆破方形
$$
\xymatrix{
E \ar[d]_\pi \ar[r]_j & X' \ar[d]^b \\
Z \ar[r]^i & X
}
$$
对具有挠上同调层的 $K \in D^+(X_\etale)$，有可区别三角
$$
K \to Ri_*(K|_Z) \oplus Rb_*(K|_{X'}) \to Rc_*(K|_E) \to K[1]
$$
，它位于 $D(X_\etale)$ 中，其中 $c = i \circ \pi = b \circ j$.
\end{lemma}

\begin{proof}
记号 $K|_{X'}$ 表示 $b_{small}^{-1}K$.
选取阿贝尔层的一个下有界复形 $\mathcal{F}^\bullet$，它表示 $K$.
注意，$i_*(\mathcal{F}^\bullet|_Z)$ 表示 $Ri_*(K|_Z)$，因为 $i_*$ 正合
（命题 \ref{etale-cohomology-proposition-finite-higher-direct-image-zero}）.
选取拟同构
$b_{small}^{-1}\mathcal{F}^\bullet \to \mathcal{I}^\bullet$
，其中 $\mathcal{I}^\bullet$ 是 $X'_\etale$ 上内射阿贝尔层的下有界复形.
此映射伴随于一个映射 $\mathcal{F}^\bullet \to b_*(\mathcal{I}^\bullet)$，
并且 $b_*(\mathcal{I}^\bullet)$ 表示 $Rb_*(K|_{X'})$.
由引理 \ref{etale-cohomology-lemma-proper-base-change-f-star} 有
$\pi_*(\mathcal{I}^\bullet|_E) = (b_*\mathcal{I}^\bullet)|_Z$，而由引理
\ref{etale-cohomology-lemma-proper-base-change}，此复形表示 $R\pi_*(K|_E)$. 因此映射
$$
Ri_*(K|_Z) \oplus Rb_*(K|_{X'}) \to Rc_*(K|_E)
$$
由下有界复形之间的满射映射
$$
i_*(\mathcal{F}^\bullet|_Z) \oplus
b_*(\mathcal{I}^\bullet)
\to
i_*\left(b_*(\mathcal{I}^\bullet)|_Z\right)
$$
表示. 为得到所需的可区别三角，只需证明典范映射
$\mathcal{F}^\bullet \to i_*(\mathcal{F}^\bullet|_Z) \oplus
b_*(\mathcal{I}^\bullet)$
拟同构地映到上面所示复形映射的核（也就是说，复形的一个短正合列在导出范畴中
确定一个可区别三角；见 Derived Categories，第
\ref{derived-section-canonical-delta-functor} 节）.
可以在一个几何点 $\overline{x}$ 处的茎上核验这一点，其中该点属于 $X$.
如果 $\overline{x}$ 不在 $Z$ 中，那么 $X' \to X$ 在 $\overline{x}$ 的一个开邻域上
是同构. 因而，在此情形下，如果用 $\overline{x}'$ 表示 $X'$ 的相应几何点，
那么必须证明
$$
\mathcal{F}^\bullet_{\overline{x}} \to \mathcal{I}^\bullet_{\overline{x}'}
$$
是拟同构. 由我们对 $\mathcal{I}^\bullet$ 的选取，这一点成立.
如果 $\overline{x}$ 在 $Z$ 中，那么
$b_(\mathcal{I}^\bullet)_{\overline{x}} \to 
i_*\left(b_*(\mathcal{I}^\bullet)|_Z\right)_{\overline{x}}$
是阿贝尔群复形之间的同构. 因而该核等于
$i_*(\mathcal{F}^\bullet|_Z)_{\overline{x}} =
\mathcal{F}^\bullet_{\overline{x}}$，正如所需.
\end{proof}

\begin{lemma}
\label{etale-cohomology-lemma-blow-up-square-etale-cohomology}
设 $X$ 是概形，并设 $K \in D^+(X_\etale)$ 具有挠上同调层.
设 $Z \subset X$ 是由有限型拟凝聚理想截出的闭子概形. 考虑相应的爆破方形
$$
\xymatrix{
E \ar[d] \ar[r] & X' \ar[d]^b \\
Z \ar[r] & X
}
$$
则有典范长正合列
$$
H^p_\etale(X, K) \to
H^p_\etale(X', K|_{X'}) \oplus
H^p_\etale(Z, K|_Z) \to
H^p_\etale(E, K|_E) \to
H^{p + 1}_\etale(X, K)
$$
\end{lemma}

\begin{proof}[第一个证明]
这立即从引理 \ref{etale-cohomology-lemma-blow-up-square-cohomological-descent} 以及等式
$$
R\Gamma(X, Rb_*(K|_{X'})) = R\Gamma(X', K|_{X'})
$$
推出（见 Cohomology on Sites，第 \ref{sites-cohomology-section-leray} 节），
其他项也类似.
\end{proof}

\begin{proof}[第二个证明]
由引理 \ref{etale-cohomology-lemma-compare-cohomology-etale-ph}，这些上同调群是
$X, X', E, Z$ 关于网站 $(\Sch/X)_{ph}$ 上某个阿贝尔层复形的上同调.
（具体来说，是导出范畴的对象 $a_X^{-1}K$；见上面的引理
\ref{etale-cohomology-lemma-describe-pullback-pi-ph}，并回忆 $K|_{X'} = b_{small}^{-1}K$.）
由 More on Flatness，引理 \ref{flat-lemma-blow-up-square-ph}，可表预层图表的
ph 层化是余笛卡儿的. 因而，把非常一般的 Cohomology on Sites，引理
\ref{sites-cohomology-lemma-square-triangle} 应用于网站 $(\Sch/X)_{ph}$
和本引理的交换图表，即得结论.
\end{proof}

\begin{lemma}
\label{etale-cohomology-lemma-blow-up-square-equivalence}
设 $X$ 是概形，并设 $Z \subset X$ 是由有限型拟凝聚理想截出的闭子概形.
考虑相应的爆破方形
$$
\xymatrix{
E \ar[d]_\pi \ar[r]_j & X' \ar[d]^b \\
Z \ar[r]^i & X
}
$$
假设给定
\begin{enumerate}
\item 一个对象 $K'$，它属于 $D^+(X'_\etale)$ 并具有挠上同调层，
\item 一个对象 $L$，它属于 $D^+(Z_\etale)$ 并具有挠上同调层，以及
\item 一个同构 $\gamma : K'|_E \to L|_E$.
\end{enumerate}
那么存在一个对象 $K$ 属于 $D^+(X_\etale)$，并存在同构
$f : K|_{X'} \to K'$、$g : K|_Z \to L$，使得
$\gamma = g|_E \circ f^{-1}|_E$. 此外，给定
\begin{enumerate}
\item 一个对象 $M$，它属于 $D^+(X_\etale)$ 并具有挠上同调层，
\item 一个态射 $\alpha : K' \to M|_{X'}$，它属于 $D(X'_\etale)$，
\item 一个态射 $\beta : L \to M|_Z$，它属于 $D(Z_\etale)$，
\end{enumerate}
使得
$$
\alpha|_E  = \beta|_E \circ \gamma.
$$
那么存在一个态射 $M \to K$ 属于 $D(X_\etale)$，它在 $X'$ 上的限制是
$a \circ f$，而在 $Z$ 上的限制是 $b \circ g$.
\end{lemma}

\begin{proof}
如果 $K$ 存在，那么引理 \ref{etale-cohomology-lemma-blow-up-square-cohomological-descent}
给出了它嵌入其中的一个可区别三角. 因此我们只需选取一个可区别三角
$$
K \to Ri_*(L) \oplus Rb_*(K') \to Rc_*(L|_E) \to K[1]
$$
，其中 $c = i \circ \pi = b \circ j$. 这里，映射
$Ri_*(L) \to Rc_*(L|_E)$ 是把 $Ri_*$ 应用于伴随映射
$E \to R\pi_*(L|_E)$ 所得的映射. 映射 $Rb_*(K') \to Rc_*(L|_E)$
是典范映射 $Rb_*(K') \to Rc_*(K'|_E)) = R$ 与 $Rc_*(\gamma)$ 的复合.
引理陈述中的映射 $g$ 和 $f$ 是这些映射的伴随. 如果把这个可区别三角限制到
$Z$，那么由基变换定理（引理 \ref{etale-cohomology-lemma-proper-base-change}），映射
$Rb_*(K) \to Rc_*(L|_E)$ 变成同构，因而映射 $g : K|_Z \to L$ 是同构.
观察该可区别三角，可见 $f : K|_{X'} \to K'$ 在
$X' \setminus E = X \setminus Z$ 上是同构. 此外，由构造有
$\gamma \circ f|_E = g|_E$. 由于 $\gamma$ 和 $g$ 是同构，遂推出 $f$
在 $E$ 的几何点处的茎上也诱导同构. 因此 $f$ 是同构.

\medskip\noindent
对于最后一个陈述，可以把 $K'$ 替换为 $K|_{X'}$，把 $L$ 替换为 $K|_Z$，
并把 $\gamma$ 替换为典范等同. 注意，$\alpha$ 和 $\beta$ 诱导交换方形
$$
\xymatrix{
K \ar[r] \ar@{..>}[d] &
Ri_*(K|_Z) \oplus Rb_*(K|_{X'}) \ar[r] \ar[d]_{\beta \oplus \alpha} &
Rc_*(K|_E) \ar[r] \ar[d]_{\alpha|_E} &
K[1] \ar@{..>}[d] \\
M \ar[r] &
Ri_*(M|_Z) \oplus Rb_*(M|_{X'}) \ar[r] &
Rc_*(M|_E) \ar[r] &
M[1]
}
$$
因此，由导出范畴的公理，得到一条虚线箭头，从而产生可区别三角之间的一个态射.
\end{proof}














\section{近似爆破方形与 h 拓扑}
\label{etale-cohomology-section-blow-up-h}

\noindent
本节继续 More on Flatness，第 \ref{flat-section-blow-up-h} 节中的讨论.
为方便读者，回忆一个近似爆破方形是交换图表
\begin{equation}
\label{etale-cohomology-equation-almost-blow-up-square}
\vcenter{
\xymatrix{
E \ar[d] \ar[r] & X' \ar[d]^b \\
Z \ar[r] & X
}
}
\end{equation}
；它由概形组成，并满足下列条件：
\begin{enumerate}
\item $Z \to X$ 是有限表现闭浸入，
\item $E = b^{-1}(Z)$ 是 $X'$ 的局部主闭子概形，
\item $b$ 是真且有限表现的，
\item 闭子概形 $X'' \subset X'$ 由 $\mathcal{O}_{X'}$ 的支集位于 $E$ 上的截面
所成的拟凝聚理想截出（Properties，引理
\ref{properties-lemma-sections-supported-on-closed-subset}）；它是 $X$ 沿 $Z$ 的爆破.
\end{enumerate}
由此可知，态射 $b$ 诱导同构 $X' \setminus E \to X \setminus Z$.

\medskip\noindent
我们将给出一个判据，判断大 fppf 网站 $(\Sch/S)_{fppf}$ 的导出范畴中的对象
是否具有``h 层性''，其中 $S$ 是概形. 在同一个底范畴上还有
第二个拓扑，即 h 拓扑（More on Flatness，第 \ref{flat-section-h} 节）.
回忆 fppf 覆盖是 h 覆盖（More on Flatness，引理
\ref{flat-lemma-zariski-h}）. 因而可以考虑态射
$$
\epsilon : (\Sch/S)_h \longrightarrow (\Sch/S)_{fppf}
$$
见 Cohomology on Sites，第 \ref{sites-cohomology-section-compare} 节.
特别地，有完全忠实函子
$$
R\epsilon_* : D((\Sch/S)_h) \longrightarrow D((\Sch/S)_{fppf})
$$
我们可以问：这个函子的本质像是什么？

\begin{lemma}
\label{etale-cohomology-lemma-blow-up-square-h}
沿用上述记号，如果 $K$ 属于 $R\epsilon_*$ 的本质像，那么 Cohomology on Sites，
引理 \ref{sites-cohomology-lemma-c-square} 中的映射 $c^K_{X, Z, X', E}$ 是拟同构.
\end{lemma}

\begin{proof}
用 ${}^\#$ 表示 h 拓扑中的层化. 我们已经在 More on Flatness，引理
\ref{flat-lemma-blow-up-square-h} 中看到
$h_X^\# = h_Z^\# \amalg_{h_E^\#} h_{X'}^\#$. 另一方面，映射
$h_E^\# \to h_{X'}^\#$ 是单射，因为 $E \to X'$ 是单态射.
所以本引理是 Cohomology on Sites，引理
\ref{sites-cohomology-lemma-descent-squares-helper} 的特殊情形（该引理本身是
Cohomology on Sites，引理 \ref{sites-cohomology-lemma-square-triangle} 的形式推论）.
\end{proof}

\begin{proposition}
\label{etale-cohomology-proposition-check-h}
设 $K$ 是 $D^+((\Sch/S)_{fppf})$ 的一个对象. 则 $K$ 属于
$R\epsilon_* : D((\Sch/S)_h) \to D((\Sch/S)_{fppf})$
的本质像，当且仅当 $c^K_{X, X', Z, E}$ 对每个近似爆破方形
(\ref{etale-cohomology-equation-almost-blow-up-square}) 都是拟同构；这里该方形位于
$(\Sch/S)_h$ 中，并且 $X$ 是仿射的.
\end{proposition}

\begin{proof}
应用 Cohomology on Sites，引理
\ref{sites-cohomology-lemma-descent-squares} 即可证明；其假设由引理
\ref{etale-cohomology-lemma-blow-up-square-h} 和 More on Flatness，命题
\ref{flat-proposition-check-h} 成立.
\end{proof}

\begin{lemma}
\label{etale-cohomology-lemma-refine-check-h}
设 $K$ 是 $D^+((\Sch/S)_{fppf})$ 的一个对象. 则 $K$ 属于
$R\epsilon_* : D((\Sch/S)_h) \to D((\Sch/S)_{fppf})$ 的本质像，当且仅当
对 More on Flatness，例 \ref{flat-example-one-generator} 和
\ref{flat-example-two-generators} 中的每个近似爆破方形，
$c^K_{X, X', Z, E}$ 都是拟同构.
\end{lemma}

\begin{proof}
应用 Cohomology on Sites，引理
\ref{sites-cohomology-lemma-descent-squares} 即可证明；其假设由引理
\ref{etale-cohomology-lemma-blow-up-square-h} 和 More on Flatness，引理
\ref{flat-lemma-refine-check-h} 成立.
\end{proof}






\section{h 拓扑中结构层的上同调}
\label{etale-cohomology-section-cohomology-O-h}

\noindent
设 $p$ 是素数. 设 $(\mathcal{C}, \mathcal{O})$ 是满足
$p\mathcal{O} = 0$ 的环化网站. 我们把 $\colim_F \mathcal{O}$ 定义为下列系统在
环层范畴中的余极限
$$
\mathcal{O} \xrightarrow{F} \mathcal{O} \xrightarrow{F}
\mathcal{O} \xrightarrow{F} \ldots
$$
，其中 $F : \mathcal{O} \to \mathcal{O}$，$f \mapsto f^p$ 是 Frobenius 端同态.

\begin{lemma}
\label{etale-cohomology-lemma-h-sheaf-colim-F}
设 $p$ 是素数. 设 $S$ 是 $\mathbf{F}_p$ 上的概形. 考虑层
$\mathcal{O}^{perf} = \colim_F \mathcal{O}$，它位于 $(\Sch/S)_{fppf}$ 上.
那么 $\mathcal{O}^{perf}$ 属于
$R\epsilon_* : D((\Sch/S)_h) \to D((\Sch/S)_{fppf})$ 的本质像.
\end{lemma}

\begin{proof}
使用引理 \ref{etale-cohomology-lemma-refine-check-h} 的判据来证明. 在核验这些条件之前，注意对
一个拟紧拟分离对象 $X$ 属于 $(\Sch/S)_{fppf}$，有
$$
H^i_{fppf}(X, \mathcal{O}^{perf}) = \colim_F H^i_{fppf}(X, \mathcal{O})
$$
见 Cohomology on Sites，引理
\ref{sites-cohomology-lemma-colim-works-over-collection}.
还将使用 $H^i_{fppf}(X, \mathcal{O}) = H^i(X, \mathcal{O})$；见 Descent，
命题 \ref{descent-proposition-same-cohomology-quasi-coherent}.

\medskip\noindent
令 $A, f, J$ 如 More on Flatness，例 \ref{flat-example-one-generator}
中所述，并考虑相伴的近似爆破方形. 因为 $X$, $X'$, $Z$, $E$ 都是仿射的，
$\mathcal{O}$ 没有高阶上同调. 因而只需核验
$$
0 \to
\mathcal{O}^{perf}(X) \to
\mathcal{O}^{perf}(X') \oplus \mathcal{O}^{perf}(Z) \to
\mathcal{O}^{perf}(E) \to 0
$$
是短正合列. 这已经在 More on Flatness，引理
\ref{flat-lemma-h-sheaf-colim-F}（的证明）中证明.

\medskip\noindent
令 $X, X', Z, E$ 如 More on Flatness，例
\ref{flat-example-two-generators} 中所述. 因为 $X$ 和 $Z$ 是仿射的，所以有
$H^p(X, \mathcal{O}_X) = H^p(Z, \mathcal{O}_X) = 0$，这里 $p > 0$.
由 More on Flatness，引理 \ref{flat-lemma-funny-blow-up}，有
$H^p(X', \mathcal{O}_{X'}) = 0$，这里 $p > 0$.
因为 $E = \mathbf{P}^1_Z$ 并且 $Z$ 是仿射的，也有
$H^p(E, \mathcal{O}_E) = 0$，这里 $p > 0$. 与上一段一样，归结为核验
$$
0 \to
\mathcal{O}^{perf}(X) \to
\mathcal{O}^{perf}(X') \oplus \mathcal{O}^{perf}(Z) \to
\mathcal{O}^{perf}(E) \to 0
$$
是短正合列. 这已经在 More on Flatness，引理
\ref{flat-lemma-h-sheaf-colim-F}（的证明）中证明.
\end{proof}

\begin{proposition}
\label{etale-cohomology-proposition-h-cohomology-structure-sheaf}
设 $p$ 是素数. 设 $S$ 是 $\mathbf{F}_p$ 上拟紧拟分离的概形. 则
$$
H^i((\Sch/S)_h, \mathcal{O}^h) =
\colim_F H^i(S, \mathcal{O})
$$
这里左端的 $\mathcal{O}^h$ 是指结构层的 h 层化.
\end{proposition}

\begin{proof}
这只是引理 \ref{etale-cohomology-lemma-h-sheaf-colim-F} 的重述. 回忆
$\mathcal{O}^h = \mathcal{O}^{perf} = \colim_F \mathcal{O}$, see
；见 More on Flatness，引理 \ref{flat-lemma-char-p}.
由引理 \ref{etale-cohomology-lemma-h-sheaf-colim-F} 可知，$\mathcal{O}^{perf}$ 作为
$D((\Sch/S)_{fppf})$ 的一个对象，具有形式 $R\epsilon_*K$，其中某个
$K \in D((\Sch/S)_h)$. 随后 $K = \epsilon^{-1}\mathcal{O}^{perf}$，
而它实际上等于 $\mathcal{O}^{perf}$，因为 $\mathcal{O}^{perf}$ 是 h 层.
见 Cohomology on Sites，第 \ref{sites-cohomology-section-compare} 节.
因而 $R\epsilon_*\mathcal{O}^{perf} = \mathcal{O}^{perf}$
（请原谅这个令人困惑的记号）. 于是本命题由 Leray
$$
R\Gamma_h(S, \mathcal{O}^{perf}) =
R\Gamma_{fppf}(S, R\epsilon_*\mathcal{O}^{perf}) =
R\Gamma_{fppf}(S, \mathcal{O}^{perf})
$$
以及等式
$$
H^i_{fppf}(S, \mathcal{O}^{perf}) =
H^i_{fppf}(S, \colim_F \mathcal{O}) =
\colim_F H^i_{fppf}(S, \mathcal{O})
$$
推出，因为 $S$ 拟紧且拟分离（见引理 \ref{etale-cohomology-lemma-h-sheaf-colim-F} 的证明）.
\end{proof}
